\documentclass[french]{book}

\usepackage[utf8]{inputenc}
\usepackage[T1]{fontenc}
\usepackage{lmodern}
\usepackage{geometry}
\usepackage{graphicx}
\usepackage{babel}
\usepackage{amsmath}
\usepackage{amsthm}
\usepackage{amssymb}
\usepackage{hyperref}
\usepackage{stmaryrd}
\usepackage{enumitem}
\usepackage{tcolorbox}
\usepackage{dsfont}
\usepackage{multirow}

\setlist[itemize]{label=$\bullet$}

\renewcommand{\P}{\mathbb{P}}
\newcommand{\N}{\mathbb{N}}
\newcommand{\Z}{\mathbb{Z}}
\newcommand{\Q}{\mathbb{Q}}
\newcommand{\R}{\mathbb{R}}
\newcommand{\C}{\mathbb{C}}
\newcommand{\K}{\mathbb{K}}
\newcommand{\U}{\mathbb{U}}
\newcommand{\E}{\mathbb{E}}
\newcommand{\V}{\mathbb{V}}
\newcommand{\st}{^\star}
\newcommand{\tq}{\middle|}
\newcommand{\inv}{^{-1}}

\renewcommand{\thesection}{\arabic{section}}
\renewcommand{\thesubsection}{\arabic{section}.\arabic{subsection}}
\renewcommand{\thesubsubsection}{\arabic{section}.\arabic{subsection}.\arabic{subsubsection}}

\theoremstyle{plain}
\newtheorem{thm}{Théorème}[chapter]
\def\thmautorefname{théorème}
\newtheorem{lem}{Lemme}[chapter]
\def\lemautorefname{lemme}
\newtheorem{prop}{Proposition}[chapter]
\def\propautorefname{propriété}
\newtheorem{cor}{Corollaire}[chapter]
\def\propautorefname{corollaire}

\theoremstyle{definition}
\newtheorem{dfn}{Définition}[chapter]
\def\dfnautorefname{définition}
\newtheorem{exa}{Exemple}[chapter]
\newtheorem{rmq}{Remarque}[chapter]
\newtheorem{met}{Méthode}[chapter]
\def\dfnautorefname{méthode}
\newtheorem{ax}{Axiome}[chapter]
\def\dfnautorefname{axiome}
\newtheorem{exe}{Exercice}[chapter]

\theoremstyle{remark}



\DeclareMathOperator{\card}{Card}
\DeclareMathOperator{\cov}{cov}
\DeclareMathOperator{\vect}{Vect}
\DeclareMathOperator{\id}{id}
\DeclareMathOperator{\ind}{\perp \!\!\! \perp}
\DeclareMathOperator{\tr}{Tr}
\DeclareMathOperator{\im}{Im}
\DeclareMathOperator{\com}{Com}
\DeclareMathOperator{\mat}{mat}

\title{Cours de mathématiques en MPSI}
\author{Eliott Perrière-Larraux}

\begin{document}
\frontmatter
\maketitle
\chapter{Introduction}
Hey ! Bienvenue à toi cher lecteur : si tu lis ceci, tu es probablement un élève de Ginette ou un petit chanceux qui a reçu ce document de mes propres mains. Je vais expliquer ici en quelques mots (voire un peu plus que "quelques") les motivations et spécificités de ce document, ainsi que les consignes d'utilisation. Si cela ne t'intéresse pas, passe directement à la suite, mais sache que vu le temps que j'ai passé à faire ce travail, j'apprécierais que tu écoutes ce que j'ai à dire. \\
\par Tout d'abord, j'ai eu l'idée de commencer ce document durant l'été entre ma sup et ma spé, face au volume colossal qu'occupe la totalité de mes cours de mathématiques. Comme je n'habite pas en région parisienne, cela m'arrangerait de pouvoir éviter de me trimbaler tous mes cours de maths de sup pendant la spé, c'est pourquoi j'ai compilé ici l'entièreté de mes cours de mathématiques de sup. Oui, tu as bien lu, l'entièreté. Le moindre résultat écrit au tableau durant les cours du matin avec M. Morlot est recensé ici. \\
\par Ce document représente au moins 200 heures de travail, réparties sur 5 semaines. J'ai tapé chaque chapitre à la main, en LaTeX (je crois que je suis bilingue LaTeX maintenant). Chaque chapitre contient toutes les définitions, tous les théorèmes (et variantes), ainsi que tous les exemples et toutes les remarques que j'ai jugés utiles et pertinents (la plupart en fait). J'y ai aussi mentionné les exercices donnés par M. Morlot en cours, avec des solutions ou des pistes de résolution pour la plupart. Les résultats hors-programme sont repérés par le sigle HP : ils sont alors toujours démontrés car tout résultat hors-programme doit être redémontré à l'écrit, et doit pouvoir l'être à la demande de l'examinateur à l'oral. Les parties entières du cours qui sont hors-programme sont aussi signalées par le sigle HP dans leur titre. Les autres démonstrations des résultats de cours ont été retirées car elles sont longues, et je souhaitais finir avant la fin des vacances d'été (ainsi que préserver le peu de santé mentale qu'il me reste). \\ 
\par Je vous demande par avance d'excuser toute faute de français (qui risque fort d'être une faute de frappe) ou toute faute relative aux mathématiques : je n'ai pas relu l'ensemble du document, j'ai d'autres bouquins de 400 pages à lire pour la rentrée. Si vous voyez une faute quelconque (mathématique ou linguistique), merci de m'envoyer un message avec la référence de la faute (chapitre, proposition par exemple) afin que je puisse la corriger. \\
\par Pour réaliser ce document, j'ai utilisé les cours que j'ai pris en note durant les cours dispensés par M. Morlot ainsi que les différents polycopiés qu'il nous a distribués (merci pour les PAPL). J'ai aussi utilisé les polycopiés de cours qu'il m'a donnés en raison de ma maladie, ce qui revient un petit peu au même que les cours pris en notes. Qu'il soit remercié ici pour tout cela. De plus, le chapitre "Compléments sur les anneaux commutatifs" est issu d'un document éponyme distribué par M. Lafitte aux Taupins, et qui a un jour atterri sur mon bureau par miracle (c'est vrai). Un grand merci à lui aussi pour ce document très intéressant. Je sais que M. Morlot n'apprécie pas que ses documents tournent (en tout cas c'est le cas pour ses polycopiés de cours), donc je vous demande de ne pas faire tourner ce document en-dehors de Ginette, et si vous n'êtes pas de Ginette, de ne le partager à personne. Je sais que ce document n'est pas à proprement parler son cours, mais il a tellement fait pour moi que je lui dois bien ça. \\
\par \textbf{Update} : le poly de Lafitte a été tapé, j'y ai rajouté quelques résultats sur les anneaux noethériens, les radicaux d'idéaux, les idéaux maximaux et les idéaux premiers. Absolument tous les PAPL et tous les cours ont été tapés. J'ai aussi rajouté un paragraphe de compléments dans le chapitre "Polynômes" qui expose brièvement les notions de contenu et de polynôme primitif, qui aboutit sur une démonstration élégante du critère d'Eisenstein. A terme, je ne pense pas rajouter quoi que ce soit sur ce document, il contient déjà bien plus que le simple programme de sup, et est déjà assez long comme ça. \\
\par Si tu utilises, comme je compte le faire personnellement, ce document pour réviser et pour ne pas avoir à traîner tout ton cours de sup, profite bien de ce document, je serai ravi s'il peut t'être utile. \\
\par Avant de vous laisser, je voudrais saluer et remercier tous mes amis de Ginette qui ont rendu cette première année magique, malgré les aspérités. En particulier, un grand merci à mes cos Théa, Célia, Alexis et Nicolas, ainsi qu'à Angélique et Alexandre que j'adore et qui m'ont permis de décompresser un peu entre deux séquences d'écriture de ce document, je vous embrasse ! Et oui j'ai ptêtre un peu pété les plombs en écrivant toute cette introduction ! \\
\par Bonne fin de vacances, ou bonne année de spé à tous. Des bisous.
\begin{flushright}
	EPL
\end{flushright}
\tableofcontents

\mainmatter
\chapter{Logique et raisonnements}
La logique, bien que son développement formel soit récent dans l'histoire des mathématiques (un siècle environ), est aujourd'hui perçue comme la base sur laquelle se fondent celles-ci. Notre objectif, modeste, est de présenter ici les principaux outils dont nous aurons besoin pour cette année. Ils nous aideront à produire une rédaction claire et débarrassée de toute ambiguïté, selon les canons actuels de la rigueur.
\section{De l'importance de bien rédiger}
Avant de passer à des considérations plus théoriques, indiquons comment bien rédiger une démonstration mathématique. On n'hésitera pas à relire les exemples qui suivent de nombreuses fois : c'est par l'imitation qui s'acquiert la technique ! Pour ceux qui ne connaîtraient pas l'usage du mot "ssi" ou des symboles $\forall$,$\exists$,$\implies$,$\iff$, tout est expliqué dans le deuxième paragraphe. \\ \par Voici quelques grands principes qu'il ne faut jamais perdre de vue lorsque vous rédigez. Pourquoi ? Tout d'abord, parce que la rigueur du raisonnement est à ce prix ! Comme toute langue étrangère, les mathématiques requièrent une pratique régulière pour être correctement écrites. Et au-delà de cela, la rigueur ainsi acquise vous sera utile dans bien d'autres domaines. En second lieu, n'oubliez pas que le jour des concours, le correcteur aura à lire beaucoup de copies en un temps très court. il est donc essentiel que vous sachiez dire les choses clairement et simplement. Dans les épreuves de mathématiques, vous rencontrerez plusieurs types de questions, que nous listons maintenant.
\subsection{Questions du type "pour tout"}
Premier type : une question commençant par "montrer que pour tout réel $x$, etc.". Dans ce cas, il importe de commencer la rédaction par \textbf{soit $x\in\R$} ou \textbf{soit $x$ un réel}. On dit qu'on a \textbf{présenté} la variable $x$.
\begin{exa}[Montrer que $\forall x\in\R,\ x^2\geqslant$]
	Soit $x\in\R$. Si $x\geqslant0$, alors on a bien $x\times x\geqslant0$. Et si $x\leqslant0$, on a aussi $x\times x\geqslant0$. Donc dans tous les cas, le résultat est prouvé.
\end{exa}
\begin{exa}[Montrer que $\forall\theta\in\R,\ e^{i\theta}\ne0$]
	Soit $\theta\in\R$. Alors $e^{i\theta}\times e^{-i\theta}=e^0=1$. En particulier, on a nécessairement $e^{i\theta}\ne0$, sans quoi on aurait $e^{i\theta}\times e^{-i\theta}=1=0$.
\end{exa}
De manière générale, \textbf{il est interdit de calculer avec une variable avant de l'avoir présentée}. Par exemple, si on demande de trouver tous les réels dont le carré vaut $1$, on ne peut pas écrire directement "supposons que $x^2=1$". En effet, dans ce cas, le correcteur serait en droit de demander "qui est $x$ ?". Il faut d'abord écrire "soit $x$ un réel".
\begin{rmq}[Portée du quantificateur "pour tout"]
	Attention à la chose suivante : si une phrase commence par "pour tout $x$" ou "$\forall x$", \textbf{la portée de la variable $x$ s'arrête à la fin de la phrase}. Autrement dit, une rédaction du type $$\text{Pour tout }x\in\R,\ x^2\geqslant0.\text{ Donc } x^2+1\geqslant0.$$ est incorrect en toute rigueur. a la place, il faudrait écrire $$\text{Pour tout }x\in\R,\ x^2\geqslant0.\text{ Donc pour tout }x\in\R,\ x^2+1\geqslant0.$$ Comme c'est un peu pénible, si on prévoit d'utiliser une même variable $x$ sur plusieurs phrases, il vaut mieux la présenter définitivement avec "soit $x$". Dans ce cas, la portée s'arrête à la fin de la question (ou au prochain "soit $x$", qui écrase en quelque sorte le précédent). Ainsi, on aurait aussi pu écrire $$\text{Soit }x\in\R. \text{ On a }x^2\geqslant0.\text{ Donc }x^2+1\geqslant0.$$
\end{rmq}
\subsection{Questions d'existence}
Supposons qu'on nous demande de montrer une \textbf{existence}. Dans ce cas, il est demandé d'\textbf{exhiber un candidat}, puis de \textbf{démontrer qu'il convient}. Souvent, il vaut mieux commencer par chercher le candidat au brouillon.
\subsection{Questions d'unicité}
Supposons qu'on nous demande de montrer une \textbf{unicité}. Dans ce cas, on se donne deux candidats, et et on montre qu'ils sont égaux.
\subsection{Équations, inéquations}
Voici la méthode de résolution.
\begin{itemize}
	\item Si ce n'est pas immédiate, on calcule le domaine de définition. Ne pas oublier de présenter l'inconnue ! Par exemple, supposons qu'elle s'appelle $x$. Si on nous demande de résoudre dans $\R$, on écrit : "Soit $x\in\R$. L'équation est définie en $x$ ssi, etc.". Appelons $\mathcal{D}$ le domaine obtenu.
	\item Pour la résolution à proprement parler, on écrit d'abord "Soit $x\in\mathcal{D}$", puis on cherche l'ensemble exact des valeurs de $x$ qui conviennent. Il faut être certain que les $x$ donnés sont solution, mais aussi qu'on n'a oublié aucune solution ! Autrement dit, il s'agit de prouver une \textbf{équivalence} logique : les $x$ donnés sont tous solution, et \textbf{réciproquement}, toute solution est une des valeurs données.
	\begin{itemize}
		\item[--] Soit on tente une résolution par équivalences successives en partant de l'équation (ou de l'inéquation). Dans ce cas, on écrit : "$x$ est solution ssi ...".
		\item[--] Soit on traite séparément chacune des deux implications : c'est très courant car il est rare de pouvoir mener une chaîne d'équivalences à son terme.
		\begin{itemize}
			\item[*] D'abord, on écrit "Supposons que $x$ est solution", et on en déduit un ensemble de valeurs potentielles de $x$ : c'est la condition \textbf{nécessaire}.
			\item[*] Puis, réciproquement, on n'oublie pas de vérifier si chacune de ces valeurs potentielles est solution ou non : c'est la condition \textbf{suffisante}. On obtient alors les solutions définitives.
		\end{itemize}
	\end{itemize}
\end{itemize}
\section{Propositions logiques}
\begin{dfn}[Proposition]
	Une \textbf{proposition} est une phrase qui est soit vraie soit fausse.
\end{dfn}
\begin{exa}
	Voici des propositions (vraies ou fausses) :
	\begin{itemize}
		\item Paris est la capitale de la France.
		\item Lyon est la capitale de la France.
		\item Les Sigmas sont les meilleurs.
	\end{itemize}
\end{exa}
\begin{dfn}[Équivalence logique]
	Deux propositions $p$ et $q$ sont dites \textbf{logiquement équivalentes} lorsqu'elles ont même table de vérité. On note alors $p\equiv q$.
\end{dfn}
\begin{dfn}[Conjonction, disjonction]
	Soit $p$ et $q$ deux propositions. A partir de $p$ et $q$, on définit les propositions "$p$ et $q$" et "$p$ ou $q$", appelées respectivement \textbf{conjonction} et \textbf{disjonction}, et notées respectivement $p\land q$ et $p\lor q$, par la table de vérité suivante : $$\begin{tabular}{|c|c|c|c|}\hline $p$&$q$&$p\land q$&$p\lor q$\\ \hline V&V&V&V \\ V&F&F&V \\ F&V&F&V \\ F&F&F&F \\ \hline\end{tabular}$$
\end{dfn}
\begin{dfn}[Négation]
	La \textbf{négation} d'une proposition $p$, notée "$\lnot p$" ou "non $p$" est définie par la table de vérité ci-contre : $$\begin{tabular}{|c|c|}\hline $p$&$\lnot p$\\ \hline V&F \\ F&V \\ \hline\end{tabular}$$ On a donc $\lnot(\lnot p)\equiv p$.
	\end{dfn}
\begin{dfn}[Quantificateurs]
	Le symbole \textbf{"pour tout"}, dit \textbf{quantificateur universel}, s'écrit $\forall$. Le symbole \textbf{"il existe"}, dit \textbf{quantificateur existentiel}, s'écrit $\exists$. Une \textbf{phrase quantifiée} s'écrit de la façon suivante : 
	\begin{itemize}
		\item D'abord l'un de ces deux symboles, appelés \textbf{quantificateurs} :
		\item Puis une relation d'appartenance d'un élément à un ensemble ;
		\item Une virgule (bien noter qu'après le symbole $\exists$, la virgule se lit "tel que") ;
		\item Une autre phrase quantifiée.
	\end{itemize}
	Ainsi, une phrase quantifiée peut se voir comme la succession de quantificateurs séparés par des virgules, et d'une affirmation finale.
\end{dfn}
\begin{rmq}
	Il existe un troisième quantificateur, qui est $\exists!$ ("il existe un unique"), mais il est moins utilisé. De toute façon, on peut le reconstituer à partir des deux premiers. Par exemple, "$\exists!x\in X,\ p(x)$" pourra se réécrire : $$\left[\exists x\in X,\ p(x)\right]\land\left[\forall(x,y)\in X^2,\ p(x)\land p(y)\implies x=y\right]$$
\end{rmq}
\begin{thm}[Négations de "et" et "ou"]
	Soit $p$ et $q$ deux propositions. On a $\lnot(p\land q)\equiv(\lnot p)\lor(\lnot q)$ et $\lnot(p\lor q)\equiv(\lnot p)\land(\lnot q)$.
\end{thm}
\begin{proof}
	Il suffit de dresser les tables de vérité.
\end{proof}
\begin{prop}[Associativité]
	Soit $p$, $q$ et $r$ trois propositions. On a :
	\begin{itemize}
		\item Associativité du "ou" : $(p\lor q)\lor r\equiv p\lor(q\lor r)$
		\item Associativité du "et" : $(p\land q)\land r\equiv p\land(q\land r)$
	\end{itemize}
\end{prop}
\begin{prop}[Commutativité]
	Soit $p$ et $q$ deux propositions. On a :
	\begin{itemize}
		\item Commutativité du "ou" : $p\lor q\equiv q\lor p$
		\item Commutativité du "et" : $p\land q\equiv q\land p$
	\end{itemize}
\end{prop}
\begin{prop}[Distributivité]
	Soit $p$, $q$ et $r$ trois propositions. On a :
	\begin{itemize}
		\item Distributivité du "ou" sur le "et" : $p\lor(q\land r)\equiv (p\lor q)\land (p\lor r)$
		\item Distributivité du "et" sur le "ou" : $p\land(q\lor r)\equiv (p\land q)\lor(p\land r)$
	\end{itemize}
\end{prop}
\begin{proof}
	Pour les trois propositions précédentes, il suffit de dresser les tables de vérité.
\end{proof}
\begin{rmq}[Portée du quantificateur $\exists$]
	La portée du quantificateur $\exists$ est aussi longue que celle de $\forall$. Si on veut réutiliser une variable énoncée après un quantificateur $\exists$, il faut écrire \textbf{"Fixons un tel ..."}. On peut alors réutiliser l'objet fixé, muni de ses propriétés supposées dans la phrase quantifiées.
\end{rmq}
\begin{thm}
	Pour écrire la négation d'une phrase quantifiée : 
	\begin{itemize}
		\item On inverse les quantificateurs : les $\forall$ deviennent des $\exists$, et les $\exists$ deviennent des $\forall$ ;
		\item Puis on nie l'affirmation finale.
	\end{itemize}
\end{thm}
\begin{proof}
	On le montre facilement par récurrence sur le nombre de quantificateurs.
\end{proof}
\begin{dfn}[Implication]
	Soit $p$ et $q$ deux propositions. L'\textbf{implication} $p\implies q$ est définie par la table de vérité exposée ci-dessous :
	$$\begin{tabular}{|c|c|c|}\hline $p$&$q$&$p\implies q$ \\ \hline V&V&V \\ V&F&F \\ F&V&V \\F&F&V \\ \hline\end{tabular}$$ Dans un énoncé, cela correspond à un raisonnement du type "\textbf{si $p$, alors $q$}".
\end{dfn}
\begin{prop}
	Soit $p$ et $q$ deux propositions. On a $(p\implies q)\equiv (\lnot p)\lor q$.
\end{prop}
\begin{proof}
	Il suffit de comparer les tables de vérité.
\end{proof}
\begin{rmq}
	Attention au fait que lorsque $p$ est fausse, $p\implies q$ est vraie. On a coutume de dire que \textbf{le faux implique n'importe quoi}. Cela dit, ce problème est relativement marginal. Il peut cependant survenir dans des initialisations de récurrences.
\end{rmq}
\begin{rmq}["Donc" $\ne \implies$]
	On ne peut pas utiliser $\implies$ comme une abréviation de "donc", ce sont deux choses différentes ! En effet, un "$p$ donc $q$" correspond à $\left[p\land \left(p \implies q\right)\right] $, qui a même valeur de vérité que $p\land q$. Un "donc" part de quelque chose de vrai pour arriver à quelque chose de vrai, alors que $p\implies q$ ne suppose pas la véracité de $p$. \par En somme, \textbf{le symbole $\implies$ est exclu dans une démonstration rédigée. Son usage n'est autorisé que dans une phrase quantifiée ou dans la résolution formelle d'une équation}.
\end{rmq}
\begin{cor}
	Soit $p$ et $q$ deux propositions logiques. On a : $$\lnot(p\implies q)\equiv p\land (\lnot q)$$
\end{cor}
\begin{proof}
	Passer à la négation la propriété précédente.
\end{proof}
\begin{dfn}[Contraposée]
	Soit $p$ et $q$ deux propositions logiques. L'implication \textbf{contraposée} de $p\implies q$ est l'implication $(\lnot q)\implies (\lnot p)$. 
\end{dfn}
\begin{prop}[Équivalence logique entre une implication et sa contraposée]
	Soit $p$ et $q$ deux propositions logiques. On a : $$p\implies q\equiv (\lnot q)\implies (\lnot p)$$
\end{prop}
\begin{proof}
	On a $(\lnot q)\implies (\lnot p)\equiv (\lnot(\lnot q))\lor (\lnot p)\equiv q\lor (\lnot p)\equiv (\lnot p)\lor q\equiv p\implies q$.
\end{proof}
\begin{thm}
	Soit $p$ et $q$ deux propositions logiques. Pour prouver $p\implies q$, il y a trois manières.
	\begin{itemize}
		\item La \textbf{manière directe} : on part de $p$ pour arriver à $q$. Dans ce cas, la bonne rédaction est "Supposons $p$, etc." puis plus loin "Donc $q$".
		\item La \textbf{contraposée} : on part de $\lnot q$ pour arriver à $\lnot p$.
		\item Le \textbf{raisonnement par l'absurde} : on suppose que $p$ est vraie, puis on suppose que $q$ est fausse pour arriver à une contradiction. Dans ce cas, après avoir supposé que $p$ est vraie, la bonne rédaction est : "Supposons que $q$ est fausse. Alors, comme on a $p$, cela impliquerait que, etc." puis plus loin "Contradiction. Donc $q$ est vraie."
	\end{itemize}
\end{thm}
\begin{proof}
	Pour le premier point, il est inutile de considérer le cas où $p$ est fausse, puisque le faux implique n'importe quoi. Il suffit donc de prouver que si $p$ est vraie alors $q$ est vraie. Pour la contraposée, on utilise la propriété précédente. Pour le raisonnement par l'absurde, on montre que $p\land (\lnot q)$ est fausse. C'est donc que sa négation, $(\lnot p)\lor q$ est vraie, soit encore que $p\implies q$ est vraie.
\end{proof}
\begin{dfn}[Équivalence]
	Soit $p$ et $q$ deux propositions logiques. Une \textbf{équivalence} entre $p$ et $q$, notée $p\iff q$, est définie par la table ci-contre : 
	$$\begin{tabular}{|c|c|c|}\hline $p$&$q$&$p\iff q$ \\ \hline V&V&V \\ V&F&F \\ F&V&F \\ F&F&V \\ \hline\end{tabular}$$ Cela signifie que $p$ et $q$ ont même valeur de vérité. Dans un énoncé, cela se traduira par une phrase du type "$p$ si, et seulement si, $q$" ou en abrégé "$p$ ssi $q$". On peut remarquer qu'on a : $$p\iff q\equiv (p\implies q)\land (q\implies p)$$
\end{dfn}
\begin{thm}
	Soit $p$ et $q$ deux propositions. Pour prouver $p\iff q$, il y a deux manières.
	\begin{itemize}
		\item Si on se sent suffisamment sûr de soi, la \textbf{manière directe} : on part de $p$, et par équivalences successives, on essaye d'arriver à $q$. C'est une méthode qui ne marche pas à tous les coups.
		\item Sinon, et c'est de loin le cas le plus fréquent, on démontre successivement $p\implies q$ et $q\implies p$. Dans la rédaction, au moment-clé où on "inverse le sens de la marche" et on passe à $q\implies p$, il est très important que le correcteur lise le mot "\textbf{réciproquement}". A la rigueur, on pourra le remplacer par un "d'une part/d'autre part".
	\end{itemize}
\end{thm}
\begin{rmq}
	Pour montrer qu'une équivalence est fausse, il suffit de prouver qu'une des deux implications est fausse.
\end{rmq}
\begin{prop}
	Soit $p$ et $q$ deux propositions logiques. On a : $$p\iff q\equiv(\lnot p)\iff(\lnot q)$$
\end{prop}
\begin{proof}
	Il suffit de passer chacune des deux implications à la contraposée.
\end{proof}
\begin{dfn}[Résolution d'une équation/inéquation]
	\textbf{Résoudre} une équation dans un ensemble $E$ consiste à trouver successivement : \begin{itemize}
		\item l'ensemble $D$ des $x\in E$ en lesquels cette équation soit définie ; 
		\item puis l'ensemble des $x\in D$ qui vérifient l'équation.
	\end{itemize}
	De même avec une inéquation. Autrement dit, à l'issue de la résolution, on doit être capable d'écrire $$\forall x\in E,\ x\text{ solution ssi }x\in S$$ où $S$ est un sous-ensemble de $D$ (lui-même sous-ensemble de $E$). $S$ est appelé \textbf{l'ensemble des solutions}.
\end{dfn}
\begin{dfn}[Raisonnement par analyse-synthèse]
	A chaque fois qu'un énoncé commencera par "trouver tous les $x\in E$ tels que ...", il est recommandé de travailler par \textbf{analyse-synthèse}. Concrètement, cela consiste à réfléchir en deux temps.
	\begin{itemize}
		\item \textbf{Analyse} : si $x\in E$ convient, alors \textbf{nécessairement}, etc. Jusqu'à ce qu'on arrive à un ensemble de valeurs possibles pour $x$, qu'on appelle l'ensemble des "candidats". On travaille donc par \textbf{implication successives}. Ne pas chercher à écrire des équivalences, cela est \textbf{inutile et dangereux}.
		\item \textbf{Synthèse} : \textbf{réciproquement}, on examine un à un tous les candidats trouvés, on garde ceux qui sont effectivement solution du problème.
	\end{itemize}
	Autrement dit, l'analyse consiste à restreindre l'ensemble possible des solutions, de sorte qu'il ne reste plus qu'un nombre limité de cas à examiner durant la synthèse.
\end{dfn}
\begin{proof}
	Formalisons un peu les choses.
	\begin{itemize}
		\item Pendant l'analyse, on détermine un ensemble $F\subset E$ tel que toutes les solutions éventuelles soient dans $F$. On a donc démontré la chose suivante : $$\forall x\in E,\ x\text{ solution}\implies x\in F$$
		\item Puis, pendant la synthèse, on détermine un ensemble $S\subset F$ tel que $$\forall x\in F,\ x\text{ solution}\iff x\in S$$
	\end{itemize}
	On vérifie alors que la conjonction des deux énoncés précédents a même valeur de vérité que $$\forall x\in E,\ x\text{ solution}\iff x\in S$$ Autrement dit, $S$ est bien l'ensemble des solutions.
\end{proof}
\begin{rmq}
	Dans les cas extrêmes (mais pas si rares !), il ne reste déjà plus qu'un seul candidat à la fin de l'analyse. Dans ce cas, cette première phase montre l'\textbf{unicité} de la solution. Ensuite la synthèse permet de montrer
	\begin{itemize}
		\item soit l'\textbf{existence} d'une solution (si le candidat trouvé répond au problème) ;
		\item soit qu'il n'y a aucune solution (si le candidat trouvé ne répond pas au problème).
	\end{itemize}
\end{rmq}
\begin{rmq}
	Il arrive que la phase d'analyse produise des conditions nécessaires si restrictives que réciproquement, toutes les valeurs trouvées soient des solutions. Dans ce cas, la synthèse est très rapide, elle se borne à une simple vérification. De manière équivalente, on peut direz qu'on a procédé par \textbf{condition nécessaire / condition suffisante} :  la phase d'analyse nous a amené à déterminer une condition nécessaire, qui s'est également révélée suffisante pendant la phase de synthèse.
\end{rmq}
\begin{dfn}[SPG]
	Enfin, citons la technique du "sans perte de généralité" ou "SPG". il s'agit, dans une démonstration dont les notations sont un peu lourdes, de se débarrasser de cette complexité en se restreignant à un sous-cas qui couvre en fait tous les cas par analogie. Cette technique fonctionne bien à chaque fois que le problème possède une certaine symétrie. Attention, la notation "SPG" est à éviter sur une copie, elle n'est pas standard.
\end{dfn}
\section{Raisonnement par récurrence}
Cette partie est d'une importance capitale pour acquérir de bons réflexes à l'écrit. La validité du raisonnement par récurrence sera établie au cours de la construction de $\N$ : ici, il s'agit seulement d'apprendre les différents types de récurrences et de savoir les rédiger proprement.
\subsection{Récurrence simple}
\begin{thm}[Principe de récurrence simple]
	Soit $(H_n)_{n\in\N}$ une suite d'assertions définies pour tout $n\in\N$. On suppose que les deux points suivants sont vrais : \begin{itemize}
		\item Initialisation : $H_0$ est vraie
		\item Hérédité : $\forall n\in\N,\ H_n\implies H_{n+1}$
	\end{itemize}
	Alors pour tout $n\in\N$, $H_n$ est vraie.
\end{thm}
\begin{rmq}
	Souvent, la phase d'initialisation est courte et facile. Pour l'hérédité, il peut être utile de se faire la main sur de petites valeurs de $n$ afin d'appréhender le fonctionnement général.
\end{rmq}
\begin{thm}[Rédaction d'une récurrence]
	Pour bien rédiger un raisonnement par récurrence : 
	\begin{itemize}
		\item Si ce n'est pas clair, on précise "soit pour tout $n\in\N$, la proposition $H_n$ : "..." ".
		\item On écrit "Initialisation :" et on prouve $H_0$.
		\item Ensuite on écrit "Hérédité : Soit $n\in\N$ tel que $H_n$ soit vraie." puis on prouve $H_{n+1}$. On peut aussi écrire "Soit $n>0$ et tel que $H_{n-1}$ puis on prouve $H_n$.
		\item A l'endroit de la preuve où on utilise le fait que $H_n$ est vraie, on écrit "par hypothèse de récurrence". A la fin de la récurrence, on écrit "ce qui achève la récurrence".
	\end{itemize}
\end{thm}
\subsection{Quelques variantes}
\begin{prop}[Récurrence à partir d'un certain rang]
	Soit $n_0\in\Z$ et soit $(P_n)_{n\geqslant n_0}$ une suite d'assertions. On suppose que les deux points suivants sont vrais : \begin{itemize}
		\item Initialisation : $P_{n_0}$ est vraie
		\item Hérédité : $\forall n\geqslant n_0,\ P_n\implies P_{n+1}$
	\end{itemize}
	Alors pour tout $n\geqslant n_0$, $P_n$ est vraie.
\end{prop}
\begin{proof}
	Il suffit de poser pour tout $n\in$ $Q_n=P_{n+n_0}$ et d'appliquer le principe de récurrence simple à $(Q_n)_{n\in\N}$. Ensuite, pour tout $n\geqslant n_0$, $n-n_0\in\N$ donc $P_n=Q_{n-n_0}$ est vraie.
\end{proof}
\begin{rmq}
	Cette technique peut se révéler utile lorsqu'on a une suite de propriétés à prouver pour $n\geqslant0$ mais que l'hérédité ne fonctionne qu'à partir d'un certain indice $n_0>0$. On traite alors les premières valeurs de $n$ (de $0$ à $n_0-1$) séparément, puis on initialise la récurrence à $n_0$.
\end{rmq}
\begin{prop}[Récurrence à deux pas]
	Soit $n_0\in\Z$ et soit $(P_n)_{n\geqslant n_0}$ une suite d'assertions. On suppose que les deux points suivants sont vrais : \begin{itemize}
		\item Initialisation : $P_{n_0}$ et $P_{n_0+1}$ sont vraies
		\item Hérédité : $\forall n\geqslant n_0,\ (P_n\land P_{n+1})\implies P_{n+2}$
	\end{itemize}
	Alors pour tout $n\geqslant n_0$, $P_n$ est vraie.
\end{prop}
\begin{proof}
	Il suffit d'appliquer le principe de récurrence à partir d'un certain rang à la suite d'assertions définies par $Q_n=(P_n\land P_{n+1})$. Comme $Q_n$ implique $P_n$, on obtient le résultat.
\end{proof}
\begin{rmq}
	Sur le même principe, on peut effectuer des récurrences à trois pas, à quatre pas... L'important est de bien le préciser sur sa copie.
\end{rmq}
\begin{prop}[Récurrence forte]
	Soit $n_0\in\Z$ et soit $(P_n)_{n\geqslant n_0}$ une suite d'assertions. On suppose que les deux points suivants sont vrais : \begin{itemize}
		\item Initialisation : $P_{n_0}$ est vraie
		\item Hérédité : $\forall n\geqslant n_0,\ (\forall k\in\llbracket n_0,n\rrbracket,\ P_k\text{ est vraie})\implies P_{n+1}$
	\end{itemize}
	Alors pour tout $n\geqslant n_0$, $P_n$ est vraie. Ainsi, une récurrence forte consiste dans sa phase d'hérédité à supposer que $tous$ les rangs $\leqslant n$ sont vrais (et non pas seulement $n$) pour en déduire $P_{n+1}$.
\end{prop}
\begin{proof}
	Cette fois, il suffit d'appliquer le principe de récurrence à partir d'un certain rang à la suite d'assertions définies par $Q_n=(\forall k\in\llbracket n_0,n\rrbracket,\ P_k)$. Comme $Q_n$ implique $P_n$, on obtient le résultat.
\end{proof}
\begin{prop}[Récurrence finie]
	Soit $(P_n)_{n_0\leqslant n\leqslant n_1}$ une suite d'assertion avec $n_0\leqslant n_1$. On suppose que les deux points suivants sont vrais : \begin{itemize}
		\item Initialisation : $P_{n_0}$ est vraie
		\item Hérédité : $\forall n\in\llbracket n_0,n_1-1\rrbracket,\ P_n\implies P_{n+1}$
	\end{itemize}
	Alors pour tout $n\in\llbracket n_0,n_1\rrbracket$, $P_n$ est vraie
\end{prop}
\begin{proof}
	Cette fois, il suffit d'appliquer le principe de récurrence à partir d'un certain rang à la suite d'assertions définies par $$Q_n=\left \{\begin{array}{ll}P_n&\text{si }n\in\llbracket n_0,n_1\rrbracket \\ \text{vrai} & \text{si }n>n_1\end{array}\right.$$ $Q_n$ est alors vraie quel que soit $n\in\N$, donc en particulier pour $n\in\llbracket n_0,n_1\rrbracket$.
\end{proof}
\begin{prop}[Récurrence descendante]
	Soit $(P_n)_{n_0\leqslant n\leqslant n_1}$ une suite d'assertion avec $n_0\leqslant n_1$. On suppose que les deux points suivants sont vrais : \begin{itemize}
		\item Initialisation : $P_{n_1}$ est vraie
		\item Hérédité : $\forall n\in\llbracket n_0+1,n_1\rrbracket,\ P_{n+1}\implies P_{n}$
	\end{itemize}
	Alors pour tout $n\in\llbracket n_0,n_1\rrbracket$, $P_n$ est vraie
\end{prop}
\begin{proof}
	Il suffit d'appliquer le principe de récurrence finie avec $Q_n=P_{n_1+n_0-n}$.
\end{proof}
\subsection{Suites définies par récurrence}
Une fois que l'on maîtrise le concept de définition par récurrence, on peut établir des \textbf{définitions} par récurrence.
\begin{dfn}
	Définir une suite $(u_n)_{n\in\N}$ par récurrence consiste à donner son premier terme $u_0$, puis à définir une relation de récurrence (assertion logique quelconque) entre $u_n$ et $u_{n+1}$.
\end{dfn}
\begin{rmq}
	On prouvera la validité d'une telle définition en construisant $\N$.
\end{rmq}
\begin{rmq}
	On peut même imaginer définir une suite par une relation de récurrence à plusieurs pas, ou même par une relation de récurrence forte.
\end{rmq}
\newpage
\chapter{Ensembles et applications} %chapitre complet, sauf compléments%
\section{Introduction à la théorie des ensembles}
Une assertion du type $\forall x\in X,\ p(x)$ est en fait une abréviation de l'assertion $\forall x,\ x\in X\implies p(x)$. Ainsi, toute assertion du type $\forall x\in\emptyset,\ p(x)$ est toujours vraie puisque le faux implique n'importe quoi. \par Une assertion du type $\exists x\in X,\ p(x)$ est en fait une abréviation de l'assertion $\forall x,\ x\in X\land p(x)$. Ainsi, toute assertion du type $\exists x\in\emptyset,\ p(x)$ est toujours fausse. \par Mais pour effectuer ces raisonnements, on admet l'existence d'un ensemble, appelé ensemble vide, qui ne contient aucun élément. Le besoin de cet ensemble nous amène donc à construire une théorie des ensembles.
\subsection{Généralités}
\begin{dfn}[Ensemble]
	Un \textbf{ensemble} $E$ est une collection d'objet. Lorsqu'un objet $x$ appartient à l'ensemble $E$, on dit que c'est un \textbf{élément} de $E$ et on note $x\in E$. Les objets de $E$ sont notés entre accolades pour décrire $E$.
\end{dfn}
\begin{rmq}
	On ne se préoccupera pas des fondements théoriques de cette définition, bien trop abstraits et éloignés du programme de classe préparatoire.
\end{rmq}
\begin{rmq}
	$E=\left\{1,2,3\right\}$ est un ensemble.
\end{rmq}
\begin{dfn}[Ensemble vide]
	On admet l'existence d'un ensemble appelé \textbf{ensemble vide} et noté $\emptyset$ qui ne contient aucun élément. Quel que soit $x$, l'assertion $x\in\emptyset$ sera toujours fausse.
\end{dfn}
\begin{rmq}
	De même, on ne se préoccupera pas des fondements de l'existence de l'ensemble vide.
\end{rmq}
\begin{dfn}[Inclusion]
	Soit $E$ et $F$ deux ensembles. $E$ est dit \textbf{inclus} dans $F$ lorsque $$\forall x\in E,\ x\in F$$ On note alors $E\subset F$ ou encore $E\subseteq F$. On dit aussi que $F$ est un \textbf{sous-ensemble} de $E$.
\end{dfn}
\begin{exa}
	On a toujours $\emptyset\subset F$.
\end{exa}
\begin{exa}
	On a toujours $E\subset E$.
\end{exa}
\begin{dfn}[Égalité ensembliste]
	Soit $E$ et $F$ deux ensembles. $E$ est dit \textbf{égal} à $F$ lorsque $$\forall x,\ x\in E\iff x\in F$$ On note alors $E=F$. 
\end{dfn}
\begin{prop}
	Soit $E$ et $F$ deux ensembles. On a : $$E=F\iff (E\subset F)\land (F\subset E)$$
\end{prop}
\begin{met}[Montrer une égalité ensembliste : double inclusion]
	On en déduit la méthode générale pour montrer que deux ensembles $E$ et $F$ sont égaux, dite de "double inclusion" :
	\begin{itemize}
		\item Soit $x\in E$. Alors ... donc $x\in F$.
		\item Réciproquement, soit $x\in F$. Alors ... donc $x\in E$.
	\end{itemize}
\end{met}
\begin{rmq}[Insensibilité à la permutation et à la répétition]
	En vertu du principe de double-inclusion, l'ordre et la répétition n'ont pas d'importance dans la définition d'un ensemble. Par exemple, on a $\left\{1,2,3\right\}=\left\{1,3,2\right\}$ et $\left\{1,2,3\right\}=\left\{1,1,2,3\right\}$.
\end{rmq}
\begin{prop}[Unicité de l'ensemble vide, HP]
	L'ensemble vide est unique.
\end{prop}
\begin{proof}
	Soit $\emptyset$ et $\emptyset'$ deux ensembles vides. On a $\emptyset\subset\emptyset'$ et $\emptyset'\subset\emptyset$ donc $\emptyset=\emptyset'$.
\end{proof}
\begin{rmq}[Note historique]
	A l'occasion de cette courte démonstration, le grand FM avait déclaré : "on passe en mode hors-programme de ta mère".
\end{rmq}
\begin{dfn}[Ensemble des parties d'un ensemble]
	Soit $E$ un ensemble. Les sous-ensembles de $E$ forment un nouvel ensemble qu'on appelle \textbf{l'ensemble des parties de $E$}. On le note $\mathcal{P}(E)$.
\end{dfn}
\begin{rmq}
	Là non plus, on ne se préoccupera pas de savoir pourquoi les sous-ensembles de $E$ forment un nouvel ensemble.
\end{rmq}
\begin{exa}
	On a $\mathcal{P}(\emptyset)=\left\{\emptyset\right\}$ et $\mathcal{P}(\mathcal{P}(\emptyset))=\left\{\emptyset,\left\{\emptyset\right\}\right\}$.
\end{exa}
\begin{dfn}[Définition d'un ensemble par compréhension]
	Soit $E$ un ensemble et $p(x)$ une assertion qui dépend de $x\in E$. On peut définir l'ensemble des $x\in E$ qui satisfont $p(x)$. On le note $\left\{x\in E\ | \ p(x)\text{ vraie}\right\}$ ou $\left\{x\in E,\ p(x)\text{ vraie}\right\}$ voire $\left\{x\in E : p(x)\text{ vraie}\right\}$. On dit qu'il s'agit d'une définition par \textbf{compréhension}.
\end{dfn}
\begin{rmq}
	Le choix de la notation pour un ensemble défini par compréhension dépend de ce qui est utilisé pour définir l'assertion $p(x)$. Par exemple, lorsqu'on travaille avec de la divisibilité, puisque le symbole "divise" s'écrit "$\mid$", on préfèrera utiliser ":" ou ",".
\end{rmq}
\begin{exa}
	$\U=\left\{z\in\C\ : \lvert z\rvert=1\right\}$
\end{exa}
\begin{met}[Montrer une égalité ensembliste bis]
	Voici une deuxième façon de montrer que $A=B$. Elle ne fonctionne que si $A$ et $B$ sont inclus dans un même sur-ensemble de référence $E$. Il suffit alors de montrer que $\forall x\in E,\ x\in A\iff x\in B$ Bien noter la différence avec la définition, on se restreint ici à un sur-ensemble et non pas à l'ensemble des éléments des mathématiques.
\end{met}
\subsection{Opérations ensemblistes}
Jusqu'à ce qu'on l'écrase, on fixe $E$ un ensemble de référence et on s'intéresse à ses parties.
\begin{dfn}[Union]
	Soit $A$ et $B$ deux sous-ensembles de $E$. L'\textbf{union de $A$ et $B$} ou \textbf{réunion}, notée $A\cup B$ est définie par compréhension comme : $$A\cup B :=\left\{x\in E\ | \ (x\in A)\lor (x\in B)\right\}$$
\end{dfn}
\begin{dfn}[Intersection]
	Soit $A$ et $B$ deux sous-ensembles de $E$. L'\textbf{intersection de $A$ et $B$}, notée $A\cap B$ est définie par compréhension comme : $$A\cap B :=\left\{x\in E\ | \ (x\in A)\land (x\in B)\right\}$$
\end{dfn}
\begin{dfn}[Ensembles disjoints]
	Soit $A$ et $B$ deux sous-ensembles de $E$. On dit que $A$ et $B$ sont \textbf{disjoints} lorsque $A\cap B=\emptyset$.
\end{dfn}
\begin{exa}
	Voici quelques cas particuliers : 
	\begin{itemize}
		\item Si $A\subset B$, alors on a $A\cup B=B$ et $A\cap B=A$.
		\item Si $A=\emptyset$, alors on a $A\cup B=\emptyset\cup B=B$ et $A\cap B=\emptyset\cap B=\emptyset$.
		\item Si $B=E$, alors on a $A\cup B=A\cup E=E$ et $A\cap B=A\cap E=A$.
	\end{itemize}
\end{exa}
\begin{rmq}
	Attention : l'assertion $\forall x\in A,\ x\in A\iff x\in B$ ne montre que $A\subset B$.
\end{rmq}
\begin{prop}[Compatibilité de l'inclusion avec l'union et l'intersection]
	Soit $A$, $B$ et $C$ des sous-ensembles de $E$. On a $$A\subset B\implies \begin{cases}
		(A\cup C)\subset(B\cup C) \\
		(A\cap C)\subset(B\cap C)
	\end{cases}$$
\end{prop}
\begin{prop}[Commutativité de l'union et de l'intersection]
	Soit $A$ et $B$ deux sous-ensembles de $E$. On a : $$A\cup B=B\cup A \quad \text{et} \quad A\cap B=B\cap A$$
\end{prop}
\begin{dfn}[Complémentaire]
	Soit $A$ un sous-ensemble de $E$. Le \textbf{complémentaire} de $A$ dans $E$ est défini par compréhension comme : $$\complement_E^A:=\left\{x\in E\ | \ x\notin A\right\}$$ On le note plus souvent $\overline{A}$ lorsqu'il n'y a pas d'ambiguïté par rapport au sur-ensemble de référence.
\end{dfn}
\begin{exa}
	Attention, on n'a pas toujours $\complement_E^A\ne A$. Cela n'est vrai que dès que $E\ne\emptyset$.
\end{exa}
\begin{exa}
	Si $A$ est un sous ensemble de $E$, on a toujours $A\cup\complement_E^A=E$ et $A\cap\complement_E^A=\emptyset$.
\end{exa}
\begin{dfn}[Différence ensembliste]
	Soit $A$ et $B$ des sous-ensembles de $E$. Plus généralement, on définit la \textbf{différence ensembliste de $A$ et $B$} par compréhension comme : $$A\setminus B:=\left\{x\in A\ | \ x\notin B\right\}$$
\end{dfn}
\begin{exa}
	Soit $A$ un sous-ensemble de $E$. On a $\complement_E^A=E\setminus A$.
\end{exa}
\begin{prop}[Expression de la différence ensembliste, utile pour les exercices]
	Soit $A$ et $B$ deux sous-ensembles de $E$. On a $A\setminus B=A\cap \overline{B}$.
\end{prop}
\begin{prop}[Décroissance du passage au complémentaire]
	Soit $A$ et $B$ deux sous-ensembles de $E$. Si $A\subset B$, alors $\overline{B}\subset\overline{A}$.
\end{prop}
\begin{prop}[Involutivité du passage au complémentaire]
	Soit $A$ un sous-ensemble de $E$. On a : $$\overline{(\overline{A})}=A$$
\end{prop}
\begin{thm}[Lois de De Morgan]
	Soit $A$ et $B$ deux sous-ensembles de $E$. On a : $$\begin{cases}
		\overline{A\cup B}=\overline{A}\cap\overline{B} \\ \overline{A\cap B}=\overline{A}\cup\overline{B}
	\end{cases}$$
\end{thm}
\begin{rmq}[Note historique]
	A propos de ce théorème, le grand FM avait déclaré : "Je crois que ça s'appelle aussi les règles de Boole pour les ensembles. J'ai déjà entendu ça quelque part, et je crois pas que j'étais bourré".
\end{rmq}
\begin{prop}[Associativité de l'union et de l'intersection]
	Soit $A$, $B$ et $C$ des sous-ensembles de $E$. On a : $$\begin{cases}
		(A\cup B)\cup C=A\cup(B\cup C) \\ (A\cap B)\cap C=A\cap(B\cap C)
	\end{cases}$$ Cela permet d'écrire $A\cup B\cup C$ et $A\cap B\cap C$ sans ambiguïté.
\end{prop}
\begin{rmq}
	Attention : on n'écrit jamais $A\cup B\cap C$ ou $A\cap B\cup C$ sans parenthèse, car il y a alors une grosse ambiguïté !
\end{rmq}
\begin{prop}[Distributivité de $\cup$ sur $\cap$ et de $\cap$ sur $\cup$]
	Soit $A$, $B$ et $C$ des sous-ensembles de $E$. On a : $$\begin{cases}
		A\cup(B\cap C)=(A\cup B)\cap (A\cup C)\\ A\cap(B\cup C)=(A\cap B)\cup(A\cap C)
	\end{cases}$$
\end{prop}
\begin{dfn}[Couple]
	Soit $E$ et $F$ des ensembles. Soit $x\in E$ et $y\in F$. A partir de $x$ et $y$, on construit le \textbf{couple} $(x,y)$.
\end{dfn}
\begin{rmq}
	Ici, on ne se préoccupe pas de la construction d'un couple.
\end{rmq}
\begin{thm}[Propriété fondamentale des couples]
	Soit $E$ et $F$ des ensembles. On a : $$\forall x\in E,\forall x'\in E,\forall y\in F,\forall y'\in F,\ (x,y)=(x',y')\iff(x=x')\land(y=y')$$
\end{thm}
\begin{rmq}
	La négation de cette propriété fondamentale s'écrit donc : $$(x,y)\ne(x',y')\iff(x\ne x')\lor(y\ne y')$$
\end{rmq}
\begin{dfn}[Produit cartésien]
	Soit $E$ et $F$ deux ensembles. L'ensemble des couples $(x,y)$ pour $x\in E$ et $y\in F$ est appelé \textbf{produit cartésien de $E$ et $F$}. On le note $E\times F$.
\end{dfn}
\begin{rmq}
	Le produit cartésien n'a aucune raison d'être commutatif !
\end{rmq}
\begin{rmq}
	Soit $E$ et $F$ deux ensembles. On a toujours $\emptyset\times F=\emptyset$ et $E\times\emptyset=\emptyset$. $A$ $contrario$, si $E$ et $F$ sont non vides, alors $E\times F$ est non vide.
\end{rmq}
On aimerait généraliser. Pour commencer, on voudrait, si l'on dispose de $E$, $F$ et $G$ des ensembles ainsi que de $x\in E$, $y\in F$ et $z\in G$, construire un objet $(x,y,z)$, dit \textbf{triplet}, qui vérifie une propriété fondamentale du type : $$(x,y,z)=(x',y',z')\iff(x=x')\land(y=y')\land(z=z')$$ On peut procéder de deux façons : 
\begin{itemize}
	\item On pose $(x,y,z):=((x,y),z)$, une sorte de couple de couple, et on vérifie alors aisément qu'il vérifie la propriété fondamentale souhaitée.
	\item On peut procéder d'une autre manière que nous verrons plus tard dans ce chapitre.
\end{itemize}
Puis, dans le cas général, si $n\geqslant 2$ et $E_1,\ldots,E_n$ sont des ensembles, à partir de $x_1\in E_1,\ldots,x_n\in E_n$, on peut construire un \textbf{$n$-uplet} (aussi dit \textbf{$n$-liste}) $(x_1,\ldots,x_n)$ qui vérifie une propriété fondamentale du type : $$(x_1,\ldots,x_n)=(x_1',\ldots,x_n')\iff(x_1=x_1')\land\ldots\land(x_n=x_n')$$
\begin{dfn}
	L'ensemble des $n$-uplets à valeurs dans $E_1,\ldots,E_n$ est noté $E_1\times\ldots\times E_n$.
\end{dfn}
\begin{dfn}
	$E^n$ est une notation standard pour le produit cartésien $E\times\ldots\times E$ où l'ensemble $E$ est présent $n$ fois.
\end{dfn}
\begin{thm}[Axiome du choix, HP]
	Lorsqu'on réalise le produit cartésien infini d'une infinité d'ensemble tous non vides, on décide que ce produit cartésien est non vide et il existe alors une \textbf{fonction de choix} qui nous permet d'en choisir un élément.
\end{thm}
\begin{rmq}
	On n'est pas obligé de mettre des parenthèses intérieures dans $n$-uplet, même si les ensembles du produit cartésien sont distincts. Formellement, on identifie $$E_1\times\ldots\times E_l\times(E_{l+1}\times\ldots\times E_m)\times E_{m+1}\times\ldots\times E_n$$ à $E_1\times\ldots\times E_n$. C'est ce qu'on peut appeler l'\textbf{associativité du produit cartésien}.
\end{rmq}

\section{Applications}
Le programme ne fait pas de distinction entre les mots "application" et "fonction" (même s'il y en a une légère en réalité) : nous considérerons donc que ces mots sont synonymes.
\subsection{Définition formelle d'une application}
\begin{dfn}[Application, fonction]
	Soit $E$ et $F$ deux ensembles. Une \textbf{application}, ou \textbf{fonction}, de $E$ dans $F$ est un objet $f$ défini par : $$f=(E,F,\Gamma)$$ où $\Gamma\subset E\times F$ est appelée \textbf{graphe de $f$} et est telle que $\forall x\in E,\ \exists!y\in F,\ (x,y)\in\Gamma$. On la note plus souvent $$\begin{array}{ccccc}f&:&E&\rightarrow&F\\&&x&\mapsto&f(x)\end{array}$$ $f$ \textbf{associe} à chaque $x\in E$ une \textbf{unique image} $f(x)$. On dit que $x$ est \textbf{un antécédent} de $f(x)$ par $f$. 
\end{dfn}
\begin{thm}
	Soit $E$, $E'$, $F$ et $F'$ des ensembles. Soit $f:E\rightarrow F$ et $g:E'\rightarrow F'$ deux applications. Alors on a : $$f=g\iff \begin{cases}
		E=E' \\ F=F' \\ \forall x\in E,\ f(x)=g(x)
	\end{cases}$$
\end{thm}
\begin{rmq}
	Le troisième point traduit l'égalité des graphes.
\end{rmq}
\begin{rmq}
	Pour représenter $f$, si $E$ et $F$ sont des parties de $\R$ on peut dessiner son graphe. Sinon, on faire un diagramme sagittal.
\end{rmq}
\begin{dfn}[Ensembles des fonctions]
	Soit $E$ et $F$ des ensembles. L'ensemble des fonctions $f:E\rightarrow F$ est noté $\mathcal{F}(E,F)$ ou $F^E$.
\end{dfn}
\begin{dfn}[Fonction indicatrice]
	Soit $E$ un ensemble et $A$ un sous-ensemble de $E$. L'\textbf{indicatrice de $A$ dans $E$} est la fonction : $$\begin{array}{ccccc}\mathds{1}_A&:&E&\rightarrow&\R \\ &&x&\mapsto&\begin{cases}
			0\text{ si }x\notin A \\ 1\text{ si }x\in A
	\end{cases}\end{array}$$
\end{dfn}
\begin{prop}[Propriétés des indicatrices]
	Soit $E$ un ensemble et $A$ et $B$ des sous-ensembles de $E$. On a : 
	\begin{itemize}
		\item $\mathds{1}_{A\cap B}=\mathds{1}_A\times\mathds{1}_B$
		\item $\mathds{1}_{\overline{A}}=1-\mathds{1}_A$
		\item $\mathds{1}_{A\cup B}=\mathds{1}_A+\mathds{1}_B-\mathds{1}_{A\cap B}=\mathds{1}_A+\mathds{1}_B-\mathds{1}_A\times\mathds{1}_B$
	\end{itemize}
\end{prop}
\begin{dfn}[Restriction]
	Soit $E$ et $F$ deux ensembles, ainsi que $A$ un sous-ensemble de $E$. Soit $f:E\rightarrow F$ une fonction. La \textbf{restriction de $f$ à $A$} est la fonction : $$\begin{array}{ccccc}f_{|A}&:&A&\rightarrow&F\\ &&x&\mapsto&f(x)\end{array}$$
\end{dfn}
\begin{dfn}[Corestriction]
	Soit $E$ et $F$ deux ensembles, ainsi que $B$ un sous-ensemble de $F$. Soit $f:E\rightarrow F$ une fonction. Si $\forall x\in E,\ f(x)\in B$, alors on peut \textbf{corestreindre} $f$ à $B$ et la \textbf{corestriction de $f$ à $B$} est la fonction : $$\begin{array}{ccccc}f^{|B}&:&E&\rightarrow&B\\ &&x&\mapsto&f(x)\end{array}$$
\end{dfn}
\begin{rmq}
	On peut même considérer $f_{|A}^{|B}$ à condition que $\forall x\in A,\ f(x)\in B$.
\end{rmq}
\begin{dfn}[Prolongement]
	Soit $E$, $E'$, $F$ et $F'$ des ensembles tels que $E\subset E'$ et $F\subset F'$. Soit $f:E\rightarrow F$. \textbf{Un prolongement} de $f$ est une fonction $g:E'\rightarrow F'$ telle que $$g_{|E}^{|F}=f$$
\end{dfn}
\begin{rmq}
	Attention : il peut exister plusieurs prolongements d'une même fonction !
\end{rmq}
\subsection{Injectivité, surjectivité, bijectivité}
Dans toute cette partie, on fixe $E$ et $F$ deux ensembles ainsi que $f:E\rightarrow F$ une fonction.
\begin{dfn}[Injectivité]
	$f$ est dite \textbf{injective} lorsque $$\forall(x,x')\in E^2,\ f(x)=f(x')\implies x=x'$$ ou de manière équivalente lorsque $$\forall (x,x')\in E^2,\ x\ne x'\implies f(x)\ne f(x')$$ En rédaction automatique, pour montrer que $f$ est injective, on écrira "Soit $(x,x')\in E^2$ tel que $f(x)=f(x')$" puis on montrera que $x=x'$ (ou alors on utilisera l'autre définition équivalente).
\end{dfn}
\begin{rmq}
	La non injectivité s'écrit donc : $\exists(x,x')\in E^2,\ (x\ne x')\land(f(x)=f(x'))$.
\end{rmq}
\begin{prop}
	Toute fonction réelle strictement monotone est injective.
\end{prop}
\begin{rmq}
	Attention : la réciproque est fausse !
\end{rmq}
\begin{rmq}
	La fonction $(\emptyset,\emptyset,\emptyset)$ est à la fois constante, croissante, décroissante, strictement croissante et strictement décroissante. Il en va de même pour les fonctions $(\emptyset,F,\emptyset)$ et $(\left\{a\right\},\R,\left\{(a,b)\right\})$.
\end{rmq}
\begin{dfn}[Injection canonique]
	Soit $E$ un ensemble et $A$ un sous-ensemble de $E$. \textbf{L'injection canonique de $A$ dans $E$} est l'application : $$\begin{array}{ccccc}\iota&:&A&\rightarrow&E\\&&x&\mapsto&x\end{array}$$ Elle est injective.
\end{dfn}
\begin{prop}
	L'injectivité s'hérite par restriction et corestriction.
\end{prop}
\begin{dfn}
	$f$ est dite \textbf{surjective} lorsque $$\forall y\in F,\ \exists x\in E,\ y=f(x)$$ En rédaction automatique, on écrira "Soit $y\in F$" puis on exhibera un candidat après avoir cherché par condition nécessaire. On pourra aussi utiliser des théorèmes d'existence comme le TVI dans $\R$.
\end{dfn}
\begin{exa}
	L'application $\begin{array}{ccc}\R&\rightarrow&\U\\ \theta&\mapsto&\exp(i\theta)\end{array}$ est surjective.
\end{exa}
Injectivité et surjectivité ne sont pas incompatibles. C'est ce qui fait l'objet de la
\begin{dfn}[Bijectivité]
	$f$ est dite \textbf{bijective} lorsqu'elle est à la fois injective et surjective. De manière équivalente : $$\forall y\in F,\ \exists!x\in E,\ y=f(x)$$
\end{dfn}
\begin{met}[Montrer une bijectivité]
	Pour montrer qu'une fonction $f:E\rightarrow F$ est bijective on peut :
	\begin{itemize}
		\item Se donner $y\in F$ et $x\in E$ et montrer que l'équation $y=f(x)$ admet une unique solution.
		\item Montrer séparément que $f$ est injective et surjective.
		\item Si on est dans $\R$, réaliser un joli tableau de variations.
	\end{itemize}
\end{met}
\begin{dfn}[Fonction identité]
	La \textbf{fonction identité de $E$} est définie par : $$\begin{array}{ccccc}\id_E&:&E&\rightarrow&E\\&&x&\mapsto&x\end{array}$$ Elle est bijective.
\end{dfn}
\begin{dfn}[Composée]
	Soit $E$, $F$ et $G$ des ensembles ainsi que $f:E\rightarrow F$ et $g:F\rightarrow G$ des fonctions. Alors, la \textbf{composée de $g$ par $f$} est la fonction définie par : $$\begin{array}{ccccc}g\circ f&:&E&\rightarrow&G\\&&x&\mapsto&g(f(x))\end{array}$$
\end{dfn}
\begin{rmq}
	Soit $E$ et $F$ des ensembles ainsi que $f:E\rightarrow F$. On a : $$\begin{cases}
		f\circ\id_E=f \\ \id_F\circ f=f
	\end{cases}$$
\end{rmq}
\begin{rmq}
	Attention : la loi $\circ$ est non commutative en général ! Déjà, elle n'est souvent même pas définie dans les deux sens pour une question de domaines, et même si les fonctions étaient définies, la coïncidence n'a aucune raison d'avoir lieu. On trouvera un contre-exemple avec une fonction constante et une fonction linéaire dans $\R$.
\end{rmq}
\begin{prop}[Associativité de la composition]
	Soit $E$, $F$, $G$ et $H$ des ensembles. Soit $f:E\rightarrow F$, $g:F\rightarrow G$ et $h:G\rightarrow H$ des fonctions. Alors, on a : $$h\circ(g\circ f)=(h\circ g)\circ f$$
\end{prop}
\begin{thm}
	Soit $E$, $F$ et $G$ des ensembles. Soit $f;E\rightarrow F$ et $g:F\rightarrow G$ des fonctions.
	\begin{itemize}
		\item Si $f$ et $g$ sont injectives, alors $g\circ f$ est injective.
		\item Si $f$ et $g$ sont surjectives, alors $g\circ f$ est surjective.
		\item Si $f$ et $g$ sont bijectives, alors $g\circ f$ est bijective.
	\end{itemize}
\end{thm}
\begin{thm}[Réciproque partielle]
	Soit $E$, $F$ et $G$ des ensembles. Soit $f:E\rightarrow F$ et $g:F\rightarrow G$ des fonctions. Réciproquement :
	\begin{itemize}
		\item Si $g\circ f$ est injective, alors $f$ est injective.
		\item Si $g\circ f$ est surjective, alors $g$ est surjective.
	\end{itemize}
\end{thm}
\begin{rmq}
	La construction de $g\circ f$ reste valable si $f:E\rightarrow F'$ et $g:F\rightarrow G$ avec $F'\subset F$. L'essentiel des résultats précédents reste valable, \textbf{sauf la composée de deux surjections (resp. bijections) qui n'a plus aucune raison d'être une surjection (resp. bijection)}.
\end{rmq}
\subsection{Bijection réciproque}
\begin{dfn}[Bijection réciproque]
	Soit $E$ et $F$ deux ensembles et $f:E\rightarrow F$ une bijection. La \textbf{bijection réciproque de $f$} est la fonction définie de la manière suivante : soit $y\in F$. $f$ est bijective donc $\exists!x\in E,\ y=f(x)$. Fixons-le. On pose alors $$f\inv(y):=x$$ La fonction $f\inv$ est alors elle aussi bijective, et elle vérifie $$\begin{cases}
		f\inv\circ f=\id_E \\ f\circ f\inv=\id_F
	\end{cases}$$
\end{dfn}
\begin{thm}
	Soit $E$ et $F$ des ensembles. Soit $f:E\rightarrow F$ une fonction quelconque. S'il existe $g:F\rightarrow E$ telle que $g\circ f =\id_E$ et $f\circ g=\id_F$ alors : 
	\begin{enumerate}
		\item $f$ est bijective
		\item $g$ est la bijection réciproque de $f$
	\end{enumerate}
\end{thm}
\begin{rmq}
	Il s'agit donc là d'une nouvelle méthode pour prouver la bijectivité d'une fonction, qui par la même occasion nous fournit sa bijection réciproque.
\end{rmq}
\begin{cor}[Involutivité du passage à la bijection réciproque]
	Soit $E$ et $F$ des ensembles. Soit $f:E\rightarrow F$ une fonction. Si $f$ est bijective, de bijection réciproque $f\inv$, alors $f\inv$ est bijective de bijection réciproque $f$.
\end{cor}
\begin{cor}[Bijection réciproque d'une composée de bijections]
	Soit $E$, $F$ et $G$ des ensembles. Soit $f:E\rightarrow F$ et $g:F\rightarrow G$ des fonctions. Si $f$ et $g$ sont bijectives, alors : 
	\begin{itemize}
		\item $g\circ f$ est bijective (on le savait déjà par un théorème précédent)
		\item La bijection réciproque de $g\circ f$ est donnée par $f\inv\circ g\inv$. Autrement dit, on a alors $(g\circ f)\inv=f\inv\circ g\inv$.
	\end{itemize}
\end{cor}
\begin{rmq}
	Attention à bien inverser $f$ et $g$ lorsqu'on passe à la bijection réciproque ! Il s'agit en fait là d'un résultat plus général que nous reverrons dans le chapitre "Introduction à l'algèbre".
\end{rmq}
\begin{dfn}[Involution]
	Soit $E$ un ensemble et $f:E\rightarrow E$. On dit que $f$ est \textbf{involutive} lorsque $f\circ f=\id_E$, $ie$ lorsque $\forall x\in E,\ f(f(x))=x$.
\end{dfn}
\begin{prop}
	Toute fonction involutive est bijection, de bijection réciproque elle-même.
\end{prop}
\begin{dfn}[Permutation]
	Soit $E$ un ensemble. Une permutation de $E$ est une bijection $E\rightarrow E$. On note $\mathcal{S}_E$ ou $\mathfrak{S}_E$ l'ensemble des permutations de $E$.
\end{dfn}
\begin{exa}
	Soit $E$ un ensemble. L'application $\begin{array}{ccc}\mathfrak{S}_E&\rightarrow&\mathfrak{S}_E\\ f&\mapsto&f\inv\end{array}$ est une involution.
\end{exa}
\begin{met}[Bonne définition d'une fonction]
	Soit $E$ et $F$ des ensembles. Voici comment montrer que $\begin{array}{ccccc}f&:&E&\rightarrow&F\\&&x&\mapsto&f(x)\end{array}$ est bien définie. Il y a trois points de vigilance, à vérifier dans l'ordre dans lequel ils sont énoncés :
	\begin{itemize}
		\item Est-ce que $x$ est un objet primaire ? Vérifier qu'il ne s'agit pas d'une formule d'un autre objet plus simple. Dans le cas où $x$ n'est pas un objet primaire et possède plusieurs représentants (rationnels, éléments d'ensembles quotients par exemple), il faut montrer que la valeur de $f(x)$ ne dépend pas du représentant choisi.
		\item Calculabilité ? Soit $x\in E$. Montrer que $f(x)$ existe.
		\item Bon ensemble d'arrivée ? Soit $x\in E$. Montrer que $f(x)\in F$.
	\end{itemize}
\end{met}
\begin{met}[Prouver une propriété sur $f\inv$]
	Pour montrer une propriété sur $f\inv$, on montre $f$ de cette propriété puis on se débrouille pour simplifier par $f$.
\end{met}
\begin{exe}
	Pour s'exercer, on pourra montrer que si $f:\R\rightarrow\R$ est bijective et impaire, alors $f\inv$ est impaire.
\end{exe}
\begin{exa}
	$\begin{array}{ccc}\R_+&\rightarrow&\R_+\\y&\mapsto&\sqrt{y}\end{array}$ est définie comme la bijection de réciproque de $\begin{array}{ccc}\R_+&\rightarrow&\R_+\\x&\mapsto&x^2\end{array}$. Or, la deuxième est strictement croissante, donc la première est aussi strictement croissante (appliquer la méthode et utiliser la contraposée de la croissance).
\end{exa}
\begin{dfn}
	Soit $E$ et $F$ deux ensembles. Les deux assertions suivantes sont équivalentes :\begin{enumerate}
		\item $\exists f:E\rightarrow F$ bijective
		\item $\exists g:F\rightarrow E$ bijective
	\end{enumerate}
	Lorsque l'une de ces deux assertions est vérifiée, on dit que $E$ et $F$ sont \textbf{équipotents}.
\end{dfn}
\begin{thm}[Théorème de Cantor-Bernstein HP]
	Soit $E$ et $F$ deux ensembles. S'il existe $f:E\rightarrow F$ injective et $g:F\rightarrow E$ injective, alors $E$ et $F$ sont équipotents.
\end{thm}
\subsection{Image directe, image réciproque}
\begin{dfn}[Image directe]
	Soit $E$ et $F$ des ensembles et $f:E\rightarrow F$. Soit $A$ un sous-ensemble de $E$. L'\textbf{image directe de $A$ par $f$} est l'ensemble défini par : $$f(A):=\left\{y\in F\ | \ \exists x\in A,\ y=f(x)\right\}$$
\end{dfn}
\begin{dfn}[Image réciproque]
	Soit $E$ et $F$ des ensembles et $f:E\rightarrow F$. Soit $B$ un sous-ensemble de $F$. L'\textbf{image réciproque de $B$ par $f$} est l'ensemble défini par : $$f\inv(B):=\left\{x\in E\ | \ f(x)\in B\right\}$$
\end{dfn}
\begin{rmq}
	Même si les fonctions précédemment définies vont d'un ensemble de parties dans un autre, on les note quand même $f$ et $f\inv$ à l'usage. \textbf{En cas de doute, vérifier le type d'objet d'entrée et de sortie}.
\end{rmq}
\begin{rmq}
	Si $f$ est bijective, $f\inv(B)$ peut être interprété comme l'image directe de $B$ par $f\inv$ ou comme l'image réciproque de $B$ par $f$. Heureusement, les deux notions coïncident alors, il s'agit d'une simple double inclusion.
\end{rmq}
\begin{exa}
	Soit $E$ et $F$ deux ensembles et $f:E\rightarrow F$. On a :
	\begin{itemize}
		\item $f(\emptyset)=\left\{y\in F\ | \ \exists x\in\emptyset, \ f(x)=y\right\}=\emptyset$
		\item $f\inv(\emptyset)=\left\{x\in E \ | \ f(x)\in\emptyset\right\}=\emptyset$
		\item $f(E)=\left\{y\in F\ | \ \exists x\in E, f(x)=y\right\}=:\im(f)$
		\item $f\inv(F)=\left\{x \in E\ | \ f(x)\in F\right\}=E$
		\item Soit $x_0\in E$. $f(\left\{x_0\right\})=\left\{y\in F\ | \ \exists\left\{x_0\right\},\ y=f(x)\right\}=\left\{f(x_0)\right\}$
	\end{itemize}
\end{exa}
\begin{prop}
	Soit $E$ et $F$ deux ensembles. $f:E\rightarrow F$ est surjective si, et seulement si, $\im(f)=F$.
\end{prop}
\begin{prop}
	Soit $E$ et $F$ deux ensembles ainsi que $f:E\rightarrow F$. Alors $f^{|\im(f)}$ existe toujours et est surjective par construction.
\end{prop}
\begin{dfn}[Ensemble défini par image directe]
	L'ensemble $f(A)$ peut aussi être noté \textbf{symboliquement} : $$\left\{f(x)\ | \ x\in A\right\}$$ C'est ce qu'on appelle la définition d'un ensemble par \textbf{image directe}.
\end{dfn}
\begin{exe}
	Soit $E$, $F$ et $G$ des ensembles. Soit $f:E\rightarrow F$ et $g:F\rightarrow G$ des fonctions.
	\begin{enumerate}
		\item Soit $A$ un sous-ensemble de $E$. Montrer que $(g\circ f)(A)=g(f(A))$.
		\item Soit $B$ un sous-ensemble de $G$. montrer que $(g\circ f)\inv(B)=f\inv(g\inv(B))$.
	\end{enumerate}
	(Procéder par double-inclusion à chaque fois)
\end{exe}
\begin{prop}[Croissance des fonctions image directe et image réciproque]
	Soit $E$ et $F$ deux ensembles et $f:E\rightarrow F$. On a : 
	\begin{itemize}
		\item $\forall A_1,A_2\subset E,\ A_1\subset A_2\implies f(A_1)\subset f(A_2)$
		\item $\forall B_1,B_2\subset F,\ B_1\subset B_2\implies f\inv(B_1)\subset f\inv(B_2)$
	\end{itemize}
\end{prop}
\begin{prop}
	Soit $E$ et $F$ deux ensembles et $f:E\rightarrow F$. On a :
	\begin{itemize}
		\item $\forall B_1,B_2\subset F,\ f\inv(B_1\cup B_2)=f\inv(B_1)\cup f\inv(B_2)$
		\item $\forall B_1,B_2\subset F,\ f\inv(B_1\cap B_2)=f\inv(B_1)\cap f\inv(B_2)$
		\item $\forall B\subset F,\ f\inv(\overline{B})=\overline{f\inv(B)}$
		\item $\forall A_1,A_2\subset E, \ f(A_1\cup A_2)=f(A_1)\cup f(A_2)$
	\end{itemize}
\end{prop}
\begin{rmq}
	Attention, on a $f(A_1\cap A_2)\ne f(A_1)\cap f(A_2)$ en général (sauf lorsque $f$ est injective) et de même on a $f(\overline{A})\ne\overline{f(A)}$ en général.
\end{rmq}
\subsection{Familles}
\begin{dfn}[Famille]
	Soit $I$ et $E$ des ensembles. Formellement, une \textbf{famille $(x_i)_{i\in I}$ indexée par $I$ et à valeurs dans $E$} est une application : $$\begin{array}{ccccc}(x_i)_{i\in I}&:&I&\rightarrow&E\\&&i&\mapsto&x_i\end{array}$$ $i$ s'appelle l'\textbf{indice}. On note donc $E^I$ l'ensemble des familles indexées par $I$ à valeurs dans $E$.
\end{dfn}
\begin{exa}
	Soit $X$ un ensemble. Une \textbf{suite à valeurs dans $X$} est une application $(x_n)_{n\in\N}\in X^\N$.
\end{exa}
\begin{dfn}[Sous-famille]
	Soit $I$ et $E$ des ensembles. Soit $(x_i)_{i\in I}\in E^I$ et $J\subset I$. Une \textbf{sous-famille} de la famille $(x_i)_{i\in I}$ est une famille du type $(x_i)_{i\in J}$.
\end{dfn}
\begin{rmq}[Retour sur la construction des $n$-uplets, HP]
	Soit $n\in\N\st$. Une façon de construire le $n$-uplet $(x_1,\ldots,x_n)$ à valeurs dans $E$ est de considérer la famille $(x_i)_{i\in\llbracket1,n\rrbracket}\in E^{\llbracket1,n\rrbracket}$ On dispos alors de la propriété fondamentale des $n$-uplets. En quelque sorte, on construit $E^n$ sous la forme $E^{\llbracket1,n\rrbracket}$. Notons que cela nécessitait avant de construire les triplets comme couples de couple, puisque les fonctions qui nous servent ici à définir les $n$-uplets sont elles-mêmes des triplets.
\end{rmq}
\begin{dfn}[Union et intersection d'une famille de parties]
	Soit $I$ et $E$ des ensembles. Soit $(A_i)_{i\in I}\in\mathcal{P}(E)^I$. On pose : $$\bigcup_{i\in I}A_i:=\left\{x\in E \ | \ \exists i\in I,\ x\in A_i\right\}$$ $$\bigcap_{i\in I}A_i:=\left\{x\in E \ | \ \forall i\in I,\ x\in A_i\right\}$$  
\end{dfn}
\begin{dfn}[Famille de parties deux à deux disjointes]
	Soit $I$ et $E$ des ensembles. Soit $(A_i)_{i\in I}\in\mathcal{P}(E)^I$. On dit que \textbf{les $A_i$ sont deux à deux disjoints} lorsque $$\forall(i,j)\in I^2,\ i\ne j\implies A_i\cap A_j=\emptyset$$
\end{dfn}
\begin{dfn}[Union carrée/union disjointe]
	Soit $I$ et $E$ des ensembles. Soit $(A_i)_{i\in I}\in\mathcal{P}(E)^I$. Si les $A_i$ sont deux à deux disjoints, $\displaystyle\bigcup_{i\in I}A_i$ sera plutôt notée $\displaystyle\bigsqcup_{i\in I}A_i$. On parle d'\textbf{union disjointe} ou d'\textbf{union carrée}.
\end{dfn}
\begin{dfn}[Recouvrement]
	Soit $I$ et $E$ des ensembles. Soit $(A_i)_{i\in I}\in\mathcal{P}(E)^I$. Si $\displaystyle\bigcup_{i\in I}A_i=E$, on dit que les $A_i$ forment un \textbf{recouvrement} de $E$.
\end{dfn}
\begin{dfn}[Recouvrement disjoint]
	Soit $I$ et $E$ des ensembles. Soit $(A_i)_{i\in I}\in\mathcal{P}(E)^I$. Si $\displaystyle\bigsqcup_{i\in I}A_i=E$ (on suppose donc implicitement que les $A_i$ sont deux à deux disjoints), on dit que les $A_i$ forment un \textbf{recouvrement disjoint} de $E$.
\end{dfn}
\begin{exe}[NE+FHP]
	Soit $A$ un ensemble appelé \textbf{alphabet}. Intuitivement, un \textbf{mot} de longueur $n\in\N$ sera un élément de $A^n$. On voudrait définir le lexique d'une langue comme une partie de $\displaystyle\bigcup_{n\in\N}A^n$. Mais cette union a-t-elle un sens ? Autrement dit, les $A^n$ peuvent-ils être vus comme appartenant à un même $\mathcal{P}(E)$.
	\begin{enumerate}
		\item Construire un ensemble $E$ tel que $\forall n\in\N,\ A^n\subset E$ (voir $A^n$ comme $A^{\llbracket1,n\rrbracket}$ puis on prendra $E=\mathcal{P}(\N)\times\left\{A\right\}\times\mathcal{P}(\N\times A)$)
		\item On voudrait définir la longueur d'un mot comme l'unique $n\in\N$ tel que ce mot appartienne à $A^n$. Mais un tel $n$ est-il bien unique ? Montrer que les $A^n$ sont deux à deux disjoints.
	\end{enumerate}
\end{exe}
Dans un ensemble, la plupart des résultat vus avec deux parties restent vraies avec une famille de parties, et nous en faisons la liste maintenant.
\begin{prop}[Lois de De Morgan]
	Soit $I$ et $E$ des ensembles. Soit $(A_i)_{i\in I}\in\mathcal{P}(E)^I$. On a $$\overline{\bigcup_{i\in I}A_i}=\bigcap_{i\in I}\overline{A_i}$$ $$\overline{\bigcap_{i\in I}A_i}=\bigcup_{i\in I}\overline{A_i}$$
\end{prop}
\begin{prop}[Distributivités]
	Soit $I$ et $E$ des ensembles. Soit $A\subset E$ et $(B_i)_{i\in I}\in\mathcal{P}(E)^I$. On a $$A\cup\left(\bigcap_{i\in I}B_i\right)=\bigcap_{i\in I}(A\cup B_i)$$ $$A\cap\left(\bigcup_{i\in I}B_i\right)=\bigcup_{i\in I}(A\cap B_i)$$
\end{prop}
\begin{prop}
	Soit $I$, $E$ et $F$ des ensembles. Soit $f:E\rightarrow F$, $(A_i)_{i\in I}\in\mathcal{P}(E)^I$ et $(B_i)_{i\in I}\in\mathcal{P}(F)^I$. On a : \begin{itemize}
		\item $\displaystyle f\inv\left(\bigcup_{i\in I}B_i\right)=\bigcup_{i\in I}f\inv(B_i)$
		\item $\displaystyle f\inv\left(\bigcap_{i\in I}B_i\right)=\bigcap_{i\in I}f\inv(B_i)$
		\item $\displaystyle f\left(\bigcup_{i\in I}A_i\right)=\bigcup_{i\in I}f(A_i)$
	\end{itemize}
\end{prop}
\begin{prop}[Cas des unions disjointes]
	Soit $I$, $E$ et $F$ des ensembles. Soit $A\subset F$ et $(B_i)_{i\in I}\in\mathcal{P}(F)^I$. Supposons que les $B_i$ sont deux à deux disjoints. Alors : \begin{itemize}
		\item $\displaystyle A\cap\left(\bigsqcup_{i\in I}B_i\right)=\bigsqcup_{i\in I}A\cap B_i$
		\item $\displaystyle f\inv\left(\bigsqcup_{i\in I}B_i\right)=\bigsqcup_{i\in I}f\inv(B_i)$
	\end{itemize}
\end{prop}
\begin{dfn}[Généralisation du produit cartésien]
	Soit $I$ et $E$ des ensembles. Soit $(A_i)_{i\in I}\in\mathcal{P}(E)^I$. On pose : $$\prod_{i\in I}A_i:=\left\{(x_i)_{i\in I}\in E^I\ | \ \forall i\in I,\ x_i\in A_i\right\}$$
\end{dfn}
\begin{rmq}
	A une famille on peut associer un ensemble et vice-versa.
	\begin{itemize}
		\item A la famille $(x_i)_{i\in I}\in E^I$ on peut associer l'ensemble $\left\{x_i\ | \ i\in I\right\}$
		\item A l'ensemble $A$, on peut associer la famille $(a)_{a\in A}$
	\end{itemize}
	A chaque fois, famille et ensemble auront les mêmes éléments. Toutefois, on perde toute information sur l'ordre ou la répétition en passant aux ensembles. \par Dans le même esprit, soit $\mathcal{A}\in\mathcal{P}(E)$ ($\mathcal{A}$ n'est plus une famille de parties de $E$ mais un ensemble de parties de $E$). On posera alors $$\bigcup_{A\in\mathcal{A}}A:=\left\{x\in E\ | \ \exists A\in\mathcal{A},\ x\in A\right\}$$ $$\bigcap_{A\in\mathcal{A}}A:=\left\{x\in E\ | \ \forall A\in\mathcal{A},\ x\in A\right\}$$
\end{rmq}

\section{Ensembles finis}
\begin{lem}[Principe des tiroirs, cas $\llbracket1,n\rrbracket$]
	Soit $m,n\in\N$ tels que $m>n$ et $f:\llbracket1,m\rrbracket\rightarrow\llbracket1,n\rrbracket$. Alors, $f$ n'est pas injective.
\end{lem}
\begin{dfn}[Ensemble fini]
	Soit $E$ un ensemble. $E$ est dit \textbf{fini} lorsqu'il existe $n\in\N$ tel que $E$ soit équipotent à $\llbracket1,n\rrbracket$.
\end{dfn}
\begin{dfn}[Cardinal d'un ensemble fini]
	Soit $E$ un ensemble fini. Alors le $n\in\N$ tel que $E$ soit équipotent à $\llbracket1,n\rrbracket$ est unique. On définit alors le cardinal de $E$ par : $$\card(E):=n$$
\end{dfn}
\begin{exa}
	Soit $a,b\in\Z$ tels que $a\leqslant b$. Alors $\llbracket a,b\rrbracket$ est fini, de cardinal $b-a+1$.
\end{exa}
\begin{dfn}[Famille finie, taille d'une famille finie]
	Une famille $(x_i)_{i\in I}\in E^I$ est dite \textbf{finie} lorsque $I$ est fini. La \textbf{taille} de la famille est alors définie comme le cardinal de $I$.
\end{dfn}
\begin{dfn}[Dénombrer]
	Dénombrer un ensemble fini consiste à calculer son cardinal.
\end{dfn}
\begin{met}[Montrer une égalité d'ensembles finis]
	Soit $E$ un ensemble finie et $A$ un sous-ensemble de $E$. Si $A\subset E$ et $\card(A)=\card(E)$, alors $A=E$.
\end{met}
\begin{prop}
	Soit $A$ et $B$ deux ensembles finis. Alors $A\cup B$ est un ensemble fini, et on a $$\card(A\cup B)\leqslant\card(A)+\card(B)$$
\end{prop}
\begin{prop}
	Soit $A_1,\ldots,A_n\subset E$ toutes finies. Alors $A_1\cup\ldots\cup A_n$ est finie et on a $$\card(A_1\cup\ldots\cup A_n)\leqslant\card(A_1)+\ldots+\card(A_n)$$
\end{prop}
\begin{prop}
	L'image directe d'un ensemble fini est finie.
\end{prop}
\begin{prop}[Caractérisation des ensembles infinis]
	Soit $E$ un ensemble. $E$ est infini si, et seulement si, il existe une suite à valeurs dans $E$ injective ($ie$ ses termes sont deux à deux distincts).
\end{prop}
\begin{exa}
	L'ensemble $[0,1]$ est infini : il suffit de considérer la suite $\displaystyle\left(\frac{1}{n+1}\right)_{n\in\N}$ qui est strictement décroissante donc injective.
\end{exa}
\begin{prop}
	Soit $E$ un ensemble fini et $F$ un ensemble quelconque. Les assertions suivantes sont équivalentes :
	\begin{enumerate}
		\item $F$ est fini de même cardinal que $E$
		\item Il existe une bijection de $E$ dans $F$
		\item Il existe une bijection de $F$ dans $E$
	\end{enumerate}
\end{prop}
\begin{thm}[Principe des tiroirs, cas général]
	Soit $E$ et $F$ des ensembles finis tels que $\card(E)>\card(F)$. Alors, il n'existe pas d'injection de $E$ dans $F$. Autrement dit, toute fonction $f:E\rightarrow F$ n'est pas injective.
\end{thm}
\begin{rmq}[Extension au cas où $E$ est infini]
	\textbf{Le résultat s'étend au cas où $E$ est infini}. En effet, s'il existait une injection de $E$ dans $F$, alors celle-ci induirait par restriction une injection d'un sous-ensemble fini de $E$ de cardinal strictement supérieur à celui de $F$ (qui existe car $E$ est infini) dans $F$, ce qui est absurde en vertu du théorème précédent. 
\end{rmq}
\begin{thm}
	Soit $E$ et $F$ des ensembles finis de même cardinal. Alors les trois assertion suivantes sont équivalentes : \begin{itemize}
	\item $f$ est injective
	\item $f$ est surjective
	\item $f$ est bijective
\end{itemize}
\end{thm}
\section{Compléments, HP}
\subsection{Théorème de Cantor-Bernstein}
Nous proposons ici une démonstration du théorème de Cantor-Bernstein.
\begin{thm}[Théorème de Cantor-Bernstein]
	S'il existe une injection de $E$ dans $F$ et et une injection de $F$ dans $E$, alors il existe une bijection de $E$ dans $F$.
\end{thm}
\begin{proof}
	Soit $f$ une injection de $E$ dans $F$ et $g$ une injection de $F$ dans $E$. Pour chaque élément $x\in E$, on regarde la liste de ses "ancêtres" successifs : on cherche s'il possède un antécédent $y$ par $g$. Si c'est le cas, on cherche si cet antécédent possède un antécédent par $g$ dans $F$, et ainsi de suite. Formellement :
	\begin{itemize}
		\item On définit $E_0:=\left\{x\in E\ | \ \forall y\in F,\ g(y)\ne x\right\}$ l'ensemble des $x\in E$ sans ancêtre.
		\item De même, on définit $F_0:=\left\{y\in F\ | \ \forall x\in E,\ f(x)\ne y\right\}$ l'ensemble des $y\in F$ sans ancêtre.
		\item Par récurrence, on définit $E_{n+1}:=g(F_n)$ et $F_{n+1}:=f(E_n)$.
		\item On pose $\displaystyle E_p:=\bigcup_{n\in\N}E_{2n}$, $\displaystyle E_i:=\bigcup_{n\in\N}E_{2n+1}$ et $\displaystyle E_{\infty}:=E\setminus(E_p\cup E_i)$.
		\item De même, on pose $\displaystyle F_p:=\bigcup_{n\in\N}F_{2n}$, $\displaystyle F_i:=\bigcup_{n\in\N}F_{2n+1}$ et $\displaystyle F_{\infty}:=F\setminus(F_p\cup F_i)$.
	\end{itemize}
	Finalement, on construit notre bijection $\varphi$ comme suit.
	\begin{itemize}
		\item Si $x\in E_p$, on pose $\varphi(x):=f(x)$. $\varphi(x)$ possède un ancêtre de plus que $x$ (car c'est son "fils"), donc $\varphi(x)\in F_i$.
		\item Si $x\in E_i$, $x$ possède un antécédent $y\in F$ puisqu'il a un nombre impair d'ancêtres (et que tout nombre naturel impair est supérieur ou égal à $1$). De plus, cet antécédent $y$ est unique car $g$ est injective. On pose alors $\varphi(x):=y$. $y$ possède un antécédent de moins que $x$ (car c'est son "père"), donc $y\in F_p$.
		\item Si $x\in E_{\infty}$, on pose $\varphi(x):=f(x)$. $x$ possède une infinité d'ancêtres, donc $\varphi(x)$ aussi car c'est son "fils".
	\end{itemize}
	On vient de montrer que $\varphi$ envoyait $E_p$ sur $F_i$, $E_i$ sur $F_p$ et $E_\infty$ sur $F_\infty$. Comme ces ensembles sont disjoints, il suffit de montrer que $\varphi$ est bijectives \textit{séparément} de $E_p$ dans $F_i$, de $E_i$ dans $F_p$ et de $E_\infty$ dans $F_\infty$.
	\begin{itemize}
		\item Par construction, $\varphi$ est injective sur $E_p$ et $E_\infty$.
		\item Sur $E_i$, si on a $\varphi(x)=\varphi(x')$ alors $x$ et $x'$ ont le même antécédent par $g$, donc comme une application ne peut avoir qu'une seule image, on a $x=x'$.
		\item Par construction, $\varphi$ est surjective sur $F_p$ (l'antécédent de $y\in F_p$ est donné par $g(y)$, c'est un élément qui possède un nombre pair d'ancêtres donc il est bien dans $E_i$).
		\item Sur $F_i$, tout élément $y$ possède par définition au moins un antécédent $x$ par $f$ (puisqu'il possède un nombre impair d'ancêtres). L'élément $x$ sera alors dans $E_p$ et s'enverra bien sur $y$ par $\varphi$.
		\item Sur $F_\infty$, tout élément $y$ possède par définition au moins un antécédent $x$ par $f$ (puisqu'il possède un nombre infini d'ancêtres). L'élément $x$ possède alors une infinité d'ancêtres et est donc dans $E_\infty$ ; par conséquent, il s'enverra bien sur $y$ par $\varphi$.
	\end{itemize}
	Ainsi, la preuve est achevée.
\end{proof}
Plusieurs variantes du théorème de Cantor-Bernstein existent. Avant toute chose, démontrons le 
\begin{lem}
	Il existe une injection de $E$ dans $F$ si, et seulement si, il existe une surjection de $F$ dans $E$. 
\end{lem}
\begin{proof}
	S'il existe une injection $i$ de $E$ dans $F$, alors on construit une surjection $s$ de $F$ dans $E$ de la manière suivante : 
	\begin{itemize}
		\item Si $y\in F$ possède un antécédent $x$ par $i$, alors $x$ est unique car $i$ est injective et on pose $s(y):=x$.
		\item On choisit un "élément-poubelle" $x_0\in E$ pour tous les $y$ qui n'ont pas d'antécédent par $i$, et on pose $s(y):=x_0$.
	\end{itemize}
	Maintenant, s'il existe une surjection $s$ de $F$ dans $E$, on construit une injection $i$ de $E$ dans $F$ comme suit : tout élément $x\in E$ possède au moins un antécédent $y\in F$ par $s$, on en choisit un et on pose $i(x):=y$.
\end{proof}
\begin{rmq}[Axiome du choix]
	Dans notre construction, nous utilisons \textbf{l'axiome du choix}, qui garantit l'existence d'une manière de choisir un élément dans un ensemble quelconque (ce qu'on appelle une \textbf{fonction de choix}). Malgré son apparente simplicité, cette question n'a rien de trivial, et les logiciens ont même démontré qu'on pouvait faire des mathématiques consistantes en niant cet axiome.
\end{rmq}
\begin{cor}[Variante du théorème de Cantor-Bernstein]
	Soit $E$ et $F$ deux ensembles. Dans les quatre cas de figure suivants, il existe une bijection de $E$ dans $F$ :
	\begin{itemize}
		\item injection de $E$ dans $F$ / injection de $F$ dans $E$
		\item injection de $E$ dans $F$ / surjection de $E$ dans $F$
		\item surjection de $F$ dans $E$ / injection de $F$ dans $E$
		\item surjection de $F$ dans $E$ / surjection de $E$ dans $F$
	\end{itemize}
\end{cor}
\begin{proof}
	Il suffit d'utiliser à chaque fois le lemme précédent et le théorème de Cantor-Bernstein.
\end{proof}
\subsection{Axiomes de Zermelo-Fraenkel}
\textit{Cette partie est excessivement ardue et abstraite. Aussi est-il formellement déconseillé de s'y aventurer si l'on n'est pas parfaitement à l'aise avec le reste du programme.} \\ \par Peu à peu, les mathématiciens se sont rendus compte qu'ils ne pouvaient plus se contenter des définitions intuitives que nous avons données dans ce chapitre. Il leur fallait une théorie beaucoup plus rigoureuse. Il s'agit d'un travail de longue haleine initié par Cantor à la fin du XIX$^e$ siècle et continué par ses successeurs au XX$^e$. Parmi eux, Zermelo et Fraenkel ont laissé leur nom au système axiomatique le plus usuel aujourd'hui. Il est constitué de plusieurs axiomes, mais tous ne nous seront pas utiles. Nous présentons donc ici seulement ceux dont nous avons absolument besoin.\\ \par Cette théorie part d'un modèle appelé \textbf{univers}, dans lequel tous les objets sont des ensembles (on suppose qu'il en existe au moins un). On peut voir l'univers comme un nuage de points, chaque point désignant un ensemble. Entre certains ensembles de notre univers existent des flèches qui les relient. Elles symbolisent \textbf{l'appartenance}. L'écriture $a\in b$ signifie que "l'ensemble $a$ appartient à l'ensemble $b$". On dit aussi que $a$ est un \textbf{élément} de $b$. Aussi étrange que cela puisse paraître au premier abord, répétons qu'absolument \textit{tous} les objets de notre univers sont des ensembles. Ainsi, les éléments d'un ensemble ne peuvent être que des ensemble à leur tour.\\ \par Pourtant, lorsqu'on considère un ensemble comme $\left\{1,2,3\right\}$ par exemple, on n'a guère envie de considérer ses trois éléments comme des ensembles. Mais ce serait oublier comment les entiers naturels sont obtenus. Nous verrons plus loin dans l'année (cf. compléments du chapitre "Réels et suites") qu'ils peuvent être vus eux-mêmes comme des ensembles. C'est d'ailleurs ce qui a fait naître le slogan : "Dieu créa l'ensemble vide, les hommes firent le reste".\\ \par Enfin, notons que l'usage de $\forall$ dans les phrases quantifiées sera légèrement différent : d'habitude, un tel quantificateur est toujours suivi du symbole $\in$ (par exemple : $\forall x\in\R$, $\forall n\in\Z$, etc.). Ici, ce n'est pas la peine, puisque tous les ensembles appartiennent à un même univers (inutile donc le préciser).
\begin{dfn}[Inclusion]
	L'\textbf{inclusion} de deux ensembles est définie par : $$\forall a,\ \forall b,\ a\subset b \iff \left(\forall c,\ c\in a\implies c\in b\right)$$
\end{dfn}
\begin{ax}[Axiome d'extensionnalité]
	Si deux ensembles ont les mêmes éléments, alors ils sont égaux : $$\forall a,\ \forall b,\ \left(\forall x,\ x\in a\iff x\in b\right)\implies a=b$$
\end{ax}
\begin{ax}[Axiome de la paire]
	Étant donnés deux ensembles $a$ et $b$, il existe un ensemble qui a pour éléments $a$ et $b$, et eux seuls : $$\forall a,\forall b,\ \exists c,\ \left(\forall x,\ x\in c\iff \left(x\in a\ \ \text{ou}\ \ x\in b\right)\right)$$
	Notons que l'ensemble $c$ est unique d'après l'axiome d'extensionnalité. On le note $\left\{a,b\right\}$. Si $a\ne b$, on dit que c'est une \textbf{paire}. A l'inverse, pour tout $a$, on pose $\left\{a\right\}:=\left\{a,a\right\}$ et on dit que c'est un \textbf{singleton}.
\end{ax}
\begin{dfn}[Couple]
	Soit $a$ et $b$ des ensembles. Le \textbf{couple} $(a,b)$ est défini par $$(a,b):=\left\{\left\{a\right\},\left\{a,b\right\}\right\}$$.
\end{dfn}
\begin{lem}
	Si $\left\{x,y\right\}=\left\{x',y'\right\}$, alors on a $\displaystyle\begin{cases}
		x=x' \\ y=y'
	\end{cases}$ ou $\displaystyle\begin{cases}
	x=y' \\ y=x'
	\end{cases}$.
\end{lem}
\begin{proof}
	Dans un premier temps, supposons que $x=y$. Alors $x'\in\left\{x',y'\right\}=\left\{x\right\}$ donc $x'=x$. De même, $y'=x$. Finalement, on a $x=x'=y=y'$ d'où le résultat. \par A partir de maintenant, supposons que $x\ne y$ et $x'\ne y'$. On a $x\in\left\{x,y\right\}=\left\{x',y'\right\}$ donc $x=x'$ ou $x=y'$. De même, $y=x'$ ou $y=y'$. Maintenant : 
	\begin{itemize}
		\item Si $x=x'$, comme $x\ne y$, on ne peut pas avoir $y=x'$ donc $y=y'$.
		\item Si $x=y'$, le raisonnement est symétrique.
	\end{itemize}
	Ainsi, la preuve est achevée.
\end{proof}
\begin{cor}
	On a $(a,b)=(a',b')$ si, et seulement si, $\displaystyle\begin{cases}
		a=a' \\ b=b'
	\end{cases}$.
\end{cor}
\begin{proof}
	Le sens réciproque est trivial. Quant au sens direct, si $(a,b)=(a',b')$, alors $\left\{\left\{a\right\},\left\{a,b\right\}\right\}=\left\{\left\{a'\right\},\left\{a',b'\right\}\right\}$. Appliquons le lemme : 
	\begin{itemize}
		\item Si $\left\{a\right\}=\left\{a'\right\}$ et $\left\{a,b\right\}=\left\{a',b'\right\}$, la première égalité fournit $a=a'$. Puis la seconde aboutit à deux sous-cas : 
		\begin{itemize}
			\item[--] Si $a=a'$ et $b=b'$ nous avons terminé.
			\item[--] Si $a=b'$ et $b=a'$, alors $a=a'=b=b'$ et nous avons terminé.
		\end{itemize}
		\item Si $\left\{a\right\}=\left\{a',b'\right\}$ et $\left\{a'\right\}=\left\{a,b\right\}$, alors $a=a'=b=b'$ et nous avons terminé.
	\end{itemize}
	Ainsi, la preuve est achevée.
\end{proof}
\begin{ax}[Axiome de la réunion]
	Si on se donne un ensemble $a$, il existe un ensemble qui est la réunion des éléments de $a$ (rappelons que les éléments de $a$ sont aussi des ensembles) : $$\forall a,\ \exists b,\ \forall x,\ \left(x\in b\iff \exists c\in a,\ \left(c\in a \ \ \text{et}\ \ x\in c\right)\right)$$ On le notera $\displaystyle\bigcup a$. Il est unique d'après l'axiome d'extensionnalité.
\end{ax}
\begin{dfn}
	Si $a$ et $b$ sont deux ensembles, on pose $a\cup b:=\displaystyle\bigcup \left\{a,b\right\}$.
\end{dfn}
\begin{dfn}
	Si $n\geqslant3$, la notation $\left\{a_1,\ldots,a_n\right\}$ désigne en fait $\left\{a_1,\ldots,a_{n-1}\right\}\cup\left\{a_n\right\}$ (il s'agit donc là d'une définition par récurrence).
\end{dfn}
\begin{ax}[Axiome de l'ensemble des parties]
	Pour tout ensemble $a$, il existe un ensemble constitué des parties de $a$ : $$\forall a,\ \exists b,\ \forall x,\ \left(x\in b\iff x\subset a\right)$$ On le notera $\mathcal{P}(a)$. il est unique d'après l'axiome d'extensionnalité.
\end{ax}
\begin{ax}[Axiome de compréhension]
	Nous en donnons une version édulcorée, mais suffisante pour notre usage. Soit $P$ une proposition logique ne faisant pas intervenir le symbole $b$. Si on se donne un ensemble $a$, alors il existe un sous-ensemble de $a$ constitué des éléments qui vérifient la proposition $P$ : $$\forall a,\ \exists b,\ \forall x,\ \left(x\in b\iff \left(x\in a \ \ \text{et}\ \ P\right)\right)$$ Encore une fois, $b$ est unique par l'axiome d'extensionnalité. On le notera $\left\{x\in a\ | \ P\right\}$. Pour être tout à fait exact, il n'y a pas \textit{un} mais \textit{des} axiomes de compréhension (un par proposition $P$, ce qui fait beaucoup !). On parlera donc de plutôt de \textbf{schéma d'axiomes de compréhension}.
\end{ax}
Comme application, on peut construire le produit cartésien de deux ensembles. En effet si $a\in A$ et $b\in B$, alors $(a,b)=\left\{\left\{a\right\},\left\{a,b\right\}\right\}\in\mathcal{P}(\mathcal{P}(A\cup B))$.
\begin{dfn}[Produit cartésien]
	Si on se donne $A$ et $B$ deux ensembles, leur \textbf{produit cartésien} est défini par : $$A\times B:=\left\{z\in\mathcal{\mathcal{P}(A\cup B)}\ | \ \exists a,\ \exists b,\ \left(a\in A \ \ \text{et}\ \ b\in B\ \ \text{et}\ \ z=(a,b)\right)\right\}$$
\end{dfn}
Comme autre application (parfois introduite comme un axiome, mais ici cela s'en déduit de l'axiome de l'existence d'un ensemble), on a la :
\begin{dfn}[Ensemble vide]
	On sait qu'il existe au moins un ensemble $a$. On pose : $$\emptyset:=\left\{x\in a\ | \ x\ne x\right\}$$ Remarquons que $\emptyset$ ne peut contenir aucun élément. Notons aussi que $\emptyset$ ne dépend pas du choix de $a$. En effet, si on se donne un autre ensemble $b$, et qu'on définit $\emptyset':=\left\{x\in b\ | \ x\ne x\right\}$, montrons que $\emptyset=\emptyset'$. Par extensionnalité, il suffit de montrer $\forall x,\ x\in\emptyset\iff x\in\emptyset'$. or, les deux propositions sont systématiquement fausses, donc elles sont bien équivalentes. On peut alors nommer $\emptyset$ \textbf{l'ensemble vide}.
\end{dfn}
\begin{dfn}
	Si $x$ et $y$ sont deux ensembles, il existe un unique ensemble $z$ tel que $$\forall t,\ \left\{t\in z\iff\left(t\in a\ \ \text{et}\ \ t\in b\right)\right\}$$ On le note $x\cap y$.
\end{dfn}
\begin{proof}
	L'unicité vient de l'axiome d'extensionnalité. Pour l'existence, on pose par exemple $x\cap y:=\left\{t\in x\ | \ t\in y\right\}$.
\end{proof}
Citons un dernier axiome, qui ne fait pas à proprement parler partie des axiomes de Zermelo-Fraenkel : 
\begin{ax}[Axiome de la fondation]
	Dans tout ensemble $x$, il existe un élément $y$ qui est disjoint de $x$ : $$\forall x,\ \exists y,\ \left(y\in x \ \ \text{et}\ \ x\cap y=\emptyset\right)$$
\end{ax}
Comme conséquence, on note qu'aucun ensemble ne peut appartenir à lui-même (si $x\in x$, appliquer l'axiome de fondation à $\left\{x\right\}$ pour obtenir une contradiction). \\ \par Il manque un axiome important, dit \textbf{axiome de l'infini}, que nous reverrons en détail dans les compléments du chapitre "Réels et suites" dans le retour sur la construction de $\N$.



\newpage
\chapter{Introduction à l'algèbre}
\section{Relations}
\begin{dfn}[Relation binaire]
	Soit $E$ un ensemble. Une \textbf{relation} (binaire) $\mathcal{R}$ indique si deux éléments sont "liés". On dira que $x$ et $y$ \textbf{sont en relation}, et on écrira $x\mathcal{R}y$.
\end{dfn}
\begin{exa}
	La relation "$x$ et $y$ ont la même parité" sur $\N$, ou la relation "$x\leqslant y$" sur $\R$.
\end{exa}
\begin{rmq}
	Formellement, on peut définir $\mathcal{R}$ comme une partie de $E^2$ qui correspond aux couples qui sont en relations. mais il importe surtout de retenir deux choses :
	\begin{itemize}
		\item A part les relations usuelles, comme celles de l'exemple ci-dessus, on peut définir des relations abstraites. Dans ce cas, typiquement, on définira $\mathcal{R}$ sous la forme :$$\forall(x,y)\in E^2,\ x\mathcal{R}y \iff p(x,y)$$ où $p(x,y)$ est une assertion qui dépend de $x$ et $y$. On trouvera quelques exemples dans la feuille d'exercices.
		\item Pour montrer que deux relations $\mathcal{R}$ et $\mathcal{R}'$ sont égales, on montre que $$\forall(x,y)\in E^2,\ x\mathcal{R}y \iff x\mathcal{R}'y$$ 
	\end{itemize}
\end{rmq}
\begin{dfn}[Relation d'ordre]
	Une \textbf{relation d'ordre} vérifie les trois axiomes fondamentaux qui suivent :
	\begin{enumerate}
		\item Elle est \textbf{réflexive} : $\forall x\in E,\ x\mathcal{R}x$.
		\item Elle est \textbf{antisymétrique} : si $x\mathcal{R}y$ et $y\mathcal{R}x$, alors $x=y$.
		\item Elle est \textbf{transitive} : si $x\mathcal{R}y$ et $y\mathcal{R}z$, alors $x\mathcal{R}z$.
	\end{enumerate}
	Si la relation permet de  comparer deux éléments quelconques - autrement dit, si $$\forall (x,y)\in E^2,\ x\mathcal{R}y \ \text{ou} \ y\mathcal{R}x$$ elle est dite \textbf{totale}. Sinon, elle est dite \textbf{partielle}.
\end{dfn}
\begin{exa}
	La relation $\leqslant$ est une relation d'ordre totale sur $\N$. En revanche, la relation $\subset$ est une relation d'ordre partielle sur $\mathcal{P}(\N)$, car par exemple, l'ensemble des entiers pairs et l'ensemble des entiers impairs ne peuvent être comparés : aucun n'est inclus dans l'autre. On parle aussi d'ensemble \textbf{totalement ordonnée} ou d'ensemble \textbf{partiellement ordonné}.
\end{exa}
\begin{rmq}
	Si $E$ est muni de la relation d'ordre $\leqslant$, on définit les relations $<$, $\geqslant$ et $>$ à partir de $\leqslant$ de la manière suivante : $$\forall(x,y)\in E^2,\ \left \{ \begin{array}{lll} x<y \iff x\leqslant y \ \text{et} \ x\ne y \\ x\geqslant y\iff y\leqslant x \\ x>y \iff y\leqslant x \ \text{et} \ x\ne y\end{array}\right. $$ $<$ et $>$ ne sont pas des relations d'ordre, puisqu'elles ne sont pas réflexives (sauf si $E=\emptyset$). En revanche, on vérifie sans peine que $\geqslant$ est une relation d'ordre. Elle est d'ailleurs total si, et seulement si, $\leqslant$ est totale.
\end{rmq}
\begin{rmq}
	On a le résultat suivant : si $x_1\leqslant x_2\leqslant\ldots\leqslant x_n$, alors $x_1\leqslant x_n$, avec égalité si, et seulement si, tous les $x_i$ son égaux (récurrence sur $n\in\N\st$). En particulier, il suffit qu'une seule des inégalités $x_i\leqslant x_{i+1}$ soit stricte pour que $x_1<x_n$.
\end{rmq}
\begin{dfn}[Partie majorée, minorée, majorant, minorant]
	Soit $E$ un ensemble muni d'une relation d'ordre qu'on notera $\leqslant$. Alors une partie $A$ de $E$ est \textbf{majorée} quand elle possède un \textbf{majorant $M$} : $$\forall x\in A,\ x\leqslant M$$ Elle est \textbf{minorée} quand elle possède un \textbf{minorant $m$} : $$\forall x\in A,\ m\leqslant x$$
\end{dfn}
\begin{rmq}
	Bien que notée $\leqslant$, cette relation d'ordre pourrait très bien être l'inclusion. En pratique, ce seront les deux seuls cas que nous considérerons.
\end{rmq}
\begin{dfn}[Maximum, minimum]
	Il faut noter que $m$ ou $M$ n'appartient pas nécessairement à $A$. Si $m\in A$, on dit que c'est un \textbf{plus petit élément}, ou un \textbf{minimum}. Si $M\in A$, on dit que c'est un \textbf{plus grand élément}, ou \textbf{maximum}.
\end{dfn}
\begin{exa}
	$\R_+\st$ possède plusieurs minorants (par exemple $0$, ou encore $-1$), mais il ne possède pas de plus petit élément. En effet, s'il possédait un plus petit élément $x$, on aurait $x\in\R_+\st$ donc $x>0$, si bien que $\displaystyle \frac{x}{2}$ serait dans $\R_+\st$, tout en étant strictement inférieur à $x$, ce qui constituerait une contradiction.
\end{exa}
\begin{prop}[Unicité du maximum et du minimum]
	S'il existe, le maximum est unique. De même avec un minimum.
\end{prop}
\begin{thm}
	Soit $(E,\leqslant)$ un ensemble totalement ordonné. Alors toute partie finie non vide de $E$ admet un minimum et un maximum.
\end{thm}
\begin{rmq}
	Si $E$ totalement ordonné est de cardinal $n$, on montre par récurrence qu'il existe une unique bijection $\sigma:\llbracket1,n\rrbracket\rightarrow E$ strictement croissante. En posant $x_i=\sigma(i)\in E$, les éléments de $E$ sont égaux à : $$x_1<\ldots<x_n$$ On dit qu'on a \textbf{ordonné} les éléments de $E$.
\end{rmq}
\begin{dfn}
	Si $x_1,\ldots,x_n$ sont des éléments de $E$ totalement ordonné (avec $n\geqslant1$), le maximum et le minimum de l'ensemble fini $\left\{x_1,\ldots,x_n\right\}$ sont usuellement notés $\max(x_1,\ldots,x_n)$ et $\min(x_1,\ldots,x_n)$. On peut aussi les noter $\displaystyle \underset{1\leqslant i\leqslant n}{\max} x_i$ et $\displaystyle \underset{1\leqslant i\leqslant n}{\min} x_i$
\end{dfn}
\begin{prop}
	Le maximum et le minimum sont associatifs au sens suivant : $$\max\left(\max(x_1,\ldots,x_n), x_{n+1}\right)=\max(x_1,\ldots,x_{n+1})$$
	$$\min\left(\min(x_1,\ldots,x_n), x_{n+1}\right)=\min(x_1,\ldots,x_{n+1})$$
\end{prop}
\begin{dfn}[Élément maximal, élément minimal - HP]
	Supposons désormais que $(E,\leqslant)$ n'est pas nécessairement totalement ordonné. On dit qu'un élément $M$ est \textbf{maximal} lorsque pour tout $x\in E$, si $x$ et $M$ sont comparables, alors $x\leqslant M$. De manière équivalente, on peut écrire : $$\forall x\in E,\ x\geqslant M \implies x=M$$ ou même $$\forall x\in E,\ x\geqslant M \implies x\leqslant M$$ Encore autrement dit, on n'a jamais $x>M$. Tout maximum est un élément maximal, mais la réciproque est fausse. De plus, un élément maximal n'est pas nécessairement unique. De même, on dit que $m$ est \textbf{minimal} lorsque : $$\forall x\in E,\ x\leqslant m\implies x=m$$
\end{dfn}
\begin{dfn}[Relation d'équivalence]
	Une \textbf{relation d'équivalence} vérifie les trois axiomes fondamentaux qui suivent :
	\begin{enumerate}
		\item Elle est \textbf{réflexive} : $\forall x\in E,\ x\mathcal{R}x$.
		\item Elle est \textbf{symétrique} : si $x\mathcal{R}y$ alors $y\mathcal{R}x$.
		\item Elle est \textbf{transitive} : si $x\mathcal{R}y$ et $y\mathcal{R}z$, alors $x\mathcal{R}z$.
	\end{enumerate}
	Pour chaque élément $x$, on peut considérer l'ensemble de tous les éléments en relation avec $x$ (notons que c'est grâce à l'axiome de symétrie que la notion de "être en relation" a du sens). On l'appelle la \textbf{classe d'équivalence} de $x$ et on la note $\overline{x}$. Une classe d'équivalence (sans préciser de quel $x$) est une partie $A$ de $E$ telle que $\exists x\in E,\ A=\overline{x}$.
\end{dfn}
\begin{prop}
	Une classe d'équivalence peut être considérée comme la classe d'équivalence de n'importe lequel de ses éléments.
\end{prop}
\begin{dfn}[Ensemble-quotient, surjection canonique - HP]
	Si $\mathcal{R}$ est une relation d'équivalence sur $E$, l'ensemble des classes d'équivalences de $E$ est appelé l'\textbf{ensemble-quotient} de $E$ par $\mathcal{R}$. On le note souvent $E/\mathcal{R}$. L'application $x\mapsto\overline{x}$ est automatiquement surjective, on l'appelle la \textbf{surjection canonique} de $E$ dans $E/\mathcal{R}$.
\end{dfn}
\begin{dfn}[Partition]
	Une \textbf{partition} d'un ensemble $E$ est un ensemble $\mathcal{A}$ de parties de $E$ tel que  :
	\begin{itemize}
		\item Aucune partie n'est vide : $\forall A\in\mathcal{A},\ A\ne\emptyset$.
		\item Deux parties distinctes sont toujours disjointes : $\forall(A,B)\in\mathcal{A}^2,\ A\ne B\implies A\cap B=\emptyset$
		\item La réunion des parties vaut $E$ : $\displaystyle E=\bigcup_{A\in \mathcal{A}} A$
	\end{itemize}
	Les deux derniers points peuvent se réécrire $$E=\bigsqcup_{A\in\mathcal{A}}A$$ Ils signifient que tout élément de $E$ appartient à un et un seul $A$. Finalement, une partition est un recouvrement disjoint dont aucun élément n'est vide.
\end{dfn}
\begin{exa}
	$\N$ est partitionné entre entiers pairs et impairs. $\R$ est partitionné entre réels strictement positifs, réels strictement négatifs et $\left\{0\right\}$.
\end{exa}
\begin{thm}
	Soit $\mathcal{R}$ une relation d'équivalence sur un ensemble $E$. Alors, ses classes d'équivalence forment une partition de $E$.
\end{thm}
\begin{rmq}
	On peut montrer par analyse-synthèse que, réciproquement, pour toute partition $X$ de $E$, il existe une unique relation d'équivalence dont les classes sont exactement les éléments de $X$.
\end{rmq}
\section{Lois de composition}
Dans toute cette partie, nous nous donnons un ensemble de base $E$. Nous verrons dans les chapitres d'algèbre que $E$ peut par exemple être un groupe, un anneau, un corps, un espace vectoriel...
\begin{dfn}[Loi de composition interne, stabilité, structure algébrique]
	Une \textbf{loi de composition interne} sur $E$ est une application 
	$$
	\begin{array}[t]{ccccc}
		*&: & E^2 & \longrightarrow & E \\
		&& (x,y) & \longmapsto & x*y
	\end{array}
	$$ La loi est dite interne parce qu'on "reste dans $E$ en mélangeant deux éléments de $E$". on dit aussi que $E$ est \textbf{stable} par $*$. Le couple $(E,*)$ est appelé une \textbf{structure algébrique}.
\end{dfn}
\begin{rmq}
	Nous verrons ultérieurement qu'une structure algébrique peut comporter plusieurs lois.
\end{rmq}
\begin{exa}
	L'addition est une loi de composition interne sur $\R_+$. La soustraction est aussi une loi de composition, mais elle n'est pas interne.
\end{exa}
\begin{dfn}[Commutativité]
	La loi de composition est \textbf{commutative} lorsqu'on peut "permuter les facteurs" : $$\forall(x,y)\in E^2,\ x*y=y*x$$
\end{dfn}
\begin{exa}
	Sur $\R$, l'addition est commutative, mais pas la soustraction.
\end{exa}
\begin{dfn}[Associativité]
	La loi de composition est \textbf{associative} lorsqu'on peut "jouer avec les parenthèses" : $$\forall(x,y,z)\in E^3,\ (x*y)*z=x*(y*z)$$
\end{dfn}
\begin{exa}
	Sur $\N$, l'addition et la multiplication sont des lois de composition internes associatives. Sur $\Z$, la soustraction est une loi de composition interne, mais elle n'est pas associative.
\end{exa}
\begin{dfn}[Élément neutre]
	Un \textbf{élément neutre} $e$ pour la loi $*$ est un élément qui "ne change rien" : $$\forall x\in E,\ x*e=e*x=x$$
\end{dfn}
\begin{prop}
	S'il existe, l'élément neutre est unique. cela permet de parler de "l'élément neutre" sans ambiguïté.
\end{prop}
\begin{exa}
	Sur $\Z$, l'élément neutre de l'addition est $0$, celui de la multiplication est $1$. Sur $E^E$, l'élément neutre de la loi $\circ$ est $\id_E$.
\end{exa}
\begin{dfn}
	Soit $(E,*)$ une structure algébrique, $A$ et $B$ deux parties de $E$ et $x\in E$. On pose respectivement : $$A*x=\left\{a*x \ | \ a\in A\right\}$$ $$x*B=\left\{x*b\ | \ b\in B\right\}$$ $$A*B=\left\{a*b\ | \ (a,b)\in A\times B\right\}$$
\end{dfn}
\begin{exa}
	Dans $(\Z,+)$, on a $0+\Z=\Z+0=\Z$. Dans $(\Z,\times)$, l'ensemble des multiples de $n$ est noté $n\Z$.
\end{exa}
\begin{exa}
	Si $*$ est associative, on a $(A*B)*C=A*(B*C)$ (chaîne d'équivalences). Si $*$ est commutative, on a $A*B=B*A$ (chaîne d'équivalences).
\end{exa}
A partir de maintenant, on suppose que $E$ est muni d'une loi interne associative $*$ et d'un élément neutre $e$.
\begin{dfn}[Inversibilité]
	On dit qu'un élément $x\in E$ est inversible lorsqu'il "possède un inverse" : $$\exists y\in E,\ x*y=y*x=e$$
\end{dfn}
\begin{prop}
	S'il existe, l'inverse de $x$ est unique. Cela permet de parler de "l'inverse de $x$" sans ambiguïté. On parle aussi du \textbf{symétrique} de $x$.
\end{prop}
\begin{rmq}
	A propos de la notation :
	\begin{itemize}
		\item Si la loi est notée $+$, le neutre sera souvent noté $0$ et l'inverse de $x$ sera souvent noté $-x$ : dans ce cas, $x-y$ désignera en fait $x+(-y)$.
		\item Dans le cas général, l'inverse de $x$ sera souvent noté $x\inv$.
	\end{itemize}
\end{rmq}
\begin{exa}
	$e$ est toujours inversible, d'inverse lui-même.
\end{exa}
\begin{cor}
	Si $x$ est inversible, alors l'inverse de $x\inv$ est $x$. Autrement dit, l'application de passage à l'inverse est involutive.
\end{cor}
\begin{cor}
	Soit $x$ et $y$ inversibles. Alors $x*y$ est inversible, d'inverse $y\inv *x\inv$.
\end{cor}
\begin{rmq}
	Attention, il ne faut alors pas oublier d'inverse l'ordre de $x$ et $y$, sinon cela devient faux dès que la loi n'est pas commutative.
\end{rmq}
\section{Construction des entiers, HP}
\subsection{Construction de $\N$}
Dans cette partie, nous construisons $\N$, ce qui constitue véritablement la pierre de fondation des mathématiques, au moins au niveau des classes préparatoires. Nous allons voir qu'à partir de trois axiomes, qui sont des "axiomes de bon sens" compte tenu de la perception pratique que nous avons de $\N$, on peut formaliser et démontrer quasiment la totalité du programme.
\begin{dfn}[Axiomes fondateurs de $\N$]
	On pose l'existence d'un ensemble $\N$ non vide, munit d'une relation d'ordre $\leqslant$, tel que :
	\begin{itemize}
		\item Toute partie non vide possède un plus petit élément.
		\item Toute partie non vide et majorée possède un plus grand élément.
		\item $\N$ n'admet pas de plus grand élément.
	\end{itemize}
\end{dfn}
\begin{exa}
	Soit $(u_n)_{n\in\N}$ une suite à valeurs dans $\N$. En considérant l'ensemble $$u(\N)=\left\{u_n\ | \ n\in\N\right\}$$ on pourra montrer que 
	\begin{itemize}
		\item si $u$ est décroissante, alors elle est stationnaire ;
		\item si $u$ est croissante et majorée, alors elle est stationnaire.
	\end{itemize}
\end{exa}
\begin{prop}
	La relation d'ordre $\leqslant$ est totale.
\end{prop}
\begin{proof}
	Pour comparer $m$ et $n$, il suffit de considérer l'ensemble $\left\{m,n\right\}$, qui est un ensemble non vide.
\end{proof}
\begin{dfn}
	On pose $\llbracket n_1,n_2\rrbracket=\left\{n\in\N\ | \ n_1\leqslant n\leqslant n_2\right\}$. En particulier, si $n_1>n_2$, on a $\llbracket n_1,n_2\rrbracket=\emptyset$.
\end{dfn}
En un certain sens, le lemme suivant formalise ce que fait un enfant lorsqu'il compte.
\begin{lem}
	Soit $n\in\N$. Alors $n$ possède un unique successeur, c'est-à-dire un entier $n^+>n$ tel que $\rrbracket n,n^+\llbracket=\emptyset$.
\end{lem}
\begin{proof}
	Considérer l'ensemble $A=\left\{k\in\N\ | \ k>n\right\}$ qui est une partie non vide de $\N$ et prendre son minimum. Pour montrer que ce minimum convient, raisonner par l'absurde pour l'ensemble discret. Pour l'unicité, prendre un autre candidat et les ordonner SPG. Conclure par un mini-RPA lié à l'ensemble vide précédemment cité.
\end{proof}
\begin{dfn}
	On note $0$ le plus petit élément de $\N$, et $1$ son successeur. On définit l'ensemble $\N\st=\N\setminus\left\{0\right\}$. On introduit également : $2=1^+$, $3=2^+$, $4=3^+$, $5=4^+$, $6=5^+$, $7=6^+$, $8=7^+$ et $9=8^+$.
\end{dfn}
\begin{thm}[Raisonnement par récurrence simple]
	Soit $P_n$ une propriété indexée sur $\N$. On suppose que \begin{itemize}
		\item Initialisation : $P_0$ est vraie
		\item Hérédité : $\forall n\in\N,\ P_n\implies P_{n+1}$
	\end{itemize}
	Alors, pour tout $n\in\N$, $P_n$ est vraie.
\end{thm}
\begin{proof}
	Considérer $A=\left\{n\in\N\ | \ P_n \ \text{est fausse}\right\}$ et raisonner par l'absurde en supposant que $A$ est non vide. Considérer alors son minimum, qui est supérieur ou égal à $1$. Le minimum moins $1$ n'est donc pas dans $A$, mais alors par hérédité le minimum n'est pas non plus dans $A$. Absurde.
\end{proof}
Le théorème qui suit est fondamental, il est la pierre de fondation d'un nombre incalculable de constructions mathématiques.
\begin{thm}[Définition d'une suite par récurrence, admis]
	Soit $E$ un ensemble, $x\in E$ et $f:E\rightarrow E$. Alors il existe une unique suite $(u_n)_{n\in\N}$ telle que $$\left \{ \begin{array}{ll}u_0=x \\ \forall n\in\N,\ u_{n^+}=f(u_n) \end{array} \right.$$
\end{thm}
Comme application, nous allons à présent définir une addition et une multiplication sur $\N$. Nous verrons par la suite comment les prolonger sur $\Z$, $\Q$ et enfin $\R$. On prolongera finalement à $\C$.
\begin{dfn}[Définition de l'addition]
	Pour $m\in\N$, on définit $m+n$ par récurrence sur $n$ : $$\left \{ \begin{array}{ll}
		m+0=m \\ \forall n\in\N,\ m+n^+=(m+n)^+
	\end{array} \right.$$
\end{dfn}
\begin{rmq}[Changement de notation]
	Notons que pour tout $n\in\N$, on a $$n^+=(n+0)^+=n+0^+=n+1$$ A partir de maintenant, on abandonnera la notation $n^+$ au profit de la notation $n+1$. Typiquement, le principe de récurrence sera écrit sous la forme $P_n\implies P_{n+1}$, et la construction d'une suite par récurrence s'écrira $u_{n+1}=f(u_n)$.
\end{rmq}
La petite proposition qui suit est d'une importance capitale dans nombre d'inégalités, en particulier en analyse.
\begin{prop}[Théorème de Roger, ou théorème de la licorne, NS]
	Si $m<n+1$, alors $m\leqslant n$.
\end{prop}
\begin{proof}
	Si on avait $m>n$, alors $m$ viendrait s'intercaler entre $n$ et $n^+=n+1$.
\end{proof}
On démontre sans trop de difficulté que l'addition est associative et commutative, et qu'elle est compatible avec $\leqslant$ au sens suivant.
\begin{proof}
	Pour tout $(m,n,p)\in\N^3$, on a : $m\leqslant n \implies m+p\leqslant n+p$.
\end{proof}
Voici maintenant la définition de la multiplication par récurrence.
\begin{dfn}[Définition de la multiplication]
	Pour $m\in\N$ fixé, on définit $m\times n$ par récurrence sur $n$ :$$\left \{ \begin{array}{ll}
		m\times0=0 \\ \forall n\in\N,\ m\times (n+1)=(m\times n)+m
	\end{array} \right.$$
\end{dfn}
Là aussi, on démontre assez facilement que la multiplication est associative, commutative et distributive sur l'addition. On montre également sa compatibilité avec $\leqslant$ au sens suivant.
\begin{prop}
	On a : $\forall (m,n,p)\in\N^3,\ m\leqslant n\implies m\times p\leqslant n\times p$
\end{prop}
\subsection{Extension à $\Z$}
Jusqu'ici, nous avons construit les entiers naturels, leur addition et leur multiplication. Mais nous savons bien que dans $\N$, aucun entier non nul n'est inversible pour l'addition. D'où l'idée de construire une structure plus riche, contenant $\N$, dans laquelle on puisse inverse ces entiers pour obtenir ce qu'on appelle des "entiers négatifs". \par Une telle construction est proposée en annexe. Donnons sans démonstration les principaux résultats qui en découlent.
\begin{itemize}
	\item Il existe un ensemble $\Z$ contenant $\N$, dans lequel on peut prolonger l'addition et la multiplication définies sur $\N$. Ces deux lois gardent leurs propriétés (associativité, commutativité, distributivité). De plus, tout élément de $\Z$ est inversible pour l'addition.
	\item La relation d'ordre $\leqslant$ définie peut être prolongée sur $\Z$. Elle reste totale, et elle reste compatible avec l'addition : $$\forall (m,n,p)\in\Z^3,\ m\leqslant n\implies m+p\leqslant n+p$$ Cela montre que dans une addition, on peut "faire changer les termes de côté" de part et d'autres du signe $\leqslant$ ou du signe $\geqslant$ à condition d'inverser le signe. Elle est également compatible avec la multiplication au sens suivant : $$\forall(m,n,p)\in\Z^3,\ \left\{\begin{array}{ll}m\leqslant n \\ p\geqslant 0\end{array}\right. \implies m\times p\leqslant n\times p$$
	\item Tout élément de $\Z$ s'écrit d'une unique manière $0$, $n$ ou $-n$ avec $n\in\N\st$. Cette écriture sous la forme $\pm n$ s'appelle usuellement un \textbf{entier relatif}. C'est celle que nous garderons systématiquement à partir de maintenant.
	\item Enfin, on peut généraliser le résultat sur les inégalités strictes d'entiers : $$\forall(m,n)\in\Z^2,\ m<n+1\implies m\leqslant n$$
\end{itemize}
\section{Sommes et produits}
\subsection{Retour sur les formules sommatoires}
\begin{dfn}
	Soit $(x_k)_{k\in\N}$ une famille de complexes, et $m\in\N$ fixé. On définit $\displaystyle \sum_{k=m}^{n}x_k$ et $\displaystyle \prod_{k=m}^{n}x_k$ par récurrence sur $n\geqslant m$.
	\begin{itemize}
		\item Si $n=m$, on les définit comme étant égal à $x_k$. 
		\item Et pour tout $n\geqslant m$, on pose $\displaystyle \sum_{k=m}^{n+1}x_k=\left(\sum_{k=m}^{n}x_k\right)+x_{n+1}$ et $\displaystyle \prod_{k=m}^{n+1}x_k=\left(\prod_{k=m}^{n}x_k\right)\times x_{n+1}$
	\end{itemize}
	Si $n=m-1$, on pose $\displaystyle \sum_{k=m}^{n}x_k=0$ et $\displaystyle \prod_{k=m}^{n}x_k=1$ par convention. Et si $n<m-1$, on évite de définir ces grandeurs.
\end{dfn}
\begin{rmq}
	Par récurrence, on peut démontrer les relations admises dans le chapitre sur les complexes : séparation et regroupement de termes, distributivité, relation de Chasles. A chaque fois, il s'agit d'une récurrence rapide.
\end{rmq}
\begin{exa}
	Soit $n\in\N$. Par une récurrence immédiate, on a $\displaystyle \sum_{k=1}^{n}1=n$ et $\displaystyle \prod_{k=1}^{n}1=1$. Soit $\lambda\in\C$. Par distributivité, on a $\displaystyle \sum_{k=1}^{n}\lambda=\lambda n$ et $\displaystyle \prod_{k=1}^{n}\lambda=\lambda^n$.
\end{exa}
Grâce à la commutativité, on se dit intuitivement que nous pouvons sommer les $c_k$ dans l'ordre que nous le souhaitons. Cela permet en particulier d'indexer les éléments sur un autre ensemble que $\llbracket1,n\rrbracket$, et de donner un sens à des formules comme $\displaystyle \sum_{\substack{k=1 \\ k \ \text{pair}}}^{n}x_k$ ou $\displaystyle \sum_{1\leqslant i<j\leqslant n}x_{i,j}$.
\begin{dfn}
	Soit $I$ un ensemble \textbf{fini} d'indices, et $(x_i)_{i\in I}$ une famille de complexes indexée sur $I$. Puisque $I$ est fini, il existe une bijection $\sigma$ d'un certain $\llbracket1,n\rrbracket$ sur $I$. On pose alors : $$\sum_{i\in I}x_i=\sum_{k=1}^{n}x_{\sigma(k)} \ \ \text{et} \ \ \prod_{i\in I}x_i=\prod_{k=1}^{n}x_{\sigma(k)}$$
\end{dfn}
\begin{rmq}
	On admet que le résultat final reste le même si on change de choix pour $\sigma$.
\end{rmq}
\begin{exa}
	Voici quelques cas très simples.
	\begin{itemize}
		\item Si $I=\emptyset$, on a $\displaystyle \sum_{i\in I}x_i=0$ et $\displaystyle \prod_{i\in I}x_i=1$.
		\item Si $I=\left\{a\right\}$, on a $\displaystyle \sum_{i\in I}x_i=x_a$ et $\displaystyle \prod_{i\in I}x_i=x_a$.
		\item Si $I=\left\{a,b\right\}$ avec $a\ne b$, on a $\displaystyle \sum_{i\in I}x_i=x_a+x_b$ et $\displaystyle \prod_{i\in I}x_i=x_a x_b$.
		\item Si $I=\llbracket1,n\rrbracket$, on a $\displaystyle \sum_{i\in I}x_i=\sum_{k=1}^{n}x_k$ et $\displaystyle \prod_{i\in I}x_i=\prod_{k=1}^{n}x_k$.
		\item Si la famille $(x_i)_{i\in I}$ est constante égale à $\lambda$, on a $\displaystyle \sum_{i\in I}x_i=\lambda\card(I)$ et $\displaystyle \prod_{i\in I}x_i=\lambda^{\card(I)}$.
	\end{itemize}
\end{exa}
\begin{thm}
	Soit $I$ est $J$ deux ensembles finis, $\varphi:J\rightarrow I$ une bijection, et $(x_i)_{i\in I}$ une famille de complexes. Alors, on a : $$\sum_{i\in I}x_i=\sum_{j\in J} x_{\varphi(j)} \ \ \text{et} \ \ \prod_{i\in I}x_i=\prod_{j\in J} x_{\varphi(j)}$$ Autrement dit, on "interprète $i$ comme un $\varphi(j)$".
\end{thm}
Les formules de séparation et regroupement de termes ainsi que de distributivité restent évidemment vraies. L'analogue de la relation de Chasles serait le 
\begin{thm}[Associativité de la somme et du produit]
	Si $I_1,\ldots,I_n$ sont des parties deux à deux disjointes de $I$, alors on a : $$\sum_{i\in I_1\cup\ldots\cup I_n}=\sum_{i\in I_1}x_i +\ldots+\sum_{i\in I_n}x_i \ \ \ \text{et} \ \ \ \prod_{i\in I_1\cup\ldots\cup I_n}=\left(\prod_{i\in I_1}x_i\right)\times\ldots\times\left(\prod_{i\in I_n}x_i\right)$$
\end{thm}
\begin{proof}
	Une démonstration se trouve dans les compléments.
\end{proof}
Un particulier est celui des sommes doubles.
\begin{cor}[Théorème de Fubini]
	Soit $\displaystyle\left(x_{i,j}\right)_{\substack{1\leqslant i\leqslant m \\ 1\leqslant j\leqslant n}}$ une famille de complexes doublement indexée. Alors : $$\sum_{\substack{1\leqslant i\leqslant m \\ 1\leqslant j\leqslant n}}x_{i,j}=\sum_{i=1}^{m}\left(\sum_{j=1}^{n}x_{i,j}\right)=\sum_{j=1}^{n}\left(\sum_{i=1}^{m}x_{i,j}\right)$$ et $$\prod_{\substack{1\leqslant i\leqslant m \\ 1\leqslant j\leqslant n}}x_{i,j}=\prod_{i=1}^{m}\left(\prod_{j=1}^{n}x_{i,j}\right)=\prod_{j=1}^{n}\left(\prod_{i=1}^{m}x_{i,j}\right)$$
\end{cor}
\begin{rmq}
	On retiendra que dans une somme double de termes qui commutent (toujours le cas dans $\C$ donc), on peut sommer les termes dans l'ordre que l'on veut. Avec le même raisonnement, on généralise immédiatement à des sommes indexées sur des ensembles finis quelconques. Idem avec le produit.
\end{rmq}
\begin{exa}
	On veut calculer $\displaystyle \sum_{1\leqslant i,j\leqslant n}\min(i,j)$. On peut choisir de sommer d'abord selon $i$ ou d'abord selon $j$. Commençons par exemple avec $i$. On obtient : \begin{align*}
		\sum_{1\leqslant i,j\leqslant n}\min(i,j) &= \sum_{i=1}^{n}\left(\sum_{j=1}^{n}\min(i,j)\right) \\
		&=\sum_{i=1}^{n}\left(\sum_{j=1}^{i}j+\sum_{j=i+1}^{n}i\right) \\
		&=\sum_{i=1}^{n}\left(\frac{i(i+1)}{2}+(n-i)i\right) \\
		&=-\frac{1}{2}\sum_{i=1}^{n}i^2+\left(n+\frac12\right)\sum_{i=1}^{n}i \\
		&= -\frac{n(n+1)(2n+1)}{6}+\frac{(2n+1)n(n+1)}{4} \\
		&= \frac{n(n+1)(2n+1)}{6}
	\end{align*}
\end{exa}
De manière imagée, on vient de sommer les termes $x_{i,j}$ pour $(i,j)$ variant dans un carré de taille $n\times n$, et le résultat est le même que l'on somme par les lignes ou par les colonnes. Le théorème de Fubini donne un résultat légèrement plus général, lorsque $(i,j)$ parcourt un rectangle de taille $m\times n$. Mais l'associativité de la somme et du produit permet de traiter des cas plus complexes, comme celui où $(i,j)$ varie dans un triangle. \par Dans de tels cas, on peut commencer par fixer l'une des deux variables, mettons $i$, puis observer dans quelle plage (en fonction de $i$) varie $j$. Il est plus que recommandé de représenter graphiquement sur un plan le domaine dans lequel varie $(i,j)$.
\begin{exa}
	On souhaite calculer $\displaystyle \sum_{1\leqslant i\leqslant j\leqslant n}ij$. On a :
	\begin{align*}
		\sum_{1\leqslant i\leqslant j\leqslant n}ij &= \sum_{i=1}^{n}\left(i\sum_{j=i}^{n}j\right) \\
		&= \sum_{i=1}^{n}\left[i\left(\frac{n(n+1)}{2}-\frac{(i-1)i}{2}\right)\right] \\
		&=\frac{n(n+1)}{2}\sum_{i=1}^{n}i-\frac{1}{2}\sum_{i=1}^{n}i^3+\frac{1}{2}\sum_{i=1}^{n}i^2 \\
		&=\left(\frac{n(n+1)}{2}\right)^2-\frac12\left(\frac{n(n+1)}{2}\right)^2+\frac{n(n+1)(2n+1)}{12} \\
		&= \frac{n(n+1)}{2}\left[\frac{n(n+1)}{4}+\frac{2n+1}{6}\right] \\
		&= \frac{n(n+1)}{2}\times\frac{3n^2+7n+2}{12}\\
		&= \frac{n(n+1)(n+1)(3n+1)}{24}		
	\end{align*}
\end{exa}
\subsection{Itération d'une loi de composition interne}
Dans toute cette partie, on fixe une structure algébrique $(E,*)$ associative, admettant un élément neutre $e$.
\begin{dfn}
	Soit $x\in E$. On définit $x^n$ par récurrence de la manière suivante : $$\left \{ \begin{array}{ll}
		x^0=e \\ \forall n\in\N,\ x^{n+1}=x^n*x
	\end{array}\right.$$
\end{dfn}
\begin{rmq}
	Plus généralement, on définira par récurrence toute expression du type $$\prod_{k=1}^{n}x_k=x_1*\ldots*x_n$$
	\begin{itemize}
		\item En théorie, il s'agit donc de $$\left(\left(\left(x_1*x_2\right)*x_3\right)*\ldots*x_{n-1}\right)*x_n$$ Mais en fait, puisque $*$ est associative, on peut placer les parenthèses "où l'on veut" ! $$x_1*\ldots*x_n=x_1*\ldots*x_k*(x_{k+1}*\ldots*x_m)*x_{m+1}*\ldots*x_n$$ Cela se démontre rigoureusement par récurrence.
		\item Autre exemple : donnons-nous $A_1,\ldots,A_n\subset E$. On peut alors considérer $$A_1*\ldots*A_n$$ au sens de la structure algébrique $(\mathcal{P}(E),*)$. Une récurrence immédiate montre que $$A_1*\ldots*A_n=\left\{x_1*\ldots*x_n\ | \ (x_1,\ldots,x_n)\in A_1\times\ldots\times A_n \right\}$$
	\end{itemize}
\end{rmq}
\begin{rmq}
	Si de plus la loi est commutative, on peut définir des expressions comme $\displaystyle \prod_{i\in I}x_i$ avec $I$ fini, et l'ensemble des résultats précédents restent vrais. Par exemple, on pourra penser à l'union ou à l'intersection d'un	sur les parties d'un ensemble.
\end{rmq}
\begin{prop}
	Soit $x\in E$. On a : $$\forall (m,n)\in\N^2,\ \left \{ \begin{array}{ll} x^m*x^n=x^{m+n} \\ (x^m)^n=x^{mn}\end{array}\right.$$
\end{prop}
\begin{proof}
	Fixer $m$ et raisonner par récurrence sur $n$.
\end{proof}
\begin{rmq}
	En revanche, attention : si $x$ et $y$ ne commutent pas, on n'a pas nécessairement $(x*y)^n=x^n*y^n$.
\end{rmq}
\begin{prop}
	Soit $(x,y)\in E^2$. Si $x*y=y*x$, alors : $$\left\{\begin{array}{ll}\forall(m,n)\in\N^2,\ x^m*y^n=y^n*x^m \\ \forall n\in\N,\ (x*y)^n=x^n*y^n \end{array}\right.$$
\end{prop}
\begin{proof}
	Pour la première, fixer $m$ et raisonner par récurrence sur $n$. Pour la seconde, raisonner par récurrence sur $n$.
\end{proof}
On généralise désormais l'inverse d'un produit d'éléments inversibles.
\begin{prop}
	Soit $x_1,\ldots,x_n$ inversibles. Alors $x_1*\ldots*x_n*$ est inversible, d'inverse $$x_n\inv*\ldots x_1\inv$$
\end{prop}
\begin{proof}
	C'est une récurrence immédiate.
\end{proof}
\begin{exa}
	Soit $x\in E$ inversible. Alors, pour tout $n\in\N\st$, $x^n$ est inversible, d'inverse $(x\inv)^n$.
\end{exa}
\begin{dfn}
	Soit $x\in E$ inversible et $n\in\N\st$. On pose $$x^{-n}=(x^n)\inv=(x\inv)^n$$
\end{dfn}
\begin{rmq}
	Les relations précédentes continuent à être vraies pour des exposants dans $\Z$. Ainsi, si $x$ est inversible, on a : $$\forall (m,n)\in\Z^2,\ \left \{ \begin{array}{ll} x^m*x^n=x^{m+n} \\ (x^m)^n=x^{mn}\end{array}\right.$$ Si $y$ est inversible et $x*y=y*x$, on a de plus : $$\left\{\begin{array}{ll}\forall(m,n)\in\Z^2,\ x^m*y^n=y^n*x^m \\ \forall n\in\Z,\ (x*y)^n=x^n*y^n \end{array}\right.$$ Les vérifications sont fastidieuses, mais faciles. Il suffit à chaque fois de revenir à la définition. 
\end{rmq}
\section{Deux théorèmes d'arithmétique}
\begin{thm}
	Il existe une \textbf{division euclidienne} sur $\Z$ : $$\forall (a,b)\in\Z\times\N\st,\ \exists!(q,r)\in\Z\times\llbracket0,b-1\rrbracket,\ a=bq+r$$ La division euclidienne sur $\Z$ est d'une utilité considérable en arithmétique et en théorie des groupes.
\end{thm}
\begin{proof}
	Pour l'existence, fixons $a$ et $b$ comme dans l'énoncé et considérons $$A=\left\{a-bq\ | \ q\in\Z\right\}\cap\N$$ Puisque $b\geqslant1$, on a $\forall q \leqslant0,\ a-bq\geqslant a-q$. Donc pour $q=\min(0,a)$, on voit que $A$ est non vide. Posons alors $r=\min(A)$ et soit $q$ associé. On a bien $a=bq+r$ et $r\geqslant0$. Par ailleurs, on a $r<b$ sans quoi on pourrait écrire $$r-b=a-b(q+1)\in A$$ ce qui contredirait la minimalité de $r$. Donc $r\leqslant b-1$. \par Pour l'unicité, si deux couples $(q,r)$ et $(q',r')$ conviennent, supposons que $q>q'$. Alors on aurait $q\geqslant q'+1$, puis $r'\geqslant r'-r=b(q-q')\geqslant b$, ce qui est absurde. Donc $q\leqslant q'$. Symétriquement, on prouve que $q'\leqslant q$, d'où $q=q'$ par antisymétrie, puis $r=r'$ en réinjectant dans l'égalité $a=a$ avec les deux divisions euclidiennes.
\end{proof}
\begin{thm}
	Soit $b$ un entier supérieur ou égal à $2$. Alors tout entier $n>0$ possède une unique \textbf{décomposition en base $b$} : $$n=\overline{a_k a_{k-1}\ldots a_1 a_0}^b=\sum_{i=0}^{k}a_ib^i$$ où $k\geqslant0$, $\forall i\in\llbracket0,k\rrbracket,\ 0\leqslant a_i<b$ et $a_k\ne0$.
\end{thm}
\begin{proof}
	Il faut raisonner par récurrence et pour l'existence, et pour l'unicité (et j'ai la flemme de le taper parce que c'est pas beau).
\end{proof}
\begin{exa}
	$2014$ s'écrit $\overline{11111011110}^2$ en base $2$, aussi appelée \textbf{numération binaire}.
\end{exa}
\begin{exa}
	Très provisoirement, notons à nouveau $9^+$ le successeur de $9$. En base $9^+$, on a donc $$\overline{10}^{9^+}=1\times9^++0\times 1 = 9^+$$ Sauf mention expresse du contraire, tous les nombre seront écrits en base $9^+$. On pourra donc omettre la barre horizontale au-dessus de $\overline{10}$, et écrire $9^+=10$ (ce qui ne serait pas nécessairement vrai dans une autre base : examiner le cas d'une base strictement supérieure à $10$ par exemple). On parlera de \textbf{base 10} ou de \textbf{base décimale}.	
\end{exa}
\section{Compléments, HP}
\subsection{Démonstration du théorème d'existence / unicité des suites définies par récurrence avec une fonction}

\begin{thm}
	Soit $E$ un ensemble, $x \in E$ fixé et $f : E \rightarrow E$. Alors il existe une unique suite $(u_n)_{n \in \mathbb{N}}$ telle que $$\begin{cases}
		u_0 = x \\
		\forall n \in \mathbb{N}, u_{n^+} = f(u_n)
	\end{cases}$$
\end{thm}

\begin{proof}
	Pour l'unicité, on suppose que deux suites $(u_n)_{n \in \mathbb{N}}$ et $(v_n)_{n \in \mathbb{N}}$ conviennent, et on démontre par récurrence sur $n$ que $u_n = v_n$. Passons à l'existence, qui est loin d'être aussi évidente qu'il n'y paraît.
	\begin{itemize}
		\item Premier point : pour $n \in \mathbb{N}$, on commence par définir l'ensemble $E_n$ (très grand !) des suites $(v_k)_{k \in \mathbb{N}}$ telles que $$\begin{cases}
			v_0 = x \\
			\forall k < n, v_{k^+} = f(v_k)
		\end{cases}$$
		Autrement dit, $E_n$ est l'ensemble des suites telles que "les termes de 0 à $n$ conviennent". En particulier, $E_0$ est l'ensemble des suites telles que $v_0 = x$.
		\item Deuxième point : aucun $E_n$ n'est vide. Prouvons le résultat par récurrence sur $n$.
		\begin{itemize}
			\item Initialisation : $E_0 \neq \emptyset$. En effet, la suite constante égale à $x$ convient.
			\item Hérédité : si $E_n \neq \emptyset$, considérons une suite $(v_k)_{k \in \mathbb{N}}$ dans cet ensemble. En définissant la suite $(v_k')_{k \in \mathbb{N}}$ par $$\begin{cases}
				\forall k \neq n^+, v_k' = v_k \\
				v_{n^+}' = f(v_n)
			\end{cases}$$
			on obtient alors un élément de $E_{n^+}$.
		\end{itemize}
		\item Troisième point : si $(v_k)_{k \in \mathbb{N}}$ et $(v_k')_{k \in \mathbb{N}}$ appartiennent à $E_n$, alors $v_n = v_n'$. Il s'agit d'une récurrence facile sur $n$.
		\item Dernier point : on construit alors notre suite $(u_n)_{n \in \mathbb{N}}$ comme suit. Pour $n \in \mathbb{N}$ fixé, on choisit $(v_k)_{k \in \mathbb{N}}$ dans $E_n$ (ce qui est possible puisque $E_n \neq \emptyset$) puis on pose $u_n = v_n$. Le point précédent montre que le résultat ne dépend pas du choix de $(v_k)$. Ainsi, $(u_n)$ est bien définie. Montrons maintenant qu'elle satisfait les conditions demandées.
		\begin{itemize}
			\item[--] Déjà, par construction on a bien $u_0 = x$.
			\item[--] Puis si $n > 0$, montrons que $f(u_n) = u_{n^+}$. Soit $(v_k)_{k \in \mathbb{N}} \in E_{n^+}$ quelconque. Par construction on a que $E_{n^+} \subset E_n$. Ainsi, pour définir $u_n$ et $u_{n^+} = v_{n+1}$, il est légitime de choisir $(v_k)_{k \in \mathbb{N}} \in E_n$, et on obtient $u_n = v_n$. Enfin, par définition de $E_{n^+}$ on a $v_{n^+} = f(v_n)$ d'où $u_{n^+} = f(u_n)$.
		\end{itemize}
	\end{itemize}
	Ainsi la preuve est achevée.
\end{proof}
À présent, précisons un peu la notion de suite définie par récurrence forte. Pour plus de simplicité dans les notations, nous indexerons nos suites sur $\mathbb{N}^*$. Intuitivement, il s'agit d'initialiser une suite $(u_n)_{n \in \mathbb{N}^*}$ à $u_1 = x \in E$, puis pour $n \in \mathbb{N}^*$ de définir $u_{n+1}$ en fonction de $u_1, ..., u_n$. De manière formelle, il faudrait pouvoir écrire $u_{n+1} = f(u_1, ..., u_n)$ pour une application $f$ bien choisie, dont le nombre d'arguments serait variable. Tout cela a un sens, car on rappelle que $\displaystyle F = \bigcup_{n \in \mathbb{N}^*} E^n$ est bien défini (cf. l'exemple sur les alphabets du chapitre "Ensembles et applications"). \par 
De plus on pourra parler de la taille d'un élément de $F$ : il s'agit de l'unique entier $n \in \mathbb{N}^*$ tel que $\varphi \in E^n$. Dans le reste du paragraphe, on le notera $T(\varphi)$. \\ \par
Revenons à notre problème: on définit une application $f: F \rightarrow E$, puis on cherche à prouver l'existence et l'unicité d'une suite $(u_n)_{n \in \mathbb{N}^*}$ vérifiant $$
\begin{cases}
	u_1 = x \\
	\forall n \in \mathbb{N}^*, u_{n+1} = f[(u_1, ..., u_n)]
\end{cases}$$
On pourra noter $f(u_1, ..., u_n)$ plutôt que $f[(u_1, ..., u_n)]$. \par
L'unicité se montre facilement par récurrence forte. Pour l'existence, commençons par définir une opération de concaténation: si $\varphi = (x_1, ..., x_n) \in F, x \in E$, on pose $$\varphi \sqcup x = (x_1, ..., x_n, x)$$
Formellement, il s'agit de l'unique prolongement de $\varphi$ à $\llbracket1, T(\varphi) + 1\rrbracket$ qui vaut $x$ en $T(\varphi) + 1$. Ensuite, appliquons le théorème démontré précédemment en posant $$\begin{cases}
	\varphi_1 = (x) \\
	\forall n \in \mathbb{N}^*, \varphi_{n+1} = \varphi_n \sqcup f(\varphi_n)
\end{cases}$$
Par une récurrence immédiate, on montre que $T(\varphi_n) = n$. Si on note $u_n = \varphi_n(n)$ le dernier terme du $n$-uplet $\varphi_n$, une autre récurrence immédiate montre que $\forall n \in \mathbb{N}^* \varphi_n = (u_1, ..., u_n)$.
On en déduit que la suite $(u_n)_{n \in \mathbb{N}^*}$ convient.

\subsection{Retour sur la construction de $\mathbb{N}$}

Dans cette partie, on donne les démonstrations détaillées des résultats évoqués dans le cours au sujet de l'addition et de la multiplication dans $\mathbb{N}$.
\begin{prop}
	L'addition est associative.
\end{prop}
\begin{proof}
	On montre par récurrence sur $k \in \mathbb{N}$ que $$\forall (m, n) \in \mathbb{N}^2, m + (n + k) = (m + n) + k$$
\begin{itemize}
	\item Initialisation: $m + (n + 0) = m + n = (m + n) + 0$.
	\item Hérédité: supposons que $m + (n + k) = (m + n) + k$. Alors $$m + (n + k^+) = m + (n + k)^+ = (m + (n + k))^+ = ((m + n) + k)^+ = (m + n) + k^+$$
\end{itemize}
Ainsi la récurrence est achevée.
\end{proof}
\begin{lem}
	Pour tout $n \in \mathbb{N}$, on a $n = 0 + n$ et $n^+ = 1 + n$
\end{lem}
\begin{proof}
	Montrons par récurrence sur $n$ que $n = 0 + n$.
\begin{itemize}
	\item $0 = 0 + 0$.
	\item $n = 0 + n$, alors $0 + n^{+} = (0 + n)^{+} = n^{+}$
\end{itemize}
Montrons par récurrence sur $n$ que $n^{+} = 1 + n$.
\begin{itemize}
	\item $0^{+} = 1 = 1 + 0$.
	\item Si $n^{+} = 1 + n$, alors $(n^{+})^{+} = (1 + n)^{+} = 1 + n^{+}$.
\end{itemize}
Ainsi la preuve est achevée.
\end{proof}
\begin{cor}
	L'addition est commutative.
\end{cor}
\begin{proof}
	Montrons par récurrence sur $k$ que $\forall n \geqslant 0, n + k = k + n$
\begin{itemize}
	\item $\forall n \geqslant 0, n + 0 = 0 + n = n$ d'après le lemme précédent.
	\item Soit $n \geqslant 0$. Alors $n + k^{+} = (n + k)^{+} = (k + n)^{+} = k + n^{+}$,
	par hypothèse de récurrence.
	Or d'après le lemme précédent, $k + n^{+} = k + (1 + n) = (k + 1) + n$ par associativité.
	Enfin: $(k + 1) + n = (k + 0^{+}) + n = (k + 0)^{+} + n = k^{+} + n$.
\end{itemize}
Ainsi la récurrence est achevée.
\end{proof}
\begin{lem}
	Si $m^{+} = n^{+}$ alors $m = n$.
\end{lem}
\begin{proof}
	Supposons que $m < n$. Alors $n$ viendrait (strictement) s'intercaler entre $m$
et $n^{+} = m^{+}$, absurde. Donc $m \geqslant n$. Symétriquement on a $n \geqslant m$, puis on conclut par antisymétrie.
\end{proof}
\begin{prop}
	$k + n = k' + n$ alors $k = k'$. Ainsi, à partir de maintenant, il sera permis de simplifier par un même terme de part et d'autre du signe dans une addition.
\end{prop}
\begin{proof}
	On montre par récurrence sur $n$ que $k + n = k' + n$ alors $k = k'$
\begin{itemize}
	\item Initialisation: c'est trivial.
	\item Hérédité: supposons que $k + n^{+} = k' + n^{+}$ alors
	$k + n^{+} = k' + n^{+} \Rightarrow (k + n)^{+} = (k' + n)^{+} \Rightarrow k + n = k' + n$, et on conclut en utilisant l'hypothèse de récurrence.
\end{itemize}
Ainsi la récurrence est achevée.
\end{proof}
\begin{lem}
	On a $m \geqslant n\iff\exists k \in \mathbb{N}, m = n + k$, et $k$ est alors unique.
\end{lem}
\begin{proof}
	Commençons par le sens direct; on fixe $n \in \mathbb{N}$ et on montre par récurrence sur $m \geqslant n$ l'existence de $k$ tel que $m = n + k$.
	\begin{itemize}
		\item Initialisation: $m = m + 0$.
		\item Hérédité: supposons que la propriété est vraie au rang $m$. Alors on écrit $m^{+} = (n + k)^{+} = n + k^{+}$.
	\end{itemize}
	Quant à l'unicité, c'est une simple conséquence de la propriété de simplification. Venons-en maintenant à la réciproque, que nous montrons par récurrence sur $k$.
	\begin{itemize}
		\item Initialisation: $m + 0 = m \geqslant m$.
		\item Hérédité: $m + k^+ = (m + k)^+ \geqslant m + k \geqslant m$.
	\end{itemize}
	Ainsi la récurrence est achevée.
\end{proof}
\begin{cor}
	On en déduit $m \leqslant n \Rightarrow m + p \leqslant n + p$.
\end{cor}
\begin{proof}
	On peut introduire $k$ tel que $n = m + k$. En utilisant l'associativité et la commutativité on en déduit: $n + p = (m + k) + p = (m + p) + k \geqslant m$.
\end{proof}
\begin{lem}
	La multiplication est distributive à droite sur l'addition: $(m + n)p = mp + np$.
\end{lem}
\begin{proof}
	On montre par récurrence sur $p \in \mathbb{N}$ que $$\forall(m,n) \in \mathbb{N}^2,\ (m + n)p = mp + np$$
\begin{itemize}
	\item Initialisation: $(m + n) \cdot 0 = 0 = 0 + 0 = m \cdot 0 + n \cdot 0$.
	\item Hérédité: $(m + n)(p + 1) = (m + n)p + (m + n) = (mp + np) + (m + n)$ par hypothèse de récurrence. Puis $(mp + np) + (m + n) = m(p + 1) + n(p + 1)$ par associativité et commutativité de de l'addition.
\end{itemize}
Ainsi la récurrence est achevée.
\end{proof}
\begin{lem}
	Pour tout $n \in \mathbb{N}$ on a $0 = 0 \cdot n$ et $n = 1 \cdot n$.
\end{lem}
\begin{proof}
	Déjà, on a $0 \cdot n = (0 + 0)n = 0 \cdot n + 0 \cdot n$ d'après le lemme précédent, et on conclut par simplification. Puis on montre par récurrence sur $n \in \mathbb{N}$ que $1 \cdot n = n$.
\begin{itemize}
	\item Initialisation : $1 \cdot 0 = 0$ par définition.
	\item Hérédité: $1 \cdot (n + 1) = 1 \cdot n + 1$ par définition, et $1 \cdot n = n$ par hypothèse de récurrence.
\end{itemize}
Ainsi la récurrence est achevée.
\end{proof}
\begin{cor}
	La multiplication est commutative.
\end{cor}
On montre par récurrence sur $n \in \mathbb{N}$ que $\forall m \in \mathbb{N}, mn = nm$.
\begin{itemize}
	\item Initialisation: $0 \cdot m = 0$ d'après la proposition précédente. Par ailleurs, on a par définition $m \cdot 0 = 0$.
	\item Hérédité: $m(n + 1) = mn + m = nm + m$ par hypothèse de récurrence.
	Puis $nm + m = nm + 1 \cdot m = (n + 1) \cdot m = (n + 1)m$ par distributivité,
	ce qui achève la récurrence.
\end{itemize}
\begin{cor}
	On en déduit la distributivité à gauche de la multiplication sur l'addition.
\end{cor}
\begin{proof}
	Il suffit d'écrire $m(n+p) = (n+p)m = nm + pm$.
\end{proof}
\begin{cor}
	La multiplication est associative.
\end{cor}
\begin{proof}
	On montre par récurrence sur $p\in\mathbb{N}$ que $\forall(m,n)\in\mathbb{N}^{2},\ (mn)p=m(np)$.
\begin{itemize}
	\item Initialisation : $(mn)\cdot0=0$ et $m(n\cdot0)=m\cdot0=0$.
	\item Hérédité: $(mn)(p+1)=(mn)p+mn=m(np)+mn$ par hypothèse de récurrence. Par distributivité à gauche, on a donc $m(np)+mn=m(np+n)=m(n(p+1))$, ce qui achève la récurrence.
\end{itemize}
\end{proof}
\begin{cor}
	On a $\forall(m,n,p)\in\mathbb{N}^{3},\ m\leqslant n\Rightarrow mp\leqslant np$.
\end{cor}
\begin{proof}
	On introduit $k$ tel que $n=m+k$ puis on a : $np=mp+kp\geqslant mp$.
\end{proof}
\begin{cor}
	On a $n_{1}...n_{k}=0\iff\exists i\in\llbracket1,k\rrbracket,\ n_{i}=0$
\end{cor}
\begin{proof}
	On montre la partie directe par récurrence sur $k\in\N\st$.
\begin{itemize}
	\item Initialisation : pour $k=1$, il n'y a rien à montrer.
	\item Hérédité: supposons $n_{1}...n_{k+1}=0$ avec $k\geqslant1$. Si $n_{k+1}=0$ c'est terminé. Sinon on a $n_{k+1}\geqslant1$. D'après le corollaire précédent, on a $0=n_{1}...n_{k+1}\ge n_{1}...n_{k}$. Donc $n_{1}\ldots n_{k}=0$, et on conclut par hypothèse de récurrence.
\end{itemize}
Pour la réciproque, il suffit d'utiliser la commutativité.
\end{proof}

\begin{cor}
	On a $$\forall(m,n,n^{\prime})\in\mathbb{N}^{*}\times\mathbb{N}^{2},\ mn=mn^{\prime}\Rightarrow n=n^{\prime}$$
Ainsi, à partir de maintenant, il sera permis de simplifier par un terme \textbf{non nul} à gauche et à droite du signe dans une multiplication.
\end{cor}
\begin{proof}
	Parmi $n$ et $n^{\prime}$, l'un des deux est inférieur ou égal à l'autre, car $\leqslant$ est une relation d'ordre totale. Sans perte de généralité, supposons que $n\leqslant n^{\prime}$ et écrivons $n^{\prime}=n+k$. Alors $mn^{\prime}=mn+mk$ donc $mk=0$ par simplification. Comme $m\ne0$, on en déduit $k=0$ puis $n=n^{\prime}$.
\end{proof}
Il nous reste à montrer les propriétés élémentaires sur les puissances.
\begin{prop}
	On a $m^{k}n^{k}=(mn)^{k}$ et $m^{k}m^{l}=m^{k+l}$.
\end{prop}
\begin{proof}
	On montre par récurrence sur $k$ que $m^{k}n^{k}=(mn)^{k}$.
\begin{itemize}
	\item Initialisation : $m^{0}n^{0}=1\cdot1=1=(mn)^{0}$
	\item Hérédité : $m^{k+1}n^{k+1}=(m^{k}m)(n^{k}n)=(mn)^{k}(mn)=(mn)^{k+1}$
\end{itemize}
On montre maintenant par récurrence sur $l$ que \(m^{k}m^{l}=m^{k+l}.\)
\begin{itemize}
	\item Initialisation: $m^{k}m^{0}=m^{k}\cdot1=m^{k}=m^{k+0}$
	\item Hérédité: $m^{k}m^{l+1}=(m^{k}m^{l})m=m^{k+l}m=m^{k+l+1}$
\end{itemize}
Ainsi la preuve est achevée.
\end{proof}
\subsection{Retour sur la construction de $\Z$}
Comment construire un entier négatif à partir des entiers positifs ? Justement en le voyant comme la différence de deux entiers positifs. Par exemple, on pourrait voir le nombre $-1$ comme le couple $(2,3)$, parce que $2 - 3 = -1$. Le problème, c'est qu'il n'y a pas unicité : on peut aussi écrire $-1 = 4 - 5$ ou $-1 = 79 - 80$. Il va donc falloir \textit{quotienter} $\N^2$ par une relation d'équivalence. \par
On voudrait donc identifier les couples $(a,b)$ et $(a',b')$ tels que $a - b = a' - b'$. Malheureusement, on n'a pas encore défini la soustraction. Mais suffit de "faire passer $b$ et $b'$ de l'autre côté" pour se ramener à des additions. On pose donc la
\begin{dfn}
	On définit sur l'ensemble $\N^2$ la relation suivante.
$$(a,b) \sim (a',b') \Leftrightarrow a + b' = b + a'$$
\end{dfn}
\begin{lem}
	C'est une relation d'équivalence.
\end{lem}
\begin{proof}
	On vérifie les trois axiomes d'une relation d'équivalence.
\begin{itemize}
	\item Réflexivité : comme l'addition sur $\N$ est commutative, on a $a + b = b + a$, soit $(a,b)\sim(a,b)$.
	\item Symétrie : si $a + b' = b + a'$, alors $a' + b = b' + a$ (par commutativité).
	\item Transitivité : supposons qu'on ait $a + b' = b + a'$ et $a' + b'' = b' + a''$. Si on avait le droit d'utiliser des soustractions, on écrirait tout simplement
	$$a - b = a' - b' = a'' - b''$$
	Mais encore une fois, nous n'avons le droit d'utiliser que l'addition. L'astuce consiste à ajouter un terme "fantôme" qui sert d'intermédiaire et que l'on simplifie à la fin. Ici, en l'occurrence il s'agit de $b$. Rappelons qu'on veut montrer $a + b'' = b + a''$. Si on ajoute $b$ de part et d'autre, en utilisant l'associativité et la commutativité de l'addition on a $$(a + b'') + b' = (a + b') + b'' = (b + a') + b'' = b + (a' + b'') = b + (b' + a'') = (b + a'') + b'$$ Il ne reste plus qu'à simplifier par b'.
\end{itemize}
Ainsi la preuve est achevée.
\end{proof}
\begin{dfn}
	On peut donc quotienter $\N^2$ par $\N$ et considérer $\overline{(a,b)}$, la classe d'équivalence de $(a,b)$. On note $\Z$ cet ensemble.
\end{dfn}
Puisque formellement, on peut écrire $\overline{(a,b)} = a - b$, on voudrait additionner deux entiers relatifs de la manière suivante : $(a - b) + (a' - b') = (a + a') - (b + b') = \overline{(a + a', b + b')}$. Pour l'instant, il ne s'agit que d'un sommes donc amenés à considérer le
\begin{lem}
	Si $(a,b)\sim(a',b')$ et $(c,d)\sim(c',d')$ alors $(a + c, b + d)\sim(a' + c', b' + d')$.
\end{lem}
\begin{proof}
	Si $a + b' = b + a'$ et $c + d' = d + c'$, alors $a + c + b' + d' = b + d + a' + c'$ (on a utilisé l'associativité et la commutativité de l'addition).
\end{proof}
\begin{cor}
	La définition suivante de l'addition sur $\Z$ est consistante : $$\overline{(a,b)} + \overline{(c,d)} = \overline{(a + c, b + d)}$$
	C'est-à-dire qu'elle ne dépend pas du représentant choisi pour chaque classe d'équivalence.
\end{cor}
Ainsi, pour additionner deux éléments de $\mathbb{Z}$, on prend deux couples d'entiers qui représentent les deux classes d'équivalence, on les additionne terme à terme, et on prend la classe du tout. On résume cela par : $$\overline{(a,b)} + \overline{(c,d)} = \overline{(a+c, b+d)}$$
Comme le procédé ne dépend pas du choix des représentants des classes, l'addition sur $\mathbb{Z}$ est bien définie.
\begin{thm}
	L'addition sur $\mathbb{Z}$ est associative et commutative.
\end{thm}
\begin{proof}
	Pour l'associativité, on a :
\begin{align*}
	\left(\overline{(a,b)} + \overline{(c,d)}\right) + \overline{(e,f)} &= \overline{(a+c, b+d)} + \overline{(e,f)} \\ &= \overline{((a+c)+e, (b+d)+f)} \\&= \overline{(a+(c+e), b+(d+f))} \\&= \overline{(a,b)} + \left(\overline{(c,d)} + \overline{(e,f)}\right)
\end{align*}
On voit qu'elle hérite directement de l'associativité de $+$ dans $\mathbb{N}$. De même, la commutativité vient de celle de $+$ dans $\mathbb{N}$.
\end{proof}
\begin{dfn}
	On choisit d'inclure $\mathbb{N}$ dans $\mathbb{Z}$ en identifiant $n \in \mathbb{N}$ avec $\overline{(n,0)} \in \mathbb{Z}$. Cette inclusion est bien injective. L'addition sur $\mathbb{Z}$ est alors un prolongement de celle sur $\mathbb{N}$.
\end{dfn}
\begin{proof}
	Si $n \neq n'$, alors $\overline{(n,0)} \neq \overline{(n',0)}$ puisque $n+0 \neq 0+n'$. De plus, pour tout couple $(n,n') \in \mathbb{N}^2$, on a $(n,0) + (n',0) = (n+n',0)$.
\end{proof}
Toujours par le même calcul formel on introduit la 
\begin{dfn}
	On définit sur $\mathbb{Z}$ l'ordre $\leqslant$ par : $\overline{(a,b)} \leqslant \overline{(c,d)} \Leftrightarrow a+d \leqslant b+c$. Cette définition est bien consistante.
\end{dfn}
\begin{proof}
	Supposons que $a+d \leqslant b+c$, $(a,b) \sim (a',b')$ et $(c,d) \sim (c',d')$. Écrivons $b+c = a+d+k$. Alors $a'+d'+k+b+c = a+d+k+b'+c' = b+c+b'+c'$. Donc $a'+d'+k = b'+c'$ par simplification. On a alors $a'+d' \leqslant b'+c'$. Réciproquement, si $a'+d' \leqslant b'+c'$ alors $a+d \leqslant b+c$ en échangeant les rôles des variables.
\end{proof}
\begin{prop}
	$\leqslant$ est une relation d'ordre totale, et elle prolonge la définition de $\leqslant$ sur $\mathbb{N}$.
\end{prop}
\begin{proof}
	On vérifie les trois axiomes d'une relation d'ordre.
	\begin{itemize}
		\item Réflexivité : $a+b = b+a$, donc $(a,b) \leqslant (a,b)$.
		\item Antisymétrie : si $(a,b) \leqslant (c,d)$ et $(c,d) \leqslant (a,b)$, alors $a+d \leqslant b+c$ et $c+b \leqslant d+a$, donc $a+d = b+c$, soit $\overline{(a,b)} = \overline{(c,d)}$.
		\item Transitivité : si $(a,b) \leqslant (c,d)$ et $(c,d) \leqslant (e,f)$, alors $a+d \leqslant b+c$ et $c+f \leqslant d+e$, donc $a+d+f \leqslant b+c+f \leqslant b+d+e$. Il existe donc $k$ tel que $b+d+e = a+d+f+k$, puis par simplification $b+e = a+f+k$, ce qui implique $a+f \leqslant b+e$, soit $(a,b) \leqslant (e,f)$.
	\end{itemize}
	Il reste à vérifier que $\leqslant$ est totale : si on n'a pas $a+d \leqslant b+c$ alors on a $b+c \leqslant a+d$ car $\leqslant$ est totale sur $\mathbb{N}$. En utilisant la commutativité dans $\mathbb{N}$ on en déduit $c+b \leqslant d+a$, soit $\overline{(c,d)} \leqslant \overline{(a,b)}$. Enfin, pour $(m,n) \in \mathbb{N}^2$ : $m \leqslant_{\mathbb{N}} n \Leftrightarrow m+0 \leqslant_{\mathbb{N}} 0+n \Leftrightarrow \overline{(m,0)} \leqslant_{\mathbb{Z}} (n,0)$.
\end{proof}
\begin{prop}
	La relation $\leqslant$ reste compatible avec l'addition : $$m \leqslant n \Rightarrow m+p \leqslant n+p$$
\end{prop}
\begin{proof}
	Supposons qu'on a $\overline{(a,b)} \leqslant \overline{(c,d)}$, et donnons-nous $p = \overline{(e,f)}$. Alors on a $a+d \leqslant b+c$. En sommant $\underline{e+f}$ et par commutativité, on en déduit $a+e+d+f \leqslant b+f+c+e$, soit encore $\overline{(a+e,b+f)} \leqslant \overline{(c+e,d+f)}$.
\end{proof}
La proposition capitale vue sur $\mathbb{N}$ reste vraie (en revanche, elle deviendra fausse dans $\mathbb{Q}$ et \textit{a fortiori} dans $\mathbb{R}$).
\begin{prop}
	Soit $m,n \in \mathbb{Z}, m < n+1$ alors $m \leqslant n$.
\end{prop}
\begin{proof}
	Si $\overline{(a,b)} < \overline{(c,d)} + 1 = \overline{(c+1,d)}$, alors $a+d \leqslant b+c+1, a+d \neq b+c+1$, donc $a+d < b+c+1$. D'après ce qui a été vu sur $\mathbb{N}$ on peut donc écrire $a+d \leqslant b+c$.
\end{proof}
\begin{prop}
	L'addition sur $\mathbb{Z}$ admet un élément neutre, et tout $x \in \mathbb{Z}$ admet un inverse.
\end{prop}
\begin{proof}
	L'élément neutre est donné par $\overline{(0,0)}$ et l'inverse de $\overline{(a,b)}$ est $\overline{(b,a)}$.
\end{proof}
On a vu que $\mathbb{N}$ pouvait être considéré comme inclus dans $\mathbb{Z}$. Ainsi, si $n \in \mathbb{N}$ l'écriture $n \in \mathbb{Z}$ un sens. D'après la proposition précédente, l'écriture $-n$ a elle aussi un sens. Réciproquement, on a la
\begin{prop}
	Tout élément de $\mathbb{Z}$ s'écrit d'une unique manière $0$, $n$ ou $-n$ avec $n \in \mathbb{N}^*$.
\end{prop}
\begin{proof}
	Pour l'existence, on se donne $(a,b) \in \mathbb{Z}$. Si $a > b$ alors on écrit $a = b+n$ avec $n \in \mathbb{N}^*$, puis on a 
$$\overline{(a,b)} = \overline{(n,0)} = n$$
De même, si $a < b$, alors on écrit $b = a+n$, puis on a
$$\overline{(a,b)} = \overline{(0,n)} = -\overline{(n,0)} = -n$$
Enfin, si $a = b$ on a $\overline{(a,b)} = \overline{(0,0)} = 0$. Pour ce qui est de l'unicité, soit $x \in \mathbb{Z}$
\begin{itemize}
	\item Si $x = 0$, alors il ne peut pas s'écrire sous la forme $\pm n$ avec $n \neq 0$ donc son écriture est unique.
	\item Si $x > 0$, alors il ne peut pas s'écrire sous la forme $-n$ avec $n \neq 0$ donc il ne peut s'écrire que sous la forme $n$ avec $n \neq 0$. $n$ est alors unique (rappelons que l'inclusion de $\mathbb{N}$ dans $\mathbb{Z}$ est injective).
	\item Si $x < 0$, alors il ne peut pas s'écrire sous la forme $n$ avec $n \neq 0$ donc il ne peut s'écrire que sous la forme $-n$ avec $n \neq 0$. Puis si $-n = -n'$, alors en ajoutant $n+n'$ de part et d'autre on obtient $n = n'$.
\end{itemize}
Ainsi la preuve est achevée.
\end{proof}
À présent, comment définir une multiplication sur $\Z$ ? Le même petit calcul formel donne $$(a-b)(c-d)=(ac+bd)-(ad+bc)$$
\begin{dfn}
	La définition suivante de la multiplication sur $\Z$ est consistante : $$\overline{(a,b)}\times\overline{(c,d)}=\overline{(ac+bd,ad+bc)}$$ De plus, elle prolonge la définition de la multiplication sur $\N$.
\end{dfn}
\begin{proof}
	Il suffit de vérifier la même indépendance par rapport au choix des représentants. Si $a+b'=b+a'$ et $c+d'=d+{c'}$, alors on veut montrer que $$(ac+bd,ad+bc)\sim(a'c'+b'd',a'd'+b'c')$$
On utilise le même genre d'astuce que dans $\N$, avec l'introduction d'un "terme fantôme". Le but est de montrer que $$ac+bd+a'd'+b'c'=ad+bc+a'c'+b'd'$$
Avec un peu plus de mal, on trouve que qu'un terme fantôme possible est $ad'+bc'+bd'+ac'$ :
\begin{align*}
	(ac+bd+a'd'+b'c')+(ad'+bc'+bd'+ac')&=a(c+d')+b(c'+d)+(a'+b)d'+(a+b')c'\\
	&=a(c'+d)+b(c+d')+(a+b')d'+(a'+b)c'\\
	&=(ad+bc+a'c'+b'd')+(ad'+bc'+bd'+ac')\\
\end{align*}
Après simplification, on obtient le résultat. De plus, on a
$$\overline{(n,0)}\cdot\overline{(n',0)}=\overline{\left(nn'+0\cdot0,n\cdot0+0\cdot n'\right)}=\overline{(nn',0)}$$
Donc la multiplication de $\Z$ prolonge bien celle de $\N$.
\end{proof}
\begin{thm}
	La multiplication sur $\Z$ est associative, commutative, et distributive sur l'addition.
\end{thm}
\begin{proof}
	Vérifions l'associativité.
\begin{align*}
	\left(\overline{(a,b)}\cdot\overline{(c,d)}\right)\cdot\overline{(e,f)}&=\overline{(ac+bd,ad+bc)}\cdot\overline{(e,f)}\\
	&=\overline{((ac+bd)e+(ad+bc)f,(ac+bd)f+(ad+bc)e)}\\
	&=\overline{(a(ce+df)+b(cf+de),a(cf+de)+b(ce+df))}\\
	&=(a,b)\cdot(ce+df,cf+de)\\
	&=\overline{(a,b)}\cdot\left(\overline{(c,d)}\cdot\overline{(e,f)}\right)
\end{align*}
Vérifions maintenant la commutativité.
\begin{align*}
	\overline{(a,b)}\cdot\overline{(c,d)}&=\overline{(ac+bd,ad+bc)}\\
	&=\overline{(ca+db,cb+da)}\\
	&=\overline{(c,d)}\cdot\overline{(a,b)}
\end{align*}
Quant à la distributivité, par commutativité il suffit de la vérifier d'un côté seulement.
\begin{align*}
	\left(\overline{(a,b)}+\overline{(c,d)}\right)\cdot\overline{(e,f)}&=\overline{(a+c,b+d)}\cdot\overline{(e,f)}\\
	&=\overline{((a+c)e+(b+d)f,(a+c)f+(b+d)e)}\\
	&=\overline{(ae+bf+ce+df,af+be+cf+de)}\\
	&=\overline{(ae+bf,af+be)}+\overline{(ce+df,cf+de)}\\
	&=\overline{(a,b)}\cdot\overline{(e,f)}+\overline{(c,d)}\cdot\overline{(e,f)}
\end{align*}
Ainsi la preuve est achevée.
\end{proof}
\begin{rmq}
	On vérifie immédiatement que 1 continue à être élément neutre.
\end{rmq}
\begin{cor}
	On a 
$$\begin{cases}
	\forall n\in\mathbb{Z},\ n\times0=0\times n=0\\
	\forall(m,n)\in\mathbb{Z}^{2},\ m(-n)=(-m)n=-(mn)
\end{cases}$$
On dit que 0 est absorbant.
\end{cor}
\begin{proof}
	Pour la première ligne, il suffit d'écrire $n=n\times1=n(1+0)=n+n\times0$, puis de simplifier par n. Le raisonnement est le même avec $0\times n$. \par 
Pour la seconde on écrit $m(-n)+mn=m(-n+n)=m\times0=0.$ Le raisonnement est le même avec $(-m)n$.
\end{proof}
\begin{prop}
	On sait que dans $\Z$, la relation $\leqslant$ reste compatible avec l'addition. Pour la multiplication, il faut être plus prudent.
$$\begin{cases}
	m\leqslant n\\
	p\geqslant0
\end{cases}\Rightarrow mp\leqslant np$$
\end{prop}
\begin{proof}
	En utilisant la distributivité et la compatibilité avec l'addition, tout revient à montrer que $n\geqslant0$ et $p\geqslant0,$ alors $np\geqslant0$ Pour cela on peut invoquer le fait que la multiplication sur $\Z$ prolonge celle sur $\N$.
\end{proof}
\begin{rmq}
	On en déduit qu'un carré (d'entier) est toujours positif ou nul.
\end{rmq}

\subsection{Formules sommatoires}
On se donne un ensemble $E$ muni d'une loi associative $*$ et d'un élément neutre $e$.
\begin{prop}
	Quels que soient $0\leqslant k\leqslant m\leqslant n$, $$x_{1}*...*x_{n}=x_{1}*...*x_{k}*(x_{k+1}*...*x_{m})*x_{m+1}*...*x_{n}$$
(avec la convention habituelle si $k=0$ ou $m=n$).
\end{prop}
\begin{proof}
	Dans un premier temps, fixons $k$ et supposons que $n>m$. Il s'agit alors simplement de montrer $x_{1}*...*x_{m}=x_{1}*...*x_{k}*(x_{k+1}*...*x_{m}).$ Si $k=0$ il n'y a rien à prouver. Sinon, effectuons une récurrence sur $m$.
\begin{itemize}
	\item Initialisation : si m=k, il n'y a rien à prouver.
	\item Hérédité: si la propriété est vraie au rang m, écrivons 
	\begin{align*}
		x_{1}*...*x_{m+1}&=(x_{1}*...*x_{m})*x_{m+1}\\
		&=(x_{1}*...,*x_{k}*(x_{k+1}*...*x_{m}))*x_{m+1}\\
		&=((x_{1}*...*x_{k})*(x_{k+1}*...*x_{m}))*x_{m+1}\\
		&=(x_{1}*...*x_{k})*((x_{k+1}*...*x_{m})*x_{m+1})\\
		&=x_{1}*...*x_{k}*(x_{k+1}*...*x_{m+1})
	\end{align*}
\end{itemize}
Montrons à présent le cas général par récurrence sur $n$.
\begin{itemize}
	\item Initialisation : si $n=m$, c'est le cas précédent.
	\item Hérédité: si la propriété est vraie au rang $n$, on a
	\begin{align*}
		x_{1}*...*x_{n+1}&=(x_{1}*...*x_{n})*x_{n+1}\\	
		&=(x_{1}*...*x_{k}*(x_{k+1}*...*x_{m})*x_{m+1}*...*x_{n})*x_{n+1}\\
		&=x_{1}*...*x_{k}*(x_{k+1}*...*x_{m})*x_{m+1}*...*x_{n+1}
	\end{align*}
\end{itemize}
On vérifie que le calcul reste valable pour $k=0$. Ainsi la preuve est achevée.
\end{proof}
À partir de maintenant, nous travaillons sur une loi associative \textbf{et} commutative. C'est pourquoi nous la noterons plutôt $+$, afin de mieux évoquer la commutativité, mais on pourrait aussi la noter $\times$ dans le cas de la multiplication. En particulier, l'élément neutre sera noté $0$, et la composée
$x_{1}*...*x_{n}$ sera notée $\displaystyle\sum_{k=1}^{n}x_{k}$.
\begin{dfn}[Permutations]
	Soit $n\in\mathbb{N}$. Une \textbf{permutation} de $\llbracket1,n\rrbracket$ est une bijection de l'ensemble $\llbracket1,n\rrbracket$ sur lui-même. L'ensemble des permutations de $\llbracket1,n\rrbracket$ est noté $\mathcal{S}_{n}$ voire plus anciennement (et quand on arrive à le calligraphier) $\mathfrak{S}_n$.
\end{dfn}
\begin{thm}
	Avec la loi $+$, on a le droit de composer les termes dans l'ordre que l'on veut.
Autrement dit, $$\forall\sigma\in\mathcal{S}_{n},x_{1}+...+x_{n}=x_{\sigma(1)}+...+x_{\sigma(n)}$$
\end{thm}
\begin{proof}
	Prouvons le résultat par récurrence sur $n$.
\begin{itemize}
	\item Initialisation : pour $n=0$, le seul élément de $\mathcal{S}_{0}$ est "l'application vide".
	\item Hérédité : supposons la propriété vraie au rang $n$ avec $n\geqslant1$ et montrons-la au rang $n+1$. Supposons dans un premier temps que $\sigma(n+1)=n+1$ Alors $\sigma$ induit une permutation de $\llbracket1,n\rrbracket$ et il suffit d'invoquer l'hypothèse de récurrence. Dans un second temps, supposons que $\sigma(n+1)\ne n+1$. Grâce à l'hypothèse de récurrence, on peut supposer sans perte de généralité que $\sigma(n)=n+1$. Alors, par un simple argument de commutativité, on permute $x_{\sigma(n)}$ et $x_{\sigma(n+1)}$ et on se ramène au cas précédent.
\end{itemize}
Ainsi la récurrence est achevée.
\end{proof}
\begin{dfn}
	Soit $I$ un ensemble d'indices fini, et $(x_{i})_{i\in I}$ une famille à valeurs dans $(E,+)$. Si on pose $n:=|I|$, il existe une bijection $\sigma$ de $\llbracket1,n\rrbracket$ sur $I$. On pose alors $$\sum_{i\in I}x_{i}=\sum_{k=1}^{n}x_{\sigma(k)}$$
	D'après le théorème précédent, la définition est bien indépendante du choix de $\sigma$. En effet, si on se donne une autre permutation $\sigma^{\prime}$, en notant $\tau:=\sigma^{-1}\circ\sigma^{\prime}$ on obtient $$\sum_{k=1}^{n}x_{\sigma(k)}=\sum_{k=1}^{n}x_{\sigma(\tau(k))}=\sum_{k=1}^{n}x_{\sigma^{\prime}(k)}$$
\end{dfn}
\begin{rmq}
	On observera que la formule $\displaystyle\sum_{i\in I}(x_{i}+y_{i})=\sum_{i\in I}x_{i}+\sum_{i\in I}y_{i}$ est une conséquence directe de l'associativité et de la commutativité de $+$.
\end{rmq}
\begin{lem}
	Si $I_{1}$ et $I_{2}$ sont des parties disjointes de $I$, alors $\displaystyle\sum_{i\in I_{1}\cup I_{2}}x_{i}=\sum_{i\in I_{1}}x_{i}+\sum_{i\in I_{2}}x_{i}$.
\end{lem}
\begin{proof}
	Raisonnons par récurrence sur $\card(I_1 \cup I_2)$.
	\begin{itemize}
		\item Initialisation : si $\card(I_1 \cup I_2) = 0$, alors $I_1 = I_2 = \emptyset$ donc il n'y a rien à prouver.
		\item Hérédité : donnons-nous une bijection $\sigma$ de $\llbracket1,\ldots, n+1\rrbracket$ sur $I_1 \cup I_2$ et supposons sans perte de généralité que $\sigma(n+1) \in I_2$.
		\begin{itemize}
			\item[--] Si $I_1 = \emptyset$, il n'y a toujours rien à prouver.
			\item[--] Sinon, introduisons $I_2' = I_2 \setminus \{\sigma(n+1)\}$. Comme $I_1 \neq \emptyset$, $I_2'$ est strictement inclus dans $I_1 \cup I_2$ et par hypothèse de récurrence on a $$\sum_{i \in I_2} x_i = \left(\sum_{i \in I_2'} x_i\right) + x_{\sigma(n+1)}$$
			De même, $I_1 \cup I_2'$ est strictement inclus dans $I_1 \cup I_2$ donc : $$\sum_{i \in I_1 \cup I_2'} x_i = \sum_{i \in I_1} x_i + \sum_{i \in I_2'} x_i$$
			Enfin, remarquons que $\sigma$ induit une bijection de $\llbracket1,\ldots,n\rrbracket$ sur $I_1 \cup I_2'$. Par associativité, on en déduit
			\begin{align*}
				\sum_{i \in I_1} x_i + \sum_{i \in I_2} x_i &= \sum_{i \in I_1} x_i + \left(\sum_{i \in I_2'} x_i + x_{\sigma(n+1)}\right)\\
				&= \left(\sum_{i \in I_1} x_i + \sum_{i \in I_2'} x_i\right) + x_{\sigma(n+1)}\\
				&= \left(\sum_{i \in I_1 \cup I_2'} x_i\right) + x_{\sigma(n+1)}\\
				&= \left(\sum_{k=1}^n x_{\sigma(k)}\right) + x_{\sigma(n+1)}\\
				&= \sum_{k=1}^{n+1} x_{\sigma(k)}
			\end{align*}
			ce qui achève la récurrence.
		\end{itemize}
	\end{itemize}
\end{proof}
\begin{thm}
	Si $I_1,\ldots, I_n$ sont des parties deux à deux disjointes de $I$, alors $$\sum_{i \in I_1 \cup ... \cup I_n} x_i = \sum_{i \in I_1} x_i + ... + \sum_{i \in I_n} x_i$$
\end{thm}
\begin{proof}
	C'est une récurrence immédiate sur $n$.
\end{proof}




\newpage
\chapter{Groupes, anneaux, corps}
\section{Groupes}
Pour commencer, rappelons que la composition des applications est associative. En revanche, elle n'a aucune raison d'être commutative.
\subsection{Premières définitions}
De manière schématique, un groupe est un ensemble dans lequel on peut toujours composer deux éléments de manière associative, qui contient un élément neutre, et tel que tout élément possède un inverse.
\begin{dfn}[Groupe]
	Soit un ensemble $G$ muni d'une loi de composition interne $*$. La structure algébrique $(G,*)$ est dite un \textbf{groupe} lorsqu'elle vérifie les trois axiomes suivants.
	\begin{itemize}
		\item $*$ est associative.
		\item $(G,*)$ est muni d'un \textbf{élément neutre}.
		\item Tous les éléments de $G$ sont \textbf{inversibles} pour $*$.
	\end{itemize}
Par un léger abus de terminologie, on dira souvent que $G$ lui-même est un groupe (pour la loi $*$). Généralement l'élément neutre de $G$ est noté $e$, et l'inverse de $x\in G$ sera noté $x\inv$. Assez rapidement, s'il n'y a pas d'ambiguïté, on omettra le symbole $*$ et on notera simplement $xy$ pour $x*y$. Si de plus, $*$ est commutative, on dit que $G$ est \textbf{commutatif} (on rencontrera parfois le mot \textbf{abélien}). Sa loi pourra alors être notée $+$, et son neutre pourra être noté $0$. Dans ce cas, on rappelle que l'inverse de $x$ est usuellement noté $-x$, et que l'écriture $x-y$ désigne en fait $x+(-y)$.
\end{dfn}
\begin{exa}
	$(\Z,+)$ est un groupe commutatif.
\end{exa}
\begin{exa}
	Jusqu'au baccalauréat, intentionnellement ou non, on rencontre plusieurs autres exemples de groupes.
	\begin{itemize}
		\item Tout d'abord, l'ensemble des nombres complexes $\C$ muni de l'addition : la somme de deux complexes est un complexe, et l'opposé d'un complexe est un complexe.
		\item On aurait aussi pu considérer $\C\st$ muni de la multiplication : le produit de deux complexes non nuls est un complexe non nul, et l'inverse d'un complexe non nul est un complexe non nul.
		\item Ensuite, l'ensemble des nombres complexes de module $1$ muni de la multiplication (autrement dit, le cercle unité) : si on multiplie $z$ et $z'$ de module $1$, alors $zz'$ est de module $1$, et l'inverse d'un nombre complexe de module $1$ est de module $1$.
		\item Ensuite, le premier groupe fini rencontré : les racines $n$-èmes de l'unité. Le produit de deux racines de l'unité est une racine $n$-ème de l'unité, et l'inverse d'une racine $n$-ème de l'unité est une racine $n$-ème de l'unité.
		\item Après, la grande famille de la géométrie : les translations, les homothéties-translations, les rotations-translations, les similitudes... Tous ces ensembles forment des groupes pour la composition des applications. Attention, la composition est notre premier exemple de loi qui ne soit pas commutative. Par exemple, deux similitudes qui n'ont pas le même centre ne commutent pas en général.
		\item On peut citer le cas de vecteurs de $\R^2$ ou de $\R^3$, voire de $\R^n$ en général, qui sont des groupes pour l'addition.
	\end{itemize}
\end{exa}
On voit à quel point les groupes sont omniprésents en mathématiques, dans des branches très variées. L'intérêt d'identifier cette structure est de mettre au point une fois pour toutes des théorèmes généraux, valables dans n'importe quel groupe, et ainsi on n'aura pas besoin de redémontrer à chaque fois le même résultat "déguisé" sous une forme différente.
\begin{prop}
	Dans un groupe, on peut "simplifier" une égalité à droite ou à gauche : $$\forall(x,y,z)\in G^3,\ \left \{ \begin{array}{ll}xy=xz\implies y=z \\ yx=zx\implies y=z\end{array}\right.$$ On dit que $x$ est \textbf{régulier}.
\end{prop}
\begin{rmq}
	Rappelons que dans un groupe, l'inverse de $xy$ est $y\inv x\inv$. Plus généralement, l'inverse de $x_1\ldots x_n$ est $x_n\inv\ldots x_1\inv$.
\end{rmq}
\begin{exa}
	L'ensemble $\mathcal{S}_X$ des permutations d'un ensemble quelconque $X$ forme un groupe pour la composition. En particulier, $\mathcal{S}_{\llbracket1,n\rrbracket}$ est un groupe non commutatif dès que $n\geqslant3$. On l'appelle \textbf{groupe symétrique d'ordre $n$} et on le note plutôt $\mathcal{S}_n$.
\end{exa}
\begin{rmq}
	Soit $G$ un groupe, pas nécessairement commutatif. Si $xy=e$, alors en composant par $x\inv$, on obtient $y=x\inv$, d'où $yx=e$.
\end{rmq}
\begin{prop}[Groupe-produit]
	Soit $(G_1,*_1)$ et $(G_2,*_2)$ deux groupes. Alors $G_1\times G_2$ muni de la loi $$(x_1,x_2)*(y_1,y_2):=(x_1*_1y_1,\ x_2*_2y_2)$$ est un groupe, dit \textbf{groupe-produit} de $G_1$ et $G_2$. On omettra souvent les signes $*$, $*_1$ et $*_2$.
\end{prop}
\begin{rmq}
	Si $G_1$ et $G_2$ sont commutatifs, alors $G_1\times G_2$ est commutatif.
\end{rmq}
\begin{rmq}
	Bien sûr, on peut étendre les résultat précédents à un produit $G_1\times\ldots\times G_n$ pour $n\geqslant2$.
\end{rmq}
\begin{rmq}
	Comme $(\R,+)$ est un groupe, on en déduit immédiatement les structures de groupe de $(\R^2,+)$, $(\R^3,+)$ et $(\R^n,+)$.
\end{rmq}
\subsection{Sous-groupes}
\begin{dfn}[Sous-groupe]
	Soit $G$ un groupe. Un \textbf{sous-groupe} $H$ de $G$ est une partie $H\subset G$ telle que :
	\begin{itemize}
		\item $e\in H$
		\item $\forall (x,y)\in H^2,\ xy\in H$ (stabilité par la l.c.i)
		\item $\forall x\in H,\ x\inv\in H$ (stabilité par passage à l'inverse)
	\end{itemize}
\end{dfn}
\begin{prop}
	Soit $G$ un groupe et $H$ un sous-groupe. Muni de la loi induite par celle de $G$, $H$ est un groupe.
\end{prop}
\begin{rmq}
	Si $G$ est commutatif, alors tous ses sous-groupes sont commutatifs.
\end{rmq}
\begin{exa}
	L'ensemble des nombre complexes de module $1$ forme un sous-groupe multiplicatif de $\C\st$. C'est donc un groupe commutatif.
\end{exa}
\begin{prop}[Recette du sous-groupe]
	Soit $G$ un groupe et $H\subset G$. Pour que $H$ soit un sous-groupe de $G$, il faut et il suffit que \begin{itemize}
		\item $H\ne\emptyset$ (par exemple $e\in H$)
		\item $\forall(x,y)\in H^2,\ xy\inv\in H$
	\end{itemize}
\end{prop}
\begin{rmq}
	En pratique, on n'utilise que la condition suffisante. En effet, bien souvent, on pourra utiliser cette recette lorsqu'un énoncé nous demandera de démontrer que tel ou tel ensemble est un groupe : en l'interprétant comme un sous-ensemble d'un groupe connu, on montrera directement qu'il s'agit d'un sous-groupe.
\end{rmq}
\begin{exa}
	$\left\{e\right\}$ est le plus petit sous-groupe de $G$.
\end{exa}
\begin{exa}
	Soit $(G,+)$ un groupe commutatif, et $A$ et $B$ deux sous-groupes. Alors $A+B$ est un sous-groupe, comme nous le vérifions maintenant.
	\begin{itemize}
		\item $0=0+0\in A+B$
		\item Soit $a+b\in A+B$ et $a'+b'\in A+B$. Alors :
		\begin{align*}
			(a+b)-(a'+b')&=(a+b)+(-b'-a') \\
			&= (a-a')+(b-b')&\text{par commutativité et associativité} \\
			&\in A+B
		\end{align*}
	\end{itemize}
\end{exa}
\begin{exa}
	Soit $G$ un groupe non nécessairement commutatif. On appelle \textbf{centre} de $G$ l'ensemble des éléments qui commutent avec tous les éléments de $G$ : $$Z(G)=\left\{x\in G\ | \ \forall y\in G,\ xy=yx\right\}$$ (le "Z" vient de l'allemand "\textit{Zentrum}", les Allemands étant omniprésents dans l'histoire de l'algèbre). $Z(G)$ est un sous-groupe de $G$, comme nous le vérifions maintenant.
	\begin{itemize}
		\item Comme $\forall x\in G,\ ex=xe=x$, on a $e\in Z(G)$.
		\item Soit $(x,y)\in\Z(G)^2$. Soit $z\in G$. On a :
		\begin{align*}
			(xy\inv)z&=x(y\inv z) & \text{par associativité} \\
			&=x(z\inv y)\inv \\
			&=x(yz\inv)\inv & \text{par} \ y\in Z(G) \\
			&=x(zy\inv)\\
			&=(xz)y\inv & \text{par associativité} \\
			&=(zx)y\inv & \text{car} \ x\in Z(G)\\
			&=z(xy\inv) & \text{par associativité}
		\end{align*}
		si bien que $xy\inv\in Z(G)$
	\end{itemize}
\end{exa}
\begin{thm}
	Soit $\mathcal{H}$ un ensemble de sous-groupes de $G$. Alors $\displaystyle \bigcap_{H\in\mathcal{H}} H$ est un sous-groupe de $G$.
\end{thm}
Ce théorème justifie la définition suivante.
\begin{dfn}[Sous-groupe engendré]
	Soit $X$ une partie de $G$. L'intersection de tous les sous-groupes de $G$ qui contiennent $X$ s'appelle le \textbf{sous-groupe engendré par $X$}. Puisqu'il est inclus dans tout sous-groupe contenant $X$, c'est le plus petit sous-groupe contenant $X$. On le note $<X>$, voire $(X)$. Si $x$ est un élément de $G$, on notera plus simplement $<x>$ plutôt que $<\left\{x\right\}>$.
\end{dfn}
\begin{exa}
	Dans $(\R,+)$, on a $<1>=\Z$.
\end{exa}
\begin{dfn}[Itération de la l.c.i.]
	Soit $G$ un groupe dont la loi est notée multiplicativement, et $x\in G$. On rappelle qu'on pose $$\left \{ \begin{array}{ll}
		x^0=e \\ \forall n\in\N,\ x^{n+1}=x^nx
	\end{array}\right.$$ Puis on rappelle qu'on étend la définition à $n\in\Z$ en posant $$\forall n\in\N\st,\ x^{-n}=(x^n)\inv$$ On rappelle qu'on a aussi qu'on a les relations suivantes : $$\forall (m,n)\in\Z^2,\ \left \{ \begin{array}{ll} x^mx^n=x^{m+n} \\ (x^m)^n=x^{mn}\end{array}\right.$$ Et si $x$ et $y$ commutent, alors : $$\left\{\begin{array}{ll}\forall(m,n)\in\Z^2,\ x^my^n=y^nx^m \\ \forall n\in\Z,\ (xy)^n=x^ny^n \end{array}\right.$$
\end{dfn}
\begin{exa}
	On a $\forall n\in\Z,\ e^n=e$.
\end{exa}
\begin{rmq}
	Dans le cas où le groupe est commutatif, il n'est pas rare que sa loi soit notée additivement. Dans ce cas, l'usage veut que $x^n$ soit plutôt noté $nx$ (pour correspondre à l'idée qu'on se fait de $x+\ldots+x$). La définition se réécrit ainsi : $$\left \{ \begin{array}{lll}0x=0 \\ \forall n\in\N, (n+1)x=nx+x\\ \forall n\in\N\st,\ (-n)x=-(nx)\end{array}\right.$$ Par ailleurs, on obtient les propriétés suivantes. $$\left \{\begin{array}{llll}\forall(m,n)\in\Z^2,\ mx+nx=(m+n)x \\ \forall(m,n)\in\Z^2,\ n(mx)=(mn)x & (=(nm)x=m(nx))\\ \forall(m,n)\in\Z^2,\ mx+ny=ny+mx & (\text{par commutativité})\\ \forall n\in\Z,\ n(x+y)=nx+ny & (\text{par commutativité})\end{array}\right.$$
\end{rmq}
\begin{prop}
	Soit $G$ un groupe et $x\in G$. Alors le sous-groupe engendré par $x$ vaut $$<x>=\left\{x^n\ | \ n\in\Z\right\}$$
\end{prop}
\begin{rmq}
	Si la loi du groupe est notée additivement, on a $<x>=\left\{nx\ | \ n\in\Z\right\}$
\end{rmq}
\begin{dfn}[Groupe monogène, groupe cyclique]
	Un groupe est \textbf{monogène} lorsqu'il est engendré par un seul élément. Il est \textbf{cyclique} lorsqu'il est monogène et fini.
\end{dfn}
\begin{exa}
	$(\Z,+)$ est monogène puisqu'il est engendré par $1$. $(\U_n,\times)$ est cyclique, engendré par $\displaystyle\omega=\exp\left(i\frac{2\pi}{n}\right)$.
\end{exa}
\begin{rmq}
	Tout groupe monogène est nécessairement commutatif, puisque tous ses éléments s'écrivent en itérant un unique élément.
\end{rmq}
Pour finir, citons le célèbre
\begin{thm}[Théorème de Lagrange, HP]
	Soit $G$ un groupe fini, et $H$ un sous-groupe de $G$. Alors, $\card(H)$ divise $\card(G)$.
\end{thm}
\begin{proof}
	On trouvera une preuve dans les compléments sur les groupes.
\end{proof}
\subsection{Morphismes}
\begin{dfn}[Morphisme de groupes]
	Soit $G$ et $G'$ deux groupes. Un \textbf{morphisme} de $G$ dans $G'$ est une application $\varphi:G\rightarrow G'$ telle que $$\forall(x,y)\in G^2,\ \varphi(xy)=\varphi(x)\varphi(y)$$ Autrement dit, $\varphi$ "transporte la loi de $G$ dans celle de $G'$". On note $Hom(G,G')$ l'ensemble des morphismes de $G$ dans $G'$.
\end{dfn}
\begin{rmq}
	La notation $Hom(G,G')$ vient de ce qu'anciennement, un morphisme s'appelait un homomorphisme.
\end{rmq}
\begin{dfn}[Isomorphisme, endomorphisme, automorphisme]
	Un \textbf{isomorphisme} est un morphisme bijectif. Un \textbf{endomorphisme} est un morphisme de $G$ dans lui-même. On note $Hom(G)$ l'ensemble des endomorphismes de $G$. Enfin, un \textbf{automorphisme} est un endomorphisme qui est aussi un isomorphisme.
\end{dfn}
\begin{prop}[Propriétés des morphismes]
	Pour tout morphisme $\varphi:G\rightarrow G'$, on a : $$\left \{ \begin{array}{ll}\varphi(e_G)=e_{G'}\\ \forall x\in G,\ \varphi(x\inv)=\varphi(x)\inv\end{array}\right.$$
\end{prop}
\begin{rmq}[Morphisme de groupes et itération]
	Plus généralement, on montre que $$\forall x\in G,\ \forall k\in\Z,\ \varphi\left(x^k\right)=\varphi(x)^k$$
\end{rmq}
\begin{dfn}[Noyau]
	Soit $\varphi$ un morphisme de $G$ dans $G'$. Le \textbf{noyau de $\varphi$} est défini par $$\ker(\varphi):=\varphi\inv\left(\left\{e_{G'}\right\}\right)=\left\{x\in G\ | \ \varphi(x)=e_{G'}\right\}$$ la notation $\ker$ vient de l'allemand "\textit{der Kern}" qui signifie "noyau".
\end{dfn}
\begin{dfn}[Image]
	Soit $\varphi$ un morphisme de $G$ dans $G'$. L'\textbf{image de $\varphi$} est défini par $$\im(\varphi):=\varphi(G)=\left\{y\in G'\ | \ \exists x\in G,\ y=\varphi(x)\right\}$$
\end{dfn}
\begin{thm}
	L'image d'un sous-groupe $H$ de $G$ par un morphisme de groupes est un sous-groupe de $G'$. L'image réciproque d'un sous-groupe $H'$ de $G'$ par un morphisme de groupes est un sous-groupe de $G$.
\end{thm}
\begin{cor}
	$\ker(\varphi)$ est un sous-groupe de $G$. $\im(\varphi)$ est est un sous-groupe de $G'$.
\end{cor}
\begin{thm}
	Soit $\varphi$ un morphisme de $G$ dans $G'$. Alors, $\varphi$ est injectif si, et seulement si, $\ker(\varphi)=\left\{e_G\right\}$.
\end{thm}
\begin{exa}
	Un exemple typique d'utilisation est celui où $G$ et $G'$ sont finis de même cardinal. Si on montre que $\ker(\varphi)=\left\{e_G\right\}$, alors $\varphi$ sera injectif, puis bijectif. On pourra ainsi démontrer très efficacement que $\varphi$ est surjectif, ce qui autrement est toujours un peu délicat.
\end{exa}
\begin{prop}
	La composée de deux morphismes est encore un morphisme. La bijection réciproque d'un isomorphisme est un isomorphisme.
\end{prop}
\begin{exa}
	$(Aut(G),\circ)$ (avec $Aut(G)$ l'ensemble des automorphismes de $G$) est un groupe.
\end{exa}
\begin{cor}
	Il existe un isomorphisme de $G$ dans $G'$ si, et seulement si, il existe un isomorphisme de $G'$ dans $G$. De manière condensée, on dira alors que $G$ et $G'$ sont \textbf{isomorphes}. On pourra écrire : $$G\cong G'$$
\end{cor}
\section{Anneaux}
\subsection{Premières définitions}
\begin{dfn}[Anneau]
	Soit $A$ un ensemble muni de deux lois de composition internes $+$ et $\times$. On dit que $(A,+,\times)$ est un anneau lorsque :
	\begin{itemize}
		\item $(A,+)$ est un groupe commutatif.
		\item $\times$ est associative et possède un élément neutre. En revanche, les éléments ne sont $a \ priori$ pas inversibles pour $\times$, et $\times$ n'est pas nécessairement commutative.
		\item $\times$ est \textbf{distributive} à gauche et à droite sur $+$ :
		$$\forall (x,y,z)\in A^3,\ \left \{\begin{array}{ll}x\times(y+z)=x\times y+x\times z \\ (x+y)\times z=x\times z +y\times z\end{array}\right.$$
	\end{itemize}
	Si $\times$ est commutative, on dit que l'anneau est \textbf{commutatif}.
\end{dfn}
\begin{rmq}
	Si on veut montrer que $A$ est un anneau commutatif, on a intérêt à montrer la commutativité avant la distributivité. Ainsi, il suffira de vérifier la distributivité d'un côté seulement.
\end{rmq}
\begin{rmq}[Notations]
	Traditionnellement, l'élément neutre pour $+$ est noté $0$, et l'élément neutre pour $\times$ est noté $1$. Rapidement, s'il n'y a pas d'ambiguïté, on ne notera même plus la loi $\times$ : $x\times y$ sera noté $xy$. S'il existe, l'inverse de $x$ pour la deuxième loi sera noté $x\inv$. Enfin, par "priorité d'opérateur", l'expression $xy+x'y'$ est définie comme étant égale à $(xy)+(x'y')$.
\end{rmq}
\begin{exa}
	$\Z$, $\Q$, $\R$ et $\C$ sont des anneaux commutatifs.
\end{exa}
\begin{exa}
	Soit $f$ et $g$ deux fonctions d'un intervalle $I\subset\R$ dans $\R$. On définit de manière évidente : $$\left \{\begin{array}{ll}f+g:x\mapsto f(x)+g(x)\\fg:x\mapsto f(x)g(x)\end{array}\right.$$ Avec cette notation, $(\R^I,+,\times)$ est un anneau (commutatif). On vérifie que chacune des propriétés s'hérite des mêmes propriétés dans l'anneau $\R$.
\end{exa}
\begin{rmq}
	Plus généralement, si $X$ est un ensemble quelconque et $(A,+,\times)$ un anneau, alors $(A^X,+,\times)$ est un anneau au même titre que $A$.
\end{rmq}
\begin{prop}
	On a : $$\left \{\begin{array}{ll}\forall x\in A,\ 0\times x=x\times0=0\\ \forall(x,y)\in A^2,\ x(-y)=(-x)y=-xy\end{array}\right.$$ On dit que $0$ est \textbf{absorbant}.
\end{prop}
\begin{dfn}[Anneau trivial, anneau non trivial]
	Supposons que $A$ est réduit à un singleton : $A=\left\{x\right\}$. Alors, on a nécessairement : $$\left \{\begin{array}{ll}x+x=x\\ x\times x=x\end{array}\right.$$ Réciproquement, tout singleton $\left\{x\right\}$ muni des deux lois exposées ci-dessus est bien un anneau. On l'appelle \textbf{anneau trivial}. Notons qu'on a nécessairement $0_A=1_A$ et $A=\left\{0_A\right\}$. C'est pourquoi on parle aussi d'\textbf{anneau nul}. Si $A$ n'est pas l'anneau trivial, on dira que $A$ est un \textbf{anneau non trivial}.
\end{dfn}
\begin{rmq}
	Dès que $A$ est non trivial, on a $0\ne1$. En effet, on choisit alors un élément $x\ne0$, on a $x\times 1=x\ne x\times0=0$.
\end{rmq}
Y a-t-il d'autres éléments dont le produit vaut 0 ? Cela dépend et fait l'objet de la définition suivante.
\begin{dfn}[Anneau intègre]
	Un anneau est dit \textbf{intègre} lorsqu'il est non trivial, commutatif, et qu'il ne contient pas de "diviseurs de $0$" : $$\forall(x,y)\in A^2,\ xy=0\implies x=0 \ \text{ou} \ y=0$$
\end{dfn}
\begin{exa}
	$\Z$ est un anneau intègre. En effet, supposons que $(\pm m)(\pm n)=0$ avec $(m,n)\in\N^2$. Alors on a $mn=0$ d'où $m=0$ ou $n=0$.
\end{exa}
\begin{prop}
	Dans un anneau intègre, tout élément non nul est régulier pour la deuxième loi.
\end{prop}
\begin{proof}
	Par commutativité, il suffit de démontrer la régularité d'un seul côté. Supposons donc que $xy=xz$ avec $x$ non nul. Alors $x(y-z)=0$, et comme $x\ne0$, on en déduit $y-z=0$ puis $y=z$.
\end{proof}
\begin{prop}[Anneau-produit]
	De même qu'on a un groupe-produit, on peut définir un anneau-produit pour $n\geqslant2$ anneaux.
\end{prop}
\begin{dfn}[Unités]
	Les \textbf{unités} de $A$ sont les éléments de $A$ inversibles pour la deuxième loi. Leur ensemble est noté $\mathcal{U}(A)$, voire $A^\times$.
\end{dfn}
\begin{prop}[Structure de groupe des unités]
	La loi $\times$ de $A$ induit une structure de groupe sur $\mathcal{U}(A)$. De manière légèrement abusive, on pourra dire que $(\mathcal{U}(A),\times)$ est un groupe.
\end{prop}
\subsection{Calculs algébriques dans un anneau}
\begin{dfn}
	On a rappelle que dans un groupe dont la loi est notée additivement, le $n$-ème itéré de $x$ est noté $nx$. Dans $(A,+)$, on pose donc : $$\left \{ \begin{array}{lll}0_\Z.x=0_A \\ \forall n\in\N, (n+1_\Z)x=nx+x\\ \forall n\in\N\st,\ (-n)x=-(nx)\end{array}\right.$$ Au passage, on notera la priorité d'opérateur : $mx+ny$ désigne en fait $(mx)+(ny)$.
\end{dfn}
\begin{exa}
	Si $A=\Z$, on montre facilement que cette définition coïncide avec la multiplication (c'est-à-dire la deuxième l.c.i.). Si $n\in\N$, c'est une simple récurrence, et si $n<0$, on passe à l'opposé. Plus généralement, ce raisonnement est valable dans tout sur-anneau de $\Z$.
\end{exa}
Cela dit, dans le cas général, lorsqu'on écrit $nx$, il ne s'agit pas de la multiplication dans $A$, mais plutôt d'une loi \textbf{externe} : on compose un élément de $\Z$ et de $A$, qui n'ont pas la même nature, et obtient un élément de $A$. Cependant, on se rend vite compte qu'en termes d'associativité et de distributivité, tout fonctionne comme si on travaillait avec la multiplication dans $A$.
\begin{prop}
	On a les relations suivantes : 
	$$\left \{\begin{array}{lllll}\forall(m,n)\in\Z^2,\ \forall x\in A,\ (mn)x=m(nx)& (\text{pseudo-associativité}) \\ \forall n\in\Z,\ \forall(x,y)\in A^2,\ n(x+y)=nx+ny & (\text{pseudo-distributivité à gauche})\\ \forall(m,n)\in\Z^2,\ \forall x\in A, (m+n)x=mx+nx &(\text{pseudo-distributivité à droite})\\ \forall x\in A,\ 1_\Z\cdot x=x &(\text{opérateur neutre})\\ \forall n\in\Z,\ \forall(x,y)\in A^2,\ (nx)\times y=x\times(ny)=n(x\times y)&(\text{compatibilité avec le produit})\end{array}\right.$$
\end{prop}
\begin{rmq}
	La dernière propriété permet d'écrire $nxy$ sans ambiguïté.
\end{rmq}
\begin{exa}
	On en déduit que le produit est bilinéaire au sens suivant : $$\left \{\begin{array}{ll}\forall(m,n)\in\Z^2,\ \forall(x,y,z)\in A^3,\ x(my+nz)=m(xy)+n(xz)\\ \forall(m,n)\in\Z^2,\ \forall(x,y,z)\in A^3,\ (mx+ny)z=m(xz)+n(yz)\end{array}\right.$$
\end{exa}
\begin{prop}[Pseudo-absorbance et pseudo-règle des signes]
	Soit $n\in\Z$ et $x\in A$. On a : $$\left \{\begin{array}{lll}0_\Z. x=n.0_A=0_A \\ (-n).x=n.(-x)=-(n.x)\\ (-n).(-x)=n.x\end{array}\right.$$
\end{prop}
\begin{dfn}[Itération de la loi $\times$]
	On rappelle qu'on définit : $$\left \{\begin{array}{ll}x^0=1_A \\ \forall n\in\N,\ x^{n+1}=x^nx\end{array}\right.$$ Mais attention, en général, on ne pourra pas définir $x^n$ pour $n<0$, sauf si $x$ est inversible pour $\times$, auquel cas on retrouve la définition de $x^n$ dans le groupe $\mathcal{U}(A)$.
\end{dfn}
\begin{dfn}[Élément nilpotent, indice de nilpotence]
	Un élément $x$ de l'anneau $A$ est \textbf{nilpotent} s'il existe $n\in\N$ tel que $x^n=0$. Dans ce cas, le plus petit $n$ vérifiant cette égalité s'appelle l'\textbf{indice de nilpotence} de $x$.
\end{dfn}
\begin{exa}
	On rencontrera les exemples les plus classiques dans le chapitre "Matrices". En attendant, en anticipant un tout petit peu sur l'arithmétique, un premier exemple est donné par l'anneau $\Z/4\Z$ des entiers comptés modulo $4$, où $2$ est nilpotent d'indice $2$.
\end{exa}
\begin{prop}[Distributivité généralisée]
	Soit $(a_i)_{1\leqslant i\leqslant m}$ et $(b_j)_{1\leqslant j\leqslant n}$ deux familles finies à valeurs dans un anneau $A$. Alors on a $$\left(\sum_{i=1}^{m}a_i\right)\left(\sum_{j=1}^{n}b_j\right)=\sum_{(i,j)\in\llbracket1,m\rrbracket\times\llbracket1,n\rrbracket}a_ib_j$$
\end{prop}
\begin{prop}
	On peut généraliser la formule précédente à un nombre quelconque (fini) de sommes, et on obtient : $$\left(\sum_{i_1\in I_1}a_{i_1,1}\right)\times\ldots\times\left(\sum_{i_n\in I_n}a_{i_n,n}\right)=\sum_{(i_1,\ldots,i_n)\in I_1\times\ldots\times I_n}a_{i_1,1}\times\ldots\times a_{i_n,n}$$
\end{prop}
\begin{proof}
	Il s'agit d'un récurrence facile sur le nombre de sommes en utilisant la propriété précédente.
\end{proof}
\begin{cor}[Développement d'un carré]
	Dans tout anneau, on a : $$\left(\sum_{i=1}^{n}a_i\right)^2=\sum_{i=1}^{n}a_i^2+\sum_{i\ne j}a_ia_j$$ Si de plus les $a_i$ commutent deux à deux, on a : $$\left(\sum_{i=1}^{n}a_i\right)^2=\sum_{i=1}^{n}a_i^2+2\sum_{i<j}a_ia_j$$
\end{cor}
Dans le cas où deux éléments $a$ et $b$ commutent, on a deux résultats importants.
\begin{thm}
	Si $a$ et $b$ des éléments de $A$ \textbf{commutent}, alors on a pour $n\in\N$:
	\begin{itemize}
		\item Formule du binôme de Newton : $$(a+b)^n=\sum_{k=0}^{n}\binom{n}{k}a^kb^{n-k}$$
		\item Formule de Bernoulli : $$a^n-b^n=(a-b)\sum_{k=0}^{n-1}a^{n-1-k}b^k=\left(\sum_{k=0}^{n-1}a^{n-1-k}b^k\right)(a-b)$$
	\end{itemize}
\end{thm}
\subsection{Sous-anneaux}
\begin{dfn}[Sous-anneau]
	Un \textbf{sous-anneau} d'un anneau $A$ est une partie $B\subset A$ telle que 
	\begin{itemize}
		\item $B$ est un sous-groupe de $A$
		\item $1\in B$
		\item $\forall(x,y)\in B^2,\ xy\in B$
	\end{itemize}
	Dans ce cas, comme on le vérifie immédiatement, $B$ muni des lois induites est lui-même un anneau.
\end{dfn}
\begin{prop}[Recette du sous-anneau]
	Soit $A$ un anneau et $B\subset A$. Pour que $B$ soit un sous-anneau de $A$, il faut et il suffit que \begin{itemize}
		\item $1\in B$
		\item $\forall(x,y)\in B^2,\ x-y\in B$
		\item $\forall(x,y)\in B^2,\ xy\in B$
	\end{itemize}
\end{prop}
\begin{rmq}
	De même que pour les sous-groupes, on utilise principalement la condition suffisante, qui permet de montrer qu'un ensemble $B$ est un anneau en l'interprétant comme un sous-anneau d'un anneau de référence.
\end{rmq}
\begin{exa}
	Tout sous-anneau d'un anneau commutatif est commutatif. Tout sous-anneau d'un anneau intègre est intègre.
\end{exa}
\begin{prop}
	Soit $\mathcal{B}$ un ensemble de sous-anneaux de $A$. Alors $\displaystyle \bigcap_{B\in\mathcal{B}}B$ est encore un sous-anneau de $A$.
\end{prop}
\begin{dfn}[Sous-anneau engendré]
	Si $X$ est une partie d'un anneau $A$, le \textbf{sous-anneau engendré par $X$} est l'intersection de tous les sous-anneaux contenant $X$. C'est le plus petit sous-anneau contenant $X$.
\end{dfn}
\begin{exa}
	Dans $\Q$, le sous-anneau engendré par $\displaystyle\frac12$ est égal à $\displaystyle\left\{\frac{k}{2^n}\ | \ (k,n)\in\Z\times\N\right\}$.
\end{exa}
\subsection{Morphismes}
\begin{dfn}[Morphisme d'anneaux]
	Un \textbf{morphisme} $\varphi$ d'un anneau $A$ dans un anneau $B$ est une application qui transporte les deux lois, ainsi que l'élément neutre de la deuxième loi : $$\left \{\begin{array}{lll}\forall(x,y)\in A^2,\ \varphi(x+y)=\varphi(x)+\varphi(y) \\ \forall(x,y)\in A^2,\ \varphi(xy)=\varphi(x)\varphi(y) \\ \varphi(1_A)=1_B\end{array}\right.$$
\end{dfn}
\begin{prop}
	$\varphi$ induit un morphisme du groupe $\mathcal{U}(A)$ dans le groupe $\mathcal{U}(B)$.
\end{prop}
\begin{proof}
	En effet, si $x\in\mathcal{U}(A)$, alors $\varphi(x)$ est inversible, d'inverse $\varphi(x\inv)$.
\end{proof}
\begin{rmq}
	En particulier, on voit que $\varphi$ est un morphisme du groupe $(A,+)$ dans le groupe $(B,+)$. Donc $\ker(\varphi)$ est toujours défini comme $\varphi\inv\left(\left\{0_B\right\}\right)$, et il pourra toujours servir à montrer que $\varphi$ est injectif (si, et seulement si, $\ker(\varphi)=\left\{0_A\right\}$). Attention à ne pas confondre les éléments neutres et à ne pas écrire $\ker(\varphi)=\varphi\inv\left(\left\{1_B\right\}\right)$ !
\end{rmq}
\begin{prop}
	Soit $\varphi:A\rightarrow B$ un morphisme d'anneaux. On a les résultats suivants.
	\begin{itemize}
		\item Itération de la loi $+$ : $$\forall x\in A,\ \forall n\in\Z,\ \varphi(nx)=n\varphi(x)$$
		\item Itération de la loi $\times$ : $$\forall x\in A,\ \forall n\in\N,\ \varphi(x^n)=\varphi(x)^n$$
	\end{itemize}
\end{prop}
\begin{proof}
	Le premier point a déjà été prouvé. Le deuxième se démontre par une récurrence facile.
\end{proof}
\begin{prop}
	L'image d'un sous-anneau par un morphisme d'anneaux est un sous-anneau. L'image réciproque d'un sous-anneau par un morphisme d'anneaux est un sous-anneau. On en déduit que le noyau et l'image d'un morphisme d'anneaux sont des sous-anneaux.
\end{prop}
\begin{prop}
	La composée de deux morphismes d'anneaux est un morphisme d'anneaux. La bijection réciproque d'un isomorphisme d'anneaux est un isomorphisme d'anneaux.
\end{prop}
\section{Corps}
\subsection{Premier exemple : construction de $\Q$, HP}
Après avoir construit $\N$, on a éprouvé le besoin de construire $\Z$ parce que $\N$ "ne passait pas à l'inverse" pour l'addition. Pourrait-on effectuer le même travail sur la multiplication ? C'est ce qu'on va voir maintenant en construisant $\Q$. De même qu'on a écrit tout entier relatif comme différence de deux entiers naturels, on va écrire tout rationnel comme quotient de deux entiers relatifs. A un rationnel, on associe donc un couple d'entiers.
\begin{prop}
	Sur l'ensemble $\Z\times\Z\st$, on considère la relation définie par $$(p,q)\sim(p',q')\iff pq'=qp'$$ C'est une relation d'équivalence.
\end{prop}
\begin{dfn}[Ensemble $\Q$]
	On note $\Q$ l'ensemble $\Z\times\Z\st$ quotienté par $\sim$. Ses éléments sont appelés les \textbf{rationnels}, il seront notés $\overline{(p,q)}$, ou encore plus simplement $\displaystyle\frac{p}{q}$. En particulier, si $k\in\Z\st$, on peut "simplifier en haut et en bas par $k$" : $$\frac{pk}{qk}=\frac{p}{q}$$
\end{dfn}
\begin{prop}[Addition sur $\Q$]
	On définit une addition sur $\Q$ par : $$\frac{p}{q}+\frac{r}{s}=\frac{ps+qr}{qs}$$ Cette opération est bien définie.
\end{prop}
\begin{rmq}
	On a $\displaystyle\frac{a}{n}+\frac{b}{n}=\frac{an+nb}{n^2}=\frac{a+b}{n}$.
\end{rmq}
\begin{prop}[Multiplication sur $\Q$]
	On définit une multiplication sur $\Q$ par : $$\frac{p}{q}\times\frac{r}{s}=\frac{pr}{qs}$$ Cette opération est bien définie.
\end{prop}
\begin{prop}[Inclusion abusive $\Z\subset\Q$]
	L'application $n\mapsto\displaystyle\frac{n}{1}$ de $\Z$ dans $\Q$ est injective. Elle permet donc de voir $\Z$ comme un sous-ensemble de $\Q$. Avec cette inclusion, l'addition et la multiplication sur $\Q$ prolongent celles sur $\Z$.
\end{prop}
\begin{thm}
	$(\Q,+,\times)$ est un anneau commutatif.
\end{thm}
La grande nouveauté de $\Q$ par rapport à $\Z$ réside dans l'inversibilité pour la deuxième loi.
\begin{thm}
	Tout rationnel non nul est inversible pour la multiplication.
\end{thm}
\begin{thm}
	On peut prolonger $\leqslant$ à $\Q$ tout entier. Le relation d'ordre obtenue reste totale, et compatible avec l'addition et la multiplication au même sens que dans $\Z$.
\end{thm}
\begin{proof}
	Ce théorème est démontré dans les compléments.
\end{proof}
\subsection{Généralisation}
\begin{dfn}[Corps]
	Un \textbf{corps} $(K,+,\times)$ est un anneau commutatif, non trivial, tel que tout élément non nul est inversible pour la deuxième loi. On note $K\st=K\setminus\left\{0\right\}$.
\end{dfn}
\begin{rmq}
	Certaines sources un peu anciennes ne mentionnent pas la commutativité. Mais progressivement, cela tend à disparaître, afin de s'aligner sur la terminologie anglo-saxonne (selon laquelle tout corps est commutatif, la traduction exacte étant le mot $field$). Un "corps" dont la deuxième loi ne serait pas nécessairement commutative est maintenant appelé un \textbf{corps gauche} ($skew$ $field$ en anglais), et la terminologie semble bien stabilisée.
	\par De toute manière, le théorème de Wedderburn montre que s'il est fini, tout corps gauche est commutatif. Il faut donc aller chercher des corps gauches infinis pour espérer rencontrer un cas de non-commutativité. Si on le souhaite, on pourra s'intéresser aux quaternions (cf. TD Matrices).
\end{rmq}
\begin{rmq}
	Réciproquement, si $x$ est inversible, alors $x\ne0$, puisque $\forall y\in K,\ 0\times y=0\ne1$. Donc $(K\st,\times)$ est égal au groupe des unités de $K$. En particulier, c'est un groupe.
\end{rmq}
\begin{exa}
	$\Q$, $\R$ et $\C$ sont des corps.
\end{exa}
\begin{dfn}[Corps des fractions, HP]
	Le travail que nous avons effectué pour passer de $\Z$ à $\Q$ peut s'effectuer avec n'importe quel anneau intègre $A$. En quotientant $A\times (A\setminus\left\{0\right\})$ par la relation définie par : $$(a,b)\sim(c,d)\iff ad=bc$$ on obtient ce qu'on appelle le \textbf{corps des fractions de $A$}, et $A$ se plonge naturellement dedans. Plus précisément, si $(a,b)\in A\times (A\setminus\left\{0\right\})$, la classe de $(a,b)$ sera plutôt notée $\displaystyle\frac{a}{b}$, et tout élément $a\in A$ pourra être identifié à $\displaystyle\frac{a}{1}$.
\end{dfn}
\begin{proof}
	Ce fait sera développé dans des compléments (de ce chapitre, ou du chapitre "Arithmétique", je ne sais pas encore)
\end{proof}
\begin{rmq}[Localisation d'un anneau, HP]
	Le corps des fractions est un cas particulier de ce que l'on appelle la \textbf{localisation d'un anneau} qui consiste à rendre inversible de force une partie stable par multiplication d'un anneau. Toutefois, le plongement n'est alors pas injectif, à moins que l'anneau de départ soit intègre. Si l'anneau de départ $A$ est intègre et qu'on effectue ce travail sur la partie $A\setminus\left\{0\right\}$, on obtient le corps des fractions de $A$.
\end{rmq}
\begin{rmq}[Calcul fractionnaire]
	Si nous prenons un anneau intègre $A$ et que nous construisons son corps des fractions $K$, rappelons que si $(a,b)\in A\times (A\setminus\left\{0\right\})$, on a $$\left \{\begin{array}{ll}\displaystyle\frac{1}{b}\times b=\frac{b}{b}=1 \ \text{donc} \ \frac{1}{b}=b\inv \\[3mm] \displaystyle\frac{a}{b}=\frac{a\times 1}{1\times b}=\frac{a}{1}\times\frac{1}{b}=ab\inv=b\inv a\end{array}\right.$$ Plus généralement, dans un corps (et notamment dans $\Q$,$\R$ ou $\C$), l'inverse d'un élément $y\ne0$ pourra être noté $\displaystyle\frac{1}{y}$, et l'écriture $\displaystyle\frac{x}{y}$ désignera indifféremment : $$x\times\frac{1}{y}=\frac{1}{y}\times x$$ On vérifie immédiatement que les formules bien connus $$\frac{x}{y}+\frac{x'}{y'}=\frac{xy'+x'y}{yy'} \ \ , \ \ \frac{x}{y}\times\frac{x'}{y'}=\frac{xx'}{yy'} \ \ \text{et} \ \ x\times\frac{x'}{y'}=\frac{xx'}{y}$$ sont justifiées. On en déduit que si $\displaystyle\frac{x}{y}\ne0$, alors $\displaystyle \left(\frac{x}{y}\right)\inv=\frac{y}{x}$.
\end{rmq}
\begin{prop}
	Tout corps est un anneau intègre.
\end{prop}
\begin{dfn}[Sous-corps]
	Un \textbf{sous-corps} d'un corps $K$ est une partie $L\in K$ telle que
	\begin{itemize}
		\item $L$ est un sous-anneau de $K$
		\item $\forall x\in L,\ x\ne0\implies x\inv\in L$ (stabilité par passage à l'inverse d'un élément non nul)
	\end{itemize}
	Dans ce cas, comme on le vérifie immédiatement, $L$ est lui-même un corps.
\end{dfn}
\begin{prop}[Recette du sous-corps]
	Soit $K$ un anneau et $L\subset K$. Pour que $L$ soit un sous-corps de $K$, il faut et il suffit que \begin{itemize}
		\item $1\in L$
		\item $\forall(x,y)\in L^2,\ x-y\in L$
		\item $\forall(x,y)\in L^2,\ y\ne0\implies xy\inv\in L$
	\end{itemize}
\end{prop}
\begin{rmq}
	De même que pour les sous-groupes, on utilise principalement la condition suffisante, qui permet de montrer qu'un ensemble $L$ est un corps en l'interprétant comme un sous-corps d'un corps de référence.
\end{rmq}
\begin{prop}
	Soit $\mathcal{L}$ un ensemble de sous-corps de $K$. Alors $\displaystyle\bigcap_{L\in\mathcal{L}}L$ est encore un sous-corps de $K$.
\end{prop}
\begin{dfn}[Sous-corps engendré]
	Si $X$ est une partie de $K$, le \textbf{sous-corps engendré par $X$} est l'intersection de tous les sous-corps contenant $X$. C'est le plus petit sous-corps contenant $X$.
\end{dfn}
\begin{exa}
	Dans $\R$, le sous corps engendré par $1$ est $\Q$. En effet, il contient $1$ donc par stabilité, une récurrence immédiate montre qu'il contient $\N$, puis $\Z$. Par passage à l'inverse, il contient donc les nombres de la forme $\displaystyle\frac{1}{n}$, puis par stabilité il contient $\Q$. Comme $\Q$ est un corps, c'est le plus petit sous-corps de $\R$ qui contient $1$.
\end{exa}
\begin{dfn}[Morphisme de corps]
	Un \textbf{morphisme} de corps est un morphisme pour les anneaux sous-jacents.
\end{dfn}
\begin{exa}
	Dans $\C$, la conjugaison est un automorphisme involutif de corps.
\end{exa}
\begin{rmq}
	On peut montrer sans difficulté que l'image d'un sous-corps par un morphisme de corps est un sous-corps. De même avec l'image réciproque.
\end{rmq}
\begin{rmq}
	Attention ! Contrairement au cas des groupes et des anneaux, le produit cartésien de deux corps $K_1$ et $K_2$ n'est pas un corps ! En effet, $(1_{K_1},0_{K_2})$ et $(0_{K_1},1_{K_2})$ ne peuvent pas être inversés.
\end{rmq}
\begin{prop}[HP, mais à connaître sur le bout des doigts]
	Tout morphisme de corps est injectif.
\end{prop}
\begin{proof}
	Soit $\varphi$ un morphisme du corps $K$ vers le corps $L$. Raisonnons par l'absurde : supposons que $\varphi$ n'est pas injectif. On peut alors fixer $x \in \ker(\varphi) \setminus \{0_K\}$. Or, $x$ est non nul, donc on a :
	$$1_L=\varphi(1_K)=\varphi(xx\inv)=\varphi(x)\varphi(x\inv)=0_L$$ par propriétés des morphismes de corps ainsi que par absorbance. Or, nos corps sont toujours supposés non triviaux, donc ceci est \textbf{absurde}. Ainsi, $\varphi$ est bien injectif, ce qui conclut.
\end{proof}
\begin{prop}[HP, mais à connaître]
	Tout anneau intègre fini est un corps.
\end{prop}
\begin{proof}
	Se donner un anneau intègre fini $A$ et $a$ un élément non nul de $A$. Considérer l'application
	$$
	\begin{array}[t]{lrcl}
		\varphi_a : & A & \longrightarrow & A \\
		& x & \longmapsto & ax
	\end{array}
	$$ Déjà, cette application est bien définie par stabilité de l'anneau par la deuxième loi. Ensuite, elle injective par intégrité de $A$. Enfin, elle est alors surjective par égalité de cardinaux. Donc on peut fixer un antécédent de $1_A$ par $\varphi_a$, ce qui montre que $a$ est inversible. Ainsi, $A$ est bien un corps car seule l'inversibilité des éléments non nuls manquait.
\end{proof}
\section{Compléments sur les groupes et les corps, HP}
\subsection{Groupes finis}
Pour commencer, voici la démonstration du théorème de Lagrange.
\begin{thm}[Théorème de Lagrange]
	Soit $G$ un groupe fini, et $H$ un sous-groupe de $G$. Alors, $\card(H)$ divise $\card(G)$.
\end{thm}
\begin{proof}
	Introduisons la relation sur $G$ : $$\forall(x,y)\in G^2,\ x\sim y \iff x\inv y \in H$$ et vérifions que c'est une relation d'équivalence. Soit $(x,y,z)\in G^3$.
	\begin{itemize}
		\item Réflexivité : puisque $e\in H$, on a bien $x\sim x$.
		\item Symétrie : si $x\inv y\in H$, alors par passage à l'inverse $(x\inv y)\inv=y\inv x\in H$
		\item Transitivité : si $x\inv y\in $ et $y\inv z\in H$, alors par stabilité $(x\inv y)(y\inv z)=x\inv z\in H$ 
	\end{itemize}
	Ainsi, $\sim$ partitionne $g$ en classes d'équivalence. Ces classes sont en nombre fini, car $G/\sim$ est l'image directe de $G$ par la surjection canonique. Notons $n$ le nombre de classes. \par Par ailleurs, si $\overline{x}$ est une de ces classes d'équivalence, montrons qu'elle est égale à $xH$. Soit $y\in G$. On a $$y\in\overline{x}\iff\exists h\in H,\ x\inv y=h \iff \exists h\in H,\ y=xh \iff y\in xH$$ Considérons alors l'application $$\begin{array}{ccccc} \varphi &:&H&\rightarrow&xH \\ &&h&\mapsto&xh\end{array}$$
	\begin{itemize}
		\item $\varphi$ est injective par régularité de $x$
		\item $\varphi$ est surjective par définition
	\end{itemize}
	Donc $\varphi$ est bijective, si bien que $\card(xH)=\card(H)$. Ainsi, toutes les classes d'équivalence ont le même cardinal, à savoir $\card(H)$. Finalement $$\card(G)=n\times\card(H)$$d'où le résultat.
\end{proof}
\subsection{Retour sur la construction de $\Q$}
Nous souhaitons maintenant prolonger $\leqslant_\Z$ à $\Q$. Pour cela, on sait bien intuitivement "quand une fraction est négative" : c'est que son numérateur et son dénominateur n'ont pas le même signe. Autrement dit, c'est quand le produit du numérateur est strictement négatif. De même, on sait "quand une fraction est plus petite qu'une autre : c'est que leur différence est négative ou nulle.
\begin{prop}
	On choisit de définir $\displaystyle \frac{p}{q}\leqslant\frac{r}{s}$ par $(ps-qr)qs\leqslant_\Z 0$. Cette définition est consistante, et $\leqslant$ est une relation d'ordre totale qui prolonge $\leqslant$ sur $\Z$. Enfin, on garde les mêmes compatibilités avec l'addition et la multiplication.
\end{prop}
\begin{proof}
	Prouvons les résultats dans l'ordre de l'énoncé.
	\begin{itemize}
		\item Premier point : supposons qu'on ait plusieurs écritures du type $\frac{p}{q}=\frac{p'}{q'}$ et $\frac{r}{s}=\frac{r'}{s'}$, et montrons que $$(ps-qr)qs\leqslant0 \iff (p's'-q'r')q's'\leqslant0$$ Par symétrie, il suffit de montrer le sens direct, le sens direct se démontrant en échangeant les variables. Supposons donc que $(ps-qr)qs\leqslant0$, alors 
		\begin{align*}
			(ps-qr)qsq'^2s'^2 &\text{donc} \ ps^2qq'^2s'^2-q^2rsq'^2s'^2\leqslant0 \\
			&\text{donc} \ p's^2q^2q's'^2-q^2r's^2q'^2s'\leqslant0 \\
			&\text{donc} \ (p's'^2q'-r's'q'^2)q^2s^2\leqslant0 \\
			&\text{donc} \ (p's'-q'r')q's'\leqslant0
		\end{align*}
		Voilà pour la consistance de la définition.
		\item Conséquence : si on veut comparer deux rationnels, en les "mettant au même dénominateur" on peut toujours se ramener à comparer $\displaystyle\frac{p}{q}$ et $\displaystyle\frac{p'}{q}$ avec $q>0$, et l'inégalité $\displaystyle\frac{p}{q}\leqslant\frac{p'}{q}$ est alors équivalente à $(pq-p'q)q^2\leqslant0$, soit encore $p\leqslant p'$.
		\item Axiomes d'une relation d'ordre totale : ils viennent directement du point précédent.
		\item Le prolongement de $\leqslant_\Z$ est immédiat : $$\frac{m}{1}\leqslant\frac{n}{1}\iff(m\times 1-1\times n)\times1\times1\leqslant0\iff m\leqslant n$$
		\item Compatibilité avec l'addition : il s'agit de montrer l'assertion du type $$\frac{p}{q}\leqslant\frac{p'}{q'}\implies \frac{p}{q}+\frac{r}{s}\leqslant\frac{p'}{q'}+\frac{r}{s}$$Là aussi, on suppose sans perte de généralité de $q=s$ et le résultat est alors immédiat.
		\item Le raisonnement est le même avec la multiplication.
	\end{itemize}
Ainsi, la preuve est achevée.
\end{proof}
\subsection{Caractéristique d'un corps}
Dans ce paragraphe, on se place dans un corps $K$, et on note $n\cdot x$ la $n$-ème itérée de $x$ au sens de l'addition (avec $n\in\Z$ et $x\in K$).
\begin{dfn}[Caractéristique]
	S'il existe, le plus petit entier $p\in\N\st$ tel que $p\cdot 1_K=0_K$ est appelé la \textbf{caractéristique de $K$}. S'il n'existe pas, on dit par convention que $K$ est de caractéristique nulle.
\end{dfn}
\begin{exa}
	$\R$ et $\C$ sont de caractéristique nulle. $\Z/p\Z$, où $p$ est premier, est un corps de caractéristique $p$ (cf. le chapitre "Arithmétique").
\end{exa}
\begin{thm}
	Un corps fini est de caractéristique non nulle.
\end{thm}
\begin{proof}
	Soit $n$ le cardinal du corps. Alors dans l'ensemble $\left\{m\cdot 1_K \ | \ m\in\llbracket0,n\rrbracket\right\}$, au moins deux éléments sont égaux par le principe des tiroirs. On peut donc écrire $m\cdot1_K=m'\cdot1_K$ avec $m<m'$, d'où on en déduit $(m'-m)\cdot1_K=0_K$. L'ensemble $\left\{q\in\N\st\ | \ q\cdot1_K=0_K\right\}$ est alors non vide, et il admet un plus petit élément. 
\end{proof}
\begin{thm}
	Si la caractéristique d'un corps est non nulle, alors c'est un nombre premier.
\end{thm}
\begin{proof}
	Notons $p$ la caractéristique du corps, et supposons que $p\ne0$. Si $p$ n'était pas premier, il suffirait d'écrire $p=ab$, $1\leqslant a,b <p$. On aurait alors $$(a\cdot1_K)\times(b\cdot1_K)=0_K$$ Comme $K$ est intègre, on aurait $a\cdot1_K=0_K$ ou $b\cdot1_K=0_K$, en contradiction avec la minimalité de $p$. Donc $p$ est un nombre premier.
\end{proof}
\begin{prop}
	Soit $K$ de caractéristique $p>0$, $x\in K\st$, et $m\in\Z$. Alors $m\cdot x=0_K$ si, et seulement si, $p$ divise $m$.
\end{prop}
\begin{proof}
	Il suffit de montrer que $x$ est d'ordre $p$ dans le groupe $(K,+)$.
	\begin{itemize}
		\item D'une part, par les opérations usuelles, on 
		\begin{align*}
			p\cdot x&=p\cdot(1_K\times x) \\
			&= (p\cdot1_K)\times x \\
			&=0_K\times x \\
			&= 0_K
		\end{align*}
		\item D'autre part, soit $m\in\llbracket1,p-1\rrbracket$. Toujours par les mêmes règles, on a $$m\cdot x=(m\cdot1_K)\times x$$ qui est non nul par intégrité.
	\end{itemize}
Ainsi, la preuve est achevée.
\end{proof}
\newpage
\chapter{Arithmétique}
L'arithmétique, ou science des nombres entiers, trouve ses origines chez les Babyloniens côté occidental et chez les Indiens et les Chinois côté oriental. C'est donc une branche très ancienne des mathématiques, qui a fourni et fournit encore certains des plus beaux problèmes (le grand théorème de Fermat, résolu seulement en 1994 par Andrew Wiles et son ancien étudiant Richard Taylor). La plupart comportent un énoncé extrêmement simple, alors que leur résolution, quand elle est connue, est souvent excessivement ardue et fait appel à des domaines $a$ $priori$ très éloignés (analyse complexe, courbes elliptiques, etc.). A titre d'exemple, on citera la célèbre "conjecture de Goldbach" qui, comme son nom l'indique, n'est toujours pas résolue à l'heure où M. Morlot écrit ce cours (été 2013) ni à celle où je tape ce cours (été 2024) : est-ce que tout nombre pair plus grand que $2$ peut s'écrire comme somme de deux nombres premiers ? Il semblerait bien que oui, en tout cas c'est ce que montrent des tests menés sur ordinateur jusqu'à des bornes impressionnantes. Mais quant à démontrer le théorème en toute généralité, il faut encore attendre... Plus généralement, les nombres premiers sont un thème central de la recherche en théorie des nombres, et leur réalité sous-jacente n'est encore que très partiellement comprise.
\section{Divisibilité}
\begin{lem}
	Les unités de $\Z$ sont $-1$ et $1$.
\end{lem}
\begin{dfn}[Divisibilité]
	Soit $(a,d)\in\Z^2$. Lorsqu'il existe $k\in\Z$ tel que $a=kd$, on dit que $a$ est un \textbf{multiple} de $d$ et que $d$ \textbf{divise} $a$. On écrit $d\mid a$.
\end{dfn}
\begin{exa}
	Les diviseurs de $1$ sont $1$ et $-1$.
\end{exa}
\begin{exa}
	Soit $a,d,\lambda\in\Z$. Si $d\mid a$, alors $\lambda d\mid\lambda a$. Si $\lambda\ne0$, la réciproque est également vraie.
\end{exa}
\begin{rmq}
	Si $d$ divise $a$ et $b$, alors $d$ divise toute $\Z$-combinaison linéaire $au+bv$ avec $(u,v)\in\Z^2$.
\end{rmq}
\begin{rmq}[Notations]
	L'ensemble des multiples de $d$ est donc égal à $d\Z$. Quant à l'ensemble des diviseurs de $a$, il est souvent noté $\mathcal{D}(a)$, et l'ensemble de ses diviseurs strictement positifs est souvent noté $\mathcal{D}^+(a)$. Cependant, il peut être prudent de rappeler ces notations avant de les utiliser.
\end{rmq}
\begin{rmq}
	Soit $d\in\Z$. $d\Z$ est égal à l'ensemble des itérés de $d$ pour la loi $+$. Il s'agit donc du sous-groupes de $(\Z,+)$ engendré par $d$.
\end{rmq}
\begin{prop}
	Sur $\N$, la divisibilité est une relation d'ordre partielle.
\end{prop}
\begin{rmq}
	Sur $\Z$, elle reste réflexive et antisymétrique.
\end{rmq}
\begin{rmq}
	$1$ est le plus petit élément pour cette relation d'ordre (car il divise tous les entiers naturels), et $0$ est le plus grand élément (car tout entier naturel divise $0$).
\end{rmq}
\begin{rmq}[Lien avec la relation d'ordre $\leqslant$]
	Dans $\N\st$, si $m\mid n$, alors $m\leqslant n$. En pratique, il suffit de vérifier que $n\in\N\st$.
\end{rmq}
Nous en venons à la définition du pgcd et du ppcm de deux entiers. Dans un premier temps, restreignons-nous au cas où les entiers sont non nuls.
\begin{dfn}[Pgcd, ppcm]
	Soit $a$ et $b$ deux entiers non nuls.
	\begin{itemize}
		\item 
		L'ensemble des diviseurs communs strictement positifs de $a$ et $b$ est non vide, puisqu'il contient au moins $1$. De plus, il est majoré (par exemple par $\lvert a\rvert$). Il admet donc un plus grand élément noté $a\wedge b$. On l'appelle le \textbf{plus grand commun diviseur} (ou \textbf{pgcd}) de $a$ et $b$.
		\item L'ensemble des multiples communs strictement positifs de $a$ et $b$ est non vide, puisqu'il contient au moins $\lvert ab\rvert$. Il admet donc un plus petit élément noté $a\vee b$. On l'appelle \textbf{plus petit commun multiple} (ou \textbf{ppcm}) de $a$ et $b$.
	\end{itemize}
\end{dfn}
On rappelle que pour tout $n\in\Z$, $n\Z$ est un sous-groupe de $\Z$ (c'est le sous-groupe engendré par $n$).
\begin{thm}[Sous-groupes de $\Z$]
	Réciproquement, pour tout sous-groupe $H$ de $\Z$, il existe un unique $n\in\N$ tel que $H=n\Z$.
\end{thm}
\begin{cor}
	Soit $a$ et $b$ des entiers non nuls. Alors $$\left \{\begin{array}{ll}a\Z+b\Z=(a\wedge b)\Z \\ a\Z\cap b\Z=(a\vee b)\Z\end{array}\right.$$
\end{cor}
Ce corollaire justifie la définition suivante.
\begin{dfn}
	On choisit de prolonger les notions de pgcd et de ppcm dans le cas où $a$ et $b$ sont deux entiers quelconques. Puisque $a\Z+b\Z$ et $a\Z\cap b\Z$ sont des sous-groupes additifs de $\Z$, on décide que le pgcd $a\wedge b$ est l'unique $\delta\in\N$ tel que $a\Z+b\Z=\delta\Z$, et que le ppcm $a\vee b$ est l'unique $\mu\in\N$ tel que $a\Z\cap b\Z=\mu\Z$.
\end{dfn}
\begin{exa}
	On pourra montrer que $\forall n \in\Z,\ n\wedge0=\lvert n\rvert$ et $n\vee0=0$. De même, $n\wedge1=1$ et $n\vee1=\lvert n\rvert$.
\end{exa}
\begin{rmq}
	En reprenant le principe des démonstrations des théorèmes précédents, on constate les faits suivants.
	\begin{itemize}
		\item D'une part, $a\wedge b$ reste un diviseur commun de $a$ et $b$, et $a\vee b$ reste un multiple commun de $a$ et $b$. En revanche, ils n'ont plus de raison d'être maximal et minimal au sens de $\leqslant$.
		\item D'autre part, par définition, $\exists(u,v)\in\Z^2,\ au+bv=a\wedge b$
	\end{itemize}
\end{rmq}
\begin{rmq}[Utile pour les exercices calculatoires]
	On ne change pas $a\wedge b$ et $a\vee b$ :
	\begin{itemize}
		\item En remplaçant $a$ par $-a$ ou $b$ par $-b$. En effet, $a\Z=(-a)\Z$ et $b\Z=(-b)\Z$ ;
		\item En échangeant $a$ et $b$. En effet, $a\Z+b\Z=b\Z+a\Z$ et $a\Z\cap b\Z=b\Z\cap a\Z$. Ainsi, le pgcd et le ppcm sont \textbf{commutatifs}.
	\end{itemize}
	De même, comme on le vérifie immédiatement, $\wedge$ et $\vee$ sont \textbf{associatifs}. C'est une conséquence immédiate de l'associativité de $+$ et $\cap$ dans $\mathcal{P}(\Z)$.
\end{rmq}
\begin{prop}
	Soit $(a,b)\in\Z^2$. Les diviseurs communs de $a$ et $b$ sont exactement les diviseurs de $a\wedge b$. Les multiples communs de $a$ et $b$ sont exactement les multiples de $a\vee b$.
\end{prop}
\begin{dfn}[Pgcd et ppcm de plus de deux entiers]
	 Soit $(a_1,\ldots,a_n)\in\Z^n$. Pour définir leur pgcd et leur ppcm, on peut le faire par récurrence en se plaçant dans les structures algébriques $(\Z,\wedge)$ et $(\Z,\vee)$. De même, on peut définir par récurrence $a_1\Z+\ldots+a_n\Z$ et $a_1\Z\cap\ldots\cap a_n\Z$ en se plaçant dans les structures algébriques $(\mathcal{P}(\Z),+)$ et $(\mathcal{P}(\Z),\cap)$. Et alors, par une récurrence immédiate, on s'aperçoit que tout revient à fonctionner de la manière suivante.
	 \begin{itemize}
	 	\item Le pgcd $a_1\wedge\ldots\wedge a_n$ est l'unique $\delta\in\N$ tel que tel que $a_1\Z+\ldots+a_n\Z=\delta\Z$. Il divise chacun des $a_i$, et l'ensemble des diviseurs communs des $a_i$ est exactement égal à l'ensemble des diviseurs de $\delta$.
	 	\item Le ppcm $a_1\vee\ldots\vee a_n$ est l'unique $\mu\in\N$ tel que tel que $a_1\Z\cap\ldots\cap a_n\Z=\mu\Z$. Il est multiple de chacun des $a_i$, et l'ensemble des multiples communs des $a_i$ est exactement égal à l'ensemble des multiples de $\mu$.
	 \end{itemize}
\end{dfn}
\begin{rmq}
	Comme précédemment, si on note $\delta$ le pgcd des $a_i$, alors $$\exists(u_1,\ldots,u_n)\in\Z^n, \ \sum_{i=1}^{n}u_ia_i=\delta$$
\end{rmq}
En tant qu'éléments inversibles, $1$ et $-1$ divisent trivialement tout entier. Réciproquement, on a la 
\begin{dfn}[Entiers premiers entre eux]
	Deux entiers sont \textbf{premiers entre eux} lorsque leurs diviseurs communs sont exactement $1$ et $-1$. Plus généralement, soit $(a_1,\ldots,a_n)\in\Z^n$. Les $a_i$ sont dits \textbf{premiers entre eux dans leur ensemble} lorsque leurs diviseurs communs sont exactement $1$ et $-1$.
\end{dfn}
\begin{rmq}
	Soit $(a,b)\in\Z^2$.
	\begin{itemize}
		\item Pour montrer que $a$ et $b$ sont premiers entre eux, il suffit donc de se donner $d\in\mathcal{D}(a)\cap\mathcal{D}(b)$ puis de montrer que $d\in\left\{-1,1\right\}$. En effet, l'inclusion réciproque est toujours vraie.
		\item On constate facilement que $a$ et $b$ sont premiers entre eux si, et seulement si, $\mathcal{D}^+(a)\cap\mathcal{D}^+(b)=\left\{1\right\}$. Une variante consiste donc à se donner $d\in\mathcal{D}^+(a)\cap\mathcal{D}^+(b)$, et à montrer que $d=1$. En effet, là aussi, l'inclusion réciproque est toujours vraie.
	\end{itemize}
\end{rmq}
\begin{exa}
	Deux entiers consécutifs sont toujours premiers entre eux.
\end{exa}
\begin{rmq}
	Soit $(a,b)\in\Z^2$. il est facile de voir qu'on a $a\wedge b=1$ si, et seulement si, $a$ et $b$ sont premiers entre eux.
	\begin{itemize}
		\item Si $a\wedge b=1$, alors tout diviseur de $a$ et $b$ doit diviser $1$, donc être une unité de $\Z$.
		\item Réciproquement, si $a$ et $b$ sont premiers entre eux, alors $a\wedge b$, en tant que diviseur commun, est égal à $\pm1$. Et comme il est dans $\N$, il est égal à $1$.
	\end{itemize}
	Plus généralement, des entiers sont premiers entre eux dans leur ensemble si, et seulement si, leur pgcd vaut $1$.
\end{rmq}
\begin{rmq}
	Si les $a_i$ sont deux à deux premiers entre eux, alors ils sont en particulier premiers entre eux dans leur ensemble (au moins si $n\geqslant2$), mais la réciproque est fausse ! Considérer par exemple le tripler $(2,3,4)$.
\end{rmq}
\begin{prop}
	Soit $(a,b)\in\Z^2$ et $k\in\Z$. Alors $$(ka)\wedge(kb)=k(a\wedge b)$$
\end{prop}
\begin{exe}
	On peut prouver le symétrique pour le ppcm. Il est recommandé de commencer traiter à part le cas $k=0$.
\end{exe}
\begin{prop}
	Soit $(a,b)\in\Z^2$ et $\delta=a\wedge b$. Alors $$\exists(a',b')\in\Z^2,\ \left \{ \begin{array}{lll}a=\delta a' \\ b=\delta b' \\ a'\wedge b'=1\end{array}\right.$$
\end{prop}
\begin{rmq}
	Ces deux dernières propositions restent vraies avec plus de deux entiers relatifs.
\end{rmq}
\begin{thm}[Théorème de Bézout]
	Soit $(a,b)\in\Z$. Alors $a$ et $b$ sont premiers entre eux si, et seulement si, $$\exists(u,v)\in\Z^2,\ au+bv=1$$
\end{thm}
\begin{cor}[Corollaire du théorème de Bézout]
	Si $a$ est premier avec chacun des $(b_i)_{1\leqslant i\leqslant n}$, alors il est premier avec leur produit.
\end{cor}
\begin{exa}
	Soit $(a_1,\ldots,a_m)\in\Z^m$ et $(b_1,\ldots,b_n)\in\Z^n$ tels que $$\forall(i,j)\in\llbracket1,m\rrbracket\times\llbracket1,n\rrbracket,\ a_i\wedge b_j=1$$ Alors, une double application de ce corollaire montre que $$(a_1\times\ldots\times a_m)\wedge(b_1\times\ldots\times b_n)=1$$ En particulier, si $a\wedge b=1$, alors $\forall(p,q)\in(\N\st)^2,\ a^p\wedge b^q=1$.
\end{exa}
\begin{prop}[Généralisation du théorème de Bézout]
	Soit $(a_1,\ldots,a_n)$ une famille finie d'entiers. Alors les $a_i$ sont premiers entre eux dans leur ensemble si, et seulement si, $$\exists(u_1,\ldots,u_n)\in\Z^n,\ \sum_{i=1}^{n}a_iu_i=1$$
\end{prop}
Étant donnés $a$ et $b$, il est bon de connaître une méthode pratique pour déterminer un couple de Bézout $(u,v)$ tel que $au+bv=a\wedge b$ (notons qu'il n'y a pas unicité). Cette méthode est fournie par \textbf{l'algorithme d'Euclide}, qui permet à la fois d'obtenir le pgcd de $a$ et $b$ et un couple $(u,v)$ correspondant. Avant tout, démontrons le
\begin{lem}[Invariance du pgcd par transvection]
	Soit $(a,b,\lambda)\in\Z^3$. Alors on a : $$a\wedge b=a\wedge (b+\lambda a)$$
\end{lem}
\begin{rmq}
	Si $a,b\in\N\st$, on en déduit une méthode pour calculer $a\wedge b$. En notant $r$ le reste de la division euclidienne de $a$ par $b$, le lemme montre que $a\wedge b=b\wedge r$. Ainsi, on se ramène à des entiers plus petits.
\end{rmq}
\begin{met}[Algorithme d'Euclide pour calculer le pgcd de $a>0$ et $b>0$]
	Soit $a>0$ et $b>0$. \begin{enumerate}
		\item On pose $a_0=a$ et $b_0=b$.
		\item $r_0$ est le reste de la division euclidienne de $a$ par $b$.
		\item Tant que $r_n\ne0$ : \begin{itemize}
			\item On pose $a_{n+1}=b_n$ et $b_{n+1}=r_n$ ;
			\item $r_{n+1}$ est le reste de la division euclidienne de $a_{n+1}$ par $b_{n+1}$.
		\end{itemize}
		\item En sortie de boucle, en notant $n_f$ le rang pour lequel $r_{n_f}=0$, le pgcd cherché est $b_{n_f}$ (ou de manière équivalente $r_{n_f-1}$).
		\item Version "étendue" de l'algorithme : si à chaque étape on exprime $r_n$ comme combinaison de $a$ et $b$, à la fin on obtient une combinaison égale à $a\wedge b$.
	\end{enumerate}
\end{met}
\begin{thm}[Lemme de Gauss]
	Soit $(a,b,c)\in\Z^3$. Si $a$ divise $bc$ et si $a$ et $b$ sont premiers entre eux, alors $a$ divise $c$.
\end{thm}
\begin{cor}
	Si les $(a_i)_{1\leqslant i\leqslant n}$ divisent tous $b$ et sont deux à deux premiers entre eux, alors leur produit divise $b$.
\end{cor}
\begin{exe}
	Montrer que si les $a_i$ sont dans $\N$ en étant deux à deux premiers entre eux, alors $$a_1\vee\ldots\vee a_n=a_1\times\ldots\times a_n$$
\end{exe}
\begin{dfn}[Forme irréductible d'un rationnel]
	Tout rationnel non nul s'écrit d'une unique façon sous la forme $\displaystyle\frac{p}{q}$ avec $(p,q)\in\Z\times\N\st$ et $p\wedge q=1$. Dans ce cas, on dit que le rationnel est écrit sous \textbf{forme irréductible}.
\end{dfn}
\section{Nombres premiers}
\begin{dfn}[Nombres premiers]
	Un entier naturel $p\geqslant2$ est un \textbf{nombre premier} lorsque ses seuls diviseurs positifs sont $1$ et lui-même. Dans la suite de ce chapitre, l'ensemble des nombres premiers sera noté $\mathcal{P}$, mais il est bon de rappeler cette notation sur une copie.
\end{dfn}
\begin{rmq}
	Par convention, $1$ n'est donc pas premier, on verra plus tard pourquoi ce choix. Ensuite, à part $2$, tous les nombres premiers sont impairs puisque tous les nombres pairs sont divisibles par $2$.
\end{rmq}
\begin{exa}
	Les premiers nombres premiers sont $2,3,5,7,11,13,\ldots$
\end{exa}
\begin{prop}
	Soit $n\geqslant2$. les assertions suivantes sont équivalentes :
	\begin{enumerate}
		\item $n$ n'est pas premier (on dit parfois que $n$ est \textbf{composé})
		\item $\exists d_1,d_2\in\N\st,\ n=d_1d_2 \ \ \text{et} \ \ d_1,d_2>1$ 
		\item $\exists d_1,d_2\in\N\st,\ n=d_1d_2 \ \ \text{et} \ \ d_1,d_2<n$
		\item $\exists d_1,d_2\in\N\st,\ n=d_1d_2 \ \ \text{et} \ \ 1<d_1,d_2<n$
	\end{enumerate}
\end{prop}
\begin{rmq}
	Soit $p$ un nombre premier et $n\in\Z$. Si $p$ ne divise pas $n$, alors $p$ est premiers avec $n$. En particulier, deux nombres premiers distincts sont toujours premiers entre eux.
\end{rmq}
\begin{prop}[Lemme d'Euclide]
	Soit $(a_1,\ldots,a_n)\in\Z^n$ et $p$ premier. Si $p$ divise le produit des $a_i$, alors $p$ divise l'un des $a_i$.
\end{prop}
\begin{exa}
	Soit $a\in\Z$, $n\in\N\st$ et $p$ premier. Alors $p$ divise $a^n$ si, et seulement si, $p$ divise $a$.
\end{exa}
\begin{lem}
	Tout entier $n\geqslant2$ admet un diviseur premier.
\end{lem}
\begin{cor}
	Si $n$ n'est divisible par aucun nombre premier de l'intervalle $\left\llbracket2,\lfloor\sqrt{n}\rfloor\right\rrbracket$, alors $n$ est premier.
\end{cor}
\begin{rmq}
	Soit $(a,b)\in\Z^2$. Si $a$ et $b$ ne sont pas premiers entre eux, alors ils admettent un diviseur premier commun. Par l'absurde, cela permet régulièrement de montrer que deux entiers sont premiers entre eux.
\end{rmq}
\begin{exe}
	Soit $a$ et $b$ deux entiers premiers entre eux. Montrer que $a^2+b^2$ et $ab$ sont premiers entre eux.
\end{exe}
\begin{thm}
	Il existe une infinité de nombres premiers.
\end{thm}
\begin{dfn}[Valuation $p$-adique]
	Soit $n\in\Z\st$ et $p$ premier. La \textbf{valuation $p$-adique} de $n$ est définie par : $$v_p(n)=\max\left\{k\in\N : p^k\mid n\right\}$$
\end{dfn}
\begin{rmq}
	Si $k\leqslant v_p(n)$, alors $p^k$ divise toujours $n$, et si $k>v_p(n)$, $p^k$ ne divise jamais $n$. En clair, $p^k$ divise $n$ si, et seulement si, $k\leqslant v_p(n)$. Autrement dit, on a : $$\left\{k\in\N : p^k\mid n\right\}=\llbracket0,v_p(n)\rrbracket$$ En particulier, $p$ divise $n$ si, et seulement si, $v_p(n)\geqslant1$, et donc $p$ ne divise pas $n$ si, et seulement si, $v_p(n)=0$.
\end{rmq}
\begin{exa}
	On a toujours $v_p(-n)=v_p(n)$.
\end{exa}
\begin{exa}
	Si $p$ est premier et $\alpha\in\N$, alors $v_p(p^\alpha)=\alpha$ (utiliser la stricte croissance de $k\mapsto p^k$).
\end{exa}
\begin{lem}
	Soit $n\in\Z\st$, $p$ premier et $k\in\N$. Alors $$v_p(n)=k \iff \exists n'\in\Z,\ \left \{ \begin{array}{ll}n=p^kn'\\ n'\wedge p=1\end{array}\right.$$
\end{lem}
\begin{prop}
	Soit $m,n\in\Z\st$. pour tout nombre premier $p$, on a : $$v_p(mn)=v_p(m)+v_p(n)$$
\end{prop}
\begin{exa}
	Soit $k\in\N\st$. Par une récurrence immédiate, on montre que $$\forall(n_1,\ldots,n_k)\in\left(\Z\st\right)^k,\ v_p\left(\prod_{i=1}^{k}n_i\right)=\sum_{i=1}^{k}v_p(n_i)$$ En particulier, $$\forall n\in\Z\st,\ v_p(n^k)=kv_p(n)$$
\end{exa}
\begin{dfn}[Décomposition en facteurs premiers]
	Soit $n\geqslant1$. Une \textbf{décomposition en facteurs premiers} (ou \textbf{DFP}) est une écriture de $n$ sous la forme $$n=\prod_{p\in I}p^{\alpha_p}$$ où $I$ est une partie finie de $\mathcal{P}$ et les $\alpha_p$ sont dans $\N\st$. Pour $n=1$, on peut poser $I=\emptyset$, et on obtient ce qu'on appelle la "DFP vide".
\end{dfn}
\begin{rmq}
	On évitera l'acronyme "DFP" sur une copie, car il n'est pas standard.
\end{rmq}
\begin{thm}
	Tout entier $n\geqslant1$ admet une décomposition en facteurs premiers.
\end{thm}
\begin{prop}
	Soit $n\geqslant1$, et une DFP de $n$ écrite sous la forme $\displaystyle n=\prod_{p\in I}p^{\alpha_p}$. Alors, pour tout nombre premier $p_0$ : 
	\begin{itemize}
		\item Si $p_0\in I$, on a $v_{p_0}(n)=\alpha_{p_0}$
		\item Sinon, on a $v_{p_0}(n)=0$
	\end{itemize}
\end{prop}
\begin{cor}
	Soit $n\geqslant1$. Alors sa DFP est unique.
\end{cor}
\begin{rmq}
	Pour cette raison (existence et unicité de la DFP), on dit que $\Z$ est un anneau \textbf{factoriel}.
\end{rmq}
\begin{rmq}
	Au passage, on a montré que $p$ divise $n$ si, et seulement si, $p\in I$.
\end{rmq}
\begin{prop}
	Soit $m,n\geqslant1$. Alors $m$ et $n$ sont premiers entre eux si, et seulement si, ils n'ont aucun facteur commun dans leurs DFP respectives.
\end{prop}
\begin{prop}
	Soit $m,n\in\Z\st$. Alors $m\mid n$ si, et seulement si, pour tout nombre premier $p$, on a $v_p(m)\leqslant v_p(n)$.
\end{prop}
\begin{exa}
	Si $m,n\in\N\st$, alors $m=n$ si, et seulement si, pour tout nombre premier $p$, on a $v_p(m)=v_p(n)$.
\end{exa}
\begin{thm}
	Les diviseurs strictement positifs de $\displaystyle n=\prod_{p\in I}p^{\alpha_p}$ sont exactement les entiers $\displaystyle m=\prod_{p\in I}p^{\beta_p}$ avec $\forall p \in I,\ \beta_p\leqslant\alpha_p$.
\end{thm}
\begin{cor}
	$\displaystyle n=\prod_{p\in I}p^{\alpha_p}$ possède $\displaystyle\prod_{p\in I}(\alpha_p+1)$ diviseurs strictement positifs.
\end{cor}
\begin{rmq}
	En numérotant les éléments de $I$, on peut écrire la DFP de $n$ sous la forme bien connue $$n=p_1^{\alpha_1}\ldots p_k^{\alpha_k}$$ On obtient alors que "tout entier $n\geqslant1$ admet une unique DFP, à l'ordre des facteurs près". En pratique, c'est cette écriture que nous garderons dans la suite, car elle est plus commode pour les calculs. On s'entraînera avec profit à reformuler les résultats précédents dans ce formalisme.
\end{rmq}
\begin{prop}
	Soit $a$ et $b$ deux entier non nuls. Pour tout nombre premier $p$, on a : $$\left \{\begin{array}{ll}v_p(a\wedge b)=\min(v_p(a),v_p(b)) \\ v_p(a\vee b)=\max(v_p(a),v_p(b))\end{array}\right.$$
\end{prop}
\begin{rmq}
	Ce résultat se généralise à plus de deux entiers.
\end{rmq}
\begin{cor}
	On a $$\forall(a,b)\in\Z^2,\ (a\wedge b)(a\vee b)=\lvert ab\rvert$$
\end{cor}
\begin{exa}
	Soit $a,b>0$, $\delta=a\wedge b$ et $\mu=a\vee b$. En écrivant $a=\delta a'$ et $b=\delta b'$ avec $a'\wedge b'=1$, on en tire $\mu=\delta ab$.
\end{exa}
\begin{exa}
	Si $a,b\in\N$ sont premiers entre eux, alors on retrouve que $a\vee b=ab$.
\end{exa}
\section{Congruences}
Dans toute cette partie, $n$ est un entier naturel strictement positif.
\begin{dfn}[Congruence modulo $n$]
	On rappelle que $a$ et $b$ sont \textbf{congrus modulo $n$}, et on écrit $a\equiv b \ [n]$ (voire $a=b \ [n]$) lorsque $n\mid b-a$ (soit encore lorsque $b-a\in n\Z$).
\end{dfn}
\begin{exa}
	Soit $a,b,k\in\Z$ avec $k\wedge n=1$. Si  $ka\equiv kb\ [n]$, alors $a\equiv b\ [n]$.
\end{exa}
\begin{thm}
	La congruence modulo $n$ est une relation d'équivalence sur $\Z$.
\end{thm}
\begin{dfn}[$\Z/n\Z$]
	L'ensemble des classes d'équivalence modulo $n$ (autrement dit l'ensemble-quotient de $\Z$ par la relation de congruence modulo $n$) est noté $\Z/n\Z$. C'est un ensemble fini de cardinal $n$, et ses éléments sont donnés par : $$\Z/n\Z=\left\{\overline{0},\overline{1},\ldots,\overline{n-1}\right\}$$
\end{dfn}
\begin{rmq}
	Ainsi, deux entiers sont congrus modulo $n$ si, et seulement si, ils sont le même reste dans la division euclidienne par $n$.
\end{rmq}
\begin{prop}
	L'addition et la multiplication de $\Z$ passent au quotient : si $a\equiv a'\ [n]$ et $b\equiv b'\ [n]$, alors $a+b\equiv a'+b'\ [n]$ et $ab\equiv a'b'\ [n]$.
\end{prop}
\begin{rmq}
	La compatibilité avec l'addition est vrai modulo n'importe quoi. Toutefois, elle tombe souvent en défaut pour la multiplication et ne fonctionne ici que parce que le modulo, \textit{ie} $n$, est entier.
\end{rmq}
\begin{exa}
	Par une récurrence immédiate, on montre que pour tout $k\in\N$, si $a\equiv b\ [n]$, alors $a^k\equiv b^k\ [n]$. Attention à ne pas écrire cela pour $k<0$, cela n'a aucun sens !
\end{exa}
Cela justifie la
\begin{dfn}
	On munit $\Z/n\Z$ d'une addition et d'une multiplication comme suit. $$\left \{\begin{array}{ll}\overline{a}+\overline{b}=\overline{a+b}\\ \overline{a}\times\overline{b}=\overline{ab}\end{array}\right.$$ Ces opérations son bien définies : elles ne dépendent pas des représentants choisis.
\end{dfn}
\begin{thm}
	$(\Z/n\Z,+,\times)$ est un anneau commutatif, d'élément neutre $\overline{1}$ pour $\times$.
\end{thm}
\begin{rmq}
	Autrement dit, la surjection canonique $\varphi:x\mapsto\overline{x}$ devient par définition un morphisme d'anneaux. En particulier,
	\begin{itemize}
		\item par itération de la loi $+$, on a $\forall x\in\Z,\ \forall k\in\Z,\ k\cdot\overline{x}=\overline{kx}$ ;
		\item par itération de la loi $\times$, on a $\forall x\in\Z,\ \forall k\in\N,\ \overline{x}^k=\overline{x^k}$.
	\end{itemize}
\end{rmq}
\begin{exa}
	L'application de $(\Z/n\Z,+)$ dans $(U_n,\times)$ qui à $\overline{k}$ associe $\exp\left(i\frac{2k\pi}{n}\right)$ est bien définie et est un isomorphisme de groupes.
\end{exa}
\begin{thm}
	$\overline{a}$ est inversible dans $\Z/n\Z$ si, et seulement si, $a\wedge n=1$.
\end{thm}
\begin{proof}
	Si $\overline{a}$ est inversible, on peut écrire $\overline{a}\overline{b}=\overline{1}$ puis $n\mid ab-1$, donc d'après le théorème de Bézout, $a$ et $n$ sont premiers entre eux. Réciproquement, utilisons à nouveau le théorème de Bézout. Si $a\wedge n=1$, on écrit $au+nv=1$, puis en passant dans $\Z/n\Z$, on obtient $\overline{a}\overline{u}=\overline{1}$.
\end{proof}
\begin{thm}
	Soit $n\in\N\st$. Les trois assertions suivantes sont équivalentes :
	\begin{enumerate}
		\item $n$ est premier
		\item $\Z/n\Z$ est un corps
		\item $\Z/n\Z$ est un anneau intègre
	\end{enumerate}
\end{thm}
\begin{proof}
	Si $n$ est premier, donnons-nous $\overline{a}\ne\overline{0}$. Alors $a\wedge n=1$, donc $\overline{a}$ est inversible. De plus, tout corps est un anneau intègre. Enfin, si $\Z/n\Z$ est intègre, supposons que $n$ n'est pas premier; On ne peut pas avoir $n=1$, sans quoi $\Z/n\Z$ serait l'anneau nul, donc $n\geqslant2$. Écrivons donc $n=ab$ avec $a,b<n$. On obtient $\overline{a}\times\overline{b}=\overline{0}$, ce qui est absurde par intégrité, donc $n$ est premier.
\end{proof}
Pour finir, citons le célèbre "petit théorème de Fermat". Avant toute chose, démontrons le
\begin{lem}
	Soit $p$ un nombre premier et $a,b\in\Z$. Alors $$(a+b)^p\equiv a^p+b^p\ [p]$$
\end{lem}
\begin{proof}
	Soit $k\in\llbracket1,p-1\rrbracket$. $p$ ne divise pas $k$, donc $p\wedge k=1$. D'après la formule du capitaine, $\displaystyle k\binom{p}{k}=p\binom{p-1}{k-1}$ puis par le lemme de Gauss, $p$ divise $\displaystyle\binom{p}{k}$. Il suffit alors d'utiliser le binôme de Newton, de passer dans $\Z/p\Z$ et d'éliminer alors les termes non extrémaux.
\end{proof}
\begin{thm}[Petit théorème de Fermat]
	Soit $p$ un nombre premier et $a\in\Z$. Alors : $$a^p\equiv a\ [p]$$
\end{thm}
\begin{proof}
	Sans perte de généralité, on se ramène au reste $r$ de $a$ modulo $p$. puis il suffit d'effectuer une récurrence finie sur $r$ en appliquant le lemme qui précède pour l'hérédité.
\end{proof}
\section{Compléments d'arithmétique, HP}
Dans ce paragraphe, on démontre bon nombre de résultats arithmétiques hors programme majeurs, qui permettent souvent de faciliter la résolution d'un grand nombre d'exercice. Il est donc important de connaître ces résultats et les grandes lignes de leurs démonstrations. \\ \par
Pour commencer, voici une application d'un théorème que nous avons vu dans le cas des groupes.
\begin{thm}[Petit théorème de Fermat - autre démonstration]
	Soit $p$ un nombre premier, et $a$ non divisible par $p$. Alors $a^{p-1}\equiv 1 \ [p]$
\end{thm}
\begin{proof}
	Puisque, $\overline{a}$ appartient au groupe $\left(\Z/p\Z\right)\st$ qui est de cardinal $p-1$. Il suffit alors d'appliquer le "théorème de Fermat" vu dans le chapitre de théorie des groupes : un élément du groupe élevé à la puissance du cardinal du groupe vaut l'élément neutre.
\end{proof}
\begin{dfn}[Indicatrice d'Euler]
	L'\textbf{indicatrice d'Euler} est la fonction définie par $$\varphi(n)=\card\left(\left\{k\in\llbracket1,n\rrbracket \ | \ k\wedge n=1\right\}\right)$$ 
\end{dfn}
Plus généralement, on peut alors démontrer le : 
\begin{thm}[Théorème de Fermat-Euler]
	Si $a\wedge n=1$, alors $a^{\varphi(n)} \equiv 1 \ [n]$.
\end{thm}
\begin{proof}
	Si $a$ est premier avec $n$, alors $\overline{a}$ appartient au groupes des unités $\mathcal{U}\left(\Z/n\Z\right)$, qui est de cardinal $\varphi(n)$. Il suffit alors d'appliquer le "théorème de Fermat" vu sur les groupes.
\end{proof}
\begin{thm}[Théorème des restes chinois]
	Si $a\wedge b=1$, alors l'application $$
	\begin{array}[t]{lrcl}
		\psi: & \Z/ab\Z& \longrightarrow & \Z/a\Z\times\Z/b\Z \\
		& \overline{x}[ab] & \longmapsto & (\overline{x}[a],\overline{x}[b])
	\end{array}
	$$ est un isomorphisme d'anneaux.
\end{thm}
\begin{proof}
	Pour commencer, $\psi$ est bien définie, car si $x\equiv x' \ [ab]$, alors en particulier $x\equiv x'\ [a]$ et $x\equiv x'\ [b]$. Ensuite, c'est évidemment un morphisme d'anneaux d'après les définitions des opérations sur les ensembles $\Z/m\Z$. Il reste à montrer que c'est un isomorphisme. \par Pour montrer que $\psi$ est injectif, supposons que $x\equiv 0\ [a]$ et  $x\equiv 0\ [b]$, alors $x\equiv 0\ [ab]$ d'après le corollaire du lemme de Gauss. \par Pour montrer que $\psi$ est surjective, il suffit alors de comparer les cardinaux. Si on souhaite une preuve plus constructive, donnons $\alpha\in\llbracket1,a\rrbracket$, $\beta\in\llbracket1,b\rrbracket$, et cherchons $x$ tel que $x\equiv\alpha\ [a]$ et $x\equiv\beta\ [b]$. Il suffit alors de considérer un couple de Bézout $(u,v)$ tel que $ua+bv=1$ puis de poser $x=ua\beta+vb\alpha$.
\end{proof}
\begin{rmq}
	Ce théorème s'appelle ainsi car il trouve son origine dans une récréation mathématique publiée dans un ancien manuscrit chinois.
\end{rmq}
\begin{cor}
	Si $a\wedge b=1$, alors $\varphi(ab)=\varphi(a)\varphi(b)$.
\end{cor}
\begin{proof}
	Un élément d'un anneau-produit est inversible si, et seulement si, ses deux composantes le sont. Or, $\varphi(ab)$ représente le nombre d'inversibles de $\Z/ab\Z$, et $\varphi(a)\varphi(b)$ celui des éléments inversibles de $\Z/a\Z\times\Z/b\Z$. On conclut par le théorème des restes chinois.
\end{proof}
Le théorème qui suit constitue un résultat important, utile dans de nombreux cas (et notamment dans le calcul du \textbf{déterminant de Smith}).
\begin{thm}
	On a l'identité suivante : $$\forall n\in\N\st,\ \sum_{d\mid n}\varphi(d)=n$$
\end{thm}
\begin{proof}
	Soit $n \in\N^*$. On définit $\displaystyle E_n= \left\{ \frac{k}{n} \mid k \in \llbracket 1,n\rrbracket\right\}$. On définit alors pour tout diviseur positif $d$ de $n$ l'ensemble $\displaystyle D_d= \left\{ \frac{k}{d} \mid k \in \llbracket 1,d\rrbracket, k \wedge d =1 \right\}$. Il suffit alors de montrer l'égalité suivante, ce qui suffira à conclure par passage aux cardinaux : $\displaystyle E_n = \bigsqcup_{d \mid n} D_d$.	En effet, chaque $D_d$ est immédiatement en bijection avec l'ensemble par lequel est défini $\varphi(d)$. Passons à la preuve de cette égalité.
	\begin{itemize}
		\item Montrons qu'il s'agit là d'une union disjointe. Soit $d_1$ et $d_2$
		des diviseurs distincts de $n$. Raisonnons par l'absurde : supposons qu'il
		existe un $x \in D_{d_1} \cap D_{d_2}$ et fixons-le. Conformément aux
		définitions de $D_{d_1}$ et $D_{d_2}$, on peut fixer $\left( k_1, k_2
		\right) \in \llbracket 1,d_1\rrbracket \times \llbracket 1,d_2\rrbracket$
		tel que	$$x=\frac{k_1}{d_1}=\frac{k_2}{d_2} \ \text{avec} \  k_1 \wedge d_1=1 \ \text{et} \ k_2 \wedge d_2=1$$
		On a alors $k_1 d_2 = k_2 d_1$ puis, d'après le lemme de Gauss, on a $k_1
		\mid k_2$ et $k_2 \mid k_1$, puis il s'ensuit que $k_1 \le k_2$ et $k_2
		\le k_1$ car $k_1$ et $k_2$ sont des entiers naturels non nuls. Par
		antisymétrie, on obtient que $k_1=k_2$, ce qui est \textbf{absurde}. L'union
		considérée est donc disjointe.
		
		\item Montrons l'inclusion directe. Soit $x \in E_n$. Fixons $k \in
		\llbracket 1, n \rrbracket$ tel que $x= \displaystyle\frac{k}{n}$. On peut
		fixer alors $k'$ et $d$ tels que $k= k' \times (k \wedge n)$ et $n= d
		\times (k \wedge n)$. Puisque $k \le n$ alors $k' \le d$ puis $x =
		\displaystyle\frac{k'}{d}$ avec $d \mid n$ donc $x \in D_d$. L'inclusion
		directe est donc bien vérifiée.
		
		\item Montrons l'inclusion réciproque. Soit $x \in \displaystyle\bigsqcup_{d
			\mid n} D_d$. Fixons $d \mid n$ tel que $x \in D_d$. Fixons $k \in
		\llbracket 1,d \rrbracket$ tel que $x= \displaystyle\frac{k}{d}$. Or, $d
		\mid n$ donc on peut fixer $l$ tel que $dl=n$. On sait que $1 \le k \le d$
		donc $1 \le l \le kl \le n$. Puis $$x= \frac{kl}{n} \in E_n$$
	\end{itemize}
	Ainsi, la preuve est achevée.
\end{proof}
\begin{thm}
	On a $$\forall n \in\N^*, \ \varphi(n) = n\prod_{p\mid n} \left(1-\frac{1}{p}\right)$$
\end{thm}

\begin{proof}
	On expose ici une démonstration probabiliste de ce résultat, moins lourde que la démonstration arithmétique pure. Soit $n \in\N^*$. On considère $\Omega = \llbracket 1,n \rrbracket$ muni de la probabilité uniforme qu'on note $\mathbb{P}$. Si $d \in \llbracket 1,n \rrbracket$ est un diviseur positif de $n$, on note$$
		A_d = \left\{ x \in \llbracket 1,n \rrbracket \mid \exists k \in \llbracket 1,n \rrbracket, \ x=kd \right\}$$
	Soit $d \in \llbracket 1,n \rrbracket$ un diviseur de $n$. On montre aisément que $\card \left(A_d\right)=\left\lfloor \displaystyle\frac{n}{d} \right\rfloor = \displaystyle\frac{n}{d}$. Il s'ensuit que $$
		\mathbb{P} \left(A_d\right) = \frac{1}{d}$$
	Ensuite, on note $p_1 < ... < p_r$ les diviseurs premiers de $n$ rangés dans l'ordre croissant.\par
	On montre tout d'abord que les événements $\left(A_{p_i}\right)_{1 \le i \le r}$ sont mutuellement indépendants. En effet, soit $I \subset \llbracket 1,r \rrbracket$. Soit $\omega \in \Omega$. On a alors, en vertu du corollaire du lemme de Gauss$$
		\omega \in \bigcap_{i \in I} A_{p_i}
		\iff \forall i \in I, \ p_i \mid \omega
		\iff \prod_{i \in I} p_i \big| \omega
		\iff \omega \in A_{\prod_{i \in I} p_i}$$
	Il s'ensuit alors, d'après notre premier résultat et puisque $\displaystyle\prod_{i \in I} p_i \big| n$ :$$
		\mathbb{P} \left(\bigcap_{i \in I} A_{p_i}\right) = \mathbb{P} \left(A_{\prod_{i \in I} p_i}\right)= \frac{1}{\displaystyle\prod_{i \in I} p_i} = \displaystyle\prod_{i \in I} \frac{1}{p_i} = \displaystyle\prod_{i \in I} \mathbb{P} \left(A_i\right)$$
	L'indépendance est démontrée. \par
	Enfin, on considère $A =\left\{ m \in \llbracket 1,n \rrbracket \mid m \wedge
	n =1 \right\} $ et on veut montrer que $\varphi(n) = n\displaystyle\prod_{p\mid n} \left(1-\frac{1}{p}\right)$.	Montrons que $\displaystyle\frac{\varphi(n)}{n} = \displaystyle\prod_{p\mid n} \left(1-\frac{1}{p}\right)$, ce qui suffira à conclure. Or, on a d'une part, puisque la probabilité est uniforme :$$
		\mathbb{P}(A)=\frac{\varphi(n)}{n}$$
	Et d'autre part, puisque $A = \displaystyle\bigcap_{i=1}^{r} \overline{A_{p_i}} $ et puisque les événements $\left(\overline{A_{p_i}}\right)_{1 \le i \le r}$ sont mutuellement indépendants, on a :$$
		\mathbb{P}(A) =\mathbb{P} \left(\bigcap_{i=1}^r A_{p_i}\right)=\prod_{i=1}^r \mathbb{P} \left( \overline{A_{p_i}}\right) =\prod_{i=1}^r \left( 1-\mathbb{P}\left({A_{p_i}}\right)\right) = \prod_{i=1}^r \left( 1-\frac{1}{p_i}\right)$$
	D'où : $$\frac{\varphi(n)}{n} =\prod_{p \mid n} \left( 1-\frac{1}{p}\right)$$
	Ce qui achève la preuve.
\end{proof}
Voici pour finir un théorème assez célèbre, un des premiers qu'on rencontre en arithmétique élémentaire.
\begin{thm}[Théorème de Wilson]
	$p$ est premier si, et seulement si, $(p-1)!\equiv -1\ [p]$.
\end{thm}
\begin{proof}
	Si $(p-1)!\equiv -1\ [p]$, alors toute classe non nulle modulo $p$ est inversible. En effet, il suffit d'écrire $$\overline{x}\times\left(-\prod_{\substack{1\leqslant i \leqslant p-1 \\ \overline{i}\ne\overline{x}}}\overline{i}\right)=\overline{1}$$ Ainsi, $\Z/p\Z$ est un corps et $p$ est premier. \par Réciproquement, supposons que $p$ est premier. Le cas $p=2$ étant immédiat, passons à $p\geqslant3$. Toute classe $\overline{x}\ne\overline{0}$ est inversible puisque $\forall x\in\llbracket1,p-1\rrbracket,\ x\wedge p=1$. Cherchons les classes qui sont leur propre inverse. On a : $$\overline{x}^2=\overline{1}\iff(\overline{x}+\overline{1})(\overline{x}-\overline{1})=\overline{0}\iff \overline{x}=\pm\overline{1}$$ par intégrité dans le corps $\Z/p\Z$. A part ces deux classes, toute autre classe non nulle a un inverse qui n'est pas elle-même. Dans le produit $\overline{1}\times\ldots\times\overline{p-1}$, mettons à part $\overline{1}$ et $\overline{p-1}=\overline{-1}$, et regroupons les autres par paires d'inverses mutuels. On obtient alors le produit de $\overline{1}$ et de $\overline{-1}$, d'où le résultat.
\end{proof}
\section{Petite généralisation :  idéaux et anneaux principaux, HP}
Dans toute cette partie, l'anneau $A$ est supposé commutatif.
\begin{dfn}[Idéal]
	Un \textbf{idéal} $I$ de $A$ est un sous-groupe additif de $A$ absorbant pour la multiplication : $$\forall(a,i)\in A\times I,\ ai\in I$$
\end{dfn}
\begin{exa}[Noyau d'un morphisme d'anneaux commutatifs]
	Le noyau d'un morphisme d'anneaux commutatifs est toujours un idéal. Il s'agit évidemment d'un sous-groupe de $A$, et il est absorbant par absorbance de $0$ dans l'anneau et par propriété de morphisme d'anneau pour le produit.
\end{exa}
\begin{prop}[Idéal engendré par un élément]
	Soit $x\in A$. Alors $xA$ est un idéal de $A$, appelé \textbf{idéal engendré par $x$}. On le note souvent $(x)$.
\end{prop}
\begin{proof}
	Montrons déjà que $xA$ est un sous-groupe additif. Déjà, $0=x\times 0\in xA$. Ensuite, soit $xa,xa'\in xA$. Alors $xa-xa'=x(a-a')\in xA$. Pour l'absorbance, la commutativité est essentielle. Soit $xa\in A$ et $b\in A$. Alors $b(xa)=x(ba)\in xA$.
\end{proof}
On peut se demander si la réciproque est vraie. C'est ce qui fait l'objet de la
\begin{dfn}[Anneau principal]
	Si réciproquement
	\begin{itemize}
		\item les seuls idéaux de $A$ sont de la forme $xA$
		\item $A$ est intègre
	\end{itemize}
	alors $A$ est dit \textbf{principal}.
\end{dfn}
\begin{exa}[$\Z$ est principal]
	Les idéaux de $\Z$ sont de la forme $n\Z$, avec $n\in\N$ (ce sont mêmes les sous-groupes additifs de $\Z$). En particulier, $\Z$ est principal.
\end{exa}
\begin{dfn}[Divisibilité]
	Soit $A$ un anneau intègre. On dit que $a$ divise $b$ et on note $a\mid b$ lorsque $b\in aA$.
\end{dfn}
\begin{dfn}[Éléments premiers entre eux]
	Soit $a$ et $b$ deux éléments de $A$. Les unités de $A$ divisent toujours $a$ et $b$. Si réciproquement, les seuls diviseurs communs de $a$ et $b$ sont les unités de $A$, on dit que $a$ et $b$ sont \textbf{premiers entre eux}.
\end{dfn}
\begin{exa}
	Dans $\Z$, les seuls unités étant $-1$ et $1$, on retrouve bien la définition usuelle selon laquelle deux entiers sont premiers entre eux lorsque leurs seuls diviseurs communs sont $-1$ et $1$.
\end{exa}
\begin{thm}
	La somme de deux idéaux est un idéal.
\end{thm}
\begin{proof}
	On sait déjà que dans un groupe commutatifs, la somme de deux sous-groupes est un sous-groupe. La propriété d'absorbance se vérifie immédiatement par distributivité dans l'anneau de départ.
\end{proof}
\begin{cor}
	Dans un anneau principal, le théorème de Bézout (y compris sa généralisation à $n$ éléments) reste vrai. $a$ et $b$ sont premiers entre eux si, et seulement si, $$\exists(u,v)\in A^2,\ ua+vb=1$$
\end{cor}
\begin{proof}
	La démonstration s'adapte sans difficulté. Si $ua+bv=1$, tout diviseur commun de $a$ et $b$ doit diviser $1$, donc doit être une unité. Réciproquement, soit $a$ et $b$ premiers entre eux. On sait que $aA+bA$ est un idéal de $A$, donc il est de la forme $\delta A$. En particulier, $a$ et $b$ sont dans $\delta A$, donc $\delta$ est une unité. On écrit ensuite $au+bv=\delta$ puis $$(\delta\inv)a+(\delta\inv)b=1$$ Ainsi, la preuve est achevée.
\end{proof}
\begin{cor}
	Dans un anneau principal, le lemme de Gauss reste vrai. De même, les corollaires du théorème de Bézout et du lemme de Gauss restent également vrais.
\end{cor}
\begin{proof}
	Les preuves sont intégralement transposables.
\end{proof}
\begin{rmq}
	En revanche, la définition d'un pgcd et d'un ppcm étant plus délicate, nous ne nous attarderons pas sur ce point ici. Nous reviendrons dessus plus en détail dans les compléments sur les anneaux commutatifs et en verrons un exemple dans le chapitre des polynômes.
\end{rmq}
\chapter{Compléments sur les anneaux commutatifs, HP}
Avant de commencer ce chapitre bien éloigné du programme, je tiens à remercier mon co taupin Nicolas qui a un jour déposé ce polycopié de compléments sur les anneaux commutatifs de M. Lafitte, duquel toute cette partie (ou presque) est tirée, sur mon bureau à ND3. Et en plus il ne m'avait même pas dit que c'était lui. Bisous mon loulou ! <3 \\ \par Dans tous ces compléments, les anneaux considérés seront commutatifs. Pour $A$ un anneau commutatif, on notera $0$ pour $0_A$ et $1$ pour $1_A$.
\section{Idéal d'un anneau commutatif}
\subsection{Diviseurs de zéro}
Certains éléments d'un anneau ont des propriétés particulières par rapport à la multiplication, ce qui justifie quelques définitions.
\begin{dfn}[Élément simplifiable, diviseur de zéro]
	Soit $A$ un anneau commutatif. On dit qu'un élément $a\in A$ non nul est \textbf{simplifiable} si $$\forall b\in A,\ ab=0\implies b=0$$ Dans le cas contraire, on dit que $a$ est un \textbf{diviseur de zéro}.
\end{dfn}
\begin{dfn}[Anneau intègre]
	On dit qu'un anneau $A$ est \textbf{intègre} si $A$ est non trivial, commutatif (ce qui est supposé dans tout le chapitre) et si tout élément non nul de $A$ est simplifiable.
\end{dfn}
\begin{dfn}[Ensemble des inversibles]
	On note $\mathcal{U}(A)$, ou $A^\times$ ou encore $A\st$ (mais cette dernière notation est parfois ambigüe), l'ensemble des éléments inversibles de l'anneau $A$. On l'appelle l'ensemble des unités de $A$.
\end{dfn}
\begin{prop}[Groupes des unités]
	L'ensemble des éléments inversibles d'un anneau $A$ est un groupe pour la multiplication. On l'appelle le \textbf{groupe des unités de $A$}.
\end{prop}
\begin{proof}
	Déjà fait dans le chapitre "Groupes, anneaux, corps".
\end{proof}
\begin{dfn}[Relation d'association]
	Pour $a$ et $b$ deux éléments d'un anneau commutatif $A$, on dit que $a$ est \textbf{associé} à $b$ s'il existe un élément inversible $u\in A^\times$ tel que $a=bu$. La relation d'association est une relation d'équivalence sur $A$.
\end{dfn}
\begin{exa}
	L'ensemble des associés à $1$ est exactement égal au groupe des unités de $A$.
\end{exa}
\begin{prop}
	Un morphisme d'anneaux $f:A\rightarrow B$ (non nécessairement commutatifs) induit par restriction et corestriction un morphisme de groupes de $A^\times$ dans $B^\times$.
\end{prop}
\begin{proof}
	Soit $f:A\rightarrow B$ un morphisme d'anneaux. Soit $a\in A^\times$ et $b$ l'inverse de $a$. Comme $ab=ba=1$, on a $1=f(1)=f(a)f(b)=f(b)f(a)$. par suite, $f(a)$ est inversible, d'inverse $f(b)$. On peut donc définir la restriction de $f$ à $A^\times$ et corestreinte à $B^\times$. C'est clairement un morphisme de groupes par héritage du morphisme d'anneaux.
\end{proof}
\begin{dfn}[Élément nilpotent, indice de nilpotence]
	Soit $A$ un anneau (non nécessairement commutatif). On dit que $a\in A$ est \textbf{nilpotent} lorsqu'il existe $n\in\N\st$ tel que $a^n=0$. On appelle alors \textbf{indice de nilpotence de $a$} le plus petit entier $p\in\N\st$ vérifiant cette propriété.
\end{dfn}
\subsection{Idéaux}
Dans tous ces compléments, on traite uniquement le cas des anneaux commutatifs. Cette hypothèse permet de simplifier de nombreuses notions et de nombreuses preuves, mais il existe une théorie des idéaux à droite et à gauche.
\begin{dfn}[Idéal]
	Un \textbf{idéal} $I$ d'un anneau $A$ est une partie de $A$ vérifiant les deux points suivants : 
	\begin{itemize}
		\item L'ensemble $I$ est un sous-groupe de $(A,+)$.
		\item Pour tout $a\in I$ et tout $b\in A$, on a $ba\in I$. On dit que $I$ est \textbf{absorbant}.
	\end{itemize}
\end{dfn}
\begin{exa}
	$\left\{0\right\}$ et $A$ sont des idéaux (triviaux) de $A$.
\end{exa}
\begin{dfn}[Idéal engendré par un élément]
	Pour $a\in A$, $(a):=\left\{ax \ | \ x\in A\right\}$ est un idéal. On dit que c'est \textbf{l'idéal engendré par $a$}. On le note aussi $aA$.
\end{dfn}
\begin{prop}[Recette de l'idéal]
	Comme $(-1_A)$ est un élément de $A$, pour prouver qu'une partie $I$ de $A$ est un idéal, il suffit de montrer les propriétés suivantes : 
	\begin{enumerate}
		\item $0\in I$
		\item $\forall (a,b)\in I^2,\ a+b\in I$
		\item $\forall (a,b)\in A\times I,\ ab\in I$
	\end{enumerate}
\end{prop}
\begin{proof}
	Le troisième point donne l'absorbance. Le troisième point donne aussi la stabilité par passage à l'opposé car $(-1_A)\in A$. Les deux premiers points, couplés à la stabilité par passage à l'inverse, donnent le caractère de sous-groupe. 
\end{proof}
\begin{exa}
	Si $\K$ est un corps, ses seuls idéaux sont les idéaux triviaux. En effet, soit $I$ un idéel à gauche de $\K$ distinct de de $\left\{0\right\}$ et soit $a$ un élément non nul de $I$. Soit $b\in\K$. Comme $a$ est non nul, on peut considérer l'élément $ba\inv$ de $\K$ et, par définition d'un idéal à gauche $(ba\inv)a\in\K$. On a donc $b\in I$ puis $I=\K$.
\end{exa}
\begin{exa}
	Plus généralement, si un idéal contient un élément inversible, alors il est égal à $A$ tout entier.
\end{exa}
\begin{prop}[Idéaux de $\Z$]
	Si $I$ est un idéal de $\Z$, alors il existe un unique $n\geqslant0$ tel que $I=(n)$.
\end{prop}
\begin{proof}
	On a déjà prouvé ceci car tout sous-groupe de $(\Z,+)$ est en fait aussi un idéal.
\end{proof}
\subsection{Opérations sur les idéaux}
On dispose d'un certain nombre d'opérations intéressantes sur les idéaux.
\begin{prop}[Intersection d'idéaux]
	Si $I$ et $J$ deux idéaux de $A$, alors l'ensemble $I\cap J$ est encore un idéal de $A$. Plus généralement, l'intersection d'une famille d'idéaux de $A$ est encore un idéal de $A$.
\end{prop}
\begin{proof}
	Il suffit de revenir à la définition. C'est exactement la même démonstration que pour l'intersection d'un nombre quelconque de sous-groupes.
\end{proof}
\begin{prop}[Idéal engendré par une partie]
	Si $S$ est une partie de $A$, il existe un plus petit idéal (pour l'inclusion) $I$ de $A$ contenant $S$, qu'on appelle \textbf{idéal engendré par $S$}. On le note $<S>$. Cela signifie que $I$ est un idéal contenant $S$ contenant $S$ et que si $J$ est un idéel contenant $S$, alors $J$ contient $I$. 
\end{prop}
\begin{proof}
	Comme dans le cas des groupes, il suffit de définir $I$ comme l'intersection des idéaux de $A$ qui contiennent $S$. C'est un idéal d'après la proposition précédente et il contient clairement $S$. C'est aussi le plus petit pour l'inclusion.
\end{proof}
\begin{dfn}[Famille presque nulle]
	Soit $(a_S)_{s\in S}$ une famille d'éléments de $A$. On dit qu'elle est \textbf{presque nulle} s'il existe une partie finie $S'$ de $S$ telle que $$\forall s\in S\setminus S',\ a_s=0$$
\end{dfn}
\begin{rmq}
	Dans le cas où $S$ est elle-même une partie finie, toutes la familles d'éléments de $A$ indexées par $S$ sont presque nulles.
\end{rmq}
\begin{prop}
	Avec les mêmes notations que dans la proposition précédente, $I=<S>$ est l'ensemble des combinaisons linéaires à coefficients dans $A$ presque nulles d'éléments de $S$, c'est-à-dire l'ensemble des éléments de la forme $$\sum_{s\in S}a_ss$$ avec $(a_s)_{s\in S}$ une famille d'éléments de $A$ presque nulle.
\end{prop}
\begin{proof}
	Notons $J$ l'ensemble des combinaisons linéaires presque nulles d'éléments de $S$ à coefficients dans $A$. Si $(a_s)_{s\in S}$ est une famille presque nulle d'éléments de $A$, alors $\displaystyle\sum_{s\in S}a_ss$ est un élément de tout idéal de $A$ contenant $S$ puisqu'un tel idéal est stable par $+$ et absorbant. Il reste donc à prouver que $J$ est un idéal de $A$. D'une part, il contient $0=\displaystyle\sum_{s\in S}0\times s$. D'autre part, si $x=\displaystyle\sum_{s\in S}a_ss$ et  $y=\displaystyle\sum_{s\in S}b_ss$ sont deux éléments de $J$, alors la famille $(a_s+b_s)_{s\in S}$ est presque nulle et $x+y\in J$. Enfin, pour $a\in A$ et $x=\displaystyle\sum_{s\in S}b_ss$ un élément de $J$, on a $ax=\displaystyle\sum_{s\in S}(ab_s)s$ ce qui prouve que $ax \in J$. Finalement, on a montré que $J$ est bien un idéal de $A$, et il est clair qu'il contient chaque élément de $S$ en prenant pour $(a_s)_{s\in S}$ la famille nulle sauf pour un élément $a_{s_0}$ qui vaut $1_A$.
\end{proof}
\begin{prop}
	Pour $a\in A$, l'idéal engendré par $\left\{a\right\}$ est noté $(a)$ et on a $(a)=aA$.
\end{prop}
\begin{proof}
	Immédiat par la proposition précédente.
\end{proof}
\begin{prop}
	Le noyau d'un morphisme d'anneaux $f:A\rightarrow B$ est un idéal.
\end{prop}
\begin{proof}
	Un morphisme d'anneaux étant un morphisme de groupes abéliens, $\ker(f)$ est un sous-groupe de $A$. De plus, pour $x\in\ker(f)$ et $a\in A$, on a $f(ax)=f(a)f(x)=f(a)\times 0=0$ donc $ax\in\ker(f)$.
\end{proof}
\begin{prop}[Image réciproque d'un idéal]
	Soit $f:A\rightarrow B$ un morphisme d'anneaux et $J$ un idéal de $B$. l'image réciproque de $J$ par $f$ est un idéal de $A$.
\end{prop}
\begin{proof}
	Comme $f(0)=0\in J$, on a $0\in f\inv(J)$. Soit $(a,b)\in f\inv(J)^2$. On a $f(a+b)=f(a)+f(b)\in J$ car $f(a)$ et $f(b)$ sont dans $J$ puisque $J$ est stable par $+$. Enfin, pour $a\in A$ et $b\in f\inv(J)$, $f(ab)=f(a)f(b)\in J$ puisque $f(b)\in J$ et $J$ est absorbant.
\end{proof}
\begin{rmq}
	En revanche, l'image d'un idéal par un morphisme d'anneaux n'est pas forcément un idéal.
\end{rmq}
\begin{prop}[Somme de deux idéaux]
	Soit $I$ et $J$ deux idéaux de $A$. L'ensemble des sommes $a+b$ avec $a\in I$ et $b\in J$ est un idéal de $A$, noté $I+J$. C'est aussi l'idéal engendré par la partie $I\cup J$.
\end{prop}
\begin{prop}[Somme quelconque d'idéaux]
	Plus généralement, si $(I_s)_{s\in S}$ est une famille d'idéaux de $A$, l'ensemble des sommes presque nulles $\displaystyle\sum_{s\in S}a_s$ où, pour tout $s\in S$, $a_s\in I_s$ est un idéal de $A$, noté $\displaystyle\sum_{s\in S}I_s$. C'est aussi l'idéal de $A$ engendré par la partie $\displaystyle\bigcup_{s\in S}I_s$ 
\end{prop}
\begin{proof}
	On note $I=\displaystyle\sum_{s\in S}I_s$.
	\begin{itemize}
		\item Comme $0=\displaystyle\sum_{s\in S} 0$ et comme $0\in I_s$ pour tout $s$, $0\in I$.
		\item Soit $\displaystyle a=\sum_{s\in S} a_s$ et $\displaystyle b=\sum_{s\in S} b_s$ deux éléments de $I$. On a $\displaystyle a+b=\sum_{s\in S} a_s+b_s$, et cette somme est, bien entendu, presque nulle. On a donc $a+b\in I$ puisque $a_s+b_s$ pour tout $s\in S$ par stabilité de $I_s$ par $+$.
		\item Soit $\displaystyle a=\sum_{s\in S} a_s$ un élément de $I$ et $b\in A$. $\displaystyle ba=\sum_{s\in S} ba_s$ avec $s\in S$, $bas\in I_s$ puisque $I_s$ est absorbant.
		\item On a montré que $I$ est un idéal de $A$.
		\item Reste à prouver $I$ est l'idéal de $A$ engendré par la partie $\displaystyle\bigcup_{s\in S}I_s$. Nous devons établir deux inclusions. Comme $I$ est un idéal de $A$ contenant la réunion des $I_s$, on a déjà $\displaystyle\left\langle\bigcup_{s\in S}I_s\right\rangle\subset I$. D'autre part, $\displaystyle\left\langle\bigcup_{s\in S}I_s\right\rangle$ est un idéal de $A$ contenant chaque $I_s$ donc il est stable par $+$ et absorbant ce qui montre que tout somme presque nulle $\displaystyle\sum_{s\in S}a_s$ où, pour tout $s\in S$, $a_s\in I_s$, est aussi un élément de $\displaystyle\left\langle\bigcup_{s\in S}I_s\right\rangle$.
	\end{itemize}
\end{proof}
\begin{dfn}[Idéal produit]
	Soit $I$ et $J$ deux idéaux d'un anneau commutatif $A$. On appelle \textbf{idéal produit de $I$ et $J$}, et on note $I.J$, l'idéal engendré par l'ensemble des $xy$ où $x\in I$ et $y\in J$. C'est donc l'ensemble des sommes $\displaystyle\sum_{k\in E}x_ky_k$ avec $E$ fini, les $x_k$ dans $I$ et les $y_k$ dans $J$. 
\end{dfn}
\begin{rmq}
	Attention ! La notation $IJ$ peut désigner l'ensemble des produits $xy$ pour $x\in I$ et $j\in J$, mais cette structure n'est pas intéressante car n'est même pas un groupe en général. On prendra donc garde à noter l'idéal produit $I.J$ avec le point entre les deux idéaux, c'est le plus petit idéal qui contient alors $IJ$.
\end{rmq}
\begin{prop}
	Soit $I$ et $J$ deux idéaux d'un anneau commutatif $A$. On a $I.J\subset I\cap J$.
\end{prop}
\begin{proof}
	Il suffit de remarquer que $IJ\subset I\cap J$ par absorbance (avec $IJ$ l'ensemble des produits $xy$ pour $x\in I$ et $y\in J$). Puis $I.J$ est le plus petit idéal contenant $IJ$ par définition, donc $I.J\subset I\cap J$.
\end{proof}
\begin{prop}
	Soit $I$ et $J$ deux idéaux d'un anneau commutatif $A$. On a : $$(I+J).(I\cap J)\subset I.J$$
\end{prop}
\begin{proof}
	Soit $a=i+j\in I+J$ et $b\in I\cap J$. Alors $ab=ib+jb\in I.J$ car $b\in J$ et $b\in I$. Puisque $(I+J).(I\cap J)$ est l'ensemble des sommes finies d'éléments de ce type et $I.J$ est un idéal, on a bien $(I+J).(I\cap J)\subset I.J$.
\end{proof}
\begin{prop}
	Soit $I$ et $J$ deux idéaux d'un anneau commutatif $A$. Si $I$ et $J$ sont comaximaux (\textit{ie} $I+J=A$, cf. le paragraphe $2$), alors $I.J=I\cap J$.
\end{prop}
\begin{proof}
	En vertu d'une proposition précédente, il suffit de prouver que $I\cap J\subset I.J$. Or, d'après la proposition précédente et l'hypothèse, on a $A.(I\cap J)\subset I.J$. Puis, si on se donne $x\in I\cap J$, alors $1x\in A.(I\cap J)$ donc $x\in I\cap J$. On a alors $I.J=I\cap J$.
\end{proof}
\begin{dfn}[Idéal quotient]
	Soit $I$ et $J$ deux idéaux d'un anneau commutatif $A$. On appelle \textbf{idéal quotient de $I$ par $J$}, et on note $I:J$ l'ensemble défini par : $$I:J=\left\{x\in A\ | \ xJ\subset I\right\}$$
\end{dfn}
\begin{proof}
	Montrons qu'il mérite bien son nom d'idéal.
	\begin{itemize}
		\item $0J=\left\{0\right\}\subset  I$ donc $0\in I:J$.
		\item Soit $(a,b)\in\left(I:J\right)^2$. Soit $x\in(a+b)J$. Fixons $y\in J$ tel que $x=(a+b)y=ay+by$. Par hypothèse, $ay\in I$ et $by\in I$, donc par stabilité de $I$ par $+$, $x=(a+b)y\in I$. Donc $(a+b)J\subset I$. Donc $a+b\in I:J$.
		\item Soit $a\in I:J$ et $b\in A$. Soit $x\in abJ$. Fixons $y\in J$ tel que $x=aby$. Comme $J$ est un idéal, $by\in J$ donc $x\in aJ$. Donc $x\in I$. Donc $abJ\subset I$. Donc $ab\in I:J$.
	\end{itemize}
	Ainsi, $I:J$ est bien un idéal.
\end{proof}
\subsection{Radical d'un idéal}
\begin{dfn}[Nilradical]
	Le \textbf{nilradical} d'un anneau commutatif est l'ensemble de ses éléments nilpotents.
\end{dfn}
\begin{dfn}[Radical d'un idéal]
	Plus généralement, on définit le \textbf{radical} d'un idéal $I$ de $A$ commutatif par $$\sqrt{I}:=\left\{a\in A\ | \ \exists n\in\N\st,\ a^n\in I\right\}$$ Le nilradical de $A$ est donc le radical de l'idéal nul.
\end{dfn}
\begin{prop}
	Le radical $\sqrt{I}$ d'un idéal $I$ de $A$ est un idéal de $A$ qui contient $I$.
\end{prop}
\begin{proof}
	On voit immédiatement qu'on a $I\subset\sqrt{I}$ puisqu'il suffit de prendre $n=1$ dans la définition de $\sqrt{I}$ pour tous les éléments de $I$. Montrons désormais que $\sqrt{I}$ est un idéal.
	\begin{itemize}
		\item Comme $0^1=0\in I$, $0\in\sqrt{I}$.
		\item Soit $a\in\sqrt{I}$ et $b\in\sqrt{I}$. Il existe $n$ et $m$ dans $\N\st$ tels que $a^n\in I$ et $b^m\in I$, fixons-les. Comme l'anneau est commutatif, on peut utiliser la formule du binôme de Newton : $$(a+b)^{n+m}=\sum_{k=0}^{n+m}\binom{n+m}{k}a^kb^{n+m-k}$$ Dans cette somme, tous les termes appartiennent à $I$. En effet, c'est vrai pour ceux d'indice $k\geqslant n$ puisque $a^k=a^na^{n-k}$ et $a^n\in I$. De même, pour ceux d'indice $k\leqslant n$, on a $n+m-k\geqslant m$ et b${n+m-k}=b^mb^{n-k}$ et $b^m\in I$. On a donc $(a+b)^{n+m}\in I$, d'où $a+b\in I$.
		\item Enfin, si $a\in\sqrt{I}$ et $b\in A$, choisissons $n\in\N\st$ tel que $a^n\in I$. Alors, $(ba)^n=b^na^n\in I$ et $ba\in\sqrt{I}$.
		\item $\sqrt{I}$ est donc bien un idéal de $A$.
	\end{itemize}
\end{proof}
\begin{exa}
	En particulier, le nilradical de d'un anneau commutatif est donc un idéal.
\end{exa}
\begin{dfn}[Idéal radiciel]
	Soit $A$ un anneau commutatif. On dit qu'un idéal $I$ de $A$ est \textbf{radiciel} lorsqu'il est égal à son radical, $ie$ $I$ est radiciel lorsque $I=\sqrt{I}$.
\end{dfn}
\begin{prop}[Radicialité des radicaux]
	Soit $A$ un anneau commutatif et $I$ un idéal de $A$. On a $\sqrt{\sqrt{I}}=\sqrt{I}$. Autrement dit, le radical d'un idéal est radiciel.
\end{prop}
\begin{proof}
	On a déjà l'inclusion $\sqrt{I}\subset\sqrt{\sqrt{I}}$. Montrons l'inclusion réciproque. Soit $x\in\sqrt{\sqrt{I}}$. Par hypothèse, il existe $n\in\N\st$ tel que $x^n\in\sqrt{I}$. Fixons $n$. Or, puisque $x^n\in\sqrt{I}$, il existe $m\in\N\st$ tel que $(x^n)^m\in I$. Alors, $nm\in\N\st$ et $x^{nm}\in I$ donc $x\in\sqrt{I}$. On a bien $\sqrt{\sqrt{I}}\subset\sqrt{I}$, d'où lr résultat.
\end{proof}
\begin{prop}[Croissance des radicaux]
	Soit $A$ un anneau commutatif et $I$ et $J$ deux idéaux de $A$. SI $I\subset J$, alors : $$\sqrt{I}\subset\sqrt{J}$$
\end{prop}
\begin{proof}
	Supposons que $I\subset J$. Soit $x\in \sqrt{I}$. Il existe $n\in\N\st$ tel que $x^n\in I$. Fixons-le. Alors $x^n\in J$ par hypothèse, et comme $n\in\N\st$, $x\in\sqrt{J}$. On a donc $\sqrt{I}\subset\sqrt{J}$.
\end{proof}
\begin{prop}
	Soit $I$ et $J$ deux idéaux d'un anneau commutatif $A$. On a : $$\sqrt{I}+\sqrt{J}\subset\sqrt{I+J}$$
\end{prop}
\begin{proof}
	Soit $x\in\sqrt{I}$ et $y\in\sqrt{J}$. Fixons $m,n\in\N\st$ tels que $x^m\in I$ et $y^n\in J$. Grâce au binôme de Newton puis en factorisant, on montre aisément que $(a+b)^{m+n}\in I+J$ donc $a+b\in\sqrt{I+J}$ ce qui conclut.
\end{proof}
\begin{prop}[Radical et intersection]
	Soit $A$ un anneau commutatif et $I$ et $J$ deux idéaux de $A$. On a : $$\sqrt{I\cap J}=\sqrt{I}\cap\sqrt{J}$$
\end{prop}
\begin{proof}
	Déjà, $I\cap J$ est un idéal tel que $I\cap J\subset I$ et $I\cap J\subset J$. Par croissance, $\sqrt{I\cap J}\subset\sqrt{I}$ et $\sqrt{I\cap J}\subset\sqrt{J}$ donc $\sqrt{I\cap J}\subset\sqrt{I}\cap\sqrt{J}$. Montrons l'inclusion réciproque. Soit $x\in\sqrt{I}\cap\sqrt{J}$. Alors il existe $m\in\N\st$ et $n\in\N\st$ tels que $x^n\in I$ et $x^m\in J$. Alors, par absorbance de $I$ et $J$, $x^{\max(m,n)}\in I\cap J$. Or, $\max(m,n)\in\N\st$, donc $x\in\sqrt{I\cap J}$. On a bien $\sqrt{I\cap J}=\sqrt{I}\cap\sqrt{J}$.
\end{proof}
\begin{prop}
	Soit $I$ et $J$ deux idéaux d'un anneau commutatif $A$. On a : $$\sqrt{I.J}=\sqrt{I\cap J}=\sqrt{I}\cap\sqrt{J}$$
\end{prop}
\begin{proof}
	Il suffit de montrer la première égalité. On sait que $I.J\subset I\cap J$ donc par croissance des radicaux, on a $\sqrt{I.J}\subset\sqrt{I\cap J}=\sqrt{I}\cap\sqrt{J}$. Réciproquement, soit $x\in\sqrt{I}\cap\sqrt{J}$. Fixons $m,n\in\N\st$; tels que $x^m\in I$ et $x^n\in J$. Alors $x^{m+n}\in IJ$ donc $x^{m+n}\in I.J$. Donc $x\in\sqrt{I.J}$. On a bien $$\sqrt{I.J}=\sqrt{I\cap J}=\sqrt{I}\cap\sqrt{J}$$
\end{proof}
\section{Anneaux quotients}
Étant donné un anneau et une relation d'équivalence convenable sur cet anneau, l'objectif est de munir l'ensemble des classes d'équivalence d'une structure d'anneau. Cela revient en fait à "rendre nuls" les éléments d'un idéal de l'anneau sans modifier les autres règles de calcul.
\subsection{Construction}
Rappelons qu'une relation $\mathcal{R}$ sur un ensemble $X$ est dite relation d'équivalence si elle est réflexive, symétrique et transitive. L'ensemble des classes d'équivalence de $X$ pour la relation $\mathcal{R}$ est noté $X/\mathcal{R}$. \\ \par Soit maintenant un anneau $A$. On peut alors chercher les relations d'équivalence sur $A$ qui sont \textbf{compatibles} avec la structure d'anneau. Autrement dit, on veut que soit vérifiée la propriété suivante : $$\forall (x,y,x',y')\in A^4,\ (x\mathcal{R}y \land x'\mathcal{R}y')\implies((x+x')\mathcal{R}(y+y')\land(xx')\mathcal{R}(yy'))$$
\begin{prop}[Analyse]
	Une relation d'équivalence compatible avec la structure d'anneau est nécessairement de la forme $\mathcal{R}_I$ avec $I$ un idéal et $$\forall (x,y)\in A^2,\ x\mathcal{R}_Iy\iff x-y\in I$$
\end{prop}
\begin{proof}
	Notons $I$ la classe de $0$. Si $x\mathcal{R}y$, comme $(-y)\mathcal{R}(-y)$, alors $x-y\mathcal{R}0$ par compatibilité avec l'addition, soit $x-y\in I$. Réciproquement, si $x-y\in I$, on a $(x-y)\mathcal{R}0$ et, puisque $y\mathcal{R}y$, il vient $x\mathcal{R}y$. Ce petit calcul montre que la relation $\mathcal{R}$ est définie par $x\mathcal{R}y$ si, et seulement si, $x-y\in I$. \\ \par Montrons par ailleurs que $I$ est un idéal de $A$. On a déjà $0\in I$. De plus, si $x\in I$ et $y\in I$, il vient $x\mathcal{R}0$ et $y\mathcal{R}0$ donc $(x+y)\mathcal{R}0$ ce qui prouve que $x+y\in I$. Enfin, si $x\in I$ et $a\in A$, on a $x\mathcal{R}0$ et $a\mathcal{R}a$ donc $(ax)\mathcal{R}(a\times 0)$ ce qui montre que $ax\in I$. Donc $I$ est bien un idéal.
\end{proof}
\begin{prop}[Synthèse]
	Soit $A$ un anneau commutatif et $I$ un idéal de $A$. La relation $\mathcal{R}$ définie sur $A$ par $$\forall (x,y)\in A^2,\ x\mathcal{R}y\iff x-y\in I$$ est une relation d'équivalence compatible avec la structure d'anneau. l'ensemble quotient $A/\mathcal{R}$ possède alors une structure canonique d'anneau. \\ \par De plus, la surjection canonique $\pi:A\rightarrow A/\mathcal{R}$ est un morphisme d'anneaux surjectif de noyau $I$. \\ \par L'anneau quotient $A/\mathcal{R}$ est noté $A/I$ et, pour $a\in A$, on note $\overline{a}$ la classe d'équivalence de $a$ modulo $I$ s'il n'y a pas d'ambiguïté.
\end{prop}
\begin{proof}
	Il suffit d'effectuer les mêmes calculs dans l'autre sens.
\end{proof}
\subsection{Idéaux comaximaux}
\begin{dfn}[Idéaux comaximaux]
	Deux idéaux $I$ et $J$ d'un anneau $A$ sont dits \textbf{comaximaux} lorsque $I+J=A$.
\end{dfn}
\begin{prop}[Théorème chinois]
	On considère un anneau $A$ ainsi que deux idéaux $I$ et $J$ de $A$ comaximaux. Dans ce cas, le morphisme $$\begin{array}{ccccc}f&:&A&\rightarrow&(A/I)\times(A/J)\\ &&a&\mapsto&(Cl_I(a),Cl_J(a))\end{array}$$ est surjectif et son noyau est l'idéal $I\cap J$. Il en résulte, par passage au quotient, un isomorphisme d'anneaux : $$\begin{array}{ccccc}\tilde{f}&:&A/(I\cap J)&\rightarrow&(A/I)\times(A/J)\\ &&Cl_{I\cap J}(a)&\mapsto&(Cl_I(a),Cl_J(a))\end{array}$$
\end{prop}
\begin{proof}
	Déjà, $f$ est clairement un morphisme d'anneaux pour la structure d'anneau produit puisque la surjection canonique est morphisme d'anneaux. Ensuite, si on se donne $Cl_I(x)\in A/I$ et $Cl_J(y)\in A/J$, alors on peut fixer $i\in I$ et $j\in J$ tel que $x-y=i+j$ puisque $I$ et $J$ sont comaximaux. On pose alors $a:=x-i=y+j$. Puisque $a-x=-i\in I$ et $a-y=j\in J$, on a bien $$f(a)=(Cl_I(x),Cl_J(y))$$ donc $f$ est surjectif. Ensuite, soit $a\in A$. $a$ appartient à $\ker(f)$ si, et seulement si, $Cl_I(a)=I$ et $Cl_J(a)=J$, donc $a$ appartient à $\ker(f)$ si, et seulement si, $a\in I\cap J$. Donc on a bien $\ker(f)=I\cap J$. \\ \par Pour $\tilde{f}$, déjà cette fonction est bien définie car deux éléments de la même classe pour $I\cap J$ ont \textit{a fortiori} même classe pour $I$ et pour $J$. Ensuite, es propriétés de morphisme et la surjectivité s'héritent de $f$. Enfin, $\tilde{f}$ est injective car son noyau est réduit à la classe de $I\cap J$, le neutre pour $+$ de l'anneau quotient $A/(I\cap J)$.
\end{proof}
\begin{cor}
	Soit $I$ et $J$ deux idéaux comaximaux d'un anneau $A$. Pour tout couple $(x,y)\in A^2$, il existe $a\in A$ tel que $a\in x+I$ et $a\in y+J$.
\end{cor}
\begin{proof}
	Il suffit d'utiliser la surjectivité du morphisme $f$.
\end{proof}
\subsection{Idéaux maximaux}
\begin{dfn}[Idéaux propres]
	Soit $A$ un anneau commutatif. On appelle \textbf{idéal propre} de $A$ tout idéal de $A$ strictement inclus dans $A$.
\end{dfn}
\begin{dfn}[Idéaux maximaux]
	Soit $A$ un anneau commutatif. Un idéal $I$ de $A$ est dit \textbf{maximal} lorsqu'il est maximal parmi les idéaux propres (au sens de l'inclusion). En d'autres termes, $I$ est un idéal propre et si $J$ est un idéal propre tel que $I\subset J$, alors $J=I$.
\end{dfn}
\begin{prop}
	Soit $A$ un anneau commutatif. Un idéal $I$ est maximal si, et seulement si, $A/I$ est un corps.
\end{prop}
\begin{proof}
	On prouve successivement les deux implications.
	\begin{itemize}
		\item $\Longrightarrow$ : Supposons que $I$ est maximal. On sait déjà que $A/I$ est un anneau commutatif, il reste donc à prouver l'inversibilité des classes non nulles. Soit $\overline{x}\in A/I$ telle que $\overline{x}\ne\overline{0}$. Autrement dit, $x\notin I$. L'ensemble $I+xA$ est alors un idéal. Puisque $I\subset I+xA$ et $x\notin I$, alors $I\subsetneq I+xA$. Par maximalité de $I$, $I+xA$ ne peut pas être un idéal propre sans quoi on aurait $x\in I$. Donc les idéaux $I$ et $xA$ sont comaximaux, $ie$ $I+xA=A$. On peut donc trouver $i\in I$ et $a\in A$ tels que $i+xa=1$. En passant au quotient, on obtient $\overline{x}\overline{a}=\overline{1}$ si bien que $\overline{x}$ est inversible. Ainsi, $A/I$ est un corps.
		\item $\Longleftarrow$ : Réciproquement, supposons que $A/I$ soit un corps et montrons que $I$ est un idéal maximal. Soit $J$ un idéal de $A$ tel que $I\subsetneq J$. Montrons que $J=A$. Soit $x\in J\setminus I$. Alors $x\notin I$ donc $\overline{x}\ne\overline{0}$. Il s'ensuit que $\overline{x}$ est inversible dans $A/I$ car $A/I$ est un corps. On peut donc fixer $y\in A$ tel que $\overline{x}\overline{y}=\overline{1}$. On peut donc fixer $z\in I$ tel que $xy-1=z$. On a donc $1=xy-z$. Or, $x\in J$ donc $xy\in J$ et $z\in I$ donc $z\in J$. Donc $1\in J$ puis $J=A$. Ainsi, nécessairement, $I$ est un idéal maximal.
	\end{itemize}
\end{proof}
\begin{exa}
	Les idéaux maximaux de $\Z$ sont exactement les $n\Z$ où $n$ est premier. En effet, soit $I$ un idéal de $\Z$. Il est de la forme $I=n\Z$ avec $n\in\N$. Or, il est maximal si, et seulement si, $\Z/n\Z$ est un corps, c'est-à-dire si, et seulement si, $n$ est un nombre premier.
\end{exa}
\subsection{Idéaux premiers}
\begin{dfn}[Idéal premier]
	Soit $A$ un anneau commutatif. Un idéal $I$ de $A$ est dit premier lorsque tout produit d'éléments n'appartenant pas à $I$ n'appartient pas à $I$, autrement dit lorsque $A\setminus I$ est stable par $\times$. Cela revient à dire que $A/I$ est intègre.
\end{dfn}
\begin{exa}
	Dans $\Z$, l'idéal $\left\{0\right\}$ est premier car $\Z$ est un anneau intègre.
\end{exa}
\begin{prop}
	Soit $A$ un anneau commutatif. Alors tout idéal maximal est premier.
\end{prop}
\begin{proof}
	Soit $I$ un idéal maximal de $A$. Soit $(x,y)\in(A\setminus I)^2$. Montrons que $xy\in A\setminus I$. On sait d'après le paragraphe précédent que $A/I$ est un corps. Or, $\overline{x}\ne\overline{0}$ et $\overline{y}\ne\overline{0}$ donc, par intégrité du corps $A/I$, $\overline{xy}=\overline{x}\overline{y}\ne\overline{0}$. Donc $xy\notin I$, ce qui achève la preuve. Plus succinctement, si $I$ est maximal alors $A/I$ est un corps donc $A/I$ est intègre donc $I$ est premier.
\end{proof}
\begin{rmq}
	Attention, la réciproque est fausse ! Par exemple, $\left\{0\right\}$ est un idéal premier de $\Z$ mais il n'est pas maximal car il est contenu strictement par l'idéal $2\Z$ qui est un idéal propre de $\Z$.
\end{rmq}
\begin{prop}
	Tout idéal premier est radiciel.
\end{prop}
\begin{proof}
	Soit $A$ un anneau commutatif et $I$ un idéal premier de $A$. On sait déjà que $I\subset\sqrt{I}$. Montrons l'inclusion réciproque. $I$ est premier donc $A/I$ est intègre. Soit $x\in\sqrt{I}$ et fixons $n\in\N\st$ tel que $x^n\in I$. Alors $\overline{x}^n=\overline{0}$ donc par intégrité de $A/I$, $\overline{x}=\overline{0}$. Donc $x\in I$, si bien que $\sqrt{I}\subset I$. Un idéal premier est donc bien radiciel.
\end{proof}
\begin{rmq}
	Attention : la réciproque est fausse !
\end{rmq}
\section{L'anneau $\Z/n\Z$}
\begin{dfn}[$\Z/n\Z$]
	Pour $n\in\N\st$, on note $\Z/n\Z$ l'anneau quotient de $\Z$ par l'idéal $n\Z$. Pour $x\in\Z$, on note $\overline{x}$ la classe de $x$. On a alors $$\Z/n\Z=\left\{\overline{k}\ | \ k\in\llbracket0,n-1\rrbracket\right\}$$ De plus, la surjection canonique $\pi:\Z\rightarrow\Z/n\Z$ est un morphisme d'anneaux.
\end{dfn}
\begin{rmq}
	De façon générale, il existe un unique morphisme d'anneaux $f$ de $\Z$ dans un anneau $A$. En effet, $\Z$ est engendré par $1$, comme $f(1)=1$, cette égalité permet de définir sans ambiguïté $f$.
\end{rmq}
\begin{prop}[Inversibles de $\Z/n\Z$]
	L'ensemble des éléments inversibles de $\Z/n\Z$ est $$(\Z/n\Z)^\times=\left\{\overline{k}\ | \ k\in\llbracket0,n-1\rrbracket,\ k\wedge n=1\right\}$$
\end{prop}
\begin{proof}
	Il suffit d'utiliser le théorème de Bézout.
\end{proof}
\begin{dfn}[Indicatrice d'Euler]
	On appelle \textbf{fonction indicatrice d'Euler} la fonction $\varphi$ de $\N\st$ dans $\N$ qui à $n$ associe le nombre d'éléments inversibles de l'anneau $\Z/n\Z$, c'est-à-dire le nombre d'entiers $p\in\llbracket0,n-1\rrbracket$ premiers avec $n$.
\end{dfn}
\begin{rmq}
	L'indicatrice d'Euler joue un grand rôle en arithmétique et en codage informatique comme par exemple dans le chiffrement RSA.
\end{rmq}
\begin{exa}
	$\varphi(2)=1$, $\varphi(3)=2$, $\varphi(4)=2$. Pour $p$ un nombre premier, on a $\varphi(p)=p-1$. On a même la réciproque pour $p\geqslant2$. Pour $\alpha\in\N\st$ et $p$ premier, on a $\varphi(p^\alpha)=p^{\alpha-1}(p-1)$.
\end{exa}
\begin{prop}
	Soit $n\in\N\st$. Les trois assertions suivantes sont équivalentes :
	\begin{enumerate}
		\item $n$ est premier
		\item $\Z/n\Z$ est un corps
		\item $\Z/n\Z$ est un anneau intègre
	\end{enumerate}
\end{prop}
\begin{proof}
	Pour $1\implies2$, se donner $m\in\llbracket1,m-1\rrbracket$. $n$ et $m$ sont alors premiers entre eux et il suffit d'appliquer le théorème de Bézout avant de passer au quotient. $2\implies3$ est trivial, cela est toujours vrai. Pour $3\implies1$, supposons par l'absurde que $n$ n'est pas premier et écrivons $n=ab$ avec $1<a,b<n$. Alors $\overline{n}=\overline{a}\overline{b}=\overline{0}$ donc $\overline{a}=0$ ou $\overline{b}=0$ par intégrité, contradiction. Donc $n$ est premier.
\end{proof}
\begin{prop}[Théorème de Fermat-Euler]
	Soit $n\in\N\st$. pour tout $a\in\Z$ tel que $a\wedge n=1$, on a : $$a^{\varphi(n)}\equiv 1\ [n]$$
\end{prop}
\begin{proof}
	Si $a$ est premier avec $n$, alors $\overline{a}$ appartient au groupes des unités $\mathcal{U}\left(\Z/n\Z\right)$, qui est de cardinal $\varphi(n)$. il suffit alors d'appliquer le "théorème de Fermat" vu sur les groupes (c'est le cas "facile" où le groupe est commutatif et où il suffit d'utiliser la bijection $x\mapsto ax$).
\end{proof}
\begin{rmq}
	Pour $a\in\Z$ tel que $a\wedge n=1$, on peut utiliser ce théorème pour déterminer un inverse modulaire de $a$ qui n'est autre que $a^{\varphi(n)-1}$.
\end{rmq}
\begin{prop}[Théorème de Wilson]
	Soit $p\in\N\st$. $p$ est premier si, et seulement si, $(p-1)!\equiv-1\ [p]$.
\end{prop}
\begin{proof}
	Si $(p-1)!\equiv -1\ [p]$, alors toute classe non nulle modulo $p$ est inversible. En effet, il suffit d'écrire $$\overline{x}\times\left(-\prod_{\substack{1\leqslant i \leqslant p-1 \\ \overline{i}\ne\overline{x}}}\overline{i}\right)=\overline{1}$$ Ainsi, $\Z/p\Z$ est un corps et $p$ est premier. \par Réciproquement, supposons que $p$ est premier. Le cas $p=2$ étant immédiat, passons à $p\geqslant3$. Toute classe $\overline{x}\ne\overline{0}$ est inversible puisque $\forall x\in\llbracket1,p-1\rrbracket,\ x\wedge p=1$. Cherchons les classes qui sont leur propre inverse. On a : $$\overline{x}^2=\overline{1}\iff(\overline{x}+\overline{1})(\overline{x}-\overline{1})=\overline{0}\iff \overline{x}=\pm\overline{1}$$ par intégrité dans le corps $\Z/p\Z$. A part ces deux classes, toute autre classe non nulle a un inverse qui n'est pas elle-même. Dans le produit $\overline{1}\times\ldots\times\overline{p-1}$, mettons à part $\overline{1}$ et $\overline{p-1}=\overline{-1}$, et regroupons les autres par paires d'inverses mutuels. On obtient alors le produit de $\overline{1}$ et de $\overline{-1}$, d'où le résultat.
\end{proof}
\begin{prop}[Théorème chinois]
	Soit $p$ et $q$ deux nombres entiers premiers entre eux ainsi que $(y,z)\in\Z^2$. Alors, il existe un entier $x\in\Z$ tel que $y\equiv y\ [p]$ et $x\equiv z\ [q]$. De plus, toutes les solutions de ce système sont congrues modulo $pq$. 
\end{prop}
\begin{proof}
	On peut démontrer ce résultat "à la main", mais il s'agit d'un résultat général. En effet, puisque $p\wedge q=1$, $p\Z+q\Z=\Z$ et ainsi les idéaux $p\Z$ et $q\Z$ sont comaximaux. Il suffit alors d'utiliser le théorème chinois vu dans la partie sur les idéaux comaximaux. On a en effet $p\Z\cap q\Z=pq\Z$ car $p$ et $q$ sont premiers entre eux.
\end{proof}
\begin{cor}
	Si $p$ et $q$ sont deux entiers premiers entre eux, alors il existe un isomorphisme d'anneaux de $Z/pq\Z$ sur $(\Z/p\Z)\times(\Z/q\Z)$ (muni de la structure produit).
\end{cor}
\begin{proof}
	Il s'agit encore là d'un résultat général, c'est la deuxième partie du théorème chinois sur les idéaux comaximaux. On peut évidemment le refaire à la main en comme dans la démontstration du théorème sur les idéaux comaximaux.
\end{proof}
\begin{cor}[Multiplicativité de $\varphi$]
	Si $p$ et $q$ sont deux entiers premiers entre eux, alors $\varphi(pq)=\varphi(p)\varphi(q)$. On dit que $\varphi$ est une fonction \textbf{arithmétique multiplicative}.
\end{cor}
\begin{proof}
	Il suffit de comparer les cardinaux de $(\Z/pq\Z)^\times$ et $((\Z/p\Z)\times(\Z/q\Z))^\times=(\Z/p\Z)^\times\times(\Z/q\Z)^\times$. En effet, on montre que la restriction et corestriction de l'isomorphisme fourni par le théorème précédent aux groupes des unités est un isomorphisme et on montre aisément l'égalité ensembliste énoncée précédemment.
\end{proof}
\begin{cor}[Formule explicite de l'indicatrice d'Euler]
	Pour $n\in\N\st$, on note $S$ la partie finie de $\mathcal{P}$ constituée des nombres premiers $p$ tel que $v_p(n)>0$. On a alors : $$\varphi(n)=\prod_{p\in S}p^{v_p(n)-1}(p-1)=n\prod_{p\in S}\left(1-\frac{1}{p}\right)$$
\end{cor}
\section{Anneaux principaux}
On considère dans ce paragraphe uniquement des anneaux intègres (et toujours commutatifs).
\subsection{Divisibilité dans un anneau intègre}
\begin{dfn}[Divisibilité]
	Soit $A$ un anneau intègre et $(a,b)\in A^2$. On dit que $a$ \textbf{divise} $b$ (ou encore que $b$ est un \textbf{multiple} de $a$) s'il existe $q\in A$ tel que $b=aq$. On note alors $a\mid b$.
\end{dfn}
La relation $\mid$ est réflexive et transitive (on dit que c'est une relation de \textbf{préordre} sur $A$) mais elle n'est pas antisymétrique en général. On a toutefois le résultat suivant :
\begin{prop}
	Soit $A$ un anneau intègre et $(a,b)\in A^2$. On a $$a\mid b\land b\mid a\iff \exists u\in A^\times,\ b=au$$ On a vu que cela revenait à dire que $a$ et $b$ sont associés. 
\end{prop}
\begin{proof}
	Le sens réciproque est immédiat. Pour le sens direct, si $a=0$ alors $b=0$ et il suffit de prendre $u=1$. Si $a\ne0$, il existe $(k,k')\in A^2$ tel que $a=bk$ et $b=ak'$. On a donc $a(1-kk')=0$ puis, par intégrité de $A$ et par commutativité, $k$ et $k'$ sont inversibles et conviennent.
\end{proof}
\begin{rmq}
	Rappelons que l'ensemble $(a)=aA$ des multiples de $a$ est l'idéal engendré par $a$. On peut alors reformuler le résultat précédent de la manière suivante : $$(a)=(b)\iff a\mid b\land b\mid a\iff \exists u\in A^\times,\ b=au$$
\end{rmq}
\begin{prop}
	Soit $(a,b,c)\in A^3$ avec $A$ un anneau commutatif et intègre. Si $c\mid a$ et $c\mid b$, alors $$\forall(u,v)\in A^2,\ c\mid au+bv$$
\end{prop}
\begin{proof}
	La vérification est immédiate.
\end{proof}
\subsection{Anneaux principaux}
\begin{dfn}[Idéal principal]
	Un idéal $I$ d'un anneau intègre $A$ est dit \textbf{principal} s'il est engendré par un seul élément, c'est-à-dire s'il existe $a\in A$ tel que $I=(a)$.
\end{dfn}
\begin{dfn}[Anneau principal]
	Un anneau intègre $A$ est dit \textbf{principal} lorsque tous ses idéaux sont principaux.
\end{dfn}
\begin{prop}
	$\Z$ est un anneau principal.
\end{prop}
\begin{proof}
	Déjà prouvé car $\Z$ est intègre et tous es idéaux sont de la forme $n\Z$.
\end{proof}
\begin{prop}
	$\K[X]$ est un anneau principal dès que le corps $\K$ est commutatif.
\end{prop}
\begin{proof}
	La preuve est réalisée dans le chapitre "Polynômes". 
\end{proof}
Ainsi, les deux objets que sont l'anneau des entiers relatifs $\Z$ et celui des polynômes à coefficients dans un corps commutatifs sont des anneaux principaux. \\ \par Dans les paragraphes suivants, on va définir le pgcd et le ppcm de deux éléments d'un anneau principal. On retrouvera alors les notions déjà vues dans $\Z$ et on utilisera la proposition précédente pour construire la même théorie pour les polynômes.
\subsection{PGCD dans un anneau principal}
\begin{dfn}[PGCD]
	Soit $A$ un anneau principal et $(a,b)\in A^2$. L'idéal $(a)+(b)$ est principal et, par conséquent, il existe $\delta\in A$ tel que $(a)+(b)=(\delta)$. L'élément $\delta$ n'est pas unique en général, mais il est unique à un coefficient multiplicatif inversible près. On dit que $\delta$ est \textbf{un pgcd} de $a$ et $b$ et on note $\delta=a\wedge b$.
\end{dfn}
\begin{prop}
	Soit $A$ un anneau principal et $(a,b)\in A^2$. On suppose que $(a,b)\ne(0,0)$. Dans ce cas, tout pgcd de $a$ et $b$ est non nul. Soit $\delta$ un tel pgcd. Toujours à un inversible près, $\delta$ est le "plus grand" (pour la relation $\mid$) diviseur commun à $a$ et $b$.
\end{prop}
\begin{proof}
	Soit $d$ un diviseur commun à $a$ et $b$. On a $(a)+(b)\subset(d)$ donc $(\delta)\subset(d)$ puis $d\mid \delta$.
\end{proof}
\begin{rmq}
	Comme on l'a vu précédemment, la relation $\mid$ est seulement une relation de préordre sur $A$. Il n'y a donc pas véritablement de plus grand élément à l'ensemble des diviseurs. On peut avoir une véritable unicité si on travaille non plus dans $A$ mais dans l'ensemble quotient $A/\mathcal{R}$ où $\mathcal{R}$ est la relation définie par $x\mathcal{R}y$ si, et seulement si, $x$ est associé avec $y$. Dans ce cas, la relation $\mid$ induit une vraie relation d'ordre sur $A/\mathcal{R}$ mais la manipulation des objets de l'ensemble quotient étant plus compliquée, on préfère dire simplement "à un inversible près".
\end{rmq}
\begin{exa}
	Dans le cas où $A=\Z$, pour avoir une véritable unicité, on fait le choix de prendre le pgcd dans $\N$ donc positif.
\end{exa}

\begin{exa}
	Dans le cas où $A=\K[X]$, pour avoir une véritable unicité, on fera le choix de prendre le polynôme unitaire (c'est-à-dire avec un coefficient dominant qui vaut $1$).
\end{exa}
\begin{dfn}[Ensemble des diviseurs]
	Soit $a$ un élément d'un anneau principal $A$. On note $\mathcal{D}(a)$ l'ensemble des diviseurs de $a$.
\end{dfn}
\begin{prop}
	Soit $a$ et $b$ deux éléments d'un anneau principal $A$. On a : $$\mathcal{D}(a)\cap\mathcal{D}(b)=\mathcal{D}(a\wedge b)$$
\end{prop}
\begin{proof}
	Comme $\delta=a\wedge b\in(a)+(b)$, il existe $(u,v)\in A^2$ tel que $au+bv=\delta$. Fixons-les. Ainsi, tout diviseur de $a$ et de $b$  est aussi un diviseur de $\delta$. De plus, $a\in(a)+(b)$ car $a=a+0$. Ainsi, $a\in(\delta)$ et $\delta\mid a$. De même, $\delta\mid b$. On a donc, par transitivité de la divisibilité, $D(\delta)\subset\mathcal{D}(a)\cap\mathcal{D}(b)$, ce qui achève la preuve.
\end{proof}
\begin{prop}
	Soit $A$ un anneau principal. A multiplication par un élément inversible près, on a :
	\begin{enumerate}
		\item $\forall (a,b)\in A^2,\ a\wedge b=b\wedge a$ (commutativité du PGCD)
		\item $\forall (a,b,c)\in A^3,\ (a\wedge b)\wedge c=a\wedge(b\wedge c)$ (associativité du PGCD)
		\item $\forall (a,b,x)\in A^3,\ (xa)\wedge(xb)=x(a\wedge b)$
	\end{enumerate}
\end{prop}
\begin{proof}
	\begin{enumerate}
		\item Clair par la proposition précédente et par commutativité de l'intersection ensembliste.
		\item Clair par la proposition précédente et par associativité de l'intersection ensembliste.
		\item On a \begin{align*}
			((xa)\wedge(xb))&=(xa)+(xb)\quad\text{(définition)}\\
			&=x[(a)+(b)]\quad\text{(double inclusion)}\\
			&=x(a\wedge b)\quad\text{(définition)}\\
			&=(x(a\wedge b))\quad\text{(double inclusion)}
		\end{align*}
	\end{enumerate}
\end{proof}
\begin{prop}
	Soit $a$ et $b$ deux éléments d'un anneau principal $A$. On a : $$\forall q\in A,\ \mathcal{D}(a)\cap\mathcal{D}(b)=\mathcal{D}(b)\mathcal{D}(a-bq)$$ Ainsi, à multiplication par un élément inversible près : $$a\wedge b=b\wedge(a-bq)$$
\end{prop}
\begin{proof}
	Exactement comme dans le cas de l'anneau $\Z$ !
\end{proof}
\begin{dfn}[Éléments premiers entre eux]
	Soit $a$ et $b$ deux éléments d'un anneau principal $A$. On dit que $a$ et $b$ sont \textbf{premiers entre eux} lorsque $a\wedge b=1$ (c'est-à-dire lorsque les seuls diviseurs communs de $a$ et $b$ sont les unités de l'anneau).
\end{dfn}
\begin{prop}[Théorème de Bézout]
	Soit $a$ et $b$ deux éléments d'un anneau principal $A$. Les assertions suivantes sont équivalentes : \begin{enumerate}
		\item $a$ et $b$ sont premiers entre eux
		\item Il existe $(u,v)\in A^2$ tel que $au+bv=1$
	\end{enumerate}
\end{prop}
\begin{proof}
	Si $a$ et $b$ sont premiers entre eux, alors $(a)+(b)=(1)=A$. Comme $1\in(1)$, il existe $(u,v)\in A^2$ tel que $au+bv=1$. Réciproquement, supposons qu'il existe $(u,v)\in A^2$ tel que $au+bv=1$ et fixons-les. On a $au+bv\in(a)+b(a)$ donc $1\in(a)+(b)$ donc $(1)\subset(a)+(b)$, c'est-à-dire $A\subset(a)+(b)$. Comme l'inclusion réciproque est toujours vraie, on a $(a)+(b)=(1)$ donc $a$ et $b$ sont premiers entre eux.
\end{proof}
\begin{cor}
	Soit $A$ un anneau principal et $(a,b,c)\in A^3$. On suppose que $a\wedge b=1$ et $a\wedge c=1$. Dans ce cas, $a\wedge bc=1$.
\end{cor}
\begin{proof}
	D'après le théorème de Bézout, on peut fixer $(u,v,u',v')\in A^4$ tel que $au+bv=1$ et $au'+cv'=1$. On a donc $a(auu'+ucv'+u'bv)+bc(vv')=1$. D'après le théorème de Bézout, cela signifie que $a\wedge bc=1$.
\end{proof}
\begin{rmq}
	On généralise sans problème par récurrence pour $n\geqslant 2$.
\end{rmq}
\begin{prop}[Lemme de Gauss]
	Soit $A$ un anneau principal et $(a,b,c)\in A^3$. Si $a\mid bc$ et $a\wedge b=1$ alors $a\mid c$.
\end{prop}
\begin{prop}
	D'après le théorème de Bézout, on fixer $(u,v)\in A^2$ tel que $au+bv=1$. On a donc $auc+bvc=c$ et $a$ divise $bc$ donc $a$ divise $c$ car divise le membre de gauche.
\end{prop}
\begin{cor}
	Soit $A$ un anneau principal et $(a,b,c)\in A^3$. Si $a\wedge b=1$, $a\mid c$ et $b\mid c$ alors $ab\mid c$. 
\end{cor}
\begin{proof}
	Comme $a\mid c$, il existe $q\in A$ tel que $c=aq$. Fixons-le. Par le lemme de Gauss, $b\mid q$ donc il existe $q'\in A$ tel que $q=bq'$. Donc $c=abq'$ puis $ab\mid c$.
\end{proof}
Dans un anneau principal quelconque, le calcul du PGCD n'est pas aisé. On s'intéresse ici à un cas particulier, celui où il existe une division euclidienne.
\begin{dfn}[Anneau euclidien, stathme]
	On considère un anneau $A$ commutatif et intègre. On dit que $A$ est un \textbf{anneau euclidien} lorsqu'il existe une application $$\begin{array}{ccccc}\nu&:&A\setminus\left\{0\right\}&\rightarrow&\N\end{array}$$vérifiant la propriété suivante : $$\forall (a,b)\in A\times(A\setminus\left\{0\right\}),\ \exists(q,r)\in A^2,\ a=bq+r\ \text{et}\ \left(r=0\ \text{ou}\ \nu(r)<\nu(b)\right)$$ On dit alors que $\nu$ est un \textbf{stathme} sur l'anneau euclidien $A$. On note $(A,\nu)$ un tel anneau euclidien, ou tout simplement $A$ quand le stathme est fixé.
\end{dfn}
\begin{rmq}[Étymologie]
	Le nom "stathme" provient du grec "stathmos" signifiant arrêt.
\end{rmq}
\begin{rmq}
	Attention, on ne demande pas l'unicité du couple $(q,r)$ !
\end{rmq}
\begin{exa}
	L'anneau $\Z$ est euclidien avec le stathme $x\mapsto \lvert x\rvert$.
\end{exa}
\begin{rmq}
	Avec les mêmes notations, on dit que $q$ est \textbf{un} quotient et $r$ est \textbf{un} reste dans la division euclidienne de $a$ par $b$.
\end{rmq}
\begin{prop}[Principalité des anneaux euclidiens]
	Tout anneau euclidien $(A,\nu)$ est principal.
\end{prop}
\begin{proof}
	Soit $(A,\nu)$ un anneau euclidien. Soit $I$ un idéal de $A$ non réduit à $\left\{0\right\}$. Considérons l'ensemble $$S:=\left\{\nu(a)\ | \ a\in I\setminus\left\{0\right\}\right\}$$ Par hypothèse, puisque $I\setminus\left\{0\right\}$ est non vide, $S$ est non vide. C'est aussi une partie de $\N$, donc elle admet un plus petit élément et il existe $a_0\in I\setminus\left\{0\right\}$ tel que $\nu(a_0)=\min(S)$. Montrons alors que $I=(a_0)$. Par absorbance de $I$, on a immédiatement $(a_0)\subset I$. Réciproquement, soit $a\in I$. Effectuons la division euclidienne de $a$ par $a_0$ : il existe $(q,r)\in A^2$ tel que $a=a_0q+r$ et $r=0$ ou $\nu(r)<\nu(a_0)$. Or, $r=a-a_0q\in I$ car $a_0q\in I$ d'après la première inclusion. Mais par minimalité de $\nu(a_0)$, on ne peut avoir $\nu(r)<\nu(a_0)$, donc nécessairement $r=0$ puis $a=a_0q\in(a_0)$. On a bien $I=(a_0)$ donc $I$ est un idéal principal. Comme l'idéal nul est évidemment principal, l'anneau $A$ est alors principal.
\end{proof}
\begin{met}[Algorithme d'Euclide]
	Soit $(A,\nu)$ un anneau euclidien. Puisque $A$ est principal, pour $q\in A$, on a, à multiplication par un élément inversible près, $a\wedge b=b\wedge(a-bq)$. On utilise cette propriété avec $q$ un quotient de $a$ par $b$ quand $b\ne0$ et il vient avec les notations habituelles : $$a\wedge b=b\wedge r$$ En itérant cette propriété, on construit une suite $(r_k)$ d'éléments de $A$ tels que $\nu(r_{k+1})<\nu(r_k)$ tant que $r_{k+1}\ne0$ en effectuant la division euclidienne de $r_{k-1}$ par $r_k$ après avoir posé $r_0=a$ et $r_1=b$. Comme il n'existe pas de suite strictement décroissante d'entiers naturels, cet algorithme se termine et le pgcd est le dernier reste non nul.
\end{met}
\subsection{PPCM dans un anneau principal}
\begin{dfn}[PPCM]
	Soit $A$ un anneau principal et $(a,b)\in A^2$. L'idéal $(a)\cap(b)$ est principal et, par conséquent, il existe $\mu\in A$ tel que $(a)\cap(b)=(\mu)$. L'élément $\delta$ n'est pas unique en général, mais il est unique à un coefficient multiplicatif inversible près. On dit que $\mu$ est \textbf{un ppcm} de $a$ et $b$ et on note $\mu=a\vee b$.
\end{dfn}
\begin{prop}
	Soit $A$ un anneau principal et $(a,b)\in A^2$. On suppose que $(a,b)\ne(0,0)$. Dans ce cas, tout ppcm de $a$ et $b$ est non nul. Soit $\mu$ un tel ppcm. Toujours à un inversible près, $\mu$ est le "plus petit" (pour la relation $\mid$) multiple commun à $a$ et $b$.
\end{prop}
\begin{proof}
	Soit $m$ un multiple commun à $a$ et $b$. On a $(m)\subset(a)\cap(b)$ donc $(m)\subset(\mu)$ puis $\mu\mid m$.
\end{proof}
\begin{prop}
Soit $A$ un anneau principal. A multiplication par un élément inversible près, on a :
\begin{enumerate}
	\item $\forall (a,b)\in A^2,\ a\vee b=b\vee a$ (commutativité du ppcm)
	\item $\forall (a,b,c)\in A^3,\ (a\vee b)\vee c=a\vee(b\vee c)$ (associativité du ppcm)
	\item $\forall (a,b,x)\in A^3,\ (xa)\vee(xb)=x(a\vee b)$
\end{enumerate}
\end{prop}
\begin{proof}
\begin{enumerate}
	\item Clair par commutativité de l'intersection ensembliste.
	\item Clair par associativité de l'intersection ensembliste.
	\item Le dernier point provient des égalités suivantes faciles à vérifier : $$(xa)\cap(xb)=x\left[(a)\cap(b)\right]=x(a\vee b)=(x(a\vee b))$$ On peut aussi se ramener au cas où $(a,b)\ne(0,0)$, poser $\mu=a\vee b$ et montrer que $x\mu$ est un multiple commun à $xa$ et $xb$ (facile) et que c'est le "plus petit" pour $\mid$. Pour cela, on considère $m\in A$ tel que $xa\mid m$ et $xb\mid m$. On a alors $x\mid m$ par transitivité et il existe $ù'\in A$ tel ue $m=xm'$. Ainsi, $a\mid m'$ et $b\mid m'$ donc $\mu\mid m'$ et il existe $u\in A$ tel que $m'=\mu u$. Finalement, $m=x\mu u$.
\end{enumerate}
\end{proof}
\begin{prop}
	Soit $A$ un anneau principal et $(a,b)\in A^2$. On a, à multiplication par un élément inversible près : $$(a\wedge b)(a\vee b)=ab$$
\end{prop}
\begin{proof}
	Comme dans le cas de l'anneau $\Z$, on pose $\delta=a\wedge b$ et il existe alors $(a',b')\in A^2$ tel que $a=\delta a'$, $b=\delta b'$ et $a'\wedge b'=1$. Il suffit alors de prouver que $a'\vee b'=a'b'$. On procède comme dans les démonstrations précédentes. On a clairement $a'\mid a'b'$ et $b'\mid a'b'$. On considère alors $m\in A$ tel que $a'\mid m$ et $b'\mid m$. Comme $a'\wedge b'=1$, on a $a'b'\mid m$. On peut en conclure que $a'b'$ est bien le plus petit multiple commun de $a'$ et $b'$, ce qui achève la preuve.
\end{proof}
\begin{rmq}
	En pratique, pour calculer un ppcm de deux éléments, on utilisera la relation $(a\wedge b)(a\vee b)=ab$ pour se ramener au calcul d'un pgcd de $a$ et $b$.
\end{rmq}
\section{Anneaux noethériens}
\begin{dfn}[Anneau noethérien]
	Soit $A$ un anneau commutatif. Les trois propriétés suivantes sont équivalentes : 
	\begin{enumerate}
		\item Tout idéal de $A$ est engendré par un nombre fini d'élément.
		\item Toute suite croissante d'idéaux de $A$ $(I_n)_{n\in\N}$ est stationnaire à partir d'un certain rang.
		\item Toute famille non vide d'idéaux admet un élément maximal pour l'inclusion.
	\end{enumerate}
	Dans ce cas, on dit que l'anneau $A$ est \textbf{noethérien}.
\end{dfn}
\begin{proof} 
	On démontre l'équivalence en faisant une circulante.
	\begin{itemize}
		\item $1\implies2$ : Supposons que tout idéal de $A$ est engendré par un nombre fini d'éléments. Soit $(I_n)_{n\in\N}$ une suite croissante d'idéaux. Alors $I:=\displaystyle\bigcup_{n\in\N}I_n$ est encore un idéal de $A$. En effet, $0\in I_0$ donc $0\in I$ et soit $(a,b)\in I^2$. Alors il existe $n\in\N$ tel que $(a,b)\in I_n^2$ par croissance de la suite. Puis $a+b\in I_n$ donc $a+b\in I$. Enfin, soit $a\in I$ et $b\in A$. Il existe $n\in\N$ tel que $a\in I_n$ puis $ab\in I_n$ par absorbance de $I_n$ donc $ab\in I$. Ensuite, $I$ est engendré par un nombre fini d'éléments $a_1,\ldots,a_m$. On peut alors trouver un $n_0\in\N$ tel que $a_1,\ldots,a_m$ soient des éléments de $I_{n_0}$. On a immédiatement $I_{n_0}\subset I$ par définition de $I$, et réciproquement, $I_{n_0}$ étant un idéal contenant tous les $a_i$, il est contenu par $I$ car $I$ est engendré par les $a_i$. Donc $I=I_{n_0}$. Soit $n\geqslant n_0$. On a $I_{n_0}\subset I_n$ par croissance de la suite. Or, $I_{n_0}=I$ donc par définition de $I$, $I_n\subset I_{n_0}$. On a donc $I_n=I_{n_0}$ et la suite $(I_n)_{n\in\N}$ est stationnaire.
		\item $2\implies3$ : On raisonne par contraposée. Supposons qu'il existe une famille $(I_s)_{s\in S}$ non vide d'idéaux de $A$ qui n'admette pas d'élément maximal pour l'inclusion. On va construire par récurrence une suite $(J_n)_{n\in\N}$ d'idéaux de $A$ strictement croissante. Pour l'initialisation, on choisit $s_0\in S$ car $S$ est non vide, et on pose $J_0:=I_{s_0}$. Désormais, supposons avoir construit $J_n$ pour $n\in\N$, tel que $\forall k\in\llbracket0,n\rrbracket$, $\exists s\in S,\ J_k=I_s$ et telle que $\forall k\in\llbracket0,n-1\rrbracket,\ J_k\subsetneq J_{k+1}$. Puisque $(I_s)_{s\in S}$ ne possède pas d'élément maximal $\exists s\in S,\ J_n\subsetneq I_s$. Fixons ce $s$ et posons $J_{n+1}:=I_s$. Ainsi, on a construit une suite $(J_n)_{n\in\N}$ d'idéaux de $A$ strictement croissante, donc il existe une suite d'idéaux croissante et non stationnaire.
		\item $3\implies1$ : Supposons que tout famille non vide d'idéaux admette un élément maximal pour l'inclusion. Soit $I$ un idéal de $A$ et considérons l'ensemble des idéaux de $I$ engendrés par un nombre fini d'éléments et qui sont inclus dans $I$. Cet ensemble est non vide car il contient l'idéal nul. Par hypothèse, il possède alors un élément maximal $J$. On a alors $J\subset I$. Soit $x\in I$. Alors $J+xA$ est un idéal, comme somme de deux idéaux, et il engendré par un nombre fini d'éléments puisqu'il est engendré par le nombre fini d'éléments qui engendrent $J$ et $x$. Or, $xA\subset I$ donc par stabilité de $I$ par $+$, $J+xA\subset I$. Or, $J\subset J+xA$, donc, par maximalité de $J$, $J=J+xA$. Or, $x=0+x\times 1\in J+xA$ donc $x\in J$. Ainsi, $J=I$, donc $I$ est bien engendré par un nombre fini d'éléments.
	\end{itemize}
	Ainsi, la preuve est achevée.
\end{proof}
\begin{prop}
	Tout anneau principal est noethérien.
\end{prop}
\begin{proof}
	Tous les idéaux d'un anneau principal sont engendrés par un seul élément, donc en particulier par un nombre fini d'éléments. Tout anneau principal vérifie donc la condition $1$, et est alors noethérien.
\end{proof}
\section{Anneaux factoriels}
\subsection{Définition et première propriétés}
\begin{dfn}[Élément irréductible]
	Soit $A$ un anneau intègre. Un élément $a\in A$ est dit \textbf{irréductible} s'il est non nul, non inversible et si ses seuls diviseurs sont les diviseurs triviaux $u$ et $ua$ pour $u\in A^\times$. Ce dernier point revient à dire qu'il vérifie la propriété suivante : $$\forall (u,v)\in A^2,\ (a=uv)\implies(u\in A^\times\ \text{ou}\ v\in A^\times)$$
\end{dfn}
\begin{exa}
	Les éléments irréductibles et positifs de $\Z$ sont les nombres premiers.
\end{exa}
\begin{exa}
	Si $\K$ est un corps, les éléments inversibles de l'anneau $\K[X]$ sont les polynômes constants non nuls. Par suite, un polynôme est irréductible si et seulement si ses seuls diviseurs sont les constantes non nulles et les polynômes qui lui sont associés.
\end{exa}
\begin{dfn}[Anneau factoriel]
	Un anneau $A$ est dit \textbf{factoriel} lorsqu'il est intègre et et que tout élément non nul admet une décomposition en produit de facteurs irréductibles et si cette décomposition est unique à l'ordre près et au produit par un inversible près. Formellement, cela signifie que :
	\begin{enumerate}
		\item Pour tout $a\in A\setminus\left\{0\right\}$, il existe $u\in A^\times$, des éléments irréductibles $p_1,\ldots,p_k$ de $A$ et des entiers naturels non nuls $\alpha_1,\ldots,\alpha_k$ tels que : $$a=up_1^{\alpha_1}\ldots p_k^{\alpha_k}=u\prod_{i=1}^{k}p_i^{\alpha_i}$$
		\item Il y a unicité de la décomposition en ce sens : si $$a=up_1^{\alpha_1}\ldots p_k^{\alpha_k}=vq_1^{\beta_1}\ldots q_m^{\alpha_m}$$ avec $\alpha_1,\ldots,\alpha_k,\beta_1,\ldots,\beta_m\in\N\st$ et où les $p_1,\ldots,p_k$ d'une part, et les $q_1,\ldots,q_m$ d'autre part, ne sont pas associés, alors $k=m$ et il existe une permutation $\sigma$ de $\llbracket1,k$ et des éléments inversibles $u_i$ tels que $q_{\sigma(i)}=u_ip_i$ et $\beta_{\sigma(i)}=\alpha_i$ pour tout $i\in\llbracket1,k\rrbracket$.
	\end{enumerate}
	On notera $(1)$ la condition d'existence et $(2)$ la condition d'unicité.
\end{dfn}
\begin{prop}
	Sous réserve que la condition $(1)$ soit vérifiée, alors la condition $(2)$ équivaut à cette nouvelle condition, notée $(2')$ : tout élément irréductible $a$ de $A$ vérifie $$\forall(b,c)\in A^2,\ a\mid bc\implies a\mid b\ \ \text{ou}\ \ a\mid c$$
\end{prop}
\begin{proof}
	Supposons que $A$ est intègre et vérifie la condition $(1)$.
	\begin{itemize}
		\item $(2)\implies(2')$ : Supposons la condition $(2)$ vérifiée par $A$. Si $a$ est irréductible et divise $bc$, on écrit $bc=ad$, puis les décompositions de $b$, $c$ et $d$ en produits de facteurs irréductibles. L'unicité de la décomposition montre alors que $a$ figure dans la décomposition de $bc$, donc obligatoirement dans celle de $b$ ou dans celle de $c$.
		\item $(2')\implies(2)$ : Supposons la condition $(2')$ vérifiée par $A$. On raisonne par récurrence sur $s\in\N\st$ où $s$ sera la somme des exposants figurant dans une des décompositions de $a$. L'hypothèse de récurrence $H(s)$ est donc la suivante : \\ \\ "Si $a\in A\setminus\left\{0\right\}$ vérifie $a=up_1^{\alpha_1}\ldots p_k^{\alpha_k}=vq_1^{\beta_1}\ldots q_m^{\alpha_m}$ avec $\alpha_1,\ldots,\alpha_k,\beta_1,\ldots,\beta_m\in\N\st$ et où les $p_1,\ldots,p_k$ d'une part, et les $q_1,\ldots,q_m$ d'autre part, ne sont pas associés, et $\alpha_1+\ldots+\alpha_k=s$, alors $k=m$ et il existe une permutation $\sigma$ de $\llbracket1,k$ et des éléments inversibles $u_i$ tels que $q_{\sigma(i)}=u_ip_i$ et $\beta_{\sigma(i)}=\alpha_i$ pour tout $i\in\llbracket1,k\rrbracket$." \\ \\ Initialisation : si $a=up_1=vq_1^{\beta_1}\ldots q_m^{\alpha_m}$ alors $q_i$ divise $p_1$ pour tout $i$ en utilisant $(2')$ et le fait que $q_i$ ne divise pas $u$ puisque $u$ est inversible (donc $u\mid q_i$) et que $q_i$ n'est pas inversible. Cela impose à $q_i$ d'âtre associé à $p_1$. Comme les $q_i$ ne sont pas associés entre eux, on aura $m=1$, $q_1$ associé à $p_1$ et $\beta_1=1$. Ainsi, $H(1)$ est vraie.\\ \\
		Hérédité : Supposons que $H(s)$ soit vraie pour un certain $s\in\N\st$ et montrons la propriété pour $(s+1)$. On considérer alors les décompositions $$a=up_1^{\alpha_1}\ldots p_k^{\alpha_k}=vq_1^{\beta_1}\ldots q_m^{\alpha_m}$$ avec $\alpha_1+\ldots+\alpha_k=s+1$. On a $q_m$ irréductible et $q_m$ divise le produit des $p_i$ et de $u$ mais $q_m$ ne divise pas $u$ car $u$ est inversible (donc $u\mid q_m$) et $q_m$ n'est pas inversible. Ainsi, $q_m$ divise l'un des $p_i$ d'après $(2')$. Quitte à permuter les $p_i$, on peut supposer que $q_m$ divise $p_k$. Comme $p_k$ est irréductible, il existe $w\in A^\times$ tel que $q_m=wp_k$ et l'on obtient, après simplification car l'anneau $A$ est intègre : $$a=up_1^{\alpha_1}\ldots p_k^{\alpha_k-1}=vq_1^{\beta_1}\ldots q_m^{\alpha_m-1}$$ Il suffit alors d'appliquer l'hypothèse de récurrence à cette décomposition en distinguant deux cas pour s'assurer que les exposants qui interviennent son vraiment tous strictement positifs : si $\alpha_k=1$, alors $\beta_m=1$ autrement $q_m$ (et donc $p_k$) diviserait l'un des $p_i$ avec $i\ne k$, ce qui est absurde. Si $\alpha_k>1$, alors $\beta_k>1$ autrement $p_k$ diviserait l'un des $q_i$ avec $i\ne m$. La propriété $H(s+1)$ est alors démontrée. \\ \\ La récurrence est alors achevée.
	\end{itemize}
	Ainsi, la preuve est achevée.
\end{proof}
\begin{prop}
	Soit $A$ un anneau factoriel. $a\in A$ est irréductible si, et seulement si, il est non nul, non inversible et vérifie la propriété $a\mid bc\implies a\mid b\ \ \text{ou}\ \ a\mid c$.
\end{prop}
\begin{prop}
	Le sens direct a déjà été prouvé lors de la démonstration précédente. Réciproquement, si $a$ vérifie cette propriété, et si $a=uv$, alors $a\mid uv$ donc $a\mid u$ ou $a\mid v$. Si par exemple $u=aa'$ alors par intégrité $1=a'v$ et $v$ sera inversible. Idem si $v=aa'$.
\end{prop}
\subsection{Anneaux atomiques}
On va donner ici la définition d'un type d'anneau intègre qui vérifie la condition $(1)$ et montrer une propriété sur ce type d'anneau qui servira pour la partie suivante.
\begin{dfn}[Anneau atomique]
	Un anneau $A$ est dit \textbf{atomique} lorsqu'il est intègre et qu'il vérifie la condition $(1)$, c'est-à-dire l'existence d'une décomposition en produit d'irréductibles pour tout élément non nul.
\end{dfn}
\begin{exa}
	Un anneau factoriel est nécessairement atomique.
\end{exa}
\begin{prop}[Atomicité des anneaux noethériens intègres]
	Tout anneau intègre et noethérien est atomique.
\end{prop}
\begin{proof}
	Soit $A$ un anneau intègre et noethérien. On raisonne par l'absurde en supposant l'existence d'un élément non nul $a$ de $A$ sans décomposition en produit de facteurs irréductibles. Alors $a$ n'est pas irréductibles donc il existe $a_1$ et $b_1$ deux éléments de $A$ non inversibles tels que $a=a_1b_1$.  Au moins l'un des deux facteurs, disons $a_1$, ne se décompose pas en produit de facteurs irréductibles, sans quoi en réunissant les deux décompositions $a$ se décomposerait en produit de facteurs irréductibles. Il existe alors $a_2$ et $b_2$ des éléments de $A$ non inversibles tels que $a_1=a_2b_2$. Par exemple, $a_2$ ne se décompose pas en produit de facteurs irréductibles, et l'on continue de la même manière par récurrence. On obtient alors une suite infinie strictement croissante d'idéaux $$(a)\subsetneq(a_1)\subsetneq(a_2)\subsetneq\ldots\subsetneq(a_n)\subsetneq\ldots$$ en contradiction avec le caractère noethérien de $A$. On notera  que la suite d'idéaux est bien strictement croissante puisque $(a_i)=(a_{i+1})$ entraîne l'existence de $q$ et $q'$ des éléments de $A$ tels que $a_i=qa_{i+1}$ et $a_{i+1}=q'a_i$ donc par intégrité $1=qq'$ et $q'$ est inversible, puis, en remplaçant dans $a_i=a_{i+1}b_{i+1}$, on trouve par intégrité $1=ub_{i+1}$ ce qui est en contradiction avec le fait que $b_{i+1}\notin A^\times$. Ainsi, $A$ est bien atomique. 
\end{proof}
\subsection{Factorialité des anneaux principaux}
On arrive finalement à la proposition fondamentale : 
\begin{prop}[Factorialité des anneaux principaux]
	Tout anneau principal est factoriel.
\end{prop}
\begin{proof}
	Soit $A$ un anneau principal et montrons que $A$ vérifie les conditions $(1)$ et $(2')$.
	\begin{itemize}
		\item $(1)$ : $A$ est principal donc $A$ est noethérien. Or, $A$ est principal donc $A$ est intègre. Comme $A$ est intègre et noethérien, il est atomique d'après le paragraphe qui précède et alors il vérifie la condition $(1)$.
		\item $(2')$ : Soit $a$ un élément irréductible de l'anneau $A$ et $(b,c)\in A^2$ tel que $a\mid bc$. Si $a\mid b$,a lors c'est bon. Si $a$ ne divise pas $b$, alors puisque $a$ est irréductible, $a$ est premier avec $b$ et d'après le lemme de Gauss il divise $c$ et la condition $(2')$ est vérifiée.
		\item Comme $(1)$ est vérifiée, alors $(2')$ est équivalente à $(2)$, et puisque $(2')$ est vérifiée, alors $(2)$ est vérifiée. Par conséquent, $A$ est bien factoriel.
	\end{itemize}
	Ainsi, la preuve est achevée. 
\end{proof}
\begin{rmq}
	Notons l'importance capitale de ce résultat. Il prouve d'un seul coup les théorèmes de décomposition en facteurs premiers dans $\Z$ et de décomposition en facteurs irréductibles dans $\K[X]$. Plus généralement, puisqu'on a vu qu'il suffisait d'être un anneau euclidien pour être un anneau principal, donc le simple fait de démontrer l'existence (et même pas l'unicité !) d'une division euclidienne sur un anneau garantit l'existence et l'unicité d'une décomposition en produit de facteurs irréductibles. Nous aurions donc pu nous arrêter aux théorèmes de division euclidienne dans $\Z$ et $\K[X]$ (lorsque $\K$ est un corps) plutôt que de démontrer tout un tas d'autres propriétés pour garantir l'existence et l'unicité de la décomposition en produit d'irréductibles.
\end{rmq}
\subsection{PGCD et PPCM dans un anneau factoriel}
Dans tout le paragraphe, $A$ est un anneau factoriel fixé. \\ \par  Soit $a$ et $b$ deux éléments de $A$. Il existe $p_1,\ldots,p_k$ des éléments irréductibles de $A$ deux à deux non associés et des éléments inversibles $u$ et $v$ tels que $$a=up_1^{\alpha_1}\ldots p_k^{\alpha_k} \ \ \text{et}\ \ b=vp_1^{\beta_1}\ldots p_k^{\beta_k}$$ avec $\alpha_1,\ldots,\alpha_k,\beta_1,\ldots\beta_k\in\N$. En effet, on regroupe les irréductibles associés et on autorise les exposants nuls de façon à avoir les mêmes irréductibles pour $a$ et $b$. \\ \par Dans tout le paragraphe, on conserve ces décompositions.
\begin{prop}[Critère de divisibilité]
	Avec les mêmes notations : $$a\mid b\iff \forall i\in\llbracket1,k\rrbracket,\ \alpha_i\leqslant\beta_i$$
\end{prop}
\begin{proof}
	Pour le sens direct, il suffit de remarquer que la décomposition de $b$ "contient" alors celle de $a$, ce qui force les inégalités. Réciproquement, si on a les inégalités, il suffit de "compléter" les exposants dans la décomposition de $a$ pour obtenir un élément $c$ tel que $ac=b$.
\end{proof}
\begin{dfn}[Expression du PGCD et du PPCM]
	Avec les mêmes notation, et toujours à un inversible près, on a : $$a\wedge b=\prod_{i=1}^{k}p_i^{\gamma_i}$$ et $$a\vee b=\prod_{i=1}^{k}p_i^{\lambda_i}$$ avec pour tout $i\in\llbracket1,k\rrbracket$, $\gamma_i:=\min(\alpha_i,\beta_i)$ et $\lambda_i:=\max(\alpha_i,\beta_i)$.
\end{dfn}
\begin{proof}
	Il suffit d'utiliser le critère de divisibilité et de montrer que les nombres donnés ici sont des multiples communs ou diviseurs communs de $a$ et $b$ et le critère de divisibilité montre qu'ils sont les "plus petit" et "plus grand" pour la divisibilité, toujours à un inversible près.
\end{proof}
\begin{rmq}
	Là encore, cette proposition effectue d'un seul coup ce que nous avons longuement prouvé dans $\Z$ et dans $\K[X]$. Remarquons que dans ces anneaux, on impose toutefois une condition supplémentaire pour avoir l'unicité du PGCD et du PPCM : dans $\Z$ on le prend positif car les unités de $\Z$ sont $-1$ et $1$, et dans $\K[X]$ on le prend unitaire car les unités de $\K[X]$ sont les polynômes constants non nuls.
\end{rmq}
\section{Corps des fractions (cas commutatif)}
Dans le deuxième paragraphe qui traite des anneaux quotients, nous avons d'une certaine façon "forcé" des éléments d'un anneau à être nuls. Nous voulons maintenant effectuer une opération "opposée" : rendre inversibles les éléments d'un anneau convenable. \\ \par Dans cette partie, l'anneau $A$ est commutatif et intègre. 
\begin{prop}
	On considère la relation $\sim$ définie sur $A\times(A\setminus\left\{0\right\})$ par : $$\forall(a,b,a',b')\in\left(A\times(A\setminus\left\{0\right\})\right)^2,\ (a,b)\sim(a',b') \iff ab'=a'b$$ La relation binaire $\sim$ est une relation d'équivalence.
\end{prop}
\begin{proof}
	Les vérifications sont immédiates. 
\end{proof}
\begin{dfn}[Notation $frac(A)$]
	On note alors $frac(A)=(A\times(A\setminus\left\{0\right\}))/\sim$.
\end{dfn}
\begin{dfn}[Addition et multiplication sur $frac(A)$]
	On définit une multiplication et une addition sur $frac(A)$ par : $$\begin{array}{ccccc}\tilde{+}&:&frac(A)&\rightarrow&frac(A)\\&&\left(\overline{(a,b)},\overline{(a',b')}\right)&\mapsto&\overline{(ab'+a'b,bb')}\end{array}$$ et $$\begin{array}{ccccc}\tilde{\times}&:&frac(A)&\rightarrow&frac(A)\\&&\left(\overline{(a,b)},\overline{(a',b')}\right)&\mapsto&\overline{(aa',bb')}\end{array}$$ Ainsi, ce que l'on veut imiter, c'est tout simplement le calcul de fractions que l'on apprend au collège.
\end{dfn}
\begin{prop}
	Les opérations ci-dessus sont bien définies sur $frac(A)$.
\end{prop}
\begin{proof}
	Soient $(a_1,b_1)$, $(a_1',b_1')$, $(a_2,b_2)$ et $(a_2',b_2')$ des éléments de $A\times(A\setminus\left\{0\right\})$ tels que $(a_1,b_1)\sim(a_1',b_1')$ et $(a_2,b_2)\sim(a_2',b_2')$. Dans ce cas, on a $$(a_1b_2+a_2b_1,b_1b_2)\sim(a_1'b_2'+a_2'b_1',b_1'b_2')$$ car \begin{align*}
		(a_1b_2+a_2b_1)b_1'b_2'&=a_1b_1'b_2b_2'+a_2b_2'b_1b_1' \\
		&= a_1'b_1b_2b_2'+a_2'b_2b_1b_1' \\
		&= (a_1'b_2'+a_2'b_1')b_1b_2
	\end{align*} Ceci montre que l'application $\tilde{+}$ est bien définie car deux classes d'équivalences ont la même images, et aussi car par intégrité de $A$ $bb'$ est non nul lorsque $b$ et $b'$ sont non nuls. De plus, on a : $$(a_1a_2,b_1b_2)\sim(a_1'a_2',b_1'b_2')$$ car \begin{align*}
	a_1a_2b_1'b_2'&=(a_1b_1')(a_2b_2')\\
	&=(a_1'b_1)(a_2'b_2)\\
	&=a_1'a_2'b_1b_2
	\end{align*} ce qui prouve que $\tilde{\times}$ est bien définie.
\end{proof}
\begin{prop}
	Muni des deux opérations définies ci-dessus, $frac(A)$ est un corps commutatif. L'élément neutre puir $\tilde{+}$ est $\overline{(0,1)}$ et l'élément neutre pour $\tilde{\times}$ est $\overline{(1,1)}$.
\end{prop}
\begin{proof}
	Vérifions les différents axiomes.
	\begin{itemize}
		\item L'addition est associative. En effet, on a 
		\begin{align*}
			\left(\overline{(a_1,b_1)}\tilde{+}\overline{(a_2,b_2)}\right)\tilde{+}\overline{(a_3,b_3)}&=\overline{(a_1b_2+a_2b_1,b_1b_1)}\tilde{+}\overline{(a_3,b_3)} \\
			&= \overline{\left((a_1b_2+a_2b_1)b_3+a_3(b_1b_2),(b_1b_2)b_3)\right)}\\
			&= \overline{\left(a_1b_2b_3+a_2b_1b_3+a_3b_1b_2,b_1b_2b_3)\right)} \\
			&= \overline{(a_1,b_1)}\tilde{+} \overline{(a_2b_3+a_3b_2,b_2b_3)}\\
			&=\overline{(a_1,b_1)}\tilde{+}\left(\overline{(a_2,b_2)}\tilde{+}\overline{(a_3,b_3)}\right)
		\end{align*}
		\item L'addition est clairement commutative par commutativité de $+$ et $\times$ sur $A$.
		\item L'élément neutre pour $\tilde{+}$ est $\overline{(0,1)}$ et l'opposé de $\overline{(a,b)}$ est $\overline{(-a,b)}$ (vérifications immédiates).
		\item La multiplication est clairement associative et commutative par associativité et commutativité de $\times$ sur $A$.
		\item L'élément neutre pour $\tilde{\times}$ est $\overline{(1,1)}$, cela s'hérite immédiatement du fait que $1$ est le neutre pour $\times$ dans $A$.
		\item Le produit est distributif sur la somme. En effet, on a \begin{align*}
			\left(\overline{(a_1,b_1)}\tilde{\times}\overline{(a_3,b_3)}\right)\tilde{+}\left(\overline{(a_2,b_2)}\tilde{\times}\overline{(a_3,b_3)}\right)&=\overline{(a_1a_3,b_1b_3)}\tilde{+}\overline{(a_2a_3,b_2b_3)}\\
			&=\overline{((a_1b_2+a_2b_1)a_3b_3,b_1b_2b_3^2)}\\
			&=\overline{((a_1b_2+a_2b_1)a_3,b_1b_2b_3)}\\
			&=\overline{(a_1b_2+a_2b_1,b_1b_2)}\tilde{\times}\overline{(a_3,b_3)} \\
			&=\left(\overline{(a_1,b_1)}\tilde{+}\overline{a_2,b_2}\right)\tilde{\times}\overline{(a_3,b_3)}
		\end{align*}
		\item Enfin, pour $\overline{(a,b)}\ne0$, on a $a\ne0$. En effet, les éléments de classe de $0$ sont tous les éléments de $frac(a)$ de la forme $(0,c)$, avec $c\ne0$. On peut donc conclure que $\overline{(a,b)}$ est inversible dans $frac(A)$ puisque son inverse n'est autre que $\overline{(b,a)}$ qui a un sens puisque $a\ne0$.
	\end{itemize}
	Ainsi, la preuve est achevée.
\end{proof}
\begin{prop}[Inclusion abusive $A\subset frac(A)$]
	L'application $$\begin{array}{ccccc}i&:&A&\longrightarrow&frac(A)\\&&a&\longmapsto&\overline{(a,1)}\end{array}$$ est un morphisme injectif d'anneaux. Ainsi, l'anneau $A$ est isomorphe à l'anneau $i(A)$ qui est un sous-anneau de $frac(A)$. Grâce à cet isomorphisme, tout calcul dans $A$ peut se voir comme un calcul dans $A$ peut se voir comme un calcul dans $i(A)$ et c'est pourquoi on considère que $A$ est inclus dans $frac(A)$ alors que, rigoureusement, ce n'est pas vraiment une inclusion mais seulement un isomorphisme entre $A$ et un sous-anneau de $frac(A)$.
\end{prop}
\begin{proof}
	On vérifie facilement que l'application $i$ est un morphisme d'anneaux puisque les applications $\tilde{+}$ et $\tilde{\times}$ sont bien définies et que le neutre pour $\tilde{\times}$ est $\overline{(1,1)}$. L'injectivité provient directement de l'étude du noyau et de l'intégrité de $A$.
\end{proof}
\begin{rmq}
	Cela revient aussi à considérer qu'un anneau n'a d'importance ou n'existe vraiment qu'à un isomorphisme près.
\end{rmq}
\begin{dfn}[Notations finales]
	Grâce à cette inclusion abusive, on note désormais $+$ et $\times$ pour $\tilde{+}$ et $\tilde{\times}$ puisque ces dernières opérations prolongent $+$ et $\times$. Souvent, quand il n'y a pas d'ambiguïté, on notera plutôt $\displaystyle\frac{a}{b}$ l'élément $\overline{(a,b)}$.
\end{dfn}
\begin{rmq}
	Le corps des fractions d'un corps est lui-même.
\end{rmq}
\begin{exa}
	Le corps des fractions de l'anneau intègre $\Z$ est le corps $\Q$.
\end{exa}
\begin{exa}
	Lorsque $\K$ est un corps, l'anneau des polynômes à coefficients dans $\K$ $\K[X]$ est intègre et son corps des fractions, noté $\K(X)$ est le corps des fractions rationnelles à coefficients dans $\K$.
\end{exa}
\begin{rmq}[Généralisation : localisation d'un anneau]
	Il existe une construction encore plus générale à l'aide d'une relation d'équivalence similaire qui consiste à rendre toute partie $S$ de $A$ stable par multiplication et contenant $1$ inversible de force. On parle de \textbf{localisation} de la partie $S$ de l'anneau $A$. Tout fonction de même, les vérifications sont plus pénibles car la relation d'équivalence est alors plus complexes. En revanche, le morphisme $i$ n'est plus nécessairement injectif si l'anneau $A$ n'est pas intègre, donc l'inclusion abusive n'a plus lieu d'être. Lorsqu'on localise $A\setminus\left\{0\right\}$ pour un anneau intègre (cette partie étant alors stable par multiplication et contenant $1$), on obtient le corps de fractions de $A$, ce qui prouve bien que la méthode de localisation est encore plus générale.
\end{rmq}
\section{Exercices}
\textit{Note : Je n'ai pas pris le temps de tenter de résoudre un seul de ces exercices, donc je ne connais pas les méthodes pour les résoudre. Si toutefois vous trouvez une solution, je serai ravi que vous me la partagiez.}
\begin{exe}
	Soit $n\geqslant2$ un entier. Déterminer les éléments nilpotents et les éléments inversibles de $\Z/n\Z$.
\end{exe}
\begin{exe}
	Soit $n$ et $m$ des entiers non nuls. Montrer que l'application canonique de $\Z/nm\Z$ dans $\Z/n\Z$ est un morphisme d'anneaux. Montrer qu'il induit une surjection de $\left(\Z/nm\Z\right)^\times$ sur $\left(Z/n\Z\right)^\times$.
\end{exe}
\begin{exe}
	Soit $n$ et $m$ des entiers non nuls. A quelle condition peut-on trouver un morphisme d'anneaux de $\Z/n\Z$ dans $\Z/m\Z$ ?
\end{exe}
\begin{exe}
	Exhiber un morphisme d'anneaux $f:A\rightarrow B$ qui soit surjectif mais tel que le morphisme de groupes de $A^\times$ dans $B^\times$ déduit par restriction de $f$ ne le soit pas.
\end{exe}
\begin{exe}
	Soit $\K$ un corps et $A$ un anneau non trivial. Montrer que tout morphisme d'anneaux de $\K$ dans $A$ est injectif.
\end{exe}
\begin{exe}
	Quel est le radical de l'idéal $(12)$ dans $\Z$ ? (Oui, c'est un poly de Taupins au départ, je le rappelle...)
\end{exe}
\begin{exe}[Généralisation du précédent]
	Soit $n$ un entier strictement positif qu'on décompose en facteurs premiers $n=\displaystyle\prod_{i=1}^{k}p_i^{\alpha_i}$ avec les $\alpha_i$ des entiers strictement positifs et les $p_i$ des nombre premiers deux à deux distincts. On pose $m:=\displaystyle\prod_{i=1}^{k}p_i$. Montrer que $\sqrt{n\Z}=m\Z$.
\end{exe}
\begin{exe}
	Soit $A$ un anneau commutatif et $(a,b)\in A^2$. Montrer que si $a$ et $b$ sont associés, alors les idéaux $(a)$ et $(b)$ sont égaux. Réciproquement, si $A$ est intègre et si $(a)=(b)$, montrer que $a$ et $b$ sont associés. Et si $A$ n'est pas intègre ?
\end{exe}
\begin{exe}
	Montrer qu'un anneau intègre ne possédant qu'un nombre fini d'idéaux est un corps. On pourra montrer que tout élément non nul $x$ est inversible en introduisant les idéaux $x^nA$ pour $n\in\N\st$ et utiliser le principe des tiroirs.
\end{exe}
\begin{exe}
	Soit $I$ et $J$ deux idéaux d'un anneau commutatif. Montrer que si $I$ et $J$ sont comaximaux (\textit{ie} si $I+J=A$), alors pour tout $n\in\N\st$, $I^n$ et $J^n$ (pas au sens du produit cartésien mais au sens de l'idéal produit) sont comaximaux (\textit{ie} $I^n+J^n=A$).
\end{exe}
\begin{exe}
	Soit $A$ un anneau non nécessairement commutatif.
	\begin{enumerate}
		\item Montrer par un contre-exemple que l'ensemble des éléments nilpotents de $A$ ne forme par un sous-groupe abélien en général. On pourra choisir $A=\mathcal{M}_2(\C)$.
		\item Soit $N$ l'ensemble des éléments $a\in A$ tel que $ax$ soit nilpotent pour tout $x\in A$. Montrer que $N$ est un idéal de $A$ dont tout élément est nilpotent.
		\item Soit $I$ un idéal de $A$ dont tout élément est nilpotent. Montrer que $I\subset N$.
	\end{enumerate}
\end{exe}
\begin{exe}
	Soit $A$ un anneau non nécessairement commutatif et soit $I$ l'idéal engendré par les $xy-yx$ pour $(x,y)\in A^2$.
	\begin{enumerate}
		\item Montrer que l'anneau $A/I$ est commutatif.
		\item Soit $J$ un idéal de $A$ tel que $A/J$ soit commutatif. Montrer que $I\subset J$.
	\end{enumerate}
\end{exe}
\begin{exe}
	 On s'intéresse au polynôme $X^2+1$.
	\begin{enumerate}
		\item Soit $\K$ un corps commutatif et $P\in\K[X]$ un polynôme à coefficients dans $\K$. Montrer que l'anneau $\K[X]/(P)$ est un corps si, et seulement si, $P$ est irréductible dans $\K[X]$.
		\item Montrer que le polynôme $X^2+1$ est irréductible dans l'anneau $\Z[X]$. L'anneau $A:=\Z[X]/(X^2+1)$ est-il un corps ? On pourra définir un isomorphisme de $A$ sur l'anneau $\Z[i]$.
		\item Soit $p$ un nombre premier. Montrer que le polynôme $X^2+1$ est irréductible dans $(\Z/p\Z)[X]$ si, et seulement si, $p\equiv3\ [4]$.
	\end{enumerate}
\end{exe}
\begin{exe}
	Soit $A$ un anneau et $I$ un idéal de $A$. On note $I[X]$ l'ensemble des polynômes $P\in A[X]$ dont tous les coefficients appartiennent à $I$.
	\begin{enumerate}
		\item Montrer que si $I$ est un idéal bilatère de $A$, alors $I[X]$ est un idéal bilatère de $A[X]$.
		\item Construir un isomorphisme de l'anneau $A[X]/I[X]$ sur l'anneau $(A/I)[X]$.
	\end{enumerate}
\end{exe}
\begin{exe}
	Soit $A$ un anneau et $I$ un idéal de $A$. On note $\mathcal{M}_n(I)$ l'ensemble des matrices de $\mathcal{M}_n(A)$ dont tous les coefficients appartiennent à $I$.
	\begin{enumerate}
		\item Montrer que $\mathcal{M}_n(I)$ est un idéal de $\mathcal{M}_n(A)$ et construire un isomorphisme d'anneaux de $\mathcal{M}_n(A)/\mathcal{M}_n(I)$ sur $\mathcal{M}_n(A/I)$.
		\item Inversement, montrer que tout idéal de $\mathcal{M}_n(A)$ est de la forme $\mathcal{M}_n(I)$, pour $I$ un idéal bilatère de $A$.
	\end{enumerate}
\end{exe}

\newpage
\chapter{Réels et suites}
La secte des Pythagoriciens possédait un joyau dont elle ne faisait don à personne d'extérieur : la démonstration du théorème de Pythagore. Ce théorème impliquait notamment que la diagonale d'un carré de côté $1$ valait $\sqrt{2}$. Lorsque les Pythagoriciens, pour qui "tout était nombre" (au sens de "tout est quotient de nombres entiers"), découvrirent que $\sqrt{2}$ était irrationnel, ils furent catastrophés : il était possible de faire coexister géométriquement des grandeurs "incommensurables" numériquement ! Finalement, cette crise des fondements ne fut résolue de manière satisfaisante que beaucoup plus tard, avec les premières constructions rigoureuses de $\R$ proposées au XIX$^e$ siècle. \par C'est une telle construction que nous proposons à la fin de ce chapitre, sans doute une des plus classiques. Comme elle n'est pas exigible au programme, on pourra admettre le résultat fondamental qui en découle, à savoir la propriété de la borne supérieure. C'est sur cette propriété fondatrice que s'appuie l'essentiel de l'analyse réelle.
\section{Propriétés fondatrices de $\R$}
Exposons sans démonstration les principales propriétés qui proviennent directement de la construction de $\R$.
\begin{thm}[Existence de $\R$ - admis]
	On peut construire $(\R,+,\times)$ un sur-corps de $\Q$, totalement ordonné par $\leqslant$. La relation $\leqslant$ est compatible avec l'addition et la multiplication au même sens que dans $\Z$ et $\Q$ : $$\forall(x,y,z)\in\R^3,\ x\leqslant y \implies x+z\leqslant y+z$$ $$\forall(x,y,z)\in\R^3,\ \left \{\begin{array}{ll}x\leqslant y \\ z\geqslant 0\end{array}\right. \implies xz\leqslant yz$$
\end{thm}
\begin{dfn}
	Soit $(a,b)\in\R^2$. A partir de $\leqslant$, on définit de la manière habituelle les ensembles $[a,b]$, $[a,b[$, $]a,b]$ et $]a,b[$. Si $a>b$, ces ensembles sont tous vides. $$\left \{ \begin{array}{llll} \left[a,b\right]=\left\{t\in\R \ | \ a\leqslant t\leqslant b\right\} \\ \left[a,b\right[=\left\{t\in\R \ | \ a\leqslant t< b\right\} \\ \left]a,b\right]=\left\{t\in\R \ | \ a<t\leqslant b\right\} \\ \left]a,b\right[=\left\{t\in\R \ | \ a<t< b\right\}\end{array}\right.$$
\end{dfn}
\begin{cor}[Variantes]
	Au niveau de l'addition, on a la variante suivante : $$\forall(x,y,z)\in\R^3,\ x\leqslant y \implies x+z\leqslant y+z$$ Au niveau de la multiplication, on a les trois variantes suivantes : $$\forall(x,y,z)\in\R^3,\ \left \{\begin{array}{ll}x\leqslant y \\ z\leqslant 0\end{array}\right. \implies xz\geqslant yz$$ $$\forall(x,y,z)\in\R^3,\ \left \{\begin{array}{ll}x<y \\ z>0\end{array}\right. \implies xz<yz$$ $$\forall(x,y,z)\in\R^3,\ \left \{\begin{array}{ll}x<y \\ z<0\end{array}\right. \implies xz>yz$$
\end{cor}
\begin{rmq}
	On retiendra que dans une inéquation, on peut faire "changer de côté" ce que l'on veut.
	\begin{enumerate}
		\item Dans une somme, à condition que les $+$ deviennent des $-$ et vice-versa.
		\item Dans un produit, à condition que le terme soit non nul, que les $\times$ deviennent des $/$ et vice-versa. Si le facteur est strictement positif le sens de l'inéquation ne change pas, sinon il est inversé.
	\end{enumerate}
\end{rmq}
\begin{rmq}[Passage à l'inverse dans une inéquation]
	La fonction $\displaystyle t\mapsto\frac{1}{t}$ envoie $\R_+\st$ dans lui-même et est strictement décroissante sur ce domaine. même remarque avec $\R_-\st$.
\end{rmq}
\begin{thm}
	Soit $x_1,\ldots,x_n\geqslant0$. Alors $\displaystyle\sum_{i=1}^{n}x_i\geqslant0$, avec égalité si, et seulement si, tous les $x_i$ sont nuls.
\end{thm}
\begin{cor}
	Soit $a_1,\ldots,a_n$ et $b_1,\ldots,b_n$ tels que $\forall i\in\llbracket1,n\rrbracket,\ a_i\leqslant b_i$. Alors $$\sum_{i=1}^{n}a_i\leqslant\sum_{i=1}^{n}b_i$$ avec égalité si, et seulement si, $\forall i\in\llbracket1,n\rrbracket,\ a_i=b_i$.
\end{cor}
\begin{dfn}[Borne supérieure]
	Soit $A$ une partie non vide de $\R$. Si elle existe, la \textbf{borne supérieure de $A$} est le plus petit des majorants de $A$. On la note $\sup(A)$.
\end{dfn}
\begin{dfn}[Borne inférieure]
	Soit $A$ une partie non vide de $\R$. Si elle existe, la \textbf{borne inférieure de $A$} est le plus grand des minorants de $A$. On la note $\inf(A)$.
\end{dfn}
\begin{rmq}
	Notons que si elles existent, $\sup(A)$ et $\inf(A)$ sont nécessairement uniques, puisque ce sont respectivement un maximum et un minimum.
\end{rmq}
\begin{rmq}
	Si $A$ admet un maximum, alors $A$ admet une borne supérieure égale à ce maximum. De même avec le minimum et la borne inférieure.
\end{rmq}
\begin{thm}[Propriété fondamentale de $\R$ - admise]
	Tel qu'il est construit, $\R$ vérifie la propriété fondamentale suivante : toute partie non vide et majorée de $\R$ admet une borne supérieure ; et toute partie non vide et minoré de $\R$ possède une borne inférieure.
\end{thm}
\begin{rmq}
	La seconde propriété peut se voir comme une conséquence de la première, et réciproquement.
\end{rmq}
\begin{exa}
	La borne supérieure de $]-\infty,1[$ vaut $1$.
\end{exa}
\begin{rmq}
	Attention ! Selon les cas, la borne supérieure (ou inférieure) peut appartenir ou non à $A$. Par exemple, la borne supérieure de $]-\infty,1[$ vaut $1$ mais la borne supérieure de $]-\infty,1]$ vaut $1$ aussi.
\end{rmq}
\begin{prop}[Passage au sup/ à l'inf dans une inégalité large]
	Soit $A\subset\R$ non vide.
	\begin{itemize}
		\item Si $\exists M\in\R, \forall x\in A,\ x\leqslant M$, alors $A$ admet une borne supérieure et on a $\sup(A)\leqslant M$.
		\item Si $\exists m\in\R, \forall x\in A,\ x\geqslant m$, alors $A$ admet une borne inférieure et on a $\inf(A)\geqslant m$.
	\end{itemize}
\end{prop}
\begin{rmq}
	$A$ $fortiori$, si $\forall x\in A,\ x<M$, alors on a encore $\sup(A)\leqslant M$. Mais \textbf{attention}, l'inégalité stricte pour le majorant est devenue large pour la borne supérieure ! Et c'est le meilleur résultat que l'on puisse obtenir, il suffit de considérer l'exemple de $]-\infty,1[$ donné précédemment.
\end{rmq}
Cette proposition est souvent  utilisée comme méthode alternative pour déterminer la borne supérieure (ou inférieure) d'une partie. On procède alors par double inégalité.
\begin{exa}
	Si $A$ est une partie non vide et majorée de $\R$, et $\lambda$ est un réel positif ou nul, on a $\sup(\lambda A)=\lambda\sup(A)$.
\end{exa}
\begin{thm}[Partie entière]
	Soit $x$ un réel. Il existe un unique $n\in\Z$ tel que $n\leqslant x< n+1$. On l'appelle la \textbf{partie entière de $x$}, et on le note $\lfloor x\rfloor$.
\end{thm}
\begin{dfn}[Partie fractionnaire]
	Soit $x$ un réel. La \textbf{partie fractionnaire de $x$}, généralement notée Frac$(x)$ (voire $\left\{x\right\}$ s'il n'y a pas d'ambiguïté), est définie par Frac$(x)=x-\lfloor x\rfloor$. Elle appartient donc à $[0,1[$.
\end{dfn}
\begin{rmq}
	Intuitivement, au moins si $x\geqslant0$, $\lfloor x\rfloor$ représente "les chiffres de $x$ avant la virgule" et Frac$(x)$ représente "les chiffres de $x$ après la virgule.
\end{rmq}
\begin{exa}
	Dans une division euclidienne $a=bq+r$, on a toujours $q=\displaystyle\left\lfloor\frac{a}{b}\right\rfloor$.
\end{exa}
\begin{prop}
	Si $x\in\R$ et $m\in\Z$, alors on a $\lfloor x+m\rfloor=\lfloor x\rfloor+m$.
\end{prop}
\begin{rmq}
	Attention, si $m\in\Z$, ne pas écrire $\lfloor mx\rfloor=m\lfloor x\rfloor$, c'est faux en général ! Prendre par exemple $x=1/2$ et $m=2$. De même, on n'a pas $\lfloor x+y\rfloor=\lfloor x\rfloor+\lfloor y\rfloor$. Prendre par exemple $x=y=1/2$.
\end{rmq}
\begin{prop}
	La fonction $x\mapsto\lfloor x\rfloor$ est croissante de $\R$ dans $\R$. La fonction $x\mapsto \left\{x\right\}$ est $1$-périodique de $\R$ dans $\R$.
\end{prop}
\begin{exa}
	Si $x\geqslant n$ avec $x\in\R$ et $n\in\Z$, alors $\lfloor x\rfloor\geqslant n$
\end{exa}
\begin{dfn}[Valeur absolue]
	La \textbf{valeur absolue} d'un réel $x$ est définie par $$\lvert x\rvert=\left \{\begin{array}{ll}x&\text{si}x\geqslant0 \\ -x&\text{si}x\leqslant0\end{array}\right.$$
\end{dfn}
Au passage, on vérifie que la définition est bien consistante lorsque $x=0$. Voici à présent quelques résultats de bon sens, que l'on admet généralement "d'instinct" mais qu'il n'est pas forcément mauvais de démontrer une fois.
\begin{prop}
	On a toujours $\lvert x\rvert \geqslant0$, avec égalité si, et seulement si, $x=0$.
\end{prop}
\begin{prop}
	On a toujours $-\lvert x\rvert\leqslant x\leqslant \lvert x\rvert$.
\end{prop}
\begin{lem}
	On a $x=\pm y$ si, et seulement si, $\lvert x\rvert =\lvert y\rvert$.
\end{lem}
\begin{cor}
	On a $\forall(x,y)\in\R^2,\ \lvert xy\rvert=\lvert x\rvert\times\lvert y\rvert$.
\end{cor}
\begin{rmq}
	De même, on montre que $$\forall y\in\R\st,\ \left\lvert\frac{1}{y}\right\rvert \ \ \text{et} \ \ \forall(x,y)\in\R\times\R\st,\ \left\lvert\frac{x}{y}\right\rvert=\frac{\lvert x\rvert}{\lvert y\rvert}$$
\end{rmq}
\begin{rmq}
	On lit couramment que $\forall x\in\R,\ \sqrt{x^2}=\lvert x\rvert$. cette assertion est vraie, mais encore faut-il avoir construit la fonction $\sqrt{\cdot}$. Cela dit, dès à présent, on peut montrer que l'application de $\R_+$ dans $\R_+$ $t\mapsto t^2$ est strictement croissante. Soit $0\leqslant t< u$. Alors $t^2\leqslant tu$ et $tu< u^2$. On conclut par transitivité. \par Une fois ce point démontré, on en déduit que tout $y\in\R_+$ admet au plus une racine carrée dans $\R_+$. En particulier, soit $x\in\R$. Alors $\lvert x\rvert$ est l'unique racine carrée de $x^2$ dans $\R_+$. En effet, on a $\lvert x\rvert^2=\lvert x^2\rvert=x^2$.
\end{rmq}
\begin{prop}
	Soit $\delta\geqslant0$ fixé. Alors $\forall x\in\R,\ -\delta\leqslant x\leqslant\delta\iff\lvert x\rvert\leqslant\delta$.
\end{prop}
\begin{rmq}
	Plus généralement, soit $a\in\R$ et $\delta\geqslant0$. On a : $$\left \{\begin{array}{ll}\forall x\in\R,\ x\in[a-\delta,a+\delta]\iff\lvert x-a\rvert\leqslant\delta\\ \forall x\in\R,\ x\in]a-\delta,a+\delta[\iff\lvert x-a\rvert<\delta \end{array}\right.$$
\end{rmq}
\begin{prop}[Formules du $\min$ et du $\max$]
	Soit $(x,y)\in\R^2$. On a $$\left \{\begin{array}{ll}\displaystyle\max(x,y)=\frac{x+y}{2}+\frac{\lvert x-y\rvert}{2} \\[4mm] \displaystyle\min(x,y)=\frac{x+y}{2}-\frac{\lvert x-y\rvert}{2}\end{array}\right.$$
\end{prop}
\begin{prop}[Inégalités triangulaires]
	Soit $(x,y)\in\R^2$. On a : $$\left \{\begin{array}{ll} \lvert x+y\rvert \leqslant\lvert x\rvert+\lvert y\rvert \\ \big|\lvert x\rvert -\lvert y\rvert\big|\leqslant\lvert x-y\rvert\end{array}\right.$$
\end{prop}
\begin{dfn}[Partie positive, partie négative]
	Soit $x\in\R$. La \textbf{partie positive} de $x$ est définie par $x^+=\max(x,0)$. La \textbf{partie négative} de $x$ est définie par $x^-=\max(-x,0)$. On a $$\left \{\begin{array}{ll}0\leqslant x^+,x^-\leqslant x \\ x=x^+-x^-\end{array}\right.$$
\end{dfn}
\begin{dfn}[Droite réelle achevée $\overline{\R}$]
	La \textbf{droite réelle achevée $\overline{\R}$} est égale à l'ensemble $\R$ auquel on adjoint deux nouveaux éléments $-\infty$ et $+\infty$. La relation $\leqslant$ se prolonge de manière évidente par $$\forall x\in\overline{\R},\ -\infty\leqslant x\leqslant+\infty$$ On vérifie sans difficulté que $\leqslant$ reste une relation d'ordre totale sur $\overline{\R}$.
\end{dfn}
\begin{rmq}
	Insistons sur le fait que $-\infty$ et $+\infty$ sont de \textbf{nouveaux} éléments. Autrement dit, ils ne sont pas réels et ils sont distincts. En particulier, on en déduit que $$\forall x\in\R,\ -\infty<x<+\infty$$
\end{rmq}
\begin{dfn}
	Soit $a,b\in\overline{\R}$. Pour généraliser le cas fini, on définit les ensembles suivants. $$\left \{ \begin{array}{llll} \left[a,b\right]=\left\{t\in\overline{\R} \ | \ a\leqslant t\leqslant b\right\} \\ \left[a,b\right[=\left\{t\in\overline{\R} \ | \ a\leqslant t< b\right\} \\ \left]a,b\right]=\left\{t\in\overline{\R} \ | \ a<t\leqslant b\right\} \\ \left]a,b\right[=\left\{t\in\overline{\R} \ | \ a<t< b\right\}\end{array}\right.$$ Dans le cas où $a$ et $b$ sont réels, on vérifie sans peine que l'ancienne définition coïncide avec l'ancienne.
\end{dfn}
\begin{exa}
	On a $\overline{\R}=\left[-\infty,+\infty\right]$.
\end{exa}
\begin{exa}
	En particulier, si $a>b$, on obtient que $\left[a,b\right]=\emptyset$.
\end{exa}
\begin{rmq}
	En revanche, l'addition et la multiplication sont plus délicates à prolonger. on ne peut le faire que partiellement en posant $$\left \{\begin{array}{ll}\forall a\in\left]-\infty,+\infty\right],\ a+(+\infty)=(+\infty)+a=+\infty \\ \forall a\in\left[-\infty,+\infty\right[,\ a+(-\infty)=(-\infty)+a=-\infty\end{array}\right.$$ $$\left \{\begin{array}{llll}\forall a\in\left]0,+\infty\right],\ (+\infty)\times a=a\times(+\infty)=+\infty \\ \forall a\in\left]0,+\infty\right],\ (-\infty)\times a=a\times(-\infty)=-\infty \\ \forall a\in\left[-\infty,0\right[,\ (+\infty)\times a=a\times(+\infty)=-\infty \\ \forall a\in\left[-\infty,0\right[,\ (-\infty)\times a=a\times(-\infty)=+\infty\end{array}\right.$$ Autrement dit, il y a ce qu'on appelle deux \textbf{formes indéterminées}, qui sont du type $\infty-\infty$ et $0\times \infty$. Nous verrons dans le chapitre "Sommabilité" un ultime prolongement, mais nous n'en faisons pas usage dans l'immédiat.
\end{rmq}
\begin{exa}
	Sur $\left]-\infty,+\infty\right]$, la loi $+$ reste associative et commutative. $0$ reste neutre, et $+\infty$ absorbant. Remarque symétrique pour $\left[-\infty,+\infty\right[$.
\end{exa}
\begin{rmq}
	Soit $A\subset\R$ non vide. On montre facilement que $A$ admet une borne supérieure au sens de $\overline{\R}$.
	\begin{itemize}
		\item Si $A$ est majorée par un réel, on montre que $\sup_{\overline{\R}}(A)=\sup(A)$
		\item Sinon, on montre que $\sup_{\overline{\R}}=+\infty$
	\end{itemize}
	Le premier point permet d'écrire $\sup(A)$ sans risque de confusion. En résumé, on peut dire que $\sup(A)\in\R\cup\left\{+\infty\right\}$, et que $A$ est majorée par un réel si, et seulement si, $\sup(A)<+\infty$. Bine sûr, on peut faire la remarque symétrique avec la borne inférieure.
\end{rmq}
\begin{dfn}[Intervalle]
	Un \textbf{intervalle} de $\R$ est une \textbf{partie convexe} de $\R$. Autrement dit, c'est une partie $I\subset\R$ telle que $$\forall (x,y)\in I^2,\ x\leqslant y\implies \left[x,y\right]\subset I$$
\end{dfn}
\begin{rmq}
	En pratique, on peut se restreindre à $x<y$ puisque si $x=y$, alors $\left[x,y\right]=\left\{x\right\}\subset I$ car $x\in I$. C'est ce que l'on fera à chaque fois dans les exercices.
\end{rmq}
\begin{thm}[Caractérisation des intervalles non vides de $\R$]
	Les intervalles non vides de $\R$ sont :
	\begin{itemize}
		\item les $\left[a,b\right]$ avec $-\infty<a\leqslant b<+\infty$
		\item les $\left]a,b\right]$ avec $-\infty\leqslant a<b<+\infty$
		\item les $\left[a,b\right[$ avec $-\infty<a<b\leqslant+\infty$
		\item les $\left]a,b\right[$ avec $-\infty\leqslant a< b\leqslant+\infty$
	\end{itemize}
\end{thm}
\begin{rmq}
	Soit $I\subset\R$ un intervalle non vide. La démonstration précédente montre qu'il suffit de choisir $a=\inf(I)$ et $b=\sup(I)$. Mais on est en droit de se demander s'il s'agit d'une condition nécessaire. La réponse est oui.
	\begin{itemize}
		\item Si $I$ s'écrit d'une des quatre manières ci-dessus, on montre que nécessairement $a=\inf(I)$ et $b=\sup(I)$. En particulier, $a$ et $b$ sont uniques. On les appelle les \textbf{bornes} de $I$. Elles sont dites \textbf{ouvertes} ou \textbf{fermées} selon que l'inégalité qui les concerne est stricte ou large.
		\item Au passage, on en déduit que les quatre cas ci-dessus sont deux à deux exclusifs : un intervalle ne peut pas être à la fois ouvert et fermé, etc.
		\item Si $a$ et $b$ l'intervalle est dit \textbf{borné}. Un intervalle fermé borné s'appelle aussi un \textbf{segment}.
	\end{itemize}
\end{rmq}
\begin{dfn}[Intérieur]
	Soit $I\subset\R$ un intervalle non vide, ainsi que $a=\inf(I)$ et $b=\sup(I)$ ses deux bornes, avec $a\leqslant b$ dans $\overline{\R}$. L'\textbf{intérieur} de $I$ est l'intervalle $]a,b[$. On le note $\mathring{I}$.
\end{dfn}
\begin{exa}
	Un intervalle non vide $I$ est d'intérieur non vide si, et seulement si, il contient au moins deux points.
\end{exa}
\begin{prop}
	L'intersection de deux intervalles est un intervalle.
\end{prop}
\begin{rmq}
	En revanche, l'union de deux intervalles n'a aucun raison d'être un intervalle en général. Considérer par exemple $\left[0,1\right]\cup\left[2,3\right]$.
\end{rmq}
\begin{dfn}[Partie dense dans $\R$]
	Soit $A$ une partie de $\R$. Les deux propositions suivantes sont équivalentes :
	\begin{itemize}
		\item $A$ rencontre tout intervalle ouvert borné non trivial :$$\forall(a,b)\in\R^2,\ a<b\implies A\cap]a,b[\ne\emptyset$$
		\item $A$ rencontre tout intervalle fermé borné non trivial :$$\forall(a,b)\in\R^2,\ a<b\implies A\cap[a,b]\ne\emptyset$$
	\end{itemize}
	Lorsque $A$ vérifie l'une de ces assertions, on dit que \textbf{$A$ est dense dans $\R$}.
\end{dfn}
On sait qu'il n'existe pas de rationnel $r$ tel que $r^2=2$. Sinon, il existerait deux entiers $p,q\in\Z\st$ tels que $2q^2=p^2$ donc on aurait $2v_2(p)=v_2(p^2)=v_2(2q^2)=1+2v_2(q)$. Or, un entier ne pouvant pas à la fois être pair et impair, on aurait une contradiction. On en tire plusieurs conséquences.
\begin{thm}
	$\Q$ et $\R\setminus\Q$ sont denses dans $\R$.
\end{thm}
\begin{rmq}
	Dans $\Q$, la propriété de la borne supérieure est fausse. En effet, il suffit de considérer l'ensemble $]0,\sqrt{2}[\cap\Q$. Supposons qu'il possède une borne supérieure $\alpha\in\Q$. Si on a $\alpha>\sqrt{2}$, alors, comme il existe un rationnel dans $]\sqrt{2},\alpha[$, ce rationnel est un majorant, ce qui est une contradiction car $\alpha$ doit être le plus petit majorant. De même, on ne peut pas avoir $\alpha<\sqrt{2}$. Donc on doit avoir $\alpha=\sqrt{2}$. Contradiction, car $\sqrt{2}$ n'est pas rationnel.
\end{rmq}
Nous verrons ultérieurement que le TVI provient essentiellement de la propriété de la borne supérieure, puis que beaucoup de théorèmes d'analyse réelle proviennent eux-mêmes du TVI. On voit donc que la propriété de la borne supérieure est véritablement fondatrice. Le fait qu'elle soit fausse dans $\Q$ limite considérablement les théorèmes que l'on peut y trouver.
\begin{dfn}[Ensemble des décimaux $\mathbb{D}$]
	L'ensemble des \textbf{nombres décimaux} est défini par $$\mathbb{D}=\left\{\frac{p}{10^n} \ | \ (p,n)\in\Z\times\N\right\}$$
\end{dfn}
\begin{exa}
	On peut montrer que $\mathbb{D}$ est un anneau intègre.
\end{exa}
\begin{dfn}[Valeur approchée par excès, par défaut]
	Soit $x\in\R$ et $n\in\N$. Alors il existe un unique $p\in\Z$ tel que $$\frac{p}{10^n}\leqslant x<\frac{p}{10^n}+\frac{1}{10^n}$$ Par définition, cet entier $p$ vaut $\lfloor10^nx\rfloor$. $\displaystyle\frac{p}{10^n}$ s'appelle la \textbf{valeur décimale approchée de $x$ à $10^{-n}$ près}. On dit que c'est la valeur approchée \textbf{par défaut}, tandis que $\displaystyle\frac{p}{10^n}+\frac{1}{10^n}$ est la valeur approchée \textbf{par excès}.
\end{dfn}
Le théorème qui suit se révèle souvent bien utile, même s'il ne figure pas au programme.
\begin{thm}[Sous-groupes additifs de $\R$, HP à connaître]
	Les sous-groupes additifs de $\R$ sont soit de la forme $a\Z$ avec $a\in\R_+$, soit denses dans $\R$.
\end{thm}
\begin{proof}
	Soit $H$ un sous-groupe additif de $\R$. Si $H=\left\{0\right\}$, alors il est bien de la forme $a\Z$ avec $a=0$. Sinon, il contient un réel non nul, et quitte à passer à l'opposé, il contient un réel strictement positif. Donc $H\cap\R_+\st$ est non vide, et il admet une borne inférieure $a$.
	\begin{itemize}
		\item Supposons que $a\in H$, et montrons alors que $H=a\Z$. \\ D'une part, on a $<a>=a\Z\subset H$ par définition du sous-groupe engendré. Réciproquement, soit $x\in H$. Comme $a\in\R_+\st$, on peut considérer $q=\displaystyle\left\lfloor\frac{x}{a}\right\rfloor\in\Z$ et $r=x-aq$. Alors, d'après la première inclusion, $r\in H$, et comme $r\in[0,a[$, on a $r=0$ par minimalité de $a$. Donc $x=aq\in a\Z$.
		\item Supposons que $a\notin H\cap\R_+\st$., et montrons alors que $H$ est dense dans $\R$. \\ Donnons-nous un intervalle $]x,y[$ de largeur $\varepsilon=y-x>0$.
		\begin{itemize}
			\item Montrons qu'on peut trouver un $h\in]0,\varepsilon[$ dans $H$. Prenons dans $H\cap\R_+\st$ un $h_1$ tel que $h_1\in[a,a+\varepsilon[$. Comme $a\notin H\cap\R_+\st$, on a même $h_1\in]a,a+\varepsilon[$. Par le même raisonnement, prenons dans $H$ un $h_2\in]a,h_1[$. il suffit alors de poser $h=h_1-h_2$.
			\item Ensuite, posons $q=\displaystyle\left\lfloor\frac{x}{h}\right\rfloor+1\in\Z$. Alors $qh$ appartient à $H$ et il vérifie $x<qh<y$.
		\end{itemize}
		\item Enfin, aucun $H=a\Z$ ne peut être dense dans $\R$ (il s'agit bien d'un "soit/soit"). En effet :
		\begin{itemize}
			\item Si $a=0$, alors l'intervalle $]0,1[$ ne contient aucun élément de $H$.
			\item Si $a>0$, alors l'intervalle $]0,a[$ ne contient aucun élément de $H$.
		\end{itemize}
	\end{itemize}
	Ainsi, la preuve est achevée.
\end{proof}

\section{Suites réelles}
\begin{dfn}[Suite réelle]
	Une \textbf{suite réelle} $(u_n)_{n\in\N}$ est une application de $\N$ dans $\R$ qui à $n$ associe $u_n$. $n$ s'appelle l'\textbf{indice}. On peut plus prosaïquement la noter $u$ s'il n'y a pas d'ambiguïté sur l'ensemble que parcourt l'indice.
\end{dfn}
\begin{dfn}[Relation $\leqslant$ sur $\R^\N$]
	On définit une relation d'ordre sur $\R^\N$ par $$u\leqslant v\iff \forall n\in\N,\ u_n\leqslant v_n$$ C'est une relation d'ordre partielle.
\end{dfn}
\begin{dfn}[Suite majorée, minorée, bornée]
	Une suite $u$ est \textbf{majorée} lorsque l'ensemble $\left\{u_n\ | \ n\in\N\right\}$ est majoré. Elle est \textbf{minorée} si cet ensemble est minoré. Elle est \textbf{bornée} si elle est à la fois majorée et minorée.
\end{dfn}
\begin{rmq}
	Pour montrer efficacement qu'une suite $(u_n)_{n\in\N}$ est bornée, on montrera de manière équivalente de $(\lvert u_n\rvert)_{n\in\N}$ est majorée (une seule opération au lieu de deux).
\end{rmq}
\begin{exa}
	Pour montrer que $(\cos(n))_{n\in\N}$ est bornée, il suffit de remarquer que $\forall n\in\N,\ \lvert\cos(n)\rvert\leqslant1$.
\end{exa}
\begin{exa}
	Les suites bornées forment un sous-anneau de $\R^\N$.
\end{exa}
\begin{dfn}[Suite croissante, décroissante, monotone, stationnaire]
	Une suite $u$ est \textbf{croissante} lorsque $$\forall(m,n)\in\N^2,\ m\leqslant n \implies u_m\leqslant u_n$$ En pratique, il suffira de vérifier que $\forall n\in\N,\ u_n\leqslant u_{n+1}$. Elle est \textbf{décroissante} lorsque $$\forall(m,n)\in\N^2,\ m\leqslant n \implies u_m\geqslant u_n$$ En pratique, il suffira de vérifier que $\forall n\in\N,\ u_n\geqslant u_{n+1}$. Elle est \textbf{monotone} si elle croissante ou décroissante. Elle est \textbf{stationnaire} si elle est constante à partir d'un certain rang.
\end{dfn}
\begin{rmq}
	On définit de même les notions de suite strictement croissante, strictement décroissante et strictement monotone. En revanche, les notions de suite strictement majorée, strictement minorée et strictement bornée sont sans intérêt, car elles sont équivalentes aux mêmes notions sans le "strictement".
\end{rmq}
\begin{dfn}[Extension]
	On peut définir de même des suites réelles indexées sur $\Z$ : $(u_n)_{n\in\Z}$ ; sur $\llbracket n_0,+\infty\llbracket$ pour $n_0\in\Z$ : $(u_n)_{n\in \llbracket n_0,+\infty\llbracket}$ ; ou même sur deux indices à la fois, c'est ce qu'on appelle une \textbf{suite double} : $(u_{m,n})_{(m,n)\in\N^2}$ ou $(u_{m,n})_{(m,n)\in\Z^2}$. \par Les relations d'ordre $\leqslant$ et $\geqslant$ sont toujours définies, ainsi que les notions de suites majorées, minorées et bornées. En revanche, si les notions de suite croissante, décroissante et monotone ont encore du sens pour $(u_n)_{n\in\Z}$ et $(u_n)_{n\in \llbracket n_0,+\infty\llbracket}$, elles n'ont plus de sens pour les suites doubles.
\end{dfn}
\begin{dfn}[Convergence, divergence]
	Soit $l\in\R$. On dit que la suite réelle $(u_n)_{n\in\N}$ \textbf{converge vers $l$} ou que \textbf{$u_n$ tend vers $l$} lorsque $$\forall\varepsilon>0,\ \exists N\in\N,\ \forall n\in\N,\ n\geqslant N\implies \lvert u_n-l\rvert\leqslant\varepsilon$$ Le nombre $l$ s'appelle \textbf{une limite de $u$} et on écrit $u_n\xrightarrow[n\to+\infty]{}l$ ou encore $u_n\rightarrow l$, voire $\displaystyle \underset{n\to+\infty}{\lim}u_n=l$. Une suite est dite \textbf{convergente} lorsqu'elle admet une limite finie : $\exists l\in\R,\ u_n\rightarrow l$. Au contraire, elle est dite \textbf{divergente} lorsqu'elle n'admet pas de limite finie :$$\forall l\in\R,\ \exists\varepsilon>0,\ \forall N\in\N,\ \exists n\in\N,\ (n\geqslant N)\land(\lvert u_n-l\rvert>\varepsilon)$$
\end{dfn}
\begin{rmq}
	On montre facilement que la définition de la convergence est équivalente si on remplace les $\leqslant\varepsilon$ par des $<\varepsilon$. On utilisera régulièrement cette remarque pour des raisons de confort lors de certains démonstrations.
\end{rmq}
\begin{rmq}
	De manière plus intuitive, on peut dire que $(u_n)_{n\in\N}$ converge vers $l$ lorsque, aussi petit soit le couloir autour de $l$, $u_n$ finit par se "laisser apprivoiser" et se retrouve alors coincé dans ce couloir au bout 'un certain temps. \par Au contraire, si la suite diverge, cela signifie qu'elle est "sauvage" (grrr Fred) et que toute tentative pour la faire tendre vers une limite échouera. Plus précisément, quel que soit $l\in\R$, $u_n$ ne se "laisse pas apprivoiser" et il existe un "couloir récalcitrant" d'épaisseur non nulle autour de $l$ tel que $u_n$ ressorte régulièrement de ce couloir, aussi loin qu'on aille chercher.
\end{rmq}
\begin{exa}
	Soit $l\in\R$. Si $(u_n)_{n\in\N}$ est constante égale à $l$, alors elle converge vers $l$. Plus généralement, si elle stationne à la valeur $l$, alors elle tend vers $l$.
\end{exa}
\begin{rmq}
	Soit $l\in\R$. Quel que soit $\varepsilon>0$, on a $\lvert u_n-l\rvert\leqslant\varepsilon$ ssi $\lvert(u_n-l)-0\rvert\leqslant\varepsilon$ ssi $\left\lvert \lvert u_n-l\rvert -0\right\rvert\leqslant\varepsilon$. Donc $$u_n\rightarrow l\iff u_n-l\rightarrow 0 \iff \lvert u_n-l\rvert\rightarrow0$$ en particulier, pour $l=0$, on obtient $$u_n\rightarrow l\iff\lvert u_n\rvert \rightarrow0$$
\end{rmq}
\begin{prop}[Unicité de la limite]
	Si $u$ possède une limite, alors celle-ci est unique. A partir de maintenant, on pourra donc parler de \textbf{la} limite d'une suite convergente.
\end{prop}
\begin{rmq}
	On utilisera souvent ce résultat pour montrer l'égalité de deux nombres, en calculant par exemple une limite de deux façons différentes et en concluant par unicité de la limite.
\end{rmq}
\begin{thm}
	Toute suite convergente est bornée.
\end{thm}
\begin{rmq}
	On ne peut pas prendre le maximum d'un nombre infini de réels, or, il y a un nombre infini de $\lvert u_n\rvert$. On voit donc toute l'importance de "circonscrire" le domaine sur lequel évolue $n$ pour prendre le maximum, ce qui nous avons pu faire grâce à l'hypothèse de convergence. \par C'est une technique extrêmement générale, qui consiste à couper la suite en deux parties (voire plus). Une première partie finie, sur laquelle on peut faire à peu près ce que l'on veut (prendre le minimum, le maximum, calculer un produit, une somme...) ; et une deuxième partie infinie sur laquelle on a besoin d'hypothèses supplémentaires (majorant, minorant, convergence, etc.) pour pouvoir agir en une fois sur une infinité d'indices. Souvent, la première partie sera plus facile à traiter que la deuxième. \par On trouvera un exemple de cette technique dans la partie des compléments sur la moyenne de Cesàro.
\end{rmq}
\begin{dfn}
	On dit que $u_n$ tend vers $+\infty$ lorsque $$\forall M\in\R,\ \exists N\in\N,\ \forall n\in\N,\ n\geqslant N\implies u_n\geqslant M$$ On note $u_n\xrightarrow[n\to+\infty]{}+\infty$. Là aussi, la définition est équivalente si on remplace $\geqslant M$ par $>M$.
\end{dfn}
\begin{rmq}
	La définition est également équivalente en remplaçant $M\in\R$ par $M\in\R_+$ ou $M\in\R_+\st$. Elle est plus généralement équivalente avec $M\in]A,+\infty[$ ou $M\in[A,+\infty[$ avec $A\in\R$ fixé.
\end{rmq}
\begin{rmq}
	De manière imagée, lorsque $u_n\rightarrow +\infty$, on peut dire que aussi haute soit l'altitude, $u_n$ finit par la dépasser définitivement à partir d'un certain temps.
\end{rmq}
\begin{dfn}
	On dit que $u_n$ tend vers $-\infty$ lorsque $$\forall m\in\R,\ \exists N\in\N,\ \forall n\in\N,\ n\geqslant N\implies u_n\leqslant m$$ On note $u_n\xrightarrow[n\to+\infty]{}-\infty$. Là aussi, la définition est équivalente si on remplace $\leqslant m$ par $<m$.
\end{dfn}
\begin{rmq}
	La définition est également équivalente en remplaçant $m\in\R$ par $m\in\R_-$ ou $m\in\R_-\st$. Elle est plus généralement équivalente avec $m\in]-\infty,A[$ ou $m\in]-\infty,A]$ avec $A\in\R$ fixé.
\end{rmq}
\begin{rmq}
	De manière imagée, lorsque $u_n\rightarrow +\infty$, on peut dire que aussi basse soit la profondeur, $u_n$ finit par la dépasser définitivement à partir d'un certain temps.
\end{rmq}
\begin{rmq}
	Si $u_n\rightarrow+\infty$, alors $u$ est divergente. De même avec $-\infty$. C'est la contraposée du théorème qui stipule que tout suite convergente est bornée.
\end{rmq}
Par ailleurs, elle ne peut pas tendre à la fois vers $-\infty$ et $+\infty$, ou vers $\pm\infty$ et un réel. Pour toutes ces raisons, si $u$ admet une limite dans $\overline{\R}$, celle-ci continue à être unique.
\begin{exa}
	La suite définie par $u_n=n$ tend vers $+\infty$. Plus généralement, pour tout $k\in\N\st$, la suite définie par $u_n=n^k$ tend vers $+\infty$.
\end{exa}
\begin{exa}
	La suite définie sur $\N\st$ par $u_n=\displaystyle\frac1n$ tend vers $0$. Plus généralement, pour tout $k\in\N\st$, la suite définie sur $\N\st$ par $u_n=\displaystyle\frac{1}{n^k}$ tend vers $0$.
\end{exa}
\begin{prop}[Principe de troncature]
	Soit $p\in\N\st$ et $l\in\overline{\R}$. La suite $(u_n)_{n\in\N}$ tend vers $l$ si, et seulement si, la suite $(u_n)_{n\geqslant p}$ tend vers $l$.
\end{prop}
\begin{exa}
	Si $(u_n)_{n\in\N}$ et $(v_n)_{n\in\N}$ coïncident à partir d'un certain rang, alors $u_n\rightarrow l$ si, et seulement si, $v_n\rightarrow l$. On dit que la limite est une notion \textbf{asymptotique}.
\end{exa}
\begin{exa}
	Soit $u\in\R^\N$ une suite qui ne s'annule pas à partir d'un certain rang. Lorsqu'on écrire une assertion comme $$\frac{1}{u_n}\rightarrow l$$ il sera inutile de se préoccuper du rang à partir duquel la suite $\displaystyle\left(\frac{1}{u_n}\right)$ a été définie (il suffit de savoir qu'il en existe au moins un).
\end{exa}
\begin{prop}[Principe de décalage d'indice]
	Soit $p\in\N\st$ et $l\in\overline{\R}$. La suite $(u_n)_{n\geqslant p}$ tend vers $l$ si, et seulement si, la suite $(u_{n+p})_{n\geqslant0}$.
\end{prop}
\begin{exa}
	Si $u_n\rightarrow l$, alors par troncature et décalage d'indice, on obtient que $u_{n+1}\rightarrow l$.
\end{exa}
\begin{thm}[Théorème d'encadrement]
	Soit $l\in\R$, et $u$, $v$ et $w$ trois suites telles que $$\left \{\begin{array}{lll}u_n\rightarrow l\\ w_n\rightarrow l \\ \forall n\in\N,\ u_n\leqslant v_n\leqslant w_n\end{array}\right.$$ Alors, $v_n\rightarrow l$.
\end{thm}
\begin{rmq}
	Ce théorème est aussi appelé "théorème des gendarmes" à cause de sa forme imagée. Si on imagine $v_n$ comme la trajectoire d'un voleur qui se fait poursuivre par deux gendarmes $u_n$ et $w_n$, l'un à sa gauche et l'autre à sa droite, si les deux gendarmes convergent vers la même limite, alors ils "coincent le voleur en le prenant en sandwich".
\end{rmq}
\begin{cor}
	Si $\lvert u_n\rvert\leqslant\alpha_n$ avec $\alpha_n\rightarrow0$, alors $u_n\rightarrow0$.
\end{cor}
\begin{exa}
	La suite de terme général $\displaystyle\frac{\cos(n)}{n}$ définie pour $n\geqslant1$ converge vers $0$.
\end{exa}
\begin{thm}[Théorème de minoration/majoration]
	Supposons que $\forall n\in\N,\ u_n\leqslant v_n$. Si $u_n\rightarrow+\infty$, alors $v_n\rightarrow+\infty$ (minoration). Si $v_n\rightarrow-\infty$, alors $u_n\rightarrow-\infty$ (majoration).
\end{thm}
\begin{rmq}
	Le théorème d'encadrement, son corollaire et le théorème de comparaison restent vrais si on n'a les inégalités qu'à partir d'un certain rang.
\end{rmq}
Le théorème suivant est à mettre en regard avec ce que nous avons vu sur les bornes supérieures et les bornes inférieures.
\begin{thm}[Passage à la limite dans une inégalité dans une inégalité large]
	Soit $u$ une suite qui converge vers une limite $l\in\R$. Soit $m$ et $M$ deux réels. \begin{itemize}
		\item Si $\forall n\in\N,\ u_n\leqslant M$, alors $l\leqslant M$.
		\item Si $\forall n\in\N,\ u_n\geqslant m$, alors $l\geqslant m$.
	\end{itemize}
\end{thm}
\begin{rmq}
	Attention, quand bien même on aurait une inégalité stricte $u_n<M$, celle-ci deviendrait large en passant à la limite !
\end{rmq}
\begin{rmq}
	Le théorème reste vrai si les inégalités n'ont lieu qu'à partir d'un certain rang. Le théorème reste également vrai si on autorise $l$ à être infinie.
\end{rmq}
\begin{exa}
	Soit $u$ et $v$ deux suites telles que $u_n\rightarrow l\in\R$ et $v_n\rightarrow l'\in\R$. Si $$\forall n\in\N,\ u_n\leqslant v_n$$ alors $l\leqslant l'$ (considérer $u_n-v_n$ et lui appliquer le théorème qui précède).
\end{exa}
Voici maintenant une sorte de réciproque partielle.
\begin{thm}
	Soit $a<b$ des réels, et $(u_n)_{n\in\N}$ une suite de réels qui converge vers $l\in\R$.
	\begin{itemize}
		\item Si $l>a$, alors $u_n>a$ à partir d'un certain rang.
		\item Si $l<b$, alors $u_n<b$ à partir d'un certain rang.
		\item Si $l\in]a,b[$, alors $u_n\in]a,b[$ à partir d'un certain rang.
	\end{itemize}
\end{thm}
\begin{exa}
	Soit $(u_n)_{n\in\N}$ qui tend vers $l\in\overline{\R}$. Si $l\ne0$, alors $u_n\ne0$ APCR.
\end{exa}
\begin{exe}
	Soit $u\in\Z^\N$ qui converge vers une limite finie $l$. Montrer que $l\in\Z$ et que $u$ est stationnaire.
\end{exe}
Nous arrivons à une séquence importante. Depuis le lycée, on effectue régulièrement des "opérations algébriques sur les limites"., sans trop savoir si on la droit de le faire. Nous nous employons maintenant à le justifier.
\begin{prop}[Combinaison linéaire]
	Soit $u$ et $v$ des suites réelles qui convergent respectivement vers $l$ et $l'$. Alors, pour tous $\lambda,\mu\in\R$, $\lambda u +\mu v$ converge vers $\lambda l+\mu l'$.
\end{prop}
\begin{exa}
	Si $u_n\rightarrow l$, alors $\lambda u_n\rightarrow \lambda l$.
\end{exa}
\begin{exa}
	Si $u$ converge, alors nécessairement $u_{n+1}-u_n\rightarrow0$.
\end{exa}
\begin{cor}
	Si $u$ est bornée et $v_n\rightarrow0$, alors $u_nv_n\rightarrow0$.
\end{cor}
\begin{prop}[Produit]
	Soit $u$ et $v$ des suites réelles qui convergent respectivement vers $l$ et $l'$. Alors $uv$ converge vers $ll'$.
\end{prop}
\begin{exa}
	Les suites convergentes forment donc un sous-anneau des suites bornées.
\end{exa}
En utilisant la deuxième inégalité triangulaire, on obtient la
\begin{prop}[Valeur absolue]
	Si $u_n\rightarrow l$, alors $\lvert u_n\rvert\rightarrow \lvert l\rvert$.
\end{prop}
Avant de traiter le passage à l'inverse, nous avons besoin du
\begin{lem}
	Si $u_n\rightarrow l\in\R_+\st$, alors $u$ est minorée par un certain $\eta>0$ à partir d'un certain rang.
\end{lem}
\begin{rmq}
	Attention à ce résultat assez subtil : ce n'est pas parce qu'on a $u_n>0$ pour tout $n$ que $u$ peut être minorée par un réel strictement positif. Prendre comme contre-exemple la suite $\displaystyle\left(\frac{1}{n}\right)_{n\in\N\st}$.
\end{rmq}
\begin{prop}[Passage à l'inverse]
	Supposons que $u$ ne s'annule par et qu'elle converge vers un réel $l\ne0$. Alors $$\frac{1}{u_n}\rightarrow\frac1l$$
\end{prop}
\begin{rmq}
	Dans un exercice, même si on ne suppose pas $a$ $priori$ que $u$ ne s'annule pas, le fait que $l\ne0$ garanti ue la suite ne s'annule pas à partir d'un certain rang, ce qui permet de conserver le résultat.
\end{rmq}
\begin{prop}
	Les limites passent encore aux opérations algébriques (combinaison linéaire, produit, valeur absolue) dans $\overline{\R}$, pourvu qu' on n'ait pas de forme indéterminée.
\end{prop}
\begin{dfn}
	On écrit $u_n\rightarrow0^+$ lorsque $u_n\rightarrow0$ et $u_n>0$ à partir d'un certain rang. On écrit $u_n\rightarrow0^-$ lorsque $u_n\rightarrow0$ et $u_n<0$ à partir d'un certain rang.
\end{dfn}
\begin{prop}
	Supposons que $u$ ne s'annule pas. Si $u_n\rightarrow0^+$, alors $\displaystyle\frac{1}{u_n}\rightarrow+\infty$, et si $u_n\rightarrow+\infty$, alors $\displaystyle\frac{1}{u_n}\rightarrow0^+$. De même, si $u_n\rightarrow0^-$, alors $\displaystyle\frac{1}{u_n}\rightarrow-\infty$, et si $u_n\rightarrow-\infty$, alors $\displaystyle\frac{1}{u_n}\rightarrow0^-
	$.
\end{prop}
\begin{rmq}
	En revanche, une limite du type $\displaystyle\frac10$ doit être considérée comme une forme indéterminée. La proposition ci-dessus montre qu'on peut obtenir $-\infty$ comme $+\infty$. On peut même imaginer que la suite $\displaystyle\frac{1}{u_n}$ n'ait pas de limite : on pourra considérer $\displaystyle u_n=\frac{(-1)^n}{n}$.
\end{rmq}
\begin{thm}
	Soit $\alpha$ un réel.
	\begin{itemize}
		\item Si $\alpha>1$, alors $\alpha^n\rightarrow+\infty$.
		\item Si $\lvert\alpha\rvert<1$, alors $\alpha^n\rightarrow0$.
	\end{itemize}
\end{thm}
Intuitivement, une suite extraite de $u$ consiste à ne garder que certains termes en en "sautant" d'autres, et en s'interdisant de "revenir en arrière". Formellement, on la 
\begin{dfn}[Suite extraite, extractrice]
	Une \textbf{suite extraite} de $u$ est une suite de la forme $(u_{\varphi(n)})_{n\in\N}$, où $\varphi$ est strictement croissante de $\N$ dans $\N$. C'est en fait l'application $u\circ\varphi$ de $\N$ dans $\R$. On dit que $\varphi$ est une \textbf{extractrice}. Intuitivement, les $\varphi(n)$ son les indices que l'on garde dans $u$.
\end{dfn}
\begin{prop}
	Si $\varphi$ est strictement croissante de $\N$ dans $\N$, alors c'est une injection et on a $$\forall n\in\N,\ \varphi(n)\geqslant n$$
\end{prop}
\begin{thm}
	Si $u$ tend vers $l\in\overline{\R}$, alors toute suite extraite de $u$ tend vers $l$.
\end{thm}
\begin{cor}
	Pour montrer qu'une suite n'admet pas de limite, il suffit d'exhiber deux suites extraites qui tendent vers des limites différentes.
\end{cor}
\begin{exa}
	Considérer les suites extraites de rang pair et impair pour montrer que la suite de terme général $u_n=(-1)^n$ n'admet pas de limite. Plus généralement, on peut montrer e cette façon qui si $\alpha\leqslant-1$, alors la suite de terme général $\alpha^n$ n'admet pas de limite.
\end{exa}
Réciproquement, on rencontre souvent le cas particulier suivant dans les exercices.
\begin{prop}
	Supposons que $u_{2n}\rightarrow m$ et $u_{2n+1}\rightarrow l$, avec $l\in\overline{\R}$. Alors, $u_n\rightarrow l$.
\end{prop}
\begin{thm}[Caractérisation séquentielle de la densité]
	Une partie $X\subset\R$ est dense dans $\R$ si, et seulement si, tout réel est limite d'une suite $(x_n)_{n\in\N}$ à valeurs dans $X$.
\end{thm}
\begin{cor}
	Tout réel est limite d'une suite de rationnels.
\end{cor}
\begin{exa}
	$\mathbb{D}$ est dense dans $\R$. En effet, si on se donne $x\in\R$, il suffit de considérer la suite de ses approximations décimales (ce que l'on avait déjà fait pour la densité de $\Q$ dans $\R$).
\end{exa}
Les deux théorèmes qui suivent constituent ce qu'on appelle la "\textbf{caractérisation séquentielle de la borne supérieure/inférieure}".
\begin{thm}
	Soit $X$ une partie non vide de $\R$. Alors il existe une suite $(x_n)_{n\in\N}$ à valeurs dans $X$ qui tend vers $\sup(X)\in\R\cup\left\{+\infty\right\}$. On a le résultat symétrique avec la borne inférieure.
\end{thm}
\begin{thm}
	Réciproquement, soit $M\in\R\cup\left\{+\infty\right\}$ un majorant de $X$, et supposons qu'il existe une suite $(x_n)_{n\in\N}$ à valeurs dans $X$ telle que $x_n\rightarrow M$. Alors $M=\sup(X)$. de même avec un minorant et la borne inférieure.
\end{thm}
\begin{exa}
	On peut alors retrouver le résultat $\sup(\lambda A)=\lambda\sup(A)$ avec cette caractérisation séquentielle.
\end{exa}
\section{Théorèmes d'existence de limites}
\begin{thm}[Théorème de la limite monotone]
	Toute suite réelle monotone admet une limite dans $\overline{\R}$.
\end{thm}
\begin{rmq}
	De manière précise, la preuve nous montre qu'on a les résultats suivants.
	\begin{itemize}
		\item Toute suite réelle croissante admet une limite dans $\R\cup\left\{+\infty\right\}$.
		\begin{itemize}
			\item Si $(u_n)_{n\in\N}$ est majorée par $M$, alors elle converge et sa limite vérifie $l\leqslant M$. Attention ! Même si tous les $u_n$ sont strictement inférieurs à $M$, on doit garder l'inégalité au sens large pour la limite.
			\item Si $(u_n)_{n\in\N}$ n'est pas majorée, alors elle tend vers $+\infty$.
		\end{itemize}
		\item Toute suite réelle décroissante admet une limite dans $\R\cup\left\{-\infty\right\}$.
		\begin{itemize}
			\item Si $(u_n)_{n\in\N}$ est minoré par $m$, alors elle converge et sa limite vérifie $l\geqslant m$.
			\item Si $(u_n)_{n\in\N}$ n'est pas minorée, alors elle tend vers $-\infty$.
		\end{itemize}
	\end{itemize}
\end{rmq}
\begin{rmq}
	On abrègera souvent le nom de ce théorème en "TLM" (NS).
\end{rmq}
\begin{dfn}[Suites adjacentes]
	Deux suites sont dites adjacentes lorsque l'une est décroissante, l'autre est croissante et leur différence tend vers $0$.
\end{dfn}
\begin{lem}
	Si $u$ et $v$ sont adjacentes (avec $u$ décroissante et $v$ croissante), on a : $$\forall n \in\N,\ u_n\geqslant v_n$$
\end{lem}
\begin{thm}
	Deux suivantes adjacentes convergent vers une même limite finie.
\end{thm}
\begin{exa}
	On montre aisément que $\forall n\in\N,\ u_n\geqslant l\geqslant v_n$.
\end{exa}
\begin{thm}[Théorème de Bolzano-Weierstrass]
	Soit $(u_n)_{n\in\N}$ une suite de réels bornée. Alors on peut trouver une suite extraite qui converge.
\end{thm}
\begin{rmq}
	On abrègera souvent le nom de ce théorème en "TBW" (NS).
\end{rmq}
\begin{dfn}[Valeur d'adhérence]
	On dit que $l$ est une valeur d'adhérence de la suite $(u_n)_{n\in\N}$ lorsqu'il existe une suite extraite de $(u_n)_{n\in\N}$ qui converge vers $l$.
\end{dfn}
\begin{rmq}
	Le théorème de Bolzano-Weierstrass peut donc être reformulé de la manière suivante : toute suite bornée admet une unique valeur d'adhérence.
\end{rmq}
\begin{prop}[HP]
	Réciproquement, tout suite réelle bornée admettant une unique valeur d'adhérence est convergente, et converge alors vers cette valeur d'adhérence.
\end{prop}
\begin{proof}
	Il faut raisonner par l'absurde. Notons $l$ l'unique valeur d'adhérence On peut alors fixer un $\varepsilon>0$ tel qu'une infinité de termes de la suite soient au moins $\varepsilon$-loin de $l$. On construit une extractrice avec ces points. Or, la suite obtenue avec cette extractrice reste bornée, donc admet une valeur d'adhérence $l'$ par le théorème de Bolzano-Weierstrass. Or, par PLIL, $l\ne l'$. Donc la suite ne possède pas une unique valeur d'adhérence. Absurde. Donc la suite converge vers l'unique valeur d'adhérence.
\end{proof}
\begin{rmq}[Compacts, HP, spé]
	Ce théorème est plus généralement valable dans tout \textbf{espace métrique compact} ($ie$ un espace muni d'une notion de convergence et tel que toute suite à valeurs dans cet espace admette au moins une valeur d'adhérence). La démonstration est rigoureusement la même que celle présentée ci-dessus.
\end{rmq}
\section{Quelques compléments}
\subsection{Suites complexes}
\begin{dfn}
	Une \textbf{suite à valeurs complexes} est une application $(u_n)_{n\N}$ de $\N$ dans $\C$. La \textbf{partie réelle de $u$} est la suite réelle $(\text{Re}(u_n))_{n\in\N}$. La \textbf{partie imaginaire de $u$} est la suite réelle $(\text{Im}(u_n))_{n\in\N}$. La \textbf{suite conjuguée de $u$} est la suite complexe $(\overline{u_n})_{n\in\N}$. La \textbf{suite des modules de $u$} est la suite réelle $(\lvert u_n\rvert)_{n\in\N}$. 
\end{dfn}
\begin{dfn}
	Une suite complexe est dite \textbf{bornée} lorsque la suite de ses modules est majorée. Si la suite est réelle, on retrouve bien l'ancienne définition.
\end{dfn}
\begin{exa}
	La suite définie par $u_n=\exp(in\pi/6)$ est bornée.
\end{exa}
\begin{dfn}[Convergence]
	La suite $(u_n)_{n\in\N}$ est \textbf{convergente de limite $l\in\C$} lorsque $$\forall\varepsilon>0,\ \exists N\in\N,\ \forall n\in\N,\ n\geqslant N\implies \lvert u_n-l\rvert\leqslant\varepsilon$$
\end{dfn}
Cette fois, il ne s'agit plus d'une valeur absolue mais d'un module. Cela dit, puisque le module est un prolongement de la valeur absolue à $\C$, la définition ci-dessus coïncide bien dans le cas des suites réelles.
\begin{rmq}
	On montre comme dans $\R$ que si la limite existe, alors elle est unique. On raisonne par l'absurde, et à la place des valeurs absolues, on utilise des modules.
\end{rmq}
\begin{rmq}
	Les principes de troncature et de décalage d'indice fonctionnent encore.
\end{rmq}
\begin{rmq}
	On a $u_n\xrightarrow[n\to+\infty]{\C} l$ ssi $u_n-l\xrightarrow[n\to+\infty]{\C} 0$ ssi $\lvert u_n-l\rvert\xrightarrow[n\to+\infty]{\R} 0$. En particulier, $u_n\xrightarrow[n\to+\infty]{\C} 0$ ssi $\lvert u_n\rvert\xrightarrow[n\to+\infty]{\C} 0$.
\end{rmq}
\begin{exa}
	Soit $(u_n)_{n\in\N}$ une suite complexe et $(v_n)_{n\in\N}$ une suite réelle. Supposons que $\forall n\in\N,\ \lvert u_n\rvert\leqslant v_n$. Si $v_n\xrightarrow[n\to+\infty]{\R}0$, alors $u_n\xrightarrow[n\to+\infty]{\C}0$.
\end{exa}
\begin{thm}[Théorème-passerelle]
	Soit $(u_n)_{n\in\N}$ une suite complexe. On a : $$u_n\xrightarrow[n\to+\infty]{\C}l=a+ib\in\C \iff \left \{\begin{array}{ll}\text{Re}(u_n)\xrightarrow[n\to+\infty]{\R}a \\ \text{Im}(u_n)\xrightarrow[n\to+\infty]{\R}b\end{array}\right.$$
\end{thm}
On vérifie alors que de nombreux résultats vus dans le cas réels restent vrais. Les voici dans l'ordre.
\begin{prop}
	Toute suite complexe convergente est bornée.
\end{prop}
\begin{exa}
	Si $\lvert u_n\rvert \xrightarrow[n\to+\infty]{\R}+\infty$, alors $(u_n)_{n\in\N}$ diverge dans $\C$.
\end{exa}
\begin{thm}[Opérations algébriques]
	Les opérations algébriques usuelles sur les limites sont valables : combinaison linéaire, produit, passage à l'inverse (pour une limite non nulle), module, conjugaison.
\end{thm}
\begin{prop}
	Si une suite complexe converge vers $l\in\C$, alors toute suite extraite converge vers $l$. Réciproquement, si $u_{2n}\xrightarrow[n\to+\infty]{}l$ et $u_{2n+1}\xrightarrow[n\to+\infty]{}l$, alors $u_{n}\xrightarrow[n\to+\infty]{}l$.
\end{prop}
\begin{thm}[Théorème de Bolzano-Weierstrass]
	Le théorème de Bolzano-Weierstrass reste vrai : toute suite complexe bornée admet une suite extraite convergente.
\end{thm}
\subsection{Moyenne de Cesàro, HP mais à connaître par cœur}
Cette partie n'est pas $stricto$ $sensu$ au programme; mais elle est tellement classique que nous la donnons comme résultat de cours. \\ \par Soit une suite complexe $(u_n)_{n\geqslant1}$. Au lieu de s'intéresser à la convergence de $u$, on va s'intéresser à la convergence des moyennes arithmétiques successives : $$\forall n\in\N\st,\ \mu_n=\frac{u_1+\ldots+u_n}{n}$$ Intuitivement, on se dit que si $u_n$ tend vers $l$, alors à la longue, les moyennes $\mu_n$ vont "gommer" les écarts initiaux de $u_n$ autour de $l$, et "lisser les irrégularités", de sorte que $\mu_n$ tendra aussi vers $l$. \\ \par Formalisons maintenant cela. Tout d'abord, on commence par démontrer un cas particulier. C'est l'objet du
\begin{lem}
	Si $u_n\xrightarrow[n\to+\infty]{}0$, alors $\mu_n\xrightarrow[n\to+\infty]{}0$.
\end{lem}
\begin{proof}
	Donnons-nous $\varepsilon>0$, et soit $N$ un rang à partir duquel on ait $\lvert u_n\rvert\leqslant\displaystyle\frac{\varepsilon}{2}$. Posons ensuite $M=\lvert u_1+\ldots+u_N\rvert$, puis donnons-nous un rang $N'$ à partir duquel on ait $\displaystyle\frac{M}{n}\leqslant\frac{\varepsilon}{2}$. Alors, pour $n\geqslant\max(N,N')$, on a
	\begin{align*}
		\lvert\mu_n\rvert&\leqslant\frac{M+\lvert u_{N+1}+\ldots+u_n}{n}\\
		&\leqslant\frac{M}{n}+\frac{n-N}{n}\frac{\varepsilon}{2}\\
		&\leqslant \frac\varepsilon2+1\times\frac\varepsilon2 \\
		&\leqslant \varepsilon
	\end{align*}
	La preuve est achevée.
\end{proof}
\begin{thm}
	Si $u_n\xrightarrow[n\to+\infty]{}l$, alors $\mu_n\xrightarrow[n\to+\infty]{}l$.
\end{thm}
\begin{proof}
	Il suffit d'appliquer le lemme précédent à $u_n-l$ qui tend alors vers $0$. On obtient que $$\frac{(u_1-l)+\ldots+(u_n-l)}{n}=\mu_n-l$$ tend vers $0$. On conclut par somme de limites.
\end{proof}
\begin{rmq}
	En revanche, la réciproque est fausse. Comme contre-exemple, prendre par exemple la suite alternée définie par $u_n=(-1)^n$. On a $\mu_n\rightarrow0$ comme on s'en convainc en extrayant les sous-suites de rangs pairs et impairs. Pourtant, on n'a pas $u_n\rightarrow0$, puisque par exemple $u_{2n}\rightarrow1$.
\end{rmq}
A partir de maintenant, on suppose que $(u_n)$ est réelle. On peut se demander si le résultat reste vrai dans le cas où la limite vaut $\pm\infty$. La réponse est oui, comme le montre le théorème suivant.
\begin{thm}
	Si $u_n\xrightarrow[n\to+\infty]{}+\infty$, alors $\mu_n\xrightarrow[n\to+\infty]{}+\infty$. De même avec $-\infty$.
\end{thm}
\begin{proof}
	Quitte à changer $u$ en $-u$, démontrons sans perte de généralité le résultat sur $+\infty$. Soit $M\in\R_+$. Au-delà d'un certain rang $N$ on aura $u_n\geqslant2M+2$. Posons alors $S=u_1+\ldots+u_N$, puis donnons-un rang $N'$ à partir duquel on ait $\displaystyle\frac{S}{n}\geqslant-1$. Enfin, donnons-nous un rang $N''$ à partir duquel on ait $\displaystyle 1-\frac{N}{n}\geqslant\frac12$. Pour $n\geqslant\max(N,N',N'')$, on peut alors écrire :
	\begin{align*}
		\mu_n&\geqslant \frac{S}{n}+\frac1n\sum_{k=N+1}^{n}(2M+2)\\
		&\geqslant -1+\frac{n-N}{n}\times(2M+2) \\
		&\geqslant -1+\frac{1}{2}\times(2M+2)\\
		&\geqslant M
	\end{align*} puisque $2M+2\geqslant0$. La preuve est achevée.
\end{proof}
\begin{rmq}
	Là aussi, on pourra montrer que la réciproque est fausse.
\end{rmq}
\subsection{Suites particulières}
\subsubsection{Suites arithmético-géométriques}
On sait calculer le terme général d'une suite arithmétique, d'une suite géométrique, mais qu'en est-il d'un "mélange des deux" ? \\ \par On considère la suite définie par $$\left \{\begin{array}{ll}u_0=x\\ \forall n\in\N,\ u_{n+1}=au_n+b\end{array}\right.$$ où $x,a$ et $b$ sont des complexes. On impose $a\ne1$, car sinon on a affaire à ue suite arithmétique que l'on connaît déjà. On cherche alors une solution constante $q$ de la relation de récurrence. Or $$\forall q\in\C,\ q=aq+b\iff q=\frac{b}{1-a}$$ On parle de \textbf{point fixe}. Ensuite, on considère la suite auxiliaire $v$ de terme général $v_n=u_n-q$. Elle vérifie $v_{n+1}=av_n$, si bien qu'elle est géométrique. On sait alors calculer son terme général, puis on se ramène à celui de $u$ en ajoutant $q$.
\subsubsection{Suites récurrentes linéaires d'ordre 2}
\begin{thm}[Cas complexe]
	Soit $(a,b)\in\C\times\C\st$, et cherchons l'ensemble des suites complexes qui vérifient la relation de récurrence double $$\forall n\in\N,\ u_{n+2}+au_{n+1}+bu_n=0$$ Tout d'abord, on cherche dans $\C$ les racines de ce qu'on appelle le "polynôme caractéristique" : $X^2+aX+b$.
	\begin{itemize}
		\item Si ce polynôme possède deux racines distinctes $r_1\ne r_2$, les suites cherchées sont exactement celles de la forme $$u_n=\lambda r_1^n+\mu r_2^n$$ avec $(\lambda,\mu)\in\C^2$.
		\item Si ce polynôme possède une racine double $r$, les suites cherchées sont exactement celles de la forme $$u_n=(\lambda+\mu n)r^n$$ avec $(\lambda,\mu)\in\C^2$.
	\end{itemize}
\end{thm}
\begin{thm}[Cas réel]
	Soit $(a,b)\in\R\times\R\st$, et cherchons l'ensemble des suites complexes qui vérifient la relation de récurrence double $$\forall n\in\N,\ u_{n+2}+au_{n+1}+bu_n=0$$ Tout d'abord, on cherche dans $\C$ les racines de ce qu'on appelle le "polynôme caractéristique" : $X^2+aX+b$.
	\begin{itemize}
		\item Si ce polynôme possède deux racines réelles distinctes $r_1\ne r_2$, les suites cherchées sont exactement celles de la forme $$u_n=\lambda r_1^n+\mu r_2^n$$ avec $(\lambda,\mu)\in\R^2$.
		\item Si ce polynôme possède une racine double $r$, les suites cherchées sont exactement celles de la forme $$u_n=(\lambda+\mu n)r^n$$ avec $(\lambda,\mu)\in\R^2$.
		\item Si ne possède pas de racine réelle, on écrit les deux racines complexes conjuguées sous la forme $\rho\exp(\pm i\theta)$, puis les suites cherchées sont exactement celles de la forme $$u_n=\rho^n\left(\lambda\cos(n\theta)+\mu\sin(n\theta)\right)$$ avec $(\lambda,\mu)\in\R^2$.
	\end{itemize}
\end{thm}
\begin{rmq}
	A chaque fois, les constantes $\lambda$ et $\mu$ pourront se déterminer à partir des premiers termes, sortes de "conditions initiales" (ksssss).
\end{rmq}
\begin{exa}
	En appliquant ceci à la suite de Fibonacci définie par $$\left \{\begin{array}{lll}F_0=0 \\ F_1=1 \\ \forall n\in\N,\ F_{n+2}=F_{n+1}+F_n\end{array}\right.$$ on obtient $$\forall n\in\N,\ F_n=\frac{1}{\sqrt{5}}\left[\left(\frac{1+\sqrt{5}}{2}\right)^n-\left(\frac{1-\sqrt{5}}{2}\right)^n\right]$$
\end{exa}
\subsubsection{Suites définies par une fonction}
Considérons une fonction $f$ \textbf{continue} à valeurs réelles et un intervalle $I\subset\R$. On rappelle que $I$ est dit \textbf{stable} par $f$ lorsque $f(I)\subset I$. Dans ce cas, on peut définir une suite de la manière suivante : $$\left \{\begin{array}{ll}u_0\in I\\ \forall n\in\N, u_{n+1}=f(u_n)\end{array}\right.$$
\begin{prop}
	Voici la méthode générale pour étudier une suite du type $u_{n+1}=f(u_n)$, avec $f$ continue. \begin{itemize}
		\item On cherche un intervalle $I$ stable par $f$ contenant $u_0$.
		\item On étudie les points fixes de $f$.
		\item On étudie la monotonie éventuelle de $(u_n)_{n_in\N}$.
		\item On étudie la limite éventuelle de $(u_n)_{n\in\N}$
	\end{itemize}
\end{prop}
Chacun des points mérite d'être commenté séparément. \\ 
\paragraph{Recherche d'un intervalle stable} D'une manière générale, on aura toujours intérêt à choisir un intervalle le plus petit possible autour de $u_0$, ainsi l'étude de la suite sera plus rapide. Si $f$ n'est pas continue sur l'ensemble de son domaine de définition, il faudra choisir $I$ de telle sorte que $f_{|I}$ soit continue. \\ Dans la recherche de $I$, on aura bien souvent intérêt à étudier les variations de $f$ et à tracer l'allure de sa courbe représentative. \\
\paragraph{Recherche des points fixes}  Pour commencer, on donne le
\begin{thm}
	Si $u_n\rightarrow l$ et si $f$ est continue en $l$, alors $f(l)=l$.
\end{thm}
Ce théorème est souvent utilisé dans le cadre d'un raisonnement par analyse-synthèse. Si $u$ admet une limite $l$ en laquelle $f$ est définie, $l$ sera nécessairement un point fixe de $f$ par continuité. Notons qu'il est inutile de rechercher les points fixes "trop loin de $I$". En pratique, notre théorème de passage à la limite dans les inégalités nous montre que si $u$ converge vers $l$, alors $l$ appartient à $I$ ou est une borne de $I$ (en effet, les inégalités d'encadrement deviennent larges en passant à la limite). Donc, \textbf{la plupart du temps, il vaudra mieux avoir défini $f$ sur un intervalle fermé}. \\ En tout cas, on comprend l'intérêt de choisir $x$ petit :  cela permet de limiter le nombre de points fixes éventuels. En pratique, on étudiera souvent $g=f-\id$. \\
\paragraph{Étude de la monotonie} \begin{itemize}
	\item Notons que $\forall n\in\N,u_n{n+1}-u_n=g(u_n)$. Ainsi, on pourra se ramener à étudier le signe de $g$. bien souvent, l'étude précédente sur les zéros de $g$ fournira en même temps le signe sur le reste du domaine à l'aide d'un tableau de variations.
	\item Parfois, on peut gagner du temps en remarquant que $f$ est croissante ou décroissante. Dans ce cas, il suffit de comparer $u_0$ et $u_1$, puis une récurrence immédiate montre que $u$ est monotone.
\end{itemize} 
\paragraph{Étude de la limite de $u$}  Une fois qu'on a trouvé une limite éventuelle $l$, pour montrer que $u_n\rightarrow l$, il existe différente méthodes. Le plus souvent, il s'agira d'utiliser le théorème de la limite monotone. Pour cela, il faudra que l'étude de la monotonie de $u$ ait été concluante. On pourra aussi utiliser l'inégalité des accroissements finis.
\begin{exa}
	En appliquant cette méthode à la suite définie par $$\left \{\begin{array}{ll}u_0\in[0,1] \\ \forall n\in\N,\ u_{n+1}=\sin(u_n)\end{array}\right.$$ on montre que $u_n\rightarrow 0$.
\end{exa}
\section{Compléments : construction de $\R$ et complétude, HP}
\subsection{Construction de $\mathbb{R}$}

Nous nous proposons ici de construire $\mathbb{R}$ comme une extension de $\mathbb{Q}$, à l'aide des coupures de Dedekind, du nom de Richard Dedekind, un très grand mathématicien allemand du XIX$^e$ siècle.

\subsubsection{Construction de l'ensemble des réels}

\begin{dfn}
	Une \textbf{section commençante ouverte} (ou SCO) de $\mathbb{Q}$ est une partie $A \subset \mathbb{Q}$ telle que
\begin{itemize}
	\item $A \neq \emptyset$ et $A \neq \mathbb{Q}$
	\item $\forall a \in A, \forall a' \in \mathbb{Q}, a' < a \Rightarrow a' \in A$
	\item $A$ ne possède pas de maximum.
\end{itemize}

L'ensemble des SCO est noté $\mathbb{R}$, et à partir de maintenant ses éléments seront plutôt appelés les \textbf{réels}.
\end{dfn}
\begin{rmq}
	À la fin de cette partie, il s'avèrera en fait que les SCO sont les ensembles de la forme $]-\infty, A[ \cap \mathbb{Q}$ avec $A \in \mathbb{R}$. Cela peut aider à les visualiser.
\end{rmq}
\begin{rmq}
	Pour éviter au maximum les confusions, on réservera les majuscules aux éléments de $\mathbb{R}$ et les minuscules aux éléments de $\mathbb{Q}$.
\end{rmq}
\begin{prop}
	L'application $r \mapsto A_r = \{r' \in \mathbb{Q} \mid r' < r\}$ va de $\mathbb{Q}$ dans $\mathbb{R}$ et elle est injective. Ainsi, elle permet d'identifier $\mathbb{Q}$ à un sous-ensemble de $\mathbb{R}$.
\end{prop}
\begin{proof}
	Soit $r \in \mathbb{Q}$. Montrons déjà que $A_r \in \mathbb{R}$:
\begin{itemize}
	\item $r-1 \in A_r$ donc $A_r \neq \emptyset$ et $r \notin A_r$ donc $A_r \neq \mathbb{Q}$
	\item Soit $s \in A_r$ et $s' < s$. Alors $s' < r$ donc $s' \in A_r$
	\item Supposons que $A_r$ admette un plus grand élément $s$. Alors on aurait $s < r$, puis $\frac{r+s}{2} \in A_r$ avec $\frac{r+s}{2} > s$, contradiction.
\end{itemize}
Soit maintenant $r' \in \mathbb{Q}$ tel que $A_r = A_{r'}$, et montrons $r = r'$. Supposons par l'absurde que $r < r'$ alors on aurait $r \in A_{r'} = A_r$, contradiction. Donc $r \geqslant r'$. En échangeant le rôle de $r$ et $r'$, on obtient $r' \geqslant r$. Donc $r = r'$.
\end{proof}
\begin{dfn}
	On définit une relation sur $\mathbb{R}$ par $A \leqslant B \Leftrightarrow A \subset B$.
\end{dfn}

\begin{prop}
	$\leqslant$ est une relation d'ordre totale qui prolonge $\leqslant$ définie sur $\mathbb{Q}$.
\end{prop}
\begin{proof}
	On sait déjà que $\leqslant$ est une relation d'ordre (partielle) sur les parties de n'importe quel ensemble. Ici, il faut vérifier qu'elle est totale. Soit $(A,B) \in \mathbb{R}^2$ et supposons qu'on n'a pas $A \subset B$. Montrons que $B \subset A$. Pour cela, prenons $a \in A \backslash B$. Soit $b \in B$ quelconque, le but est de montrer que $b \in A$. Or on ne peut pas avoir $b > a$ sans quoi on aurait $a \in B$. On en déduit $b \leqslant a$. Ensuite, si $b = a$ alors c'est terminé. Si $b < a$, alors on a aussi $b \in A$. \par
	Pour finir, soit $(r,r') \in \mathbb{Q}^2$ et montrons que $r \leqslant r' \Leftrightarrow A_r \subset A_{r'}$. Le sens direct est trivial. Réciproquement, si $A_r \subset A_{r'}$, supposons $r > r'$, alors on a $r' \in A_r$ donc $r' \in A_{r'}$, absurde. Donc $r \leqslant r'$.
\end{proof}
Voici à présent un petit lemme dont nous ferons régulièrement usage.
\begin{lem}
	Soit $A > 0$, alors $A$ contient un rationnel strictement positif.
\end{lem}
\begin{proof}
	On n'a pas $A \leqslant 0$, c'est-à-dire que $A$ n'est pas inclus dans $A_0$. Autrement dit, $A$ contient un rationnel $r \geqslant 0$. Si $r = 0$, il suffit de remarquer que par hypothèse, $A$ ne possède pas de maximum.
\end{proof}

\subsubsection{Construction de l'addition}
\begin{prop}
	L'opération $A + B = \{a + b \mid (a, b) \in A \times B\}$ est bien définie, et elle prolonge l'addition sur $\mathbb{Q}$.
\end{prop}
\begin{proof}
	Soit $(A, B) \in \mathbb{R}^2$. Commençons par vérifier que $A + B$ est bien une SCO.
	\begin{itemize}
		\item $A \neq \emptyset$ donc il existe $a \in A$. De même il existe $b \in B$. Finalement, $A + B \neq \emptyset$. De plus, soit $x \in \mathbb{Q} \backslash A$ et $y \in \mathbb{Q} \backslash B$. Quel que soit $a \in A$, on ne peut pas avoir $a > x$ sans quoi on aurait $x \in A$. De même on ne peut pas avoir $a = x$, donc on a $a < x$. De même, tout $b \in B$ vérifie $b < y$. Finalement, tout $a + b \in A + B$ vérifie $a + b < x + y$, si bien que
		$A + B \neq \mathbb{Q}$		
		\item Prenons $a + b \in A + B$ et $c < a + b$. En notant $\varepsilon = a + b - c > 0$, on a $$c = \underbrace{\left(a - \frac{\varepsilon}{2}\right)}_{\in A} + \underbrace{\left(b - \frac{\varepsilon}{2}\right)}_{\in B} \in A + B$$
		\item Enfin, montrons que $A + B$ ne possède pas de maximum; si on se donne $a + b \in A + B$ alors il suffit de choisir $a' > a$ dans $A$, $b' > b$ dans $B$, et de considérer $a' + b'$.
	\end{itemize}
	Maintenant, soit $(r, r') \in \mathbb{Q}^2$ et montrons que $A_r + A_{r'} = A_{r + r'}$.
	\begin{itemize}
		\item D'une part, si on se donne $r' + s' \in A_r + A_{s'}$ alors $r' < r$ et $s' < s$ donc $r' + s' < r + s$.
		\item Réciproquement, si on se donne $t \in A_{r + s}$, alors en notant $\varepsilon = r + s - t > 0$ on écrit de la même manière que tout à l'heure $$t = \underbrace{\left(r - \frac{\varepsilon}{2}\right)}_{\in A_r} + \underbrace{\left(s - \frac{\varepsilon}{2}\right)}_{\in A_{s}} \in A_r + A_{s'}$$
	\end{itemize}
	Ainsi la preuve est achevée.
\end{proof}
Avant de continuer, nous avons besoin de démontrer le
\begin{lem}
	Soit $\frac{p}{q} \in \mathbb{Q}$. Alors il existe $n \in \mathbb{N}$ tel que $n >\frac{p}{q}$.
\end{lem}
\begin{proof}
	Supposons sans perte de généralité que $q \in \mathbb{N}^*$. Si $p \leqslant 0$ alors $n = 1$ convient, et sinon $n = 2p$ convient. En effet, on a $q \geqslant 1$ d'où on tire $$2p > p \geqslant \frac{p}{q}$$ ce qui achève la preuve.
\end{proof}
\begin{thm}
	$(\mathbb{R},+)$ est un groupe commutatif, et l'addition est compatible avec $\leqslant$ au même sens que $\mathbb{Q}$.
\end{thm}
\begin{proof}
	Vérifions les différents axiomes d'un groupe commutatif.
	\begin{itemize}
		\item Associativité et commutativité: elles s'héritent directement de $\mathbb{Q}$.
		\item Élément neutre: c'est $0:=A_{0}$
		\item Opposé: soit $A\in\mathbb{R}$
		\begin{itemize}
			\item[--]Si A est de la forme $A_{r}$ alors on sait que $A_{r}+A_{-r}=A_{0}=0$.
			\item[--]Si A ne peut pas s'écrire sous la forme $A_{r}$ montrons que $B=\{-b^{\prime}|b^{\prime}\notin A\}$ convient.
			Pour commencer, montrons que $B\in\mathbb{R}$.
			\begin{itemize}
				\item[*]Comme $A\ne\emptyset$ on peut choisir $a\in A$ et on aura $-a\notin B$ Donc $B\ne\mathbb{Q}$. \\
				Comme $A\ne\mathbb{Q}$ on peut choisir $b^{\prime}\in\mathbb{Q}\backslash A$ et on aura $-y\in B$. Donc $B\ne\emptyset$.
				\item[*]Soit $b\in B c<b$ Ecrivons $b=-b^{\prime}$ avec $b^{\prime}\notin A c=-c^{\prime}$. Alors $c^{\prime}>b^{\prime}$ donc $c^{\prime}\notin A$ sans quoi on aurait $b^{\prime}\in A$ Donc $c\in B$.
				\item[*] Soit $b\in B$ et écrivons $b=-b^{\prime}$. Par hypothèse on a $A\ne A_{b^{\prime}}$. Supposons $A>A_{b^{\prime}}$ c'est-à-dire $A\supseteq A_{b^{\prime}}$ alors on pourrait trouver un $a\in A\backslash A_{b'}$ On aurait donc $b^{\prime}\le a$ d'où on déduirait $b^{\prime}\in A$ dans tous les cas, contradiction. Ainsi on a $A<A_{b'}$ Autrement dit, il existe un $c^{\prime}<b^{\prime}$ avec $c^{\prime}\notin A$. On a alors $-c^{\prime}>-b^{\prime}$ avec $-c^{\prime}\in B$. Donc B ne possède pas de maximum.
			\end{itemize}
			Ensuite, montrons que $A+B=A_{0}$.
			\begin{itemize}
				\item[*]Soit $c\in A+B$, écrivons $c=a-b^{\prime}$ avec $a\in A$ et $b^{\prime}\notin A$. On a donc $b'>a$ (sans quoi on aurait $b^{\prime}\in A$ dans tous les cas), d'où on déduit $c\in A_{0}$.
				\item[*]Réciproquement soit $c<0$, et supposons par l'absurde que $\forall a\in A,\ a-c\in A$. Si on choisit $a_{0}\in A$ fixé, une récurrence immédiate montre qu'on aurait $a_{0}-nc\in A$ pour tout $n\in\mathbb{N}$. Or si on prend $x\in\mathbb{Q}\backslash A$, on peut trouver $n\in\N$ tel que $\displaystyle n>\frac{a_{0}-x}{c}$ par le lemme. On en déduirait $nc<a_{0}-x$ d'où $x<a_{0}-nc$ et finalement $x\in A$, contradiction. Donc on peut trouver $a\in A$ vérifiant $b^{\prime}=a-c\notin A$. Il suffit alors de poser $b=-b^{\prime}\in B$ et de remarquer que $c=a+b$.
			\end{itemize}
		\end{itemize}
	\end{itemize}
	Pour finir, soit $(A,B,C)\in\mathbb{R}^{3}$. Il faut vérifier la compatibilité de $\le$ avec $+$ au sens suivant :$$A\le B\Rightarrow A+C\le B+C$$ En fait, c'est à peu près immédiat. Si on se donne $a+c\in A+C$, alors $a$ $fortiori$ on $a\in B$, d'où le résultat.
\end{proof}
\subsubsection{Construction de la multiplication}
Nous en venons maintenant à la définition de la multiplication, qui est un peu plus compliquée que celle de l'addition.
\begin{prop}
	L'opération \(A\cdot B=\{c\in\mathbb{Q}\ |\ \exists(a,b)\in A\times B, a\ge0,\ b\ge0 ,c\le ab\}\) est bien définie sur \(\mathbb{R}_{+}^{*}\) et elle prolonge la multiplication sur \(\mathbb{Q}_{+}^{*}\).
\end{prop}
\begin{proof}
	Soit \(A,B\in\mathbb{Q}_{+}^{*}\) et commençons par vérifier que \(A\cdot B\in\mathbb{R}_{+}^{*}\)
	\begin{itemize}
		\item Comme on n'a pas $A \subset A_0$, c'est qu'on peut trouver $a \in A$ avec $a \geqslant 0$. De même, on peut trouver $b \in B$ avec $b \geqslant 0$. Alors $ab \in A \cdot B$ donc $A \cdot B \neq \emptyset$. De plus, considérons $x \in \mathbb{Q} \backslash A$ et $y \in \mathbb{Q} \backslash B$. Rappelle que $a \in A$ vérifie $a < x$ et tout $b \in B$ vérifie $b < y$. Donc si $a, b \geqslant 0$ on a $xy > ay \geqslant ab$. Ainsi $xy \notin A \cdot B$ et $A \cdot B \neq \mathbb{Q}$.
		\item Soit $c \in A \cdot B$ et $c' < c$. Écrivons $c \leqslant ab$, avec $a \geqslant 0$ et $b \geqslant 0$. Alors $a$ $fortiori$ on a bien $c' \leqslant ab$ et $(a, b) \in A \times B$.
		\item Enfin, soit $c \in A \cdot B$ et écrivons $c \leqslant ab$ de la même manière. On peut trouver $a' > a$ dans $A$ et $b' > b$ dans $B$, si bien que $a', b' > 0$ puis $a'b' > a'b \geqslant ab$. Donc $a'b' > c$ avec $a'b' \in A \cdot B$.
		\item Le point précédent montre que $A \cdot B$ n'est pas inclus dans $A_0$, d'où on tire $A \cdot B > 0$.
	\end{itemize}
	Maintenant, soit $r, s \in \mathbb{Q}_{+}^{*}$ et montrons que $A_r \cdot A_s = A_{rs}$.
	\begin{itemize}
		\item D'une part, soit $t \in A_r \cdot A_s$. Alors on peut écrire $t \leqslant r's'$, avec $(r', s') \in A_r \times A_s$ et $r', s' \geqslant 0$. Comme précédemment, on en déduit $r's' \leqslant r's < rs$ donc $t \in A_{rs}$.
		\item Réciproquement, soit $t \in A_{rs}$, alors $t < rs$. Introduisons alors $\varepsilon = \min(r, s, \frac{rs-t}{r+s})$ de telle sorte que $0 < \varepsilon \leqslant r, s, \frac{rs-t}{r+s}$. On a alors $(r-\varepsilon)(s-\varepsilon) = rs - \varepsilon(r+s) + \varepsilon^2 > rs - \varepsilon(r+s) \geqslant t$, d'où $t \in A_{\varepsilon} \cdot A_s$.
	\end{itemize}
	Ainsi la preuve est achevée.
\end{proof}
\begin{dfn}
	On prolonge $A \cdot B$ à tout entier de la manière suivante :
	\begin{itemize}
	\item Si $A = 0$ ou $B = 0$, on pose $A \cdot B = 0$.
	\item Si $A < 0$ et $B > 0$, on pose $A \cdot B = -[(-A) \cdot B]$.
	\item Si $A > 0$ et $B < 0$, on pose $A \cdot B = -[A \cdot (-B)]$.
	\item Si $A < 0$ et $B < 0$, on pose $A \cdot B = (-A) \cdot (-B)$.
\end{itemize}

\end{dfn}
\begin{prop}
	Ce prolongement coïncide bien avec la multiplication sur tout entier.
\end{prop}
\begin{proof}
	C'est immédiat en remarquant que $\forall r \in \mathbb{Q}, -A_r = A_{-r}$. À titre d'exemple, soit $r < 0$ et $s > 0$, On a $A_r \cdot A_s = -[(-A_r) \cdot A_s] = -[A_{-r} \cdot A_s] = -[A_{(-r)s}] = -A_{-rs} = A_{rs}$ et ainsi de suite.
\end{proof}
Ce petit lemme nous aidera pour les vérifications sur l'associativité et la distributivité.
\begin{lem}
	On a $\forall (A,B) \in \mathbb{R}^2, (-A) \cdot B = A \cdot (-B) = -(A \cdot B)$.
\end{lem}
\begin{proof}
	Soit $(A,B) \in \mathbb{R}^2$. Si A ou B est nul, le résultat est immédiat. Sinon, on a quatre cas.
	\begin{itemize}
		\item Si $A > 0$ et $B > 0$, alors $(-A) \cdot B = -(A \cdot B)$ par définition, et de même avec le deuxième terme.
		\item Si $A < 0$ et $B < 0$, alors $(-A) \cdot B = -[(-A) \cdot (-B)] = -(A \cdot B)$, et de même avec le deuxième terme.
		\item Si $A > 0$ et $B < 0$ alors $(-A) \cdot B = A \cdot (-B)$ par définition, et $A \cdot B = -[A \cdot (-B)]$ d'où $-(A \cdot B) = A \cdot (-B)$.
		\item Si $A < 0$ et $B > 0$ la preuve est symétrique.
	\end{itemize} 
	Ainsi la preuve est achevée.
\end{proof}
Là aussi, nous avons besoin d'un autre lemme préliminaire, symétrique du lemme pour le groupe pour $+$.
\begin{lem}
	Soit $\alpha > 1$ dans $\mathbb{R}$ et $m \in \mathbb{Q}$. Alors on peut trouver $n \in \mathbb{N}$ tel que $\alpha^n > m$.
\end{lem}
\begin{proof}
	Pour $n \in \mathbb{N}^*$, on a $$\alpha^n - 1 = (\alpha - 1)(\alpha^{n-1} + ... + 1) \geqslant n(\alpha - 1)$$ d'où $\alpha^n \geqslant 1 + n(\alpha - 1)$. Il suffit ensuite de prendre $n \in \mathbb{N}^*$ tel que $\displaystyle n > \frac{m - 1}{\alpha - 1}$ à un lemme précédent
\end{proof}
\begin{thm}
	$(\mathbb{R}, +, \cdot)$ est un corps.
\end{thm}
\begin{proof}
	Nous avons déjà vérifié que c'était un groupe commutatif pour l'addition, et ce n'est pas l'anneau trivial car il contient $\mathbb{Q}$. Vérifions donc les axiomes sur la multiplication.
	\begin{itemize}
		\item \textbf{Associativité} : Soit $(A,B,C) \in \mathbb{R}^3$. Il faut montrer que $(A \cdot B) \cdot C = A \cdot (B \cdot C)$. Si l'un des trois réels est nul c'est immédiat. Sinon, sans perte de généralité on peut se ramener au cas où $A,B,C>0$, les autres cas s'en déduisant immédiatement à partir du lemme.
		\begin{itemize}
			\item[--]Soit $r \in (A \cdot B) \cdot C$, et écrivons $r \leqslant sc$ et $s \leqslant ab$ avec les conditions habituelles.
			On a alors $r \leqslant a(bc)$ avec $bc \in B \cdot C$ et $a \in A$, d'où $r \in A \cdot (B \cdot C)$.
			\item[--]L'inclusion réciproque se démontre de la même manière.
		\end{itemize}
		\item \textbf{Commutativité} : c'est une conséquence de la définition, qui est symétrique en A et B.
		\item \textbf{Élément neutre} : c'est $1 = A_1$. En effet, soit $A \in \mathbb{R}$ et montrons que $A \cdot A_1 = A$ (par commutativité, ce sera suffisant). Si $A = 0$ le résultat est évident, et si $A < 0$ on se ramène à $A > 0$ par le lemme 8.69. Supposons donc $A > 0$ et commençons par montrer $A \cdot A_1 \subset A$. Soit $r \in A \cdot A_1$ et écrivons $r \leqslant aa_1$ avec les conditions habituelles. Comme $a \geqslant 0$ et $a_1 < 1$ nous avons $r < a$ donc $r \in A$. Réciproquement, soit $a \in A$. Si on choisit $a' > a$ dans A (puisque A ne possède pas de maximum), alors $a'>0$ puis $\displaystyle\frac{a}{a'} \in A_1$ avec $\displaystyle\frac{a}{a'} \geqslant 0$. Il ne reste plus qu'à écrire $a = a' \cdot \frac{a}{a'} \in A \cdot A_1$.
		\item Passage à l'inverse: soit $A \neq 0$. Grâce au lemme précédent, supposons même $A > 0$.
		\begin{itemize}
			\item[--]Si $A$ est de la forme $A_r$ alors $r > 0$ puis on sait que $A_r \cdot A_{1/r} = A_1 = 1$.
			\item[--]Si $A$ ne peut pas s'écrire sous la forme $A_{r_i}$ montrons que $\displaystyle B = \mathbb{R}_- \cup \{\frac{1}{b'}\ |\ b' \notin A\}$	convient. Pour commencer, montrons que $B$ est bien définie et que $B \in \mathbb{R}$.
			\begin{itemize}
				\item[*]$A$ n'est pas égal à tout entier, donc il existe des $b' \in \mathbb{Q} \backslash A$. Par ailleurs $A$	contient un rationnel strictement positif par un lemme précédent, d'où on déduit que	tous les $b' \in \mathbb{Q} \backslash A$ sont strictement positifs. Donc $\displaystyle\{\frac{1}{b'} | b' \notin A\}$ est bien défini, non vide et inclus dans $\mathbb{Q}_+^*$. Conclusion: $B$ est bien défini et possède au moins un rationnel strictement positif.
				\item[*]Comme $\mathbb{R}_- \subset B$ on a $B \neq \emptyset$ par ailleurs on peut trouver $\alpha > 0$ dans $A$, d'où	on déduit $\frac{1}{a} \notin B$ puis $B \neq \mathbb{Q}$.
				\item[*]Soit $b \in B, c < b, c \leqslant 0$ alors $c \in B$ sinon c'est qu'on a $0 < c < b$. Écrivons $\displaystyle c = \frac{1}{c'}$ et $\displaystyle b = \frac{1}{b'}$ avec $0 < b' < c'$.
				$b \notin A$. On en	déduit $c' \notin A$, d'où $c \in B$.
				\item[*] Soit $b \in B$ et trouvons un $c > b$ dans $B$. Si $b \leqslant 0$ nous avons vu que $B$ contenait un rationnel strictement positif. Sinon, écrivons $\displaystyle b =\frac{1}{b'}$ avec $b' \notin A$ et $b' > 0$.	Par hypothèse on a $A \neq A_{b'}$. Supposons $A > A_{b'}$ -à-dire $A \supseteq A_{b'}$. Alors	on pourrait trouver un $a \in A \backslash A_{b'}$ puis on aurait $a \geqslant b'$. On en déduirait donc $b' \in A$ dans tous les cas, contradiction. Ainsi on a $A <A_{b'}$. Autrement dit, il existe un $c' < b'$ avec $c' \notin A$. On en déduit $c' > 0$ puis $\displaystyle \frac{1}{c'} > \frac{1}{b'}$. Il suffit alors de poser $\displaystyle c := \frac{1}{c'}$.
			\end{itemize}
			Montrons maintenant que $A \cdot B = A_1$
			\begin{itemize}
				\item[*] Soit $c \in A \cdot B$ écrivons $c \leqslant ab$ avec les conditions habituelles et montrons que $ab < 1$. Si $b = 0$ le résultat est immédiat, sinon écrivons $b = \frac{1}{b'}$ avec $b' > 0$ et $b' \notin A$. On a nécessairement $a < b'$, d'où le résultat.
				\item[*]Réciproquement soit $c \in A_1$, et montrons que $c \in A \cdot B$. Si $c \leqslant 0$ il suffit de remarquer que $0 \in A \cap B$, donc considérons à de maintenant que $c > 0$. Supposons que pour tout $a > 0$ dans $A$, on ait $\frac{a}{c} \in A$. Si nous choisissons $a_0 > 0$ fixé dans $A$ (ce qui est possible par un lemme précédent), une récurrence immédiate montre qu'on aurait $\displaystyle\frac{a_0}{c^n} \in A$ pour tout $n \in \mathbb{N}$. Or si on prend $x \in \mathbb{Q} \backslash A$, on peut trouver $n \in \mathbb{N}$ tel que $\displaystyle (\frac{1}{c})^n > \frac{x}{a_0}$ un lemme. On en déduirait $x < \frac{a_0}{c^n}$ et finalement $x \in A$ contradiction.	Donc il existe $a > 0$ dans $A$ tel que $\displaystyle b' = \frac{a}{c} \notin A$. Si on pose $\displaystyle b= \frac{1}{b'} \in B$, on obtient alors $c = ab \in A \cdot B$.
			\end{itemize}
		\end{itemize}
		\item \textbf{Distributivité} : Par commutativité il suffit de la vérifier d'un côté seulement. Soit $(A,B,C) \in \mathbb{R}^3$ et montrons par exemple que $A \cdot (B+C) = A \cdot B + A \cdot C$. Là aussi, si $A$, $B$ ou $C$ est nul le résultat est immédiat; sinon, quitte à changer $A$ en $-A$ et $(B,C)$ en $(-B,-C)$, un lemme montre qu'on peut supposer $A, B > 0$. Commençons par supposer aussi $C > 0$ et montrons le résultat.
		\begin{itemize}
			\item[--] Soit $r \in A \cdot (B+C)$, et écrivons $r \leqslant a(b+c)$ avec $a \geqslant 0$ et $b+c \geqslant 0$. On a $B > 0$ donc $B$ contient un rationnel strictement positif. On en déduit $b_{2} = max(0,b) \in B$ Or $ab \leqslant ab_{2}$ d'où $ab \in A \cdot B$. De même on montre $ac \in A \cdot C$. On a donc $a(b+c) \in A \cdot B + A \cdot C$ et de même pour $r$.
			\item[--]Réciproquement, soit $r \in A \cdot B + A \cdot C$ et écrivons $r = r_{1} + r_{2}$ avec $r_{1} \leqslant a_{1}b$, $r_{2} \leqslant a_{2}c$ et les conditions habituelles Si on pose $a = max(a_{1},a_{2}) \in A$ on a $r \leqslant a(b+c)$ (et bien sûr et $a \geqslant 0$ et $b+c \geqslant 0$) d'où $r \in A \cdot (B+C)$.		
		\end{itemize}
		Maintenant, si $C < 0$ montrons plutôt $A \cdot (B-C) = A \cdot B - A \cdot C$ avec  $C > 0$.
		\begin{itemize}
			\item[--]Si $C < B$ on peut appliquer le point précédent, et on obtient
			\begin{align*}
				A \cdot (B-C) + A \cdot C &= A \cdot B + A \cdot (-C) + A \cdot C \\
				&= A \cdot B - A \cdot C + A \cdot C \\
				&= A \cdot B
			\end{align*}
			\item[--]Si $C = B$ le résultat est immédiat.
			\item[--]Si $C > B$ on applique un lemme pour se ramener au premier cas.
			\begin{align*}
				A \cdot (B-C) &= -[A \cdot (C-B)] \\
				&= -[A \cdot C - A \cdot B] \\
				&= A \cdot B - A \cdot C
			\end{align*}
		\end{itemize}
	\end{itemize}
	Ainsi la preuve est achevée.
\end{proof}

\begin{prop}
	On garde la compatibilité de $\leqslant$ avec la multiplication.
\end{prop}
\begin{proof}
	De manière précise, montrons que $$\forall(A,B,C)\in\mathbb{R}^{3} [A\le B \text{ et } C\ge0] \Rightarrow AC\le BC$$ Par distributivité et compatibilité avec l'addition, tout revient à se donner $A,B\ge0$ et à montrer que $A \cdot B\ge0$.
	\begin{itemize}
		\item Si $A$ ou $B$ est nul, le résultat est immédiat.
		\item Si $A,B>0$, montrons que $A \cdot B\supset A_{0}$. Soit donc $x<0$. Comme on peut trouver $a>0$ dans A et $b>0$ dans $B$, il vient immédiatement $x\le ab$.
	\end{itemize}
	Ainsi la preuve est achevée.
\end{proof}
\begin{thm}
	Soit $X$ une partie non vide et majorée de $\R$. Alors elle admet une borne supérieure.
\end{thm}
\begin{proof}
	Posons $\displaystyle M = \bigcup_{A\in X}A\subset\mathbb{Q}$, et montrons que M est la borne supérieure cherchée. Pour cette démonstration, il est particulièrement important de se rappeler que les minuscules
désignent des rationnels, donc des éléments de Q, et les majuscules des réels, donc des parties de $\R$.
\begin{itemize}
	\item Montrons déjà que $M \in \mathbb{R}$
	\begin{itemize}
		\item[--]$X$ est non vide, donc on peut choisir $A \in X$. Ensuite on a $M \supset A \neq \emptyset$, donc $M \neq \emptyset$. Par ailleurs, si $r$ est majorant de $X$, prenons $r \in \mathbb{Q} \backslash M_1$ et montrons que $r \notin M$. Ainsi aurons $M \neq \mathbb{Q}$. Soit donc $A \in X$ quelconque, et montrons que $r \notin A$. Comme $A \subset M_1$, si on avait $r \in A$ on aurait $r \in M_1$, contradiction.
		\item[--]Soit maintenant $m \in M$ et $m' < m$. Alors on peut trouver $A \in X$ tel que $m \in A$ puis on a $m' \in A$ d'où $m' \in M$.
		\item[--]Enfin, soit $m \in M$. Et trouvons $m' > m$ dans $M$. De même, on prend $A \in X$ tel que $m \in A$, puis comme $A$ ne possède de maximum, on peut trouver $m' > m$ dans $A$. On a alors $m' \in M$.
	\end{itemize}
	\item Montrons maintenant que $M = sup(X)$.
	\begin{itemize}
		\item[--]Par construction, $M$ contient tous les $A$, donc c'est un majorant de $X$.
		\item[--]Donnons-nous $M_2 < M$, et trouvons $A > M_2$ dans $X$. Comme $M$ contient strictement $M_2$ on peut trouver $m \in M \backslash M_2$. Soit $A \in X$ tel que $m \in A$. Puisque $m \notin M_2$, $A$ n'est pas inclus dans $M_2$. Et comme sur $\mathbb{R}$, $\leqslant$ est une relation d'ordre totale, c'est donc que $A > M_2$.
	\end{itemize}
\end{itemize}
Ainsi la preuve est achevée.
\end{proof}
\begin{cor}
	Soit $X$ une partie non vide et minorée de $\mathbb{R}$. Alors elle admet une borne inférieure.
\end{cor}
\begin{proof}
	Considérons $-X = \{-A | A \in X\}$. Cette partie est non vide et majorée, donc elle admet une borne supérieure qu'on note $m$. On en déduit que $m$ est un minorant de $X$. Puis si on se donne $m_1 > m$, alors on a $-m_1 < -m$, donc $-m_1$ n'est pas un majorant de $-X$. Ainsi on peut trouver $-A > -m_1$ avec $A \in X$. Autrement dit on a $A < m_1$, ce qui montre que $m_1$ n'est pas un minorant de $X$. Donc $m = inf(X)$.
\end{proof}

\subsection{Complétude de $\mathbb{R}$}

\begin{dfn}
	La suite $(u_n)_{n \in \mathbb{N}}$ est une \textbf{suite de Cauchy }lorsque "ses termes sont de plus en plus proches les uns des autres" : $$\forall \varepsilon > 0, \exists N \in \mathbb{N}, \forall m, n \geqslant N, |u_m - u_n| \leqslant \varepsilon$$
\end{dfn}
\begin{prop}
	Toute suite convergente est de Cauchy.
\end{prop}
\begin{proof}
	Soit $(u_n)_{n \in \mathbb{N}}$ une suite de limite $l$, et soit $\varepsilon > 0$. Alors il existe $N \in \mathbb{N}$ tel que pour $n \geqslant N$ on ait $\displaystyle|u_n - l| \leqslant \frac{\varepsilon}{2}$. Pour $m, n \geqslant N$ on aura bien $|u_m - u_n| \leqslant \varepsilon$ par l'inégalité triangulaire.
\end{proof}
Plus globalement, ce dernier résultat est vrai dans beaucoup d'espaces, dès que l'on peut y définir une distance. En revanche, la question de la réciproque est éminemment délicate dans le cas général. C'est même elle qui motive la construction de $\mathbb{R}$, puisqu'elle est fausse dans $\mathbb{Q}$ et vraie dans $\mathbb{R}$. \par Avant toute chose, démontrons le

\begin{lem}
	Si une suite de Cauchy possède une suite extraite convergente, alors elle est convergente.
\end{lem}
\begin{proof}
	Soit $(u_{n})_{n\in \mathbb{N}}$ une suite de Cauchy telle que $(u_{\varphi(n)})_{n\in\mathbb{N}}$ converge vers une limite $l$, et donnons-nous $\varepsilon>0$. Pour $m,n\ge N_{1}$ on a $|u_{m}-u_{n}|\le\frac{\varepsilon}{2}$; pour $n\ge N_{2}$ on a$\displaystyle|u_{\varphi(n)}-l|\le\frac{\varepsilon}{2}$. \\
De plus, rappelons qu'on a $\varphi(n)\ge n$. Donc pour $n\ge max(N_{1},N_{2})$ on a
$$|u_{n}-l|\le|u_{n}-u_{\varphi(n)}|+|u_{\varphi(n)}-l|\le\varepsilon_{i}$$
ce qui achève la preuve.
\end{proof}
\begin{thm}[Complétude de $\R$]
	Toute suite de Cauchy converge dans $\mathbb{R}$. On dit que $\mathbb{R}$ est \textbf{complet}.
\end{thm}
\begin{proof}
	Soit $(u_{n})_{n\in\mathbb{N}}$ une suite de Cauchy, et $N\in\mathbb{N}$ tel que $\forall m, n\ge N |u_{m}-u_{n}|\le1$. En particulier, pour $n\ge N$ a $$|u_{n}|\le|u_{n}-u_{N}|+|u_{N}|\le1+|u_{N}|$$ Donc l'ensemble des $\lvert u_n\rvert $ est majoré par $\max(|u_{0}|,...,|u_{N-1}|,1+|u_{N}|)$ et $(u_{n})_{n\in\mathbb{N}}$ est bornée.
D'après le théorème de Bolzano-Weierstrass extraite convergente, et on conclut par le lemme.
\end{proof}
\begin{thm}
	$\mathbb{Q}$ n'est pas complet.
\end{thm}
\begin{proof}
Avant toute chose, précisons ce que nous entendons par là. Une suite de
rationnels $(r_{n})_{n\in\mathbb{N}}$ converge (dans $\mathbb{Q}$) vers $l\in\mathbb{Q}$ lorsque $$\forall\varepsilon\in\mathbb{Q}_{+}^{*},\ \exists N\in\mathbb{N},\ \forall n\ge N, |r_{n}-l|\le\varepsilon$$.
On vérifie que les résultats suivants, déjà prouvés sur $\mathbb{R}$ dans le cours, restent valables sans modification.
\begin{itemize}
	\item Lorsqu'elle existe, la limite est unique.
	\item On peut effectuer les mêmes opérations algébriques (somme, produit, etc.).
	\item La suite $\displaystyle(\frac{1}{2^{n}})_{n\in\mathbb{N}}$ tend vers 0 dans $\Q$ (conséquence d'un lemme).
\end{itemize}
De même, une suite de rationnels $(r_{n})_{n\in\mathbb{N}}$ est une suite de Cauchy (dans $\mathbb{Q}$) lorsque $$\forall\varepsilon\in\mathbb{Q}_{+}^{*},\ \exists N\in\mathbb{N},\ \forall m,n\ge N,\ |r_{m}-r_{n}|\le\varepsilon$$
Ensuite, si on veut passer par $\mathbb{R}$, il suffit d'écrire $\sqrt{2}$ comme limite d'une suite de rationnels. Mais si on cherche une démonstration plus intrinsèque qui passe uniquement par les rationnels, en voici une. La preuve fonctionne en plusieurs étapes.
\begin{itemize}
	\item Première étape: on construit une suite $(r_{n})$ telle que $r_{n}>0$ et $r_{n}^{2}\rightarrow2$ dans $\mathbb{Q}$. Soit la suite $(p_{n},q_{n})$ à valeurs dans $\mathbb{N}^{2}$ définie par $$\begin{cases}(p_{0},q_{0})=(1,0)\\ \forall n\ge0,(p_{n+1},q_{n+1})=(p_{n}+2q_{n},p_{n}+q_{n})\end{cases}$$
	On montre alors par récurrence la propriété suivante: $p_{n}\ge1, q_{n}\ge n p_{n}^{2}-2q_{n}^{2}=(-1)^{n}$.
	\begin{itemize}
		\item[--]Initialisation : $p_0 \geqslant 1$, $q_0 \geqslant 0$ et $p_0^2 - 2q_0^2 = 1 = (-1)^0$.
		\item[--]Hérédité : soit $n \in \mathbb{N}$ tel que la propriété soit vraie.
		\begin{itemize}
			\item[*] On a $p_{n+1} = p_n + 2q_n \geqslant 1 + 2n \geqslant 1$.
			\item[*] Par ailleurs, $q_{n+1} = p_n + q_n \geqslant 1 + n$.
			\item[*] Enfin, on a \begin{align*}
				p_{n+1}^2 - 2q_{n+1}^2 = (p_n + 2q_n)^2 - 2(p_n + q_n)^2 &= -p_n^2 + 2q_n^2\\
				&= -(-1)^n\\
				&= (-1)^{n+1}
			\end{align*}
		\end{itemize}
	\end{itemize}
	On en déduit $\forall n > 0$, $q_n \neq 0$ et $\displaystyle\frac{p_n^2}{q_n^2} = 2 + \frac{(-1)^n}{q_n^2}$ Or $\displaystyle\left\lvert\frac{(-1)^n}{q_n^2}\right\rvert = \frac{1}{q_n^2} \leqslant \frac{1}{n^2} \leqslant \frac{1}{2^n} \rightarrow 0$. Donc $\frac{p_n^2}{q_n^2} \rightarrow 2$
	\item Deuxième étape : la suite $(r_n)$ est de Cauchy dans $\mathbb{Q}$. En effet, soit $\varepsilon \in \mathbb{Q}_{+}^{*}$. À partir d'un certain rang $N_1$, on a $r_n^2 \geqslant 1$ (puisque ce nombre tend vers 2), $r_n \geqslant 1$. À partir d'un certain rang $N_2$, $|r_n^2 - 2| \leqslant \varepsilon$. Pour $m, n \geqslant max(N_1, N_2)$, alors $$|r_n - r_m| = \frac{|r_n^2 - r_m^2|}{r_n + r_m} \leqslant \frac{2\varepsilon}{2} = \varepsilon$$
	\item Troisième étape : supposons que la suite $(r_n)$ converge vers $r \in \mathbb{Q}$. Alors $(r_n^2)$ convergerait vers $r^2$ puis par unicité de la limite on obtiendrait $r^2 = 2$, contradiction.
\end{itemize}
Conclusion : $\mathbb{Q}$ n'est pas complet !
\end{proof}
\begin{rmq}
	Nous connaissons donc deux propriétés de $\R$ que ne possède pas $\Q$ : la propriété de la borne supérieure et la complétude. En fait, ces deux propriétés ne sont que les deux versants d'une même montagne. \par En construisant $\R$ comme nous l'avons fait, nous avons mis l'accent sur la propriété de la borne supérieure, et la complétude en a découlé. Mais il existe une autre construction de $\R$, qui consiste à quotienter l'anneau des suites rationnelles de Cauchy par l'idéal des suites rationnelles qui tendent vers $0$ (cf. DM Avancé). Cette construction (qui peut être généralisée à d'autres ensembles) rend $de$ $facto$ $\R$ complet, et la propriété de la borne supérieure est alors une conséquence aisément démontrable (pas sûr du call là...) : on commence par démontrer le théorème des suites adjacentes, et le reste s'en déduit.
\end{rmq}
En définissant la notion de suite de Cauchy dans $\C$ comme dans $\R$ (à condition de remplacer les valeurs absolues par des modules), on obtient le 
\begin{thm}
	$\C$ est complet.
\end{thm}
\begin{proof}
	Il suffit de remarquer que si la suite complexe $(a_n+ib_n)_{n\in\N}$ est de Cauchy, alors $(a_n)_{n\in\N}$ et $(b_n)_{n\in\N}$ le sont aussi en utilisant l'inégalité sur les modules des parties réelles et imaginaires.
\end{proof}



\section{Compléments : retour sur la théorie des ensembles, HP}
Nous arrivons maintenant au terme de notre voyage ensembliste. Après avoir posé axiomatiquement l'existence de $\mathbb{N}$, nous avons construit $\mathbb{Z}$, $\mathbb{Q}$, $\mathbb{R}$, $\mathbb{C}$... puis toute l'analyse s'en déduit, comme le montre la suite du cours. Il est assez merveilleux de contempler la façon dont quelques axiomes de base très simples sur $\mathbb{N}$ nous ont permis de construire un socle inamovible pour une grande partie des mathématiques... Mais on peut aller plus loin...\\

Rappelons qu'à partir de $\mathbb{N}$, les éléments de $\mathbb{Z}$ ont été définis comme des classes d'équivalence sur $\mathbb{N}^2$, les éléments de $\mathbb{Q}$ comme des classes d'équivalence sur $\mathbb{Z} \times \mathbb{Z}^*$, les éléments de $\mathbb{R}$ comme des parties de $\mathbb{Q}$, et enfin les éléments de $\mathbb{C}$ comme des couples dans $\mathbb{R}^2$. Bref, à chaque fois, les nouveaux nombres que nous introduisions étaient des ensembles. De même pour l'analyse, une fonction ou une application est un ensemble. D'où l'idée un peu folle qui a germé dans le cerveau des mathématiciens : et si tout était ensemble ? Il ne reste plus qu'à construire les éléments de $\mathbb{N}$ eux-mêmes comme des ensembles, et ce sera effectivement vrai. Nous allons voir qu'une telle construction est possible. \\ \par

Il y a quelques chapitres, nous avions posé axiomatiquement l'existence d'un ensemble non vide $\mathbb{N}$, muni d'une relation d'ordre, et qui vérifiait trois propriétés fondamentales. Pour mémoire :
\begin{itemize}
	\item toute partie non vide admet un minimum ;
	\item toute partie non vide et majorée admet un maximum ;
	\item $\mathbb{N}$ n'admet pas de maximum.
\end{itemize}
Nous en avions notamment déduit la validité du principe de récurrence, qui est à la base de l'arithmétique (définition par récurrence des lois de $\mathbb{N}$, division euclidienne). Mais le point de vue inverse est possible : il existe une construction de $\mathbb{N}$, due au mathématicien italien Giuseppe Peano, qui postule d'emblée la validité du principe de récurrence (c'est un axiome), pour en déduire les trois propriétés fondamentales sus-citées. Les deux constructions sont donc équivalentes.\\

Or quel que soit le point de vue adopté (à partir des trois propriétés fondamentales ou à partir des axiomes de Peano), on peut se demander si un tel ensemble $\mathbb{N}$ doit vraiment être posé axiomatiquement. En clair : peut-on démontrer l'existence et l'unicité de $\mathbb{N}$ ? \\ \par

La question de l'unicité est assez simple. Par exemple, si deux ensembles ordonnés non vides vérifient les trois propriétés fondamentales, il est facile de construire une bijection de l'un dans l'autre qui respecte la relation d'ordre (on peut le faire par récurrence). On peut donc dire que, sous réserve d'existence, $\mathbb{N}$ est unique "à isomorphisme près".\\

Reste la question de l'existence. Il faudrait pouvoir construire $\mathbb{N}$ de manière effective à partir des axiomes de Zermelo-Fraenkel. Ainsi, les trois propriétés fondamentales (ou les axiomes de Peano selon le point de vue) perdraient leur statut d'axiome, pour devenir des théorèmes. Sur le plan mathématique ce serait satisfaisant, car on est toujours un peu gêné d'introduire des axiomes supplémentaires : par souci d'économie et d'esthétique, on peut se demander si tout cela n'est pas redondant. Cette construction est presque possible, à ceci près que nous allons devoir introduire un dernier axiome, dit "axiome de l'infini".\\ \par

Construisons donc $\mathbb{N}$ de manière effective. Ses éléments seront eux-mêmes des ensembles. Intuitivement, on choisit pour 0 le "plus petit" ensemble possible, c'est-à-dire qu'on pose 0 = $\emptyset$. Une fois qu'on a construit $n$ (comme un ensemble, rappelons-le), il est légitime de chercher un ensemble "juste un peu plus grand". Le mieux est de réunir $n$ avec un singleton $\{m\}$, tel que
$m \notin n$. L'axiome de fondation nous suggère alors de choisir $m=n$. Ainsi, pour tout $n$, on
pose $n+1 = n \cup \{n\} = \{0,1,...,n\}$. Par exemple, on a
$$\begin{cases}
	0=\emptyset\\
	1 = \{\emptyset\}\\
	2 = \{\emptyset, \{\emptyset\}\}\\
	3 = \{\emptyset, \{\emptyset\}, \{\emptyset, \{\emptyset\}\}\}\\
	\vdots
\end{cases}$$

Il reste à formaliser cela.

\begin{dfn}
	Si $X$ est un ensemble, on note $\mathcal{Q}(X)$ la propriété $$\begin{cases}
	\alpha \in X \\
	\forall x \in X, x \cup \{x\} \in X
\end{cases}$$
\end{dfn}

$X$ est donc un ensemble qui contient notre $0$, et qui est stable par l'opération "successeur".
On pose alors l'axiome suivant.

\begin{ax}[Axiome de l'infini]
	Il existe au moins un ensemble $X$ qui vérifie $\mathcal{Q}(X)$.
\end{ax}

Maintenant qu'il existe au moins un tel ensemble $X$, on va définir $\mathbb{N}$ comme le plus petit de ces ensembles au sens de l'inclusion. Cela peut se faire par compréhension. Soit $X$ un ensemble qui vérifie $\mathcal{Q}(X)$ et posons la

\begin{dfn}[$\N$]
	On définit $\mathbb{N} = \{x \in X\ | \ \forall Y, \mathcal{Q}(Y) \Rightarrow x \in Y\}$.
\end{dfn}
On remarque tout de suite que $\mathbb{N}$ vérifie $\mathcal{Q}(\mathbb{N})$. Notons aussi que la définition de $\mathbb{N}$ est indépendante du choix de $X$. En effet, si nous choisissons un autre ensemble $\tilde{X}$, dont dérive un certain $\tilde{\mathbb{N}}$ alors quel que soit $x \in \mathbb{N}$ on a $\mathbb{Q}(\mathbb{N})$ donc $x \in \mathbb{N}$. Autrement dit, on a $\mathbb{N} \subset \tilde{N}$. Par symétrie on a l'inclusion réciproque, puis conclut par l'axiome d'extensionnalité.\par 
Tel que nous l'avons contruit, $\mathbb{N}$ vérifie automatiquement un certain principe de récurrence, mais d'un genre bien particulier. Une récurrence simple, de type ensembliste, avec "initialisation à $\emptyset$". Nous en déduirons les trois propriétés fondatrices de $\mathbb{N}$. Pour commencer, on a le

\begin{thm}[Principe de récurrence ensembliste]
	Soit $A$ une partie de $\mathbb{N}$ vérifiant les deux points suivants.
\begin{itemize}
	\item Initialisation: $\emptyset \in A$
	\item Hérédité: $\forall n,\ n \in A \Rightarrow n \cup \{n\} \in A$.
\end{itemize}
Alors $A = \mathbb{N}$.
\end{thm}
\begin{proof}
	Il suffit de remarquer que $A$ vérifie $\mathcal{Q}(A)$. Ainsi, si $n \in \mathbb{N}$ en posant $Y = A$ on a $n \in A$. On en déduit que $\mathbb{N} \subset A$, et comme $A \subset \mathbb{N}$ a $A = \mathbb{N}$ par extensionnalité.
\end{proof}

\begin{dfn}
	Sur $\mathbb{N}$, on définit la relation par $m \leqslant n$ (m est inférieur ou égal à n).
\end{dfn}

Nous allons bien sûr démontrer que $\leqslant$ est une relation d'ordre, mais pour nous simplifier la tâche, nous allons utiliser un petit lemme à l'énoncé pour le moins surprenant.

\begin{lem}
	Soit $m, n \in \mathbb{N}$. Si $m \in n$ alors $m \subset n$ (non, il n'y a pas de faute de frappe!).
\end{lem}
\begin{proof}
	On démontre le résultat par récurrence sur $n$. Plus précisément, définissons par compréhension $$A := \{n \in \mathbb{N} \mid \forall m, [m \in n \Rightarrow m \subset n] \}$$
	\begin{itemize}
		\item Initialisation : si $n=\emptyset$, alors on ne peut pas avoir $m\in n$ et comme le faux implique n'importe quoi, on en déduit que $n\in A$
		\item Hérédité : soit $n\in A$ et montrons que $n\cup\left\{n\right\}\in A$. Soit donc $m\in n\cup\left\{n\right\}$/
		\begin{itemize}
			\item[--] Si $m\in n$, alors $m\subset n$ car $n\in A$ (on pourrait dire qu'on applique l'hypothèse de récurrence) et $a$ $fortiori$ on a $m\subset n\cup\left\{n\right\}$.
			\item[--] Si $m\in\left\{n\right\}$, alors $m=n$ donc $m\subset n$ et on se ramène au cas précédent.
		\end{itemize}
	\end{itemize}
	Ainsi la récurrence est achevée.
\end{proof}
\begin{cor}
	Soit $m, n \in \mathbb{N}$. $m \leqslant n$ alors $m \subset n$
\end{cor}
\begin{proof}
	Si $m\in n$ c'est une conséquence du lemme précédent, et si $m=n$ le résultat est immédiat.
\end{proof}
\begin{cor}
	La relation $\leqslant$ est une relation d'ordre sur $\N$.
\end{cor}
\begin{proof}
	Vérifions les axiomes : 
	\begin{itemize}
		\item Réflexivité : soit $n\in\N$, on a $n=n$ donc $n\leqslant n$.
		\item Symétrie : soit $m,n\in\N$ tels que $m\leqslant n$ et $n\leqslant m$. Alors $m\subset n$ et $n\subset m$ donc $m=n$ par extensionnalité.
		\item Transitivité : soit $m,n,p\in\N$ tels que $m\leqslant n$ et $n\leqslant p$. Si $m=n$, on a bien $m\leqslant p$. Sinon, on a $m\in n$, et comme $b\subset p$, on a $m\in p$, donc $m\leqslant p$.
	\end{itemize}
	Ainsi, la preuve est achevée.
\end{proof}
\begin{lem}
	Soit $m,n\in\N$. Si $m\in n$, alors $m\cup\left\{m\right\}\leqslant n$.
\end{lem}
\begin{proof}
	On montre encore le résultat par récurrence sur $n$. Définissons par compréhension $$A:=\left\{n\in\N\ | \ \forall m,\ (m\in\N\ \ \text{et}\ \ m\in n)\implies m\cup\left\{m\right\}\leqslant n\right\}$$
	\begin{itemize}
		\item Initialisation : quel que soit $m$, on n'a jamais $m\in\emptyset$, et comme le faux implique n'importe quoi, on en déduit que $\emptyset\in A$.
		\item Hérédité : supposons que $n\in A$, et soit $m\in\N$ tel que $m\in n\cup\left\{n\right\}$.
		\begin{itemize}
			\item[--] Si $m\in n$, comme $n\in A$, on a $m\cup\left\{m\right\}\leqslant n$ et comme $n\leqslant n\cup\left\{n\right\}$ (car $n\in n\cup\left\{n\right\}$), on conclut par transitivité.
			\item[--] Si $m\in\left\{n\right\}$, alors $m=n$ donc $m\cup\left\{m\right\}=n\cup\left\{n\right\}$ et on conclut par réflexivité.
		\end{itemize}
	\end{itemize}
	Ainsi, la récurrence est achevée.
\end{proof}
Nous en arrivons au théorème final.
\begin{thm}
	$\N$ ainsi construit est non vide, et il vérifie les trois propriétés suivantes.
	\begin{itemize}
		\item N n'admet pas de maximum.
		\item Toute partie non vide et majorée admet un maximum.
		\item Toute partie non vide admet un minimum.
	\end{itemize}
\end{thm}
\begin{proof}
	Montrons que $\mathbb{N}$ vérifie ce qu'on lui demande.
	\begin{itemize}
		\item $\mathbb{N}$ est non vide. En effet, il contient $\emptyset$.
		\item $\mathbb{N}$ n'admet pas de maximum. Supposons par l'absurde que $\mathbb{N}$ admette un maximum $n$. Alors comme $n \cup \{n\} \in \mathbb{N}$, on aurait $n \cup \{n\} \subseteq n$ et en particulier on aurait $n \in n$, contradiction d'après l'axiome de fondation.
		\item Toute partie non vide et majorée de $\mathbb{N}$ admet un maximum. Soit $B$ non vide et majorée par $a \in \mathbb{N}$, et montrons par récurrence sur $a$ que $B$ admet un maximum. Plus précisément, nous introduisons par compréhension $$A := \{a \in \mathbb{N}\ |\ \forall B, B \neq \emptyset \text{ et } B \text{ majorée par } a \Rightarrow B \text{ admet un maximum} \}$$
		\begin{itemize}
			\item[--]Initialisation: si $B$ est majorée par $\emptyset$ alors tous les éléments de $B$ sont $\leqslant \emptyset$ donc $\subseteq \emptyset$ et ils sont donc $=\emptyset$. $B$ étant non vide, elle admet donc un maximum, qui est $\emptyset$. Ainsi, $\emptyset \in A$.
			\item[--]Hérédité: supposons que $a \in A$ et montrons que $a \cup \{a\} \in A$. Soit $B$ non vide et majorée par $a \cup \{a\}$. Si $a \cup \{a\} \in B$, alors $B$ admet un maximum. Sinon, tous les éléments de $B$ appartiennent à $a \cup \{a\}$ donc ils appartiennent à $a$ ou ils sont égaux à $a$. Dans tous les cas ils sont $\leqslant a$, on conclut par hypothèse de récurrence.
		\end{itemize}
		\item Toute partie non vide de $\mathbb{N}$ admet un minimum. Soit $A$ non vide. Alors $A$ admet au moins un minorant, qui est $\emptyset$. Donc si on définit par compréhension l'ensemble de ses minorants $$B := \{b \in \mathbb{N}\ |\ \forall a, a \in A \Rightarrow b \leqslant a\}$$ alors $B$ est non vide. De plus, $A$ est non vide donc elle admet au moins un élément, et cet élément majore $B$. Ainsi, d'après le point précédent $B$ admet un maximum $b$. Montrons que $b$ est le minimum de $A$. Déjà il minore bien $A$, donc il ne reste plus qu'à montrer qu'il appartient à $A$.\\ Supposons que $b \notin A$. Alors pour tout $a \in A$ on pourrait écrire $b \in a$. Mais alors on aurait $b \cup \{b\} \leqslant a$ le lemme, donc $b \cup \{b\} \in B$. En particulier, on aurait $b \cup \{b\} \leqslant b$, donc $b \cup \{b\} \subseteq b$; puis $b \in b$, contradiction l'axiome de fondation.
	\end{itemize}
	Ainsi la preuve est achevée.
\end{proof}
\begin{rmq}
	On rappelle que $\leqslant$ est automatiquement totale, puisque si $(m,n) \in \mathbb{N}^2$, l'ensemble $\{m,n\}$ est non vide donc il admet un minimum.
\end{rmq}

\newpage
\chapter{Nombres complexes} %chapitre complet%
\section{Premières définitions}
\subsection{Construction de $\C$, HP}
\begin{dfn}[Ensemble des nombres complexes et i]
	On définit un nombre complexe comme un couple de réels $(a,b) \in \R^2$. On note $\C$ l'ensemble des nombres complexes et on définit le nombre complexe $i$ par $i=(0,1)$.
\end{dfn}
\begin{rmq}L'application de $\R$ dans $\C$ qui à $x$ associe $(x,0)$ est injective, ce qui permet de considérer abusivement que $\R \subset \C$. \end{rmq} 
\begin{dfn}[Lois de $\C$]
	On définit les lois $+$ et $\times$ sur $\C$ de la manière suivante :
	\begin{enumerate}
		\item $\forall ((a,b),(c,d))\in\C^2, \ (a,b)+(c,d)=(a+c,c+d)$
		\item $\forall ((a,b),(c,d))\in\C^2, \ (a,b)\times(c,d)=(ac-bd,ad+bc)$
	\end{enumerate}
	On vérifie immédiatement que $\C$ est stable par ces lois, et que ces lois sont associatives et commutatives. On vérifie de même que $\times$ est distributive sur $+$. Leurs éléments neutres respectifs sont $(0,0)$ et $(1,0)$. $(0,0)$ est absorbant pour $\times$. Ces lois prolongent les lois usuelles de $\R$. 
\end{dfn}
\begin{thm}[Unicité de la forme algébrique]
	$$\forall z \in \C, \ \exists!(a,b)\in\R^2, \ z=a+ib$$
	C'est ce que on appelle la forme \textbf{algébrique} de $z$.
\end{thm}
\begin{prop}
	$(\C,+,\times)$ est un corps.
\end{prop}
\begin{rmq}
	Attention : il n'y a pas de relation d'ordre convenable sur $\C$.
\end{rmq}
\subsection{Conjugaison}
\begin{dfn}[Conjugué]
	Soit $z \in \C$ qu'on écrit sous forme algébrique $z=a+ib$. Le \textbf{conjugué} de $z$, noté $\overline{z}$, est défini par : $$\overline{z}:=a-ib$$
\end{dfn}
\begin{prop}[Involutivité du passage au conjugué]
	On a $\forall z\in\C,\ \overline{(\overline{z})}=z$.
\end{prop}
\begin{prop}
	On a :
	$$\forall z\in\C,\ \begin{cases}
		\displaystyle\text{Re}(f)=\frac{z+\overline{z}}{2} \\
		\displaystyle\text{Im}(f)=\frac{z-\overline{z}}{2i}
	\end{cases}$$
\end{prop}
\begin{prop}[Caractérisation de $\R$ et $i\R$ par la conjugaison]
	Soit $z\in\C$.$$z\in\R\iff \overline{z}=z$$ $$z\in i\R\iff \overline{z}=-z$$
\end{prop}
\begin{rmq}
	On pourra souvent utiliser cette proposition pour montrer qu'un complexe appartient à $\R$ ou $i\R$, en particulier en présence de nombres complexes de module $1$, avec lesquels le conjugué fait bon ménage.
\end{rmq}
\begin{prop}
	$$\forall (z,z')\in\C^2,\ \begin{cases}
		\overline{z+z'}=\overline{z}+\overline{z'} \\
		\overline{zz'}=\overline{z}\overline{z'}
	\end{cases}$$
\end{prop}
\begin{prop}
	Soit $z\in\C$. On a $\overline{-z}=-\overline{z}$. De plus, supposons que $z\ne0$. Alors $\overline{z}\ne 0$, et on a $\displaystyle\overline{\left(\frac{1}{z}\right)}=\frac{1}{\overline{z}}$.
\end{prop}

\begin{exa}
	$$\forall(z,z')\in\C^2,\ \overline{z-z'}=\overline{z}-\overline{z'}$$
	$$\forall(z,z')\in\C\times\C\st,\ \overline{\left(\frac{z}{z'}\right)}=\frac{\overline{z}}{\overline{z'}}$$
\end{exa}
\begin{prop}[$\R$-linéarité des parties réelle et imaginaire]
	Soit $(z,z')\in\C^2$ et $(\lambda,\mu)\in\R^2$. Alors : $$\begin{cases}
		\text{Re}(\lambda z+\mu z')=\lambda\text{Re}(z)+\mu\text{Re}(z') \\
		\text{Im}(\lambda z+\mu z')=\lambda\text{Im}(z)+\mu\text{Im}(z')
	\end{cases}$$
\end{prop}
\subsection{Module}
\begin{dfn}
	Soit $z \in \C$ qu'on écrit sous forme algébrique $z=a+ib$. Le \textbf{module} de $z$, noté $\lvert z\rvert$, est défini par : $$\lvert z\rvert:=\sqrt{a^2+b^2}$$
\end{dfn}
\begin{rmq}
	Le module prolonge la valeur absolue réelle.
\end{rmq}
\begin{dfn}[Cercle unité $\U$]
	On pose $$\U:=\left\{z\in\C\ | \ \lvert z\rvert=1\right\}$$ C'est l'ensemble des \textbf{unimodulaires}, aussi appelé \textbf{cercle unité} ou \textbf{cercle trigonométrique}.
\end{dfn}
\begin{prop}
	$\forall z\in\C,\ z\overline{z}=\lvert z\rvert^2$
\end{prop}
\begin{met}[Calculs impliquant des modules]
	En présence de calculs de modules, penser à élever au carré, ce qui fait apparaître des conjugués très pratiques par la proposition précédente.
\end{met}
\begin{prop}
	$\forall z\in\U,\ \overline{z}=\displaystyle\frac{1}{z}$
\end{prop}
\begin{prop}[Séparation]
	$\forall z\in\C,\ \lvert z\rvert=0\iff z=0$
\end{prop}
\begin{prop}
	$\forall z\in\C,\ \lvert z\rvert=\lvert \overline{z}\rvert=\lvert -z\rvert=\lvert -\overline{z}\rvert$
\end{prop}
\begin{prop}[Vers l'inégalité triangulaire]
	On a : $$\forall z\in\C,\ \begin{cases}
		\left\lvert\text{Re}(z)\right\rvert\leqslant\lvert z\rvert \\
		\left\lvert\text{Im}(z)\right\rvert\leqslant\lvert z\rvert
	\end{cases}$$
\end{prop}
\begin{met}[Montrer que deux complexes sont différents]
	Pour montrer que deux complexes sont différents, il suffit de montrer que leurs parties réelles sont différentes, que leurs parties imaginaires sont différentes ou que leurs modules sont différents.
\end{met}
\begin{met}[Factorisation d'une somme de carrés dans $\C$]
	Ne pas oublier dans $\C$ la dernière identité remarquable $a^2+b^2=(a+ib)(a-ib)$.
\end{met}
\begin{prop}
	On a :
	$$\begin{cases}
		\forall (z,z')\in\C^2,\ \lvert zz'\rvert=\lvert z\rvert \lvert z'\rvert \\
		\forall(z,z')\in\C\times\C\st,\ \displaystyle\left\lvert\frac{z}{z'}\right\rvert=\frac{\lvert z\rvert}{\lvert z'\rvert}
	\end{cases}$$
\end{prop}
\begin{thm}[Inégalité triangulaire]
	$$\forall(z,z')\in\C,\ \lvert z+z'\rvert\leqslant\lvert z\rvert+\lvert z'\rvert$$
\end{thm}
\begin{cor}[Inégalité triangulaire généralisée]
	$$\forall n\geqslant2,\ \forall (z_1,\ldots,z_n)\in\C^n,\ \left\lvert\sum_{i=1}^{n}z_i\right\rvert\leqslant\sum_{i=1}^{n}\lvert z_i\rvert$$
\end{cor}
\begin{cor}[Inégalité triangulaire renversée]
	$$\forall(z,z')\in\C^2,\ \big|\lvert z\rvert-\lvert z'\rvert\big|\leqslant\lvert z-z'\rvert$$
\end{cor}
\begin{dfn}[Colinéarité directe]
	Soit $(z,z')\in\C^2$. Les trois assertions suivantes sont équivalentes :
	\begin{enumerate}
		\item $\displaystyle\begin{cases}
			\exists\lambda\in\R_+,\ z'=\lambda z \\
			\text{ou}\\
			\exists\mu\in\R_+,\ z=\mu z'
		\end{cases}$
		\item $\displaystyle\begin{cases}
			\exists\lambda\in\R_+,\ z'=\lambda z \\
			\text{ou}\\
			z=0
		\end{cases}$
		\item $\displaystyle\begin{cases}
			z'=0\\
			\text{ou}\\
			\exists\mu\in\R_+,\ z=\mu z'
		\end{cases}$
	\end{enumerate} Lorsque ces assertions sont vérifiées, on dit que $z$ et $z'$ sont \textbf{colinéaires directs}.
\end{dfn}
\begin{thm}[Cas d'égalité de l'inégalité triangulaire]
	Soit $(z,z')\in\C^2$. On a $\lvert z+z'\rvert=\lvert z\rvert+\lvert z'\rvert$ si, et seulement si, $z$ et $z'$ sont colinéaires directs.
\end{thm}
\section{Exponentielle complexe, trigonométrie}
\subsection{Exponentielle complexe}
\begin{dfn}[Prolongement de l'exponentielle à $\C$]
	La fonction $\exp:\R\rightarrow\R_+\st$ admet un prolongement, encore noté $\exp$ et donné par : $$\begin{array}{ccccc}\exp&:&\C&\rightarrow&\C\st \\&&z&\mapsto&\displaystyle \sum_{n=0}^{+\infty}\frac{z^n}{n}\end{array}$$ $\exp(z)$ peut aussi être noté $e^z$. On admettra que : $$\begin{cases}
		\forall(z,z')\in\C^2,\ e^{z+z'}=e^ze^{z'}\\
		\forall z\in\C,\ e^{\overline{z}}=\overline{e^z}
	\end{cases}$$
\end{dfn}
\begin{prop}
	On a : \begin{enumerate}
		\item $\forall\theta\in\R,\ \lvert e^{i\theta}\rvert=1$
		\item $\forall(a,b)\in\R^2,\ \lvert e^{a+ib}\rvert=e^a$
	\end{enumerate}
\end{prop}
\begin{dfn}[Cosinus et sinus]
	Les fonctions \textbf{cosinus} et \textbf{sinus}, notées respectivement $\cos$ et $\sin$ sont définies par : $$\begin{array}{ccccc}\cos&:&\:&\rightarrow&\R\\&&\theta&\mapsto&\text{Re}(e^{i\theta})\end{array}$$ et
	$$\begin{array}{ccccc}\sin&:&\R&\rightarrow&\R\\&&\theta&\mapsto&\text{Im}(e^{i\theta})\end{array}$$
\end{dfn}
\begin{prop}
	$\forall\theta\in\R,\ -1\leqslant \cos(\theta),\sin(\theta)\leqslant 1$
\end{prop}
\begin{prop}
	$\forall\theta\in\R,\ \cos(\theta)^2+\sin(\theta)^2=1$
\end{prop}
\begin{prop}[Formules d'Euler]
	On a les formules suivantes, dites \textbf{formules d'Euler} :
	$$\forall\theta\in\R,\ \begin{cases}
		\displaystyle\cos(\theta)=\frac{e^{i\theta}+e^{-i\theta}}{2} \\
		\displaystyle\sin(\theta)=\frac{e^{i\theta}-e^{-i\theta}}{2i}
	\end{cases}$$
\end{prop}
\begin{thm}[Admis]
	$\sin$ et $\cos$ sont dérivables et on a $\sin'=\cos$ et $\cos'=-\sin$.
\end{thm}
\begin{thm}[Technique de l'arc-moitié]
	Soit $(a,b)\in\R^2$. On a les formules dites d'\textbf{arc-moitié} $$\begin{cases}
		e^{ia}+e^{ib}=\displaystyle 2\cos\left(\frac{a-b}{2}\right)\exp\left(i\frac{a+b}{2}\right)\\
		e^{ia}-e^{ib}=\displaystyle 2i\sin\left(\frac{a-b}{2}\right)\exp\left(i\frac{a+b}{2}\right)
	\end{cases}$$
\end{thm}
\begin{rmq}
	Penser au cas où $b\equiv0\ [2\pi]$, $ie$ au cas où $e^{ib}=1$.
\end{rmq}
\subsection{Développements, linéarisations}
\begin{thm}[Formule de Moivre]
	Soit $\theta\in\R$. On a la \textbf{formule de Moivre} :$$\forall n\in\Z,\ (\cos(\theta)+i\sin(\theta))^n=\cos(n\theta)+i\sin(n\theta)$$ Autrement dit, $\forall n\in\Z,\ (e^{i\theta})^n=e^{in\theta}$.
\end{thm}
\begin{met}[Développement de $\cos(n\theta)$ ou $\sin(m\theta)$]
	On souhaite développer $\cos(n\theta)$ ou $\sin(m\theta)$, c'est-à-dire les exprimer comme des polynômes en $\cos(\theta)$ ou en $\sin(\theta)$. Pour cela : \begin{itemize}
		\item On écrit ce cosinus ou ce sinus comme une partie réelle ou imaginaire d'une exponentielle complexe.
		\item On utilise la formule de Moivre sur cette exponentielle complexe.
		\item On développe par le binôme de Newton après avoir écrit l'exponentielle sous la puissance sous forme algébrique.
		\item Enfin, on passe à la partie réelle ou imaginaire selon celle qu'on avait au départ pour obtenir le résultat.
	\end{itemize}
\end{met}
\begin{met}[Linéarisation de $\cos(\theta)^p\sin(\theta)^q$]
	On souhaite linéariser $\cos(\theta)^p\sin(\theta)^q$, c'est-à-dire le transformer en combinaison linéaire de $\cos(n\theta)$ et de $\sin(n\theta)$. Pour cela : 
	\begin{itemize}
		\item On exprime le cosinus et / ou le sinus à partir des formules d'Euler.
		\item Ensuite, on développe le tout avec le binôme de Newton, en mettant la puissance de $2$ en facteur pour simplifier les calculs. Si on en en présence d'un produit d'un cosinus et d'un sinus, il faut développer normalement après avoir utilisé le binôme.
		\item On utilise la formule de Moivre après avoir développé.
		\item On regroupe les termes conjuguées grâce aux formules d'Euler pour former des $\cos(n\theta)$ et des $\sin(n\theta)$<;
	\end{itemize}
\end{met}
\subsection{Le nombre $\pi$}
\begin{dfn}[$\pi$, HP]
	L'ensemble $\left\{\theta>0\ | \ \cos(\theta)=0\right\}$ admet un plus petit élément. On pose alors : $$\pi:=2\min\left\{\theta>0\ | \ \cos(\theta)=0\right\}$$
\end{dfn}
\begin{thm}[Admis]
	On a $e^{i\frac{\pi}{2}}=i$ et $e^{i\pi}=-1$.
\end{thm}
\begin{prop}
	Soit $\theta\in\R$.
	$$\begin{cases}
		\cos(-\theta)=\cos(\theta)\\
		\sin(-\theta)=-\sin(\theta)
	\end{cases}$$
	$$\begin{cases}
		\cos(\pi-\theta)=-\cos(\theta)\\
		\sin(\pi-\theta)=\sin(\theta)
	\end{cases}$$
	$$\begin{cases}
		\cos(\pi+\theta)=-\cos(\theta)\\
		\sin(\pi+\theta)=-\sin(\theta)
	\end{cases}$$
	$$\begin{cases}
		\cos\left(\frac{\pi}{2}-\theta\right)=\sin(\theta)\\
		\sin\left(\frac{\pi}{2}-\theta\right)=\cos(\theta)
	\end{cases}$$
	$$\begin{cases}
		\cos\left(\frac{\pi}{2}+\theta\right)=-\sin(\theta)\\
		\sin\left(\frac{\pi}{2}+\theta\right)=\cos(\theta)
	\end{cases}$$
\end{prop}
On trouvera en \textbf{fin de chapitre} un \textbf{formulaire de trigonométrie} réunissant toutes les formules de trigonométries à connaître.
\begin{prop}[Parité et périodicité]
	$\cos$ est paire et $2\pi$-périodique. $\sin$ est impaire et $2\pi$-périodique.
\end{prop}
\begin{thm}[Valeurs particulières]
	Voici quelques valeurs particulières des fonctions cosinus et sinus : $$\begin{tabular}{|c|c|c|}
		\hline $\theta$&$\cos(\theta)$&$\sin(\theta)$\\[3mm]
		\hline $0$&$1$&$0$ \\[4mm]
		\hline$\displaystyle\frac{\pi}{6}$&$\displaystyle\frac{\sqrt{3}}{2}$&$\displaystyle\frac{1}{2}$ \\[4mm]
		\hline$\displaystyle\frac{\pi}{4}$&$\displaystyle\frac{\sqrt{2}}{2}$&$\displaystyle\frac{\sqrt{2}}{2}$ \\[4mm]
		\hline$\displaystyle\frac{\pi}{3}$&$\displaystyle\frac{1}{2}$&$\displaystyle\frac{\sqrt{3}}{2}$ \\[4mm]
		\hline$\displaystyle\frac{\pi}{2}$&$0$&$1$ \\[4mm]
		\hline
	\end{tabular}$$
\end{thm}
\begin{dfn}[Congruence modulo $\alpha$]
	Soit $\alpha>0$ (en pratique, souvent $\pi$ ou $2\pi$ dans ce chapitre). Soit $(x,y)\in\R^2$. On dira que \textbf{$x$ et $y$ sont congrus modulo $\alpha$}, et on notera $x\equiv y\ [\alpha]$ lorsque : $$\exists k\in\Z,\ x=y+k\alpha$$
\end{dfn}
\begin{prop}
	La congruence modulo $\alpha>0$ est une relation d'équivalence sur $\R$.
\end{prop}
\begin{prop}[Quelques propriétés immédiates]
	Soit $\alpha>0$.
	\begin{itemize}
		\item $\forall(x,y)\in\R^2,\ x\equiv y\ [\alpha]\iff -x\equiv-y\ [\alpha]$
		\item $\forall(x,y,z)\in\R^3,\ x\equiv y\ [\alpha]\iff x+z\equiv y+z\ [\alpha]$
		\item $\forall \lambda>0,\ \forall(x,y)\in\R^2,\ x\equiv y\ [\alpha]\iff \lambda x\equiv\lambda y\ [\lambda\alpha]$
	\end{itemize}
\end{prop}
\begin{prop}
	Soit $\alpha>0$. Les lois $+$ et $-$ de $\R$ passent au quotient de la relation de congruence modulo $\alpha$. Précisément :
	$$\forall(x,x',y,y')\in\R^4,\ \begin{cases}
		x\equiv x'\ [\alpha] \\
		y\equiv y'\ [\alpha]
	\end{cases} \implies \begin{cases}
		x+y\equiv x'+y'\ [\alpha] \\
		x-y\equiv x'-y'\ [\alpha]
	\end{cases}$$
\end{prop}
\begin{rmq}
	Attention, cela ne fonctionne pas pour la multiplication en général, sauf pour les $\alpha\in\Z$.
\end{rmq}
\begin{rmq}
	Grâce à une partie entière, il est toujours possible d'encadrer un $\theta\in\R$ par deux multiples de $2\pi$ consécutifs, avec un des encadrements strict.
\end{rmq}
\begin{thm}[Quelques équations trigonométriques]
	Soit $(\theta,\psi)\in\R^2$. On a : $$\cos(\theta)=1\iff \theta\equiv 0\ [2\pi]$$ $$\cos(\theta)=0\iff \theta\equiv\displaystyle\frac{\pi}{2}\ [\pi]$$ $$\cos(\theta)=-1\iff \theta\equiv \pi\ [2\pi]$$ $$\cos(\theta)=\cos(\psi)\iff \theta\equiv\psi\ [2\pi] \ \ \text{ou}\ \ \theta\equiv-\psi\ [2\pi]$$ et $$\sin(\theta)=1\iff \theta\equiv \displaystyle\frac{\pi}{2}\ [2\pi]$$ $$\sin(\theta)=0\iff \theta\equiv0\ [\pi]$$ $$\sin(\theta)=-1\iff \theta\equiv\displaystyle-\frac{\pi}{2} \ [2\pi]$$ $$\sin(\theta)=\sin(\psi)\iff \theta\equiv\psi\ [2\pi] \ \ \text{ou}\ \ \theta\equiv\pi-\psi\ [2\pi]$$
\end{thm}
\begin{rmq}
	Bien souvent, un dessin d'un cercle trigonométrique pourra servir de preuve pour des inéquations trigonométriques, en résolvant le cas d'égalité avec le théorème qui précède.
\end{rmq}
\subsection{La fonction tangente}
\begin{dfn}[Tangente]
	On définit la fonction \textbf{tangente}, notée $\tan$, par : $$\begin{array}{ccccc}\tan&:&\R\setminus\left\{\displaystyle\frac{\pi}{2}+k\pi\ | \ k\in\Z\right\}&\rightarrow&\R\\&&\theta&\mapsto&\displaystyle\frac{\sin(\theta)}{\cos(\theta)}\end{array}$$
\end{dfn}
\begin{thm}[Étude de la fonction tangente]
	$\tan$ est impaire et $\pi$-périodique. Elle est dérivable en tout point de son domaine de définition et on a : $$\tan'=1+\tan^2=\frac{1}{\cos^2}$$ Elle est strictement croissante sur $\displaystyle\left]-\frac{\pi}{2},\frac{\pi}{2}\right[$ et vérifie $$\tan(x)\xrightarrow[x\to\displaystyle-\frac{\pi}{2}^+]{}-\infty$$ et $$\tan(x)\xrightarrow[x\to\displaystyle\frac{\pi}{2}^-]{}+\infty$$
\end{thm}
\begin{prop}[Formules impliquant la tangente de l'arc-moitié]
	Soit $\theta\in\R\setminus\left\{\pi+2k\pi\ | \ k\in\Z\right\}$ et posons $\displaystyle t:=\frac{\theta}{2}$. On a : $$\begin{cases}
		\displaystyle\cos(\theta)=\frac{1-t^2}{1+t^2}\\[3mm]
		\displaystyle\sin(\theta)=\frac{2t}{1+t^2}\\[3mm]
		\displaystyle\tan(\theta)=\frac{2t}{1-t^2}\\
	\end{cases}$$
\end{prop}
\begin{rmq}
	Ces formules se révèleront capitales en intégration.
\end{rmq}
\begin{rmq}
	Pour retenir ces formules, penser aux parités et au fait que la tangente n'est pas toujours définie, alors que le cosinus et le sinus si.
\end{rmq}
\begin{thm}[Valeurs particulières]
	Voici quelques valeurs particulières de la fonction tangente: $$\begin{tabular}{|c|c|}
		\hline $\theta$&$\tan(\theta)$\\[3mm]
		\hline $0$&$0$\\[4mm]
		\hline$\displaystyle\frac{\pi}{6}$&$\displaystyle\frac{1}{\sqrt{3}}$\\[4mm]
		\hline$\displaystyle\frac{\pi}{4}$&$1$\\[4mm]
		\hline$\displaystyle\frac{\pi}{3}$&$\sqrt{3}$\\[4mm]
		\hline
	\end{tabular}$$
\end{thm}
\begin{dfn}[Cotangente, HP]
	On définit la fonction \textbf{cotangente}, notée $\text{cotan}$, par : $$\begin{array}{ccccc}\text{cotan}&:&\R\setminus\left\{k\pi\ | \ k\in\Z\right\}&\rightarrow&\R\\&&\theta&\mapsto&\displaystyle\frac{\cos(\theta)}{\sin(\theta)}\end{array}$$
\end{dfn}
\begin{rmq}[Lien cotangente/tangente]
	Si $\theta\in\R\setminus\left\{\displaystyle k\frac{\pi}{2}\ | \ k\in\Z\right\}$, alors $\displaystyle\text{cotan}(\theta)=\frac{1}{\tan(\theta)}$
\end{rmq}
\begin{thm}[Équations trigonométriques]
	Soit $\theta,\psi\in\R\setminus\left\{\displaystyle\frac{\pi}{2}+k\pi\ | \ k\in\Z\right\}$. On a : $$\tan(\theta)=0\iff \theta\equiv0\ [\pi]$$ $$\tan(\theta)=\tan(\psi)\iff \theta\equiv\psi\ [\pi]$$
\end{thm}
\section{Première incursion dans les formules sommatoires}
\subsection{Préliminaires}
\begin{dfn}[Sommes et produits]
	Soit $(a_k)_{k\in\N}$ une suite de nombres complexes. Soit $n\in\N$.
	\begin{itemize}
		\item La \textbf{somme pour $k$ allant de $0$ à $n$}, notée $\displaystyle\sum_{k=0}^{n}$, désigne en fait $a_0+a_1+\ldots+a_n$.
		\item Le \textbf{produit pour $k$ allant de $0$ à $n$}, noté $\displaystyle\prod_{k=0}^{n}$, désigne en fait $a_0\times a_1\times\ldots\times a_n$.
		\item $k$ (ou la \textbf{variable muette} utilisée à la place de $k$) s'appelle \textbf{l'indice courant}. Par défaut, l'indice courant augmente de $1$ en $1$.
	\end{itemize}
	On définit de même des sommes et des produits pour $k$ allant de $n_0$ à $n_1\geqslant n_0$. Par convention, on pose : \begin{itemize}
		\item $\displaystyle\sum_{k=n}^{n-1}a_k=0$ ("somme vide")
		\item $\displaystyle\prod_{k=n}^{n-1}a_k=1$ ("produit vide")
	\end{itemize}
	Si $n_1\leqslant n_0-2$, on évite tout simplement de définir ces expressions.
\end{dfn}
\begin{rmq}[Définition formelle]
	En réalité, ces expressions sont définitions par récurrence et la convention sur les sommes et produits vides en est l'initialisation. Pour plus de détails, se référer à la quatrième partie "Sommes et produits" du chapitre "Introduction à l'algèbre".
\end{rmq}
\begin{rmq}
	Il est fondamental de maîtriser la technique de changement de variable ou de changement d'indice : il faut savoir passer d'une somme indexée par $k$ à une somme indexée par $k+1$, $2k$, etc.
\end{rmq}
\begin{thm}[Séparation / regroupement de termes]
	Soit $(a_k)_{k\in\N}$ et $(b_k)_{k\in\N}$ deux suites de nombres complexes. On a : $$\forall n\in\N,\ \sum_{k=1}^{n}(a_k+b_k)=\left(\sum_{k=1}^{n}a_k\right)+\left(\sum_{k=1}^{n}b_k\right)$$
	$$\forall n\in\N,\ \prod_{k=1}^{n}(a_k\times b_k)=\left(\prod_{k=1}^{n}a_k\right)\times\left(\prod_{k=1}^{n}b_k\right)$$
\end{thm}
\begin{thm}[Distributivité]
	Soit $(a_k)_{k\in\N}$ une suite de nombres complexes. On a :$$\forall n\in\N,\ \forall\lambda\in\C,\ \sum_{k=1}^{n}(\lambda a_k)=\lambda\left(\sum_{k=1}^{n}a_k\right)$$
	$$\forall n\in\N,\ \forall\lambda\in\C,\ \prod_{k=1}^{n}(\lambda a_k)=\lambda^n\left(\prod_{k=1}^{n}a_k\right)$$
\end{thm}
\begin{thm}[Relation de Chasles]
	Soit $(a_k)_{k\in\N}$ une suite de nombres complexes. On a : $$\forall(n_0,n_1,n_2)\in\N^3,\ \begin{cases}
		n_0\leqslant n_1\leqslant n_2 \\ \text{ou}\\ n_0\leqslant n_1+1\leqslant n_2
	\end{cases}\implies \left(\sum_{k=n_0}^{n_1}a_k\right)+\left(\sum_{k=n_1+1}^{n_2}a_k\right)=\sum_{k=n_0}^{n_2}a_k$$
	$$\forall(n_0,n_1,n_2)\in\N^3,\ \begin{cases}
		n_0\leqslant n_1\leqslant n_2 \\ \text{ou}\\ n_0\leqslant n_1+1\leqslant n_2
	\end{cases}\implies \left(\prod_{k=n_0}^{n_1}a_k\right)\times\left(\prod_{k=n_1+1}^{n_2}a_k\right)=\prod_{k=n_0}^{n_2}a_k$$
\end{thm}
\begin{prop}[Sommes usuelles]
	Soit $n\in\N$. On a : $$\sum_{k=0}^{n}k=\frac{n(n+1)}{2}$$ $$\sum_{k=0}^{n}k^2=\frac{n(n+1)(2n+1)}{6}$$ $$\sum_{k=0}^{n}k^3=\frac{n^2(n+1)^2}{4}=\left(\sum_{k=0}^{n}k\right)^2$$
\end{prop}
\begin{dfn}[Suite arithmétique de raison $r$]
	Soit $(u_n)_{n\in\N}\in\C^\N$ et $r\in\C$. $(u_n)_{n\in\N}$ est dite \textbf{arithmétique de raison $r$} lorsque $\forall n\in\N,\ u_{n+1}=u_n+r$. Dans ce cas, on a : $$\forall n\in\N,\ u_n=u_0+nr$$
\end{dfn}
\begin{dfn}[Suite géométrique de raison $q$]
	Soit $(v_n)_{n\in\N}\in\C^\N$ et $q\in\C$. $(v_n)_{n\in\N}$ est dite \textbf{géométrique de raison $q$} lorsque $\forall n\in\N,\ v_{n+1}=v_n\times q$. Dans ce cas, on a : $$\forall n\in\N,\ v_n=v_0\times q^n$$
\end{dfn}
\begin{thm}[Somme des termes d'une suite arithmétique]
	Soit $(u_n)_{n\in\N}\in\C^\N$ une suite arithmétique raison $r\in\C$. On a : $$\forall n\in\N,\ \sum_{k=0}^{n}u_k=(n+1)\frac{u_0+u_n}{2}$$
\end{thm}
\begin{thm}[Somme des termes d'une suite géométrique de raison $q\ne1$]
	Soit $(v_n)_{n\in\N}\in\C^\N$ une suite géométrique raison $q\in\C\setminus\left\{1\right\}$. On a : $$\forall n\in\N,\ \sum_{k=0}^{n}v_k=v_0\frac{1-q^{n+1}}{1-q}$$
\end{thm}
\begin{thm}[Formule de Bernoulli]
	Soit $(a,b)\in\C^2$ et $n\in\N$. Alors on a la \textbf{formule de Bernoulli} : $$a^n-b^n=(a-b)\sum_{k=0}^{n-1}a^{n-1-k}b^k$$
\end{thm}
\begin{rmq}[Cas $n$ impair]
	Lorsque $n$ est impair, on peut obtenir une formule pour $a^n+b^n$ en effectuant $b\leftarrow -b$ dans la formule de Bernoulli : $$a^n+b^n=(a+b)\sum_{k=0}^{n-1}(-1)^ka^{n-1-k}b^k$$
\end{rmq}
\begin{exa}
	$a^2-b^2=(a-b)(a+b)$
\end{exa}
\begin{exa}
	$a^3-b^3=(a-b)(a^2+ab+b^2)$
\end{exa}
\begin{exa}
	$a^3+b^3=(a+b)(a^2-ab+b^2)$
\end{exa}
\begin{rmq}
	On a aussi $a^n-b^n=\displaystyle\prod_{\omega\in\U_n}(a-\omega b)$ (on verra dans la suite du chapitre que $\U_n$ désigne l'ensemble des racines $n$-èmes de l'unité).
\end{rmq}
\subsection{Coefficients binomiaux}
\begin{dfn}[Factorielle]
	On pose $\forall n\in\N\st,\ n!:=1\times\ldots\times n$ et $0!:=1$ par convention. Plus rigoureusement, cette \textbf{factorielle} est définie par récurrence : $$\begin{cases}
		0!:=1 \\ \forall n\in\N, (n+1)!:=(n!)\times(n+1)
	\end{cases}$$
\end{dfn}
\begin{dfn}[Coefficients binomiaux]
	Soit $k\in\Z$ et $n\in\N$. On définit le \textbf{coefficient binomial} $\displaystyle\binom{n}{k}$ par : $$\binom{n}{k}:=\begin{cases}
		0&\text{si }k<0 \\
		\displaystyle\frac{n!}{k!(n-k)!}&\text{si }k\in\llbracket0,n\rrbracket \\
		0&\text{si }k>n
	\end{cases}$$
\end{dfn}
\begin{exa}[Cas particuliers fréquents]
	Voici quelques cas particuliers fréquents à connaitre par cœur :
	\begin{itemize}
		\item $\displaystyle\forall n\in\N,\ \binom{n}{0}=\binom{n}{n}=1$
		\item $\displaystyle\forall n\in\N\st,\ \binom{n}{1}=\binom{n}{n-1}=n$
		\item $\displaystyle\forall n\geqslant2,\ \binom{n}{2}=\binom{n}{n-2}=\frac{n(n-1)}{2}$
	\end{itemize}
\end{exa}
\begin{prop}[Relation de Pascal]
	On a la \textbf{relation de Pascal} :
	$$\forall n\in\N,\ \forall k\in\Z,\ \binom{n}{k}+\binom{n}{k+1}=\binom{n+1}{k+1}$$
\end{prop}
\begin{prop}[Symétrie des coefficients binomiaux]
	Les coefficients binomiaux sont symétriques au sens suivant : $$\forall n\in\N,\ \forall k\in\Z,\ \binom{n}{k}=\binom{n}{n-k}$$
\end{prop}
\begin{prop}[Formule du capitaine]
	On a la formule dite \textbf{du capitaine} :
	$$\forall(k,n)\in(\N\st)^2,\ k\binom{n}{k}=n\binom{n-1}{k-1}$$
\end{prop}
\begin{thm}[Binôme de Newton]
	On a la \textbf{formule du binôme de Newton}
	$$\forall n\in\N,\ \forall(a,b)\in\C^2,\ (a+b)^n=\sum_{k=0}^{n}\binom{n}{k}a^kb^{n-k}=\sum_{k=0}^{n}\binom{n}{k}a^{n-k}b^{k}$$
\end{thm}
\begin{thm}[Généralisation : formule du multinôme de Newton, HP]
	Il existe une formule plus générale qui généralise le binôme, appelée \textbf{formule du multinôme de Newton} : $$\forall n\in\N,\ \forall p\in\N\st,\ \forall(a_i)_{1\leqslant i\leqslant p}\in\C^p,\ \left(\sum_{i=1}^{p}a_i\right)^n=\sum_{k_1+\ldots+k_p=n}\frac{n!}{k_1!\ldots k_p!}a_1^{k_1}\ldots a_p^{k_p}$$
\end{thm}
\begin{proof}
	La preuve peut se faire par récurrence sur $p$ avec une hypothèse en $\forall n\in\N$et en utilisant le binôme. Elle peut aussi se faire par dénombrement.
\end{proof}
\begin{exa}
	$\forall n\in\N,\ \displaystyle\sum_{k=0}^{n}\binom{n}{k}=2^n$
\end{exa}
\begin{exa}
	$\forall n\in\N,\ \displaystyle\sum_{k=0}^{n}(-1)^k\binom{n}{k}=0^n=\delta_{n,0}$
\end{exa}
\begin{exa}
	$\forall n\in\N,\ \forall x\in\C,\ \displaystyle\sum_{k=0}^{n}\binom{n}{k}x^k=(x+1)^n$
\end{exa}
\begin{exa}
	$\forall n\in\N,\ \forall x\in\C,\ \displaystyle\sum_{k=0}^{n}\binom{n}{k}x^k(1-x)^{n-k}=1$
\end{exa}
\begin{met}[Calculs de sommes]
	Voici différentes manières de calculer des sommes :
	\begin{itemize}
		\item On pourra utiliser les formules de Bernoulli, du binôme de Newton, du multinôme de Newton, des sommes de termes d'une suite arithmétique ou géométrique.
		\item On pourra faire apparaître un télescopage.
		\item On pourra utiliser les formules du capitaine ou la relation de Pascal pour se ramener à des cas usuels ou des télescopages.
		\item Si ce sont des sommes qui mettent en jeu des $\cos(k\theta)$ ou des $\sin(k\theta)$, on peut les interpréter comme des parties réelles ou imaginaires de sommes d'exponentielles complexes, puis appliquer le binôme de Newton, la formule de Moivre ou la formule des sommes géométriques à ces sommes exponentielles. On pourra finir par un arc-moitié et il ne faut pas oublier de repasser à la partie réelle ou imaginaire.
		\item Si la somme dépend d'un paramètre $\theta$ défini sur un intervalle, on peut interpréter la somme comme la dérivée d'une autre sommes. On travaille alors sur la somme "primitivée", puis le résultat d'obtient par dérivation. On peut utiliser les autres méthodes pour calculer cette somme "primitivée". On pourra considérer la sommes des $k\sin(k\theta)$ pour $\theta\in]0,2\pi[$ en guise d'exemple.
	\end{itemize}
\end{met}
\section{Racines de l'unité}
\subsection{Écriture trigonométrique d'un nombre complexe}
\begin{lem}
	Soit $(a,b)\in\R^2$ tel que $a^2+b^2=1$. Alors : $$\exists\theta\in\R,\ \begin{cases}
		a=\cos(\theta)\\ b=\sin(\theta)
	\end{cases}$$
\end{lem}
\begin{cor}
	Soit $z\in\U$. Alors $\exists\theta\in\R,\ z=e^{i\theta}$. Autrement dit, la fonction $$\begin{array}{ccc}\R&\rightarrow&\U\\ \theta&\mapsto&e^{i\theta}\end{array}$$ est surjective.
\end{cor}
\begin{dfn}[Forme trigonométrique]
	Toute nombre complexe $z\in\C$ admet une \textbf{écriture trigonométrique}, ou \textbf{forme trigonométrique}, $ie$ une écriture de la forme $z=\rho e^{i\theta}$ avec $\rho\in\R_+$ et $\theta\in\R$.
\end{dfn}
\begin{exa}
	$1+i=\sqrt{2}e^{i\frac{\pi}{4}}$
\end{exa}
\begin{exa}
	$-2024=2024e^{i\pi}$
\end{exa}
\begin{dfn}[Arguments]
	Soit $z\in\C\st$. \textbf{Un argument} de $z$ est un $\theta\in\R$ tel qu'il existe $\rho\in\R_+$ tel qu'on ait $z=\rho e^{i\theta}$. Remarquons-qu'on a alors nécessairement $\rho\in\R_+\st$.
\end{dfn}
\begin{rmq}
	Par convention, $0$ n'a pas d'argument.
\end{rmq}
\begin{lem}
	Soit $(\alpha,\beta)\in\R^2$. On a $$e^{i\alpha}=e^{i\beta}\iff\alpha\equiv\beta\ [2\pi]$$
\end{lem}
\begin{thm}[Pseudo-unicité des écritures trigonométriques]
	L'ensemble des écritures trigonométriques de $0$ est l'ensemble ds couples $(0,\theta)$ pour $\theta\in\R$. Soit $z\in\C\st$. $z$ possède un argument $\theta_0$. L'ensemble des écritures trigonométriques de $z$ est l'ensemble des couples $(\lvert z\rvert,\theta_0+2k\pi)$ pour $k\in\Z$
\end{thm}
\begin{prop}
	Soit $(z,z')\in\C^2$. On a : $$e^z=e^{z'}\iff z-z'\in 2i\pi\Z$$
\end{prop}
\begin{exa}
	Soit $a\in\C$. Si $a=0$, l'équation $e^z=a$ n'a pas de solution. Si $a\ne0$, fixons $\theta_a$ un argument de $a$. L'ensemble des solutions vaut alors $\left\{\ln\lvert a\rvert +i(\theta_a+2k\pi)\ | \ k\in\Z\right\}$.
\end{exa}
\begin{dfn}[Argument principal]
	Soit $z\in\C\st$. \textbf{L'argument principal de $z$}, noté $\arg(z)$ (NS), est l'unique argument $\theta$ de $z$ tel que $\theta\in]-\pi,\pi]$.
\end{dfn}
\begin{rmq}
	Toujours bien préciser ce que signifie cette notation quand on l'utilise (elle est NS), et aussi préciser qu'on choisit et argument dans $]-\pi,\pi]$ car d'autres le prennent dans $[0,2\pi[$.
\end{rmq}
\begin{rmq}
	$0$ n'a pas d'argument, donc $a$ $fortiori$ par d'argument principal.
\end{rmq}
\begin{prop}[Propriété logarithmiques de l'argument principal]
	Soit $(z,z')\in(\C\st)^2$. On a : \begin{itemize}
		\item $\arg(zz')\equiv\arg(z)+\arg(z')\ [2\pi]$
		\item $\displaystyle\arg\left(\frac{z}{z'}\right)\equiv\arg(z)-\arg(z')\ [2\pi]$
		\item $\forall n\in\Z,\ \arg(z^n)\equiv n\arg(z)\ [2\pi]$
	\end{itemize}
\end{prop}
\begin{rmq}
	Attention, il faut bien mettre des $\equiv$ et non des $=$. En effet, par exemple, $n\arg(z)$ à très peu de chances de rester dans $]-\pi,\pi]$ lorsque $n$ est grand.
\end{rmq}
\begin{thm}[Superposition de deux ondes]
	Soit $(\lambda,\mu)\in\R^2\setminus\left\{(0,0)\right\}$. Alors : $$\exists(A,\varphi)\in\R_+\st\times\R,\ \forall t\in\R,\ \lambda\cos(t)+\mu\sin(t)=A\cos(t-\varphi)$$
\end{thm}
\begin{rmq}
	On peut utiliser ce théorème pour résoudre des inéquations ou équations trigonométriques (mais il faut alors expliciter $A$ et $\varphi$). Ce théorème sert aussi en physique (ksssss).
\end{rmq}
\subsection{Racines $n$-èmes}
Dans tout le paragraphe, on fixe $n\in\N\st$.
\begin{thm}[Racines $n$-èmes de l'unité]
	L'équation $z^n=1$ admet exactement $n$ solutions dans $\C$. Ce sont les $ e^{i\frac{2k\pi}{n}}$ pour $k\in\llbracket0,n-1\rrbracket$, qu'on appelle les \textbf{racines $n$-èmes de l'unité}. On note alors : $$\U_n:=\displaystyle\left\{e^{i\frac{2k\pi}{n}}\ | \ k\in\llbracket0,n-1\rrbracket\right\}$$
\end{thm}
\begin{rmq}[Interprétation géométrique]
	Géométriquement, dans le plan complexe, les racines $n$-èmes de l'unité sont exactement les sommets d'un $n$-gone régulier présent sur le cercle unité (dès que $n\geqslant2$ pour que la notion de $n$-gone ait du sens).
\end{rmq}
\begin{rmq}
	On a aussi : $\U_n:=\displaystyle\left\{e^{i\frac{2k\pi}{n}}\ | \ k\in\llbracket1,n\rrbracket\right\}$
\end{rmq}
\begin{exa}
	On a $\U_4=\left\{1,i,-1,-i\right\}$.
\end{exa}
\begin{exa}
	Soit $n\geqslant2$. On a $\displaystyle\sum_{\omega\in\U_n}\omega=0$.
\end{exa}
\begin{dfn}[Nombre $j$]
	On pose $j:=e^{i\frac{2\pi}{3}}$.
\end{dfn}
\begin{exa}
	On a $\U_3=\left\{1,j,j^2\right\}$.
\end{exa}
\begin{prop}[Quelques propriétés de $j$]
	On a $1+j+j^2=0$ et $j^2=\overline{j}=\displaystyle\frac{1}{j}$.
\end{prop}
\begin{met}
	Soit $u\in\C\st$ et $n\in\N\st$. Voici comment trouver les racines $n$-èmes de $u$ dans $\C$.
	\begin{itemize}
		\item On cherche une \textbf{solution particulière}. Pour cela, on écrit $u$ sous forme trigonométrique $u=\rho e^{i\theta}$ et alors $z_0:=\sqrt[n]{\rho}e^{i\frac{\theta}{n}}$ convient comme solution particulière.
		\item On en déduit toutes les solutions : $z^n=u$ ssi $z^n=z_0^n$ ssi $\displaystyle\left(\frac{z}{z}\right)^n=1$. Donc $z$ est solution si, et seulement si, $$\exists k\in\llbracket0,n-1\rrbracket,\ z=z_0e^{i\frac{k\theta}{n}}$$
		\item Conclusion : $0$ admet une racine $n$-ème dans $\C$ qui est $0$, et $u$ admet exactement $n$ racines $n$-èmes dans $\C$.
	\end{itemize} 
\end{met}
\begin{rmq}[Interprétation géométrique]
	De même, les $n$ racines $n$-èmes d'un complexe non nul sont exactement les sommets d'un $n$-gone régulier centré en $(0,0)$ (du moins dès que $n\geqslant2$).
\end{rmq}
\begin{met}[Recherche de racines carrées sous forme algébrique]
	Voici comment chercher les racines carrées sous forme algébrique d'un complexe non nul. On raisonne par analyse-synthèse.
	\begin{itemize}
		\item Analyse : On écrit $z=a+ib$ une solution sous forme algébrique. En mettant au carré et en égalant avec le nombre complexe dont on cherche la racine, on peut obtenir une expression de $a^2-b^2$ en passant aux parties réelles, une expression de $a^2+b^2$ en passant aux modules et une expression de $ab$ en passant aux parties imaginaires. On combine les deux premières par demi-somme et demi-différence pour obtenir $a$ et $b$ au signe près. On conclut sur le signe de $a$ et $b$. On obtient alors deux solutions opposées.
		\item Synthèse : Réciproquement, puisque le nombre complexe dont on cherche les racines carrées est non nul, il admet deux racines carrées dans $\C$, qui sont nécessairement celles obtenues à la fin de l'analyse.
	\end{itemize}
\end{met}
\subsection{Équations du second degré}
\begin{thm}[Résolution d'une équation du second degré à coefficients complexes]
	Soit $(a,b,c)\in\C\st\times\C^2$. On considère l'équation $az^2+bz+c=0$ d'inconnue $z\in\C$. Notons $S_{a,b,c}$ l'ensemble des solutions. On pose $\Delta:=b^2-4ac$, dit \textbf{discriminant} de cette équation. Soit $\delta$ une racine carrée de $\Delta$ dans $\C$ (si $\Delta$ en possède $2$, on prend celle que l'on veut cela n'a pas d'importance). Alors, $$S_{a,b,c}=\left\{\frac{-b-\delta}{2a},\frac{-b+\delta}{2a}\right\}$$
\end{thm}
\begin{rmq}
	En particulier, si $\Delta=0$ alors $\delta=0$ et $\displaystyle S_{a,b,c}=\left\{\frac{-b}{2a}\right\}$.
\end{rmq}
\begin{rmq}[Notations et formulations pour $\delta$]
	Pour $\delta$, bien dire que $\delta$ est \textbf{une} racine carrée de $\Delta$ dans $\C$ ou que $\delta^2=\Delta$.
\end{rmq}
\begin{rmq}[Notations et formulations pour l'ensemble des solutions]
	Voici les trois \textbf{seules} bonnes manières de conclure la résolution : \begin{itemize}
		\item Soit $z\in\C$. $z$ est solution ssi $z=\ldots$ ou $z=\ldots$.
		\item Les solutions sont données par : $z_1:=\ldots$ et $z_2:=\ldots$.
		\item L'ensemble des solutions vaut : $\left\{\ldots;\ldots\right\}$.
	\end{itemize}
\end{rmq}
\begin{rmq}
	On retiendra qu'il y a toujours au moins une solution. Précisément, si $\Delta=0$, une seule solution, et si $\Delta\ne0$, exactement deux solutions.
\end{rmq}
\begin{cor}
	Soit $(a,b,c)\in\C\st\times\C^2$. On considère l'équation $az^2+bz+c=0$ d'inconnue $z\in\C$. Soit $\delta$ une racine carrée du discriminant de cette équation et posons $\displaystyle z_1:=\frac{-b-\delta}{2}$ et $\displaystyle z_2:=\frac{-b+\delta}{2}$. Alors on a : $$\begin{cases}
		z_1+z_2=\displaystyle\frac{-b}{a}\\ z_1z_2=\displaystyle\frac{c}{a}
	\end{cases}$$
\end{cor}
\begin{rmq}
	On a coutume de dire que "la somme des racines comptées avec multiplicité vaut $\displaystyle\frac{-b}{a}$" et que "le produit des racines comptées avec multiplicité vaut $\displaystyle\frac{c}{a}$". 
\end{rmq}
\begin{cor}
	Soit $(S,P,u,v)\in\C^4$. Alors $u+v=S$ et $uv=P$ si, et seulement si, $u$ et $v$ sont les deux racines (éventuellement confondues) de l'équation $z^2-Sz+P=0$ d'inconnue $z\in\C$.
\end{cor}
\begin{rmq}
	Rigoureusement, c'est si, et seulement si, l'ensemble $\left\{u,v\right\}$ est égal à l'ensemble des solutions de l'équation $z^2-Sz+P=0$ d'inconnue $z\in\C$.
\end{rmq}
\begin{met}[Résoudre une équation polynomiale de degré $>2$]
	Pour résoudre une équation du type $P(z)=0$ d'inconnue $z\in\C$ avec $P$ un polynôme à coefficients complexes de degré strictement supérieure à deux, voici la méthode : 
	\begin{itemize}
		\item On teste des valeurs simples à la recherche d'une racine $\lambda\in\C$. Dans ce cas, on rappelle que $P(z)$ est factorisable par $z-\lambda$.
		\item Cela nous permet de nous ramener à une équation de degré inférieur (strictement). On itère jusqu'à arriver à une équation de degré $2$, puis on applique les théorèmes qui précèdent.
	\end{itemize}
\end{met}
\section{Application à la géométrie}
Dans cette partie, on considère le plan affine usuel $\mathcal{P}$, dit "plan des points". On le munit d'un repère orthonormal direct $(O;\overrightarrow{e_1},\overrightarrow{e_2})$. On note alors $\overrightarrow{\mathcal{P}}$ le "plan des vecteurs" correspondant.
\subsection{Affixe}
\begin{dfn}[Affixe d'un point]
	Soit $M\in\mathcal{P}$ de coordonnées $(x,y)$ dans notre repère. L'\textbf{affixe de $M$} est le complexe : $$z_M:=x+iy$$
\end{dfn}
\begin{dfn}[Affixe d'un vecteur]
	Soit $\overrightarrow{u}$ un vecteur de coordonnées $\displaystyle\begin{pmatrix}a\\ b\end{pmatrix}$ dans notre repère. L'\textbf{affixe de $\overrightarrow{u}$} est le complexe : $$z_{\overrightarrow{u}}:=a+ib$$
\end{dfn}
\begin{exa}
	Soit $(A,B)\in\mathcal{P}^2$. On a : $z_{\overrightarrow{AB}}=z_B-z_A$.
\end{exa}
\begin{prop}
	Soit $\overrightarrow{u},\overrightarrow{v}\in\overrightarrow{\mathcal{P}}$ et $\lambda\in\R$. On a : $$\begin{cases}
		z_{\overrightarrow{u}+\overrightarrow{v}}=z_{\overrightarrow{u}}+z_{\overrightarrow{v}} \\ z_{\lambda\overrightarrow{u}}=\lambda z_{\overrightarrow{u}}
	\end{cases}$$
\end{prop}
\begin{rmq}
	Attention, $\lambda$ \textbf{doit} être un \textbf{réel} !
\end{rmq}
\begin{prop}[Relation de Chasles]
	Soit $(A,B,C)\in\mathcal{P}^3$. On a $z_{\overrightarrow{AB}+\overrightarrow{BC}}=z_{\overrightarrow{AC}}$, soit la \textbf{relation de Chasles} $\overrightarrow{AB}+\overrightarrow{BC}=\overrightarrow{AC}$.
\end{prop}
\subsection{Interprétation géométrique du module}
\begin{prop}
	Soit $\overrightarrow{u}\in\overrightarrow{\mathcal{P}}$. On a : $\lvert z_{\overrightarrow{u}}\rvert =\lVert\overrightarrow{u}\rVert$.
\end{prop}
\begin{exa}
	Pour tout $M\in\mathcal{P}$, $\lvert z_M\rvert=OM$.
\end{exa}
\begin{dfn}[Cercle]
	Soit $\Omega\in\mathcal{P}$ et $\R\geqslant0$. Le \textbf{cercle de centre $\Omega$ et de rayon $R$} est l'ensemble : $$\left\{M(z)\in\mathcal{P}: \lvert z-z_{\Omega}\rvert =R\right\}$$ On dit qu'il "a pour équation" $\lvert z-z_{\Omega}\rvert=R$.
\end{dfn}
\begin{dfn}[Disque fermé]
	Soit $\Omega\in\mathcal{P}$ et $\R\geqslant0$. Le \textbf{disque fermé de centre $\Omega$ et de rayon $R$} est l'ensemble : $$\left\{M(z)\in\mathcal{P}: \lvert z-z_{\Omega}\rvert \leqslant R\right\}$$ On dit qu'il "a pour équation" $\lvert z-z_{\Omega}\rvert\leqslant R$.
\end{dfn}
\begin{dfn}[Disque ouvert]
	Soit $\Omega\in\mathcal{P}$ et $\R\geqslant0$. Le \textbf{disque ouvert de centre $\Omega$ et de rayon $R$} est l'ensemble : $$\left\{M(z)\in\mathcal{P}: \lvert z-z_{\Omega}\rvert <R\right\}$$ On dit qu'il "a pour équation" $\lvert z-z_{\Omega}\rvert<R$.
\end{dfn}
\subsection{Interprétation géométrique des arguments}
\begin{thm}
	Soit $\overrightarrow{u},\overrightarrow{v}$ des vecteurs \textbf{non nuls}. Une mesure particulière de l'angle orienté $(\widehat{\overrightarrow{u},\overrightarrow{v}})$ est donnée par $\displaystyle\arg\left(\frac{z_{\overrightarrow{v}}}{z_{\overrightarrow{v}}}\right)$ où $\arg$ désigne l'argument principal. Formellement : $$(\widehat{\overrightarrow{u},\overrightarrow{v}}) \equiv \arg\left(\frac{z_{\overrightarrow{v}}}{z_{\overrightarrow{v}}}\right) \ [2\pi]$$ Autrement dit, les mesures de $(\widehat{\overrightarrow{u},\overrightarrow{v}})$ sont exactement les arguments de $\displaystyle\frac{z_{\overrightarrow{v}}}{\overrightarrow{v}}$.
\end{thm}
\begin{cor}[Formule de l'emmerdement maximal]
	Soit $(A,B,C,D)\in\mathcal{P}^4$. Alors, une mesure de $(\widehat{\overrightarrow{AB},\overrightarrow{CD}})$ est donnée par $$\arg\left(\frac{z_D-z_C}{z_B-z_A}\right)$$
\end{cor}
\begin{rmq}[Origine du nom du théorème]
	Le nom quelque peu vulgaire de ce théorème provient de mon professeur de terminale et permet de retenir la formule : les lettres sont à l'envers par rapport à l'ordre de départ, et il y a des signes $-$.
\end{rmq}
\begin{exa}
	Soit $M\in\mathcal{P}$. Alors $\arg(z_M)$ est une mesure de $(\widehat{\overrightarrow{e_1},\overrightarrow{OM}})$.
\end{exa}
\begin{prop}[Relation de Chasles pour les mesures d'angles orientés]
	Soit $\overrightarrow{u},\overrightarrow{v}$ et $\overrightarrow{w}$ des vecteurs non nuls. On peut écrire symboliquement : $$(\widehat{\overrightarrow{u},\overrightarrow{v}})+(\widehat{\overrightarrow{v},\overrightarrow{w}})\equiv (\widehat{\overrightarrow{u},\overrightarrow{w}})\ [2\pi]$$
\end{prop}
\begin{dfn}[Vecteurs colinéaires]
	Soit $\overrightarrow{u},\overrightarrow{v}\in\overrightarrow{\mathcal{P}}$. On dira que $\overrightarrow{u}$ et $\overrightarrow{v}$ sont \textbf{colinéaires} lorsque : \begin{itemize}
		\item $\overrightarrow{u}$ ou $\overrightarrow{v}$ est le vecteur nul
		\item $\overrightarrow{u}$ et $\overline{v}$ sont non nuls et $(\widehat{\overrightarrow{u},\overrightarrow{v}})\equiv 0\ [2\pi]$
	\end{itemize}
\end{dfn}
\begin{rmq}
	La définition donnée coïncide avec les définitions analogues à celles données pour les nombres complexes.
\end{rmq}
\begin{dfn}[Vecteurs orthogonaux]
	Soit $\overrightarrow{u},\overrightarrow{v}\in\overrightarrow{\mathcal{P}}$. On dira que $\overrightarrow{u}$ et $\overrightarrow{v}$ sont \textbf{orthogonaux} lorsque : \begin{itemize}
		\item $\overrightarrow{u}$ ou $\overrightarrow{v}$ est le vecteur nul
		\item $\overrightarrow{u}$ et $\overline{v}$ sont non nuls et $(\widehat{\overrightarrow{u},\overrightarrow{v}})\equiv \displaystyle-\frac{\pi}{2} \ \ \text{ou}\ \ \frac{\pi}{2} \ [2\pi]$
	\end{itemize}
\end{dfn}
\begin{thm}[CNS de colinéarité et d'orthogonalité]
	Soit $\overrightarrow{u},\overrightarrow{v}\in\overrightarrow{\mathcal{P}}$. On a : \begin{itemize}
		\item $\overrightarrow{u}$ et $\overrightarrow{v}$ sont colinéaires si, et seulement si, $\overline{z_{\overrightarrow{u}}}\times z_{\overrightarrow{v}}\in\R$
		\item $\overrightarrow{u}$ et $\overrightarrow{v}$ sont orthogonaux si, et seulement si, $\overline{z_{\overrightarrow{u}}}\times z_{\overrightarrow{v}}\in i\R$
	\end{itemize}
\end{thm}
\begin{rmq}
	En travaillant dans le plan complexe, on peut prouver que les expressions du cosinus, du sinus et de la tangente données en $3^e$ pour le triangle rectangle sont bien valables au vu de nos définitions actuelles de ces fonctions. Il suffit de placer le sommet à l'angle droit au centre du plan complexe et d'utiliser les formules sur les angles, puis d'utiliser les définitions de cosinus et sinus.
\end{rmq}
\subsection{Quelques transformations du plan}
Dans toute cette partie, on exclut les homothéties de rapport $0$.
\begin{dfn}[Translation]
	Soit $\overrightarrow{u}\overrightarrow{\mathcal{P}}$. Soit $M\in\mathcal{P}$. Alors il existe un unique point $M'\in\mathcal{P}$ tel que $\overrightarrow{MM'}=\overrightarrow{u}$. La \textbf{translation de vecteur $\overrightarrow{u}$} est alors la fonction : $$\begin{array}{ccccc}t_{\overrightarrow{u}}&:&\mathcal{P}&\rightarrow&\mathcal{P}\\&&M&\mapsto&M'\end{array}$$ Soit $M(z)\in\mathcal{P}$ et $M'(z')\in\mathcal{P}$ son image. On a : $z'=z+z_{\overrightarrow{u}}$
\end{dfn}
\begin{dfn}[Homothétie]
	Soit $\Omega\in\mathcal{P}$ et $\lambda\in\R\st$. Soit $M\in\mathcal{P}$. Alors il existe un unique point $M'\in\mathcal{P}$ tel que $\overrightarrow{\Omega M'}=\lambda\overrightarrow{\Omega M}$. L'\textbf{homothétie de centre $\Omega$ et de rapport $\lambda$} est alors la fonction : $$\begin{array}{ccccc}h_{\Omega,\lambda}&:&\mathcal{P}&\rightarrow&\mathcal{P}\\&&M&\mapsto&M'\end{array}$$ Soit $M(z)\in\mathcal{P}$ et $M'(z')\in\mathcal{P}$ son image. On a : $z'-z_{\Omega}=\lambda(z-z_{\Omega})$
\end{dfn}
\begin{dfn}[Rotation]
	Soit $\Omega\in\mathcal{P}$ et $\theta\in\R\st$. Soit $M\in\mathcal{P}$. Alors il existe un unique point $M'\in\mathcal{P}$ tel que $\Omega M'=\Omega M$ et $(\widehat{\overrightarrow{\Omega M},\overrightarrow{\Omega M'}})\equiv\theta\ [2\pi]$. La \textbf{rotation de centre $\Omega$ et d'angle $\theta$} est alors la fonction : $$\begin{array}{ccccc}r_{\Omega,\theta}&:&\mathcal{P}&\rightarrow&\mathcal{P}\\&&M&\mapsto&M'\end{array}$$ Soit $M(z)\in\mathcal{P}$ et $M'(z')\in\mathcal{P}$ son image. On a : $z'-z_{\Omega}=e^{i\theta}(z-z_{\Omega})$
\end{dfn}
\begin{rmq}[Application affine, partie linéaire, HP]
	Soit $f$ une des trois transformations précédentes et $\overrightarrow{v}\in\overrightarrow{\mathcal{P}}$ Pour savoir comment $f$ transforme $\overrightarrow{v}$, on voudrait écrire $f(\overrightarrow{v})$. Problème : $\overrightarrow{v}\notin\mathcal{P}$ A défaut, pourrait-on créer $\varphi :\overrightarrow{\mathcal{P}}:\rightarrow\overrightarrow{\mathcal{P}}$ qui serait le pendant de $f$ côté $\overrightarrow{\mathcal{P}}$ ? Pour cela, il faudrait : \begin{itemize}
		\item Fixer $(A,B)\in\mathcal{P}^2$ tel que $\overrightarrow{AB}=\overrightarrow{v}$
		\item Calculer $A':=f(A)$ et $B':=f(B)$
		\item Poser $\varphi(\overrightarrow{v}):=\overrightarrow{A'B'}$ et vérifier que cette définition ne dépend pas du couple $(A,B)$ choisi.
	\end{itemize}
	On appelle alors $\varphi$ la \textbf{partie linéaire} de $f$ et on dit que $f$ est une \textbf{application affine}.
\end{rmq}
\begin{dfn}[Similitude directe]
	Soit $(a,b)\in\C\st\times\C$. Une \textbf{similitude directe} est une application de la forme : $$\begin{array}{ccc}\mathcal{P}&\rightarrow&\mathcal{P}\\ M(z)&\mapsto&M'(az+b)\end{array}$$
\end{dfn}
\begin{exa}
	Les translations, homothéties et rotations sont des cas particuliers similitudes directes.
	\begin{itemize}
		\item Translation : $a=1$ et $b=z_{\overrightarrow{u}}$
		\item Homothétie : $a=\lambda$ et $b=z_{\Omega}-\lambda z_{\Omega}$
		\item Rotation : $a=e^{i\theta}$ et $b=z_{\Omega}-e^{i\theta}z_{\Omega}$
	\end{itemize}
\end{exa}
\begin{lem}
	Soit $h$ une homothétie de rapport $\lambda\in\R\st$ et $r$ une rotation d'angle $\theta\in\R$ de \textbf{même centre $\Omega$}. Alors, $h$ et $r$ commutent, c'est-à-dire que $$\forall M\in\mathcal{P},\ h(r(M))=r(h(M))$$
\end{lem}
\begin{dfn}
	$h\circ r=r\circ h$ s'appelle la \textbf{similitude directe de centre $\Omega$, de rapport $\lambda$ et d'angle $\theta$}.
\end{dfn}
\begin{met}[Étude d'une similitude directe]
	Soit $(a,b)\in\C\st\times\C$. Soit $s$ la similitude directe d'équation $z'=az+b$. Voici comment étudier la similitude $s$.
	\begin{itemize}
		\item Premier cas : $a=1$. On reconnaît une translation selon un vecteur d'affixe $b$.
		\item Second cas : $a\ne1$. 
		\begin{itemize}
			\item[--] On cherche un point fixe. Soit $M(z)\in\mathcal{P}$. On a $s(M)=M$ ssi $\displaystyle z=\frac{b}{1-a}$. Notons $\Omega$ le point d'affixe $\displaystyle \frac{b}{1-a}$.
			\item[--] Soit $M(z)\in\mathcal{P}$ et $M'(z')\in\mathcal{P}$ son image par $s$. On a $$\begin{cases}
				z'=az+b \\ z_{\Omega}=az_{\Omega}+b
			\end{cases}$$ donc $z'-z_{\Omega}=a(z-z_{\Omega})$
			\item[--] On écrit $a$ sous forme trigonométrique $a=\lambda e^{i\theta}$ et on obtient $z'-z_{\Omega}=\lambda e^{i\theta}(z-z_{\Omega})$. Ainsi, $s$ est la similitude directe de centre $\Omega$, de rapport $\lambda$ et d'angle $\theta$
		\end{itemize}
	\end{itemize}
\end{met}
\begin{rmq}
	Dans le lemme qui précède, on avait vu un cas particulier de similitudes directes. Désormais, le point méthode montre que, mises à part les translations, on est toujours dans le cas $h\circ r=r\circ h$. On obtient le centre avec le point fixe et le rapport et la l'angle avec la forme trigonométrique.
\end{rmq}
\begin{rmq}[Mise en œuvre pratique du point méthode]
	En pratique, on pourra ne pas appliquer la deuxième étape du second cas du point méthode et directement passer à la troisième étape après la première, puisque l'on sait que la deuxième étape fournit toujours le même résultat.
\end{rmq}
\begin{exa}[L'ensemble des similitudes directes est stable par composition]
	Donnons-nous deux similitudes directes $s_1$ et $s_2$ d'équations respectives $z'=a_1z+b_1$ et $z'=a_2z+b_2$. Étudions $s_1\circ s_2$. Soit $M\in\mathcal{P}$. On a : $z_{(s_1\circ s_2)(M)}=a_1a_2z_M+(a_1b_2+b_1)$ Or, $(a_1a_2,a_1b_2+b_1)\in\C\st\times\C$, donc $s_1\circ s_2$ est une similitude directe. Attention, on a $z_{(s_2\circ s_1)(M)}=a_2a_1z_M+(a_2b_1+b_2)$ donc les similitudes n'ont aucune raison de commuter : les deux composées ont même rapport et même angle, mais probablement pas même centre. 
\end{exa}
\begin{dfn}[Symétries d'axes $(O;\overrightarrow{e_1})$ et $(O;\overrightarrow{e_2})$]
	La symétrie d'axe $(O;\overrightarrow{e_1})$ a pour équation $z'=\overline{z}$. La symétrie d'axe $(O;\overrightarrow{e_2})$ a pour équation $z'=-\overline{z}$.
\end{dfn}
\begin{dfn}[Similitude indirecte, HP]
	La composée d'une similitude directe et d'une symétrie axiale s'appelle une \textbf{similitude indirecte}.
\end{dfn}
\section{Formulaire de trigonométrie circulaire}
Voici un formulaire de trigonométrie circulaire à connaître sur le bout de doigts. On rappelle que s'en déduit le formulaire de trigonométrie hyperbolique par la "règle d'Osborn" (cf. chapitre "Fonctions usuelles", paragraphe "$3.2$ Trigonométrie hyperbolique").
$$\begin{cases}
	\cos(-\theta)=\cos(\theta)\\
	\sin(-\theta)=-\sin(\theta)\\
\end{cases}$$
$$\begin{cases}
	\cos(\pi-\theta)=-\cos(\theta)\\
	\sin(\pi-\theta)=\sin(\theta)\\
	\cos(\pi+\theta)=-\cos(\theta)\\
	\sin(\pi+\theta)=-\sin(\theta)\\
\end{cases}$$
$$\begin{cases}
	\displaystyle\cos\left(\frac{\pi}{2}-\theta\right)=\sin(\theta)\\[3mm]
	\displaystyle\sin\left(\frac{\pi}{2}-\theta\right)=\cos(\theta)\\[3mm]
	\displaystyle\cos\left(\frac{\pi}{2}+\theta\right)=-\sin(\theta)\\[3mm]
	\displaystyle\sin\left(\frac{\pi}{2}+\theta\right)=\cos(\theta)
\end{cases}$$
$$\begin{cases}
	\cos(a+b)=\cos(a)\cos(b)-\sin(a)\sin(b)\\
	\cos(a-b)=\cos(a)\cos(b)+\sin(a)\sin(b)\\
	\sin(a+b)=\sin(a)\cos(b)+\cos(a)\sin(b)\\
	\sin(a-b)=\sin(a)\cos(b)-\cos(a)\sin(b)
\end{cases}$$
$$\begin{cases}
	\displaystyle\cos(a)\cos(b)=\frac{1}{2}\left(\cos(a-b)+\cos(a+b)\right)\\[3mm]
	\displaystyle\sin(a)\sin(b)=\frac{1}{2}\left(\cos(a-b)-\cos(a+b)\right)\\[3mm]
	\displaystyle\sin(a)\cos(b)=\frac{1}{2}\left(\sin(a-b)+\sin(a+b)\right)
\end{cases}$$
$$\begin{cases}
	\displaystyle\cos(p)+\cos(q)=2\cos\left(\frac{p+q}{2}\right)\cos\left(\frac{p-q}{2}\right)\\[3mm]
	\displaystyle\cos(p)-\cos(q)=-2\sin\left(\frac{p+q}{2}\right)\sin\left(\frac{p-q}{2}\right)\\[3mm]
	\displaystyle\sin(p)+\sin(q)=2\sin\left(\frac{p+q}{2}\right)\cos\left(\frac{p-q}{2}\right)\\[3mm]
	\displaystyle\cos(p)-\cos(q)=2\cos\left(\frac{p+q}{2}\right)\sin\left(\frac{p-q}{2}\right)
\end{cases}$$
$$\begin{cases}
	\cos(2a)=2\cos(a)^2-1\\
	\cos(2a)=\cos(a)^2-\sin(a)^2\\
	\cos(2a)=1-2\sin(a)^2\\
	\sin(2a)=2\sin(a)\cos(a)
\end{cases}$$
$$\begin{cases}
	\displaystyle\cos(a)^2=\frac{1+\cos(2a)}{2}\\[3mm]
	\displaystyle\sin(a)^2=\frac{1-\cos(2a)}{2}
\end{cases}$$
$$\begin{cases}
	\displaystyle\tan(a+b)=\frac{\tan(a)+\tan(b)}{1-\tan(a)\tan(b)}\\[3mm]
	\displaystyle\tan(a-b)=\frac{\tan(a)-\tan(b)}{1+\tan(a)\tan(b)}\\[3mm]
	\displaystyle\tan(2a)=\frac{2\tan(a)}{1-\tan(a)^2}
\end{cases}$$
$$t:=\displaystyle\tan\left(\frac{\theta}{2}\right)\implies\begin{cases}
	\displaystyle\cos(\theta)=\frac{1-t^2}{1+t^2}\\[3mm]
	\displaystyle\sin(\theta)=\frac{2t}{1+t^2}\\[3mm]
	\displaystyle\tan(\theta)=\frac{2t}{1-t^2}
\end{cases}$$
\section{Compléments : retour sur la notion d'angle orienté, HP}
La définition que le lycée donne de la mesure d'un angle orientée peut paraître un peu floue. De fait, elle est plus difficile à mettre en place qu'il n'y paraît. C'est ce à quoi nous nous attelons maintenant, avant de démontrer le théorème qui donne une mesure d'un angle en fonction de l'argument d'un certain complexe. \\ \par
Considérons le cercle unité, c'est-à-dire le cercle de centre O et de rayon 1. Nous avons vu précédemment que tout point du cercle unité pouvait s'écrire $z = e^{i\theta}$. De plus, nous avons vu que la fonction $\theta \mapsto e^{i\theta}$ était continue, $2\pi$-périodique, et nous avons étudié le sens de variation de $\theta \mapsto \operatorname{Re}(e^{i\theta}) = \cos(\theta)$ et $\theta \mapsto \operatorname{Im}(e^{i\theta}) = \sin(\theta)$ sur $[0, 2\pi[$. De tout cela, on déduit que lorsque $\theta$ croît de 0 à $2\pi$, le point d'affixe $e^{i\theta}$ parcourt (continûment) le cercle trigonométrique en partant de 1 dans le "sens inverse des aiguilles d'une montre", aussi appelé sens direct.

\begin{thm}
	Un cercle de rayon $r$ a pour périmètre $2\pi r$.
\end{thm}
\begin{proof}
	Posons un repère au centre du cercle. Par définition, le périmètre du cercle est la limite de la longueur d'une ligne brisée de plus en plus fine le long de sa circonférence. Considérons donc $n$ sommets ($n \geqslant 2$) équirépartis (au sens de la distance) autour du cercle, avec les affixes $e^{i\frac{2k\pi}{n}} (0 \leqslant k \leqslant n-1)$. La longueur de la ligne brisée vaut $$\sum_{k=0}^{n-1} \left| re^{i\frac{2(k+1)\pi}{n}} - re^{i\frac{2k\pi}{n}} \right| = \sum_{k=0}^{n-1} 2r \sin\left(\frac{\pi}{n}\right) = 2rn \sin\left(\frac{\pi}{n}\right)$$
	On remarque ensuite que $$2rn \sin\left(\frac{\pi}{n}\right) = 2\pi r \cdot \frac{\sin(\frac{\pi}{n})}{\frac{\pi}{n}} \xrightarrow[n \to +\infty]{} 2\pi r \cdot \sin'(0) = 2\pi r$$
Ainsi la preuve est achevée.
\end{proof}
Maintenant, on va prouver le lien entre angle orienté et argument. Mais au juste, que signifie "mesurer" un angle orienté ? Tout commence avec la notion d'angle géométrique que l'on apprend au collège, voire en primaire : il s'agit d'un couple de deux demi-droites issues du même sommet. Ensuite, comment définit-on la mesure d'un angle géométrique lorsqu'on est écolier ? Intuitivement, c'est "quelque chose qui se mesure avec un rapporteur". Et un rapporteur, ce n'est ni plus ni moins qu'un cercle dont le pourtour est "régulièrement" gradué, au sens de la circonférence. On en arrive donc à la définition suivante :

\begin{dfn}
	Soit $(AB)$ et $(AC)$ deux demi-droites issues de A. Quel que soit le cercle de centre A, les demi-droites le coupent en deux arcs. La longueur du petit arc (au sens large) divisée par la longueur du rayon est constante (comme on s'en rend compte en appliquant une homothétie). C'est une grandeur sans dimension comprise entre 0 et $\pi$, que l'on appelle la mesure de l'angle $\widehat{BAC}$ (en radians). On peut partir du cercle unité, ce qui donne directement la mesure de l'angle sans avoir à diviser par le rayon.
\end{dfn}

\begin{rmq}
	Dans le cas d'un angle plat, les deux arcs de cercle ont la même longueur, donc peu importe le choix de l'arc, on obtiendra la même mesure.
\end{rmq}
\begin{exa}
Un angle droit vaut $\frac{\pi}{2}$ radians.
\end{exa}
\begin{dfn}
	On se donne deux vecteurs $\overrightarrow{u}$ et $\overrightarrow{v}$ non nuls. Pour mesurer l'angle orienté $(\overrightarrow{u}, \overrightarrow{v})$, on mesure l'angle géométrique qu'ils forment, c'est-à-dire qu'on calcule la longueur du petit arc de cercle trigonométrique qui joint $\overrightarrow{u}$ et $\overrightarrow{v}$. On obtient un réel $\theta \in [0, \pi]$. Puis, selon que l'on tourne dans le sens direct ou indirect pour aller de $\overrightarrow{u}$ à $\overrightarrow{v}$ le long de l'arc, on pose $\alpha = \pm \theta$. On obtient $\alpha \in [-\pi, \pi]$, et comme dans le cas de l'argument, on décide que la mesure d'un angle orienté est définie à $2\pi$ près. On dit donc que $\alpha$ est une mesure de l'angle.
\end{dfn}
\begin{rmq}
	À vrai dire, pour aller de $\overrightarrow{u}$ à $\overrightarrow{v}$, il est possible de tourner dans les deux sens. Supposons que l'angle géométrique associé au petit arc vaut $\theta$ et que l'angle orienté mesure $\alpha = \pm \theta$. Sans perte de généralité, mettons $\alpha = \theta$ pour fixer les idées. \par
On peut généraliser la mesure d'un angle géométrique aux angles plus grands que $\pi$, en convenant que l'angle géométrique associé au grand arc mesure $\theta' = 2\pi - \theta$. En passant aux angles orientés, on remarque que le sens de parcours a été inversé. On obtient donc une nouvelle mesure $\alpha' = -(2\pi - \theta) = \theta - 2\pi$. Ainsi, on retrouve bien la même mesure à $2\pi$ près.\par
Et même possible de faire plus d'un tour pour aller de $\overrightarrow{u}$ à $\overrightarrow{v}$. Dans tous les cas, on trouvera un résultat de la forme $\theta + 2k\pi (k \in \mathbb{Z})$.
\end{rmq}

Rappelons que la mesure géométrique d'un angle n'est rien d'autre que la longueur d'arc parcourue sur le cercle trigonométrique. Si on convient de noter la longueur avec un signe moins lorsqu'elle est parcourue dans le sens indirect, on obtient la
\begin{dfn}[Définition équivalente]
	Une mesure de l'angle orienté $(\widehat{\overrightarrow{u},\overrightarrow{v}})$ est la distance (comptée en valeur algébrique) d'un arc de cercle trigonométrique qui joint $\overrightarrow{u}$ à $\overrightarrow{v}$, y compris en faisant plusieurs tours.
\end{dfn}

On peut alors prouver le
\begin{thm}
	Donnons-nous un angle orienté du plan $(\widehat{AB,CD})$ (avec $A \neq B$ et $C \neq D$). Alors $\displaystyle\arg\left(\frac{z_{D}-z_{C}}{z_{B}-z_{A}}\right)$ est une mesure de $(\widehat{AB,CD})$
\end{thm}
\begin{proof}
	Pour commencer, on effectue les deux opérations suivantes.
	\begin{itemize}
		\item On ramène $\overrightarrow{AB}$ et $\overrightarrow{CD}$ sur l'origine (c'est-à-dire qu'on suppose sans perte de généralité que $O=A=C$). Ainsi, on se ramène à calculer $\arg(z_{D}/z_{B})$.
		\item Et on renormalise nos deux vecteurs (c'est-à-dire qu'on suppose $OB=OD=1$). Pour cela, il suffit de remplacer $z_{B}$ par $z_{B}/|z_{B}|$ et $z_{D}$ par $z_{D}/|z_{D}|$. Cela ne change ni l'angle orienté $(\displaystyle\widehat{\overrightarrow{OB},\overrightarrow{OD}})$ ni la valeur de $\arg(z_{D}/z_{B})$.
	\end{itemize}
	On obtient ainsi deux complexes $z_{B}=e^{i\theta_{1}}$ $z_{D}=e^{i\theta_{2}}$ sur le cercle trigonométrique. Quitte à ajouter un multiple de $2\pi$ à $\theta_{2}$, ce qui ne changera pas la position de $D$ et laissera donc l'angle $(\widehat{\overrightarrow{OB},\overrightarrow{OD}})$ inchangé, on suppose $\theta_{1}\leq\theta_{2}<\theta_{1}+2\pi$. Ainsi, on peut compter la longueur de l'arc de cercle positivement. Puis comme précédemment, il s'agit de calculer la limite d'une ligne brisée qui joint $e^{i\theta_{1}}$ à $e^{i\theta_{2}}$.
	\begin{align*}
		\sum_{k=0}^{n-1}\left\lvert e^{i\left(\theta_1+\frac{(k+1)(\theta_2-\theta_1)}{n}\right)}-e^{i\left(\theta_1+\frac{k(\theta_2-\theta_1)}{n}\right)}\right\rvert&= \sum_{k=0}^{n-1} 2 \left\lvert\sin\left(\frac{\theta_{2}-\theta_{1}}{2n}\right)\right\rvert \\
		&= 2n \left\lvert\sin\left(\frac{\theta_{2}-\theta_{1}}{2n}\right)\right\rvert \\
		&= 2n\sin\left(\frac{\theta_{2}-\theta_{1}}{2n}\right) \\
		&\xrightarrow[n\to\infty]{} \theta_{2}-\theta_{1}
	\end{align*}
	Donc une mesure de l'angle orienté est donnée par $\theta_{2}-\theta_{1}$ et par ailleurs on a bien
	$$\arg(z_{D}/z_{B}) = \arg\left(e^{i(\theta_{2}-\theta_{1})}\right) \equiv \theta_{2}-\theta_{1}\ [2\pi]$$
	Ainsi la preuve est achevée.
\end{proof}

\newpage
\chapter{Espaces vectoriels} %chapitre complet%
Nous considérons dans tout le chapitre un corps $\K$. En pratique, ce corps sera souvent égal à $\R$ ou $\C$, mais la quasi-totalité des résultats restent vrais pour un corps quelconque.
\section{Premières définitions} %paragraphe complet%
Au lycée, on apprend que $\R^2$ ou $\R^3$ peut être vu comme un ensemble de $vecteurs$, sur lesquels on peut effectuer des $combinaisons$ $linéaires$. Plus généralement, nous avons la 
\begin{dfn}[$\K$-espace vectoriel]
	Un \textbf{$\K$-espace vectoriel} ou \textbf{$\K$-ev} est un ensemble $E$ muni d'une loi de composition interne $+$ et d'une loi de composition externe :
	$$\begin{array}{ccc}
		\K\times E&\rightarrow&E \\
		(\lambda, x)&\mapsto&\lambda x
	\end{array}$$ vérifiant les axiomes suivants. \begin{itemize}
	\item $(E,+)$ est un groupe commutatif.
	\item On a une pseudo-distributivité : $$\left \{\begin{array}{ll}\forall(\lambda,x,y)\in\K\times E^2,\ \lambda(x+y)=\lambda x+\lambda y\\ \forall(\lambda,\mu,x)\in\K^2\times E,\ (\lambda+\mu)x=\lambda x+\mu x\end{array}\right.$$
	\item On a une pseudo-associativité : $$\forall(\lambda,\mu,x)\in\K^2\times E,\ (\lambda\mu)x=\lambda(\mu x)$$
	\item On a une propriété d'opérateur neutre : $$\forall x\in E,\ 1x=x$$ 
	\end{itemize}
	Les éléments de $E$ sont appelés les \textbf{vecteurs}, et ceux de $\K$ les \textbf{scalaires}.
\end{dfn}
\begin{rmq}
	À part dans certains cas lorsqu'il pourra y avoir ambiguïté, on cessera de noter les vecteurs avec des flèches.
\end{rmq}
\begin{rmq}
	Si $E$ est un $\mathbb{C}$-ev, alors $a$ $fortiori$ c'est un $\mathbb{R}$-ev, et même un $\mathbb{Q}$-ev.
\end{rmq}
\begin{rmq}[$\Z$-module, HP]
	Ces quatre axiomes ont déjà été vus lorsque nous avons introduit la loi externe $n \cdot x$ dans un anneau $A$ (avec $n \in \mathbb{Z}$ et $x \in A$). Plus généralement, ils restent vrais sur $\mathbb{Z} \times E$ dès que $(E, +)$ est un groupe commutatif. La seule différence est que $\mathbb{Z}$ n'est pas un corps, mais un anneau. C'est pour cela qu'on ne parle pas de $\mathbb{Z}$-ev mais de $\mathbb{Z}$-module.
\end{rmq}
\begin{dfn}[Combinaison linéaire]
	Soit $x_1, ..., x_n$ des vecteurs de $E$. Une \textbf{combinaison linéaire} des $x_i$ est une expression de la forme $\displaystyle\sum_{i=1}^{n} \lambda_i x_i$,
	où les $\lambda_i$ sont des scalaires.
\end{dfn}
\begin{rmq}[Linéarité de la somme]
	Par une récurrence immédiate, on montre que la pseudo-distributivité se généralise :$$	\lambda \sum_{i=1}^{n} x_i + \mu \sum_{i=1}^{n} y_i = \sum_{i=1}^{n} (\lambda x_i + \mu y_i)$$
\end{rmq}
On montre tout de suite quelques propriétés élémentaires.
\begin{exa}
	Si $x = \lambda y$ avec $\lambda \neq 0$, alors $\displaystyle y = \frac{1}{\lambda}x$
\end{exa}
\begin{prop}[Pseudo-absorbance, pseudo-intégrité]
	$\forall (\lambda, x) \in \mathbb{K} \times E,\ \lambda x = 0 \Leftrightarrow \lambda = 0_{\mathbb{K}}$ ou $x = 0_{E}$.
\end{prop}
\begin{cor}
	$\forall (\lambda, x) \in \mathbb{K} \times E,\ (-\lambda)x = \lambda(-x) = -(\lambda x)$
\end{cor}
\begin{exa}
	Si $\lambda u=\mu u$ avec $u\ne0_E$, alors $\lambda=\mu$.
\end{exa}
\begin{rmq}[Itération de la loi $+$ et l.c.e.]
	On pourra montrer le résultat suivant:$$	\forall (n, \lambda, x) \in \mathbb{Z} \times \mathbb{K} \times E,\	(n\lambda)x = \lambda(nx) = n(\lambda x)$$ En particulier, on pourra écrire $n\lambda x$ sans risque d'être ambigu. De même, si $\mathbb{K}$ est un sur-ensemble de $\mathbb{Z}$, on pourra vérifier que la notation $\lambda x$ désigne indifféremment l'itération du $+$ de $E$, ou la l.c.e.
\end{rmq}
\begin{dfn}[Colinéarité]
	Soit $x, y$ deux vecteurs. Les trois assertions suivantes sont équivalentes:
	\begin{enumerate}
		\item $\exists \lambda \in \mathbb{K}, y = \lambda x$ ou $\exists \mu \in \mathbb{K}, x = \mu y$
		\item $\exists \lambda \in \mathbb{K}, y = \lambda x$ ou $x = 0$,
		\item $\exists \mu \in \mathbb{K}, x = \mu y$ ou $y = 0$.
	\end{enumerate} Dans ce cas, on dit que $x$ et $y$ sont \textbf{colinéaires}. Bien sûr, par symétrie de définition, $x$ et $y$ sont colinéaires ssi $y$ et $x$ sont colinéaires.
\end{dfn}
\begin{exa}
	Le vecteur nul est colinéaire à tout vecteur.
\end{exa}
\begin{rmq}
	Attention : à moins de travailler avec $\K=\Q$ ou $\R$, cela n'a aucun sens de parler de colinéarité directe.
\end{rmq}
\begin{exa}
	$\K$ lui-même est un $\K$-ev.
\end{exa}
\begin{prop}[Espace vectoriel produit]
	Le produit cartésien \(E_{1}\times E_{2}\) de deux espaces vectoriels $E_{1}$ et $E_{2}$ muni des lois suivantes:
	$$\begin{cases}(x_{1},x_{2})+(x_{1}^{\prime},x_{2}^{\prime})=(x_{1}+x_{1}^{\prime},x_{2}+x_{2}^{\prime})\\ \lambda(x_{1},x_{2})=(\lambda x_{1},\lambda x_{2})\end{cases}$$
est un espace vectoriel.
\end{prop}
\begin{rmq}
	On peut généraliser ce résultat à un produit quelconque (même infini) d'espaces vectoriels.
\end{rmq}
\begin{exa}[Structure canonique de $\K$-ev de $\K^n$]
	Pour $n\in\mathbb{N}^{*}$. $\mathbb{K}^{n}$ est un espace vectoriel pour les lois suivantes.
$$\left(\begin{matrix}x_{1}\\ x_{2}\\ \vdots\\ x_{n}\end{matrix}\right)+\left(\begin{matrix}x_{1}^{\prime}\\ x_{2}^{\prime}\\ \vdots\\ x_{n}^{\prime}\end{matrix}\right)=\left(\begin{matrix}x_{1}+x_{1}^{\prime}\\ x_{2}+x_{2}^{\prime}\\ \vdots\\ x_{n}+x_{n}^{\prime}\end{matrix}\right)$$ et $$\lambda\left(\begin{matrix}x_{1}\\ x_{2}\\ \vdots\\ x_{n}\end{matrix}\right)=\left(\begin{matrix}\lambda x_{1}\\ \lambda x_{2}\\ \vdots\\ x_{n}\end{matrix}\right)$$
Il s'agit simplement d'une application de la remarque précédente et du fait que $\K$ est lui-même un $\K$-ev. On dit qu'on a muni $\mathbb{K}^{n}$ de sa \textbf{structure canonique} d'espace vectoriel.
\end{exa}
\begin{exa}
	Soit X un ensemble quelconque et E un espace vectoriel. Un autre exemple très important est fourni par \(E^{X}\) muni des lois suivantes.
$$\begin{array}{ccccc}f+g&:&X&\mapsto&E\\ &&x&\rightarrow&f(x)+g(x)\end{array}$$ et $$\begin{array}{ccccc}\lambda f&:&X&\mapsto&E\\ &&x&\rightarrow&\lambda f(x)\end{array}$$
Vérifications sur la loi +.
\begin{itemize}
\item L'associativité et la commutativité s'héritent de celles dans \((E,+).\)
\item Elément neutre: c'est l'application \(x\mapsto0,\)
\item Opposé de : c'est l'application \(x\mapsto-f(x),\)
\end{itemize}
Vérifications sur la loi externe.
\begin{itemize}
	\item La pseudo-distributivité et la pseudo-associativité s'héritent de celles dans E.
	\item L'opérateur neutre est $1_{\mathbb{K}}$
\end{itemize}
\end{exa}
\begin{rmq}
	On pourra utiliser largement ce dernier exemple lorsqu'il s'agira de montrer que tel ensemble de fonctions est un espace vectoriel. Bien souvent, il suffira de montrer que c'est un sous-espace vectoriel d'un certain $E^{X}$ ce qui sera toujours plus rapide.
\end{rmq}
\begin{rmq}
	Un cas particulier est donné par $X=\mathbb{N}$. $E^{X}$ est alors l'ensemble des suites à valeurs dans $E$, et c'est donc un espace vectoriel. Pour $X=\llbracket1,n\rrbracket$, on retrouve $\mathbb{K}^{n}$ muni de sa structure canonique.
\end{rmq}
\section{Sous-espaces} %paragraphe complet%
\begin{dfn}[Sous-espace vectoriel]
	Un \textbf{sous-espace vectoriel} ou \textbf{sev} de $E$ est une partie $F \subset E$ telle que
\begin{itemize}
	\item $(F,+)$ est un sous-groupe de $(E,+)$
	\item $\forall (\lambda, x) \in F, \lambda x \in F$ (stabilité par la l.c.e.).
\end{itemize}
\end{dfn}
\begin{rmq}
	Soit $F$ un sev de $E$ et $n \in \mathbb{N}$, par une récurrence immédiate, $F$ est stable par toute combinaison linéaire de $n$ vecteurs.
\end{rmq}
\begin{prop}[Recette pratique pour montrer que $F$ est un sev de $E$]
	Soit $E$ un $K$-ev et $F \subset E$. Pour que $F$ soit un sev de $E$, il faut et il suffit que
\begin{itemize}
	\item $F \neq \emptyset$ (par exemple $0 \in F$)
	\item $\forall (\lambda, \mu, x, y) \in \mathbb{K}^2 \times F^2, \lambda x + \mu y \in F$ (stabilité par combinaison linéaire de deux vecteurs)
\end{itemize}
\end{prop}
\begin{rmq}
	En pratique, on utilise la condition suffisante pour montrer que $F$ est un sev de $E$ avec $E$ un espace vectoriel de référence, ce qui prouve que $F$ un espace vectoriel.
\end{rmq}
\begin{rmq}
	De manière équivalente, le deuxième point de la recette peut être remplacé par $$\forall(\lambda,x,y)\in\K\times F^2,\ \lambda x+y\in F$$
\end{rmq}
\begin{exa}
	Dans le $\mathbb{R}$-ev des suites réelles, les suites bornées forment un sev, à l'intérieur duquel les suites convergentes forment un sev, à l'intérieur duquel les suites qui convergent vers 0 forment un autre sev.
\end{exa}
\begin{prop}
	Soit $\mathcal{F}$ un ensemble de sev de $E$. Alors $\displaystyle \bigcap_{F \in \mathcal{F}} F$ est un sev de $E$.
\end{prop}
\begin{dfn}[Sous-espace vectoriel engendré]
	Soit $A$ une partie de $E$. Alors le \textbf{sous-espace engendré par $A$}, noté $\text{Vect}(A)$, est l'intersection de tous les sev de $E$ contenant $A$. C'est aussi le plus petit sev contenant $A$. Donc si un sev contient $A$, alors il contient $\vect(A)$.
\end{dfn}
\begin{rmq}
	Si $A = \{x\}$, $\text{Vect}(\{x\})$ pourra plus simplement être noté $\text{Vect}(x)$.
\end{rmq}
\begin{exa}
	On a $\text{Vect}(0) = \{0\}$ et $\text{Vect}(E) = E$. Si $F$ est un sev de $E$ on a $\text{Vect}(F) = F$.
\end{exa}
\begin{prop}[$\vect$ d'une partie finie]
	Si $A = \{x_1, ..., x_n\}$, alors $\vect(A)$ est l'ensemble des combinaisons linéaires des éléments de $A$ : $$\vect(A) = \left\{\lambda_1 x_1 + ... + \lambda_n x_n \mid \lambda_1, ..., \lambda_n \in \mathbb{K}\right\}$$
\end{prop}
\begin{exa}
	Voici quelques cas très simples.
	\begin{itemize}
		\item Soit $x \in E$. On a $\vect(x) = \{\lambda x \mid \lambda \in \mathbb{K}\}$. On pourra le noter $\mathbb{K}x$. Si de plus $x \neq 0$, $\vect(x)$ est appelé une \textbf{droite vectorielle}. Géométriquement, on peut le représenter comme une droite qui passe par l'origine.
		\item Soit $(x, y) \in E^2$. On a $\vect(\{x, y\}) = \{\lambda x + \mu y \mid (\lambda, \mu) \in \mathbb{K}^2\}$. On pourra le noter $\mathbb{K}x + \mathbb{K}y$. Si $x$ et $y$ sont non colinéaires, $\vect(\{x, y\})$ est appelé un \textbf{plan vectoriel}. Géométriquement, on peut le représenter comme un plan qui passe par l'origine.
	\end{itemize}
\end{exa}
\section{Applications linéaires}
De même qu'un morphisme de groupe préserve la loi sous-jacente, ou qu'un morphisme d'anneaux ou de corps préserve les lois sous-jacentes, on a la

\begin{dfn}[Application linéaire, ensemble $\mathcal{L}(E,F$)]
	Une \textbf{application linéaire} est un morphisme $f$ d'un espace vectoriel $E$ dans un espace vectoriel $F$, c'est-à-dire qu'elle vérifie	$$\forall (\lambda, \mu, x, y) \in \mathbb{K}^2 \times E^2,\ f(\lambda x + \mu y) = \lambda f(x) + \mu f(y)$$
	Si $f$ est bijective, on rappelle qu'elle s'appelle un \textbf{isomorphisme}. Si $E = F$, on rappelle que $f$ s'appelle un \textbf{endomorphisme}. Et si $f$ est à la fois un isomorphisme et un endomorphisme, c'est un \textbf{automorphisme}. \par L'ensemble des applications linéaires de $E$ dans $F$ est noté $\mathcal{L}(E,F)$. L'ensemble des endomorphismes de $E$ est plus simplement noté $\mathcal{L}(E)$.
\end{dfn}
\begin{rmq}
	De manière équivalente, on peut demander que $f$ vérifie $$\forall (\lambda, x, y) \in \mathbb{K} \times E^2,\ f(\lambda x + y) = \lambda f(x) + f(y)$$
\end{rmq}
\begin{rmq}
	Soit $f \in \mathcal{L}(E,F)$. Par une récurrence immédiate, on vérifie que, plus généralement $$\forall (\lambda_1, ..., \lambda_n) \in \mathbb{K}^n,\ \forall (x_1, ..., x_n) \in E^n,\ f\left(\sum_{i=1}^n \lambda_i x_i\right) = \sum_{i=1}^n \lambda_i f(x_i)$$
\end{rmq}
\begin{exa}
	Prenons $E = F = \mathbb{R}^2$. Chaque vecteur de $\mathbb{R}^2$ s'exprime d'une unique manière comme combinaison linéaire de $e_1 = \begin{pmatrix} 1 \\ 0 \end{pmatrix}$ et $\displaystyle e_2 = \begin{pmatrix} 0 \\ 1 \end{pmatrix}$ : $$x=\lambda e_1+\mu e_2$$
	On en déduit $f(x) = \lambda f(e_1) + \mu f(e_2)$. En particulier, le maillage $\{me_1 + ne_2 \mid (m,n) \in \mathbb{Z}^2\}$ se transforme en $\{mf(e_1) + nf(e_2) \mid (m,n) \in \mathbb{Z}^2\}$. Plus généralement, il faut s'imaginer que $f$ "déforme" continûment un maillage infiniment fin. $f$ ne préserve pas forcément les longueurs ni les angles, mais elle préserve au moins les parallélismes.
\end{exa}
\begin{rmq}
	On a les deux résultats suivants.
	\begin{itemize}
		\item Si $f \in \mathcal{L}(E, F)$ et si on considère $0 \in \mathcal{L}(E, E)$, alors $f \circ 0 = 0$.
		\item Et bien sûr, si $f \in \mathcal{L}(E, F)$ et si on considère $0 \in \mathcal{L}(F, G)$, alors $0 \circ f = 0$ (mais ce dernier résultat n'utilise pas la linéarité).
	\end{itemize}	
\end{rmq}
\begin{rmq}
En particulier, $f$ est un morphisme du groupe $(E,+)$ dans le groupe $(F,+)$. Cela permet de définir $\ker(f)$ et de caractériser l'injectivité de $f$. On en déduit aussi que $$\begin{cases}
	f(0) = 0 \\
	\forall x \in E, f(-x) = -f(x)
\end{cases}$$
(mais de toute façon, on le voit déjà par linéarité).
\end{rmq}
\begin{prop}
Soit $f \in \mathcal{L}(E, F)$. L'image d'un sous-espace vectoriel de $E$ est un sous-espace vectoriel de $F$, et l'image réciproque d'un sous-espace vectoriel de $F$ est un sous-espace vectoriel de $E$.
\end{prop}
\begin{cor}
Soit $f \in \mathcal{L}(E, F)$. Alors $\ker(f)$ et $\im(f)$ sont des sous-espaces vectoriels de $E$ et $F$ respectivement.
\end{cor}
\begin{prop}
$\mathcal{L}(E, F)$ est un $\mathbb{K}$-espace vectoriel. En particulier, $\mathcal{L}(E)$ est un $\mathbb{K}$-espace vectoriel.
\end{prop}
\begin{lem}
Si $f \in \mathcal{L}(F, G)$ et $g \in \mathcal{L}(E, F)$, alors $f \circ g \in \mathcal{L}(E, G)$. De plus, la composition est bilinéaire au sens suivant : $$\begin{cases}
	\forall (\lambda, \mu) \in \mathbb{K}^2,\ \forall (f, g, h) \in \mathcal{L}(F, G)^2 \times \mathcal{L}(E, F),\ (\lambda f + \mu g) \circ h = \lambda (f \circ h) + \mu (g \circ h) \\ \forall (\lambda, \mu) \in \mathbb{K}^2,\ \forall (f, g, h) \in \mathcal{L}(F, G)\times\mathcal{L}(E,F)^2,\ f \circ (\lambda g + \mu h) = \lambda (f \circ g) + \mu (f \circ h)
\end{cases}$$
\end{lem}
\begin{exa}
	En particulier, on a $(\lambda f) \circ g = f \circ (\lambda g) = \lambda(f \circ g)$.
\end{exa}
\begin{prop}
	On pourra montrer qu'on a toujours 
	$$\begin{cases}
	Im(f \circ g) \subset Im(f) \\
	Ker(f \circ g) \supset Ker(g)
	\end{cases}$$
De même, si $f \circ g = 0$ on montrera que
$$Im(g) \subset Ker(f)$$
\end{prop}
\begin{cor}
	$(\mathcal{L}(E), +, \circ)$ est un anneau (non commutatif en général).
\end{cor}
\begin{rmq}
	On s'entraînera avec profit à réaliser directement des calculs sur les endo- morphismes, à la manière des éléments d'un anneau, sans passer par leurs images élément par élément. Dans ce contexte, $f^n$ désignera la n-ième itérée de $f$ au sens de la loi $\circ$. \\ \par Si $f$ et $g$ commutent, on rappelle que la formule du binôme est toujours valable :$$(f+g)^m = \sum_{k=0}^m \binom{m}{k} f^k g^{m-k}$$ ainsi que la formule de Bernoulli :$$f^m - g^m = (f-g)(\sum_{k=0}^{m-1} f^k g^{m-k-1}) = (\sum_{k=0}^{m-1} f^k g^{m-k-1})(f-g)$$ En particulier, on utilisera régulièrement ces formules avec $g = \id$.
\end{rmq}

\begin{exa}[Suites récurrentes linéaires d'ordre 2]
	Considérons les suites $(u_n)_{n \in \mathbb{N}}$ à valeurs complexes qui vérifient une relation de récurrence du type $$u_{n+2} + au_{n+1} + bu_n = 0$$
avec $b \neq 0$. Alors l'ensemble $E$ de ces suites forme un sev de $\mathbb{C}^{\mathbb{N}}$ comme on peut le vérifier directement par le calcul. Mais une autre manière de voir les choses consiste à considérer l'application
$$\begin{array}{ccc}\varphi:\C^\N&\rightarrow&\C^\N \\ (u_n)_{n\in\N}&\mapsto&(u_{n+1})_{n\in\N}\end{array}$$
qui est linéaire. $E$ est alors le noyau de l'application $\varphi^2 + a\varphi + b\id$, qui est linéaire.
\end{exa}
\begin{dfn}[$\K$-algèbre]
	On dit que $(A, +, \cdot, \times)$ est une \textbf{algèbre} sur le corps $\K$, ou $\K$-algèbre lorsque
\begin{itemize}
	\item$(A, +, \cdot)$ est un K-ev;
	\item$(A, +, \times)$ est un anneau:
	\item la deuxième loi de l'anneau est compatible avec la loi externe : $$\lambda(x \times y) = (\lambda x) \times y = x \times (\lambda y)$$
\end{itemize}
\end{dfn}
\begin{exa}
	$(\mathcal{L}(E),+,\cdot,\times)$ est une $\K$-algèbre. $(\K,+,\cdot,\times)$ est une $\K$-algèbre. On verra d'autres exemples de $\K$-algèbres dans les chapitres sur les polynômes et les matrices.
\end{exa}
\begin{dfn}[Sous-algèbre]
	Si A est une algèbre, une \textbf{sous-algèbre} de $A$ sera donc une partie de $A$ qui est à la fois un sev et un sous-anneau de $A$. Il suffira de vérifier que c'est un sev (par la recette), qu'elle est stable par produit et qu'elle contient 1. Comme c'est un sev elle sera stable par différence, et ce sera donc aussi un sous-anneau. La compatibilité s'héritera automatiquement.
\end{dfn}
\begin{rmq}[$\Z$-algèbre, HP]
	Si $(A, +, \times)$ est un anneau, en considérant la loi externe $n \cdot x$ sur $\Z\times A$ nous avons vu qu'elle était également compatible (cf. chapitre "Groupes, anneaux, corps). Une fois de plus, la seule différence est que $\Z$ n'est pas un corps, mais un anneau. A cause de cela, on dira que $A$ est une algèbre sur l'anneau $\Z$, ou $\Z$-algèbre. On peut résumer les choses dans le petit tableau ci-contre.
$$\displaystyle\begin{tabular}{|c|c|c|c}\hline &$\substack{\text{pseudo-associativité} \\ \text{pseudo-distributivité à gauche et à droite} \\ \text{opérateur neutre} }$&axiomes précédents + compatibilité \\ \hline $\Z$ & $\substack{\Z\text{-module} \\ \text{exemple : groupe} (G, +) \text{commutatif} }$ &$\substack{\Z\text{-algèbre} \\ \text{exemple: anneau} (A, +, \times)}$ \\ \hline $\K$&$\K$-ev&$\K$-algèbre \\ \hline
\end{tabular}$$
\end{rmq}
\begin{lem}
	Si $f$ est un isomorphisme de $E$ dans $F$, alors $f^{-1}$ est un isomorphisme de $F$ dans $E$.
\end{lem}
\begin{dfn}[Espaces vectoriels isomorphes]
	Soit $E$ et $F$ deux $\K$-ev. Alors il existe un isomorphisme de E dans F ssi il existe un isomorphisme de F dans E. Dans ce cas, on dit que $E$ et $F$ sont \textbf{isomorphes}. On pourra écrire $E \simeq F$.
\end{dfn}
\begin{dfn}[Groupe linéaire]
	L'ensemble des automorphismes de E est noté $GL(E)$. C'est un groupe
pour la composition. On l'appelle le \textbf{groupe linéaire} de $E$.
\end{dfn}
\section{Familles de vecteurs}
\subsection{Cas particulier des familles finies}
Afin de donner une meilleure intuition à l'étudiant, nous commençons par traiter le cas de familles finies de vecteurs $(e_i)_{1 \leqslant i \leqslant n}$. Nous verrons ensuite ce qu'il en est du cas général. On cherche maintenant à exprimer tout vecteur de $E$ comme combinaison linéaire d'une famille donnée, afin de voir si celle-ci engendre tout l'espace ou non. Soit donc une famille finie $(e_i)_{1 \leqslant i \leqslant n}$, et considérons
$$\begin{array}{ccc}\varphi:\K^n&\rightarrow&E\\ (\lambda_i)_{1\leqslant i\leqslant n}&\mapsto&\displaystyle\sum_{i=1}^{n}\lambda_ie_i\end{array}$$
On va s'intéresser de près à la surjectivité et/ou l'injectivité de $\varphi$. Pour cela, il est plus simple de vérifier d'abord le
\begin{lem}
	$\varphi$ est linéaire.
\end{lem}
\begin{dfn}[Famille génératrice]
	La famille $(e_i)_{1 \leqslant i \leqslant n}$ est \textbf{génératrice} si tout vecteur de $E$ peut s'écrire comme une combinaison linéaire des $e_i$ : $$\forall x\in E,\ \exists(\lambda_1,\ldots,\lambda_n)\in\K^n,\ x=\sum_{i=1}^{n}\lambda_ie_i$$Autrement dit, $\varphi$ est surjective.
\end{dfn}
\begin{exa}
	Si $u$ est un vecteur non nul, $(u)$ est une famille génératrice de la droite vectorielle $Ku$. Si $u_1$ et $u_2$ sont deux vecteurs non colinéaires de l'espace, $(u_1, u_2)$ est une famille génératrice du plan vectoriel $\K u_1 + \K u_2$. La famille vide engendre $\{0\}$.
\end{exa}
On s'intéresse maintenant à la question de savoir si l'écriture comme une combinaison linéaire des $e_i$ est unique.
\begin{dfn}[Famille libre]
	La famille $(e_i)_{1 \leqslant i \leqslant n}$ est libre lorsqu'il y a unicité d'écriture dans $\displaystyle\sum_{i} \lambda_i e_i$ : $$\forall(\lambda_1,\ldots,\lambda)\in\K^n,\ \forall(\mu_1,\ldots,\mu_n)\in\K^n,\ \left(\sum_{i=1}^{n}\lambda_ie_i=\sum_{i=1}^{n}\mu_ie_i\implies \forall i\in\llbracket1,n\rrbracket, \lambda_i=\mu_i\right)$$	Autrement dit, $\varphi$ est injective. On dit aussi que les $e_i$ sont \textbf{linéairement indépendants}. Une famille qui n'est pas libre est dite \textbf{liée}.
\end{dfn}
\begin{rmq}
	Attention, lorsque la famille est libre, il y a unicité mais pas forcément existence!
\end{rmq}
\begin{prop}[Caractérisation des familles libres]
	La famille $(e_i)_{1 \leqslant i \leqslant n}$ est libre si et seulement si $$\forall (\lambda_i) \in \mathbb{K}^n,\ \left( \sum_{i=1}^n \lambda_i e_i = 0 \Rightarrow \forall i \in \llbracket1,n\rrbracket,\ \lambda_i = 0 \right)$$
	Au contraire, la famille est liée si et seulement si $$\exists (\lambda_i) \in \mathbb{K}^n \backslash \{0\},\ \sum_{i=1}^n \lambda_i e_i = 0$$
\end{prop}
\begin{exa}
	La famille $(x, y)$ est liée si et seulement si $x$ et $y$ sont colinéaires.
\end{exa}
\begin{rmq}
	Dès qu'une famille contient le vecteur nul, elle est liée.
\end{rmq}
\begin{prop}[Caractérisation des familles liées]
	La famille $(e_i)_{1 \leqslant i \leqslant n}$ est liée si et seulement si l'un de ses vecteurs peut s'écrire combinaison linéaire des autres.
\end{prop}
\begin{rmq}
	En un certain sens, une famille liée est une famille ``redondante''. En effet, supposons qu'on ait par exemple $$e_{i_0} = \sum_{i \neq i_0}^{n} \frac{\lambda_i}{\lambda_{i_0}} e_i$$
Alors toute combinaison linéaire des $(e_i)_{1 \leqslant i \leqslant n}$ peut se récrire comme combinaison linéaire des $(e_i)_{i \neq i_0}$ en redistribuant le coefficient de $e_{i_0}$ sur les autres $e_i$.$$\sum_{i=1}^{n} \mu_i e_i = \sum_{i \neq i_0} \left(\mu_i - \frac{\mu_{i_0} \lambda_i}{\lambda_{i_0}}\right) e_i$$
Ainsi, $e_{i_0}$ est ``inutile''.
\end{rmq}

\begin{center}
	\framebox{
		\parbox{0.9\textwidth}{
			Le caractère générateur est donc une affaire \textit{d'existence}, le caractère libre une affaire \textit{d'unicité} et le caractère lié une affaire de \textit{redondance}.
		}
	}
\end{center}
\begin{dfn}[Base]
	Si $(e_i)_{1 \leqslant i \leqslant n}$ est à la fois libre et génératrice, alors c'est une \textbf{base} de $E$. De manière équivalente, on peut dire que tout vecteur s'écrit d'une unique manière comme combinaison linéaire des $e_i$ : $$\forall x \in E, \exists!(\lambda_1, \dots, \lambda_n) \in \mathbb{K}^n, \, x = \sum_{i=1}^{n} \lambda_i e_i$$ De manière encore équivalente, $\varphi$ est bijective. Les $\lambda_i$ s'appellent les \textbf{coordonnées de $x$ dans la base $(e_i)_{1 \leqslant i \leqslant n}$}.
\end{dfn}
\begin{exa}
	Dans $\mathbb{R}^2$, $(e_1, e_2)$ est une base.
\end{exa}
Plu généralement, on a la
\begin{dfn}[Base canonique de $\K^n$]
	La \textbf{base canonique} de $\mathbb{K}^n$ est définie par : $$e_1 = \begin{pmatrix} 1 \\ 0 \\ \vdots \\ 0 \end{pmatrix}, \quad e_2 = \begin{pmatrix} 0 \\ 1 \\ \vdots \\ 0 \end{pmatrix}, \quad \dots, \quad e_n = \begin{pmatrix} 0 \\ 0 \\ \vdots \\ 1 \end{pmatrix}$$
\end{dfn}
\begin{rmq}
	Puisque $\varphi$ est un isomorphisme, on en déduit que $\varphi^{-1}$ est aussi un isomorphisme, et en particulier c'est une application linéaire. Ainsi, si les coordonnées de $x$ sont $(\lambda_i)$ et celles de $y$ sont $(\mu_i)$, alors celles de $\alpha x + \beta y$ sont $(\alpha \lambda_i + \beta \mu_i)$. Nous reviendrons plus longuement sur ce point au moment des formes linéaires.
\end{rmq}


\subsection{Généralisation aux familles quelconques} %sous-paragraphe complet%
\begin{dfn}[Famille presque nulle, support, somme presque nulle]
	Une famille de vecteurs $(x_i)_{i\in I}$ est \textbf{presque nulle} lorsque tous les $x_i$ sont nuls sauf un nombre fini d'entre eux. L'ensemble $$S=\left\{i\in I\ | \ x_i\ne0\right\}$$ s'appelle le \textbf{support} de la famille $(x_i)_{i\in I}$. Autrement dit, une famille est presque nulle lorsque son support est fini. Dans ce cas, la \textbf{somme presque nulle} $\displaystyle\sum_{i\in I}x_i$ désigne en fait $\displaystyle\sum_{i\in S}x_i$.
\end{dfn}
\begin{rmq}[Ensemble intermédiaire fini]
	Soit $S$ le support de $(x_i)_{i\in I}$ et $T$ un ensemble fini tel que $S\subset T\subset I$. Alors $$\displaystyle\sum_{i\in I}x_i=\sum_{i\in T}x_i$$ En effet, $(x_i)_{i\in I}$ et $(x_i)_{i\in T}$ ont alors même support.
\end{rmq}
\begin{dfn}[Ensemble des familles presque nulle]
	L'ensemble des familles presque nulles de vecteurs de $E$ indexées sur $I$ est noté $E^{(I)}$.
\end{dfn}
\begin{prop}
	$E^{(I)}$ est un sev de $E^I$.
\end{prop}
\begin{exa}
	$\K^{(I)}$ est un sev de $\K^I$.
\end{exa}
\begin{thm}[Opérations sur les sommes presque nulles]
	Toutes les opération usuelles sur les sommes finies restent licites avec les sommes presque nulles :
	\begin{itemize}
		\item Linéarité de la somme
		\item Changement de variable
		\item Associativité de la somme
		\item Théorème de Fubini
		\item Distributivité généralisée dans une $\K$-algèbre
	\end{itemize}
\end{thm}
\begin{rmq}
	En revanche, attention : pour ce qui est des sommes qui ne sont pas presque nulles, la théorie est beaucoup plus contraignante. On pourra se référer au chapitre "Sommabilité".
\end{rmq}
\begin{rmq}
	Plus généralement, on peut définir des sommes presque nulles dans tout groupe commutatif $(G,+)$, voire même dans tout monoïde (structure algébrique associative munie d'un élément neutre) commutatif.
\end{rmq}
\begin{dfn}[Combinaison linéaire presque nulle]
	Soit $(e_i)_{i\in I}\in E^I$ et $(\lambda_i)_{i\in I}\in\K^I$. Alors $(\lambda_ie_i)_{i\in I}$ est presque nulle, donc on peut considérer $$\sum_{i\in I}\lambda_ie_i$$ qu'on appelle encore une \textbf{combinaison linéaire des $e_i$}. Si on note $S$ le support de $(\lambda_i)_{i\in I}$, le support de $(\lambda_ie_i)_{i\in I}$ est inclus dans $S$, si bien que :$$\sum_{i\in I}\lambda_ie_i=\sum_{i\in S}\lambda_ie_i$$
\end{dfn}
\begin{rmq}
	Lorsque $I=\llbracket1,n\rrbracket$, il s'agit d'une combinaison linéaire au sens usuel, puisque toute famille indexée sur $\llbracket1,n\rrbracket$ est presque nulle. Autrement dit, on a alors $\K^{(I)}=\K^I$, et $$\sum_{i\in I}\lambda_ie_i=\sum_{i=1}^{n}\lambda_ie_i$$ car $\llbracket1,n\rrbracket$ est un sur-ensemble fini du support. Il s'agit donc bien là d'une généralisation du paragraphe précédent.
\end{rmq}
\begin{rmq}
	Un sev reste par combinaison linéaire quelconque, puisqu'il s'agit en fait d'une combinaison linéaire finie.
\end{rmq}
\begin{prop}
	Si $f$ est linéaire, alors on a toujours la relation $$f\left(\sum_{i\in I}\lambda_ie_i\right)=\sum_{i\in I}\lambda_if(e_i)$$
\end{prop}
Fixons une famille quelconque $(e_i)_{i\in I}$, et considérons la nouvelle application $$\begin{array}{cccc}\varphi&\K^{(I)}&\rightarrow&E\\&(\lambda_i)_{i\in I}&\mapsto&\sum_{i\in I}\lambda_i e_i\end{array}$$ Par linéarité de la somme, $\varphi$ reste linéaire (le principe est exactement le même qu'avec les familles finies).
\begin{dfn}
	On garde les mêmes définitions pour une famille génératrice, libre et une base.
	\begin{itemize}
		\item En utilisant l'injectivité de $\varphi$, la caractérisation des familles libres reste valable.
		\item De même, les deux caractérisations des familles liées restent valables.
		\item Dès qu'une famille contient le vecteur nul, ou qu'elle contient une répétition, elle est liée.
	\end{itemize}
\end{dfn}
\begin{thm}
	Toute sur-famille d'une famille génératrice est génératrice. Toute sous-famille d'une famille libre est libre.
\end{thm}
\begin{thm}[HP]
	Une famille est libre si, et seulement si, toutes ses sous-familles finies sont libres.
\end{thm}
\begin{proof}
	Le sens direct est immédiat d'après le théorème précédent. Pour le sens réciproque, se ramener au support de la famille lorsqu'on utilise la caractérisation.
\end{proof}
\begin{rmq}
	Si $F$ est un sev de $E$ et $(e_i)_{i\in I}$ une famille de vecteurs de $F$. Alors $(e_i)_{i\in I}$ est libre en tant que famille de vecteurs de $F$ si, et seulement si, elle est libre en tant que famille de vecteurs de $E$. \par En revanche, attention : les choses changent si on modifie le corps de base. Par exemple, on pourra montre que dans l'espace vectoriel $\R$, la famille $(1,\sqrt{2})$ est $\Q$-libre mais $\R$-liée.
\end{rmq}
\begin{dfn}
	Si $(e_i)_{i\in I}$ est une famille presque nulle de vecteurs de $E$, l'ensemble $\vect((e_i)_{i\in I})$ est défini comme l'ensemble des combinaisons linéaires de $e_i$. On le notera plus simplement $\vect(e_i)_{i\in I}$. C'est un sev de $E$ (puisqu'il est égal à $\im(\varphi)$). C'est le plus petit sev contenant chacun des $e_i$.
\end{dfn}
\begin{exa}
	Par définition, $(e_i)_{i\in I}$ est une famille génératrice de $\vect(e_i)_{i\in I}$. \begin{itemize}
		\item Si elle est libre, alors c'est une base de $\vect(e_i)_{i\in I}$. La réciproque est trivialement vraie.
		\item Elle est génératrice si, et seulement si, $\vect(e_i)_{i\in I}=E$.
	\end{itemize}
\end{exa}
\begin{exa}
	Posons $$E=\left\{x\mapsto\sum_{k=0}^{n}a_kx^k \ | \ n\in\N,\ (a_0,\ldots,a_n)\in\R^{n+1}\right\}$$ l'ensemble des fonctions polynomiales réelles, et considérons la famille $(f_k)_{k\in\N}$ définie par $$ \forall k\in\N,\ f_k:x\mapsto x^k$$ On vérifie facilement que dans le $\R$-ev $\R^\R$, on a $E=\vect(f_k)_{k\in\N}$. Donc $E$ est un $\R$-ev et $(f_k)_{k\in\N}$ en est une famille génératrice. Puis on montre que $(f_k)_{k\in\N}$ est libre, si bien que c'est une base.
\end{exa}
\begin{exa}
	Dans l'espace vectoriel $\R^\N$, définissons pour tout $k\in\N$ la suite $(e_n^k)_{n\in\N}$ par $$\forall n\in\N,\ e_n^k=\delta_{n,k}$$ Autrement dit, la suite $e^k$ est la suite qui vaut $0$ partout, sauf $1$ en $k$-ème position. On vérifie facilement que dans $\R^\N$, on a $\R^{(\N)}=\vect(e_n^k)_{k\in\N}$. Donc $(e^k)$ est une famille génératrice de $\R^{(\N)}$, puis on montre qu'elle est libre, si bien qu'elle est une base de $\R^{(\N)}$.
\end{exa}
\begin{exa}
	On peut généraliser l'exemple précédent à $\K^{(I)}$.
\end{exa}
\begin{rmq}
	Supposons que l'un des $e_i$ soit combinaison linéaire des autres. Comme dans le cas fini, on ne modifie pas le $\vect$ de la famille en le retirant.
\end{rmq}
\begin{rmq}
	On rappelle que les notions de famille et d'ensemble communiquent à travers les opérations suivantes.
	\begin{itemize}
		\item A une famille $(e_i)_{i\in I}$, on peut associer l'ensemble $\left\{e_i\ | \ i\in I\right\}$.
		\item Réciproquement, à un ensemble $A$ on peut associer la famille $(a)_{a\in A}$.
	\end{itemize}
	A chaque fois, la famille et l'ensemble ont exactement les même éléments (il faut juste prendre gade aux répétitions). Donc un sev contient les éléments de l'ensemble si, et seulement si, il contient les éléments de la famille. On en déduit que l'ensemble et la famille ont le même $\vect$.
\end{rmq}
\begin{prop}[$\vect$ d'une partie quelconque]
	Soit $A$ une partie quelconque de $E$. Alors $\vect(A)$ est égal à l'ensemble des combinaisons linéaires des vecteurs de $A$ : $$\vect(A)=\left\{\sum_{a\in A}\lambda_aa\ | \ (\lambda_a)_{a\in A}\in\K^{(A)}\right\}$$
\end{prop}

\subsection{Applications linéaires et bases} %sous-paragraphe complet%
\begin{thm}[Premier théorème de caractérisation : image d'une base]
	Soit $(e_i)_{i\in I}$ une base de $E$ et $F$ un espace vectoriel sur le même corps que $E$. Alors : $$\forall (f_i)_{i\in I}\in F^I,\ \exists!u\in\mathcal{L}(E,F),\ \forall i\in I,\ u(e_i)=f_i$$
\end{thm}
\begin{prop}
	Si $u\in\mathcal{L}(E,F)$ et $(e_i)_{i\in I}$ est une famille de vecteurs de $E$, alors on a : $$\vect(u(e_i))_{i\in I}=u\left(\vect(e_i)_{i\in I}\right)$$
\end{prop}
\begin{exa}
	Si $(e_i)_{i\in I}$ est génératrice, on a $\im(u)=\vect(u(e_i))_{i\in I}$.
\end{exa}
\begin{thm}
	Soit $u\in\mathcal{L}(E,F)$, et $(e_i)_{i\in I}$ une base de $E$. Alors
	\begin{itemize}
		\item $u$ est injective si, et seulement si, $(u(e_i))_{i\in I}$ est libre
		\item $u$ est surjective si, et seulement si, $(u(e_i))_{i\in I}$ est génératrice
		\item $u$ est un isomorphisme si, et seulement si, $(u(e_i))_{i\in I}$ une base
	\end{itemize}
\end{thm}
\section{Somme de sous-espaces} %paragraphe complet%
\begin{prop}
	Soit $E_1$ et $E_2$ deux sev de $E$. Alors leur somme $$E_1+E_2=\left\{x\in E\ | \ \exists(x_1,x_2)\in E_1\times E_2,\ x=x_1+x_2\right\}$$ est encore un sev.
\end{prop}
\begin{exa}
	On a $E_1+\left\{0\right\}=E_1$ et $\left\{0\right\}+E_2=E_2$.
\end{exa}
\begin{prop}
	Soit $(e_i)_{i\in I}$ une famille de vecteurs de $E$, et $I_1,I_2\subset I$ tels que $I=I_1\sqcup I_2$. Alors $$\vect(e_i)_{i\in I}=\vect(e_i)_{i\in I_1}+\vect(e_i)_{i\in I_2}$$
\end{prop}
\begin{exa}
	On a $\vect(E_1\cup E_2)=E_1+E_2$.
\end{exa}
On voit que tout élément de $E_1+E_2$ s'écrit comme la somme d'un élément de $E_1$ et d'un élément de $E_2$, mais y a-t-il unicité de l'écriture dans cette somme ? La petite définition suivante donne un moyen commode de le savoir.
\begin{dfn}[Somme directe]
	Les deux assertions suivantes sont équivalentes.
	\begin{enumerate}
		\item $\forall x\in E_1+E_2,\ \exists!(x_1,x_2)\in E_1\times E_2,\ x=x_1+x_2$
		\item $\forall(x_1,x_2)\in E_1\times E_2,\ x_1+x_2=0\implies x_1=x_2=0$
	\end{enumerate}
	On dit alors que $E_1$ et $E_2$ sont en \textbf{somme directe}, et on note leur somme $E_1\oplus E_2$.
\end{dfn}
\begin{prop}[Commutativité de la somme directe]
	Si $E_1$ et $E_2$ sont en somme directe, alors $E_2$ et $E_1$ sont en somme directe et on a : $$E_1\oplus E_2=E_2\oplus E_1$$
\end{prop}
Voici à présent une caractérisation extrêmement commode.
\begin{prop}
	$E_1$ et $E_2$ sont en somme directe si, et seulement si, $E_1\cap E_2=\left\{0_E\right\}$.
\end{prop}
\begin{exa}
	Soit $F$ un sev de $E$ et $x\in E\setminus F$. Alors $F$ et $\K x$ sont en somme directe. Il suffit de se donner $y\in F\cap \K x$ et de l'écrire sous la forme $\lambda x$. On a nécessairement $\lambda=0$ sans quoi $x$ appartiendrait à $F$ en dilatant de $\lambda\inv$, donc $x=0_E$ et on conclut par la caractérisation précédente.
\end{exa}
\begin{prop}[Associativité de la somme directe]
	Soit $F$, $G$ et $H$ des sev de $E$. Supposons que $F$ et $G$ soient en somme directe, et que $F\oplus G$ et $H$ soient en somme directe. Alors, $G$ et $H$ sont en somme directe, et $F$ et $G\oplus H$ sont en somme directe, et on a naturellement : $$(F\oplus G)\oplus H=F\oplus(G\oplus H)$$
\end{prop}
\begin{dfn}[Sous-espaces supplémentaires]
	Si de plus, $E_1\oplus E_2=E$, on dit que $E_1$ et $E_2$ sont \textbf{supplémentaires} (on pourra aussi dire que $E_1$ est un supplémentaire de $E_2$ et vice-versa).
\end{dfn}
\begin{rmq}
	Montrer que $E_1$ et $E_2$ sont supplémentaires correspond donc à un problème d'existence-unicité.
	\begin{itemize}
		\item Existence pour pouvoir écrire $x=x_1+x_2$ quel que soit $x\in E$
		\item Unicité pour que cette écriture soit unique
	\end{itemize}
	Comme toujours, l'unicité sera plus facile à montrer que l'existence (en l'occurrence, il suffit de montrer que $E_1\cap E_2=\left\{0_E\right\}$). On pourra souvent raisonner par analyse-synthèse pour montrer le caractère supplémentaire de deux sev.
\end{rmq}
\begin{rmq}
	Si $E_1$ est un sev de $E$, on peut montrer qu'il admet un supplémentaire : la preuve en dimension finie sera donnée dans le chapitre "Dimension finie", et la preuve en dimension infinie, plus délicate, nécessite le lemme de Zorn. En revanche, ce supplémentaire n'a aucune raison d'être unique.
\end{rmq}
\begin{exa}
	Soit $I$ un intervalle de $\R$ centré autour de $0$. Alors les fonctions paires et les fonctions impaires de $I$ dans $\R$ forment deux sev supplémentaires de $\R^I$. En effet, on montre aisément que ce sont des sev de $\R^I$ avec la recette, puis on montre la supplémentarité avec une analyse-synthèse qui fait apparaître ce qu'on appelle les parties paire et impaire de la fonction. Par exemple, $\cosh$ et $\sinh$ son respectivement les parties paire et impaire de $\exp$.
\end{exa}
\begin{thm}[Deuxième théorème de caractérisation : supplémentaires]
	Soit $E_1$ et $E_2$ deux sev de $E$ supplémentaires. Soit $u_1\in\mathcal{L}(E_1,F)$ et $u_2\in\mathcal{L}(E_2,F)$. Alors $$\exists! u\in\mathcal{L}(E,F),\ \left \{\begin{array}{ll}u_{|E_1}=u_1\\ u_{|E_2}=u_2\end{array}\right.$$
\end{thm}
\begin{thm}[Isomorphisme induit d'un supplémentaire du noyau sur l'image]
	Soit $f\in\mathcal{L}(E,F)$. Alors $f$ induit un isomorphisme de tout supplémentaire de $\ker(f)$ sur $\im(f)$.
\end{thm}
\begin{thm}[Théorème de la base adaptée]
	Soit $(e_i)_{i\in I}$ une base de $E$, supposons qu'on ait $I=I_1\sqcup I_2$. Alors $E_1=\vect(e_i)_{i\in I_1}$ et $E_2=\vect(e_i)_{i\in I_2}$ sont supplémentaires dans $E$. \par Réciproquement, si $E_1$ et $E_2$ sont supplémentaires, alors en concaténant une base de $E_1$ avec une base de $E_2$, on obtient une base de $E$. On dit que cette base de $E$ est \textbf{adaptée à la décomposition} $E_1\oplus E_2$.
\end{thm}
\begin{exa}
	Soit $(e_i)_{i\in I}$ une famille libre et $I_1,I_2\subset I$ tels que $I=I_1\sqcup I_2$. Alors $$\vect(e_i)_{i\in I}=\vect(e_i)_{i\in I_1}\oplus\vect(e_i)_{i\in I_2}$$
\end{exa}
\begin{exa}
	Dans le $\R$-ev $E$ des fonctions polynomiales réelles, on rappelle que la famille $f_k:x\mapsto x^k$ forme une base. Si on pose $E_1=\vect(f_{2k})_{k\in\N}$ et $E_2=\vect(f_{2k+1})_{k\in\N}$, on obtient : $$E=E_1\oplus E_2$$ On peut vérifier que $E_1$ (resp. $E_2$) est exactement l'ensemble des fonctions polynomiales paires (resp impaires).
\end{exa}

\section{Applications linéaires importantes} %paragraphe complet%
\begin{dfn}[Homothétie de rapport $\lambda$]
	Soit $E$ un $\K$-ev et $\lambda\in\K$. L'\textbf{homothétie de rapport $\lambda$} est l'endomorphisme de $E$ $x\mapsto\lambda x$.
\end{dfn}
\begin{prop}
	L'application $\lambda\mapsto\lambda\id$ est un morphisme de $(\K\st,\times)$ dans $(GL(E),\circ)$.
\end{prop}
\begin{dfn}[Projecteur]
	Soit $F$ et $G$ deux espaces supplémentaires $ E=F\oplus G$. Le \textbf{projecteur sur $F$ parallèlement à $G$} est défini par $$p:x=f+g\in F\oplus G\mapsto f$$ Comme il y a unicité de l'écriture $x=f+g$, $p$ est bien défini.
\end{dfn}
\begin{prop}
	$p$ est un endomorphisme de $E$ et il vérifie $p^2=p$. On dit qu'il est \textbf{idempotent}.
\end{prop}
\begin{rmq}
	Attention ! L'exposant $2$ dans $p^2=p$ est à interpréter au sens de la composition : $p^2=p\circ p$.
\end{rmq}
\begin{exa}
	Le seul projecteur inversible est l'identité (raisonner par analyse-synthèse et composer par l'inverse de ce projecteur dans l'égalité $p^2=p$).
\end{exa}
\begin{rmq}
	Si $p$ est le projecteur sur $F$ parallèlement à $G$, alors $\displaystyle p_{|F}^{|F}=\id_F$ et $\displaystyle p_{|G}^{|G}=0$.
\end{rmq}
\begin{prop}
	Soit $p$ le projecteur sur $F$ parallèlement à $G$. Alors on a : $$\left \{\begin{array}{ll}F=\im(p)\\ G=\ker(p)\end{array}\right.$$
\end{prop}
\begin{rmq}
	On en déduit que $F$ et $G$ sont déterminés de manière unique par $p$.
\end{rmq}
De manière intéressante, la réciproque est également vraie.
\begin{prop}
	Si $p\in\mathcal{L}(E)$ vérifie $p^2=p$, alors c'est le projecteur sur $\im(p)$ parallèlement à $\ker(p)$.
\end{prop}
\begin{dfn}[Projecteur associé]
	SI $p$ est le projecteur sur $F$ parallèlement à $G$, alors $q=\id-p$ est le projecteur sur $G$ parallèlement à $F$. On a donc $$\left \{\begin{array}{ll}\ker(q)=\im(p)\\ \im(q)=\ker(p)\end{array}\right.$$ $q$ est le \textbf{projecteur associé} de $p$. Bien sûr, réciproquement, $p$ est le projecteur associé de $q$. De manière générique, on dit que $p$ et $q$ sont associés.
\end{dfn}
\begin{rmq}
	Si $p$ et $q$ sont associés, on voit immédiatement que $pq=qp=0$.
\end{rmq}
\begin{dfn}[Symétrie]
	Si $E=F\oplus G$, la \textbf{symétrie par rapport à $F$ parallèlement à $G$} est définie par $$s:x=f+g\in F\oplus G\mapsto f-g$$
\end{dfn}
\begin{prop}
	$s$ est un endomorphisme de $E$, et il vérifie $s^2=\id$.
\end{prop}
\begin{rmq}
	$s$ est une involution, et par conséquent une bijection.
\end{rmq}
\begin{prop}
	Si $p$ est le projecteur sur $F$ parallèlement à $G$, alors on a $$s=2p-\id_E$$ De manière équivalente, si $q$ est le projecteur associé de $p$, on a $$s=p-q$$
\end{prop}
\begin{rmq}
	On peut alors retrouver le fait que $s$ est un endomorphisme et un calcul rapide montre que $s^2=\id_E$.
\end{rmq}
\begin{rmq}
	Si $s$ est la symétrie par rapport $F$ parallèlement à $G$, alors $\displaystyle p_{|F}^{|F}=\id_F$ et $\displaystyle p_{|G}^{|G}=-\id_G$.
\end{rmq}
\begin{prop}
Soit $s$ la symétrie par rapport $F$ parallèlement à $G$. Alors on a : $$\left \{\begin{array}{ll}F=\ker(s-\id)\\ G=\ker(s+\id)\end{array}\right.$$
\end{prop}
\begin{rmq}
	On en déduit que $F$ et $G$ sont déterminés de manière unique par $s$.
\end{rmq}
Là aussi, la réciproque est également vraie.
\begin{prop}
	Si $s\in\mathcal{L}(E)$ vérifie $s^2=\id$, alors c'est la symétrie parallèlement à $\ker(s-\id)$ parallèlement à $\ker(s+\id)$.
\end{prop}
\begin{rmq}[Symétries en caractéristique $2$, HP]
	Attention, les deux propositions précédentes deviennent fausses lorsque $\K$ est de caractéristique $2$ ! En effet, dans ce cas, on ne peut plus effectuer de divisions par $2$, puisque $2=0$. \par Voici d'ailleurs un contre-exemple. Si $\K=\Z/2\Z$, on considère le $\K$-ev $\K^2$ muni de la base canonique $(e_1,e_2)$. Alors $s$ définie par $$\left \{\begin{array}{ll}s(e_1)=e_1 \\ s(e_2)=e_1+e_2\end{array}\right.$$ vérifie $s^2=\id$ et pourtant $s\ne \id$. Or, si $s$ était une symétrie, on aurait $s(f+g)=f-g=f+g$ d'où $s=\id$. 
\end{rmq}

\section{Hyperplans et formes linéaires} %paragraphe complet%
\begin{dfn}[Forme linéaire, espace dual]
	On appelle \textbf{forme linéaire} une application $f\in\mathcal{L}(E,\K)$. $\mathcal{L}(E,\K)$ est noté $E\st$, on l'appelle \textbf{espace dual} de $E$.
\end{dfn}
\begin{prop}
	Toute forme linéaire non nulle est surjective.
\end{prop}
\begin{dfn}[Hyperplan]
	Un supplémentaire d'une droite vectorielle s'appelle un \textbf{hyperplan} de $E$.
\end{dfn}
\begin{rmq}
	Le vocabulaire provient manifestement de la dimension $3$, où le supplémentaire d'une droite vectorielle est un plan vectoriel comme nous le verrons dans le chapitre "Dimension finie". En dimension supérieure, l'analogue d'un plan s'appelle un "hyperplan" pour suggérer qu'on est monté en dimension (de même, on peut parler d'hypercube, d'hypersphère...).
\end{rmq}
\begin{thm}[Caractérisation des hyperplans par les formes linéaires]
	$H$ est un hyperplan de $E$ si, et seulement si, c'est le noyau d'une forme linéaire non nulle.
\end{thm}
\begin{cor}
	Si $H$ est un hyperplan de $E$ et $u\notin H$, alors on a toujours $H\oplus\K u=E$.
\end{cor}
\begin{rmq}
	En d'autres termes, un hyperplan a été défini comme un supplémentaire d'\textbf{une} droite vectorielle (qu'il ne contient nécessairement pas). Et $a$ $posteriori$, il s'avère qu'il est supplémentaire de \textbf{toute} droite vectorielle qu'il ne contient pas.
\end{rmq}
\begin{dfn}[Formes coordonnées associées à une base]
	Soit $(e_i)_{i\in I}$ une base de $E$. Les \textbf{formes coordonnées} associées à la base $(e_i)_{i\in I}$ sont les formes linéaires $(e_i\st)_{i\in I}$ définies par $$\forall(i,j)\in I^2,\ \delta_{i,j}=\left \{\begin{array}{ll}1&\text{si }i=j \\ 0&\text{sinon}\end{array}\right.$$ 
\end{dfn}
\begin{rmq}
	Notons que les formes coordonnées sont bien définies, car une application linéaire est fixée de manière unique par l'image d'une base.
\end{rmq}
\begin{rmq}
	Même si la notation ne le suggère pas (et cela est malheureux, mais rendrait la notation trop lourde), les formes coordonnées peuvent toutes changer si on changer un seul vecteur de la base pour obtenir une autre base. Cependant, les autres formes coordonnées ne changent pas de notation, bien qu'elles puissent changer puisque la base a changé.
\end{rmq}
\begin{dfn}[Symbole de Kronecker]
	Le \textbf{symbole de Kronecker}, ou \textbf{delta de Kronecker}, est défini par $$\forall(i,j)\in I^2,\ \delta_{i,j}=\left \{\begin{array}{ll}1&\text{si }i=j \\ 0&\text{sinon}\end{array}\right.$$ 
\end{dfn}
\begin{rmq}
	Avec cette notation, la définition des formes coordonnées devient $$\forall(i,j)\in I^2,\ e_i\st(e_j)=\delta_{i,j}$$
\end{rmq}
Le théorème suivant justifie l'appellation de "forme coordonnée".
\begin{thm}
	Soit $(e_i)_{i\in I}$ une base de $E$ et $(e_i\st)_{i\in I}$ les formes coordonnées associées. Alors, pour tout $x\in E$, la coordonnée selon $e_i$ vaut $e_i\st(x)$.
\end{thm}
\begin{exa}
	Dans le $\R$-ev $\C$, les fonctions $\text{Re}(\cdot)$ et $\text{Im}(\cdot)$ sont les formes coordonnées associées à la base $(1,i)$.
\end{exa}
\section{Translations, sous-espaces affines} %paragraphe complet%
\subsection{Notions de base}
\begin{dfn}[Translation]
	Soit $E$ un $\K$-ev et $a\in E$ fixé. La \textbf{translation de vecteur $a$} est l'application $x\mapsto x+a$. Attention, en général, ce n'est \textbf{pas} un endomorphisme (sauf si $a\ne0$).
\end{dfn}
Jusqu'ici, nous avons vu que tout sev contenait au moins $0$. Ainsi, les droites vectorielles, les plans vectoriels, et plus généralement tous les sev "passent par l'origine" au sens géométrique du terme. Or, en géométrie élémentaire, on sait qu'une droite ou un plan peut très bien ne pas passer par l'origine, tout en étant parallèle à une droite ou un plan qui passe par l'origine. Il nous faudrait donc une généralisation de la notion de sev. Ce ce qu'on appelle un sea, o sous-espace affine. Il consiste à se donner un point de l'espace comme origine, puis à définir un sous-espace qui passe par ce point et qui est parallèle à un sev donné.
\begin{dfn}[Sous-espace affine]
	Soit $F$ un sev de $E$ et $a\in E$. Alors, le \textbf{sous-espace affine de $E$ passant par $a$ et parallèle à $F$} est défini par $$a+F=\left\{a+f\ | \ f\in F\right\}$$ C'est l'image de $F$ par la translation de vecteur $a$. Les éléments de $a+f$ sont appelés \textbf{points} de $a+F$. $F$ s'appelle \textbf{la direction} du sea, et $a$ s'appelle \textbf{une origine} du sea.
\end{dfn}
\begin{prop}
	La direction d'un sea est unique. En revanche, n'importe lequel de ses points peut servir de direction. De manière précise, soit $\mathcal{F}=a+F$ un sea, et soit $a'\in E$ et $F'$ un sev de $E$. Alors $\mathcal{F}=a'+F'$ si, et seulement si, $F=F'$ et $a'\in\mathcal{F}$.
\end{prop}
\begin{exa}[Équation affine]
	Soit $f\in\mathcal{L}(E)$, et cherchons les solutions de l'équation $f(x)=b$, où $b\in E$ est fixé. Si l'équation possède au moins une solution $x_0$, alors l'ensemble des solutions est le sea passant par $x_0$ et de direction $\ker(f)$.
\end{exa}
\begin{exa}
	Soit $\varphi$ une forme linéaire non nulle, et soit $y\in\K$ fixé. Puisque $\varphi$ est surjective, l'équation $\varphi(x)=y$ admet au moins une solution. Et donc, son ensemble de solutions est un hyperplan affine dirigé par $\ker(\varphi)$.
\end{exa}
\subsection{Notions plus avancées}
\begin{dfn}[Parallélisme]
	Deux sea $\mathcal{A}$ et $\mathcal{B}$ sont dits \textbf{parallèles} lorsqu'ils ont même direction. Le parallélisme est une relation d'équivalence sur l'ensemble des sea.
\end{dfn}
\begin{prop}[Intersection de sea]
	Soit $\mathcal{A}$ un sea de direction $F$ et $\mathcal{B}$ un sea de direction $G$. Si $\mathcal{A}\cap\mathcal{B}\ne\emptyset$, alors $\mathcal{A}\cap\mathcal{B}$ est un sea de direction $F\cap G$.
\end{prop}
\begin{rmq}
	Attention aux fausses intuitions ! Dans le plan, on sait que deux droites non parallèles se coupent toujours, mais dans l'espace (et au-delà), ce résultat devient faux. Par exemple, dans $\R^3$, la droite $y=z=0$ n'a aucun point commun avec la droite $\displaystyle\left \{\begin{array}{ll}y=0 \\ z=1 \end{array}\right.$.
\end{rmq}
\begin{exa}
	Dans $\R^3$, on peut montrer que l'intersection de deux hyperplans non parallèles est une droite. Pour l'instant, on se contentera de faire un dessin ; pour une démonstration élégante, on préfèrera attendre le chapitre sur la dimension finie.
\end{exa}
\begin{rmq}
	Plus généralement, si on se donne un ensemble de sea d'intersection non vide, alors cette intersection est un sea, de direction l'intersection des directions. La démonstration est rigoureusement la même.
\end{rmq}
\begin{dfn}[Sea engendré par une partie]
	Soit $X$ une partie non vide de $E$. Alors l'intersection de tous les sea contenant $X$ est non vide (puisqu'elle contient $X$). C'est donc un sea, qu'on appelle \textbf{sea engendré par $X$}. C'est le plus petit sea contenant $X$.
\end{dfn}
\begin{exa}
	Le sea engendré par deux points non confondus est une droite affine. Le sea engendré par trois points non alignés est un plan affine.
\end{exa}
\newpage
\chapter{Dimension finie} %chapitre complet%
Dans tout le chapitre, on fixe $\K$ un corps et $E$ un $\K$-espace vectoriel.
\section{Dimension}
\begin{dfn}[Espace de dimension finie]
	$E$ dit de \textbf{dimension finie} lorsqu'il admet une famille génératrice finie.
\end{dfn}
\begin{exa}
	$\left\{0\right\}$ est de dimension finie car engendré par la famille vide qui est finie.
\end{exa}
\begin{lem}[Lemme d'adjonction, NS]
	Soit $(e_1,\ldots,e_p)$ une famille libre de vecteurs de $E$ et $x\in E\setminus\vect(e_1,\ldots,e_p)$. Alors, la famille $(e_1,\ldots,e_p,x)$ est encore libre.
\end{lem}
\begin{rmq}
	Le lemme reste vrai pour tout autre ensemble d'indexation que $\llbracket1,p\rrbracket$, même si cet ensemble est infini.
\end{rmq}
\begin{thm}[Théorème de la base incomplète, cas fini]
	Soit $\mathcal{L}$ une famille libre finie de vecteurs de $E$ et $\mathcal{G}$ une famille génératrice finie de $E$, qui va nous servir de "réservoir". Alors $\mathcal{L}$ peut être complétée en une base à l'aide de certains vecteurs de $\mathcal{G}$. 
\end{thm}
\begin{thm}[Théorème de la base extraite, cas fini]
	Soit $\mathcal{G}$ une famille génératrice finie de $E$. Alors, il existe une sous-famille de $\mathcal{G}$ qui est une base de $E$, et une base de ce type est appelée \textbf{base extraite}.
\end{thm}
\begin{thm}
	Tout espace de dimension finie admet une base finie.
\end{thm}
\begin{rmq}
	En fait, on remarque qu'il y a équivalence entre le fait de posséder une famille génératrice finie et une base finie.
\end{rmq}
\begin{lem}[Lemme du pivot]
	Supposons que $E$ admette une famille génératrice finie de taille $n\in\N$. Alors, $n+1$ vecteurs de $E$ sont toujours liés.
\end{lem}
\begin{exa}
	Trois vecteurs du plan $\R^2$ sont toujours liés. Quatre vecteurs de l'espace $\R^3$ sont toujours liés.
\end{exa}
\begin{rmq}
	Plus généralement, on montre, puisque toute sur-famille d'une famille liée est liée, que si $E$ une famille génératrice de $n$ vecteurs, alors toute famille de $m\geqslant n+1$ vecteurs est liée. Cela reste donc vrai si on a une famille infinie de vecteurs.
\end{rmq}
\begin{dfn}[Dimension d'un espace de dimension finie]
	Supposons que $E$ est de dimension finie. Alors toutes les bases de $E$ ont la même taille, notons-la $n\in\N$ par exemple. On pose alors la \textbf{dimension de $E$} comme étant égale à : $$\dim(E):=n$$
\end{dfn}
\begin{exa}[$\K^n$ pour $n\in\N\st$]
	Soit $n\in\N\st$. Alors $\K^n$ est un $\K$-espace vectoriel de dimension finie, et $\dim(\K^n)=n$ (prendre la base canonique).
\end{exa}
\begin{exa}[Droite vectorielle]
	Soit une droite vectorielle $\K u$ avec $u\ne0_E$. $(u)$ est génératrice de $\K u$ et libre, donc c'est une base de $\K u$. Puis $\dim(\K u)=1$.
\end{exa}
\begin{exa}[Plan vectoriel]
	Soit un plan vectoriel $\K u+\K v$ avec $u$ et $v$ non colinéaires. $(u,v)$ en et une base, donc $\dim(\K u+\K v)=2$.
\end{exa}
\begin{exa}[Espace nul]
	La famille vide est une base de l'espace nul, donc $\dim(\left\{0\right\})=0$.
\end{exa}
\begin{thm}[Lien entre taille et caractères libre et générateur]
	Supposons que $E$ est de dimension $n\in\N$. Alors
	\begin{enumerate}
		\item Pour toute famille libre $\mathcal{L}$ de vecteurs de $E$, on a $\text{Taille}(\mathcal{L})\leqslant n$ avec égalité si, et seulement si $\text{Taille}(\mathcal{L})=n$.
		\item Pour toute famille génératrice $\mathcal{G}$ de vecteurs de $E$, on a $\text{Taille}(\mathcal{G})\geqslant n$ (possiblement infinie) avec égalité si, et seulement si $\text{Taille}(\mathcal{G})=n$.
	\end{enumerate}
\end{thm}
\begin{lem}
	Supposons que $E$ est de dimension finie. Alors toute famille génératrice $\mathcal{G}$ de $E$ possède une sous-famille génératrice finie.
\end{lem}
\begin{thm}[Théorème de la base incomplète, cas général]
	Supposons que $E$ est de dimension finie. Soit $\mathcal{L}$ une famille libre de vecteurs de $E$ quelconque et $\mathcal{G}$ une famille génératrice de $E$ quelconque. Alors, $\mathcal{L}$ peut être complétée en une base de $E$ à partir de certains vecteurs de $\mathcal{G}$.
\end{thm}
\begin{thm}[Théorème de la base extraite, cas général]
	Supposons que $E$ est de dimension finie. Alors on peut extraire une base de $E$ de toute famille génératrice $\mathcal{G}$ de $E$.
\end{thm}
\begin{thm}[Dimension et isomorphisme]
	Supposons que $E$ est de dimension finie et que $F$ est un $\K$-espace vectoriel quelconque. Alors : \begin{itemize}
		\item Si $E$ et $F$ sont isomorphes, alors $F$ est de dimension finie, et on a $\dim(F)=\dim(F)$.
		\item Réciproquement, si $F$ est de dimension finie et si $\dim(F)=\dim(E)$, alors $E$ et $F$ sont isomorphes.
	\end{itemize}
\end{thm}
\begin{exa}
	Tout $\K$-espace vectoriel de dimension $n\in\N$ est isomorphe à $\K^n$, avec la convention $\K^0=\left\{0\right\}$.
\end{exa}
\begin{thm}[Dimension d'un produit cartésien]
	Soit $E$ un $\K$-ev de dimension finie égale à $p\in\N$ et $F$ un $\K$-ev de dimension finie égale à $n$. Alors $E\times F$ est de dimension finie égale à $p+n$, $ie$ on a alors $$\dim(E\times F)=\dim(E)+\dim(F)$$
\end{thm}
\begin{thm}[Dimension de $\mathcal{L}(E,F)$]
	Soit $E$ un $\K$-ev de dimension finie égale à $p\in\N$ et $F$ un $\K$-ev de dimension finie égale à $n$. Alors $\mathcal{L}(E,F)$ est de dimension finie égale à $p\times n$, $ie$ on a alors $$\dim(\mathcal{L}(E,F))=\dim(E)\times\dim(F)$$
\end{thm}
\begin{exa}
	Si $E$ est de dimension finie, alors $\mathcal{L}(E)$ est de dimension finie et on a : $$\dim(\mathcal{L}(E))\dim(E)^2$$
\end{exa}
\begin{exa}
	Si $E$ est de dimension finie, alors $E\st$ est de dimension finie et on a : $$\dim(E\st)=\dim(E)$$
\end{exa}
\section{Dimension d'un sous-espace vectoriel}
\begin{thm}[Dimension et sev]
	Soit $E$ de dimension finie et $F$ un sev de $E$. Alors, $F$ est de dimension finie et on a $\dim(F)\leqslant\dim(E)$, avec égalité si, et seulement si, $F=E$.
\end{thm}
\begin{rmq}
	On pourra souvent se servir du cas d'égalité pour montrer une égalité d'ensembles. Retenir que inclusion et égalité de dimension donne égalité d'ensembles.
\end{rmq}
\begin{exa}
	On retrouve que les seuls sous-espaces vectoriels de $\K$ sont $\left\{0\right\}$ et $\K$.
\end{exa}
\begin{exa}
	On retrouve le fait que toute forme linéaire non nulle est surjective (en utilisant l'exemple précédent sur l'image de cette forme linéaire).
\end{exa}
\begin{dfn}[Rang d'une famille de vecteurs]
	Soit $(x_1,\ldots,x_p)\in E^p$ avec $p\in\N$. Le \textbf{rang de la famille} $(x_j)_{1\leqslant j\leqslant p}$ est défini par : $$rg(x_j)_{1\leqslant j\leqslant p}:=\dim\left(\vect(x_j)_{1\leqslant j\leqslant p}\right)$$
\end{dfn}
\begin{prop}
	Soit $(x_1,\ldots,x_p)\in E^p$ avec $p\in\N$. On a $rg(x_j)_{1\leqslant j\leqslant p}\leqslant p$ avec égalité si, et seulement si, $(x_1,\ldots,x_p)$ est libre.
\end{prop}
\begin{thm}[Existence du supplémentaire en dimension finie]
	Soit $E$ un espace de dimension finie et $F$ un sev de $E$. Alors $F$ admet un supplémentaire dans $E$.
\end{thm}
\begin{rmq}
	Attention : on rappelle que ce supplémentaire n'a \textbf{aucune raison} d'être unique !
\end{rmq}
\begin{rmq}[Cas de la dimension infinie, HP]
	On rappelle qu'en dimension, ce résultat reste vrai, mais il faut montrer que l'ensemble des sev en somme directe avec $F$ est un ensemble inductif puis lui appliquer le lemme de Zorn. Il faut alors montrer la somme totale. ON peut par exemple raisonner par l'absurde et utiliser la maximalité offerte par le lemme de Zorn.
\end{rmq}
\begin{lem}[Lemme pré-Grassmann, NS]
	Soit $E$ un espace vectoriel quelconque, et $F$ et $G$ des sous-espaces vectoriels de $E$ de dimension finie. Si $F$ et $G$sont en somme directe, alors $F\oplus G$ est de dimension finie et on a : $$\dim(F\oplus G)=\dim(F)+\dim(G)$$
\end{lem}
\begin{exa}[Cas fondamental]
	Supposons que $E$ est de dimension finie et que $F$ et $G$ sont en somme directe. Si $\dim(F)+\dim(G)=\dim(E)$ (comme on pourra souvent le montrer plus tard par le \textbf{théorème du rang}), alors $E=F\oplus G$, $ie$ $F$ et $G$ sont supplémentaires dans $E$.
\end{exa}
\begin{thm}[Formule de Grassmann]
	Soit $E$ un espace vectoriel quelconque, et $F$ et $G$ des sous-espaces vectoriels de $E$ de dimension finie. Alors $F+G$ est de dimension finie et on a : $$\dim(F+G)=\dim(F)+\dim(G)-\dim(F\cap G)$$
\end{thm}
\begin{exa}
	Soit $E$ un espace vectoriel quelconque, et $F$ et $G$ des sous-espaces vectoriels de $E$ de dimension finie. Alors $F+G$ est de dimension finie et on a : $$\dim(F+G)\leqslant\dim(F)+\dim(G)$$
\end{exa}
\begin{exa}
	On peut généraliser par récurrence : si $F_1,\dots,F_n$ sont des sev de dimension finie de $E$, alors $F_1+\ldots+F_n$ est de dimension finie et on a : $$\dim\left(\sum_{i=1}^{n}F_i\right)\leqslant\sum_{i=1}^{n}\dim(F_i)$$
\end{exa}
\section{Rang d'une application linéaire}
\begin{dfn}[Rang d'une application linéaire]
	Soit $u\in\mathcal{L}(E,F)$. $u$ est dite de \textbf{rang fini} lorsque $\im(u)$ est de dimension finie. Dans ce cas, le \textbf{rang de $u$} est défini par : $$rg(u):=\dim\left(\im(u)\right)$$
\end{dfn}
\begin{rmq}
	Si $E$ ou $F$ est de dimension finie, il est inutile de prendre toutes ces précautions puisqu'alors $\im(u)$ est de dimension finie. En effet, si $F$ est de dimension finie alors $\im(u)$ est un sev de $F$ donc de dimension finie. Si $E$ est de dimension finie, $u$ envoie une base de $E$ sur une famille génératrice de $\im(u)$, donc $\im(u)$ est de dimension finie.
\end{rmq}
\begin{thm}[Théorème du rang]
	Soit $E$ un $\K$-ev de dimension finie, $F$ un $\K$-ev quelconque et $u\in\mathcal{L}(E,F)$. Alors : $$\dim(E)=\dim(\ker(u))+rg(u)$$
\end{thm}
\begin{cor}
	Soit $E$ et $F$ des $\K$-ev de dimension finie tels que $\dim(E)=\dim(F)$. Soit $u\in\mathcal{L}(E,F)$. Alors les trois assertions suivantes sont équivalentes :\begin{enumerate}
		\item $u$ est injective
		\item $u$ est surjective
		\item $u$ est bijective
	\end{enumerate}
\end{cor}
\begin{rmq}
	L'exemple typique d'utilisation est lorsque $E$ est de dimension finie et $u\in\mathcal{L}(E)$.
\end{rmq}
\begin{cor}
	Soit $E$ de dimension finie et $u,v\in\mathcal{L}(E)$. Si $u\circ v=\id_E$, alors $v\circ u=\id_E$.
\end{cor}
\begin{rmq}
	Plus généralement, soit $E$ et $F$ de même dimension (finie) ainsi que $u\in\mathcal{L}(E,F)$ et $v\in\mathcal{L}(F,E)$. Si $u\circ v=\id_E$, alors $v\circ u=\id_F$.
\end{rmq}
\begin{prop}
	Soit $E$, $F$ et $G$ de dimension finie. Soit $u\in\mathcal{L}(E,F)$ et $v\in\mathcal{L}(F,G)$. Alors : $$rg(v\circ u)\leqslant\min(rg(u),rg(v))$$
\end{prop}
\begin{rmq}
	Par une récurrence immédiate, on a plus généralement des inégalités du type : $$rg(u_1\circ\ldots\circ u_p)\leqslant\underset{1\leqslant i\leqslant p}{\min}(rg(u_i))$$
\end{rmq}
\begin{prop}[Invariance du rang par composition par un isomorphisme]
	Soit $E$ et $F$ des $\K$-espaces vectoriels quelconques et soit $u\in\mathcal{L}(E,F)$ de rang fini. Alors : \begin{itemize}
		\item Soit $F'$ un espace vectoriel quelconque. Pour tout isomorphisme $\varphi:F\rightarrow F'$, $\varphi\circ u$ est de rang fini et on a : $$rg(\varphi\circ u)=rg(u)$$
		\item Soit $E'$ un espace vectoriel quelconque. Pour tout isomorphisme $\psi:E'\rightarrow E$, $u\circ\psi$ est de rang fini et on a : $$rg(u\circ\psi)=rg(u)$$
	\end{itemize}
\end{prop}
\begin{exa}
	Soit $u\in\mathcal{L}(E,F)$ de rang fini. Alors $-u$ est de rang fini et on a $rg(-u)=rg(u)$ (prendre $-\id_E$ comme isomorphisme).
\end{exa}
\section{Hyperplans et dualité}
Dans tout le paragraphe, on fixe $E$ un $\K$-espace vectoriel de dimension finie $n\in\N$ et $(e_1,\ldots,e_n)$ une base de $E$.
\begin{dfn}[Base duale]
	$(e_1\st,\ldots,e_n\st)$ est une base de $E\st$, appelée \textbf{base duale} de $(e_1,\ldots,e_n)$.
\end{dfn}
\begin{rmq}[Résolution d'exercices en dualité]
	Dans les exercices de dualité, deux techniques générales se distinguent.
	\begin{itemize}
		\item Si on dispose d'une égalité de vecteurs, il est judicieux d'appliquer une forme linéaire des deux côtés de cette égalité (le plus souvent, ce sera une $e_i\st$).
		\item Si on dispose d'une égalité de formes linéaires, il est judicieux d'évaluer cette égalité en un vecteur précis (le plus souvent, ce sera en un des $e_i$).
	\end{itemize}
\end{rmq}
\begin{rmq}[Note historique]
	A propos de cette remarque, le grand FM avait déclaré : "En prépa, beaucoup d'élèves ont peur de la dualité. C'est normal, parce qu'ils ont des mauvais profs. Mais moi, j'suis un bon prof !" (d'une voix enjouée pour la dernière phrase).
\end{rmq}
\begin{prop}[Expression d'une forme linéaire dans la base duale]
	Soit $\varphi\in E\st$. Ses coordonnées dans la base $(e_i\st)_{1\leqslant i\leqslant n}$ sont $(\varphi(e_i))_{1\leqslant i\leqslant n}$. Autrement dit, on a : $$\varphi=\sum_{i=1}^{n}\varphi(e_i)e_i\st$$
\end{prop}
\begin{thm}[Caractérisation des hyperplans par les équations]
	On a la caractérisation suivante des hyperplans par les équations en dimension finie.
	\begin{itemize}
		\item Soit $\displaystyle\begin{pmatrix}a_1\\ \vdots \\ a_n\end{pmatrix}\in\K^n\setminus\begin{pmatrix}0\\ \vdots\\ 0\end{pmatrix}$ et considérons l'ensemble $\displaystyle H_{a_1,\ldots,a_n} : \sum_{i=1}^{n}a_ix_i=0_\K$. Alors $\displaystyle H_{a_1,\ldots,a_n}$ est un hyperplan.
		\item Réciproquement, soit $H$ un hyperplan de $E$. Alors $\displaystyle\exists\displaystyle\begin{pmatrix}a_1\\ \vdots \\ a_n\end{pmatrix}\in\K^n\setminus\begin{pmatrix}0\\ \vdots\\ 0\end{pmatrix},\ H=H_{a_1,\ldots,a_n}$.
		\item De plus, soit $\displaystyle\displaystyle\begin{pmatrix}a_1\\ \vdots \\ a_n\end{pmatrix}\in\K^n\setminus\begin{pmatrix}0\\ \vdots\\ 0\end{pmatrix}$ et $\displaystyle\displaystyle\begin{pmatrix}b_1\\ \vdots \\ b_n\end{pmatrix}\in\K^n\setminus\begin{pmatrix}0\\ \vdots\\ 0\end{pmatrix}$ tels que $H=H_{a_1,\ldots,a_n}=H_{b_1,\ldots,b_n}$. Alors $$\exists\lambda\in\K\st,\ \begin{pmatrix}b_1\\ \vdots \\ b_n\end{pmatrix}=\lambda \begin{pmatrix}a_1\\ \vdots \\ a_n\end{pmatrix}$$
	\end{itemize}
\end{thm}
\begin{thm}[Caractérisation des hyperplans par la dimension en dimension finie]
	En dimension finie, les hyperplans de $E$ sont exactement les sev de dimension $\dim(E)-1$.
\end{thm}
\begin{thm}
	Soit $p\in\llbracket1,n\rrbracket$. 
	\begin{itemize}
		\item Soit $H_1,\ldots,H_p$ des hyperplans. Alors : $$\dim(H_1\cap\ldots\cap H_p)\geqslant n_p$$
		\item Réciproquement, soit $F$ un sev de dimension $n-p$. Alors il existe des hyperplans $H_1,\ldots,H$ tels que $$F=H_1\cap\ldots\cap H_n$$
	\end{itemize}
\end{thm}

\newpage
\chapter{Matrices et systèmes linéaires} %chapitre complet%
Dans tout le chapitre, on fixe $\K$ un corps.
\section{Calcul matriciel}
\subsection{Structure de $\K$-espace vectoriel}
\begin{dfn}[Matrices, ensemble $\mathcal{M}_{n,p}(\K)$]
	Soit $n,p\in\N\st$. Intuitivement, une matrice à $n$ lignes et $p$ colonnes est un tableau de taille $n\times p$ à coefficients das $\K$. \par Formellement, une matrice à $n$ lignes et $p$ colonnes à coefficients dans $\K$ est un triplet $\displaystyle\left(n,p,(a_{i,j})_{\substack{1\leqslant i\leqslant n \\ 1\leqslant j\leqslant p}}\right)$. \par Les $a_{i,j}$ sont des scalaires appelés \textbf{termes}, ou \textbf{coefficients} de la matrice. \par L'ensemble des matrices à $n$ lignes, $p$ colonnes et à coefficients dans $\K$ est noté $\mathcal{M}_{n,p}(\K)$.
\end{dfn}
\begin{dfn}[Matrices-colonne, matrices-ligne]
	Soit $n,p\in\N\st$. $\mathcal{M}_{n,1}(\K)$ est l'ensemble des \textbf{matrices-colonne} de taille $n$ à coefficients dans $\K$. $\mathcal{M}_{1,p}(\K)$ est l'ensemble des \textbf{matrices-ligne} de taille $p$ à coefficients dans $\K$.
\end{dfn}
\begin{dfn}[Lois de $\K$-ev de $\mathcal{M}_{n,p}(\K)$]
	On munit $\mathcal{M}_{n,p}(\K)$ des lois suivantes : 
	\begin{itemize}
		\item Un lci notée $+$: $$\left(n,p,(a_{i,j})_{\substack{1\leqslant i\leqslant n \\ 1\leqslant j\leqslant p}}\right)+\left(n,p,(b_{i,j})_{\substack{1\leqslant i\leqslant n \\ 1\leqslant j\leqslant p}}\right):=\left(n,p,(a_{i,j}+b_{i,j})_{\substack{1\leqslant i\leqslant n \\ 1\leqslant j\leqslant p}}\right)$$
		\item Une lce notée $\cdot$ : $$\lambda.\left(n,p,(a_{i,j})_{\substack{1\leqslant i\leqslant n \\ 1\leqslant j\leqslant p}}\right)=\left(n,p,(\lambda a_{i,j})_{\substack{1\leqslant i\leqslant n \\ 1\leqslant j\leqslant p}}\right)$$
	\end{itemize}
\end{dfn}
\begin{prop}[Structure de $\K$-ev de $\mathcal{M}_{n,p}(\K)$]
	Muni de ces deux lois, $\mathcal{M}_{n,p}(\K)$ est un $\K$-ev.
\end{prop}
\begin{dfn}[Matrices élémentaires]
	Soit $n,p\in\N\st$. Les \textbf{matrices élémentaires} de $\mathcal{M}_{n,p}(\K)$ sont définies par : $$\forall(i_0,j_0)\in\llbracket1,n\rrbracket\times\llbracket1,p\rrbracket,\ E_{i_0,j_0}:=\left(n,p,(\delta_{i,i_0}\delta_{j,j_0})_{\substack{1\leqslant i\leqslant n\\ 1\leqslant j\leqslant p}}\right)$$ Si le contexte est ambigu, on pourra les noter $E_{i_0,j_0}^{(n,p)}$.
\end{dfn}
\begin{thm}[Base canonique de $\mathcal{M}_{n,p}(\K)$, dimension de $\mathcal{M}_{n,p}(\K)$]
	La famille $\displaystyle\left(E_{i,j}^{(n,p)}\right)_{\substack{1\leqslant i\leqslant n\\ 1\leqslant j\leqslant p}}$ est une base de $\mathcal{M}_{n,p}(\K)$, appelée \textbf{base canonique} de $\mathcal{M}_{n,p}(\K)$. En particulier, $\mathcal{M}_{n,p}(\K)$ est un $\K$-espace vectoriel de dimension finie égale à $n\times p$.
\end{thm}
\begin{rmq}
	Dès lors, on identifiera $\K^n$ à $\mathcal{M}_{n,1}(\K)$ car ces espaces sont isomorphes.
\end{rmq}
\subsection{Produit matriciel}
\begin{dfn}[Produit matriciel]
	On définit une loi $\times$ sur les matrices par : $$\left(n,p,(a_{i,j})_{\substack{1\leqslant i\leqslant n \\ 1\leqslant j\leqslant p}}\right)\times\left(p,q,(b_{i,j})_{\substack{1\leqslant i\leqslant p \\ 1\leqslant j\leqslant q}}\right):=\left(n,q,\left(\sum_{k=1}^{p}a_{i,k}b_{k,j}\right)_{\substack{1\leqslant i\leqslant n \\ 1\leqslant j\leqslant q}}\right)$$
\end{dfn}
\begin{rmq}
	Attention à bien prendre garde à la compatibilité des tailles ! Le nombre de lignes de la deuxième matrice doit être égal au nombre de colonnes de la première.
\end{rmq}
\begin{exa}
	Soit $A\in\mathcal{M}_{n,p}(\K)$ et $B\in\mathcal{M}_{p,q}(\K)$. On a : $$\begin{cases}
		0_{n,p}\times B=0_{n,q} \\ A\times 0_{p,q}=0_{n,q}
	\end{cases}$$
\end{exa}
\begin{exa}[Produit d'une matrice-ligne et d'une matrice-colonne]
	Soit $p\in\N\st$. On a : $$(\lambda_1,\ldots,\lambda_p)\times\begin{pmatrix}	\mu_1\\ \vdots \\ \mu_p\end{pmatrix}=\left(\sum_{k=1}^{p}\lambda_k\mu_k\right)$$ et $$\begin{pmatrix}	\mu_1\\ \vdots \\ \mu_p\end{pmatrix}\times(\lambda_1,\ldots,\lambda_p)=\left(p,p,\left(\mu_i\lambda_j\right)_{\substack{1\leqslant i\leqslant p\\ 1\leqslant j\leqslant p}}\right)$$
\end{exa}
\begin{exa}[$\mathcal{M}_{1,1}(\K)$]
	On a $(\lambda)\times(\mu)=(\lambda\times\mu)$. Comme de plus $(\lambda)+(\mu)=(\lambda+\mu)$, on pourra identifier $\K$ et $\mathcal{M}_{1,1}(\K)$.
\end{exa}
\begin{exa}[Utile]
	Soit $\lambda\in\K$ et $\displaystyle X=\begin{pmatrix}x_1\\ \vdots \\ x_n\end{pmatrix}\in\K^n$. Alors on a : $$\lambda.\begin{pmatrix}x_1\\ \vdots \\ x_n\end{pmatrix}=\begin{pmatrix}x_1\\ \vdots \\ x_n\end{pmatrix}\times(\lambda)$$
\end{exa}
\begin{thm}[Produit de matrices élémentaires]
	Soit $n,p,q\in\N\st$. Soit $(i,j)\in\llbracket1,n\rrbracket\times\llbracket1,p\rrbracket$. et $(k,l)\in\llbracket1,p\rrbracket\times\llbracket1,q\rrbracket$. On a : $$E_{i,j}^{(n,p)}\times E_{k,l}^{(p,q)}=\delta_{j,k}E_{i,l}^{(n,q)}$$
\end{thm}
\begin{thm}
	La loi $\times$ est associative et bilinéaire.
\end{thm}
\begin{exa}
	On en déduit la compatibilité du produit matriciel avec la loi externe er la distributivité du produit matriciel sur $+$.
\end{exa}
\begin{thm}[Produit par blocs]
	Soit $n_1,n_2,p_1,p_2,q_1,q_2\in\N\st$, $A\in\mathcal{M}_{n_1,p_1}(\K)$, $B\in\mathcal{M}_{n_1,p_2}(\K)$, $C\in\mathcal{M}_{n_2,p_1}(\K)$, $D\in\mathcal{M}_{n_2,p_2}(\K)$, $A'\in\mathcal{M}_{p_1,q_1}(\K)$, $B'\in\mathcal{M}_{p_1,q_2}(\K)$, $C'\in\mathcal{M}_{p_2,q_1}(\K)$ et $D'\in\mathcal{M}_{p_2,q_2}(\K)$ (en gros, il faut que les tailles de blocs soient compatibles). Alors on a : $$\begin{pmatrix}A&B\\C&D\end{pmatrix}\times\begin{pmatrix}A'&B'\\C'&D'\end{pmatrix}=\begin{pmatrix}AA'+BC'&AB'+BD'\\CA'+DC'&CB'+DD'\end{pmatrix}$$ On retiendra qu'on peut multiplier les blocs comme si c'étaient les coefficients d'une matrice normale (du moment que les blocs sont compatibles).
\end{thm}
\begin{exa}
	Voici trois cas particuliers à connaître sur le bout des doigts : \begin{itemize}
		\item Colonne $\times$ ligne : $$\begin{pmatrix}\\C\\  \\ \end{pmatrix}\times\begin{pmatrix}\lambda_1&|&\ldots&|&\lambda_n\end{pmatrix}=\begin{pmatrix}&&&&\\ C\times(\lambda_1)&|&\ldots&|&C\times(\lambda_n)\\ &&&&\end{pmatrix}=\begin{pmatrix}&&&&\\ \lambda_1 C&|&\ldots&|&\lambda_n C\\ &&&&\end{pmatrix}$$
		\item Matrice $\times$ colonne : $$\begin{pmatrix}&&&&\\C_1&|&\ldots&|&C_p\\ &&&&\end{pmatrix}\times\begin{pmatrix}\lambda_1\\ \vdots\\ \lambda_p\end{pmatrix}=\begin{pmatrix}\\C_1\times(\lambda_1)+\ldots+C_p\times(\lambda_p)\\ \\ \end{pmatrix}=\begin{pmatrix}\\ \displaystyle\sum_{i=1}^{p}\lambda_iC_i\\ \\ \end{pmatrix}$$
		\item Matrice $\times$ matrice : $$\begin{pmatrix}&&\\ &A&\\ &&\end{pmatrix}\times\begin{pmatrix}&&&&\\C_1&|&\ldots&|&C_q\\&&&&\end{pmatrix}=\begin{pmatrix}&&&&\\AC_1&|&\ldots&|&AC_q\\ &&&&\end{pmatrix}$$
	\end{itemize}
\end{exa}
\subsection{Matrices carrées}
\begin{dfn}[Matrices carrées]
	$\mathcal{M}_{n,n}(\K)$ est noté plus simplement $\mathcal{M}_n(\K)$. C'est l'ensemble des \textbf{matrices carrées d'ordre $n$ à coefficients dans $\K$}. Il s'agit donc d'un $\K$-ev de dimension $n^2$.
\end{dfn}
\begin{dfn}[Matrice identité]
	La \textbf{matrice identité d'ordre $n$} est la matrice carrée d'ordre $n$ définie par : $$I_n:=\begin{pmatrix}1&&(0)\\&\ddots&\\ (0)&&1\end{pmatrix}$$
\end{dfn}
\begin{thm}
	$(\mathcal{M}_n(\K),+,\cdot,\times)$ est une $\K$-algèbre, d'élément neutre $I_n$ pour $\times$.
\end{thm}
\begin{rmq}
	Plus généralement, soit $A\in\mathcal{M}_{n,p}(\K)$. On a : $$\begin{cases}
		A\times I_p=A\\ I_n\times A=A
	\end{cases}$$
\end{rmq}
\begin{rmq}
	$\mathcal{M}_n(\K)$ pour $n\in\N\st$ est commutatif si, et seulement si $n=1$. Dès que $n\geqslant2$, il suffit de considérer les produits $E_{1,2}\times E_{2,1}$ et $E_{2,1}\times E_{1,2}$.
\end{rmq}
\begin{rmq}
	Dès que $n\geqslant2$, $\mathcal{M}_n(\K)$ n'est pas intègre car $E_{1,2}\times E_{1,2}=0_n$.
\end{rmq}
\begin{dfn}[Groupe linéaire d'ordre $n$]
	On appelle \textbf{groupe linéaire d'ordre $n$} le groupe des unités de l'anneau $\mathcal{M}_n(\K)$. On le note $GL_n(\K)$ et on a donc : $$GL_n(\K):=\mathcal{U}\left(\mathcal{M}_n(\K)\right)=\left\{A\in\mathcal{M}_n(\K)\ | \ \exists B\in\mathcal{M}_n(K),\ AB=BA=I_n\right\}$$
\end{dfn}
\begin{exa}
	On rappelle que d'après le chapitre sur les groupes, on a alors : $$\forall(A,B)\in GL_n(\K)^2,\ (AB)\inv=B\inv A\inv$$
\end{exa}
\begin{rmq}
	Deux matrices de $GL_n(\K)$ n'ont aucune raison de commuter dès que $n\geqslant2$. Pour $n=2$, il suffit de considérer $\displaystyle\begin{pmatrix}1&1\\0&2\end{pmatrix}$ et $\displaystyle\begin{pmatrix}1&2\\0&1\end{pmatrix}$. Pour l'ordre $n\geqslant3$, il suffit de considérer les mêmes matrices mais complétées par un bloc $I_{n-2}$ et des zéros partout ailleurs.
\end{rmq}
\begin{met}
	Pour trouver un contre-exemple à l'ordre $n\geqslant2$, on peut souvent trouver un contre-exemple pur $n=2$ puis le généraliser pour $n\geqslant2$ avec des blocs. \\ \par Par exemple, si $P$ et $Q$ sont des matrices inversibles, alors la matrice $\displaystyle\begin{pmatrix}P&(0)\\ (0)&Q\end{pmatrix}$ est inversible d'inverse $\displaystyle\begin{pmatrix}P\inv&(0)\\ (0)&Q\inv\end{pmatrix}$, comme on le vérifie immédiatement par un calcul par blocs. 
\end{met}
\begin{prop}
	Soit $A,B\in\mathcal{M}_n(\K)$. Si $AB=BA$. Alors, on a : \begin{itemize}
		\item La \textbf{formule du binôme de Newton} : $$\forall p\in\N,\ (A+B)^p=\sum_{k=0}^{p}\binom{p}{k}A^kB^{p-k}$$
		\item La \textbf{formule de Bernoulli} : $$\forall p\in\N\st,\ A^p-B^p=(A-B)\left(\sum_{k=0}^{p_1}A^{p-1-k}B^k\right)=\left(\sum_{k=0}^{p_1}A^{p-1-k}B^k\right)(A-B)$$
	\end{itemize}
\end{prop}
\begin{rmq}
	On pensera souvent au cas particulier $B=I_n$ qui fonctionne toujours, puisque $I_n$ commute avec toute matrice en tant que neutre.
\end{rmq}
\begin{dfn}[Trace]
	Soit $A=(a_{i,j})_{\substack{1\leqslant i\leqslant n\\ 1\leqslant j\leqslant n}}\in\mathcal{M}_n(\K)$. La \textbf{trace} de $A$ est le \textbf{scalaire} défini par : $$\tr(A):=\sum_{i=1}^{n}a_{i,i}$$ Autrement dit, la trace d'une matrice est la somme de ses coefficients diagonaux.
\end{dfn}
\begin{rmq}
	Attention : la trace n'est définie que pour les matrices carrées !
\end{rmq}
\begin{thm}
	L'application $\tr(\cdot)$ est une forme linéaire non-nulle sur $\mathcal{M}_n(\K)$.
\end{thm}
\begin{exa}
	Ainsi, $\ker(\tr(\cdot))$ est un hyperplan de $\mathcal{M}_n(\K)$. C'est donc un espace de dimension $n^2-1$.
\end{exa}
\begin{thm}[Exercice classique / Utile pour exos difficiles, HP]
	Soit $\varphi\in\mathcal{M}_n(\K)\st$. Alors, $$\exists!A\in\mathcal{M}_n(\K),\ \varphi=\tr(A\times\cdot)$$
\end{thm}
\begin{proof}
	Pour tout $A\in\mathcal{M}_n(\K)$, on considère l'application $$\begin{array}{ccccc}\varphi_A&:&\mathcal{M}_n(\K)&\rightarrow&\K\\&&M&\mapsto&\tr(AM)\end{array}$$ puis on considère $$\begin{array}{ccccc}\theta&:&\mathcal{M}_n(\K)&\rightarrow&\mathcal{M}_n(\K)\st\\&&A&\mapsto&\varphi_A\end{array}$$ \par Déjà, $\theta$ est bien définie puisque toutes les application $\varphi_A$ sont des formes linéaires par linéarité à droite du produit matriciel et par linéarité de la trace. \par Ensuite, $\theta$ est linéaire par linéarité à gauche du produit matriciel et par linéarité de la trace.\par De plus, $\theta$ est injective. En effet, puisque $\theta$ est linéaire, soit $A\in\mathcal{M}_n(\K)$ telle que $\varphi_A$ soit la forme linéaire nulle. Alors pour tous $(i,j)\in\llbracket1,n\rrbracket^2$, on a $\varphi_A(E_{j,i})=0_\K$, ce qui après calcul, puisque \begin{align*}
		\varphi_A(E_{j,i})&=\tr(AE_{j,i})\\
		&=\sum_{k=1}^{n}(AE_{j,i})_{k,k}\\
		&=\sum_{k=1}^{n}\sum_{l=1}^{n}a_{k,l}\delta_{j,l}\delta_{i,k}\\
		&=a_{i,j}
	\end{align*} donne $a_{i,j}=0_\K$. Ainsi, $A$ est la matrice nulle, donc $\theta$ est injective. \par Enfin, par égalité de dimensions, $\theta$ est bijective, ce qui conclut.
\end{proof}
\begin{exe}[Oral X/ENS]
	On peut utiliser ce théorème pour prouver que pour $n\geqslant2$, tout hyperplan de $\mathcal{M}_n(\K)$ rencontre $GL_n(\K)$ en écrivant cet hyperplan comme le noyau d'une $\tr(A\times\cdot)$ avec $A$ non nulle. Il faut alors raisonner sur le rang de $A$ en l'écrivant comme une matrice équivalente à une matrice $J_r$. On pourra utiliser la matrice $$N:=\begin{pmatrix}0&&&0&1\\1&\ddots&&(0)&0\\&\ddots&\ddots&&\\&&\ddots&\ddots&\\ (0)&&&1&0\end{pmatrix}$$
\end{exe}
\begin{thm}
	Soit $A,B\in\mathcal{M}_n(\K)$. On a : $$\tr(AB)=\tr(BA)$$
\end{thm}
\begin{rmq}
	Si $A\in\mathcal{M}_{n,p}(\K)$ et $B\in\mathcal{M}_{p,n}(\K)$, on montre de même que $\tr(AB)=\tr(BA)$.
\end{rmq}
\subsection{Matrices carrées particulières}
Dans tout le paragraphe, on fixe $n\in\N\st$ et on s'intéresse à des matrices particulières de $\mathcal{M}_n(\K)$
\begin{dfn}[Matrice diagonale]
	Une matrice $A\in\mathcal{M}_n(\K)$ est dite \textbf{diagonale} lorsqu'elle est de la forme $$\begin{pmatrix}\lambda_1&&(0)\\&\ddots&\\ (0)&&\lambda_n\end{pmatrix}$$ On la note alors $diag(\lambda_1,\ldots,\lambda_n)$ Formellement, $A$ est diagonale lorsque tous ses coefficients non diagonaux sont nuls, $ie$ lorsque $$\forall(i,j)\in\llbracket1,n\rrbracket^2,\ i\ne j\implies a_{i,j}=0_\K$$ On notera $D_n(\K)$ l'ensemble des matrices diagonales d'ordre $n$ (NS).
\end{dfn}
\begin{dfn}[Matrice scalaire]
	Une matrice $A\in\mathcal{M}_n(\K)$ est dite \textbf{scalaire} lorsque $$\exists\lambda\in\K,\ A=\lambda I_n$$ Les matrices scalaires sont donc en particulier des matrices diagonales.
\end{dfn}
\begin{prop}[Produit de matrices diagonales]
	Le produit de deux matrices diagonales est une matrice diagonale dont le coefficient en position $(i,i)$ est le produit des coefficients en position $(i,i)$ des deux matrices diagonales de départ.
\end{prop}
\begin{thm}[Produit diagonal par blocs]
	Soit $A_1$, $A_2$, $B_1$ et $B_2$ des matrices compatibles en taille. On a : $$\begin{pmatrix}A_1&(0)\\ (0)&A_2\end{pmatrix}\times\begin{pmatrix}B_1&(0)\\ (0)&B_2\end{pmatrix}=\begin{pmatrix}A_1B_1&(0)\\ (0)&A_2B_2\end{pmatrix}$$
\end{thm}
\begin{thm}
	$D_n(\K)$ est une sous-algèbre commutative de $\mathcal{M}_n(\K)$ de dimension $n$. En effet, on a $D_n(\K)=\vect(E_{i,i})_{1\leqslant i\leqslant n}$.
\end{thm}
\begin{prop}
	$D_n(\K)\setminus\left\{0_{n,n}\right\}$ est un sous-groupe commutatif de 
\end{prop}
\begin{dfn}[Matrice triangulaire supérieure]
	Une matrice $A\in\mathcal{M}_n(\K)$ est dite \textbf{triangulaire supérieure} lorsqu'elle est de la forme $$\begin{pmatrix}*&&(*)\\&\ddots&\\ (0)&&*\end{pmatrix}$$ Formellement, $A$ est triangulaire supérieure lorsque $$\forall(i,j)\in\llbracket1,n\rrbracket^2,\ i>j\implies a_{i,j}=0_\K$$ On notera $T_n^+(\K)$ l'ensemble des matrices triangulaires supérieures d'ordre $n$ (NS).
\end{dfn}
\begin{dfn}[Matrice triangulaire supérieure stricte]
	Une matrice $A\in\mathcal{M}_n(\K)$ est dite \textbf{triangulaire supérieure stricte} lorsqu'elle est de la forme $$\begin{pmatrix}0&&(*)\\&\ddots&\\ (0)&&0\end{pmatrix}$$ Formellement, $A$ est triangulaire supérieure stricte lorsque $$\forall(i,j)\in\llbracket1,n\rrbracket^2,\ i\geqslant j\implies a_{i,j}=0_\K$$ On notera $T_n^{++}(\K)$ l'ensemble des matrices triangulaires supérieures strictes d'ordre $n$ (NS).
\end{dfn}
\begin{exa}
	On a $T_n^{++}(\K)\subset T_n^+(\K)$.
\end{exa}
\begin{dfn}[Matrice triangulaire inférieure]
	Une matrice $A\in\mathcal{M}_n(\K)$ est dite \textbf{triangulaire inférieure} lorsqu'elle est de la forme $$\begin{pmatrix}*&&(0)\\&\ddots&\\ (*)&&*\end{pmatrix}$$ Formellement, $A$ est triangulaire inférieure lorsque $$\forall(i,j)\in\llbracket1,n\rrbracket^2,\ i<j\implies a_{i,j}=0_\K$$ On notera $T_n^-(\K)$ l'ensemble des matrices triangulaires inférieures d'ordre $n$ (NS).
\end{dfn}
\begin{dfn}[Matrice triangulaire inférieure stricte]
	Une matrice $A\in\mathcal{M}_n(\K)$ est dite \textbf{triangulaire inférieure stricte} lorsqu'elle est de la forme $$\begin{pmatrix}0&&(0)\\&\ddots&\\ (*)&&0\end{pmatrix}$$ Formellement, $A$ est triangulaire inférieure stricte lorsque $$\forall(i,j)\in\llbracket1,n\rrbracket^2,\ i\leqslant j\implies a_{i,j}=0_\K$$ On notera $T_n^{--}(\K)$ l'ensemble des matrices triangulaires inférieures strictes d'ordre $n$ (NS).
\end{dfn}
\begin{exa}
	On a $T_n^{--}(\K)\subset T_n^-(\K)$.
\end{exa}
\begin{prop}
	Le produit de deux matrices triangulaires supérieures (resp. inférieures) est une matrice triangulaire supérieure (resp. inférieure) dont les termes diagonaux sont le produit des termes diagonaux correspondants.
\end{prop}
\begin{thm}
	$T_n^+(\K)$ et $T_n^-(\K)$ sont des sous-algèbres de $\mathcal{M}_n(\K)$ de dimension $\displaystyle n\times\frac{n+1}{2}$.
\end{thm}
\begin{thm}
	$T_n^{++}(\K)$ et $T_n^{--}(\K)$ sont des sous-espaces vectoriels de $\mathcal{M}_n(\K)$ de dimension $\displaystyle n\times\frac{n-1}{2}$.
\end{thm}
\begin{rmq}
	On a $T_n^+(\K)\oplus T_n^{--}(\K)=\mathcal{M}_n(\K)$ et $T_n^-(\K)\oplus T_n^{++}(\K)=\mathcal{M}_n(\K)$.
\end{rmq}
\begin{prop}
	Soit $T\in T_n^+(\K)\cap GL_n(\K)$. Alors $T\inv\in T_n^+(\K)$. Idem avec $T_n^-(\K)$.
\end{prop}
\begin{thm}[HP]
	Plus généralement, soit $\mathcal{A}$ une sous-algèbre de $\mathcal{M}_n(\K)$. Alors $$\forall A\in\mathcal{A}\cap GL_n(\K),\ A\inv\in\mathcal{A}$$
\end{thm}
\begin{proof}
	La preuve du théorème précédent s'adapte. Soit $A\in\mathcal{A}\cap GL_n(\K)$. Posons $$\begin{array}{ccccc}f_A&:&\mathcal{A}&\rightarrow&\mathcal{A}\\&&B&\mapsto&AB\end{array}$$ $f_A$ est bien définie par stabilité de $\mathcal{A}$ par $\times$. Elle est linéaire par linéarité à droite par linéarité du produit matriciel, et elle est injective car $A\in GL_n(\K)$ (il suffit de considérer un élément $B$ du noyau de $f_A$ puis multiplier $AB=0_n$ par $A\inv$ pour obtenir $B=0_n$). Par égalité de dimensions, $f_A$ est surjective donc il existe $B\in\mathcal{A}$ telle que $AB=I_n$. Donc $A\inv=B\in\mathcal{A}$.
\end{proof}
\begin{exe}
	Soit $A\in GL_n(\K)$. Montrer que $A\inv$ peut s'écrire comme un polynôme en $A$. Indication : considérer $\vect(A^k)_{k\in\N}$ et montrer que c'est une sous-algèbre de $\mathcal{M}_n(\K)$.
\end{exe}
\subsection{Transposition}
Dans tout le paragraphe, on fixe $n,p\in\N\st$.
\begin{dfn}[Transposée]
	Soit $A\in\mathcal{M}_{n,p}(\K)$. La \textbf{transposée de $A$}, notée $A^\top$ ou anciennement $^tA$, est la matrice de $\mathcal{M}_{p,n}(\K)$ définie par : $$A^\top:=\begin{pmatrix}a_{1,1}&\ldots&a_{n,1}\\ \vdots&&\vdots\\ a_{1,p}&\ldots&a_{n,p}\end{pmatrix}=(a_{j,i})_{\substack{1\leqslant j\leqslant p \\ 1\leqslant i\leqslant n}}$$
\end{dfn}
\begin{prop}[Propriétés de la transposée]
	La transposée vérifie les propriétés suivantes :\begin{itemize}
		\item $\forall A\in\mathcal{M}_{n,p}(\K),\ (A^\top)^\top=A$ (involutivité)
		\item $\forall(\lambda,\mu)\in\K^2,\ \forall (A,B)\in\mathcal{M}_n(\K)^2,\ (\lambda A+\mu B)^\top=\lambda A^\top+\mu B^\top$ (linéarité)
		\item $\forall(A,B)\in\mathcal{M}_{n,p}(\K)\times\mathcal{M}_{p,q}(\K),\ (AB)^\top=B^\top A^\top$ (transposée d'un produit)
		\item $\forall A\in\ GL_n(\K),\ A^\top\in GL_n(\K)\land (A^\top)\inv=(A\inv)^\top$ (inversibilité)
		\item $\forall A\in\mathcal{M}_n(\K),\ \tr(A^\top)=\tr(A)$ (invariance de la trace)
	\end{itemize}
\end{prop}
\begin{rmq}
	L'application de $\mathcal{M}_{n,p}(\K)$ dans $\mathcal{M}_{p,n}(\K)$ qui à une matrice associe sa transposée est bijective. En effet, on peut exhiber sa bijection réciproque, qui est la même fonction mais avec les domaines inversés.
\end{rmq}
\begin{dfn}[Matrice symétrique]
	Soit $A\in\mathcal{M}_n(\K)$. $A$ est dite \textbf{symétrique} lorsque $A^\top=A$. On note $S_n(\K)$ l'ensemble des matrices symétriques
\end{dfn}
\begin{dfn}[Matrice antisymétrique]
	Soit $A\in\mathcal{M}_n(\K)$. $A$ est dite \textbf{antisymétrique} lorsque $A^\top=-A$. On note $A_n(\K)$ l'ensemble des matrices antisymétriques
\end{dfn}
\begin{prop}
	$S_n(\K)$ et $A_n(\K)$ sont des sev de $\mathcal{M}_n(\K)$.
\end{prop}
\begin{thm}
	On a $$S_n(\K)\oplus A_n(\K)=\mathcal{M}_n(\K)$$ En effet, l'application "transposition" (qu'on notera ici pour quelques instants $s$)est une symétrie vectorielle de $\mathcal{M}_n(\K)$ et $S_n(\K)$ et $A_n(\K)$ sont respectivement $\ker(s-\id)$ et $\ker(s+\id)$.
\end{thm}
\begin{thm}
	On a $\dim(S_n(\K))=n\times\displaystyle\frac{n+1}{2}$ et $\dim(A_n(\K))=n\times\displaystyle\frac{n-1}{2}$.
\end{thm}
\begin{proof}
	Soit $(i,j)\in\llbracket1,n\rrbracket^2$. Si $i=j$, on pose $F_{i,j}:=E_{i,j}$. Si $i<j$, on pose $F_{i,j}:=E_{i,j}+E_{j,i}$. $(F_{i,j})_{1\leqslant i\leqslant j\leqslant n}$ est alors une famille libre de $S_n(\K)$ donc $\dim(S_n(\K))\geqslant\displaystyle n\times\frac{n+1}{2}$. \par Aussi, la famille $(E_{i,j}-E_{j,i})_{1\leqslant<j\leqslant n}$ est une famille libre de $A_n(\K)$ donc $\dim(A_n(\K))\geqslant n\times\frac{n-1}{2}$. \par En sommant, par supplémentarité, on remarque que ces inégalités doivent être des égalités. Donc on a bien les égalités de dimensions souhaitées, et les familles que nous avons considérées sont en réalités des bases de leurs espaces respectifs.
\end{proof}
\section{Matrices et applications linéaires}
\subsection{Principe de correspondance}
\begin{dfn}[Matrice d'un vecteur dans une base]
	Soit $E$ de dimension $n\in\N\st$ et $(e_1,\ldots,e_n)$ une base de $E$. A tout vecteur $x\in E$, on peut faire correspondre une matrice de $\mathcal{M}_{n,1}(\K)$ appelée \textbf{matrice du vecteur $x$ dans la base $(e_i)$}, notée $\mat_{(e_i)}(x)$, et définie par : $$\mat_{(e_i)}(x):=\begin{pmatrix}e_1\st(x)\\ \vdots\\ e_n\st(x)\end{pmatrix}$$
\end{dfn}
\begin{prop}
	L'application $$\begin{array}{ccccc}\mat_{(e_i)}&:&E&\rightarrow&\mathcal{M}_{n,1}(\K)\\&&x&\mapsto&\mat_{(e_i)}(x)\end{array}$$ est un isomorphisme de $\K$-espaces vectoriels.
\end{prop}
\begin{exa}
	Si on prend $E=\K^n$ et $(e_1,\ldots,e_n)$ la base canonique de $\K^n$, puisqu'on identifie $\K^n$ et $\mathcal{M}_{n,1}(\K)$, on a alors $x=\mat_{(e_i)}(x)$.
\end{exa}
\begin{rmq}
	Désormais, les lettres X, Y et Z seront réservées pour des matrices-colonne, tandis que A, B et C seront plutôt utilisées pour des matrices moins spécifiques.
\end{rmq}
\begin{dfn}[Matrice d'une famille dans une base]
	Soit $E$ de dimension $n\in\N\st$ et $(e_1,\ldots,e_n)$ une base de $E$. Soit $p\in\N\st$ et $(x_1,\ldots,x_p)\in E^p$. La matrice de la famille $(x_j)_{1\leqslant j\leqslant p}$ dans la base $(e_i)_{1\leqslant i\leqslant n}$ est définie en juxtaposant les matrices-colonne de $x_j$ dans la base $(e_i)$ : $$\mat_{(e_i)}\left[(x_1,\ldots,x_p)\right]:=\begin{pmatrix}&\vline&&\vline&\\\mat_{(e_i)}(x_1)&\vline&\ldots&\vline&\mat_{(e_i)}(x_p)\\&\vline&&\vline&\end{pmatrix}$$ Autrement dit, $\mat_{(e_i)}\left[(x_1,\ldots,x_p)\right]$ est la matrice des $e_i\st(x_j)$.
\end{dfn}
\begin{exa}
	En reprenant $E=\K^n$ et $(e_i)$ la base canonique, on a : $$\mat_{(e_i)}\left[(x_1,\ldots,x_p)\right]=\begin{pmatrix}&\vline&&\vline&\\x_1&\vline&\ldots&\vline&x_p\\&\vline&&\vline&\end{pmatrix}$$
\end{exa}
\begin{dfn}[Matrice d'un application linéaire relativement à une base de départ et à une base d'arrivée]
	Soit $E$ un espace vectoriel de dimension finie $p\in\N\st$ et $(e_1,\ldots,e_p)$ une base de $E$ ; et $F$ un espace vectoriel de dimension finie $n\in\N\st$ et $(f_1,\ldots,f_n)$ une base de $F$. Soit $u\in\mathcal{L}(E,F)$. On rappelle que $u$ est entièrement caractérisée par l'image de la base $(e_j)$ par $u$. La \textbf{matrice de $u$ relativement aux base $(e_j)$ et $(f_i)$} est la matrice définie par : $$\mat_{(e_j),(f_i)}(u):=\mat_{(f_i)}\left[(u(e_1),\ldots,u(e_p))\right]$$ Autrement dit, $\mat_{(e_j),(f_i)}(u)$ est la matrice des $f_i\st(u(e_j))$. La base $(e_j)$ est dite \textbf{base de départ} et la base $(f_i)$ est dite \textbf{base d'arrivée}.
\end{dfn}
\begin{thm}
	L'application $$\begin{array}{ccccc}\mat_{(e_j),(f_i)}&:&\mathcal{L}(E,F)&\rightarrow&\mathcal{M}_{n,p}(\K)\\&&x&\mapsto&\mat_{(e_j),(f_i)}(x)\end{array}$$ est un isomorphisme de $\K$-espaces vectoriels.
\end{thm}
\begin{rmq}[Cas particulier des endomorphismes]
	Dans le cas d'un endomorphisme (on prend $E=F$), et si on prend $(e_i)=(f_i)$, la matrice $\mat_{(e_i),(e_i)}(u)$ sera plutôt notée $\mat_{(e_i)}(u)$.
\end{rmq}
\begin{exa}[Homothéties]
	Soit $\lambda\in\K$ et $(e_i)_{1\leqslant i\leqslant n}$ une base de $E$. On a : $$\mat_{(e_i)}(\lambda\id_E)=\lambda I_n$$
\end{exa}
\begin{exa}[Formes linéaires]
	Soit $(e_1,\ldots,e_p)$ une base de $E$ et choisissons $(1_\K)$ comme base du $\K$-ev $\K$. Soit $\psi\in E\st$. On a : $$\mat_{(e_j),(1_\K)}(\psi)=\begin{pmatrix}\psi(e_1)&\ldots&\psi(e_p)\end{pmatrix}\in\mathcal{M}_{1,p}(\K)$$ On réserve donc les matrices-ligne aux formes linéaires.
\end{exa}
\begin{exa}[Similitude directe linéaire]
	On voit $\C$ comme un $\R$-ev de dimension $2$, dont on choisit la base $(1,i)$. Soit $a\in\C\st$ et on considère la similitude directe $s:z\mapsto az$ qui est un endomorphisme de $\C$. On a : $$\mat_{(1,i)}(s)=\begin{pmatrix}\text{Re}(a)&-\text{Im}(a)\\ \text{Im}(a)&\text{Re}(a)\end{pmatrix}$$
\end{exa}
\begin{cor}[Dimension de $\mathcal{L}(E,F)$]
	$\mathcal{L}(E,F)$ est de dimension finie égale à $\dim(E)\times\dim(F)$.
\end{cor}
\begin{rmq}
	On avait décidé dans le chapitre "Dimension finie", d'admettre ce fait car la démonstration est beaucoup plus aisée lorsqu'on dispose du théorème qui précède.
\end{rmq}
\begin{dfn}[Application linéaire canoniquement associée à une matrice]
	Soit $A\in\mathcal{M}_{n,p}(\K)$. \textbf{L'application linéaire canoniquement associée à $A$} (abrégée en \textbf{ALCA}) est l'unique $u\in\mathcal{L}(\K^p,\K^n)$ telle que $$\mat_{(e_j),(f_i)}(u)=A$$ où $(e_j)$ désigne la base canonique de $\K^p$ et $(f_i)$ désigne la base canonique de $\K^n$.
\end{dfn}
\begin{dfn}[Image d'une matrice]
	Soit $A$ une matrice. L'\textbf{image de $A$}, notée $\im(A)$, est définie comme étant égale à $\im(u)$ où $u$ est l'application linéaire canoniquement associée à $A$.
\end{dfn}
\begin{dfn}[Noyau d'une matrice]
	Soit $A$ une matrice. Le \textbf{noyau de $A$}, notée $\ker(A)$, est défini comme étant égal à $\ker(u)$ où $u$ est l'application linéaire canoniquement associée à $A$.
\end{dfn}
\begin{thm}[Traduction matricielle d'une équation linéaire]
	"L'équation $u(x)=y$ se traduit matriciellement $AX=Y$." Formellement, soit $x\in E$, $u\in\mathcal{L}(E,F)$ et posons $y:=u(x)$. Soit $(e_1,\ldots,e_p)$ une base de $E$ et $(f_1,\ldots,f_n)$ une base de $f$. Posons $X:=\mat_{(e_j)}(x)\in\mathcal{M}_{p,1}(\K)$, $A:=\mat_{(e_j),(f_i)}(u)\in\mathcal{M}_{n,p}(\K)$ et $Y:=\mat_{(f_i)}(y)\in\mathcal{M}_{n,1}(\K)$. On a alors : $$AX=Y$$
\end{thm}
\begin{cor}[Formule explicite de l'ALCA]
	L'application linéaire canoniquement associée à $A\in\mathcal{M}_{n,p}(\K)$ est l'application : $$\begin{array}{ccc}\K^p&\rightarrow&\K^n\\X&\mapsto &AX\end{array}$$
\end{cor}
\begin{rmq}[Détermination pratique de l'image de $A$]
	L'image de $A$ est le sous-espace engendré par ses colonnes vues comme des vecteurs de $\K^n$.
\end{rmq}
\begin{rmq}[Détermination pratique du noyau de $A$]
	Les lignes de $A$ fournissent un système d'équations du noyau de $A$.
\end{rmq}
\begin{met}
	Soit $A\in\mathcal{M}_n(\K)$. On voudrait savoir si $A$ est inversible, et si oui, on voudrait calculer son inverse. \begin{enumerate}
		\item On résout $AX=Y$ pour $X,Y\in\K^n$. Si on obtient une unique solution, alors $A$ est inversible. Si on n'obtient pas une unique solution quel que soit $Y$, alors $A$ n'est pas inversible.
		\item On se place dans le cas où $A$ est inversible. Voici comment obtenir $A\inv$. Par hypothèse, $$\forall Y\in\K^n,\ \exists!X\in\K^n,\ AX=Y$$ On remplace alors successivement $Y$ par $e_1$, ..., $e_n$ la base canonique de $\K^n$. On obtient alors à chaque fois une unique solution $X_1$, ..., $X_n$. On a alors : $$A\inv=\begin{pmatrix}&\vline&&\vline&\\X_1&\vline&\ldots&\vline&X_n\\&\vline&&\vline&\end{pmatrix}$$
	\end{enumerate}
\end{met}
\begin{thm}
	"Composer les applications linéaires revient à multiplier les matrices." \par Formellement, soit $E$, $F$ et $G$ des $\K$-ev de dimension finie, $u\in\mathcal{L}(E,F)$ et $v\in\mathcal{L}(F,G)$, $\mathcal{B}_E$, $\mathcal{B}_F$ et $\mathcal{B}_G$ des bases respectives de $E$, $F$ et $G$. Alors on a : $$\mat_{\mathcal{B}_E,\mathcal{B}_G}(v\circ u)=\mat_{\mathcal{B}_F,\mathcal{B}_G}(v)\times\mat_{\mathcal{B}_E,\mathcal{B}_F}(u)$$
\end{thm}
\begin{cor}
	Soit $(e_1,\ldots,e_n)$ une base de $E$. L'application : $$\begin{array}{ccccc}\mat_{(e_i)}&:&\mathcal{L}(E)&\rightarrow&\mathcal{M}_n(\K)\\&&u&\mapsto&\mat_{(e_i)}(u)\end{array}$$ est un isomorphisme de $\K$-algèbres.
\end{cor}
\begin{thm}
	Une matrice $A\in\mathcal{M}_n(\K)$ est inversible si, et seulement si, son ALCA est bijective.
\end{thm}
\begin{thm}[Interprétation géométrique des matrices par blocs]
	Soit $u\in\mathcal{L}(E)$ et $E_1$ et $E_2$ des sev supplémentaires dans $E$. Notons $(e_1,\ldots,e_k)$ une base de $E_1$ et $(e_{k+1},\ldots,e_n)$ une base de $E_2$. Alors $(e_1,\ldots,e_n)$ est une base adaptée à la décomposition $E=E_1\oplus E_2$. Écrivons $\mat_{(e_i)}(u)$ sous la forme : $$\mat_{(e_i)}(u)=\begin{pmatrix}A&\vline&B\\ \hline C&\vline&D\end{pmatrix}$$ Alors $B=0_{k,n-k}$ si et seulement si, $E_2$ est stable par $u$, et $C=0_{n-k,k}$ si, et seulement si, $E_1$ est stable par $u$.
\end{thm}
\subsection{Cas des matrices carrées}
\begin{thm}[Lien matrice-base]
	Soit $(e_1,\ldots,e_n)$ une base de $E$ et $(x_1,\ldots,x_n)\in E^n$ une famille quelconque. Posons $A:=\mat_{(e_i)}\left[(x_1,\ldots,x_n)\right]$. Alors : $A$ est inversible si, et seulement si, la famille $(x_1,\ldots,x_n)$ est une base de $E$.
\end{thm}
\begin{exa}[Important]
	Une matrice carrée de taille $n$ est donc inversible si, et seulement si, ses colonnes forment une base de $\K^n$. Cela est aussi équivalent, par égalité entre la taille de la famille et la dimension, au fait que ses colonnes forment une famille libre ou génératrice de $\K^n$.
\end{exa}
\begin{thm}[Lien matrice-isomorphisme]
	Soit $E$ et $F$ de même dimension $n\in\N\st$, $(e_1,\ldots,e_n)$ une base de $E$ et $(f_1,\ldots,f_n)$ une base de $F$. Soit $u\in\mathcal{L}(E,F)$, et posons $A:=\mat_{(e_j),(f_i)}(u)\in\mathcal{M}_n(\K)$. Alors $A$ est inversible si, et seulement si, $u$ est un isomorphisme.
\end{thm}
\begin{rmq}
	Dans ce cas, on a $\mat_{(f_i),(e_j)}(u\inv)=A\inv$.
\end{rmq}
\begin{exa}[Cas des endomorphismes]
	Si $E=F$ et $(e_1,\ldots,e_n)=(f_1,\ldots,f_n)$, on obtient que $u\in GL(E)$ si, et seulement si, $\mat_{(e_i)}(u)\in GL_n(\K)$.
\end{exa}
\begin{thm}[Inversibilité des matrices triangulaires]
	Soit $T$ une matrice triangulaire (supérieure ou inférieure). Alors $T$ est inversible si, et seulement si, tous ses coefficients diagonaux sont non nuls.
\end{thm}
\begin{thm}[Nilpotence des triangulaires strictes]
	Soit $T$ une matrice triangulaire stricte (supérieure stricte ou inférieure stricte) de taille $n\times n$. Alors $T$ est nilpotente, d'indice de nilpotence inférieur ou égal à $n$.
\end{thm}
\begin{exa}[Calcul de puissances de matrices par le binôme de Newton et les matrices nilpotentes]
	Lorsqu'on peut écrire $A=\lambda I_n+T$ avec $T$ triangulaire stricte, on peut utiliser la nilpotence de $T$ et le binôme de Newton pour calculer les puissances de $A$. Cela fonctionne plus généralement si on remplace $T$ par une matrice nilpotente $N$.
\end{exa}
\section{Équivalence et similitude}
\subsection{Cas général}
\begin{dfn}[Matrice de passage]
	Soit $E$ de dimension $n\in\N$, $(e_1,\ldots,e_n)$ une base de $E$, et $(e_1',\ldots,e_n')$ une base de $E$. On appelle \textbf{matrice de passage de la base $(e_i)$ vers la base $(e_i')$} la matrice : $$P:=\mat_{(e_i)}\left[\left(e_1',\ldots,e_n'\right)\right]$$ $P$ est inversible (car $(e_1,\ldots_en)$ est une base) et on a $P=\mat_{(e_i'),(e_i)}(\id_E)$.
\end{dfn}
\begin{rmq}
	$P\inv=\mat_{(e_i),(e_i')}(\id_E)$ est la matrice de passage de la base $(e_i')$ vers la base $(e_i)$.
\end{rmq}
\begin{thm}[Changement de base pour un vecteur]
		Soit $E$ de dimension $n\in\N$, $(e_1,\ldots,e_n)$ une base de $E$, $(e_1',\ldots,e_n')$ une base de $E$ et $P$ la matrice de passage de la base $(e_i)$ vers la base $(e_i')$. Soit $x\in E$ et posons $X:=\mat_{(e_i)}(x)$ et $X':=\mat_{(e_i')}(x)$. Alors on a : $$X=PX'$$
\end{thm}
\begin{thm}[Changement de base pour une application linéaire]
	Soit $E$ de dimension $p\in\N$, $(e_1,\ldots,e_p)$ une base de $E$, $(e_1',\ldots,e_p')$ une base de $E$ et $P$ la matrice de passage de la base $(e_j)$ vers la base $(e_j')$. Soit $F$ de dimension $n\in\N$, $(f_1,\ldots,f_n)$ une base de $F$, $(f_1',\ldots,f_n')$ une base de $F$ et $Q$ matrice de passage de la base $(f_i)$ vers la base $(f_i')$. Soit $u\in\mathcal{L}(E,F)$. Posons $A:=\mat_{(e_j),(f_i)}(u)$ et $A':=\mat_{(e_j'),(f_i')}(u)$. Alors on a : $$A'=Q\inv A P$$
\end{thm}
\begin{dfn}[Matrices équivalentes]
	Soit $A,A'\in\mathcal{M}_{n,p}(\K)$. $A$ et $A'$ sont dites \textbf{matriciellement équivalentes} lorsque : $$\exists(M,N)\in GL_n(\K)\times GL_p(\K),\ A'=MAN$$
\end{dfn}
\begin{thm}
	L'équivalence matricielle est une relation d'équivalence sur $\mathcal{M}_{n,p}(\K)$.
\end{thm}
\begin{thm}["2 matrices sont équivalentes ssi elles représentent une même application linéaire]
	Formellement, soit $E$ de dimension $p\in\N$, $(e_1,\ldots,e_p)$ une base de $E$ et $F$ de dimension $n\in\N$, $(f_1,\ldots,f_n)$ une base de $F$. Soit $u\in\mathcal{L}(E,F)$ et posons $A:=\mat_{(e_j),(f_i)}(u)\in\mathcal{M}_{n,p}(\K)$. \par Soit $A'\in\mathcal{M}_{n,p}(\K)$. Alors $A'$ est équivalente à $A$ si, et seulement si, il existe $(e_1',\ldots,e_p')$ une base de $E$ et $(f_1',\ldots,f_n')$ une base de $F$ telles que $$A'=\mat_{(e_j'),(f_i')}(u)$$
\end{thm}
\subsection{Cas des matrices carrées et des endomorphismes}
\begin{thm}
	Soit $E$ de dimension $p\in\N$, $(e_1,\ldots,e_p)$ une base de $E$, $(e_1',\ldots,e_p')$ une base de $E$ et $P$ la matrice de passage de la base $(e_i)$ vers la base $(e_i')$. Soit $u\in\mathcal{L}(E)$. Posons $A:=\mat_{(e_i)}(u)$ et $A':=\mat_{(e_i')}(u)$. Alors on a : $$A'=P\inv A P$$
\end{thm}
\begin{dfn}[Matrices semblables]
	Deux matrices $A;B\in\mathcal{M}_n(\K)$ sont dites \textbf{semblables} lorsque : $$\exists P\in GL_n(\K),\ B=P\inv AP$$
\end{dfn}
\begin{thm}["2 matrices sont semblables ssi elles représentent un même endomorphisme]
	Formellement, soit $E$ de dimension $n\in\N\st$ et $(e_1,\ldots,e_n)$ une base de $E$. Soit $u\in\mathcal{L}(E)$, et posons $A:=\mat_{(e_i)}(u)\in\mathcal{M}_n(\K)$. Soit $A'\in\mathcal{M}_n(\K)$. Alors $A'$ est semblable à $A$ si, et seulement si, il existe une base $(e_1',\ldots,e_n')$ de $E$ telle que : $$A'=\mat_{(e_i')}(u)$$
\end{thm}
\begin{exa}
	Soit $\lambda\in\K$. La seule matrice semblable à la matrice $\lambda I_n$ est $\lambda I_n$ (il s'agit d'une analyse-synthèse rapide qui utilise le fait que $\lambda I_n$ commute avec toute matrice). Il suffit donc de prendre des $\lambda I_n$ et $\mu I_n$ avec $\lambda\ne\mu$ pour trouver deux matrices non semblables.
\end{exa}
\begin{prop}[Semblables $\implies$ équivalentes]
	Si deux matrices sont semblables, alors elles sont équivalentes.
\end{prop}
\begin{rmq}
	Attention : la réciproque est fausse ! Déjà, ne serait-ce que pour une question de taille, mais même lorsque les tailles sont compatibles, on a des contre-exemples. Par exemple, $2I_n$ et $I_n$ sont équivalentes car $2I_n=(\sqrt{2}I_n)I_n(\sqrt{2}I_n)$, mais ne sont pas semblables d'après l'exemple précédent.
\end{rmq}
\begin{prop}[Puissances naturelles de matrices semblables]
	Soit $A,B\in\mathcal{M}_n(\K)$ semblables. Fixons $P\in GL_n(\K)$ telle que $B=P\inv AP$. Par une récurrence immédiate, on a : $$\forall k\in\N,\ B^k=P\inv A^kP$$
\end{prop}
\begin{rmq}[Diagonalisation, HP, spé]
	Ce résultat est à la base de ce que l'on appelle la \textbf{diagonalisation} : on cherche si $A$ peut être une matrice diagonale. En effet, on sait alors calculer très aisément les puissances de $A$, et donc celles de $B$. La diagonalisation sera développée dans le cours de spé.
\end{rmq}
\begin{exe}
	Montrer que toute matrice nilpotente est semblable à une matrice triangulaire stricte (supérieure stricte par exemple). On pourra procéder par récurrence sur la taille de la matrice, et on pourra utiliser l'ALCA dans l'hérédité pour se ramener à une matrice dont la première colonne est nulle en utilisant l'hypothèse de récurrence.
\end{exe}
\begin{exe}
	Notons $\mathcal{N}_n$ l'ensemble des matrices carrées de taille $n\in\N\st$ nilpotentes. Montre que la famille $(A)_{A\in\mathcal{N}_n}$ n'est pas génératrice de $\mathcal{M}_n(\K)$.
\end{exe}
\begin{prop}[Invariance de la trace par similitude]
	Deux matrices semblables ont même trace.
\end{prop}
\begin{rmq}[Important]
	Par contraposée, il suffit de montrer que deux matrices n'ont pas même trace pour montrer qu'elles ne sont pas semblables.
\end{rmq}
\begin{dfn}[Trace d'un endomorphisme en dimension finie]
	Soit $E$ de dimension finie $n\in\N\st$ et $u\in\mathcal{L}(E)$. Le scalaire $\tr(\mat_{(e_i)}(u))$ est indépendant de la base $(e_i)$ de $E$ choisie. Soit $(e_1,\ldots,e_n)$ une base quelconque de $E$. La \textbf{trace de $u$} est définie par : $$\tr(u):=\tr(\mat_{(e_i)}(u))$$
\end{dfn}
\begin{prop}
	Soit $E$ un $\K$-ev de dimension finie. Alors $\tr(\cdot):\mathcal{L}(E)\rightarrow\K$ est une forme linéaire.
\end{prop}
\begin{prop}
	On a : $$\forall(u,v)\in\mathcal{L}(E),\ \tr(v\circ u)=\tr(u\circ v)$$
\end{prop}
\begin{prop}
	Soit $E$ de dimension finie et $p\in\mathcal{L}(E)$ un projecteur. Alors on a : $$\tr(p)=rg(p)$$
\end{prop}
\begin{rmq}[Utile pour exos X/ENS]
	Pour montrer que la trace d'une application linéaire est dans $\N$, on peut se demander si cette application linéaire n'est pas une combinaison linéaire de projecteurs ou une composée de projecteurs. Réciproquement, si on travaille sur le rang d'une somme de projecteurs, il peut être utile de passer par la trace, notamment si on a des égalités du type $rg(p+q)=rg(p)+rg(q)$.
\end{rmq}
\begin{prop}[Trace et rang d'une symétrie, HP]
	Soit $E$ de dimension finie $n\in\N\st$ et $s\in\mathcal{L}(E)$ une symétrie de rang $r$. On a : $$\tr(s)=2rg(s)-n$$
\end{prop}
\begin{proof}
	Soit $(e_1,\ldots,e_r)$ une base de $\ker(s-\id_E)$ et $(e_{r+1},\ldots,e_n)$ une base de $\ker(s+\id_E)$. Alors $(e_1,\ldots,e_n)$ est une base adaptée à $E=\ker(s-\id_E)\oplus\ker(s+\id_E)$, et on a dans cette base : $$\mat_{(e_i)}(s)=\begin{pmatrix}I_r&\vline&(0)\\ \hline (0)&\vline&-I_{n-r}\end{pmatrix}$$ Si bien que $$\tr(s)=2r-n=2rg(s)-n$$
\end{proof}
\begin{exa}
	Toute matrice $A\in\mathcal{M}_n(\K)$ telle que $A^2=A$ est semblable à une matrice par blocs de la forme : $$\begin{pmatrix}I_r&\vline&(0)\\ \hline (0)&\vline&(0)\end{pmatrix}$$ avec $r\in\llbracket0,n\rrbracket$. En effet, son ALCA est alors un projecteur, et on utilise le théorème "2 matrices sont équivalentes ssi elles représentent un même endomorphisme".
\end{exa}
\section{Rang d'une matrice}
\subsection{Définition et premières propriétés}
\begin{dfn}[Rang d'une matrice]
	Soit $A\in\mathcal{M}_{n,p}(\K)$. Le \textbf{rang de la matrice $A$}, noté $rg(A)$, est défini des deux façons équivalentes qui suivent : \begin{enumerate}
		\item Le rang de $A$ est défini comme étant égal au rang de la famille de ses colonnes interprétées comme des vecteurs de $\K^n$.
		\item Le rang de $A$ est défini comme étant égal au rang de son ALCA.
	\end{enumerate}
\end{dfn}
\begin{exa}
	Soit $(i,j)\in\llbracket1,n\rrbracket\times\llbracket1,p\rrbracket$. On a $rg(E_{i,j})=1$.
\end{exa}
\begin{prop}[Minoration du rang par les tailles]
	On a : $$rg(A)\leqslant\min(n,p)$$
\end{prop}
\begin{thm}[CNS d'inversibilité sur le rang, rang plein]
	Soit $A\in\mathcal{M}_n(\K)$. On a : $$A\in GL_n(\K)\iff rg(A)=n$$ On dit alors que $A$ est de \textbf{rang plein}.
\end{thm}
\begin{thm}["Le rang d'une famille de vecteurs est égal au rang de sa matrice dans n'importe quelle base"]
	Soit $E$ de dimension de $n\in\N\st$, $(e_1,\ldots,e_n)$ une base de $E$ et $(x,1,\ldots,x_p)\in E^p$. Posons $A:=\mat_{(e_i)}\left[\left(x_1,\ldots,x_p\right)\right]\in\mathcal{M}_{n,p}(\K)$. Alors on a : $$rg(A)=rg(x_j)_{1\leqslant j\leqslant p}$$
\end{thm}
\begin{thm}["Le rang d'une application linéaire est égal au rang de sa matrice dans n'importe quel couple de bases"]
	Soit $E$ de dimension de $p\in\N\st$, $(e_1,\ldots,e_p)$ une base de $E$ et $F$ de dimension $n\in\N\st$, $(f_1,\ldots,f_n)$ une base de $F$. Soit $u\in\mathcal{L}(E,F)$. Posons $A:=\mat_{(e_j),(f_i)}(u)$. Alors on a : $$rg(A)=rg(u)$$
\end{thm}
Dans la séquence qui suit, on traduit matriciellement des résultats du chapitre "Dimension finie".
\begin{thm}
	Soit $A,B\in\mathcal{M}_n(\K)$. Si $AB=I_n$, alors $BA=I_n$ puis $A\in GL_n(\K)$ et $A\inv=B$.
\end{thm}
\begin{rmq}
	En particulier, il suffit de vérifier l'inverse d'un seul côté.
\end{rmq}
\begin{thm}
	Soit $(A,B)\in\mathcal{M}_{n,p}(\K)\times\mathcal{M}_{p,q}(\K)$. Alors on a : $$rg(AB)\leqslant\min(rg(A),rg(B))$$
\end{thm}
\begin{thm}[Invariance du rang par multiplication par une matrice inversible]
	Soit $A\in\mathcal{M}_{n,p}(\K)$. On a : \begin{itemize}
		\item $\forall P\in GL_p(\K),\ rg(AP)=rg(A)$
		\item $\forall Q\in GL_n(\K),\ rg(QA)=rg(A)$
	\end{itemize}
\end{thm}
\begin{exa}
	Deux matrices équivalentes (et aussi deux matrices semblables), ont donc même rang.
\end{exa}
\begin{exa}
	On a : $\forall \lambda\in\K\st,\ rg(\lambda A)=rg(A)$.
\end{exa}
\subsection{Rang et équivalence}
\begin{dfn}[Matrice $J_r$]
	Soit $n,p\in\N\st$ et $r\in\llbracket0,\min(n,p)\rrbracket$. On note $(e_1,\ldots,e_n)$ la base canonique de $\K^n$. On pose : $$J_r^{(n,p)}:=\mat_{(e_i)}\left[\left(e_1,\ldots,e_r,0,\ldots,0\right)\right]\in\mathcal{M}_{n,p}(\K)$$ Autrement dit, $$J_r^{(n,p)}:=\begin{pmatrix}I_r&\vline&0_{r,p-r}\\ \hline 0_{n-r,r}&\vline&0_{n-r,p-r}\end{pmatrix}$$ Si le contexte sur les tailles est clair, on pourra la noter $J_r$.
\end{dfn}
\begin{exa}
	On a : \begin{itemize}
		\item $J_1^{(n,p)}=E_{1,1}^{(n,p)}$
		\item $J_n^{(n,n)}=I_n$
		\item $J_0^{(n,p)}=0_{n,p}$
	\end{itemize}
\end{exa}
\begin{rmq}
	On a toujours $rg(J_r)=r$.
\end{rmq}
\begin{lem}
	Soit $E$ et $F$ de dimensions respectives $p\in\N\st$ et $n\in\N\st$. Soit $u\in\mathcal{L}(E,F)$ et posons $r:=rg(u)$. Alors, il existe $(e_1,\ldots,e_p)$ une base de $E$ et $(f_1,\ldots,f_n)$ une base de $F$ telles que : $$\mat_{(e_j),(f_i)}(u)=J_r^{(n,p)}$$
\end{lem}
\begin{thm}
	Soit $A\in\mathcal{M}_{n,p}(\K)$ et $r\in\llbracket0,\min(n,p)\rrbracket$. Alors $rg(A)=r$ si, et seulement si, $A$ est \textbf{matriciellement équivalente} à $J_r^{(n,p)}$.
\end{thm}
\begin{rmq}
	Attention : ce théorème utilise bien la notion d'équivalence matricielle et non pas de similitude ! Ne pas dire cela avec la similitude, cela est faux. Il y a beaucoup plus de classes d'équivalences pour la relation de similitude, c'est l'objet de la réduction de Jordan (spé).
\end{rmq}
\begin{cor}
	Soit $A,B\in\mathcal{M}_{n,p}(\K)$. Alors $A$ et $B$ sont équivalentes si, et seulement si, elles ont même rang.
\end{cor}
\begin{exa}
	Il y a donc exactement $1+\min(n,p)$ classes d'équivalence pour la relation d'équivalence matricielle sur $\mathcal{M}_{n,p}(\K)$, dont un système de représentant est fourni par les $J_r^{(n,p)}$ pour $r\in\llbracket0,\min(n,p)\rrbracket$.
\end{exa}
\begin{cor}
	On a : $$\forall A\in\mathcal{M}_{n,p}(\K),\ rg(A^\top)=rg(A)$$
\end{cor}
\begin{cor}
	Le rang d'une matrice est donc aussi égal au rang de la famille constituée de ses lignes vues comme des éléments de $\mathcal{M}_{1,p}(\K)$.
\end{cor}
\begin{exa}
	Soit $A\in\mathcal{M}_n(\K)$. Alors $A\in Gl_n(\K)$ si, et seulement si, ses lignes forment une base de $\mathcal{M}_{1,n}(\K)$ (ce qui est aussi équivalent au fait qu'elles en forment une famille libre ou génératrice).
\end{exa}
\subsection{Matrices extraites}
\begin{dfn}[Matrice extraite/sous-matrice]
	Soit $A\in\mathcal{M}_{n,p}(\K)$. Une \textbf{matrice extraite}, ou \textbf{sous-matrice}, de $A$ est une matrice obtenue en sélectionnant certaines lignes et certaines colonnes de $A$ (pas nécessairement adjacentes).
\end{dfn}
\begin{thm}
	Soit $A\in\mathcal{M}_{n,p}(\K)$ et $r\in\llbracket0,\min(n,p)\rrbracket$. Les deux assertions suivantes sont équivalentes : \begin{enumerate}
		\item $A$ admet une sous-matrice $(r,r)$ inversible
		\item $rg(A)\geqslant r$
	\end{enumerate}
	En particulier, $rg(A)$ est le plus grand entier $r$ tel que l'on puisse trouver une matrice extraite de $A$ de taille $(r,r)$ inversible.
\end{thm}
\begin{cor}
	Soit $A\in\mathcal{M}_{n,p}(\K)$ et $B$ une sous-matrice de $A$. Alors on a : $$rg(B)\leqslant rg(A)$$
\end{cor}
\subsection{Opérations élémentaires}
\begin{dfn}[Opérations élémentaires]
	On a trois types d'\textbf{opérations élémentaires sur les lignes} (abrégées en \textbf{OEL}), et de même symétriquement sur les colonnes : 
	\begin{itemize}
		\item Transvections : $L_i\leftarrow L_i+\lambda L_j$ avec $\lambda\in\K\st$ et $i\ne j$
		\item Échanges de lignes : $L_i\leftrightarrow L_j$
		\item Dilatations : $L_i\leftarrow \lambda L_i$ avec $\lambda\in\K\st$
	\end{itemize}
\end{dfn}
\begin{lem}[Effet de la multiplication par une $E_{i,j}$]
	Soit $(i,j)\in\llbracket1,n\rrbracket^2$. Dans le produit $E_{i,j}^{(n,n)}\times A$ : 
	\begin{itemize}
		\item $L_i\leftarrow L_j$
		\item Toutes les autres lignes sont mises à $0$
	\end{itemize}
	Dans le produit $A\times E_{i,j}^{(n,n)}$ : 
	\begin{itemize}
		\item $C_j\leftarrow C_i$
		\item Toutes les autres colonnes sont mises à $0$
	\end{itemize}
\end{lem}
\begin{thm}[Traduction matricielle des OE]
	Une OE sur les lignes revient à multiplier à gauche par une matrice inversible. Une OE sur les colonnes revient à multiplier à droite par une matrice inversible. Chacune des 6 OE laisse le rang invariant.
\end{thm}
\begin{met}[Détermination du rang d'une matrice par opérations élémentaires]
	Voici comment calculer le rang de $A\in\mathcal{M}_{n,p}(\K)$.
	\begin{itemize}
		\item Si $A=0_{n,p}$, alors le rang est nul.
		\item Sinon, on peut trouver un terme $a_{i,j}\ne0_\K$.
		\item Si $A$ est une matrice-ligne ou une matrice-colonne, on sait alors que son rang vaut $1$.
		\item Sinon, on ramène $a_{i,j}$ en position $(1,1)$ par une permutation de lignes et / ou une permutation de colonnes.
		\item Ce terme nous sert alors de \textbf{pivot}. En effectuant les transvections $L_i\leftarrow L_i-\displaystyle\frac{a_{i,1}}{a_{1,1}}L_1$ (pour $2\leqslant i\leqslant n$), on annule le reste de la première colonne.
		\item Par le lemme sur le calcul du rang, on se ramène alors à une matrice de taille $(n-1,p-1)$ et on itère jusqu'à arriver à une matrice nulle ou une matrice-ligne ou une matrice-colonne, dont on connaît toujours le rang.
	\end{itemize}
\end{met}
\begin{thm}[Matrices d'opérations élémentaires]
	Voici les matrices qui correspondent aux OE : \begin{itemize}
		\item Transvection $L_i\leftarrow L_i+\lambda L_j$ ou $C_j\leftarrow C_j+\lambda C_i$ avec $i\ne j$ et $\lambda\in\K\st$ : $$T_{i,j}(\lambda):=I_n+\lambda E_{i,j}$$
		\item Échange $L_i\leftrightarrow L_j$ ou $C_i\leftarrow C_j$ : $$Ech_{i,j}:=I_n-E_{i,i}-E_{j,j}+E_{i,j}+E_{j,i}$$
		\item Dilatation $L_i\leftarrow \lambda L_i$ ou $C_i\leftarrow\lambda C_i$ avec $\lambda\in\K\star$ : $$D_i(\lambda):=I_n+(\lambda-1)E_{i,i}$$
	\end{itemize}
\end{thm}
\begin{rmq}
	Toute matrice inversible reste donc inversible après multiplication par une matrice d'OE.
\end{rmq}
\begin{lem}
	Soit $n,p\geqslant2$ et $A\in\mathcal{M}_{n,p}(\K)$ de la forme : $$A=\begin{pmatrix}\lambda&\vline&(*)\\ \hline (0)&\vline&B\end{pmatrix}$$
	ou $$A=\begin{pmatrix}\lambda&\vline&(0)\\ \hline (*)&\vline&B\end{pmatrix}$$ avec $\lambda\in\K$ et $B\in\mathcal{M}_{n-1,p-1}(\K)$. Si $\lambda\ne0_\K$, alors $$rg(A)=1+rg(B)$$
\end{lem}
\begin{prop}
	Soit $A\in GL_n(\K)$. Par OEL ($ie$ en multipliant $A$ par des matrices d'OEL), on peut se ramener à une matrice triangulaire supérieure.
\end{prop}
\begin{met}[Algorithme de Gauss-Jordan]
	Voici comment inverser une matrice, ou du moins apprendre qu'elle n'est pas inversible.
	\begin{itemize}
		\item On trace deux colonnes : en haut de la colonne gauche, on écrit $A$, en haut de celle de droite, $I_n$.
		\item A partir de là, toutes les fois que l'on fera une opération élémentaire sur les lignes de $A$, on effectuera la même opération sur celles de la matrice de la colonne de droite.
		\item On s'intéresse d'abord à la première colonne de $A$.
		\begin{itemize}
			\item Si tous les $a_{i,1}$ sont nuls, alors $A$ n'est pas inversible et on s'arrête là.
			\item Sinon, on choisit un pivot $a_{i,1}\ne0_\K$ puis on permute $L_1$ et $L_i$.
		\end{itemize}
		\item Par transvections, on annule le reste de la première colonne. 
		\item On recommence en travaillant sur la matrice extraite $$\begin{pmatrix}a_{2,2}&\ldots&a_{2,n}\\ \vdots&&\vdots\\ a_{n,2}&\ldots&a_{n,n}\end{pmatrix}$$
		\begin{itemize}
			\item[--] Si à un moment on ne peut pas trouver de pivot non nul, alors la matrice n'est pas inversible et on s'arrête là.
			\item[--] Sinon, on arrive à la fin à une matrice triangulaire supérieure dont la diagonale ne s'annule pas.
		\end{itemize}
		\item Puis on "repart dans l'autre sens". Comme $a_{n,n}\ne0_\K$, on l'utilise comme pivot pour annuler tous les $a_{i,n}$ (pour $1\leqslant i\leqslant n-1$) par transvection.
		\item On recommence alors sur la matrice extraite $$\begin{pmatrix}
			a_{1,1}&\ldots&a_{1,n-1}\\ &\ddots&\vdots\\0&&a_{n-1,n-1}
		\end{pmatrix}$$et ainsi de suite. A chaque fois, on tombe directement sur un pivot non nul (c'est le terme diagonal), il n'y a plus besoin de permuter les lignes.
		\item A la fin, on obtient une matrice diagonale, dont la diagonale ne s'annule pas.
		\item On dilate chaque ligne pour ramener les termes diagonaux à $1$ (on peut aussi effectuer les dilatations avant les transvections si cela simplifie les calculs).
		\item Alors, la matrice écrite dans la colonne de droite est égale à $A\inv$.
	\end{itemize}
\end{met}
\begin{prop}[Les OEL engendrent $GL_n(\K)$, HP, exo classique]
	Les matrices d'OEL engendrent $GL_n(\K)$. Cela fonctionne encore sans les échanges de ligne. 
\end{prop}
\begin{proof}
	En effet, soit $A\in GL_n(\K)$. Puisque $A\inv\in GL_n(\K)$, on peut passer de $A\inv$ à $I_n$ par multiplication à gauche par un nombre fini de matrices d'OEL : $M_N\times\ldots\times M_1\times A\inv=I_n$. Donc $A=M_N\times\ldots\times M_1$ puis les matrices d'OEL engendrent $GL_n(\K)$.Ensuite, on peut émuler l'échange de lignes $L_i\leftrightarrow L_j$ à partir des autres OEL en effectuant successivement : $L_i\leftarrow L_i+L_j$, $L_j'\leftarrow L_j'-L_i'$, $L_i''\leftarrow L_i''+L_j''$ et enfin $L_j'''\leftarrow -L_j'''$.
\end{proof}

\section{Systèmes linéaires}
\subsection{Premières définitions}
\begin{dfn}[Système linéaire, inconnues]
	Un \textbf{système linéaire de $n$ équations à $p$ inconnues} (dans $\K$) est un système de la forme : $$E :\left \{ \begin{array}{ccccccc}
		a_{1,1}x_1&+&\ldots&+&a_{1,p}x_p&=&b_1 \\
		\vdots&&&&\vdots&\vdots&\vdots\\
		a_{n,1}x_1&+&\ldots&+&a_{n,p}x_p&=&b_n
	\end{array}\right.$$ Les $a_{i,j}$ et les $b_i$ sont des scalaires fixés, ce sont les données du problème. Formellement, l'\textbf{inconnue} est le $p$-uplet $\displaystyle\begin{pmatrix}x_1\\ \vdots\\ x_p\end{pmatrix}\in\K^p$, mais on a l'habitude de dire que les $x_j$ sont les inconnues.
\end{dfn}
\begin{rmq}[Traduction matricielle]
	Posons $$A:=\begin{pmatrix}a_{1,1}&\ldots&a_{1,p}\\ \vdots&&\cdots&\\a_{n,1}&\ldots&a_{n,p}\end{pmatrix} \in\mathcal{M}_{n,p}(\K)$$ et $B:=\displaystyle\begin{pmatrix}b_1\\ \vdots\\ b_n\end{pmatrix}\in\K^n$. Soit $X:=\displaystyle\begin{pmatrix}x_1\\ \vdots\\ x_p\end{pmatrix}\in\K^p$. Alors $X$ est solution du système si, et seulement si, $AX=B$.
\end{rmq}
\begin{dfn}[Rang d'un système linéaire]
	Le \textbf{rang du système} est défini par $rg((E)):=rg(A)$.
\end{dfn}
\begin{dfn}[Système homogène]
	Le \textbf{système homogène} $(E_0)$ associé au système $(E)$ est le système : $$E :\left \{ \begin{array}{ccccccc}
		a_{1,1}x_1&+&\ldots&+&a_{1,p}x_p&=&0_\K \\
		\vdots&&&&\vdots&\vdots&\vdots\\
		a_{n,1}x_1&+&\ldots&+&a_{n,p}x_p&=&0_\K
	\end{array}\right.$$ Matriciellement, il est équivalent à $AX=0_{\K^n}$.
\end{dfn}
\begin{thm}[Structure de l'ensemble des solutions d'un système linéaire]
	Notons $u$ l'ALCA à A, $\mathcal{S}$ l'ensemble des solutions de $(E)$ et $\mathcal{S}_0$ l'ensemble des solutions de $(E_0)$ (donc $\mathcal{S}_0=\ker(u$)). On rappelle que \begin{itemize}
		\item Soit $\mathcal{S}=\emptyset$.
		\item Soit $\mathcal{S}$ est un sea dirigé par $\mathcal{S}_0$, où $\mathcal{S}_0$ est de dimension $p-rg((E))$ par le théorème du rang.
	\end{itemize}
	En quelque sorte, $p-rg((E))$ mesure le nombre de \textbf{degrés de liberté} du système $(E)$.
\end{thm}
\subsection{Différentes interprétations de $A$}
Une première interprétation serait de savoir si $B\in\im(A)$. \\ \par Une deuxième interprétation peut se faire en écrivant $A$ sous la forme $$\begin{pmatrix}&\vline&&\vline&\\C_1&\vline&\ldots&\vline&C_p\\&\vline&&\vline&\end{pmatrix}$$ Alors, le système est équivalent à $\displaystyle\sum_{j=1}^{p}x_j\cdot C_j=B$ (produit par blocs). On cherche donc à savoir si $B$ peut s'écrire comme combinaison linéaire des $C_j$, et les $x_j$ sont précisément les coefficients de cette combinaison linéaire. \\ \par Une troisième façon de voir les choses est d'écrire $A$ sous la forme $$\begin{pmatrix}&L_1&\\ \hline &\vdots& \\ \hline &L_n&\end{pmatrix}$$ Alors, le système est équivalent à $$\begin{cases}
	L_1X=&b_1\\
	\vdots&\vdots\\
	L_nX=&b_n
\end{cases}$$ Soit $i\in\llbracket1,n\rrbracket$ et notons $\varphi_i$ l'ALCA à $L_i$Le système homogène se réécrit : $$\begin{cases}
\varphi_1(X)=&b_1\\
\vdots&\vdots\\
\varphi_n(X)=&b_n
\end{cases}$$ $ie$ $X\in\displaystyle\bigcap_{i=1}^n\ker(\varphi_i)$. A nouveau, soit $i\in\llbracket1,n\rrbracket$ : \begin{itemize}
	\item Si $L_i=(0,\ldots,0)$, soit $b_i\ne0_\K$, et alors on a une \textbf{incompatibilité} et $\mathcal{S}=\emptyset$ ; soit $b_i=0_\K$ et il est équivalent de supprimer la ligne.
	\item Si $L_i\ne(0,\ldots,0)$, alors $\varphi_i$ est non nulle, et comme c'est une forme linéaire, elle est surjective. Puis l'équation $L_iX=b_i$ admet au moins une solution puis l'ensemble des solutions de cette équation est un hyperplan affine dirigé par $\ker(\varphi_i)$.
\end{itemize}
En somme, si $\mathcal{S}\ne\emptyset$, il s'agit d'une intersection d'hyperplan affine.
\subsection{Résolution}
Dans toute cette partie, posons $r:=rg(A)$.
\begin{dfn}[Système de Cramer]
	Le système $(E)$ est dit \textbf{de Cramer} lorsque $n=p=r$ (dans ce cas, $A\in GL_n(\K)$). On a alors $AX=B\iff X=A\inv B$ donc $\mathcal{S}=\left\{A\inv B\right\}$.
\end{dfn}
\begin{met}[Résolution d'un système linéaire par le pivot de Gauss]
	\begin{itemize}
		\item On travaille sur les lignes du système par opérations élémentaires.
		\item On part du coin en haut à gauche de $A$ et tant que le scalaire $a_{j,j}$ sur la diagonale est non nul (c'est le \textbf{pivot}), on effectue des transvections $L_i\leftarrow L_i-\displaystyle\frac{a_{i,j}}{a_{i,i}}L_j$ pour $i>j$ afin d'annuler la $j$-ème colonne en-dessous de $a_{j,j}$.
		\item Si à l'étape $j$ on ne peut pas trouver de pivot non nul, c'est qu'on a : $$A=\begin{pmatrix}
			q_1&&&&&\ldots&\ldots&\ldots \\
			0&\ddots&&&&\ldots&\ldots&\ldots\\
			&\ddots&\ddots&&&\ldots&\ldots&\ldots\\
			&&\ddots&q_{j-1}&&\ldots&\ldots&\ldots\\
			0&\ldots&\ldots&0&0&\ldots&\ldots&\ldots\\
			\vdots&&&\vdots&\vdots&\ldots&\ldots&\ldots\\
			0&\ldots&\ldots&0&0&\ldots&\ldots&\ldots\\
		\end{pmatrix}$$ où les $q_1,\ldots,q_{j-1}$ sont des pivots non nuls. On essaye alors de permuter la $j$-ème colonne avec une autre colonne à sa droite (il suffit d'échanger deux inconnues), à la recherche d'un pivot non nul quelque par dans la sous matrice en bas à droite, qui commence juste une case en-dessous et à droite de $q_{j-1}$. Puis on continue.
		\item Si à un moment ce n'est pas possible, alors c'est que pour toutes les lignes pour $j\leqslant i\leqslant n$ dans le membre de gauche de l'équation du système sont nulles. Notamment, on arrivera, dans cette situation à chaque fois que $n>p$, c'est-à-dire s'il y a plus d'équations que d'inconnues. On est alors face à ce qu'on appelle un \textbf{risque d'incompatibilité} :
		\begin{itemize}
			\item[--] Si le second membre correspondant n'est pas nul, alors on a une contradiction, et il n'y a pas de solution.
			\item[--] Sinon, on est face à des lignes du type "$0=0"$, et il est équivalent de les supprimer.
		\end{itemize}
		\item On arrive finalement à une situation du type : $$\begin{pmatrix}
			q_1&\ldots&\ldots&\ldots&\ldots&\ldots\\&\ddots&&&&\\&&\ddots&&&\\ (0)&&&q_r&\ldots&\ldots
		\end{pmatrix}$$ avec $q_1,\ldots,q_r\ne0_\K$ et $r$ le rang du système initial.
		\item Si $p=r$, on a un système de Cramer \textbf{échelonné} qu'on résout de bas en haut (on peut trouver la dernière inconnue, puis on réinjecte en remontant pour obtenir toutes les inconnues une par une). Si $p>r$, les $r$ première variables sont alors appelées \textbf{variables principales}, et les $p-r$ restantes \textbf{variables auxiliaires}. On fait passer à droites les variables auxiliaires, qui deviennent des paramètres, et on se ramène au cas précédent.
	\end{itemize}
\end{met}
\subsection{Récapitulatif}
Récapitulons les fois où nous nous sommes servis d'un pivot de Gauss ou d'une de ses variantes. \begin{itemize}
	\item Trouver une relation de liaison entre des vecteurs de $\K^n$, ou montrer qu'une famille de vecteurs de $K^n$ est libre (ce qui est essentiellement la même chose)
	\item Inverser $A\in GL_n(\K)$
	\item Calculer le rang de $A\in\mathcal{M}_{n,p}(\K)$
	\item Passer de l'écriture d'un sev de $\K^n$ sous forme d'un système d'équations à une écriture sous la forme d'un $\vect$, et vice-versa
	\item Extraire une base du $\vect$ d'une famille de vecteurs de $\K^n$. Profitons-en pour rappeler la méthode, qui n'est pas décrite dans ce document jusqu'à présent : 
\end{itemize}
\begin{met}[Extraire d'une famille de vecteurs une base de son $\vect$]
	Soit $(C_1,\ldots,C_p)$ une famille de vecteurs de $\K^n$. Posons $F:=\vect(C_j)_{1\leqslant j\leqslant p}$ et cherchons une base de $F$ sous la forme $(C_{j_1},C_{j_2},\ldots)$ (qui existe d'après le théorème de la base extraite). Notons $r:=rg(C_j)_{1\leqslant j\leqslant p}$ ($ie$ $r:=\dim(F)$).
	\begin{itemize}
		\item On résout le système $\displaystyle\sum_{j=1}^{p}x_jC_j=0$ $ie$ $\displaystyle A\times\begin{pmatrix}x_1\\ \vdots \\ x_p\end{pmatrix}=0$, avec $$A=\begin{pmatrix}&\vline&&\vline&\\C_1&\vline&\ldots&\vline&C_p\\&\vline&&\vline&\end{pmatrix}$$
		\item On repère les "indices principaux".
	\end{itemize}
	On rappelle maintenant la justification de cette méthode, à reproduire au cas par cas une copie :
	\begin{itemize}
		\item SPG, supposons que les variables principales sont $x_1,\ldots,x_r$ (quitte à les renommer). On a alors $rg(A)=r$. En effet, $rg(A)$ est alors égal à $r$ en appliquant l'algorithme de détermination du rang. \textbf{Note} : dans un exercice calculatoire, ce point est inutile.
		\item Montrons que $(C_1,\ldots,C_r)$ est une base de $F$.
		\begin{itemize}
			\item[--] Liberté : soit $x_1,\ldots,x_r\in\K$ tels que $x_1C_1+\ldots+x_rC_r=0_{\K^n}$. On pose artificiellement $x_{r+1}=\ldots=x_p=0_\K$. Le vecteur $\displaystyle\begin{pmatrix}x_1\\ \vdots\\ x_p\end{pmatrix}$ est alors solution du système homogène résolut précédemment. Mais puisque $x_1,\ldots,x_r$ s'expriment de façon linéaire en fonction des $x_{r+1},\ldots,x_p$, on a $x_1=\ldots=x_p=0_\K$. La famille est libre.
			\item[--] Caractère générateur : jusqu'à présent, on montrait que $\dim(F)=r$, ce qui suffisait à conclure. Pour ce faire, on montrait que $\forall j\in\llbracket r+1,p\rrbracket,\ C_j\in \vect(C_l)_{1\leqslant l\leqslant r}$ en posant $x_j=-1$ et $0$ pour les autres et on réinjectait dans le système. Mais désormais, cela est inutile en vertu de l'algorithme de calcul du rang qui nous assure immédiatement cette égalité puisque nous avons résolu le système.
		\end{itemize}
	\end{itemize}
\end{met}

\newpage
\chapter{Déterminants}
\section{Groupe symétrique}
\subsection{Généralités}
\begin{dfn}[Groupe symétrique d'ordre $n$]
	Les permutations de $\llbracket1,n\rrbracket$ forment un groupe pour la composition. On l'appelle le \textbf{groupe symétrique d'ordre $n$} et on le note $\mathcal{S}_n$. Assez rapidement, on omettra le symbole $\circ$.
\end{dfn}
\begin{dfn}[Notation d'une permutation de $\llbracket1,n\rrbracket$]
	Souvent, une permutation $\sigma$ de $\mathcal{S}_n$ sera notée sous la forme d'un tableau à deux lignes : en haut, les entiers de $1$ à $n$ ; en bas, leurs images par $\sigma$. Par exemple, dans $\mathcal{S}_5$, on peut avoir $$\sigma =\begin{pmatrix}1&2&3&4&5\\3&4&1&5&2\end{pmatrix}$$
\end{dfn}
\begin{dfn}[Support, points fixes]
	Soit $\sigma\in\mathcal{S}_n$. L'ensemble des $x\in\llbracket1,n\rrbracket$ tels que $\sigma(x)\ne x$ s'appelle le \textbf{support} de $\sigma$. Les $x$ tels que $\sigma(x)=x$ sont les \textbf{points fixes} de $\sigma$.
\end{dfn}
\begin{rmq} Puisque $\sigma$ est injective, elle stabilise son support. \end{rmq}
\begin{prop}
	Deux permutations à supports disjoints commutent.
\end{prop}
\begin{dfn}[Orbite]
	Soit $\sigma\in\mathcal{S}_n$. La relation binaire $$\forall (x,y)\in\llbracket1,n\rrbracket,\ x\sim y \iff \exists k\in\Z,\ y=\sigma^k(x)$$ est une relation d'équivalence. La classe de $x$ est appelée \textbf{orbite} de $x$. De plus, il existe un plus petit $p>0$ tel que $\sigma^p(x)=x$, et l'orbite de $x$ est alors égale à $\left\{x,\sigma(x),\ldots,\sigma^{p-1}(x)\right\}$, avec les $\sigma^k(x)$ deux à deux distincts pour $0\leqslant k\leqslant p-1$.
\end{dfn}
\begin{exa}
	Si $x$ est un point fixe, son orbite est réduite à $\left\{x\right\}$.
\end{exa}
\begin{dfn}[$p$-cycle]
	Soit $p\geqslant 2$ et $a_1,\ldots,a_p$ des entiers deux à deux distincts dans $\llbracket1,n\rrbracket$. L'application $\sigma$ définie sur $\llbracket1,n\rrbracket$ définie par :
	\begin{enumerate}
		\item $\forall x \notin\left\{a_1,\ldots,a_p\right\},\ \sigma(x)=x$
		\item $\forall i\in\llbracket1,p-1\rrbracket,\ \sigma(a_i)=a_{i+1}$
		\item $\sigma(a_p)=a_1$
	\end{enumerate}
	est une permutation de $\llbracket1,n\rrbracket$ notée $(a_1 \ldots a_p)$. Une telle permutation est appelée \textbf{$p$-cycle} ou \textbf{cycle d'ordre $p$}.
\end{dfn}
\begin{rmq}De manière imagée, on peut imaginer que les $a_i$ sont placés sur le tour d'un cadran. Appliquer la permutation revient alors à "faire tourner le cadran" d'un cran. C'est ce qui explique qu'un cycle s'appelle aussi une \textbf{permutation circulaire}.  \end{rmq} 
\begin{rmq} Le support du cycle $(a_1 \ldots \ a_p)$ est exactement égal à $\left\{a_1,\ldots,a_p\right\}$  \end{rmq} 
\begin{rmq} On a $(a_1\ldots \ a_p)=(a_2\ldots \ a_p a_1)=\ldots=(a_p a_1\ldots \ a_{p-1})$ \end{rmq}
\begin{rmq}  On a $(a_1\ldots \ a_p)\inv=(a_p a_{p-1}\ldots \ a_1)$ \end{rmq}
\begin{thm}[Décomposition en produit de cycles à supports disjoints]
	A l'ordre près des facteurs, toute permutation admet une unique décomposition en produit de cycles à supports disjoints.
\end{thm}
\begin{rmq} Puisque les cycles sont à supports disjoints, ils commutent entre eux. Cela justifie l'expression "à l'ordre des facteurs près". \end{rmq}
\begin{exa}
	Reprenons $\displaystyle \sigma =\begin{pmatrix}1&2&3&4&5\\3&4&1&5&2\end{pmatrix}$. $1$ s'envoie sur $3$, $3$ s'envoie sur $1$. Voilà notre premier cycle, c'est $(1\ 3)$ Puis $2$ s'envoie sur $4$, qui s'envoie sur $5$, qui s'envoie sur $2$. Finalement, $\sigma=(1\ 3)(2\ 4\ 5)$.
\end{exa}
\subsection{Signature}
\begin{dfn}[Transposition]
	Un $2$-cycle s'appelle plus couramment une \textbf{transposition}. La transposition $(a\ b)$ consistant à "changer $a$ et $b$, elle est involutive.
\end{dfn}
\begin{rmq} Pour $n\geqslant 3$, $\mathcal{S}_n$ n'est pas commutatif. Il suffit de prendre $a$, $b$ et $c$ distincts et de considérer $(a\ b)$ et $(b\ c)$ puis de tester les deux compositions qui ne sont alors pas égales. En revanche, pour $n=1$ ou $n=2$, on observe immédiatement que $\mathcal{S}_n$ est bien commutatif. \end{rmq}
\begin{thm}[Décomposition en produit de transpositions]
	Toute permutation s'écrit comme produit de transpositions : $\sigma=t_1\ldots t_k$
\end{thm}
\begin{thm}[Signature, permutations paires et impaires]
	Il existe un unique morphisme de groupes $\varepsilon:\mathcal{S}_n\rightarrow\left\{-1,1\right\}$ qui vaut $-1$ sur les transpositions. Le nombre $\varepsilon(\sigma)$ s'appelle la \textbf{signature} de $\sigma$. Si $\varepsilon(\sigma)=1$ (resp. $\varepsilon(\sigma)=-1$), on dit que la permutation $\sigma$ est \textbf{paire} (resp. \textbf{impaire}).
\end{thm}
\begin{rmq} $\left\{-1,1\right\}$ est bien un groupe en tant que groupe des unités de l'anneau $(\Z,+,\times)$. \end{rmq}
\begin{cor}
	Dans la décomposition de $\sigma\in\mathcal{S}_n$ en produit de transpositions $\sigma=t_1\ldots t_k$, la parité de $k$ est indépendante de la décomposition (en effet, il suffit de passer à la signature).
\end{cor}
\begin{rmq} Un $p$-cycle a une signature qui vaut $(-1)^{p-1}$. En pratique, on pourra l'utiliser pour déterminer la signature d'une permutation après l'avoir décomposée en produit de cycles. \end{rmq}
\begin{dfn}[Groupe alterné d'ordre $n$, HP]
	Les permutations paires forment un sous-groupe de $\mathcal{S}_n$, puisqu'il s'agit de $\ker(\varepsilon)$. Ce sous-groupe est appelé \textbf{groupe alterné d'ordre $n$}, on le note $\mathcal{A}_n$.
\end{dfn}
\begin{prop}[Cardinal de $\mathcal{A}_n$ dès que $n\geqslant2$, HP]
	Supposons que $n\geqslant 2$. Si $\tau$ est une transposition fixée, l'application $\sigma\mapsto\sigma\tau$ envoie $\mathcal{A}_n$ sur $\mathcal{S}_n\setminus\mathcal{A}_n$ comme on le constate par un simple calcul de signature. De plus, elle est bijective car involutive. On en déduit alors que $$\card(\mathcal{A}_n)=\frac{\card(\mathcal{S}_n)}{2}=\frac{n!}{2}$$
\end{prop}
\section{Déterminant}
Dans toute cette partie, $\K$ désigne un corps et $E$ un $\K$-espace vectoriel de dimension $n$.
\subsection{Déterminant d'une famille de vecteurs}
\begin{dfn}[Formes $p$-linéaires]
	Une application $f:E^p\rightarrow\K$ est une forme $p$-linéaire lorsqu'elle est linéaire en chacune de ses variables, ie lorsque l'on fige $p-1$ variables, l'application est linéaire par rapport à la $p$-ème variable.
\end{dfn}
\begin{rmq} Dès que l'un des $x_j$ est nul, on a $f(x_1,\ldots,x_p)=0$ \end{rmq}
 \begin{rmq} Les formes $p$-linéaires forment un sous-espace vectoriel de $\displaystyle \K^{E^p}$. \end{rmq}
\begin{exa}
	Pour $p=1$, on retrouve la définition d'une forme linéaire.
\end{exa}
\begin{exa}
	Dans $\mathcal{M}_n(\K)$, l'application $(A_1,\ldots,A_p)\mapsto\tr(A_1\times\ldots\times A_p)$ est une forme $p$-linéaire.
\end{exa}
\begin{dfn}[HP]
	Si $f_1,\ldots,f_p$ sont des formes linéaires de $E$, alors l'application
	$$
	\begin{array}[t]{lrcl}
		f_1\otimes\ldots\otimes f_p: & E^p & \longrightarrow & \K \\
		& (x_1,\ldots,x_p) & \longmapsto & f_1(x_1)\times\ldots\times f_p(x_p)
	\end{array}
	$$ est une forme $p$-linéaire.
\end{dfn}
\begin{lem}
	Soit $(e_i)_{1\leqslant i\leqslant n}$ une base de $E$, $f$ une forme $p$-linéaire et $x_1,\ldots,x_p$ des vecteurs qu'on écrit sous la forme $$x_j=\sum_{i=1}^{n}\lambda_{i,j}e_i$$ On a alors $$f(x_1,\ldots,x_p)=\sum_{1\leqslant i_1,\ldots,i_p\leqslant n}\lambda_{i_1,1}\ldots\lambda_{i_p,p}f(e_{i_1},\ldots,e_{i_p})$$
\end{lem}
\begin{rmq} On pourrait généraliser le lemme : si $(u_{i,1})_{i\in I_1},\ldots,(u_{i,p})_{i\in I_p}$ sont des familles finies de $E$, alors on a $$f\left(\sum_{i\in I_1}\lambda_{i,1}u_{i,1},\ldots,\sum_{i\in I_1}\lambda_{i,1}u_{i,1}\right)=\sum_{(i_1,\ldots,i_p)\in I_1\times\ldots\times I_p}\lambda_{i_1,1}\ldots\lambda_{i_p,p}f(u_{i_1,1},\ldots,u_{i_p,p})$$ \end{rmq}
A partir de maintenant, on suppose que $\displaystyle\fbox{p=n}$
\begin{dfn}[Forme $n$-linéaire alternée]
	Une forme $n$-linéaire est dite \textbf{alternée} lorsqu'elle s'annule dès que sont égaux : $$\forall i\ne j,\ \left(x_i=x_j\implies f(x_1,\ldots,x_n)=0\right)$$
\end{dfn}
\begin{rmq} Les formes $n$ linéaires alternées forment un sous-espace vectoriel des formes $n$-linéaires. \end{rmq}
\begin{thm}[Invariance par transvection des formes $n$-linéaires alternées]
	Si $f$ est une forme $n$-linéaire alternée, alors le scalaire $f(x_1,\ldots,x_n)$ est invariant si on ajoute à un $x_j$ une combinaison linéaire des autres $x_i$.
\end{thm}
\begin{thm}[Échange de termes, forme antisymétrique]
	Si $f$ est une forme $n$-linéaire alternée, alors elle vérifie $$\forall i\ne j,\ f(x_1,\ldots,x_i,\ldots,x_j,\ldots,x_n)=-f(x_1,\ldots,x_j,\ldots,x_i,\ldots,x_n)$$ où on a échangé les positions de $x_i$ et $x_j$. Une telle forme est dite \textbf{antisymétrique}.
\end{thm}
\begin{rmq}[HP] En caractéristique différente de $2$, la réciproque est vraie : toute forme antisymétrique est alternée. Il suffit d'échanger les deux termes identiques puis de conclure par pseudo-intégrité, car dans ce cas, $2\ne 0$. \end{rmq}
\begin{cor}
	Si $f$ est alternée, alors elle vérifie $$\forall\sigma\in\mathcal{S}_n,\ f(x_{\sigma(1),\ldots,x_{\sigma(n)}})=\varepsilon(\sigma)f(x_1,\ldots,x_n)$$
\end{cor}
\begin{lem}
	Soit $(e_i)_{1\leqslant i\leqslant n}$ une base de $E$, $f$ une forme $n$-linéaire alternée et $x_1,\ldots,x_n$ des vecteurs qu'on écrit sous la forme $$x_j=\sum_{i=1}^{n}\lambda_{i,j}e_i$$ Alors on a : $$f(x_1,\ldots,x_n)=\left(\sum_{\sigma\in\mathcal{S}_n}\varepsilon(\sigma)\prod_{i=1}^{n}\lambda_{\sigma(i),i}\right)f(e_1,\ldots,e_n)$$
\end{lem}
\begin{dfn}[Déterminant dans une base]
	Soit $(e_i)_{1\leqslant i\leqslant n}$ une base de $E$. Alors, il existe une unique forme $n$-linéaire alternée $f$ telle que $f(e_1,\ldots,e_n)=1$. On l'appelle \textbf{déterminant dans la base $(e_i)_{1\leqslant i\leqslant n}$}, et on la note $\det\nolimits_{(e_i)}$. Elle vérifie :$$\det\nolimits_{(e_i)}(x_1,\ldots,x_n)=\sum_{\sigma\in\mathcal{S}_n}\varepsilon(\sigma)\prod_{i=1}^{n}\lambda_{\sigma(i),i}$$
\end{dfn}
\begin{cor}
	Toute forme $n$-linéaire alternée est donc proportionnelle à $\det\nolimits_{(e_i)}$. Le sous-espace vectoriel des formes $n$-linéaires alternées est donc de dimension $1$, engendré par $\det\nolimits_{(e_i)}$.
\end{cor}
\begin{cor}[Relation de Chasles pour les déterminants]
	Soit $A$ et $B$ deux bases de $E$. Alors $$\det\nolimits_{A}(\cdot)=\det\nolimits_{A}(B)\times\det\nolimits_{B}(\cdot)$$
\end{cor}
\begin{thm}[Lien entre déterminants et bases]
	Soit $(e_i)_{1\leqslant i\leqslant n}$ une base de $E$ et $(x_1,\ldots,x_n)$ une famille de vecteurs de $E$. Alors $(x_1,\ldots,x_n)$ est une base de $E$ si, et seulement si, $$\det\nolimits_{(e_i)}(x_1,\ldots,x_n)\ne 0$$
\end{thm}
\begin{rmq} C'est là la principale utilité du déterminant : savoir si une famille est une base, ou savoir si une matrice est inversible (ce qui revient au même, le second point étant une traduction du premier).\textbf{ Se rappeler de montrer que la matrice d'une famille dans une certaine base est inversible permet donc de montrer efficacement que cette famille est une base, sans passer par les caractères libre, générateur ou des questions de taille et de dimension. C'est une méthode très importante à retenir.} \end{rmq}
\subsection{Déterminant d'un endomorphisme}
\begin{lem}
	Soit $u\in\mathcal{L}(E)$ et $(e_i)_{1\leqslant i\leqslant n}$ une base de $E$. On a $$\forall(x_1,\ldots,x_n)\in E^n,\ \det\nolimits_{(e_i)}(u(x_j))=\det\nolimits_{(e_i)}(u(e_i))\times\det\nolimits_{(e_i)}(x_j)$$
\end{lem}
\begin{dfn}[Déterminant d'un endomorphisme]
	Soit $u\in\mathcal{L}(E)$ et $(e_i)_{1\leqslant i\leqslant n}$ une base de $E$. Le scalaire $\det\nolimits_{(e_i)}(u(e_i))$ est indépendant du choix de la base $(e_i)_{1\leqslant i\leqslant n}$. On l'appelle le \textbf{déterminant de $u$} et on le note $\det(u)$.
\end{dfn}
\begin{prop}
	Avec les mêmes hypothèses, on en tire l'identité : $$\forall(x_1,\ldots,x_n)\in E^n,\ \det\nolimits_{(e_i)}(u(x_1),\ldots,u(x_n))=\det(u)\det\nolimits_{(e_i)}(x_1,\ldots,x_n)$$
\end{prop}
\begin{exa}
	En particulier, on a $\det(\id)=1$.
\end{exa}
\begin{thm}
	$u$ est un automorphisme de $E$ si, et seulement si, $\det(u)\ne 0$.
\end{thm}
\begin{thm}[Propriétés du déterminant des endomorphismes]
	On a les relations suivantes :
	\begin{enumerate}
		\item $\forall\lambda\in\K,\ \forall u\in\mathcal{L}(E),\ \det(\lambda u)=\lambda^n\det(u)$
		\item $\forall(u,v)\in\mathcal{L}(E)^2,\ \det(u\circ v)=\det(u)\det(v)$
	\end{enumerate}
\end{thm}
\begin{cor}
	Si $u$ est inversible, alors on a $\displaystyle\det(u\inv)=\frac{1}{\det(u)}$.
\end{cor}
\begin{rmq} Attention, on a aucune raison d'avoir $\det(u+v)=\det(u)+\det(v)$ ! \end{rmq}
\subsection{Déterminant d'une matrice carrée}
\begin{dfn}[Déterminant d'une matrice carré - Formule des signatures]
	Soit $A\in\mathcal{M}_n(\K)$. Le \textbf{déterminant} de $A$, noté $\det(A)$, est défini comme le déterminant des colonnes de $A$ dans la base canonique de $\K^n$ : $$\det(A)=\sum_{\sigma\in\mathcal{S}_n}\varepsilon(\sigma)\prod_{i=1}^{n}a_{\sigma(i),i}$$
\end{dfn}
\begin{rmq} Attention le déterminant d'une matrice non carrée n'a aucun sens ! \end{rmq}
\begin{exa}
	Si $A=(\lambda)\in\mathcal{M}_1(\K)$, alors $\det(A)=\lambda$. Si $n=0$, les différentes conventions sur l'ensemble vide montrent que $\det(A)=1$.
\end{exa}
\begin{prop}
	Soit $(e_i)_{i\leqslant i\leqslant n}$ une base de $E$ et $(x_i)_{1\leqslant i\leqslant n}$ une famille de vecteurs. Alors :
	$$\det\nolimits_{(e_i)}(x_j)=\det(\mat_{(e_i)}(x_j))$$ De même, si $u\in\mathcal{L}(E)$, on a $$\det(u)=\det(\mat_{(e_i)}(u))$$
\end{prop}
\begin{thm}[Règle du gamma - Calcul d'un déterminant $2\times2$]
	On a la formule $$\begin{vmatrix}x_1&x_2\\y_1&y_2\end{vmatrix}=x_1y_2-y_1x_2$$
\end{thm}
\begin{thm}[Règle de Sarrus - Calcul d'un déterminant $3\times3$]
	On a la formule $$\begin{vmatrix}x_1&x_2&x_3\\y_1&y_2&y_3\\z_1&z_2&z_3\end{vmatrix}=x_1y_2z_3+y_1z_2x_3+z_1x_2y_3-x_1z_2y_3-y_1x_2z_3-z_1y_2x_3$$ Pour la retenir, on fait toutes les diagonales. Celles dans le sens de la première bissectrice du repère prennent un moins, celles dans le sens de la seconde bissectrice du repère prennent un plus.
\end{thm}
\begin{thm}[Propriétés du déterminant matriciel]
	On a les relations suivantes :
	\begin{enumerate}
		\item $\forall\lambda\in\K,\ \forall A\in\mathcal{M}_n(\K),\ \det(\lambda A)=\lambda^n\det(A)$
		\item $\forall(A,B)\in\mathcal{M}_n(\K)^2,\ \det(AB)=\det(A)\det(B)$
	\end{enumerate}
	De plus, $A$ est inversible si, et seulement si, $\det(A)\ne 0$, et dans ce cas on a $$\det(A\inv)=\frac{1}{\det(A)}$$
\end{thm}
\begin{rmq} Attention, on a aucune raison d'avoir $\det(A+B)=\det(A)+\det(B)$ ! \end{rmq}
\begin{prop}[Invariance par similitude du déterminant]
	Le déterminant est un invariant de similitude.
\end{prop}
\begin{prop}[Formules de Cramer, HP]
	Soit un système de Cramer $AX=B$. On dispose des formules de Cramer $$\forall i \in\llbracket1,n\rrbracket,\ x_i=\frac{\det\left(A_1|\ldots|A_{i-1}|B|A_{i+1}|\ldots|A_n\right)}{\det(A)}$$
\end{prop}
\begin{proof}
	En effet, on sait qu'une des interprétations possibles du système de Cramer consiste à chercher $B$ comme combinaison linéaire des colonnes de $A$. Ici, on cherche donc $x_1,\ldots,x_n$ tels que $B=x_1A_1+\ldots+x_nA_n$. Il suffit alors d'utiliser la $n$-linéarité et le caractère alterné à la matrice dont on a remplacé $A_i$ par $B$.
\end{proof}
\begin{rmq} Cette technique n'est pratique qu'à l'ordre 2 ou 3 en général. \end{rmq}
\begin{thm}[Invariance du déterminant par transposition]
	Le déterminant est invariant par transposition : $$\det(A)=\det(A^\top)$$
\end{thm}
\begin{cor}
	$\det(A)$ n'est pas seulement $n$-linéaire alterné par rapport aux colonnes de $A$, mais aussi par rapport aux lignes de $A$.
\end{cor}
\subsection{Calcul pratique}
\begin{thm}[Effet des opérations élémentaires sur le déterminant d'une matrice]
	Qu'elles soient sur les lignes ou sur les colonnes, les opérations élémentaires ont l'effet suivant sur le déterminant d'une matrice.
	\begin{enumerate}
		\item Transvection : Aucun effet
		\item Dilatation d'un facteur $\lambda$ : Multiplication du déterminant par $\lambda$
		\item Échange : Multiplication du déterminant par $-1$
	\end{enumerate}
\end{thm}
\begin{lem}[Calcul d'un déterminant par pivot]
	Soit $n\geqslant2$ et $A\in\mathcal{M}_n(\K)$ une matrice de la forme $$A=\left(\begin{array}{c|ccc}
		\lambda&*&\ldots&*\\
		\hline
		0&&& \\
		\vdots&&\tilde{A}& \\
		0&&&
	\end{array}\right)$$
	Alors on a $\det(A)=\lambda\det(A)$
\end{lem}
\begin{rmq} On pourra utiliser ce lemme pour calculer très efficacement un déterminant. \end{rmq}
\begin{cor}[Déterminant d'une matrice triangulaire]
	Le déterminant d'une matrice triangulaire est égal au produit de ses termes diagonaux.
\end{cor}
\begin{cor}[Déterminant d'une matrice triangulaire]
	Le déterminant d'une matrice diagonale est égal au produit de ses termes diagonaux.
\end{cor}
\begin{dfn}[Mineur, cofacteur]
	Soit $n\geqslant2$ et $A=(a_{i,j})\in\mathcal{M}_n(\K)$. On appelle \textbf{mineur} du terme $a_{i,j}$ le déterminant noté $\Delta_{i,j}$ obtenu en supprimant la $i$-ème ligne et la $j$-ème colonne (concrètement, on supprime la ligne et la colonne du terme $a_{i,j}$ et on calcule le déterminant de "ce qu'il reste"). Le scalaire $A_{i,j}=(-1)^{i+j}\Delta_{i,j}$ s'appelle le \textbf{cofacteur} du terme $a_{i,j}$. 
\end{dfn}
\begin{thm}[Développement selon une ligne ou une colonne]
	Soit $n\geqslant2$ et $A\in\mathcal{M}_n(\K)$. En notant $A_{i,j}$ les cofacteurs de $A$, on a :
		\begin{enumerate}
			\item $\displaystyle\forall i\in\llbracket1,n\rrbracket,\ \sum_{j=1}^{n}a_{i,j}A_{i,j}=\det(A)$ (développement selon la $i$-ème ligne)
			\item $\displaystyle\forall j\in\llbracket1,n\rrbracket,\ \sum_{i=1}^{n}a_{i,j}A_{i,j}=\det(A)$ (développement selon la $j$-ème colonne)
		\end{enumerate}
\end{thm}
\begin{thm}[Déterminant d'une matrice triangulaire par blocs]
	Soit $A$ une matrice triangulaire par blocs. Alors le déterminant de $A$ est égal au produit des déterminants des blocs diagonaux.
\end{thm}
\begin{cor}[Déterminant d'une matrice diagonale par blocs]
	Si $A$ est diagonale par blocs, son déterminant vaut le produit des déterminants des blocs diagonaux.
\end{cor}
\begin{dfn}[Comatrice]
	Si $n\geqslant2$, la matrice des cofacteurs de $A$ s'appelle la \textbf{comatrice} de $A$. On la note $\com(A)$.
\end{dfn}
\begin{thm}[Formule de la comatrice]
	On a la relation $$A\com(A)^\top=\com(A)^\top A=\det(A)I_n$$
\end{thm}
\begin{rmq} Si $n=1$, les différentes conventions sur l'ensemble vide montrent que $\com(A)=1$ et la formule reste valable. \end{rmq}
\begin{rmq} Dans le cas où $A$ est inversible, cette dernière formule peut éventuellement servir à calculer $$A\inv=\frac{1}{\det(A)}\com(A)^\top $$ mais elle a un intérêt plus théorique que pratique, car elle devient rapidement inexploitable. On lui préfèrera l'utilisation du pivot de Gauss. Toutefois, cette formule est \textbf{très utile} dans les exercices faisant appel à la comatrice ou à la densité de de $\mathcal{GL}_n(\K)$ dans $\mathcal{M}_n(\K)$. \end{rmq}

\begin{thm}(Déterminant de Vandermonde)
	Soit $a_1,\ldots,a_n$ des scalaires de $\K$. Alors $$\begin{vmatrix}
		1&a_1&\ldots&a_1^{n-1}\\
		1&a_2&\ldots&a_2^{n-1}\\
		\vdots&\vdots&&\vdots \\
		1&a_n&\ldots&a_n^{n-1}
	\end{vmatrix} = \prod_{i<j}(a_j-a_i)$$
	En particulier, le déterminant est non nul si, et seulement, les $a_i$ sont deux à deux distincts.
\end{thm}
\begin{rmq} On rencontrera parfois ce déterminant sous sa forme transposée, ce qui ne change rien car le déterminant est invariant par transposition. \end{rmq}
\newpage
\chapter{Polynômes}
Dans tout le chapitre, on fixe $\K$ un corps.
\section{Définition formelle}
\begin{dfn}[Polynôme à une indéterminée sur le corps $\K$]
	Un \textbf{polynôme à une une indéterminée sur le corps $\K$} est une suite presque nulle $(a_k)_{k\in\N}\in\K^{(\N)}$. \\ \par En particulier, deux polynômes sont égaux si, et seulement si, ils ont les mêmes coefficients.
\end{dfn}
\begin{prop}[Rappels : structure de $\K$-espace vectoriel et base canonique]
	On rappelle que $\K^{(\N)}$ est un $\K$-espace vectoriel pour les lois suivantes :
	\begin{enumerate}
		\item $(a_k)_{k\in\N}+(b_k)_{k\in\N}=(a_k+b_k)_{k\in\N}$
		\item $\lambda.(a_k)_{k\in\N}=(\lambda\times a_k)_{k\in\N}$
	\end{enumerate}
	Aussi, on pose $$\forall k\in\N,\ e_k=(\delta_{k,l})_{l\in\N}\in\K^{(\N)}$$ La famille $(e_k)_{k\in\N}$ est une base de $\K^{(\N)}$, appelée \textbf{base canonique} de $\K^{(\N)}$.
\end{prop}
\begin{rmq} Soit $(a_k)_{k\in\N}\in\K^\N$. Les deux assertions suivantes sont équivalentes :
	\begin{enumerate}
		\item $(a_k)_{k\in\N}\in\K^{(\N)}$
		\item $\exists n\in\N,\ \forall k>n,\ a_k=0_\K$
\end{enumerate} \end{rmq}
\begin{dfn}[Produit de polynômes]
	On munit $\K^{(\N)}$ d'une loi $\times$ de la manière suivante. Soit $P=(a_k)_{k\in\N}\in\K^{(\N)}$ et $Q=(b_k)_{k\in\N}\in\K^{(\N)}$. On définit alors le polynôme $P\times Q$ comme le polynôme $P\times Q=(c_k)_{k\in\N}\in\K^{(\N)}$ où on a posé $$\forall k\in\N,\ c_k=\sum_{i+j=k}a_ib_j=\sum_{i=0}^{k}a_ib_{k-i}=\sum_{j=0}^{k}a_{k-j}b_j$$
\end{dfn}
\begin{exa}
	Le terme constant d'un produit de polynômes est égal au produit des termes constants.
\end{exa}
\begin{prop}[Propriétés de $\times$]
	Sur $\K^{(\N)}$, la loi $\times$ définie comme précédemment est :
	\begin{enumerate}
		\item \textbf{associative}
		\item \textbf{commutative}
		\item \textbf{bilinéaire} ; c'est-à-dire linéaire par rapport au premier facteur et linéaire par rapport au deuxième facteur
	\end{enumerate}
\end{prop}
\begin{exa}
	En particulier, $\times$ est ainsi compatible avec $\cdot$ et distributive sur $+$.
\end{exa}
\begin{cor}[Structure de $\K$-algèbre]
	$(\K^{(\N)},+,\times,\cdot)$ est une $\K$-algèbre commutative, d'élément neutre pour $\times$ le polynôme $e_0=(\delta_{j,0})_{j\in\N}$
\end{cor}
\begin{dfn}[Indéterminée]
	On pose $$X=(\delta_{k,1})_{k\in\N}$$ On appelle $X$ l'\textbf{indéterminée} (même si $X$ est parfaitement déterminé et est une constante dans l'ensemble des polynômes, ie ce polynôme ne change pas, mais ce n'est pas non plus un polynôme constant).
\end{dfn}
\begin{lem}
	On a $$\forall k\in\N,\ X^k=(\delta_{j,k})_{j\in\N}$$
\end{lem}
\begin{dfn}[Notation]
	Dans ce contexte, on notera plutôt $\K^{(\N)}$ de la façon suivante : $$\K[X]$$ qu'on lira "$\K$ crochet $X$" ou plus exactement "l'ensemble des polynômes à une indéterminée sur le corps $\K$".
\end{dfn}
\begin{cor}
	Ainsi, $(X^k)_{k\in\N}$ est donc la \textbf{base canonique} de $\K[X]$ et tout polynôme $P=(a_k)_{k\in\N}\in\K[X]$ se décompose sur cette base sous la forme $$P=\sum_{k\in\N}a_kX^k$$
\end{cor}
\begin{dfn}[Polynôme constant]
	Soit $\lambda\in\K$. le \textbf{polynôme constant égal à $\lambda$} est le polynôme $\lambda.X^0$. C'est donc la suite $(\lambda,0,0,\ldots)$. En pratique, on pourra l'écrire $\lambda$ car l'application
	$$\begin{array}[t]{lrcl}
		\K & \longrightarrow & \K[X] \\
		 \lambda & \longmapsto & \lambda.X^0
	\end{array}$$ est un morphisme injectif de $\K$-algèbres. On identifiera alors $\K$ à l'ensemble des polynômes constants.
\end{dfn}
\begin{dfn}[Degré]
	Soit $P=\displaystyle\sum_{k\in\N}a_kX^k\in\K[X]$. On définit son \textbf{degré} de la manière suivante :
	\begin{enumerate}
		\item Si $P=0_{\K[X]}$, on pose $\deg(P)=-\infty$
		\item Si $P\ne0_{\K[X]}$, alors le support de la famille $(a_k)_{k\in\N}$ est non vide et fini par hypothèse, donc on pose alors $\deg(P)$ comme étant le maximum de ce support. C'est donc alors le plus grand $n\in\N$ tel que $a_n=0_\K$.
	\end{enumerate} 
\end{dfn}
\begin{dfn}[Coefficient dominant, terme dominant]
	Lorsque $P=\displaystyle\sum_{k\in\N}a_kX^k\in\K[X]$ est non nul de degré $\deg(P)=n\in\N$, on appelle \textbf{coefficient dominant de $P$} le scalaire $a_n$, et on appelle \textbf{terme dominant de $P$} le polynôme $a_nX^n$.
\end{dfn}
\begin{dfn}[Polynôme unitaire]
	Si $P=\displaystyle\sum_{k\in\N}a_kX^k\in\K[X]$ est non nul de degré $\deg(P)=n\in\N$, on dit que $P$ est \textbf{unitaire} lorsque son coefficient dominant vaut $1_k\K$, c'est-à-dire lorsque $a_n=1_\K$.
\end{dfn}
\begin{exa}
	Pour tout $n\in\N$, $X^n$ est unitaire de degré $n$.
\end{exa}
\begin{dfn}
	Soit $n\in\N$. On pose : $$\K_n[X]=\left\{P\in\K[X]\ | \ \deg(P)\leqslant n\right\}$$
\end{dfn}
\begin{rmq}
	On a $$\bigcup_{n\in\N}\K_n[X]=\K[X]$$ et $$\forall n\in\N,\ \K_n[X]\subset\K_{n+1}[X]$$
\end{rmq}
\begin{prop}[Propriété utile]
	Soit $n\in\N$ et $P\in\K_[X]$. On a :
	\begin{enumerate}
		\item $\displaystyle\deg(P)\leqslant n \iff \exists(a_0,\ldots,a_n)\in\K^{n+1}, \ P=\sum_{k=0}^{n}a_kX^k$
		\item $\displaystyle\deg(P)=n \iff \exists(a_0,\ldots,a_n)\in\K^n\times\K,\ P=\sum_{k=0}^{n}a_kX^k$
	\end{enumerate}
\end{prop}
\begin{prop}
	Soit $n\in\N$. On a les propriétés suivantes pour $\K_n[X]$ :
	\begin{enumerate}
		\item $\K_n[X]$ est un sous-espace vectoriel de $\K[X]$
		\item $\K_n[X]$ est de dimension finie $n+1$
		\item $(X^0,\ldots,X^n)$ est une base de $\K_n[X]$, appelée \textbf{base canonique} de $\K_n[X]$
	\end{enumerate}
\end{prop}
\begin{thm}[Degré d'un produit et d'une somme]
	Soit $P,Q\in\K[X]$. On a les propriétés suivantes sur le degré de $P+Q$ :
	\begin{enumerate}
		\item $\deg(P+Q)\leqslant\max(\deg(P),\deg(Q))$
		\item Si $\deg(P)\ne\deg(Q)$, alors $\deg(P+Q)=\max(\deg(P),\deg(Q))$ et le coefficient dominant de $P+Q$ est égal au coefficient du polynôme de plus haut degré.
	\end{enumerate}
	On a les propriétés suivantes sur le degré de $PQ$ :
	\begin{enumerate}
		\item $\deg(PQ)=\deg(P)+\deg(Q)$
		\item Si $P$ et $Q$ sont tous deux non nuls, alors $PQ$ est non nul et le coefficient dominant de $PQ$ est égal au produit des coefficients dominant de $P$ et de $Q$.
	\end{enumerate}
\end{thm}
\begin{rmq}
	On peut généraliser ces propriétés à un nombre fini de termes ou de facteurs en effectuant une récurrence.
\end{rmq}
\begin{cor}[Intégrité de l'anneau des polynômes]
	$(\K[X],+,\times)$ est un anneau intègre.
\end{cor}
\begin{prop}[HP]
	Soit $P\in\K[X]$ et $n\in\N$. On a $$\deg(P)=n\iff\exists(\lambda,Q)\in\K\st\times\K[X], \ P=\lambda X^n+Q \ \ \text{et} \ \ \deg(Q)<n$$
\end{prop}
\begin{proof}
	Pour $\Rightarrow$, utiliser la "propriété utile" et distinguer les cas $n=0$ et $n\geqslant$. Pour $\Leftarrow$, fixer les objets dont on suppose l'existence et utiliser les propriétés du degré.
\end{proof}
\begin{cor}[Renormalisé]
	Supposons que $P$ est un polynôme non nul. En notant $CD(P)$ le coefficient dominant de $P$, on a le fait que $$\frac{1}{CD(P)}.P$$ est unitaire. On dit qu'on a \textbf{renormalisé} $P$.
\end{cor}
\begin{thm}[Groupe des unités de l'anneau des polynômes]
	L'ensemble des unités de l'anneau $(\K[X],+,\times)$ vaut : $$\mathcal{U}(\K[X])=\left\{P\in\K[X]\ | \ \deg(P)=0\right\}=\K_0[X]\setminus\left\{0_{\K[X]}\right\}$$
\end{thm}
\begin{rmq}[Importante pour X/ENS]
	Si $(A,+,\times)$ est un anneau \textbf{intègre}, tout le début du chapitre fonctionne de même, à l'exception du théorème précédent. Dans ce cas, les unités de $A[X]$ sont exactement les polynômes constants égaux aux unités de l'anneau $A$.
\end{rmq}
\begin{dfn}[Polynôme composé]
	Soit $P,Q\in\K[X]$. Le \textbf{polynôme composé $P\circ Q$} est défini de la manière suivante : on écrit $P$ sous la forme $\displaystyle \sum_{k\in\N}a_kX^k$ puis on pose $$P\circ Q=\sum_{k\in\N}a_kQ^k$$
\end{dfn}
\begin{prop}[Degré de $P\circ Q$ avec $Q$ non constant, HP]
	Si $Q$ est un polynôme non constant (c'est-à-dire que $\deg(Q)\geqslant1$), alors on a : $$\deg(P\circ Q)=\deg(P)\times\deg(Q)$$
\end{prop}
\begin{proof}
	Si $P$ est nul, on conclut immédiatement. Si $P$ est non nul,on pose $n=\deg(P)$ et on écrit $$P\circ Q=a_nQ^n+\sum_{k=0}^{n-1}a_kQ^k$$ Or, on a $\deg(a_nQ^n)=n\deg(Q)$ et si $n=0$, le deuxième polynôme est nul, et si $n$ est non nul, le deuxième polynôme est de degré inférieur à $\max(\deg(Q^0),\ldots,\deg(Q^{n-1}))$ donc de degré strictement inférieur à $n\deg(Q)$. On conclut alors par les propriétés sur le degré d'une somme.
\end{proof}
\section{Arithmétique dans $\K[X]$}
\begin{dfn}[Divisibilité (rappel)]
	Soit $A,B\in\K[X]$. On dit que $B$ \textbf{divise} $A$, et on note $B\mid A$ lorsque $$\exists P\in\K[X],\ A=BP$$ c'est-à-dire lorsque $A\in B\K[X]$ ou encore lorsque $A\K[X]\subset B\K[X]$.
\end{dfn}
\begin{prop}[$\mid$ est un préordre]
	La relation $\mid$ est réflexive et transitive, mais pas antisymétrique. On dit que c'est un \textbf{préordre}.
\end{prop}
\begin{prop}(Propriétés de $\mid$)
	La relation $\mid$ vérifie les propriétés suivantes : 
	\begin{enumerate}
		\item Soit $D$ un diviseur de $A$ et $B$. Alors $$\forall(U,V)\in\K[X]^2,\ D\mid AU+BV$$
		\item Soit $Q\in\K[X]$. Si $D\mid A$, alors $QD\mid QA$. Réciproquement, si $QD\mid QA$ \textbf{et} si $Q\ne 0$, alors $D\mid A$.
		\item Si $D\mid A$ et $A\ne 0$, alors $\deg(D)\leqslant\deg(A)$
	\end{enumerate}
\end{prop}
\begin{dfn}[Association]
	Soit $A,B\in\K[X]$. Les deux assertion suivantes sont équivalentes :
	\begin{enumerate}
		\item $A\mid B$ et $B\mid A$
		\item $\exists\lambda\in\K\st,\ B=\lambda A$
	\end{enumerate}
	Dans ce cas, $A$ et $B$ sont dits \textbf{associés}. La relation d'\textbf{association} est une relation d'équivalence. Nous la noterons ici $\sim$ de façon non standard.
\end{dfn}
\begin{thm}
	Chaque classe d'équivalence pour la relation d'association admet un unique représentant unitaire ou nul (qu'on abrègera en UON, non standard).
\end{thm}
\begin{prop}
	Deux polynômes associés ont le même degré. En revanche la réciproque est fausse, on pourra considérer $X$ et $X+1$.
\end{prop}
\begin{exa}
	$A\sim B \iff A\K[X]=B\K[X]$ (utiliser une des définitions alternatives de la divisibilité)
\end{exa}
\begin{rmq}[Transformation d'un préordre en ordre, HP]
	On rappelle que la divisibilité n'est pas antisymétrique. Toutefois, on peut la rendre antisymétrique en créant une divisibilité induite sur les classes d'association. \par De façon générale, si on dispose d'un préordre $\mathcal{R}$ sur un ensemble $E$, on peut "transformer" ce préordre en ordre en quotientant $E$ par le relation définie par $x\sim y \iff (x\mathcal{R}y\land y\mathcal{R}x)$ puis en faisant passer le préordre au quotient. Il s'agira alors toujours d'une relation d'ordre (cf. exercice étoilé d'introduction à l'algèbre).
\end{rmq}
\begin{thm}[Division euclidienne]
	Soit $(A,B)\in\K[X]\times(\K[X]\setminus\left\{0_{\K[X]}\right\})$. Alors : $$\exists!(Q,R)\in\K[X]^2,\ A=BQ+R \ \ \text{et} \ \ \deg(R)<\deg(B)$$
\end{thm}
\begin{thm}[Classification des idéaux de l'anneau des polynômes]
	Soit $\mathcal{I}$ un idéal de $\K[X]$. Alors, il existe un unique polynôme $P\in\K[X]$ unitaire ou nul tel que $$\mathcal{I}=P\K[X]$$
\end{thm}
\begin{cor}[Principalité de l'anneau des polynômes]
	L'anneau $(\K[X],+,\times)$ est principal.
\end{cor}
\begin{rmq}[HP]
	On aurait déjà pu obtenir la principalité de $\K[X]$ puisque $\K[X]$ dispose d'une division euclidienne. En effet, tout anneau euclidien est principal (voir compléments sur les anneaux).
\end{rmq}
\begin{cor}
	Les résultats arithmétiques usuels s'appliquent dans $\K[X]$ :
	\begin{enumerate}
		\item Soit $A,B\in\K[X]$. Les unités de $\K[X]$ (ie les $\lambda\in\K\st$) divisent toujours $A$ et $B$. Réciproquement, si ce sont les seuls diviseurs communs de $A$ et $B$, $A$ et $B$ sont dits \textbf{premiers entre eux}.
		\item On définit de même les polynômes $A_1,\ldots,A_n$ \textbf{premiers entre eux dans leur ensemble}.
		\item \textbf{Théorème de Bézout} : Soit $(A,B)\in\K[X]^2$. $A$ et $B$ sont premiers ente eux si, et seulement si, $\exists(U,V)\in\K[X]^2, \ AU+BV=1$
		\item \textbf{Généralisation du théorème de Bézout} : Soit $n\geqslant2$ et $(A_1,\ldots,A_n)\in\K[X]^n$. $A_1,\ldots,A_n$ sont premiers entre eux dans leur ensemble si, et seulement si, $$\exists(U_1,\ldots,U_n)\in\K[X]^n,\ \sum_{i=1}^{n}A_iU_i=1$$
		\item\textbf{ Corollaire du théorème de Bézout} : Soit $n\geqslant2$ et $(A,B_1,\ldots,B_n)\in\K[X]^{n+1}$. Si $A$ est premier avec chacun des $B_i$, alors $A$ est premier avec leur produit.  On peut par exemple appliquer deux fois ce résultat si des $A_1,\ldots,A_m$ sont tous premiers avec chacun des $B_i$, alors le produit des $A_i$ est premier avec le produit des $B_i$.
		\item \textbf{Lemme de Gauss }: Soit $(A,B,C)\in\K[X]^3$. Si $A\mid BC$ et $A$ est premier avec $B$, alors $A\mid C$.
		\item \textbf{Corollaire du lemme de Gauss} : Soit $n\geqslant2$ et $(A_1,\ldots,A_n,B)\in\K[X]^{n+1}$. Si les $A_i$ sont deux à deux premiers entre eux et si $\forall i\in\llbracket1,n\rrbracket,\ A_i\mid B$, alors $$A_1\times\ldots\times A_n\mid B$$
	\end{enumerate}
\end{cor}
\begin{dfn}[\textbf{Un} PGCD, \textbf{un} PPCM]
	Soit $A$ et $B$ des polynômes tous deux non nuls.
	\begin{enumerate}
		\item On appelle \textbf{un PGCD} de $A$ et $B$ tout diviseur commun non nul de $A$ et $B$ de degré maximal.
		\item On appelle \textbf{un PPCM} de $A$ et $B$ tout multiple commun non nul de $A$ et $B$ de degré minimal.
	\end{enumerate}
\end{dfn}
\begin{thm}
	Soit $A$ et $B$ des polynômes tous deux non nuls.
	\begin{enumerate}
		\item Les PGCD de $A$ et $B$ forment une classe d'association.
		\item Les PPCM de $A$ et $B$ forment une classe d'association.
	\end{enumerate}
\end{thm}
\begin{dfn}[\textbf{Le} PGCD, \textbf{le} PPCM]
	Soit $A$ et $B$ des polynômes tous deux non nuls.
	\begin{enumerate}
		\item La classe d'association des PGCD de $A$ et $B$ admet un unique représentant UON. On l'appelle \textbf{le PGCD} de $A$ et $B$, et on le note $A\wedge B$. Il vérifie : $$A\K[X]+B\K[X]=(A\wedge B)\K[X]$$
		\item La classe d'association des PPCM de $A$ et $B$ admet un unique représentant UON. On l'appelle \textbf{le PPCM} de $A$ et $B$, et on le note $A\vee B$. Il vérifie : $$A\K[X]\cap B\K[X]=(A\vee B)\K[X]$$
	\end{enumerate}
\end{dfn}
\begin{dfn}[Extension du PGCD et du PPCM]
	Soit $A$ et $B$ des polynômes quelconques (éventuellement nuls).
	\begin{enumerate}
		\item $A\K[X]+B\K[X]$ est un idéal de $\K[X]$ dont $$\exists!P\ \text{UON},\ A\K[X]+B\K[X]=P\K[X]$$ On pose alors $A\wedge B=P$.
		\item $A\K[X]\cap B\K[X]$ est un idéal de $\K[X]$ dont $$\exists!P\ \text{UON},\ A\K[X]\cap B\K[X]=P\K[X]$$ On pose alors $A\vee B=P$.
	\end{enumerate}
	Cette définition prolonge celle donnée précédemment, et on conserve les appellations "le PGCD" et "le PPCM".
\end{dfn}
\begin{prop}[Théorème de Bézout généralisé]
	Soit $A$ et $B$ deux polynômes quelconques. Alors $$\exists (U,V)\in\K[X]^2,\ AU+BV=A\wedge B$$
\end{prop}
\begin{rmq}
	L'ensemble des diviseurs communs de $A$ et $B$ est égal à l'ensemble des diviseurs de $A\wedge B$. L'ensemble des multiples communs de $A$ et $B$ est égal à l'ensemble des multiples de $A\vee B$.
\end{rmq}
\begin{rmq}
	Dans l'ensemble des polynômes UON, la divisibilité induit une relation d'ordre : en effet, deux polynômes UON associés sont égaux. Alors, $A\wedge B$ est le maximum pour la divisibilité induite des polynômes UON divisant $A$ et $B$, et $A\vee B$ est le minimum pour la divisibilité induite des polynômes UON multiples de $A$ et de $B$.
\end{rmq}
\begin{dfn}[PGCD et PPCM de plusieurs polynômes]
	Soit $n\geqslant2$ et $A_1,\ldots,A_n$ des polynômes quelconques (éventuellement nuls).
	\begin{enumerate}
		\item $A_1\K[X]+\ldots+A_n\K[X]$ est un idéal de $\K[X]$ dont $$\exists!P\ \text{UON},\ A_1\K[X]+\ldots+A_n\K[X]=P\K[X]$$ On pose alors $A_1\wedge\ldots\wedge A_n=P$. On dit que $P$ est le PGCD de $A_1,\ldots,A_n$.
		\item $A_1\K[X]\cap\ldots\cap A_n\K[X]$ est un idéal de $\K[X]$ dont $$\exists!P\ \text{UON},\ A_1\K[X]\cap\ldots\cap A_n\K[X]=P\K[X]$$ On pose alors $A_1\vee\ldots\vee A_n=P$. On dit que $P$ est le PPCM de $A_1,\ldots,A_n$.
	\end{enumerate}
\end{dfn}
\begin{prop}[Associativité et commutativité du PGCD et du PPCM]
	Le PGCD et le PPCM sont associatifs et commutatifs.
\end{prop}
\begin{prop}[invariance par association du PGCD et du PPCM]
	Le PGCD et le PPCM sont invariants si on remplace $A$ ou $B$ par un de leurs associés.
\end{prop}
\begin{prop}[Distributivité du PGCD]
	Le PGCD est distributif : $$\forall (A,B)\in\K[X]^2,\ \forall P \ \text{UON},\ (PA)\wedge(PB)=P(A\wedge B)$$
\end{prop}
\begin{prop}
	$A$ et $B$ sont premiers entre eux si, et seulement si, $A\wedge B=1$. 	$A_1,\ldots,A_n$ sont premiers entre eux dans leur ensemble si, et seulement si, $A_1\wedge\ldots\wedge A_n=1$
\end{prop}
\begin{prop}[Propriété utile]
	Soit $(A,B)\in\K[X]^2$. On note $\Delta=A\wedge B$? Alors, il existe $(\tilde{A},\tilde{B})\in\K[X]^2$ tel que $A=\Delta\tilde{A}$, $B=\Delta\tilde{B}$ et $\tilde{A}\wedge\tilde{B}=1$.
\end{prop}
\begin{thm}[Invariance par transvection]
	Soit $(A,B)\in\K[X]^2$. On a $$\forall Q\in\K[X],\ A\wedge B = (A-BQ)\wedge B$$ En conséquence, l'algorithme d'Euclide est valable et peut être appliqué pour déterminer un PGCD.
\end{thm}
\begin{dfn}[Congruence modulo $P$]
	Soit $P\in\K[X]$ non nul. On dit que $A$ est \textbf{congru à $B$ modulo $P$}, et on note $A\equiv B \ [P]$ lorsque $P\mid A-B$. Les propriétés usuelles restent valables :
	\begin{enumerate}
		\item $\cdot\equiv\cdot\ [P]$ est une relation d'équivalence
		\item Si $A_1\equiv B_1 \ [P]$ et $A_2\equiv B_2 \ [P]$ alors $A_1+A_2\equiv B_1+B_2 \ [P]$ et $A_1A_2\equiv B_1B_2 \ [P]$.
		\item Si $A\equiv B \ [P]$, alors $\forall n\in\N,\ A^n\equiv B^n \ [P]$
	\end{enumerate}
\end{dfn}
\section{Lien avec les fonctions}
\subsection{Évaluation d'un polynôme}
\begin{dfn}[Évaluation en un point de $\K$]
	Soit $P\in\K[X]$ et $x\in\K$. Pour évaluer $P$ en $x$, on écrit $P$ sous la forme $P=\displaystyle\sum_{k\in\N}a_kX^k$ puis on pose : $$P(x)=\sum_{k\in\N} a_kx^k$$ On dit aussi qu'on a \textbf{substitué $x$ à $X$ dans $P$}.
\end{dfn}
\begin{thm}[Évaluation et opérations algébriques]
	Soit $(\lambda,\mu)\in\K^2$ et $(P,Q)\in\K[X]^2$. Soit $x\in\K$. On a :
	\begin{enumerate}
		\item $(\lambda P+\mu Q)(x)=\lambda P(x)+\mu Q(x)$
		\item $(PQ)(x)=P(x)Q(x)$
		\item $(P\circ Q)(x)=P(Q(x))$
	\end{enumerate}
\end{thm}
\begin{rmq}
	En fait, on a montré que $\begin{array}[t]{lrcl}
		\varphi_x: & \K[X] & \longrightarrow & \K \\
		& A & \longmapsto & A(x)
	\end{array}$ est un morphisme de $\K$-algèbres. Et on peut généraliser, ce qui fait l'objet du théorème suivant.
\end{rmq}
\begin{thm}[Évaluation en un point d'une $\K$-algèbre, HP, spé]
	Soit $\mathcal{A}$ une $\K$-algèbre quelconque et $P\in\K[X]$ qu'on écrit sous la forme $P=\displaystyle\sum_{k\in\N}a_kX^k$. Alors on peut encore poser $$\forall x\in\mathcal{A},\ P(x)=\sum_{k\in\N}a_kX^k$$ Alors, pour tout $x\in\mathcal{A}$, l'application $$\begin{array}[t]{lrcl}
		\varphi_x: & \K[X] & \longrightarrow & \mathcal{A} \\
		& Q & \longmapsto & Q(x)
	\end{array}$$ est un morphisme de $\K$-algèbres et le théorème qui précède reste valable. \par On pensera notamment aux $\K$-algèbres suivantes : $\mathcal{L}(E)$ où $E$ est un $\K$-espace vectoriel, $\mathcal{M}_n(\K)$ ou encore $\K[X]$.
\end{thm}
\begin{thm}[Lemme des noyaux, HP, vu en spé]
	Pour tout $u\in\mathcal{L}(E)$, pour tous polynômes $P$ et $Q$ premiers entre eux, on a $$\ker((PQ)(u))=\ker(P(u))\oplus\ker(Q(u))$$
\end{thm}
\begin{proof}
	Utiliser le théorème de Bézout et évaluer en $u$. Si on note $PA+QB=1$ cette relation de Bézout, montrer ensuite que $\im(P(u)\circ A(u))\subset\ker(Q(u))$ et $\im(Q(u)\circ B(u))\subset\ker(P(u))$ en se plaçant dans $\ker((PQ)(u))$. La somme est alors totale en évaluant en $x$ le Bézout évalué en $u$. On montre que la somme est directe en remarquant que $P(u)$ et $A(u)$ commutent, de même que $Q(u)$ et $B(u)$, car ce sont des polynômes en $u$.
\end{proof}
\subsection{Racines}
\begin{dfn}[Racine]
	Soit $P\in\K[X]$. une \textbf{racine} de $P$ (ou un \textbf{zéro} de $P$) est un scalaire $x\in\K$ tel que $P(x)=0_\K$. On notera ici de façon non standard $\mathcal{R}(P)$ l'ensemble des racines d'un polynôme $P$.
\end{dfn}
\begin{rmq}
	Si on cherche les racines communes de $A$ et $B$ avec $B$ non nul, alors elles sont aussi des racines du reste dans la division euclidienne de $A$ par $B$ (écrire cette division et l'évaluer en cette racine). Pn peut alors itérer pour se ramener à des polynômes de faible degré.
\end{rmq}
\begin{thm}
	Soit $x\in\K$ et $P\in\K[X]$. On a $$x\in\mathcal{R}(P) \Longleftrightarrow X-x\mid P$$
\end{thm}
\begin{rmq}[Très utile pour les exercices X/ENS]
	En utilisant la formule de Bernoulli, on obtient le résultat suivant : $$\forall P\in\Z[X],\ \forall (a,b)\in\Z^2,\ a-b\mid P(a)-P(b)$$
\end{rmq}
\begin{rmq}
	On reprend l'exemple des racines communes de $A$ et $B$ avec $B$ non nul.On a  $X-x\mid A$ et $X-x\mid B$ si, et seulement si, $X-x\mid A\wedge B$, donc les racines communes de $A$ et $B$ sont exactement les racines de $A\wedge B$.
\end{rmq}
\begin{prop}
	Soit $P\in\K[X]$ et $x_1,\ldots,x_n$ des scalaires de $\K$ \textbf{deux à deux distincts}. Les deux assertions suivantes sont équivalentes :
	\begin{enumerate}
		\item $\forall i\in\llbracket1,n\rrbracket,\ P(x_i)=0_\K$
		\item $(X-x_1)\times\ldots\times(X-x_n)\mid P$
	\end{enumerate}
\end{prop}
\begin{met}[Factorisation d'un polynôme aux racines distinctes] Voici comment on peut factoriser un polynôme $P$ dont les racines sont deux à deux distinctes.
	\begin{enumerate}
		\item On cherche les racines.
		\item On montre que celles-ci sont deux à deux distinctes (si c'est le cas évidemment, sinon on ne peut pas appliquer la méthode) A ce moment, le produit des $X-$racines divise $P$. On fixe alors le polynôme $Q$ correspondant à cette divisibilité.
		\item On examine le degré de $Q$ : s'il y a autant de racines distinctes que le degré de $P$, $Q$ est une constante.
		\item On examine le coefficient dominant de $P$ : il permet de déterminer $Q$ si $Q$ est une constante.
	\end{enumerate}
\end{met}
\begin{exa}
	On a par exemple la factorisation suivante dans $\C[X]$ : $$\forall n\in\N\st,\ X^n-1=\prod_{i=0}^{n-1}\left(X-\exp\left(i\frac{2k\pi}{n}\right)\right)$$ dont on en déduit l'identité : $$\forall n\in\N\st,\ \forall (a,b)\in\C^2,\ a^n-b^n=\prod_{i=0}^{n-1}\left(a-\exp\left(i\frac{2k\pi}{n}\right)b\right)$$
\end{exa}
\begin{cor}
	Soit $n\in\N$ et $P\in\K_n[X]$. Si $P$ admet au moins $n+1$ racines $2$ à $2$ distinctes, alors $P=0$.
\end{cor}
\begin{cor}
	On en déduit les deux résultats qui suivent :
	\begin{enumerate}
		\item Soit $n\in\N$ et $(P,Q)_in\K_n[X]^2$. Si $P$ et $Q$ coïncident (au moins) en $n+1$ points deux à deux distincts, alors $P=Q$.
		\item Soit $(P,Q)\in\K[X]$. Si $P$ et $Q$ coïncident en une infinité de points, alors $P=Q$.
	\end{enumerate}
\end{cor}
\begin{exa}
	Soit $n\in\N\st$ et $(A,B)\in\C[X]^2$. Alors on a $$A^n-B^n=\prod_{i=0}^{n-1}\left(A-\exp\left(i\frac{2k\pi}{n}\right)B\right)$$
\end{exa}
\begin{exa}
	Soit $P\in\C[X]$. Si $\forall x\in\R,\ P(x)\in\R$, alors $P\in\R[X]$. En effet, $P$ et son "polynôme conjugué" coïncident sur $\R$ qui est une partie infinie de $\C$. Donc tous les coefficients de $P$ sont égaux à leur conjugué, donc sont tous réels. Donc $P\in\R[X]$.
\end{exa}
\begin{dfn}[Ordre de multiplicité]
	Soit $P\in\K[X]$ un polynôme \textbf{non nul} et $x\in\K$. On appelle \textbf{ordre de multiplicité de $x$ en tant que racine de $P$} l'entier noté de façon non standard $m_P(x)$ défini par : $$m_P(x)=\max\left\{k\in\N : (X-x)^k\mid P\right\}$$
\end{dfn}
\begin{rmq}
	On a $$x\in\mathcal{R}(P)\Longleftrightarrow m_P(x)\geqslant1$$ De même, on a $$x\notin\mathcal{R}(P)\Longleftrightarrow m_P(x)=0$$ On peut alors parler de "racine d'ordre 0".
\end{rmq}
\begin{dfn}[Racine simple, double, triple]
	Si $m_P(x)=1$, on parle de racine \textbf{simple}. Si $m_P(x)=2$, on parle de racine \textbf{double}. Si $m_P(x)=3$, on parle de racine \textbf{triple}.
\end{dfn}
\begin{rmq}
	Comme en arithmétique, on montre que $$\forall k\in\N,\ (X-x)^k\mid P \Longleftrightarrow k\leqslant m_P(x)$$
\end{rmq}
\begin{cor}
	Soit $P$ de degré $n\in\N$. Alors $P$ admet au plus $n$ racines comptées avec leur multiplicité. Formellement, on a : $$\sum_{x\in\mathcal{R}(P)}m_P(x)\leqslant n$$
\end{cor}
\subsection{Fonctions polynomiales}
\begin{dfn}[Fonction polynomiale associée à un polynôme]
	Soit $P\in\K[X]$. La \textbf{fonction polynomiale associée à $P$} est la fonction : 
	$$\begin{array}[t]{lrcl}
		\tilde{P}: &\K & \longrightarrow & \K \\
		& x& \longmapsto & P(x)
	\end{array}$$
\end{dfn}
\begin{rmq}
	Attention, on ne peut identifier un polynôme et sa fonction polynômiale associée \textbf{que} si le corps de base $\K$ est infini. Dès qu'il est fini, il suffit de considérer le polynôme $\displaystyle\prod_{x\in\K}(X-x)$ qui est non nul en tant que polynôme, mais dont la fonction polynômiale associée est nulle.
\end{rmq}
\subsection{Dérivation formelle}
\begin{dfn}[Dérivée formelle]
	Soit $p\in\K[X]$. Écrivons-le sous la forme $P=\sum_{k\in\N}a_kX^k$. On définit alors la \textbf{dérivée formelle de $P$} par : $$P'=\sum_{k\in\N\st}ka_kX^{k-1}$$
\end{dfn}
\begin{rmq}[HP, caractéristique du corps et dérivation]
	Supposons que $k\ne0_\Z$ et $a_k\ne0_\K$. On n'a pas nécessairement $k.a_k\ne0_\K$. En effet, si $\K$ est de caractéristique $k$, rien ne va plus. Ainsi, \textbf{on suppose dans ce paragraphe que $\K$ est de caractéristique nulle}.
\end{rmq}
\begin{rmq}
	Insistons bien sur le fait qu'il s'agit d'une dérivation \textbf{formelle}. Cela dit, si $\K=\R$, on retrouve la dérivation usuelle des fonctions polynomiales.
\end{rmq}
\begin{exa}
	$(X^0)'=0$ et $\forall n\in\N\st, (X^n)'=nX^{n-1}$
\end{exa}
\begin{prop}[Combinaison linéaire]
	Soit $(\lambda,\mu)\in\K^2$ et $(P,Q)\in\K[X]^2$. On a : $$(\lambda P+\mu Q)'=\lambda P'+\mu Q'$$
\end{prop}
\begin{exa}
	Ainsi, $\begin{array}[t]{lrcl}
		D: &\K[X] & \longrightarrow & \K[X] \\
		& P& \longmapsto & P'
	\end{array}$ est un endomorphisme du $\K$-espace vectoriel $\K[X]$. A la réflexion, on aurait en fait pu le caractériser comme l'unique endomorphisme $D$ de $\K[X]$ vérifiant $D(X^0)=0_{\K[X]}$ et $\forall n\in\N\st,\ D(X^n)=nX^{n-1}$
\end{exa}
\begin{exa}
	Si on souhaite dériver $P=\displaystyle\sum_{k=0}^{n}a_kX^k$, la définition n'est en théorie pas applicable. Mais en utilisant la linéarité de $D$, on obtient tout de même : $$P'=\sum_{k=1}^{n}a_kkX^{k-1}$$
\end{exa}
\begin{prop}[Produit]
	Soit $(P,Q)\in\K[X]^2$. On a : $$(P\times Q)'=P'\times Q+P\times Q'$$
\end{prop}
\begin{prop}[Composée]
	Soit $(P,Q)\in\K[X]^2$. On a : $$(P\circ Q)'=(P'\circ Q)\times Q'$$
\end{prop}
\begin{prop}
	Soit $P\in\K[X]$ non constant. Alors : $$\deg(P')=\deg(P)-1$$
\end{prop}
\begin{rmq}
	Si $P$ est constant, on a $\deg(P')=-\infty$. Ainsi, dans tous les cas, on a : $$\deg(P')\leqslant\deg(P)-1$$ On peut aller plus loin : $P$ est constant si, et seulement si, $P'=0$
\end{rmq}
\begin{dfn}[$n$-ème dérivée formelle]
	Soit $P\in\K[X]$. Sa $n$-ème dérivée formelle est définie par récurrence : on pose $P^{(0)}=P$ et $\forall n\in\N,\ P^{(n+1)}=(P^{(n)})'$
\end{dfn}
\begin{exa}
	Soit $d\in\N$. On a $$\forall n\in\llbracket0,d\rrbracket,\ (X^d)^{(n)}=\frac{d!}{(d-n)!}X^{d-n}$$ et $$\forall n>d,\ (X^d)^{(n)}=0$$
\end{exa}
\begin{exa}
	Soit $P\in\K[X]$. On a $$\forall (m,n)\in\N^2,\ (P^{(m)})^{(n)}=P^{(m+n)}$$
\end{exa}
\begin{exa}
	La famille $\displaystyle\left(\frac{X^0}{0!},\ldots,\frac{X^d}{d!}\right)$ est une base de $\K_d[X]$. Dans cette base, la matrice de l'endomorphisme dérivation est une matrice comportant des $0$ partout sauf sur la surdiagonale qui ne comporte que des $1$.
\end{exa}
\begin{prop}[généralisation de la propriété précédente]
	Soit $P\in\K[X]$ et $n\in\N$. Si $\deg(P)\geqslant n$, alors $\deg(P^{(n)})=\deg(P)-n$
\end{prop}
\begin{rmq}
	Si $\deg(P)<n$ alors on vérifie que $P^{(n)}=0$. Ainsi, dans tous les cas, on a toujours : $$\deg(P^{(n)})\leqslant\deg(P)-n$$
\end{rmq}
\begin{rmq}
	On a $$\forall n\in\N,\ \forall P\in\K[X],\ \deg(P)\leqslant n \Longleftrightarrow P^{(n+1)}=0$$ On en déduit que $K_n[X]$ est un sev en tant que noyau de $D^{n+1}$.
\end{rmq}
\begin{exe}
	Dans le cas où $\deg(P)\geqslant n$, calculer $CD(P^{(n)})$ en fonction de $CD(P)$ et de $\deg(P)$
\end{exe}
\begin{prop}[Combinaison linéaire]
	Soit $n\in\N$, $(\lambda,\mu)\in\K^2$ et $(P,Q)\in\K[X]^2$. On a : $$(\lambda P+\mu Q)^{(n)}=\lambda P{(n)}+\mu Q^{(n)}$$
\end{prop}
\begin{prop}[Produit, formule de Leibniz]
	Soit $n\in\N$ et $(P,Q)\in\K[X]^2$. On a la formule de Leibniz : $$(PQ)^{(n)}=\sum_{k=0}^{n}\binom{n}{k}P^{(k)}Q^{(n-k)}$$
\end{prop}
\begin{lem}
	Soit $P\in\K[X]\setminus\left\{0_{\K[X]}\right\}$, $x\in\K$ et $\alpha\in\N$. On a : $$m_P(x)=\alpha \Longleftrightarrow\exists Q\in\K[X],\ P=(X-x)^\alpha Q \ \ \text{et} \ \ Q(x)\ne 0$$
\end{lem}
\begin{lem}
	Soit $P\in\K[X]\setminus\left\{0_{\K[X]}\right\}$, $x\in\K$ et $\alpha\in\N\st$. Supposons que $m_P(x)=\alpha$. Alors : 
	\begin{enumerate}
		\item $P'\ne0_{\K[X]}$
		\item $m_{P'}(x)=\alpha-1$
	\end{enumerate}
\end{lem}
\begin{rmq}
	Toute racine de multiplicité $\geqslant2$ de $P$ est donc racine de $P'$. Ainsi, pour rendre un polynôme $P$ à racines simples de force, il suffit de lui substituer $\displaystyle \frac{P}{P\wedge P'}$.
\end{rmq}
\begin{thm}[CNS de multiplicité]
	Soit $P\in\K[X]\setminus\left\{0_{\K[X]}\right\}$, $x\in\K$ et $\alpha\in\N$. On a : $$m_P(x)=\alpha \Longleftrightarrow (\forall k\in\llbracket0,\alpha-1\rrbracket,\ P^{(k)}(x)=0) \ \ \text{et} \ \ P^{(\alpha)}(x)\ne 0$$
\end{thm}
\begin{rmq}
	Au-delà de $\alpha$, on ne peut rien dire ! Même si APCR cela sera toujours nul, on ne sait rien de ce qu'il se passe entre $_alpha$ et ce rang.
\end{rmq}
\section{Bases de $\K_n[X]$}
\subsection{Base de Taylor}
\begin{thm}
	Soit $(P_i)_{i\in I}$ une famille de polynômes de $\K[X]$ \textbf{tous non nuls} ayant des degrés deux à deux distincts. Alors, la famille $(P_i)_{i\in I}$ est \textbf{libre}.
\end{thm}
\begin{cor}[Familles échelonnées en degré]
	On a les deux résultats suivants :
	\begin{enumerate}
		\item Soit $n\in\N$ et $(P_k)_{0\leqslant k\leqslant n}\in\K[X]^{n+1}$ telle que $\forall k \in\llbracket0,n\rrbracket, \ \deg(P_k)=k$. Alors, $n\in\N$ et $(P_k)_{0\leqslant k\leqslant n}$ est une base de $\K_n[X]$.
		\item Soit $(P_k)_{k\in\N}\in\K[X]^\N$ telle que $\forall k\in\N,\ \deg(P_k)=k$. Alors $(P_k)_{k\in\N}$ est une base de $\K[X]$.
	\end{enumerate}
\end{cor}
\begin{thm}
	Soit $a\in\K$ et $n\in\N$. D'après le corollaire précédent, $((X-a)^k)_{0\leqslant k\leqslant n}$ est une base de $\K_n[X]$. Soit $P\in\K_n[X]$. Sur cette base, $P$ se décompose en : $$P(X)=\sum_{k=0}^{n}\frac{P^{(k)}(a)}{k!}(X-a)^k$$ Pa substitution, on a aussi : $$P(X+a)=\sum_{k=0}^{n}\frac{P^{(k)}(a)}{k!}X^k$$
\end{thm}
\begin{rmq}
	Par $k!$, on entend $(1_\Z.1_\K)\times\ldots\times(k_\Z.1_\K)$. On travaille donc en caractéristique nulle.
\end{rmq}
\begin{rmq}
	Puisque toutes les dérivées suivantes sont nulles, on peut en fait écrire :$$P(X)=\sum_{k\in\N}\frac{P^{(k)}(a)}{k!}(X-a)^k$$
\end{rmq}
\subsection{Base de Lagrange}
Dans tous le paragraphe, on fixe $n\in\N$ et $x_0,\ldots,x_n$ $n+1$ scalaires deux à deux distincts.
\begin{dfn}[Polynômes interpolateurs de Lagrange]
	Les \textbf{polynômes interpolateurs de Lagrange associés aux $x_k$} sont les polynômes $L_0,\ldots,L_n$ définis par : $$\forall k\in\llbracket0,n\rrbracket,\ L_k=\prod_{\substack{0\leqslant i\leqslant n \\ i\ne k}}\frac{X-x_i}{x_k-x_i}$$
\end{dfn}
\begin{prop}[Propriété fondamentale des polynômes interpolateurs de Lagrange]
	Les polynômes interpolateurs de Lagrange sont des éléments de $\K_n[X]$ et vérifient la propriété fondamentale : $$\forall (k,j)\in\llbracket0,n\rrbracket^2,\ L_k(x_j)=\delta_{k,j}$$
\end{prop}
\begin{thm}
	La famille $(L_k)_{0\leqslant k\leqslant n}$ est une base de $K_n[X]$. Soit $P\in\K_n[X]$. Sur cette base, $P$ se décompose en : $$P=\sum_{k=0}^{n}P(x_k)L_k$$
\end{thm}
\begin{rmq}
	Soit $k\in\llbracket0,n\rrbracket$. A quoi est égale la $k$-ème forme coordonnée $L_k\st$ associée à la base $(L_k)_{0\leqslant k\leqslant n}$ ? Pour tout $P\in\K_n[X]$, on a $$L_k\st(P)=P(x_k)$$ donc en fait, $L_k\st$ est le morphisme d'évaluation $P\mapsto P(x_k)$.
\end{rmq}
\begin{cor}
	Soit $(y_0,\ldots,y_n)\in\K^{n+1}$ quelconque. Alors : $$\exists!P\in\N_n[X],\ \forall k\in\llbracket0,n\rrbracket,\ P(x_k)=y_k$$ Ce polynôme est donné par : $$P=\sum_{k=0}^{n}y_kL_k$$
\end{cor}
\subsection{Application à l'interpolation}
On veut approximer une fonction $f:I\rightarrow\K$. parfois, le calcul de $f$ est tellement coûteux qu'on préfère remplacer $f$ par une approximation polynomiale. Généralement, on demande à ce polynôme de passer par des points obligés appelés \textbf{points d'interpolation} qui sont deux à deux distincts. D'après le paragraphe précédent, cela est toujours possible.
\par Formellement, posons $\forall k\in\llbracket0,n\rrbracket, y_k=f(x_k)$. Puis on pose $$
\begin{array}[t]{lrcl}
	\varphi : & \K[X] & \longrightarrow & \K^{n+1} \\
	& P & \longmapsto & \begin{pmatrix}
		P(x_0) \\ \vdots \\ P(x_n)
	\end{pmatrix}
\end{array}
$$ Le but est donc formellement de résoudre l'équation $$\varphi(P)=\begin{pmatrix}
	y_0 \\ \vdots \\ y_n
\end{pmatrix}$$
Or $\varphi$ est linéaire et on sait qu'il existe au moins une solution $$P_{\text{fond}}=\sum_{k=0}^{n}y_kL_k$$ Ainsi, l'ensemble des solutions vaut $P_{\text{fond}}+\ker(\varphi)$ : c'est un sous-espace affine dirigé par $\ker(\varphi)$. Or, en notant $Q=\displaystyle\prod_{i=0}^{n}(X-x_k)$, on obtient que $P\in\K_n[X]$ appartient à $\ker(\varphi)$ si, et seulement si, $Q\mid P$. \par En somme, $P$ interpole correctement $f$ si, et seulement si, $$P\equiv P_{\text{fond}} \ [Q]$$
\section{Polynômes irréductibles}
\subsection{Cas général}
\begin{dfn}[Polynôme irréductible]
	Soit $P\in\K[X]$ \textbf{non constant}. On dit que $P$ est \textbf{irréductible} lorsque ses seuls diviseurs sont $1_{\K[X]}$, $P$ et leurs associés. Sinon, $P$ est dit \textbf{réductible}.
\end{dfn}
\begin{prop}
	Soit $P$ irréductible et $Q\in\K[X]$.
	\begin{enumerate}
		\item Si $P\nmid Q$, alors $P\wedge Q=1$.
		\item Supposons que $Q$ est irréductible. Si $P\nsim Q$, alors $P\wedge Q=1$.
	\end{enumerate}
\end{prop}
\begin{prop}[Lemme d'Euclide]
	Soit $P$ irréductible, $n\in\N\st$ et $A_1,\ldots,A_n\in\K[X]$. Si $P\mid A_1\times\ldots\times A_n$, alors $$\exists i\in\llbracket1,n\rrbracket,\ P\mid A_i$$
\end{prop}
\begin{prop}
	Soit $P\in\K[X]$ non constant. Posons $n=\deg(P)\in\N\st$. Les quatre assertion suivantes sont équivalentes :
	\begin{enumerate}
		\item $P$ est réductible
		\item $\exists(A,B)\in\K[X]^2,\ P=AB \ \ \text{et} \ \ \deg(A),\deg(B)>0$
		\item $\exists(A,B)\in\K[X]^2,\ P=AB \ \ \text{et} \ \ \deg(A),\deg(B)<n$
		\item $\exists(A,B)\in\K[X]^2,\ P=AB \ \ \text{et} \ \ 0<\deg(A),\deg(B)<0$
	\end{enumerate}
\end{prop}
\begin{thm}[Lien entre racines et irréductibilité]
	Soit $P\in\K[X]$ non constant. \par Supposons que $P$ est irréductible. Déjà, tout polynôme de degré $1$ est irréductible. En revanche, dès que $\deg(P)\geqslant2$, alors $P$ ne possède pas de racine. \par Réciproquement, supposons que $P$ ne possède pas de racine. Déjà, $P$ n'est pas de degré $1$. Supposons que $P$ est de degré $2$ ou $3$. En raisonnant par l'absurde et pour une question de degré, $P$ ne peut pas être réductible (sans quoi un polynôme de degré $1$ le diviserait et il possèderait alors une racine).  Dès que $\deg(P)\geqslant4$, tout peut arriver. Par exemple, dans $\R[X]$, $X^4+2X^2+1=(X^2+1)^2$ est réductible mais ne possède pas de racine. En revanche, dans $\Q[X]$, $X^4+1$ n'a pas de racine et est irréductible.
\end{thm}
Jusqu'à la fin du chapitre, on notera de façon \textbf{non standard} $\mathcal{P}$ l'ensemble des polynômes irréductibles unitaires.
\begin{thm}[Existence de la DFI]
	Soit $A\in\K[X]\setminus\left\{0\right\}$. Alors $A$ admet une décomposition en facteurs irréductibles (DFI), c'est-à-dire que $$A=\lambda \prod_{P\in I}P^{\alpha_P}$$ où $\lambda\in\K\st$, $I$ est une partie finie de $\mathcal{P}$ et $\forall P\in I,\ \alpha_P\in\N\st$.
\end{thm}
\begin{exa}
	On a nécessairement $\lambda=CD(A)$ puisque tous les $P\in I$ sont unitaires.
\end{exa}
\begin{exa}
	La forme suivante est beaucoup plus opérationnelle : $$A=\lambda P_1^{\alpha_1}\times\ldots\times P_n^{\alpha_n}$$ avec $\lambda\in\K\st$, les $P_i\in\mathcal{P}$ deux à deux distincts et les $\alpha_i$ dans $\N\st$.
\end{exa}
\begin{dfn}[Valuation $P$-adique]
	Soit $P\in\mathcal{P}$ et $A\in\K[X]\setminus\left\{0\right\}$. La \textbf{valuation $P$-adique de $A$} est l'entier défini par : $$v_P(A)=\max\left\{k\in\N : P^k\mid A\right\}$$
\end{dfn}
\begin{exa}
	Soit $x\in\K$. On a $X-x\in\mathcal{P}$. Puis $v_{X-x}(A)=m_A(x)$
\end{exa}
\begin{rmq}
	On a la même phénomène que dans $\Z$ avec la valuation $p$-adique et que dans $\K[X]$ avec la multiplicité. Soit $P\in\mathcal{P}$. On a $$\forall k\in\N,\ P^k \mid A \Longleftrightarrow k\leqslant v_P(A)$$ d'où on tire $$\left\{k\in\N : P^k\mid A\right\}=\llbracket0,v_P(A)\rrbracket$$
\end{rmq}
\begin{exa}
	Soit $A$ et $B$ des polynômes non nuls et $P$ irréductible unitaire. Si $A\sim B$, alors $v_P(A)=v_P(B)$
\end{exa}
\begin{exa}
	Soit $\alpha\in\N$ et $P$ irréductible unitaire. On a $v_P(P^\alpha)=\alpha$.
\end{exa}
\begin{lem}
	Soit $A$ un polynôme non nul, $P$ irréductible unitaire et $\alpha\in\N$. On a $$v_P(A)=\alpha \Longleftrightarrow\exists A_1\in\K[X],\ A=P^\alpha A_1 \ \ \text{et} \ \ P\wedge A_1=1$$
\end{lem}
\begin{cor}
	Soit $A$ et $B$ des polynômes non nuls et $P$ un polynôme irréductible unitaire. On a : $$v_P(AB)=v_P(A)+v_P(B)$$
\end{cor}
\begin{rmq}
	Par une récurrence immédiate, on généralise à un produit fini de polynômes non nuls.
\end{rmq}
\begin{cor}
	Soit $A\in\K[X]\setminus\left\{0\right\}$ dont on considère une DFI $A=\lambda\displaystyle\prod_{P\in I}P^{\alpha_P}$. Soit $Q\in\mathcal{P}$. On a : \begin{enumerate}
		\item Si $Q\in I$, alors $v_Q(A)=\alpha_Q$
		\item Si $Q\notin I$, alors $v_Q(A)=0$
	\end{enumerate}
\end{cor}
\begin{cor}[Unicité de la DFI]
	Dans $\K[X]\setminus\left\{0\right\}$, la DFI est unique.
\end{cor}
\begin{rmq}
	Un polynôme irréductible unitaire apparaît dans la DFI de $A$ si, et seulement si, il divise $A$.
\end{rmq}
\begin{prop}
	$A$ et $B$ non nuls sont premiers entre eux si, et seulement si, leurs DFI sont disjointes.
\end{prop}
\begin{prop}
	Soit $A$ et $B$ des polynômes non nuls. On a : 
	\begin{enumerate}
		\item $A\mid B \Longleftrightarrow \forall P\in\mathcal{P},\ v_P(A)\leqslant v_P(B)$
		\item Supposons de plus que $A$ et $B$ sont unitaires. Alors $A=B \Longleftrightarrow \forall P\in\mathcal{P},\ v_P(A)=v_P(B)$
	\end{enumerate}
\end{prop}
\begin{prop}[Valuation $P$-adique du PPCM et du PGCD]
	Soit $A$ et $B$ deux polynômes non nuls et $P$ un polynôme irréductible unitaire. On a : $$v_P(A\wedge b)=\min(v_P(A),v_P(B)) \ \ \text{et} \ \ v_P(A\vee B)=\max(v_P(A),v_P(B))$$
\end{prop}
\begin{cor}
	Soit $A$ et $B$ des polynômes \textbf{unitaires}. On a : $$AB=(A\wedge B)(A\vee B)$$
\end{cor}
\begin{dfn}
	Soit $A\in\K[X]\setminus\left\{0\right\}$. $A$ est dit \textbf{scindé} lorsque, dans sa DFI, tous les facteurs sont de degré $1$. En pratique, on pourra donc écrire : $$A=\lambda\prod_{i=1}^{r}(X-x_i)^{\alpha_i}$$ où $\lambda\in\K\st$, $r\in\N$, les $x_i$ sont deux à deux distincts et les $\alpha_i$ sont dans $\N\st$.
\end{dfn}
\begin{exa}
	Dans ce cas, on a $\forall i\in\llbracket1,r\rrbracket,\ \alpha_i=v_{X-x_i}=m_A(x_i)$
\end{exa}
\begin{rmq}
	Dans la suite, pour un polynôme scindé, on pourra écrire $A=\lambda(X-x_1)\ldots(X-x_n)$ où les $x_i$ ne seront plus nécessairement supposés deux à deux distincts.
\end{rmq}
\begin{thm}[Formules de Viète]
	Soit $A=\lambda(X-x_1)\times\ldots\times(X-x_n)$. Écrivons-le sous la forme $A=\displaystyle\sum_{k=0}^{n}a_kX^k$. Alors, on a : $$\forall k\in\llbracket0,n\rrbracket, \ \sum_{1\leqslant i_1<\ldots<i_k\leqslant n}x_{i_1}\times\ldots\times x_{i_k}=(-1)^k\frac{a_{n-k}}{a_n}$$
\end{thm}
\begin{dfn}[Fonctions symétriques élémentaires, HP]
	$$(x_1,\ldots,x_n)\mapsto \displaystyle\sum_{1\leqslant i_1<\ldots<i_k\leqslant n}x_{i_1}\times\ldots\times x_{i_k}$$ est appelée \textbf{$n$-ème fonction symétrique élémentaire d'ordre $k$}. On la note plus souvent $\sigma_{n,k}$.  Ces fonctions vérifient une propriété fondamentale de décomposition des polynômes à $n$ indéterminées.
\end{dfn}
\subsection{Cas de $\R[X]$ et de $\C[X]$}
\begin{thm}[Théorème de D'Alembert-Gauss]
	Dans $\C[X]$, tout polynôme non constant admet au moins une racine. On dit que $\C$ est \textbf{algébriquement clos}.
\end{thm}
\begin{prop}
	Soit $P\in\R[X]\setminus\left\{0\right\}$ et $z\in\C$. Alors : $$m_P(z)=m_P(\overline{z})$$
\end{prop}
\begin{exa}
	En particulier, $z$ est racine de $P$ à coefficients réels non nul si, et seulement si, $\overline{z}$ est racine de $P$.
\end{exa}
\begin{lem}
	Soit $(A,B)\in\R[X]^2$ tel que $B$ divise $A$ dans $\C[X]$? Alors, $B$ divise $A$ dans $\R[X]$.
\end{lem}
\begin{proof}
	Ce théorème consiste à effectuer la division euclidienne dans $\R[X]$ puis dans $\C[X]$ avant de plonger celle de $\R[X]$ dans $\C[X]$ puis d'invoquer l'unicité de la division euclidienne dans $\C[X]$.
\end{proof}
\begin{rmq}
	Plus généralement, ce lemme fonctionne dans deux corps quelconques $\mathbb{L}$ et $\K$ dès que $\mathbb{L}$ est une extension de corps de $\K$.
\end{rmq}
\begin{thm}[Classification des polynômes irréductibles à coefficients complexes]
	Dans $\C[X]$, les polynômes irréductibles sont exactement les polynômes de degré $1$.
\end{thm}
\begin{rmq}
	Contrairement aux autres corps où on ne dispose que d'une inégalité, si on prend $P\in\C[X]\setminus\left\{0\right\}$ de degré $n\in\N$, alors on a : $$\sum_{x\in\mathcal{R}(P)}m_P(x)=n$$
\end{rmq}
\begin{thm}[Classification des polynômes irréductibles à coefficients réels]
	Dans $\R[X]$, les polynômes irréductibles sont exactement les polynômes de degré $1$ et les trinômes du second degré de discriminant strictement négatif.
\end{thm}
\begin{rmq}
	Pour factoriser un polynôme dans $\R[X]$, on peut le factoriser dans $\C[X]$ puis regrouper les racines complexes conjuguées.
\end{rmq}
\begin{cor}
	Soit $A$ et $B$ des polynômes à coefficients complexes tous deux différents du polynôme nul. On a : $$B\mid A \Longleftrightarrow\forall z\in \C,\ m_B(z)\leqslant m_A(z)$$
\end{cor}
\section{Compléments, HP}
\subsection{Relations coefficients-racines}
Voici la démonstration formelle du résultat évoqué en cours.

\begin{thm}
	Soit $\displaystyle P=\lambda(X-\lambda_{1})(X-\lambda_{2})...(X-\lambda_{n})=\sum_{k=0}^{n}a_{k}X^{k}$ un polynôme scindé. Alors $$\forall k\in\llbracket0,n\rrbracket,\ \sum_{1\le i_{1}<...<i_{k}\le n}\lambda_{i_{1}}...\lambda_{i_{k}}=(-1)^{k}\frac{a_{n-k}}{a_{n}}$$
\end{thm}
\begin{proof}
	Pour $i\in\llbracket1,n\rrbracket$ posons $t_{i,0}=X$ et $t_{i,1}=-\lambda_{i\cdot}$.
	\begin{itemize}
		\item Les \(t_{i,j}\) sont à valeurs dans l'anneau \(\mathbb{K}[X]\) donc la formule de distributivité généralisée permet d'écrire $$P=\lambda\prod_{i=1}^{n}\left(\sum_{j\in\{0,1\}}t_{i,j}\right)=\lambda\sum_{(j_{1},...,j_{n})\in\{0,1\}^{n}}t_{1,j_{1}}...t_{n,j_{n}}$$ Pour chaque n-uplet \((j_{1},...,j_{n})\) considérons son support $S$. L'application $$\begin{array}{ccc}
			\{0,1\}^{n}&\rightarrow&\mathcal{P}(\llbracket1,n\rrbracket)\\
			(j_{1},...,j_{n})&\mapsto& S
		\end{array}$$
		est bijective. Pour s'en convaincre, le mieux est d'exhiber sa bijection réciproque :
		$$
		\begin{array}{ccc}
			\mathcal{P}(\llbracket1,n\rrbracket)&\rightarrow&\{0,1\}^{n}\\
			S&\mapsto&(\mathds{1}_{s}(1),...,\mathds{1}_{s}(n))
		\end{array}$$
		\item Par changement de variable, puis par associativité, on en déduit
		$$P=\lambda\sum_{S\subset\{1,n\}}\left(\prod_{i\in S}(-\lambda_{i})\prod_{i\in S}X\right)=\lambda\sum_{k=0}^{n}(-1)^{k}\left[\sum_{S\subset\{1,n\}}\left(\prod_{i\in S}\lambda_{i}\right)\right]X^{n-k}$$
		Pour $k=0$, on obtient $a_{n}=\lambda$ Plus généralement, fixons $k\in\llbracket0,n\rrbracket$. Tout revient à montrer que
		$$\sum_{1\le i_{1}<...<i_{k}\le n}\lambda_{i_{1}}...\lambda_{i_{k}}=\sum_{S\subset\llbracket1,n\rrbracket}\left(\prod_{i\in S}\lambda_{i}\right)$$
		\item Posons $A=\{(i_{1},...,i_{k})\in\llbracket1,n\rrbracket^{k}\ |\ i_{1}<...<i_{k}\}$ et $B=\{S\subset\llbracket1,n\rrbracket:|S|=k\}$. L'application $$
		\begin{array}{ccccc}
			\varphi&:&A&\rightarrow&B\\
			&&(i_{1},...,i_{k})&\mapsto&\{i_{1},...,i_{k}\}
		\end{array}$$ est bien définie. En effet, les $\{i_{j}\}$ sont deux à deux disjoints donc le cardinal de leur union vaut bien la somme de leurs cardinaux, c'est-à-dire k. Ensuite elle est bijective. Là aussi, pour s'en convaincre, le mieux est encore d'exhiber sa bijection réciproque; celle-ci consiste à se donner un ensemble $S$ de cardinal $k$, à ordonner ses éléments $i_{1}<...<i_{k}$, puis à retourner le $k$-uplet $(i_{1},...,i_{k})$. Par un dernier changement de variable, on en
		déduit le résultat final.
	\end{itemize}
Ainsi la preuve est achevée.
\end{proof}
\subsection{Théorème de d'Alembert-Gauss}
\begin{thm}
	Tout polynôme non constant de $\mathbb{C}[X]$ admet une racine.
\end{thm}
\begin{proof} Il s'agit d'un des théorèmes les plus difficiles de l'année, et sa preuve est très clairement hors-programme. Nous la donnons pour les plus valeureux. Il s'agit peu ou prou de la démonstration originale de Cauchy. Sans perte de généralité, on suppose que $P$ est unitaire.
\begin{enumerate}
	\item Si on écrit $P(z)=z^{n}+a_{n-1}z^{n-1}+...+a_{0}$ avec $n\geqslant1$, alors
	$$\forall z\neq0,\ |a_{n-1}z^{n-1}+...+a_{0}|=|z|^{n}\cdot\left\lvert\frac{a_{n-1}}{z}+...+\frac{a_{0}}{z^{n}}\right\rvert\leq|z|^{n}(\frac{|a_{n-1}|}{|z|}+...+\frac{|a_{0}|}{|z|^{n}})$$
	Puisque$\displaystyle\frac{|a_{n-1}|}{R}+...+\frac{|a_{0}|}{R^{n}}\underset{R\rightarrow+\infty}{\longrightarrow}0,$
	on peut choisir $R>0$ tel que
	$$\forall z\in\mathbb{C},|z|\geqslant R\Rightarrow\frac{|a_{n-1}|}{|z|}+...+\frac{|a_{0}|}{|z|^{n}}\leq\frac{1}{2}$$
	Soit $z\in\mathbb{C}.$ Par la deuxième inégalité triangulaire, on a
	$$|z|\geqslant R\Rightarrow|P(z)|\geq|z^{n}|-|a_{n-1}z^{n-1}+...+a_{0}|\geq\frac{|z|^{n}}{2}\geq\frac{R^{n}}{2}$$
	Or on a $\displaystyle\frac{|a_{0}|}{R^{n}}\leq\frac{1}{2}$, d'où finalement, $|z|\geqslant R\Rightarrow|P(z)|\geq|a_{0}|=|P(0)|$.
	\item Considérons alors le disque fermé $D=\{z\in\mathbb{C}:|z|\leqslant R\}.$ L'ensemble $\{|P(z)||z\in D\}$ étant non vide et minoré, il admet une borne inf $m$ et il existe une suite à valeurs dans $D$ telle que $|P(z_{k})|\rightarrow m$. Puisque $(z_k)$ est bornée, quitte à considérer une suite extraite, on peut supposer que $z_{k}\rightarrow\tilde{z}.$ Par passage au module, on a alors $|z_{k}|\rightarrow|\overline{z}|$ puis par passage à la limite on a $|\tilde{z}|\leqslant R$. \\
	Soit $k\in\mathbb{N}.$ On a ensuite 
	\begin{align*}
		\big||P(\tilde{z})|-|P(z_{k})|\big| &\leq|P(\tilde{z})-P(z_{k})|\\
		&\leq|\tilde{z}^{n}-z_{k}^{n}|+|a_{n-1}|\cdot|\tilde{z}^{n-1}-z_{k}^{n-1}|+...+|a_{1}|\cdot|\tilde{z}-z_{k}|
	\end{align*}
	Si on peut prouver que chacun des $|\tilde{z}^{j}-z_{k}^{j}|$ tend vers 0 lorsque $k\rightarrow+\infty,$ on en déduira par passage à la limite que $\big||P(\vec{z})|-m\big|=0.$ Or on a
	$$|\tilde{z}^{j}-z_{k}^{j}|\leq|\tilde{z}-z_{k}|\cdot(|\tilde{z}|^{j-1}+|\tilde{z}|^{j-2}|z_{k}|+...+|z_{k}|^{j-1})\leq|\tilde{z}-z_{k}|\cdot jR^{j-1}$$ ce qui permet de conclure.
	\item Par le point $2.$, $|P(z)|$ admet un minimum $m$ sur $D$, qu'il atteint en $\tilde{z}$. Et d'après le point $1.$, ce minimum est en fait un minimum global sur $\C$ Supposons maintenant que $m>0$ et aboutissons à une contradiction.\\
	Considérons la fonction $z\mapsto P(\overline{z}+z).$ D'après la formule de Taylor on peut écrire
	$\forall z\in\mathbb{C},\ P(\overline{z}+z)=b_{n}z^{n}+...+b_{1}z+b_{0}$
	avec $|b_{0}|=|P(\tilde{z})|=m\neq0.$ Puisque P n'est pas constant, le polynôme $P(\overline{z}+X)$ n'est pas constant non plus, et il existe donc $k\in\llbracket1,n\rrbracket$ tel que $b_{k}\neq0$. Considérons $k$ minimal pour cette propriété. On a $$\forall z\in\mathbb{C},\ P(\tilde{z}+z)=b_{n}z^{n}+...+b_{k}z^{k}+b_{0}$$
	Soit alors $c$ une racine $k$-ème de $-\overline{b_{k}}b_{0}$. On a
	\begin{align*}
		\forall t>0, |P(\overline{z}+ct)| &\leqslant \left|\sum_{k<j\leqslant n}b_{j}c^{j}t^{j}\right| + \left|-\left|b_{k}\right|^{2}b_{0}t^{k}+b_{0}\right| \\
		&\leqslant t^{k+1}\sum_{0\leqslant j\leqslant n-k}|b_{j+k+1}c^{j+k+1}|\cdot t^{j} + m\left|1-\left|b_{k}\right|^{2}t^{k}\right|.
	\end{align*}
	Soit maintenant $t\in[0,1]$. Alors la somme est majorée par un certain $M>0$ fixé (notons que si elle est vide, il n'y a rien à faire). Quitte à diminuer encore $t$, supposons que $1-\left|b_{k}\right|^{2}t^{k}\geqslant0$, si bien que $$|P(\tilde{z}+ct)|\leqslant t^{k+1}M+m\left(1-\left|b_{k}\right|^{2}t^{k}\right)$$
	Quitte à diminuer une dernière fois $t$, supposons que $|t|M\leq\frac{1}{2}m\left|b_{k}\right|^{2}$, ce qui est possible car on a supposé $m>0$ (c'est le point-clé!). On a alors $$|P(\overline{z}+ct)|\leqslant m+t^{k}\left(tM-m\left|b_{k}\right|^{2}\right)\leqslant m-\frac{1}{2}m\left|b_{k}\right|^{2}t^{k}<m$$
	Contradiction.
\end{enumerate}
Ainsi la preuve est achevée.
\end{proof}
\subsection{Polynômes à coefficients dans $\Z$ : contenu et critère d'Eisenstein}
On présente ici la notion de contenu d'un polynôme à coefficients dans $\Z$ et de polynôme primitif. Cela nous sert à démontrer le célèbre \textbf{critère d'Eisenstein}, grand classique d'X/ENS. On étend enfin la définition du contenu aux polynômes à coefficients dans $\Q$. Ce travail peut être généralisé plus généralement avec un anneau factoriel et son corps des fractions.
\begin{dfn}
	On rappelle que l'ensemble des polynômes à coefficients dans $\Z$, noté $\Z[X]$, muni des lois usuelles, est un anneau intègre puisque $\Z$ est intègre.
\end{dfn}
\begin{prop}
	Les éléments inversibles de $\Z[X]$ sont les polynômes constants égaux à $-1$ et $1$.
\end{prop}
\begin{proof}
	Comme pour $\K[X]$, on raisonne d'abord sur le degré, mais il faut ensuite raisonner sur le coefficient dominant qui doit être une unité de $\Z$.
\end{proof}
\begin{dfn}[Contenu d'un polynôme à coefficients dans $\Z$]
	Soit $P$ un polynôme de $\Z[X]\setminus\left\{0\right\}$ qu'on écrit sous la forme $P=\displaystyle\sum_{k=0}^{n}a_kX^k$ avec $a_n\ne0$. On appelle \textbf{contenu de $P$}, et on note $c(P)$, le PGCD des coefficients de $P$ : $$c(P):=a_0\wedge\ldots\wedge a_n$$
\end{dfn}
\begin{dfn}[Polynôme primitif]
	On dit que $P\in\Z[X]\setminus\left\{0\right\}$ est \textbf{primitif} lorsque $c(P)=1$.
\end{dfn}
\begin{prop}
	On a $$\forall \lambda\in\Z\st,\ \forall P\in\Z[X]\setminus\left\{0\right\},\ c(\lambda P)=\lvert\lambda\rvert c(P)$$
\end{prop}
\begin{proof}
	C'est une simple réécriture de la distributivité du PGCD.
\end{proof}
\begin{lem}
	Soit $P\in\Z[X]\setminus\left\{0\right\}$. Alors le polynôme $\displaystyle\frac{P}{c(P)}$ est primitif.
\end{lem}
\begin{proof}
	On a $\displaystyle c(P)=c\left(c(P)\cdot\frac{P}{c(P)}\right)=c(P)c\left(\frac{P}{c(P)}\right)$ donc $\displaystyle c\left(\frac{P}{c(P)}\right)=1$.
\end{proof}
\begin{prop}
	Le produit de deux polynômes primitifs est primitif.
\end{prop}
\begin{proof}
	Soit $P\in\Z[X]\setminus\left\{0\right\}$ et $Q\in\Z[X]\setminus\left\{0\right\}$ tels que $c(P)=c(Q)=1$. Déjà, par intégrité de $\Z[X]$, $PQ\ne0$ donc $c(PQ)$ est bien défini. \par Raisonnons par l'absurde et supposons que $c(PQ)\ne1$. Alors $c(PQ)>1$ donc $c(PQ)$ possède un diviseur premier $p$. En passant dans $\Z/p\Z[X]$, on a $\overline{PQ}=\overline{P}\times\overline{Q}=\overline{0}$, donc puisque $\Z/p\Z$ est un corps, $\Z/pZ[X]$ est un anneau intègre et donc $\overline{P}=\overline{0}$ ou $\overline{Q}=\overline{Q}$. Ainsi, les coefficients de $P$ ou de $Q$ sont tous divisibles par $p$. Contradiction avec le fait que $P$ et $Q$ sont primitifs. Donc $PQ$ est primitif, \textit{ie} $c(PQ)=1$.
\end{proof}
\begin{cor}[Contenu d'un produit de polynôme à coefficients dans $\Z$]
	Soit $P\in\Z[X]\setminus\left\{0\right\}$ et $Q\in\Z[X]\setminus\left\{0\right\}$. On a $$c(PQ)=c(P)c(Q)$$
\end{cor}
\begin{proof}
	On a $$c(PQ)=c\left(c(P)c(Q)\cdot\frac{PQ}{c(P)c(Q))}\right)=c(P)c(Q)c\left(\frac{PQ}{c(P)c(Q)}\right)=c(P)c(Q)$$ puisque $\displaystyle\frac{PQ}{c(P)c(Q)}$ est primitif. 
\end{proof}
On démontre un dernier lemme avant d'arriver au critère d'Eisenstein.
\begin{lem}
	Soit $A\in\Z[X]$ non irréductible dans $\Q[X]$. Alors il existe $B$ et $C$ dans $\Z[X]$ tels que $\deg(B)<\deg(A)$, $\deg(C)<\deg(A)$ et $A=BC$.
\end{lem}
\begin{proof}
	On pose $A_1:=\displaystyle\frac{A}{c(A)}$. Alors $A_1$ n'est pas irréductible dans $\Q[X]$ non plus et il existe $B_1$ et $C_1$ dans $\Q[X]$ tels que $\deg(B_1)<\deg(A_1)$, $\deg(C_1)<\deg(A_1)$ et $A_1=B_1C_1$. Notons $\beta$ le produit des dénominateurs de $B_1$ et $\gamma$ le produit des dénominateurs de $C_1$. Alors $B_2:=\beta B_1\in\Z[X]$ et $C_2:=\gamma C_1\in\Z[X]$. Puis : $$\beta\gamma A_1=\beta\gamma B_1c_1=B_2C_2$$ En passant au contenus : $$\beta\gamma =c(\beta\gamma A_1)=c(B_2C_2)=c(B_2)c(C_2)$$ Puis $$A=c(A)A'=c(A)\frac{BC}{\beta\gamma}=\left(c(A)\frac{B_2}{c(B_2)}\right)\left(\frac{C_2}{C_2}\right)$$ Ainsi, $\displaystyle B:=c(A)\frac{B_2}{c(B_2)}$ et $\displaystyle C:=\frac{C_2}{c(C_2)}$ conviennent.
\end{proof}
\begin{thm}[Critère d'Eisenstein]
	Soit $A=\sum_{k=0}^{n}a_kX^k\in\Z[X]$ et $p$ un nombre premier. Supposons que $$\begin{cases}
		p \nmid a_n \\
		\forall k\in\llbracket0,n-1\rrbracket,\ p\mid a_k \\
		p^2\nmid a_0
	\end{cases}$$ alors $A$ est irréductible dans $\Q[X]$.
\end{thm}
\begin{proof}
	Raisonnons par l'absurde et supposons que $A$ ne soit pas irréductible dans $\Q[X]$. On peut alors fixer $B$ et $C$ dans $\Z[X]$ tels que $\deg(B),\deg(C)<\deg(A)$ et $A=BC$. Écrivons $B=\displaystyle\sum_{k=0}^{i}b_kX^k$ et $C=\displaystyle\sum_{k=0}^{j}c_kX^k$. \par En passant dans $\Z/p\Z[X]$, on a $\overline{A}=\overline{a_n}X^n$. Or, $a_n=b_ic_j$ donc $\overline{a_n}=\overline{b_i}\times\overline{c_j}$. Par intégrité de $\Z/p\Z$, $\overline{b_i}\ne\overline{0}$ et $\overline{c_j}\ne\overline{0}$ donc $\overline{B}$ et $\overline{C}$ sont de degrés respectifs $i$ et $j$ dans $\Z/p\Z$. On a donc : $$\overline{A}=\overline{b_i}X^i\times\overline{c_j}X^j=\overline{B}\times\overline{C}$$ Puisque $X$ est irréductible dans $\Z/p\Z[X]$, en comparant les DFI et les degrés, on a $\overline{B}=\overline{b_i}X^i$ et $\overline{C}=\overline{c_j}X^j$. Or, par hypothèse, on a nécessairement $i>0$ et $j>0$. Donc $p$ divise $b_0$ et $c_0$. Or, $a_0=b_0c_0$ donc $p^2$ divise $a_0$. Absurde d'après le troisième point du critère. Donc $A$ est irréductible dans $\Q[X]$.
\end{proof}
\begin{exa}[Application]
	Comme application, on montre qu'il existe des polynômes à coefficient dans $\Z$ d'un degré quelconque irréductibles dans $\Q[X]$. Il suffit de considérer $X^n-2$ et de lui appliquer le critère d'Eisenstein avec $p=2$. On peut utiliser ce fait et un autre résultat qui découle du critère pour montrer que $\R$ est un $\Q$-ev de dimension infinie.
\end{exa}
On peut étendre la définition du contenu aux polynômes à coefficients dans $\Q$ comme suit.
\begin{dfn}[Contenu d'un polynôme à coefficients dans $\Q$]
	Soit $P\in\Q[X]\setminus\left\{0\right\}$. Il existe alors au moins un $\lambda\in\N\st$ tel que $\lambda P\in\Z[X]$. On pose alors $$c(P):=\frac{c(\lambda P)}{\lambda}$$
\end{dfn}
\begin{proof}
	Pour $\lambda$, il suffit de prendre le PPCM des dénominateurs des coefficients de $P$ ou la valeur absolue du produit des dénominateurs des coefficients de $P$. \par Ensuite, montrons que la définition de $c(P)$ ne dépend pas du choix de $\lambda$. Soit $\lambda\in\N\st$ et $\mu\in\N\st$ tels que $\lambda P\in\Z[X]$ et $\mu P\in\Z[X]$. Alors $\lambda\mu P\in\Z[X]$ et on a $c(\lambda\mu P)=\lambda c(\mu P)=\mu c(\lambda P)$ donc $\displaystyle \frac{c(\lambda P)}{\lambda}=\frac{c(\mu P)}{\mu}$. \par Enfin vérifions que la définition prolongée celle donnée dans $\Z[X]$. Supposons que $P\in\Z[X]\setminus\left\{0\right\}$. Alors $c_\Q(P)=\frac{c_\Z(1P)}{1}=c_\Z(P)$. La définition prolonge bien celle de $\Z$.
\end{proof}
\begin{cor}[Contenu d'un produit de polynôme à coefficients dans $\Q$]
	Soit $P\in\Q[X]\setminus\left\{0\right\}$ et $Q\in\Q[X]\setminus\left\{0\right\}$. On a $$c(PQ)=c(P)c(Q)$$
\end{cor}
\begin{proof}
	Fixons $\lambda\in\N\st$ tel que $\lambda P\in\Z[X]$ et $\mu\in\N\st$ tel que $\mu Q\in\Z[X]$. Alors $\lambda\mu PQ\in\Z[X]$ et on a : $$c(PQ)=\frac{c(\lambda\mu PQ)}{\lambda\mu}=\frac{c(\lambda P)c(\mu Q)}{\lambda \mu}=c(P)c(Q)$$ d'après le cas de $\Z[X]$.
\end{proof}
\begin{prop}
	Soit $P\in\Q[X]\setminus\left\{0\right\}$. Si $c(P)=1$, alors $P\in\Z[X]\setminus\left\{0\right\}$.
\end{prop}
\begin{proof}
	Écrivons $P=\displaystyle\sum_{k=0}^{n}a_kX^k$ avec $a_n\ne0$. Soit $\lambda\in\N\st$ tel que $\lambda P\in\Z[X]$. On a alors $c(\lambda P)=\lambda$. En particulier, $\lambda$ divise chacun des entiers $\lambda a_k$.  Soit $k\in\llbracket0,n\rrbracket$. On peut donc écrire $\lambda a_k=\lambda q_k$ avec $q_k\in\Z$. Comme $\lambda\ne0$, on en déduit que $a_k=q_k\in\Z$.
\end{proof}
\begin{rmq}[Extension de la définition de primitif]
	On peut alors étendre la notion de polynôme primitif à $\Q$ sans problème, puisqu'alors tout polynôme primitif est nécessairement à coefficients dans $\Z$.
\end{rmq}
\begin{rmq}[Contenu et anneau factoriel]
	Plus généralement, on peut réaliser le même travail de définition du contenu dans n'importe quel anneau factoriel ($\Z$ est factoriel car principal) et son corps des fractions ($\Q$ ici). En revanche, il sera alors défini "à un inversible près" (là où on impose la positivité du PGCD dans $\Z$). Pour plus de détail sur les anneaux factoriels, se référer au chapitre "Compléments sur les anneaux commutatifs".
\end{rmq}

\newpage
\chapter{Fractions rationnelles}
Dans tout le chapitre, on fixe $\K$ un corps. En pratique, cela sera souvent $\R$ ou $\C$, voire $\Q$.
\section{Premières définitions}
\begin{dfn}[Corps des fractions rationnelles]
	$(\K[X],+,\times)$ étant un anneau intègre, on peut considérer son corps des fractions. Il sera noté $K(X)$ et, conformément à l'usage, ses éléments seront notés $\displaystyle \frac{P}{Q}$ (où $(P,Q)\in\K[X]\times(\K[X]\setminus\left\{0\right\})$). Pour mémoire, $\displaystyle \frac{P}{Q}$ désigne la classe d'équivalence du couple $(P,Q)$.
\end{dfn}
\begin{rmq} On vérifie que la structure d'espace vectoriel de $\K[X]$ se prolonge à $\K(X)$. \end{rmq}
\begin{dfn}[Degré]
	Soit $F\in\K[X]$. On fixe $(P,Q)\in\K[X]\times(\K[X]\setminus\left\{0\right\})$ tel que $\displaystyle F=\frac{P}{Q}$ puis on définit le degré de $F$ par $$\deg(F)=\deg(P)-\deg(Q)$$ Cette définition est indépendante du choix du représentant de $F$.
\end{dfn}
\begin{rmq} Attention, contrairement au cas de $\K[X]$, les fractions rationnelles de degré 1 ne sont pas nécessairement constantes ! Par exemple, on pourra considérer $\displaystyle \frac{X+1}{X}$. \end{rmq}
\begin{prop}[Propriétés du degré]
	Soit $F_1,F_2\in\K(X)$. On a :
	\begin{enumerate}
		\item $\deg(F_1+F_2)\leqslant\max(\deg(F_1),\deg(F_2))$
		\item $\deg(F_1F_2)=\deg(F_1)+\deg(F_2)$
	\end{enumerate}
	La propriété sur la somme reste vraie pour une combinaison linéaire de $F_1$ et $F_2$.
\end{prop}
\begin{dfn}
	Soit $F\in\K(X)$. Il existe un unique couple $(P,Q)\in\K[X]\times(\K[X]\setminus\left\{0\right\})$ tel que $\displaystyle F=\frac{P}{Q}$, $P\wedge Q=1$ et $Q$ est unitaire. $\displaystyle \frac{P}{Q}$ s'appelle alors \textbf{l'écriture irréductible} de $F$.
\end{dfn}
\begin{dfn}[Racines, pôles]
	Soit $F\in\K(X)$ qu'on écrit sous forme irréductible $\displaystyle \frac{P}{Q}$. On appelle \textbf{racine} de $F$ toute racine de $P$. On appelle \textbf{pôle} de $F$ toute racine de $Q$.
\end{dfn}
\begin{dfn}[Évaluation]
	Soit $F\in\K(X)$ qu'on écrit sous forme irréductible $\displaystyle \frac{P}{Q}$. On note $\mathcal{P}$ l'ensemble de ses pôles. Pour tout $x\in\K\setminus\mathcal{P}$, on pose : $$F(x)=\frac{P(x)}{Q(x)}$$
\end{dfn}
\begin{dfn}[Fonction rationnelle]
	La fonction $\begin{array}[t]{lrcl}
		\tilde{F}: & \K\setminus\mathcal{P} & \longrightarrow & \K \\
		& x & \longmapsto & F(x)
	\end{array}
	$ s'appelle la \textbf{fonction rationnelle} associée à $F$.
\end{dfn}
\begin{prop}
	Soit $F\in\K(X)$ et $\displaystyle\frac{A}{B}$ une écriture quelconque de $F$. Soit $x\in\K$. Si $B(x)\ne0_\K$, alors $x\notin\mathcal{P}$ donc $F(x)$ existe, et on a $$F(x)=\frac{A(x)}{B(x)}$$
\end{prop}
\begin{prop}[Évaluation et opérations algébriques]
	Du moment que $x$ n'est ni un pôle de $F_1$, ni un pôle de $F_2$, alors $x$ n'est pas un pôle de toute combinaison linéaire de $F_1$ et $F_2$ ni du produit $F_1F_2$ et l'évaluation passe alors à ces opérations : $$(\lambda F_1+\mu F_2)(x)=\lambda F_1(x)+\mu F_2(x)$$ et $$(F_1F_2)(x)=F_1(x)F_2(x)$$
\end{prop}
\begin{cor}
	Si $F_1$ et $F_2$ coïncident en une infinité de points, alors $F_1=F_2$.
\end{cor}
\begin{dfn}
	Soit $F\in\K(X)$ qu'on écrit sous forme quelconque $\displaystyle \frac{A}{B}$. On pose alors : $$F(-X)=\frac{A(-X)}{B(-X)}$$
\end{dfn}
\begin{dfn}[Parité, imparité]
	$F$ est dite \textbf{paire} lorsque $F(-X)=F$. $F$ est dite \textbf{paire} lorsque $F(-X)=-F$.
\end{dfn}
\begin{prop}
	On a $$(\lambda F_1+\mu F_2)(-X)=\lambda F_1(-X)+\mu F_2(-X)$$ et $$(F_1F_2)(-X)=F_1(-X)F_2(-X)$$
\end{prop}
\begin{dfn}[Généralisation]
	Plus généralement, soit $U\in\K(X)$ non constant. $\displaystyle F(U)=\frac{A(U)}{B(U)}$ est bien défini et ne dépend pas du représentant $\displaystyle\frac{A}{B}$ et les résultats de la proposition précédente restent vrais $$(\lambda F_1+\mu F_2)(U)=\lambda F_1(U)+\mu F_2(U)$$ et $$(F_1F_2)(U)=F_1(U)F_2(U)$$ Ainsi, on pourra passer de $F$ à $F(X+1)$, $F(X^2)$ et même $F(X)$ !
\end{dfn}
\section{Décomposition en éléments simples}
\begin{dfn}[Partie entière]
	Soit $F\in\K(X)$. Alors il existe un unique couple $(E,G)\in\K[X]\times\K(X)$ tel que $F=E+G$ et $\deg(G)<0$. $E$ s'appelle la \textbf{partie entière} de $F$.
\end{dfn}
\begin{dfn}
	On pose $$\K_{-1}=\left\{G\in\K(X)\ | \ \deg(G)\leqslant-1\right\}$$ C'est un sous-espace vectoriel de $\K(X)$ qui vérifie $$\K(X)=\K[X]\oplus\K_{-1}(X)$$
\end{dfn}
\begin{rmq} Après avoir effectué cette première décomposition, on peut se demander si $G$  ne peut pas se décomposer en éléments plus simples. C'est précisément l'objet de la suite du chapitre. \end{rmq}
\begin{dfn}
	Une \textbf{élément simple} de $\K(X)$ est une fraction rationnelle $\displaystyle\frac{P}{Q^j}$ avec $j\in\N\st$, $Q$ irréductible unitaire et $\deg(P)<\deg(Q)$
\end{dfn}
\begin{thm}
	La famille $\displaystyle\left(\frac{X^k}{Q^j}\right)$ avec $Q$ irréductible unitaire, $j\in\N\st$ et $0\leqslant k<\deg(Q)$ est une base de $\K_{-1}(X)$. De plus, si $\displaystyle \frac{P}{Q}$ est une écriture irréductible de $F$ et si $\lambda Q_1^{\alpha_1}\ldots Q_n^{\alpha_n}$ est la DFI de $Q$, alors la décomposition peut s'écrire $$\frac{P}{Q}=E+\sum_{i=1}^{n}\left(\sum_{j=1}^{\alpha_i}\frac{P_{i,j}}{Q_i^j}\right)$$ Autrement dit, on peut se restreindre aux éléments simples dont le dénominateur est de la forme $Q_i^j$ avec $1\leqslant j\leqslant \alpha_i$.
\end{thm}
\begin{rmq} Le premier point est utile pour les exercices théoriques. Pour les exercices calculatoires, utiliser la forme irréductible. \end{rmq}
\section{Méthodes pratiques}
\begin{met}[Méthode 1 - Identification]
	On effectue la DES théorique, puis on met tout au même dénominateur que $F$ avant de multiplier par ce dénominateur puis d'identifier les coefficients du polynôme obtenu. On obtient alors un système linéaire à résoudre.
\end{met}
\begin{thm}[DES de $\displaystyle \frac{P'}{P}$ pour $P$ scindé]
	Soit $P\in\K[X]$ un polynôme scindé dont on écrit la DFI sous la forme $\displaystyle P=\lambda\prod_{i=1}^{r}(X-x_i)^{\alpha_i}$. Alors, la DES de $\displaystyle \frac{P'}{P}$ est donnée par $$\frac{P'}{P}=\sum_{i=1}^{r}\frac{\alpha_i}{X-x_i}$$
\end{thm}
\begin{thm}[Coefficient pour un pôle simple]
	Soit $(P,Q)\in\K[X]\times(\K[X]\setminus\left\{0\right\})$ et $\lambda$ une racine simple de $Q$. Fixons $Q_1\in\K[X]$ tel que $Q=(X-\lambda)Q_1$ et $Q_1(\lambda)\ne 0$. Alors, dans la DES de $\displaystyle\frac{P}{Q}$, le terme en $\displaystyle\frac{a}{X-\lambda}$ est donné par : $$a=\frac{P(\lambda)}{Q_1(\lambda)}=\frac{P(\lambda)}{Q'(\lambda)}$$
\end{thm}
\begin{thm}[Coefficient pour un pôle de multiplicité quelconque, HP]
	Soit $(P,Q)\in\K[X]\times(\K[X]\setminus\left\{0\right\})$ et $\lambda$ une racine de multiplicité $p\in\N\st$ de $Q$. Fixons $Q_p\in\K[X]$ tel que $Q=(X-\lambda)^pQ_p$ et $Q_p(\lambda)\ne 0$. Alors, dans la DES de $\displaystyle\frac{P}{Q}$, le terme en $\displaystyle\frac{a}{(X-\lambda)^p}$ est donné par : $$a=\frac{P(\lambda)}{Q_p(\lambda)}=\frac{p!P(\lambda)}{Q^{(p)}(\lambda)}$$
\end{thm}
\begin{met}[Méthode 2 - "Évaluations" par les théorèmes qui précèdent]
	On utilise les deux théorèmes précédents et on itère. On peut utiliser les premiers résultats obtenus sur certains coefficients pour se ramener à des calculs plus simples pour les coefficients suivants.
\end{met}
\begin{rmq} On peut toujours évaluer la DES en un point qui n'est pas un pôle de la fraction rationnelle de départ. \end{rmq}
\begin{met}[Méthode 3 - Évaluation en un point qui n'est pas un pôle]
	On peut évaluer en un point qui n'est pas un pôle pour obtenir une équation linéaire reliant les différentes inconnues. En couplant cela avec les autres méthodes, cela peut accélérer les calculs finaux.
\end{met}
\begin{thm}[Passage dans $\C$]
	Soit $F\in\C(X)$ et supposons qu'il existe $A$ et $B$ des polynômes à coefficients réels tels que $\displaystyle F=\frac{A}{B}$. La DFI de $B$ peut s'écrire sous la forme $$B=\lambda\prod_{i=1}^{r}(X-\lambda_i)^{\alpha_i}\prod_{j=1}^{s}(\text{Trinôme irréductible unitaire})^{\beta_j}$$ avec les $\lambda_i\in\R$. Or, chaque trinôme irréductible unitaire peut se factoriser sous la forme $(X-\mu_j)(X-\overline{\mu_j})$ avec les $\mu_j\in\C\setminus\R$. Puis on a alors $$B=\lambda\prod_{i=1}^{r}(X-\lambda_i)^{\alpha_i}\prod_{j=1}^{s}(X-\mu_j)^{\beta_j}(X-\overline{\mu_j})^{\beta_j}$$ Dans $\C(X)$, la DES de $\displaystyle F=\frac{A}{B}$ va donc être de la forme : $$F=E+\sum_i\sum_k\frac{a_{i,k}}{(X-\lambda_i)^k}+\sum_j\sum_l\left(\frac{b_{j,l}}{(X-\mu_j)^l}+\frac{b'_{j,l}}{(X-\overline{\mu_j})^l}\right)$$ Alors, dans cette DES :
	\begin{enumerate}
		\item Tous les $a_{i,k}$ sont réels
		\item Chaque $b'_{j,l}$ vaut $\overline{b_{j,l}}$
	\end{enumerate}
\end{thm}
\begin{met}[Méthode 4 - Écriture à coefficients réels et passage dans $\C(X)$]
	Si $F$ admet une écriture à coefficients réels, utiliser le théorème précédent.
\end{met}
\begin{met}[Méthode 5 - Parité et imparité]
	Supposons que $F$  est paire ou impaire. On commence par écrire sa DES $$F=E+\sum_{i,j}\frac{a_{i,j}}{(X-\lambda_i)^j}$$ Or, on a $$F(-X)=E(-X)+\sum_{i,j}\frac{a_{i,j}}{(-X-\lambda_i)^j}$$ En utilisant l'unicité de la DES, on peut obtenir des relations sur les coefficients.
\end{met}
\begin{met}[Méthode 6 - Passage à la limite dans $\R(X)$]
	Soit $d\in\N$. on peut écrire $x\mapsto x^dF(x)$ de 2façons différentes avec la forme initiale de $F$ ou avec sa DES. Lorsqu'on fait tendre $x$ vers $+\infty$, on obtient des relations entre les coefficients par unicité de la limite puisque l'on connaît bien les limites de fractions rationnelles en $+\infty$. Pour que la méthode ait un intérêt, il faut que $d$ ne soit pas trop grand. Régulièrement, on prend $d=-\deg(F)$, mais on peut le prendre différemment parfois. On peut aussi faire des limites en $-\infty$.
\end{met}
\begin{met}[Méthode 7 - Plongement dans $\C(X)$]
	Soit $F\in\R(X)$. on l'interprète temporairement comme une fraction rationnelle de $\C(X)$ puis on regroupe les termes conjugués pour revenir à $\R(X)$.
\end{met}
\section{Primitives de fonctions rationnelles}
On considère $f\in\K(X)$ avec $\K=\R$ ou $\K=\C$. On note $\mathcal{P}$ l'ensemble de ses pôles. On cherche des primitives $\R\setminus\mathcal{P}\rightarrow\K$ \par Pour cela, on considère la DES de $F$. La partie entière étant polynomiale, on ne s'y intéressera pas et on considèrera uniquement $\K_{-1}(X)$.
\subsection{Cas de $\C(X)$}
Les éléments simples sont, à un scalaire près et transformées en fonction rationnelles, les $$x\mapsto \frac{1}{(x-\lambda)^p}$$ avec $\lambda\in\C$ et $p\in\N\st$.
\begin{thm}
	Il faut distinguer les cas :
	\begin{enumerate}
		\item Supposons que $p\geqslant2$. Une primitive de $x\mapsto (x-\lambda)^{-p}$ est donnée par $$x\mapsto \frac{1}{1-p}(x-\lambda)^{1-p}$$ Elle est définie sur $]-\infty,\lambda[$ ou $]\lambda,+\infty[$ si $\lambda\in\R$ et sur $\R$ si $\lambda\in\C\setminus\R$.
		\item Supposons que $p=1$. \textbf{Attention}, pas de $x\mapsto\ln(\lvert x-\lambda\rvert)$, cela \textbf{ne fonctionne pas} dès que $\lambda\in\C\setminus\R$. Il faut passer par la \textbf{quantité conjuguée} et on se ramène à du $\ln$ réel et une $\arctan$ qui peut être affectée d'un scalaire $i$.
	\end{enumerate}
\end{thm}
\begin{prop}[Cas général pour $p=1$, HP]
	On cherche à primitive $\displaystyle x\mapsto\frac{1}{x-(\alpha+i\beta)}$ avec $\alpha\in\R$ et $\beta\in\R\st$. La primitive est donnée par $$x\mapsto \frac{1}{2}\ln\left((x-\alpha)^2+\beta^2\right)+i\arctan\left(\frac{x-\alpha}{\beta}\right)$$ cette primitive est alors définie sur $\R$.
\end{prop}
\begin{proof}
	A savoir retrouver vite. On peut le retrouver rapidement en passant par la quantité conjuguée et en effectuant les manipulations usuelles de primitivation.
\end{proof}
\subsection{Cas de $\R(X)$}
On peut passer dans $\R(X)$ et adapter les résultats du paragraphe précédent. Pour les éléments simples de première espèce réels, on le résultat est toujours valables. Pour ceux complexes non réels, on les regroupe directement avec leur conjugué.
\begin{prop}[Cas général, HP]
	Supposons que $\lambda=\alpha+i\beta\notin\R$ et cherchons une primitive de $$x\mapsto\frac{a}{x-\lambda}+\frac{\overline{a}}{x-\overline{\lambda}}$$ Une primitive est donnée par $$x\mapsto \frac{a+\overline{a}}{2}\ln\left((x-\alpha)^2+\beta^2\right)+i(a-\overline{a})\arctan\left(\frac{x-\alpha}{\beta}\right)$$
\end{prop}
\begin{proof}
	Idem, à savoir retrouver vite. Pour le retrouver, mettre tout au même dénominateur puis effectuer les manipulations usuelles.
\end{proof}
\begin{exe}
	Trouver une primitive de $\displaystyle x\mapsto\frac{1}{(1+x+x^2)^3}$
\\ 1) Par un changement de variable affine, se ramener à une fonction de la forme $\displaystyle t\mapsto\frac{1}{(1+t^2)^3}$
\\ 2) 1ère méthode : Écrire $t$ sous la forme $\tan(u)$
\\ 3) 2ème méthode : Par une IPP, trouver par récurrence une primitive de $\displaystyle t\mapsto\frac{1}{(1+t^2)^p}$
\end{exe}
\section{Compléments, HP}
\subsection{Démonstration du théorème de décomposition en éléments simples}
\begin{thm}
	La famille $\displaystyle(\frac{X^{k}}{Q^{j}})$ avec $Q$ irréductible unitaire, $j\in\mathbb{N}^{*}$ et $0\leqslant k<deg(Q)$ est une base de $\mathbb{K}_{-1}(X)$. De plus, si $\frac{P}{Q}$ est une écriture de $F$ et si $\lambda Q_{1}^{\alpha_{1}}\ldots Q_{n}^{\alpha_{n}}$ est la DFI de $Q$, alors la décomposition peut s'écrire $$\frac{P}{Q}=E+\sum_{i=1}^{n}\left(\sum_{j=1}^{\alpha_{i}}\frac{P_{i,j}}{Q_{i}^{j}}\right)$$ Autrement dit, on peut restreindre la somme aux éléments simples dont le dénominateur est de la forme $Q_{i}^{j}$ avec $1\leqslant j\leqslant\alpha_{i}$.
\end{thm}
\begin{proof}
	Montrons d'abord le caractère générateur. Si $Q$ est constant il n'y a rien à faire. Sinon, tout revient à montrer que si $\deg(P)<\deg(Q)$ alors on peut écrire $$\frac{P}{Q}=\sum_{i=1}^{n}(\sum_{j=1}^{\alpha_{i}}\frac{P_{i,j}}{Q_{i}^{j}})$$ Pour chaque $i$, considérons le polynôme $R_{i}=\displaystyle\prod_{\substack{k=1\\ k\ne i}}^{n}Q_{k}^{\alpha_{k}}$. \\
	Si un polynôme $\Delta$ divise tous les $R_{i}$, alors il divise \textit{a fortiori} $Q$ donc sa DFI s'écrit $\mu Q_{1}^{\beta_{1}}...Q_{n}^{\beta_{n}}$. Par ailleurs, pour $i$ fixé, comme $\Delta$ divise \(R_{i}\) on doit avoir \(\beta_{i}\le0\) Ainsi, $\Delta$ est constant. On en déduit que les $R_{i}$ sont premiers entre eux dans leur ensemble, et le théorème de Bézout généralisé permet d'écrire $$U_{1}R_{1}+\ldots+U_{n}R_{n}=1 \ \ \text{puis} \frac{P}{Q}=\sum_{i=1}^{n}\frac{PU_{i}}{Q_{i}^{\alpha_{i}}}$$ On est donc ramené au cas d'une fraction rationnelle $\frac{P}{Q^{\alpha}}$ avec $Q$ irréductible. Quitte à prendre à nouveau la partie entière, on peut supposer que $\deg(P)<\deg(Q^\alpha)$. On effectue alors une récurrence sur $\alpha$.
	\begin{itemize}
		\item Initialisation : pour $\alpha=1$, il n'y a rien à faire
		\item Hérédité: soit $\alpha\ge2$ supposons la propriété vraie au rang $\alpha-1$ et montrons-la au rang $\alpha$. Soit donc $P$ tel que $\deg(P)<\deg(Q)$, et écrivons la division euclidienne de $P$ par $Q$ : $P=BQ+R$ avec $\deg(R)<\deg(Q)$. Alors
		$\displaystyle\frac{P}{Q^{\alpha}}=\frac{B}{Q^{\alpha-1}}+\frac{R}{Q^{\alpha}}$ Comme $$\deg(B)=\deg(P-R)-deg(Q)\leqslant \deg(P)-deg(Q)<deg(Q^{\alpha-1})$$ il n'y a plus qu'à invoquer l'hypothèse de récurrence pour décomposer $\displaystyle\frac{B}{Q^{\alpha-1}}$.
	\end{itemize}
	Passons maintenant à la liberté, et montrons que si $$\sum_{i=1}^{n}\left(\sum_{j=1}^{a_{i}}\frac{P_{ij}}{Q_{i}^{j}}\right)=0$$ alors tous les $P_{i,j}$ sont nuls. Supposons que pour un certain $i_0$, les $P_{i_0,j}$ soient non tous nuls, et soit $j_{0}$ maximal tel que $P_{i_{0},j_{0}}\ne0$. On multiplie alors l'égalité par $\displaystyle Q_{i_{0}}^{j_{0}}\prod_{\substack{k=1 \\ k\ne i_{0}}}^{n}Q_{k}^{\alpha_{k}}$
	\begin{itemize}
		\item Chaque $\displaystyle\frac{P_{i,j}}{Q_{i}^{j}}$ pour $i\ne i_{0}$ devient un polynôme multiple de $Q_{i_{0}}^{j_{0}}$ donc de $Q_{i_{0}}$.
		\item Quant aux termes de la forme $\displaystyle\frac{P_{i_{0},j}}{Q_{i_{0}}^{j}}$ :
		\begin{itemize}
			\item[--] si $j>j_{0}$ ils sont nuls;
			\item[--] si $j\leqslant j_{0}$ ils deviennent des polynômes multiples de $P_{i_{0},j}Q_{i_{0}}^{j_{0}-j}$. En particulier, si $j<j_{0}$ ils sont encore multiples de $Q_{i_{0}}$.
		\end{itemize}
	\end{itemize}	
	Finalement, modulo $Q_{i_{0}}$ il ne reste plus que $\displaystyle P_{i_{0},j_{0}}\prod_{\substack{k=1\\ k\ne i_0}}^{n}Q_{k}^{\alpha_{k}}$. Or $Q_{i_{0}}$ est premier avec chacun des $Q_{k}$ pour $k\ne i_{0}$ donc d'après le corollaire du théorème de Bézout il est premier avec $\displaystyle\prod_{\substack{k=1\\ k\ne i_0}}^{n}Q_{k}^{\alpha_{k}}$. Par
	le théorème de Gauss il divise donc $P_{i_{0},j_{0}}$ et comme $\deg(P_{i_0,j_0})<\deg(Q_{i_0})$ on a finalement $P_{i_{0},j_{0}}=0$. Contradiction.
\end{proof}
\subsection{Plongement de $\mathbb{R}(X)$ dans $\mathbb{C}(X)$}
Régulièrement, nous souhaitons interpréter une fraction rationnelle de $\mathbb{R}(X)$ comme une fraction rationnelle de $\mathbb{C}(X)$ à coefficients réels. Mais tout cela a-t-il un	sens ? Afin de ne pas nous y perdre, pour noter la classe d'équivalence de $(P,Q)\in\mathbb{K}[X]\times(\mathbb{K}[X]\backslash\{0\})$ nous allons momentanément abandonner la notation $\displaystyle\frac{P}{Q}$	au profit de $\displaystyle\frac{P}{Q}_{|\K}$. \\ \par	
La question est alors de savoir si, lorsque $P$ et $Q$ sont des polynômes réels, on peut identifier $\displaystyle\frac{P}{Q}_{|\R}$ et $\displaystyle \frac{P}{Q}_{|\C}$. On considère donc $$\begin{array}{ccccc}\varphi&:&\mathbb{R}(X)&\rightarrow&\mathbb{C}(X) \\
	&&\displaystyle\frac{P}{Q}_{|\R}&\mapsto&\displaystyle\frac{P}{Q}_{|\C}
\end{array}$$
\begin{itemize}
	\item Premier point: si $\displaystyle\frac{P_{1}}{Q_{1}}_{|\R}=\frac{P_{2}}{Q_{2}}_{|\R}$ alors $P_{1}Q_{2}=P_{2}Q_{1}$ dans $\mathbb{R}[X]$. Mais comme les lois	de $\mathbb{C}[X]$ prolongent celles de $\mathbb{R}[X]$, on a aussi $P_{1}Q_{2}=P_{2}Q_{1}$ dans $\mathbb{C}[X]$. Ainsi, on a $\displaystyle\frac{P_{1}}{Q_{1}}_{|\C}=\frac{P_{2}}{Q_{2}}_{|\C}$ et $\varphi$ est bien définie.
	\item Deuxième point: on a $\displaystyle\varphi\left(\frac{P}{Q}_{|\R}+\frac{R}{S}_{|\R}\right)=\varphi\left(\frac{PS+QR}{QS}_{|\R}\right)=\frac{PS+QR}{QS}_{|\C}$. De même, puisque les lois de $\mathbb{C}[X]$ prolongent celles de $\mathbb{R}[X]$, on $\displaystyle\frac{PS+QR}{QS}_{|\C}=\frac{P}{Q}_{|\C}+\frac{R}{S}_{|\C}$ donc $\varphi$ est un morphisme pour l'addition.
	\item De la même manière, on montre que $\varphi$ est un morphisme pour la multiplication. Par ailleurs, on a bien sûr $\varphi(1_{|\R})=1_{|\C}$.
	\item Dernier point : en tant que morphisme de corps, $\varphi$ est injectif.
\end{itemize}
Le dernier point montre que $\mathbb{R}(X)$ peut légitimement être considéré comme une sous-partie de $\mathbb{C}(X)$ puis le fait que $\varphi$ soit morphisme montre que dans ce cas, les définitions des lois (cas réel et cas complexe) coïncident bien.

\subsection{Division selon les puissances croissantes}
Dans le cas d'un pôle multiple, les calculs de décomposition en éléments simples peuvent être rapidement compliqués. Voici une technique très puissante qui permet d'obtenir d'un coup l'ensemble des coefficients associés à un pôle.
\begin{thm}[Division selon les puissances croissantes]
	Soit $A$ et $B$ deux polynômes de $\mathbb{K}[X]$. On suppose que le terme constant de $B$ est non nul. Alors $$\forall p\in\mathbb{N},\ \exists!(Q,R)\in\mathbb{K}[X]^{2},\ \begin{cases}A=BQ+X^{p+1}R\\ \deg(Q)\le p\end{cases}$$
\end{thm}
\begin{proof}
	Commençons par montrer l'unicité. Si $(Q,R)$ et $(Q',R')$ conviennent, alors $B(Q-Q')=X^{p+1}(R'-R)$. Comme $0$ n'est pas racine de $B$, celui-ci est premier avec $X^{p+1}$, qui divise donc $Q-Q'$ d'après le lemme de Gauss. Comme $\deg(Q-Q')\leqslant p$, on a $Q=Q'$ donc $R=R'$. \\ \par Pour l'existence à présent, on la prouve par récurrence sur $p\in\N$.
	\begin{itemize}
		\item Initialisation : pour $p=0$, notons $\lambda$ le coefficient constant de $A$ et $\mu\neq 0$ celui de $B$. Il suffit alors de choisir $\displaystyle Q=\frac{\lambda}{\mu}$. Ainsi $A-BQ$ n'a pas de coefficient constant et on a bien $A-BQ=XR$.
		\item Hérédité : soit $p\geqslant 1$ supposons la propriété vraie au rang $p-1$ et écrivons par hypothèse de récurrence $$\begin{cases}A=BQ+X^{p}R\\ \deg(Q)\le p-1\end{cases}$$
		Notons toujours $\mu\neq 0$ coefficient constant de $B$ : $B=\mu+XB_{2}$. Notons également $\nu$ le coefficient constant de $R$ : $R=\nu+XR_{2}$. On a alors
		\begin{align*}
			A&=BQ+\nu X^{p}+X^{p+1}R_{2} \\
			&=B\left(Q+\frac{\nu}{\mu}X^{p}\right)+X^{p+1}\left(R_{2}-\frac{\nu}{\mu}B_{2}\right)
		\end{align*} avec $\displaystyle\deg\left(Q+\frac{\nu}{\mu}X^{p}\right)\leqslant p$.
	\end{itemize}
	Ainsi, la récurrence est achevée.
\end{proof}
\begin{exa}[Pratique de la division selon les puissances croissantes]
	Effectuons la division de $X^{3}+X$ par $X^{2}+X+1$ à l'ordre $p=4$. On écrit les polynômes dans l'ordre des puissances croissantes. À part cela, la technique la même que pour la division euclidienne.
	$$\begin{array}{ccccc|ccc}
		X&&+X^3&&&1&+X&+X^2 \\
		\cline{6-8}
		-X&-X^2&-X^3&&&X&-X^2&+X^3 \\
		\cline{1-3}
		&-X^2&&&&&& \\
		&X^2&+X^3&+X^4&&&& \\
		\cline{2-4}
		&&X^3&+X^4&&&& \\
		&&-X^3&-X^4&-X^5&&& \\
		\cline{3-5}
		&&&&-X^5&&&
	\end{array}$$
	On a donc $X^{3}+X=(X^{2}+X+1)(X^{3}-X^{2}+X)-X^{5}$.
\end{exa}

On en déduit une technique très efficace de décomposition en éléments simples, exposée ci-contre.
\begin{met}[Décomposition en éléments simples pour un pôle multiple]
	Voici une méthode de décomposition en éléments simples pour un pôle multiple qui utilise la division selon les puissances croissantes.
	\begin{itemize}
		\item On note toujours $\displaystyle\frac{P}{Q}$	notre fraction rationnelle, avec $P\wedge Q=1$ et $Q$ unitaire.
		\item Quitte à effectuer une division euclidienne, on suppose que $\deg(P) < \deg(Q)$ (cela permet d'alléger les calculs).
		\item Soit $\lambda$ un pôle d'ordre $p$. On a donc $Q=(X-\lambda)^pQ_p$ avec $Q_p(\lambda)\neq 0$. On effectue le changement de variable $X=Y+\lambda$. On a alors $$F=\frac{P(Y+\lambda)}{Y^pQ_p(Y+\lambda)}=\frac{\tilde{P}(Y)}{Y^p\tilde{Q_p}(Y)}$$
		\item La division selon les puissances croissantes de $\overline{P}$ par $\tilde{Q_p}$ à l'ordre $p-1$ donne $$\tilde{P}(Y)=\tilde{Q_p}(Y)S(Y)+Y^pR(Y)$$
		avec $\deg(S) < p$, ce qui permet d'écrire $\displaystyle S(Y)=\sum_{i=0}^{p-1}s_iY^i$.
		\item On écrit alors $\displaystyle F=\frac{S(Y)}{Y^p}+\frac{R(Y)}{\tilde{Q_p}(Y)}=\sum_{i=0}^{p-1}\frac{s_i}{Y^{p-i}}+\frac{R(Y)}{\tilde{Q_p}(Y)}$.
		\item On revient à $Y=X-\lambda$ : on a décomposé le pôle $\lambda$.
	\end{itemize}
\end{met}
\begin{proof}
	Justifions cette méthode.
\begin{itemize}
	\item Lorsqu'on parle de "changement de variable", tout se passe comme si $Y$ devenait la
	nouvelle inconnue à la place de $X$. En fait, en toute rigueur il faut plutôt considérer qu'on introduit un polynôme $Y=X-\lambda\in\mathbb{K}[X]$.
	\item La formule du binôme montre que $P(Y+\lambda)$ peut s'écrire $\tilde{P}(Y)$ et que $Q_{p}(Y+\lambda)$ peut s'écrire $\tilde{Q_{p}}(Y)$. De plus, cette même formule du binôme montre que le terme constant de $\tilde{Q_{p}}$ est égal à $Q_{p}(\lambda)\ne0$. On peut donc bien effectuer une division selon les puissances croissantes.
	\item La méthode donne l'impression qu'on effectue la division selon les puissances croissantes avec l'inconnue $Y$. En fait, en toute rigueur il faut d'abord l'effectuer avec $X$ : $$\tilde{P}=\tilde{Q_{p}}S+X^{p}R$$ puis dans un second temps, on évalue cette égalité polynomiale en $Y$. Cela dit, il est vrai que formellement tout se passe comme si on travaillait avec $Y$ dès le départ.
	\item Enfin, il faut justifier que les termes $\displaystyle\sum_{i=0}^{p-1}\frac{s_{i}}{Y^{p-i}}$ correspondent bien à la décomposition du pôle $\lambda$. Notons qu'on a $$F=\sum_{i=0}^{p-1}\frac{s_{i}}{(X-\lambda)^{p-i}}+\frac{R(X-\lambda)}{\tilde{Q_{p}}(Y+\lambda)}=\sum_{i=0}^{p-1}\frac{s_{i}}{(X-\lambda)^{p-i}}+\frac{R(X-\lambda)}{Q_{p}}$$ $R(X-\lambda)$ est un polynôme de $\mathbb{K}[X]$ et $Q_{p}$ est un polynôme dont les facteurs irréductibles sont exactement ceux de $Q$, sauf $X-\lambda$. En effectuant la décomposition en éléments simples de la fraction rationnelle $\displaystyle\frac{R(X-\lambda)}{Q_{p}}$ et en lui ajoutant $\displaystyle\sum_{i=0}^{p-1}\frac{s_{i}}{(X-\lambda)^{p-i}}$, on conclut par unicité.
\end{itemize}
\end{proof}
\begin{exa}[Décomposition en éléments simples de $\displaystyle\frac{X^{2}+1}{(X-1)^{3}X^{2}})$]
	Ce n'est pas la peine d'effectuer une division euclidienne, le numérateur est déjà de degré strictement inférieur à celui du dénominateur. On commence par le pôle $0$ d'ordre $0$, pas besoin d'effectuer un changement de variable. La division selon les puissances croissantes de $X^2+1$ par $(X-1)^3=X^3-3X^2+3X-1$ donne $X^2+1=(X-1)^3(-3X-1)+X^2(3X^2-8X+7)$, d'où $$F=-\frac{1}{X^2}-\frac{3}{X}+\frac{3X^2-8X+7}{(X-1)^3}$$ Le changement de variable $Y=X-1$ donne $$\frac{3X^2-8X+7}{(X-1)^3}=\frac{3(Y+1)^2-8(Y+1)+7}{Y^3}=\frac{3Y^2-2Y+2}{Y^3}=\frac{2}{Y^3}-\frac{2}{Y^2}+\frac{3}{Y}$$ Finalement $$\frac{X^2+1}{(X-1)^3X^2}=-\frac{1}{X^2}-\frac{3}{X}+\frac{2}{(X-1)^3}-\frac{2}{(X-1)^2}+\frac{3}{X-1}$$
\end{exa}
\begin{exa}[Décomposition de $\displaystyle\frac{3X^4-4X^3-2X^2+4X+3}{(X^2+1)^2(X-1)})$]
	Comme 1 est un pôle simple, son coefficient est donné par $\displaystyle\frac{P(1)}{Q_2(1)}=\frac{4}{2^2}=1$. Pour $(X^2+1)^2=(X-i)^2(X+i)^2$; il suffit de calculer les coefficients de $(X-i)(X+i)^2$, les autres seront conjugués. Posons donc $Y=X-i$, ce qui donne \begin{align*}
		F&=\frac{3(Y+i)^4-4(Y+i)^3-2(Y+i)^2+4(Y+i)+3}{Y^2(Y+2i)^2(Y+i-1)}\\
		&=\frac{3Y^4+(-4+12i)Y^3-(20+12i)Y^2+(16-16i)Y+(8+8i)}{Y^2(Y^3+(-1+5i)Y^2+(-8-4i)Y+(4-4i))}
	\end{align*}
	et effectuons notre division selon les puissances croissantes.
	\begin{align*}
		3Y^4 + (-4+12i)Y^3 - (20+12i)Y^2 &+ (16-16i)Y + (8+8i)\\&= (Y^3 + (-1+5i)Y^2 + (-8-4i)Y + (4-4i))(2i + (1+i)Y)\\ &\ \ \ + Y^2((-6+2i) + (2+6i)Y + (2-i)Y^2)
	\end{align*}
	On obtient $$F = \frac{2i}{Y^2} + \frac{1+i}{Y} + \ldots$$ et finalement
	\begin{align*}
		F &= \frac{1}{X-1} + \frac{1+i}{X-i} + \frac{2i}{(X-i)^2} + \frac{1-i}{X+i} + \frac{-2i}{(X+i)^2} \\
		&= \frac{1}{X-1} + \frac{2X-2}{X^2+1} - \frac{8X}{(X^2+1)^2}
	\end{align*}
\end{exa}




\newpage
\chapter{Espaces euclidiens}
Dans tout le chapitre, les corps de base des espaces vectoriels est expressément fixé à $\K=\R$.
\section{Produit scalaire}
\begin{dfn}[Produit scalaire]
	Soit $E$ un espace vectoriel réel. Un \textbf{produit scalaire} sur $E$ est une application $\left\langle\cdot,\cdot\right\rangle:E^2\rightarrow\R$ qui vérifie les 4 axiomes suivants :
	\begin{enumerate}
		\item Bilinéarité : Pour tout $x\in E$, $y\mapsto\left\langle x,y\right\rangle$ est linéaire et pour tout $y\in E$, $x\mapsto\left\langle x,y\right\rangle$ est linéaire
		\item Symétrie : $\forall(x,y)\in E^2,\ \left\langle x,y\right\rangle=\left\langle y,x\right\rangle$
		\item Caractère positif : $\forall x\in E,\ \left\langle x,x\right\rangle \geqslant 0$
		\item Caractère défini : $\forall x\in E,\ \left\langle x,x\right\rangle=0 \Longrightarrow x=0_E$
	\end{enumerate} On résume cela en disant qu'un produit scalaire est une \textbf{forme bilinéaire symétrique définie positive}. Le produit scalaire peut aussi être noté $<x|y>$, $(x|y)$ voire $x\cdot y$.
\end{dfn}
\begin{rmq} Si on prouve la symétrie d'abord, il suffit de prouver la linéarité par rapport à l'une des variables pour prouver la bilinéarité. \end{rmq}
\begin{dfn}[Espace préhilbertien réel, espace euclidien]
	Un \textbf{espace préhilbertien réel} est un $\R$-espace vectoriel muni d'un produit scalaire. Si cet espace vectoriel est de dimension finie sur $\R$, on dit alors qu'il s'agit d'un \textbf{espace euclidien}.
\end{dfn}
\begin{rmq} Si $E$ est préhilbertien, tout sev de $E$ est naturellement muni d'une structure préhilbertienne en restreignant le produit scalaire. De même, tout sev d'un espace euclidien est naturellement muni d'une structure euclidienne. \end{rmq}
\begin{dfn}[Produit scalaire canonique de $\R^n$]
	On définit le \textbf{produit scalaire canonique de $\R^n$} par $$\left\langle \begin{pmatrix}x_1 \\ \vdots \\ x_n\end{pmatrix}, \begin{pmatrix}y_1 \\ \vdots \\ y_n\end{pmatrix}\right\rangle=\sum_{i=1}^{n}x_iy_i$$ Matriciellement, on peut écrire $$\left\langle X,Y\right\rangle=X^\top Y$$
\end{dfn}
\begin{dfn}[Produit scalaire canonique de $\mathcal{M}_{n,p}(\R)$]
	$\mathcal{M}_{n,p}(\R)$ est euclidien pour son produit scalaire canonique $$\left\langle A,B\right\rangle =\tr(A^\top B)$$
\end{dfn}
\begin{prop}[Un produit scalaire sur les polynômes]
	Dans $\R[X]$, l'application $$\left\langle \sum_{k\in\N}a_kX^k,\sum_{k\in\N}b_kX^k\right\rangle = \sum_{k\in\N}a_kb_k$$ est un produit scalaire qui munit $\R[X]$ d'un structure préhilbertienne.
\end{prop}
\begin{thm}[Produit scalaire sur les fonctions continues]
	Soit $a<b$ des réels. Dans le $\R$-espace vectoriel $\mathcal{C}^0([a,b],\R)$, l'application définie par $$\left\langle f,g\right\rangle = \int_{[a,b]} fg$$ est un produit scalaire.
\end{thm}
\begin{dfn}[Norme euclidienne associée à un produit scalaire]
	La \textbf{norme euclidienne} associée au produit scalaire $\left\langle\cdot,\cdot\right\rangle$ est définie par $$\lVert x\rVert=\sqrt{\langle x,x\rangle}$$ Elle vérifie $\lVert \lambda.x\rVert=\lvert\lambda\rvert\times\lVert x\rVert$ et $\lVert x\rVert =0 \iff x=0_E$
\end{dfn}
\begin{dfn}(Distance euclidienne associée à un produit scalaire)
	La \textbf{distance euclidienne} associée au produit scalaire $\left\langle\cdot,\cdot\right\rangle$ est définie par $$d(x,y)=\lVert x-y\rVert$$ où $\lVert\cdot\rVert$ est la norme euclidienne associée au produit scalaire $\left\langle\cdot,\cdot\right\rangle$. Elle vérifie la propriété de symétrie $$d(x,y)=d(y,x)$$
\end{dfn}
\begin{prop}[Identités remarquables]
	Soit $(E,\left\langle\cdot,\cdot\right\rangle)$ un espace préhilbertien réel. On a les \textbf{identités remarquables suivantes} :
	\begin{enumerate}
		\item $\forall (x,y)\in E^2,\ \lVert x+y\rVert^2=\lVert x\rVert^2+\lVert y\rVert^2+2\langle x,y\rangle$
		\item $\forall (x,y)\in E^2,\ \lVert x-y\rVert^2=\lVert x\rVert^2+\lVert y\rVert^2-2\langle x,y\rangle$
		\item $\forall (x,y)\in E^2,\ \langle x+y,x-y\rangle=\lVert x\rVert^2-\lVert y\rVert^2$
	\end{enumerate}
\end{prop}
\begin{prop}[Identité du parallélogramme]
	Soit $(E,\left\langle\cdot,\cdot\right\rangle)$ un espace préhilbertien réel. On a $$\forall (x,y)\in E^2,\ \lVert x+y\rVert^2+\lVert x-y\rVert^2=2\lVert x\rVert^2+2\lVert y\rVert^2$$
\end{prop}
\begin{cor}[Identités de polarisation]
	Soit $(E,\left\langle\cdot,\cdot\right\rangle)$ un espace préhilbertien réel et $\lVert\cdot\rVert$ la norme euclidienne associée. on a alors les identités de polarisation :
	\begin{enumerate}
		\item $\displaystyle \langle x,y\rangle =\frac12\left(\lVert x+y\rVert^2-\lVert x\rVert^2-\lVert y\rVert^2\right)$
		\item $\displaystyle \langle x,y\rangle =\frac12\left(\lVert x\rVert^2+\lVert y\rVert^2-\lVert x-y\rVert^2\right)$
		\item $\displaystyle \langle x,y\rangle =\frac14\left(\lVert x+y\rVert^2-\lVert x-y\rVert^2\right)$
	\end{enumerate}
\end{cor}
\begin{cor}
	Une  norme euclidienne est associée à un unique produit scalaire.
\end{cor}
\begin{dfn}[Vecteur unitaire]
	Un vecteur $x\in E$ est dit \textbf{unitaire} lorsque $\lVert x\rVert=1$.
\end{dfn}
\begin{dfn}[Normalisation]
	Lorsque $x\ne 0_E$, le vecteur $$\frac{x}{\lVert x\rVert}=\lVert x\rVert\inv.x$$ est unitaire. On l'appelle le \textbf{renormalisé} de $x$.
\end{dfn}
\begin{rmq} Plus généralement, pour $x\in E$ et $\lambda\in\R\st$, on s'autorisera à noter $\displaystyle \frac{x}{\lambda}$ à la place de $\lambda\inv.x$ \end{rmq}
\begin{thm}[Inégalité de Cauchy-Schwarz]
	Soit $(E,\left\langle\cdot,\cdot\right\rangle)$ un espace préhilbertien réel. On a $$\forall(x,y)\in E^2,\ \lvert\langle x,y\rangle\rvert \leqslant\lVert x\rVert \lVert y\rVert$$ avec égalité si, et seulement si, $x$ et $y$ sont colinéaires.
\end{thm}
\begin{rmq} L'inégalité reste vraie sans le caractère défini (mais pas le cas d'égalité). On pourra donc appliquer l'inégalité dans le cas de formes bilinéaires symétriques positives. \end{rmq}
\begin{cor}[Inégalité triangulaire]
	La norme euclidienne associée à $\langle\cdot,\cdot\rangle$ vérifie l'inégalité triangulaire : $$\forall (x,y)\in E^2,\ \lVert x+y\rVert\leqslant\lVert x\rVert+\lVert y\rVert$$ avec égalité si, et seulement si, $x$ et $y$ sont colinéaires directs. 
\end{cor}
\begin{cor}[Seconde inégalité triangulaire]
	On en déduit la second inégalité triangulaire : $$\forall (x,y)\in E^2,\ \big|\lVert x\rVert - \lVert y\rVert\big| \leqslant \lVert x-y\rVert$$
\end{cor}
\begin{dfn}[Norme]
	Une \textbf{norme} sur un espace vectoriel $E$ est une application $\lVert\cdot\rVert:E\rightarrow\R_+$ qui vérifie les axiomes suivants :
	\begin{enumerate}
		\item Séparation : $\forall x\in E,\ \lVert x\rVert=0 \Rightarrow x=0$
		\item Homogénéité : $\forall x\in E,\  \forall \lambda\in\R,\ \lVert \lambda.x\rVert=\lvert\lambda\rvert\times\lVert x\rVert$
		\item Inégalité triangulaire : $\forall (x,y)\in E^2,\ \lVert x+y\rVert\leqslant\lVert x\rVert+\lVert y\rVert$
	\end{enumerate}
\end{dfn}
\begin{rmq} Toutes les normes ne dérivent pas d'un produit scalaire. Pour ces normes, dites non euclidiennes et étudiées en deuxième année, on ne peut plus rien dire sur le cas d'égalité de l'inégalité triangulaire. \end{rmq}

\section{Orthogonalité}
Dans toute cette partie, $(E,\langle\cdot,\cdot\rangle)$ est un espace préhilbertien réel.
\subsection{Définitions de base}
\begin{dfn}[Orthogonalité]
	Deux vecteurs $x$ et $y$ sont dits \textbf{orthogonaux} lorsque $\langle x,y\rangle=0$. Dans ce cas, on écrit $x\perp y$. Un vecteur $x$ est dit \textbf{orthogonal} à une partie $B$ lorsqu'il est orthogonal à tout vecteur de $B$ : $$\forall y \in B, \ \langle x,y\rangle=0$$ Dans ce cas, on écrit $x\perp B$. Deux parties $A$ et $B$ sont dites \textbf{orthogonales} lorsque tout vecteur de l'une est orthogonal à tout vecteur de l'autre : $$\forall (x,y)\in A\times B,\ \langle x,y\rangle=0$$ Dans ce cas, on écrit $A\perp B$.
\end{dfn}
\begin{dfn}[Orthogonal d'une partie]
	Soit $A$ une partie de $E$. L'ensemble des vecteurs orthogonaux à tous les vecteurs de $A$ est un sous-espace vectoriel de $E$, appelé \textbf{orthogonal} de $A$. On le note $$A^\perp =\left\{x\in E \ | \ \forall a\in A, \ \langle x,a\rangle=0\right\}$$ C'est aussi la plus grande partie orthogonale à $A$ au sens de l'inclusion.
\end{dfn}
\begin{exa}
	On a $E^\perp=\left\{0\right\}$ et $\left\{0\right\}^\perp=E$.
\end{exa}
\begin{prop}[Décroissance de l'othogonal]
	Si $A\subset B$ alors $B^\perp\subset A^\perp$.
\end{prop}
\begin{cor}
	Soit $A$ une partie de $E$. Alors $A^\perp=\vect(A)^\perp$.
\end{cor}
\subsection{Familles orthogonales}
\begin{dfn}[Famille orthogonale, famille orthonormale]
	Une famille $(x_i)_{i\in I}$ de vecteurs de $E$ est dite \textbf{orthogonale} si les $x_i$ sont deux à deux orthogonaux. Si de plus les $x_i$ sont tous unitaires, la famille est dite \textbf{orthonormale} (ou \textbf{orthonormée}).
\end{dfn}
\begin{exa}
	Dans $\R^n$ muni de son produit scalaire canonique, la base canonique est orthonormale.
\end{exa}
\begin{exa}
	Dans $\R[X]$ muni de son produit scalaire usuel, la base canonique est orthonormale.
\end{exa}
\begin{prop}
	Toute sous-famille d'une famille orthogonale est orthogonale. Toute sous-famille d'une famille orthonormale est orthonormale. 
\end{prop}
\begin{prop}
	Toute famille orthogonale de vecteurs \textbf{non nuls} est libre.
\end{prop}
\begin{cor}
	Toute famille orthonormale est libre.
\end{cor}
\begin{thm}[Théorème de Pythagore]
	Deux vecteurs $x$ et $y$ sont orthogonaux si, et seulement si, $$\lVert x+y\rVert^2=\lVert x\rVert^2+\lVert y\rVert^2$$
\end{thm}
\begin{thm}
	Soit $(x_i)_{1\leqslant i\leqslant n}$ une famille orthogonale finie. Alors $$\left\lVert \sum_{i=1}^{n}x_i\right\rVert^2=\sum_{i=1}^{n}\lVert x_i\rVert^2$$
\end{thm}
\begin{rmq} Attention, la réciproque est fausse dès que $n\geqslant 3$. On pourra chercher un contre-exemple dans $\R^3$ muni de son produit scalaire canonique. \end{rmq}
\begin{thm}[Existence d'une base orthonormale pour un espace euclidien]
	Tout espace euclidien admet une base orthonormale.
\end{thm}
\begin{prop}[Expression dans une base orthonormale]
	Soit $(e_i)_{1\leqslant i\leqslant n}$ une base orthonormale de $E$ euclidien. Alors tout vecteur s'exprime dans cette base par : $$x=\sum_{i=1}^{n}\langle x,e_i\rangle e_i$$ De plus, si on écrit $\displaystyle x=\sum_{i=1}^{n}x_ie_i$ et $\displaystyle y=\sum_{i=1}^{n}y_ie_i$, on a : $$\langle x,y\rangle =\sum_{i=1}^{n}x_iy_i$$ et $$\lVert x\rVert^2=\sum_{i=1}^{n}x_i^2$$
\end{prop}
\subsection{Supplémentaire orthogonal}
\begin{prop}
	L'orthogonalité est une condition suffisante (mais pas nécessaire !) pour que deux sous-espaces vectoriels $F$ et $G$ soient en somme directe.
\end{prop}
\begin{dfn}[Supplémentaires orthogonaux]
	Soit $F$ et $G$ deux sous-espaces vectoriels à la fois supplémentaires et orthogonaux. On dit alors que $F$ et $G$ sont \textbf{supplémentaires orthogonaux} et on écrit $$E=F\overset{\perp}{\oplus} G$$ On dit que $G$ est un supplémentaire orthogonal $F$, et vice-versa.
\end{dfn}
\begin{thm}[Théorème de la base adaptée, version orthonormale]
	Soit $(e_i)_{i\in I}$ une base orthonormale de $E$ et supposons que $I=I_1\sqcup I_2$. Alors $E_1=\vect(e_i)_{i\in I_1}$ et $E_2=\vect(e_i)_{i\in I_2}$ sont supplémentaires orthogonaux. Réciproquement, si $E_1$ et $E_2$ sont supplémentaires orthogonaux, alors en concaténant des bases orthonormales de $E_1$ et $E_2$, on obtient une base orthonormale de $E$.
\end{thm}
\begin{prop}
	Si $G$ est un supplémentaire orthogonal de $F$, alors on a nécessairement $G=F^\perp$. S'il existe, le supplémentaire orthogonal est donc unique.
\end{prop}
\begin{cor}
	Si $F$ admet un supplémentaire orthogonal, alors $F^\perp$ admet un supplémentaire orthogonal et on a $(F^\perp)^\perp=F$. 
\end{cor}
\begin{dfn}
	Si $F$ admet un supplémentaire orthogonal $G$, le \textbf{projecteur orthogonal} est le projecteur $p$ sur $F$ parallèlement à G. $p(x)$ s'appelle le \textbf{projeté orthogonal} de $x$ sur $F$.
\end{dfn}
\begin{rmq} Le projecteur associé $q=\id-p$ est donc le projecteur orthogonal sur $F^\perp$. \end{rmq}
\begin{dfn}[Symétrie orthogonale]
	Si $F$ admet un supplémentaire orthogonal $G$, la \textbf{symétrie orthogonale d'axe $F$} est la symétrie par rapport à $F$ parallèlement à $G$. Si on l'appelle $s$, on rappelle qu'on a $$s=2p-\id$$
\end{dfn}
\begin{dfn}[Réflexion]
	Une symétrie orthogonale par rapport à un hyperplan s'appelle une \textbf{réflexion}.
\end{dfn}
\begin{prop}[Caractérisation de l'orthogonalité pour les projecteurs et les symétries]
	Un projecteur $p$ est orthogonal si, et seulement si, $\ker(p)$ et $\im(p)$ sont orthogonaux. Une symétrie $s$ est orthogonale si, et seulement si, $\ker(s-\id)$ et $\ker(s+\id)$ sont orthogonaux.
\end{prop}
\begin{dfn}
	Soit $F$ un sous-espace vectoriel de $E$ et $x\in E$. La \textbf{distance} de $x$ à $F$ est définie par $$d(x,F)=\underset{y\in F}{\inf}\lVert x-y\rVert$$ Elle et bien définie car l'ensemble est non vide car contient $\lVert x-0\rVert$ et minoré par $0$ en tant qu'ensemble de normes euclidiennes.
\end{dfn}
\begin{thm}
	Supposons que $F$ admette un supplémentaire orthogonal, et soit $p$ le projecteur orthogonal sur $F$. Alors, pour tout $x\in E$, $d(x,F)$ est un minimum et $p(x)$ est l'unique élément de $F$ qui le réalise.
\end{thm}
\begin{thm}[Synthèse partielle - cas fondamental]
	Si $F$ est de dimension finie, alors on a $E=F\overset{\perp}{\oplus}F^\perp$. Dans ce cas, $F^\perp$ admet un supplémentaire orthogonal égal à $F$ et on pourra toujours définir le projecteur orthogonal sur $F$ et la symétrie orthogonale par rapport à $F$.
\end{thm}
\begin{rmq} Attention on a des contre-exemples en dimension infinie. On pourra considérer $\vect(X^n-1)_{n\in\N\st}$ dans $\R[X]$ muni de son produit scalaire usuel, dont l'orthogonal vaut $\left\{0\right\}$. Sinon, on pourra considérer l'ensemble $$H=\left\{f\in\mathcal{C}^0([0,1],\R) \ | \ f(0)=0\right\}$$ dans $\mathcal{C}^0([0,1],\R)$ muni de son produit scalaire usuel, qui est un hyperplan car le noyau d'un morphisme d'évaluation non-nul. Son orthogonal vaut alors $\left\{0\right\}$ (le prouver en prenant $f\in H^\perp$ puis en effectuant le produit scalaire avec $t\mapsto tf(t) \in H$ et conclure par signe et intégrale nulle pour les points autre que 0, et par continuité en 0). \end{rmq}
\begin{cor}
	Si $E$ est euclidien et $F$ est un sous-espace vectoriel de $E$, on a toujours $$\dim(F^\perp)=\dim(E)-\dim(F)$$
\end{cor}
\begin{cor}[Expression du projecteur orthogonal en dimension finie]
	Si $F$ est de dimension finie et $(e_i)_{1\leqslant i\leqslant k}$ est une base orthonormale de $F$, le projecteur orthogonal sur $F$ est donné par $$\forall x\in E,\ p(x)=\sum_{i=1}^{k}\langle x,e_i\rangle e_i$$
\end{cor}
\begin{rmq} Si on ne dispose pas d'une base orthonormale, mais simplement d'une base $(e_i)_{1\leqslant i\leqslant k}$, on peut soit appliquer l'algorithme de Gram-Schmidt expliqué après, soit raisonner de la manière suivante. \end{rmq}
\par On sait qu'il existe des scalaires $(\lambda_1,\ldots,\lambda_k)\in\R^k$ tels que $\displaystyle p(x)=\sum_{j=1}^{k}\lambda_je_j$. Ensuite, on sait que $x-p(x)\in F^\perp$, donc on obtient $$\forall i\in\llbracket1,k\rrbracket,\ \langle x,e_i\rangle = \sum_{j=1}^{k}\lambda_j\langle e_j,e_i\rangle$$ On obtient ainsi un système linéaire de $k$ équations à $k$ inconnues (dont on peut démontrer qu'il est de Cramer) qu'on peut alors toujours résoudre pour obtenir les coefficients.
\begin{prop}[Théorème de la base incomplète, version orthonormale]
	Si $E$ est euclidien, alors toute famille orthonormale de $E$ peut être complétée en une base orthonormale.
\end{prop}
\begin{thm}[Orthonormalisation de Gram-Schmidt]
	Soit $(x_1,\ldots,x_n)$ une famille libre de $E$. Alors, il existe une famille orthonormale $(\varepsilon_1,\ldots,\varepsilon_n)$ telle que $$\forall k\in\llbracket1,n\rrbracket,\ \vect(x_1,\ldots,x_k)=\vect(\varepsilon_1,\ldots,\varepsilon_k)$$
\end{thm}
\begin{rmq}[Mise en œuvre pratique du procédé d'orthonormalisation] On construit la famille par récurrence finie sur $k$. \\ Initialisation : On pose $\displaystyle \varepsilon_1=\frac{x_1}{\lVert x_1\rVert}$ \\ Hérédité : Soit $k\in\llbracket2,n\rrbracket$ et supposons avoir construit $(\varepsilon_1,\ldots,\varepsilon_{k-1})$ On pose alors $$\tilde{\varepsilon}_k=x_k-\sum_{i=1}^{k-1}\langle x_k,\varepsilon_i\rangle\varepsilon_i$$ puis on pose alors $\displaystyle \varepsilon_k=\frac{\tilde{\varepsilon}_k}{\lVert \tilde{\varepsilon}_k\rVert}$ \end{rmq}
\begin{rmq} Le théorème fonctionne encore avec une famille dénombrable, c'est-à-dire une famille de la forme $(x_k)_{k\in\N}$. On pensera notamment aux polynômes. \end{rmq}
\begin{rmq} On peut montrer par analyse-synthèse que la famille $(\varepsilon_1,\ldots,\varepsilon_n)$ est la seule famille qui vérifie la condition supplémentaire : $$\forall k\in\llbracket1,n\rrbracket,\ \langle x_k,\varepsilon_k\rangle >0$$ \end{rmq}
\section{Compléments : introduction à la géométrie affine, HP}
On se donne un espace euclidien $E$ (mais on généraliserait sans problème à un $\K$-ev quelconque).
\\ \par 
Jusqu'ici, nous avons considéré $\mathbb{R}^2$ ou $\mathbb{R}^3$ comme des espaces vectoriels. Leurs éléments sont des vecteurs, et intuitivement on les assimile à des déplacements. Mais en géométrie élémentaire, lorsqu'on "pose un repère", on sait que $\mathbb{R}^2$ et $\mathbb{R}^3$ peuvent être aussi vus comme des ensembles de points. Cette fois, leurs éléments sont plutôt assimilés à des positions fixes. Cela n'a pas tellement de sens d'additionner deux positions. En revanche, leur différence correspond à un déplacement, et on peut lui associer un vecteur.
\\ \par
Pour bien faire la différence entre $\mathbb{R}^2$ espace de points et $\mathbb{R}^2$ espace de vecteurs, on peut noter les points en ligne : $(x_1, x_2)$ et les vecteurs en colonne :$\displaystyle\begin{pmatrix}
	x_1 \\
	x_2
\end{pmatrix}$
De même avec $\mathbb{R}^3$. Formalisons
\begin{dfn}[Espace affine]
	Soit $\mathcal{E}$ un ensemble non vide. $\mathcal{E}$ est un \textbf{espace affine} de \textbf{direction} $E$ lorsqu'il existe une application $(A,B) \in \mathcal{E}^2 \mapsto \overrightarrow{AB} \in E$ vérifiant les deux axiomes suivants.
\begin{itemize}
	\item $\forall A \in \mathcal{E},\ \forall u \in E,\ \exists! B \in \mathcal{E},\ \overrightarrow{AB} = u.$ On notera $B = A + u.$
	\item Relation de Chasles : $\forall(A,B,C) \in \mathcal{E}^3,\ \overrightarrow{AB} + \overrightarrow{BC} = \overrightarrow{AC}$
\end{itemize}
\end{dfn}
\begin{rmq}
	Par souci de clarté, on essayera au maximum de garder les lettres majuscules pour les points de l'espace affine, et les lettres minuscules ou les flèches pour les vecteurs de l'espace vectoriel.
\end{rmq}
\begin{exa}
	$\mathbb{R}^n$ peut être considéré comme un espace affine : à deux $n$-uplets $(x_1, ..., x_n)$ et $(y_1, ..., y_n)$ on associe le vecteur $\displaystyle\begin{pmatrix}
	y_1 - x_1 \\
	\vdots \\
	y_n - x_n
\end{pmatrix}$
\end{exa}

\begin{exa}
	À deux points $A$ et $B$ de $E$, faisons correspondre le vecteur $B - A.$ On vérifie immédiatement que $E$ est un espace affine de direction lui-même. La notation $B = A + u$ reste parfaitement valable.
\end{exa}

\begin{dfn}[Translation]
	Si $u \in E,$ la translation de vecteur $u$ est l'application $t_u : A \mapsto A + u.$
\end{dfn}

\begin{prop}
	On a les propriétés suivantes.
\begin{itemize}
	\item $\forall(A,B) \in \mathcal{E}^2,\ A = B \Leftrightarrow \overrightarrow{AB} = \overrightarrow{0}$
	\item $\forall(A,B) \in \mathcal{E}^2, \overrightarrow{BA} = -\overrightarrow{AB}.$
\end{itemize}
\end{prop}

\begin{proof}
	Soit $A$ point de $\mathcal{E}$ fixé, et écrivons $\overrightarrow{AA} = \overrightarrow{AA} + \overrightarrow{AA},$ si bien que $\overrightarrow{AA} = \overrightarrow{0}.$ Donc si $A = B$ alors $\overrightarrow{AB} = \overrightarrow{0}.$ Réciproquement, si $\overrightarrow{AB} = \overrightarrow{0},$ alors comme $\overrightarrow{AA} = \overrightarrow{0},$ par ... on en déduit $A = B$
Pour le deuxième point, il suffit d'écrire $\overrightarrow{AB} + \overrightarrow{BA} = \overrightarrow{AA} = \overrightarrow{0}.$
\end{proof}

\begin{dfn}[Barycentre]
	Soit $A_1, ..., A_n$ des points de $E$ $(n \geqslant 1),$ et $\alpha_1, ..., \alpha_n$ des scalaires de somme non nulle. Alors il existe un unique $G$ tel que $\sum_i \alpha_i \overrightarrow{GA_i} = \overrightarrow{0}.$ C'est le \textbf{barycentre} de $((A_1, \alpha_1), ..., (A_n, \alpha_n))$. De plus on a $$\forall O \in \mathcal{E},\ \overrightarrow{OG} = \frac{\sum_i \alpha_i \overrightarrow{OA_i}}{\sum_i \alpha_i}.$$ Plus généralement, on définit de la même manière le barycentre de $(A_i, \alpha_i)_{i \in I}$ avec $I$ fini.
\end{dfn}
\begin{proof}
	Existence: utilisons la fin de l'énoncé avec $O=A_{1}$ et montrons que
$$G=A_{1}+\frac{\sum_{i}\alpha_{i}\overrightarrow{A_{1}A_{i}}}{\sum_{i}\alpha_{i}}$$
convient. On a $$\sum_{i}\alpha_{i}\overrightarrow{GA_{i}}=\sum_{i}\alpha_{i}(\overrightarrow{GA_{1}}+\overrightarrow{A_{1}A_{i}})=\overrightarrow{GA_{1}}\sum_{i}\alpha_{i}+\sum_{i}\alpha_{i}\overrightarrow{A_{1}A_{i}}=(\sum_{i}\alpha_{i})(\overrightarrow{GA_{1}}+\overrightarrow{A_{1}G})=\overrightarrow{0}$$
Unicité : si $G_{1}$ et $G_{2}$ conviennent, alors en soustrayant $\displaystyle\sum_{i}\alpha_{i}\overrightarrow{G_{1}A_{i}}=\overrightarrow{0}$ et $\displaystyle\sum_{i}\alpha_{i}\overrightarrow{G_{2}A_{i}}=\overrightarrow{0}$ on obtient $\displaystyle\sum_{i}\alpha_{i}\overrightarrow{G_{1}G_{2}}=\overrightarrow{0}$ et on conclut en remarquant que $\displaystyle\sum_{i}\alpha_{i}\ne0$.
Enfin, si on se donne $O\in\mathcal{E}$ quelconque, alors par la relation de Chasles on a$$\frac{\sum_{i}\alpha_{i}\overrightarrow{OA_{i}}}{\sum_{i}\alpha_{i}}=\frac{\sum_{i}\alpha_{i}(\overrightarrow{OG}+\overrightarrow{GA}_{i})}{\sum_{i}\alpha_{i}}=\overrightarrow{OG}_{i}$$
ce qui achève la preuve.
\end{proof}
\begin{prop}[Commutativité du barycentre]
	Pour toute permutation $\sigma\in\mathcal{S}_{n}$ le barycentre de $((A_{1},\alpha_{1}),...,(A_{n},\alpha_{n}))$ est aussi le barycentre de $((A_{\sigma(1)},\alpha_{\sigma(1)}),...,(A_{\sigma(n)},\alpha_{\sigma(n)}))$
\end{prop}
\begin{proof}
	C'est une conséquence immédiate de la commutativité de la somme.
\end{proof}

\begin{prop}[Homogénéité du barycentre]
	Si $\lambda$ est un scalaire non nul, alors le barycentre de $((A_{1},\alpha_{1}),...,(A_{n},\alpha_{n}))$ est aussi le barycentre de $((A_{1},\lambda\alpha_{1}),...,(A_{n},\lambda\alpha_{n}))$.
\end{prop}
\begin{proof}
	Il suffit de multiplier l'égalité $\displaystyle\sum_{i}\alpha_{i}\overrightarrow{GA_{i}}=\overrightarrow{0}$ par $\lambda$.
\end{proof}

\begin{prop}[Associativité du barycentre]
	Soit $\{I_{1},...,I_{p}\}$ une partition de $\llbracket1,n\rrbracket$ telle que $$\forall r\in \llbracket1,p\rrbracket,\ \sigma_{r}=\sum_{i\in I_{r}}\alpha_{i}\ne0$$

Notons $G_{r}$ le barycentre de $((A_{1},\alpha_{i}))_{i\in I_{r}}$ (pour tout $r\in\llbracket1,p\rrbracket$). Alors le barycentre de $((A_{1},\alpha_{1}),...,(A_{n},\alpha_{n}))$ est aussi barycentre de $((G_{1},\sigma_{1}),...,(G_{p},\sigma_{p}))$.
\end{prop}

\begin{proof}
	Notons G le barycentre de $((A_{1},\alpha_{1}),...,(A_{n},\alpha_{n}))$ et notons provisoirement $G^{\prime}$ le barycentre de $((G_{1},\sigma_{1}),...,(G_{p},\sigma_{p})).$ Il suffit d'écrire 
	\begin{align*}
		\sum_{i\in I}\alpha_{i}\overrightarrow{G^{\prime}A_{i}}&=\sum_{r=1}^{p}\left(\sum_{i\in I_{r}}\alpha_{i}\overrightarrow{G^{\prime}G_{r}}+\sum_{i\in I_{r}}\alpha_{i}\overrightarrow{G_{r}A_{i}}\right)\\
		&=\sum_{r=1}^{p}\sigma_{r}\overrightarrow{G^{\prime}G_{r}}\\
		&=\overrightarrow{0}
	\end{align*} Donc par unicité on a $G=G^{\prime}$
\end{proof}

\begin{dfn}[Repère affine]
	Un \textbf{repère affine} de $E$ est la donnée d'un point $O$ et d'une
base $(e_{i})_{1\le i\le n}$ de $E$. Le point $O$ est l'\textbf{origine} du repère. $n$ s'appelle la \textbf{dimension} de $E$. Le repère est \textbf{orthonormal} lorsque la base $(e_{i})_{1\le i\le n}$ est orthonormale.
\end{dfn}
\begin{prop}
	L'application 
	$$\begin{array}{ccccc}
	\varphi&:&\mathbb{R}^{n}&\longrightarrow& E \\
	&&(x_{1},...,x_{n}) &\longmapsto& M = O + \displaystyle\sum_{i}x_{i}e_{i}
	\end{array}$$
est une bijection.
\end{prop}
\begin{proof}
	Si $\displaystyle M = O + \sum_{i}x_{i}e_{i} = M^{\prime} = O + \sum_{i}x_{i}^{\prime}e_{i_{i}}$ alors $\displaystyle\sum_{i}x_{i}e_{i}=\sum_{i}x_{i}^{\prime}e_{i},$ donc est injective. Pour la surjectivité, si $M\in E$ alors il existe $x_{1},...,x_{n}$ tels que $\displaystyle\sum_{i}x_{i}e_{i}=\overrightarrow{OM}$.
\end{proof}
\begin{dfn}[Coordonnées dans un repère affine]
	Les scalaires $x_{1},...,x_{n}$ s'appellent les \textbf{coordonnées} de $M$ dans le repère. En particulier, le point $O$ a toutes ses coordonnées nulles.
\end{dfn}
\begin{exa}
	Si $A$ a pour coordonnées $(x_{1}^{A},...,x_{n}^{A})$ dans le repère et $B$ a pour coordonnées $(x_{1}^{B},...,x_{n}^{B})$ dans le repère, alors $\overrightarrow{AB}$ a pour coordonnées $\displaystyle\begin{pmatrix}x_{1}^{B}-x_{1}^{A}\\ \vdots\\ x_{n}^{B}-x_{n}^{A}\end{pmatrix}$ dans la base $(e_{i})_{1\le i\le n}$.
\end{exa}
\begin{exa}
	Si $G$ est le barycentre de $((A_{1},\alpha_{1}),...,(A_{p},\alpha_{p}))$ alors ses coordonnées sont les moyennes des coordonnées des $A$, pondérées par les $\alpha_{i}$.
\end{exa}
\begin{exa}
	Dans $\mathbb{R}^{n}$ vu comme un espace affine, on pose souvent $O=(0,...,0)$ et
$(e_{i})_{1\le i\le n}$ la base canonique (on l'appelle le repère canonique). Dans ce cas, si on se donne un élément $M-(x_{1},...,x_{n})\in\mathbb{R}^{n},$ ses coordonnées dans le repère canonique sont exactement les $x_{i}$.
\end{exa}
\begin{dfn}
	Un \textbf{sous-espace affine} (ou sea) de $E$ est une partie de $E$ de la forme $\mathcal{F}=A+F$ où $A$ est un point de $E$ et $F$ un sev de $E$. $F$ est lui-même un espace affine de direction $F$. $A$ est appelé une origine de $F$.
\end{dfn}
\begin{rmq}
	Lorsqu'on considère $E$ comme un espace affine de direction lui-même, on
retrouve bien la définition d'un sea vue précédemment. Il s'agit donc d'une généralisation.
\end{rmq}
\begin{rmq}
	N'importe quel point de $F$ peut servir d'origine. En revanche, la direction
est bien fixée de manière unique. La démonstration est la même que dans le chapitre "Espaces vectoriels".
\end{rmq}
\begin{dfn}[Droite affine]
	Si $A$ et $B$ sont deux points distincts, la \textbf{droite affine} $(AB)$ est le sous-espace affine $A+\mathbb{R}\overrightarrow{AB}.$ On vérifie que si $C$ et $D$ sont deux points distincts de $(AB)$, alors $(CD)=(AB)$.
\end{dfn}
\begin{dfn}
	Si $A$ et $B$ sont deux points distincts, la \textbf{demi-droite} $[AB)$ est l'ensemble $A+\mathbb{R}_{+}\overrightarrow{AB}$ ce n'est plus un espace affine).
\end{dfn}
 Soit $\mathcal{E}$ un espace affine de direction $E$. Une \textbf{application affine} de $\mathcal{E}$ est une application $f:\mathcal{E}\rightarrow\mathcal{E}$ telle qu'il existe une application linéaire $\varphi\in\mathcal{L}(E)$ vérifiant
$$\forall(A,B)\in\mathcal{E}^{2},\overrightarrow{f(A)f(B)}=\varphi(\overrightarrow{AB})$$
On voit que si elle existe, $\varphi$ est nécessairement unique. On l'appelle \textbf{partie linéaire} de $f$.
À ce titre, on relira avec profit la remarque de la partie Géométrie du chapitre "Nombres complexes".
\begin{prop}
	La composée de deux applications affines est une application affine.
\end{prop}
\begin{proof}
	Soit $f$ et $g$ nos deux applications affines, avec $\varphi$ et $\psi$ associées. Alors $$\forall(A,B)\in\mathcal{E}^{2},\overrightarrow{f(g(A))f(g(B))}=\varphi(\overrightarrow{g(A)g(B)})=\varphi(\psi(\overrightarrow{AB}))$$
Ainsi la preuve est achevée.
\end{proof}
\begin{rmq}
	Au passage, on voit que lorsqu'on compose les applications affines, on compose les parties linéaires.
\end{rmq}




\newpage
\chapter{Fonctions usuelles}
\section{Révisions de lycée (et plus si affinités...)}
\subsection{Équations, inéquations}
Pour les rédactions à adopter afin de résoudre des équations et des inéquations, on pourra se référer au chapitre "Logique et raisonnements". Pour les méthodes de résolution, nous n'en donnons ici que les noms : équivalences successives, analyse-synthèse, discussion selon un paramètre, tableau de signes, théorèmes de cours explicitant un ensemble de solutions (équations trigonométriques par exemple)... \\ \par A ce propos, on rappelle que dans un tableau de signe, le signe $+$ signifie \textbf{strictement positif}, et le signe $-$ signifie \textbf{strictement négatif} !
\subsection{Tableaux de variations}
On rappelle que dans un tableau de variations, le symbole $\nearrow$ signifie \textbf{strictement croissante et continue}, et le symbole $\searrow$ signifie \textbf{strictement décroissante et continue}.
\begin{met}[Tableau de variations]
	Voici comment réaliser le tableau de variations. 
	\begin{itemize}
		\item On définit le domaine d'étude $D$ et on étudie la une fonction définie sur $D$.
		\item On justifie la dérivabilité de la fonction : combinaison linéaire, produit, quotient dont le dénominateur ne s'annule pas, composée... de fonctions dérivables. Si ces théorèmes généraux ne suffisent pas, on peut être forcé de revenir à la définition de la dérivabilité sur des points précis du domaine.
		\item On calcule la dérivée. Pour le rédiger proprement, on écrit "$\forall x\in D,\ f(x)=\ldots$ donc $\forall x\in D,\ f'(x)=\ldots$.
		\item On étudie le signe de $f'$. En théorie, on a besoin de savoir le signe au sens strict de $f'$, donc trois informations (stricte positivité, stricte négativité, nullité), mais en pratique seules deux suffisent car elles permettent de déterminer la troisième avec exactitude. Le plus souvent, le signe est facile (fonctions polynomiales de degré $1$ ou $2$), mais lorsque $f'$ est plus complexe, on pourra utiliser la rédaction suivante, plus efficace : "Soit $x\in D$. $f'(x)\geqslant 0$ ssi $\ldots$, avec égalité ssi $\ldots$.
		\item De cette étude du signe de $f'$, on en déduit les variations de $f$ grâce à divers théorèmes. Pour l'énoncé de ces théorèmes et de leurs raffinements, on se réfèrera au chapitre "Dérivation".
	\end{itemize}
\end{met}
\subsection{Transformations de graphes}
\begin{dfn}[Courbe représentative]
	Si $\mathcal{P}$ désigne le plan usuel muni d'un repère orthonormal $(O ; \overrightarrow{e_1},\overrightarrow{e_2})$, on peut dessiner la \textbf{courbe représentative} d'une fonction $f:I\rightarrow\R$ : $$\mathcal{C}_f:=\left\{M(x,y)\in\mathcal{P}\ | \ x\in I,\ y=f(x)\right\}$$
\end{dfn}
\begin{prop}[Translation horizontale]
	Soit $a\in\R$ et $$\begin{array}{ccccc}t_{h,a}&:&\mathcal{P}&\rightarrow&\mathcal{P}\\&&(x,y)&\mapsto&(x+a,y)\end{array}$$ Alors $t_{h,a}(\mathcal{C}_f)$ a pour équation $y=f(x-a)$, $ie$ $$t_{h,a}(\mathcal{C}_f)=\left\{M'(x,y)\ | \ x\in I+a,\ y=f(x-a)\right\}$$
\end{prop}
\begin{prop}[Dilatation d'un facteur $\lambda$ selon $Ox$]
		Soit $\lambda\in\R\st$ et $$\begin{array}{ccccc}d_{h,\lambda}&:&\mathcal{P}&\rightarrow&\mathcal{P}\\&&(x,y)&\mapsto&(\lambda x,y)\end{array}$$ Alors $d_{h,\lambda}(\mathcal{C}_f)$ a pour équation $y=f\left(\frac{x}{\lambda}\right)$, , $ie$ $$d_{h,\lambda}(\mathcal{C}_f)=\left\{M'(x,y)\ | \ x\in \lambda I,\ y=f\left(\frac{x}{\lambda}\right)\right\}$$
\end{prop}
\begin{prop}[Symétrie d'axe la droite d'équation $x=a$]
	Soit $a\in\R$ et $$\begin{array}{ccccc}s_{h,a}&:&\mathcal{P}&\rightarrow&\mathcal{P}\\&&(x,y)&\mapsto&(2a-x,y)\end{array}$$ Alors $s_{h,a}(\mathcal{C}_f)$ a pour équation $y=f(2a-x)$, $ie$ $$s_{h,a}(\mathcal{C}_f)=\left\{M'(x,y)\ | \ x\in 2a-I,\ y=f(2a-x)\right\}$$
\end{prop}
\begin{prop}[Translation verticale]
	Soit $a\in\R$ et $$\begin{array}{ccccc}t_{v,a}&:&\mathcal{P}&\rightarrow&\mathcal{P}\\&&(x,y)&\mapsto&(x,y+a)\end{array}$$ Alors $t_{v,a}(\mathcal{C}_f)$ a pour équation $y=f(x)+a$, $ie$ $$t_{v,a}(\mathcal{C}_f)=\left\{M'(x,y)\ | \ x\in I,\ y=f(x)+a\right\}$$
\end{prop}\begin{prop}[Dilatation d'un facteur $\lambda$ selon $Oy$]
Soit $\lambda\in\R\st$ et $$\begin{array}{ccccc}d_{v,\lambda}&:&\mathcal{P}&\rightarrow&\mathcal{P}\\&&(x,y)&\mapsto&(x,\lambda y)\end{array}$$ Alors $d_{v,\lambda}(\mathcal{C}_f)$ a pour équation $y=\lambda f\left(x\right)$, , $ie$ $$d_{h,\lambda}(\mathcal{C}_f)=\left\{M'(x,y)\ | \ x\in I,\ y=\lambda f\left(x\right)\right\}$$
\end{prop}
\begin{prop}[Symétrie d'axe la droite d'équation $y=a$]
Soit $a\in\R$ et $$\begin{array}{ccccc}s_{h,a}&:&\mathcal{P}&\rightarrow&\mathcal{P}\\&&(x,y)&\mapsto&(x,2a-y)\end{array}$$ Alors $s_{v,a}(\mathcal{C}_f)$ a pour équation $y=2a-f(x)$, $ie$ $$s_{v,a}(\mathcal{C}_f)=\left\{M'(x,y)\ | \ x\in I,\ y=2a-f(x)\right\}$$
\end{prop}
\begin{prop}[Symétrie d'axe la droite d'équation $y=x$ lorsque $f$ est bijective]
	Supposons que $f:\R\mapsto\R$ est bijective et soit $$\begin{array}{ccccc}s&:&\mathcal{P}&\rightarrow&\mathcal{P}\\&&(x,y)&\mapsto&(y,x)\end{array}$$ Alors $s(\mathcal{C}_f)$ a pour équation $y=f\inv(x)$, $ie$ $$s(\mathcal{C}_f)=\left\{M'(x,y)\ | \ x\in \R,\ y=f\inv(x)\right\}$$
\end{prop}
\begin{dfn}[Équations de droites du plan]
	On dispose de deux types d'équations de droites du plan :
	\begin{itemize}
		\item Celles de la forme $y=ax+b$. Le \textbf{seul} inconvénient de ce type d'équations (et c'est le seul, car elles sont beaucoup plus manipulables que le type suivant) est qu'elles ne fournissent pas les droites verticales.
		\item Celles de la forme $ax+by+c=0$, dites "\textbf{équations cartésienne}" avec $(a,b)\ne(0,0)$. Dans ce cas, on rappelle qu'un \textbf{vecteur directeur} de la droite est donné par $\displaystyle\begin{pmatrix}
			-b\\ a
		\end{pmatrix}$ et qu'un \textbf{vecteur normal} de la droite est donné par $\displaystyle\begin{pmatrix}
		a\\ b
		\end{pmatrix}$.
	\end{itemize}
\end{dfn}
\subsection{Systèmes linéaires}
On se réfèrera à la dernière partie du chapitre "Matrices" qui traite en détail les systèmes linéaires. La méthode du pivot de Gauss doit être maîtrisée sur le bout des doigts.
\subsection{Deux théorèmes admis}
Dans ce paragraphe, on se donne $I$ et $J$ deux intervalles non triviaux de $\R$ et $f:I\rightarrow J$ une bijection. On pose $g:=f\inv:J\rightarrow I$. On donne ici deux résultat admis qui seront démontrés dans la suite de chapitres d'analyse.
\begin{thm}[Théorème de la bijection réciproque, cas continu]
	Si $f$ (avec les hypothèses du paragraphe) est strictement monotone, alors $g$ est continue.
\end{thm}
\begin{rmq}
	On montre aisément que $g$ a mêmes variations que $f$, et on en déduit alors que $f$ elle-même est continue en appliquant le théorème à $g$.
\end{rmq}
\begin{thm}
	Soit $a\in I$ et $b:=f(a)$. Si $f$ est strictement monotone et dérivable en $a$, alors $g$ est dérivable en $b$ si, et seulement si, $f'(a)\ne0$, et dans ce cas, on a $g'(b)=\displaystyle\frac{1}{f'(a)}$.
\end{thm}
\begin{rmq}
	Lorsque $f'$ ne s'annule pas, on en déduit une généralisation : $g$ est dérivable, et on a $g'=\displaystyle\frac{1}{f'\circ f}$.
\end{rmq}
\section{Logarithme, exponentielle}
\subsection{Logarithme népérien, logarithme en base $a$}
\begin{dfn}[Logarithme népérien]
	La fonction \textbf{logarithme népérien}, notée $\ln$ est l'unique primitive s'annulant en $1$ de la fonction $$\begin{array}{ccc}\R_+\st&\rightarrow&\R_+\st\\x&\mapsto&\displaystyle\frac{1}{x}\end{array}$$ Autrement dit, $\ln$ est caractérisée par le fait que $\ln(1)=0$, que $\ln$ est dérivable et que $\forall x>0,\ \displaystyle \ln'(x)=\frac1x$.
\end{dfn}
\begin{rmq}
	La justification de cette définition est admise pour l'instant et prendra son sens avec le théorème fondamental de l'analyse vu dans le chapitre "Intégration".
\end{rmq}
\begin{exa}
	En particulier, $\ln$ est continue.
\end{exa}
\begin{exa}
	On déduit de la dérivabilité en $1$ que $\displaystyle \frac{\ln(1+h)}{h}\xrightarrow[h\to 0]{}1$
\end{exa}
\begin{thm}[Primitive de $\displaystyle\frac{u'}{u}$]
	Soit $I\subset\R$ un intervalle non trivial et $u\in\mathcal{D}^1(I,\R)$ telle que $\forall x\in I,\ u(x)\ne0$. Alors, une primitive de $\displaystyle\frac{u'}{u}$ est donnée par $\ln\lvert u\rvert$.
\end{thm}
\begin{rmq}
	Attention à ne pas oublier les valeurs absolues qui sont importantes ! Toujours les mettre, puis les enlever si l'intérieure est positif.
\end{rmq}
\begin{thm}[Propriété fonctionnelle du logarithme népérien]
	On a : $$\forall x,y>0,\ \ln(xy)=\ln(x)+\ln(y)$$
\end{thm}
\begin{exa}
	Ainsi, $\ln$ est un morphisme du groupe $(\R_+\st,\times)$ dans le groupe $(\R,+)$.
\end{exa}
\begin{cor}[Autres propriétés fonctionnelles du logarithme népérien]
	On a les propriétés fonctionnelles suivantes pour le $\ln$ : 
	\begin{enumerate}
		\item $\forall x>0,\ \ln\left(\frac1x\right)=-\ln(x)$
		\item $\forall x,y>0,\ \ln\left(\frac{x}{y}\right)=\ln(x)-\ln(y)$
		\item $\forall x>0,\ \forall n\in\Z,\ \ln(x^n)=n\ln(x)$
	\end{enumerate}
\end{cor}
\begin{prop}[Limites du logarithme népérien]
	On a : $$\begin{cases}
		\ln(x)\xrightarrow[x\to +\infty]{}+\infty\\
		\ln(x)\xrightarrow[x\to 0^+]{}-\infty
	\end{cases}$$
\end{prop}
\begin{prop}[Inégalité de concavité pour le logarithme népérien]
	On a : $$\forall x>-1,\ \ln(1+x)\leqslant x$$ avec égalité si, et seulement si $x=0$.
\end{prop}
\begin{exa}
	En appliquant l'inégalité de concavité à $\frac{-x}{1+x}$, on montre que $$\forall x>-1,\ \ln(1+x)\geqslant\frac{x}{1+x}$$
\end{exa}
\begin{dfn}[Logarithme en base $a$]
	Soit $a\in\R_+\st\setminus\left\{1\right\}$. Le \textbf{logarithme en base $a$} est la fonction définie par : $$\begin{array}{ccccc}\log_a&:&\R_+\st&\rightarrow&\R\\&&x&\mapsto&\displaystyle\frac{\ln(x)}{\ln(a)}\end{array}$$
\end{dfn}
\subsection{Exponentielle}
\begin{dfn}[Exponentielle]
	On définit la \textbf{fonction exponentielle}, notée $\exp$ par : $$\exp=\ln\inv$$
\end{dfn}
\begin{exa}
	$\exp(0)=1$
\end{exa}
\begin{dfn}[Constante $e$]
	On pose $e:=\exp(1)$. On a $e\approx 2,71828$.
\end{dfn}
\begin{exa}
	On a $\ln(e)=1$ et $\log_e=\ln$.
\end{exa}
\begin{prop}[Dérivabilité et dérivée de $\exp$]
	$\exp$ est dérivable et vérifie $\exp'=\exp$.
\end{prop}
\begin{cor}
	$\exp$ est strictement croissante.
\end{cor}
\begin{exa}
	De la dérivabilité de $\exp$ en $0$ on déduit $\displaystyle\frac{\exp(h)-1}{h}\xrightarrow[h\to 0]{}1$.
\end{exa}
\begin{thm}[Propriétés fonctionnelles de l'exponentielle]
	On a : 
	\begin{enumerate}
		\item $\forall (x,y)\in\R^2,\ \exp(x+y)=\exp(x)\exp(y)$
		\item $\forall x\in\R,\ \exp(-x)=\displaystyle\frac{1}{\exp(x)}$
		\item $\forall (x,y)\in\R^2,\ \exp(x-y)=\displaystyle\frac{\exp(x)}{\exp(y)}$
		\item $\forall x\in\R,\ \forall n\in\Z,\ \exp(nx)=\exp(x)^n$
	\end{enumerate}
\end{thm}
\begin{exa}
	Ainsi, $\exp$ réalise un morphisme du groupe $(\R,+)$ dans le groupe $(\R_+\st,\times)$.
\end{exa}
\begin{prop}[Limites de l'exponentielle]
	On a : $$\begin{cases}
		\exp(x)\xrightarrow[x\to +\infty]{}+\infty\\
		\exp(x)\xrightarrow[x\to -\infty]{}0
	\end{cases}$$
\end{prop}
\begin{prop}[Inégalité de convexité pour l'exponentielle]
	On a : $$\forall x\in\R,\ \exp(x)\geqslant 1+x$$ avec égalité si, et seulement si $x=0$.
\end{prop}
\subsection{Exponentielle quelconque}
\begin{dfn}[Puissance quelconque d'un réel strictement positif]
	Soit $a\in\R_+\st$ et $x\in\R$. On \textbf{pose} : $$a^x:=\exp(x\ln(a))$$
\end{dfn}
\begin{exa}
	On a donc $\forall x\in\R,\ \exp(x)=e^x$. On rencontrera souvent cette forme, plus courte que la notation $\exp(x)$.
\end{exa}
\begin{rmq}
	On peut prolonger cette définition à $a\in\R_+\st$ et $x\in\C$.
\end{rmq}
\begin{exa}
	Soit $a\in\R_+\st\setminus\left\{1\right\}$ et $b\in\R_+\st$. On a : $$\forall x\in\R,\ a^x=b\iff x=\log_a(b)$$
\end{exa}
\begin{exa}
	Soit $a\in\R_+\st$ et $x\in\R$. On a $\ln(a^x)=x\ln(a)$.
\end{exa}
\begin{rmq}
	On vérifie que dans le cas où $x\in\Z$, la définition donnée ci-dessus (analytique) coïncide bien avec la définition algébrique (itérative) de $a^x$.
\end{rmq}
\begin{prop}[Propriétés fonctionnelles]
	Soit $a,b>0$ et $x,y\in\R$ (voire $\C$). On a :
	\begin{itemize}
		\item $a^x\times a^y=a^{x+y}$
		\item $a^{-x}=\displaystyle\frac{1}{a^x}$
		\item $a^{x-y}=\displaystyle\frac{a^x}{b^y}$
		\item $a^x\times b^x=(ab)^x$
		\item $(a^x)^y=a^{xy}$
	\end{itemize}
\end{prop}
\begin{prop}[Croissances comparées du logarithme népérien]
	On a $$\frac{\ln(x)}{x}\xrightarrow[x\to +\infty]{}0$$ et $$x\ln(x)\xrightarrow[x\to 0^+]{}0$$
\end{prop}
\begin{cor}[Généralisation]
	Soit $\alpha,\beta>0$. On a $$\frac{\ln(x)^\alpha}{x^\beta}\xrightarrow[x\to +\infty]{}0$$ et $$x^\beta\lvert\ln(x)\rvert^\alpha\xrightarrow[x\to 0^+]{}0$$
\end{cor}
\begin{prop}[Croissances comparées de l'exponentielle]
	On a $$\frac{x}{\exp(x)}\xrightarrow[x\to +\infty]{}0$$ et $$x\exp(x)\xrightarrow[x\to -\infty]{}0$$
\end{prop}
\begin{cor}[Généralisation]
	Soit $\beta,\gamma>0$. On a $$\frac{x^\beta}{\exp(x)^\gamma}\xrightarrow[x\to +\infty]{}0$$ et $$\lvert x\rvert^\beta\exp(x)^\gamma\xrightarrow[x\to -\infty]{}0$$
\end{cor}
\begin{exa}
	Soit $x>0$ fixé. On a : $a^x\xrightarrow[a\to 0^+]{}0$.
\end{exa}
\begin{dfn}[$0^x$]
	Par \textbf{convention}, on pose : $$\forall x>0,\ 0^x=0$$ On rappelle que $0^0=1$ par \textbf{définition} et que pour $x<0$, $0^x$ \textbf{n'existe pas}.
\end{dfn}
\begin{rmq}
	Les identités restent vraies, du moment qu'elles sont définies. Pour être précis, il s'agit de :\begin{itemize}
		\item $\forall x,y>0,\ 0^x\times 0^y=0^{x+y}$
		\item $\forall x>0,\ \forall b>0,\ 0^x\times b^x=(0\times b)^x$
		\item $\forall x,y>0,\ (0^x)^y=0^{xy}$
	\end{itemize}
\end{rmq}
\begin{exa}
	On a $$\forall a\in\R_+\st,\ \forall n\in\N\st,\ \sqrt[n]{a}=a^{\frac1n}$$
\end{exa}
\section{Trigonométrie}
\subsection{Trigonométrie circulaire (cosinus, sinus et tangente)}
\begin{prop}[Une inégalité utile]
	On a : $$\forall x\in\R,\ \lvert\sin(x)\rvert\leqslant\lvert x\rvert$$
\end{prop}
\begin{dfn}[Fonction arc-cosinus]
	La fonction $\cos_{|[0,\pi]}^{[-1,1]}$ est bijective. On appelle \textbf{arc-cosinus}, et on note $\arccos:[-1,1]\rightarrow[0,\pi]$ sa bijection réciproque.
\end{dfn}
\begin{dfn}[Fonction arc-sinus]
	La fonction $\sin_{|[-\frac{\pi}{2},\frac{\pi}{2}]}^{[-1,1]}$ est bijective. On appelle \textbf{arc-sinus}, et on note $\arcsin:[-1,1]\rightarrow[-\frac{\pi}{2},\frac{\pi}{2}]$ sa bijection réciproque.
\end{dfn}
\begin{dfn}[Fonction arc-tangente]
	La fonction $\tan_{|]-\frac{\pi}{2},\frac{\pi}{2}[}$ est bijective. On appelle \textbf{arc-tangente}, et on note $\arctan:\R\rightarrow]-\frac{\pi}{2},\frac{\pi}{2}[$ sa bijection réciproque.
\end{dfn}
\begin{lem}
	On a : $$\forall x\in[-1,1],\ \cos(\arcsin(x))=\sin(\arccos(x))=\sqrt{1-x^2}$$
\end{lem}
\begin{thm}[Étude de $\arccos$]
	$\arccos$ est continue et strictement décroissante. Elle est dérivable en tout point de $]-1,1[$ et on a : $$\forall x\in]-1,1[,\ \arccos'(x)=\frac{-1}{\sqrt{1-x^2}}$$
\end{thm}
\begin{thm}[Étude de $\arcsin$]
	$\arcsin$ est impaire, continue et strictement croissante. Elle est dérivable en tout point de $]-1,1[$ et on a : $$\forall x\in]-1,1[,\ \arcsin'(x)=\frac{1}{\sqrt{1-x^2}}$$
\end{thm}
\begin{thm}[Étude de $\arctan$]
	$\arctan$ est impaire, continue et strictement croissante. Elle est dérivable en tout point de $\R$ et on a : $$\forall x\in\R,\ \arctan'(x)=\frac{1}{1+x^2}$$
\end{thm}
\begin{prop}[Trois relations utiles]
	On a les relations :\begin{itemize}
		\item $\forall x\in ]-1,1[,\ \arccos(x)+\arcsin(x)=\displaystyle\frac{\pi}{2}$
		\item $\displaystyle\forall x>0,\ \arctan(x)+\arctan\left(\frac{1}{x}\right)=\frac{\pi}{2}$ Cette relation s'avèrera très utile pour certains DL.
		\item $\displaystyle\forall x<0,\ \arctan(x)+\arctan\left(\frac{1}{x}\right)=-\frac{\pi}{2}$
	\end{itemize}
\end{prop}
\subsection{Trigonométrie hyperbolique}
\begin{dfn}[Partie paire, partie impaire]
	Soit $f:\R\rightarrow\R$. Il existe un unique couple de fonctions $(f_1,f_2)\in\left(\R^\R\right)^2$ tel que $f_1$ soit paire, $f_2$ soit impaire et on ait $f=f_1+f_2$. $f_1$ s'appelle la \textbf{partie paire} de $f$ et $f_2$ s'appelle la \textbf{partie impaire} de $f$.
\end{dfn}
\begin{rmq}
	Le résultat et la définition fonctionnent encore pour $f:I\rightarrow\R$ avec $I$ centré en $0$ ($ie$ $\forall x\in I, -x\in I$)
\end{rmq}
\begin{dfn}[Cosinus hyperbolique]
	La fonction \textbf{cosinus hyperbolique}, notée $\cosh$ ou ch est définie comme étant la partie paire de la fonction $\exp$ : $$\begin{array}{ccccc}\cosh&:&\R:\rightarrow&\R\\&&x&\mapsto&\displaystyle\frac{e^x+e^{-x}}{2}\end{array}$$
\end{dfn}
\begin{dfn}[Sinus hyperbolique]
	La fonction \textbf{sinus hyperbolique}, notée $\sinh$ ou sh est définie comme étant la partie impaire de la fonction $\exp$ : $$\begin{array}{ccccc}\sinh&:&\R:\rightarrow&\R\\&&x&\mapsto&\displaystyle\frac{e^x-e^{-x}}{2}\end{array}$$
\end{dfn}
\begin{dfn}[Tangente hyperbolique]
	La fonction \textbf{cosinus hyperbolique}, notée $\tanh$ ou th est définie par: $$\begin{array}{ccccc}\tanh&:&\R:\rightarrow&\R\\&&x&\mapsto&\displaystyle\frac{\sinh(x)}{\cosh(x)}\end{array}$$
\end{dfn}
\begin{thm}
	On a : $$\forall x\in\R,\ \cosh(x)^2-\sinh(x)^2=1$$
\end{thm}
\begin{rmq}
	C'est cette relation qui explique le nom de "trigonométrie hyperbolique" : elle définit le paramétrage d'une hyperbole.
\end{rmq}
\begin{thm}[Étude de $\cosh$]
	$\cosh$ est paire, strictement décroissante sur $]-\infty,0]$ et strictement croissante sur $[0,+\infty[$. Elle est dérivable et vérifie $\cosh'=\sinh$. On a $\cosh(0)=1$, $\displaystyle\cosh(x)\xrightarrow[x\to +\infty]{}+\infty$ et $\displaystyle\cosh(x)\xrightarrow[x\to -\infty]{}+\infty$.
\end{thm}
\begin{thm}[Étude de $\sinh$]
	$\sinh$ est impaire et strictement croissante. Elle est dérivable et vérifie $\sinh'=\cosh$. On a $\sinh(0)=0$, $\displaystyle\sinh(x)\xrightarrow[x\to +\infty]{}+\infty$ et $\displaystyle\sinh(x)\xrightarrow[x\to -\infty]{}-\infty$.
\end{thm}
\begin{thm}[Étude de $\tanh$]
	$\tanh$ est impaire et strictement croissante. Elle est dérivable et vérifie $\tanh'=1-\tanh^2=\displaystyle\frac{1}{\cosh^2}$. On a $\tanh(0)=0$, $\displaystyle\tanh(x)\xrightarrow[x\to +\infty]{}1$ et $\displaystyle\tanh(x)\xrightarrow[x\to -\infty]{}-1$.
\end{thm}
\begin{met}[Règle d'Osborn pour les formules de trigonométrie hyperbolique, HP]
	Il existe des formules de trigonométrie hyperbolique, analogues aux formules de trigonométrie circulaire. Ces formules ne sont pas au programme et pas à connaître, mais elles peuvent parfois servir. Il existe une façon simple de retenir ces formules, qu'on appelle \textbf{règle d'Osborn} : une formule de trigonométrie hyperbolique est la même qu'une formule de trigonométrie circulaire en remplaçant les $\cos$ par des $\cosh$, les $\sin$ par des $\sinh$ et les $\tan$ par des $\tanh$, à ceci près qu'un $\sin(a)\sin(b)$ doit être remplacé par un $-\sinh(a)\sinh(b)$.
\end{met}
\begin{exa}
	La formule $\cos(a+b)=\cos(a)\cos(b)-\sin(a)\sin(b)$ devient $$\cosh(a+b)=\cosh(a)\cosh(b)+\sinh(a)\sinh(b)$$ (on pourra le vérifier en revenant à la définition si on souhaite vérifier la validité de cette sorcellerie). La formule $\cos(x)^2+\sin(x)^2=1$ devient $\cosh(x)^2-\sinh(x)^2=1$, ce que nous avions prouvé par des calculs explicites.
\end{exa}
\begin{rmq}[Justification de la règle d'Osborn, HP]
	Ceci est dû au fait qu'on peut prolonger les fonctions $\sin$ et $\cos$ à $\C$ grâce aux formules d'Euler, et alors $\sinh$ et $\cosh$ ne sont que des cas particuliers de sinus et de cosinus complexes. Par exemple, on a $\forall x\in\R,\ \sinh(x)=-i\sin(ix)$ et $\forall x\in\R,\ \cosh(x)=\cos(ix)$ ce qui explique l'apparition du signe $-$ devant ce qui était auparavant un produit de 2 sinus. On peut en fait généraliser les formules de trigonométrie circulaire aux sinus et cosinus prolongés sur $\C$.
\end{rmq}
\section{Quelques rappels sur le calcul de primitives}
On se réfèrera au chapitre "Intégration", et plus précisément à la partie "Intégration et dérivation", où tout est développé. \\ \par Voici un tableau des primitives usuelles à connaître par cœur :$$
\begin{tabular}{|c|c|}
	\hline $f(x)$&$\displaystyle\int f$ \\
	\hline $x^\alpha \quad (\alpha\in\R\setminus\left\{-1\right\})$&$\displaystyle\frac{x^{\alpha+1}}{\alpha+1}$ \\
	$\displaystyle\frac{1}{x}$&$\ln \lvert x\rvert$ \\
	$e^{ax} \quad (a\in\R\st)$&$\displaystyle\frac{e^{ax}}{a}$ \\
	\hline $\cos(x)$&$\sin(x)$\\
	$\sin(x)$&$-\cos(x)$\\
	$\tan(x)$&$-\ln\lvert\cos(x)\rvert$ \\
	\hline $\displaystyle\frac{1}{\cos(x)}$&$\displaystyle\ln\left\lvert\tan\left(\frac{x}{2}+\frac{\pi}{4}\right)\right\rvert\ \ (\text{NE})$\\
	$\displaystyle\frac{1}{\sin(x)}$&$\displaystyle\ln\left\lvert\tan\left(\frac{x}{2}\right)\right\rvert\ \ (\text{NE})$\\
	\hline $\displaystyle\frac{1}{1+x^2}$&$\arctan(x)$\\
	$\displaystyle\frac{1}{1-x^2}$&argth$(x)$\\
	$\displaystyle\frac{1}{\sqrt{1+x^2}}$&argsh$(x)$\\
	$\displaystyle\frac{1}{\sqrt{x^2-1}}$&argch$(x)$\\
	$\displaystyle\frac{1}{\sqrt{1-x^2}}$&$\arcsin(x)$\\
	$\displaystyle\frac{-1}{\sqrt{1-x^2}}$&$\arccos(x)$\\
	\hline
\end{tabular}$$

\newpage
\chapter{Limites et continuité} %chapitre complet%
\section{Préliminaires}
Dans tout le paragraphe, on fixe $X\subset\R$.
\subsection{Initiation à la topologie}
On introduit ici quelques notions de topologie simples qui nous permettront d'unifier la suite du chapitre, et en particulier de ne pas à avoir à donner différentes définitions pour les différentes limites.
\begin{dfn}[Voisinages fondamentaux dans $\R$]
	Soit $a\in\R$. Un \textbf{voisinage fondamental de $a$} est un intervalle de la forme $[a-\delta,a+\delta]$ où $\delta\in\R_+\st$. On notera ici $\mathcal{V}(a)$ (NS) l'ensemble des voisinages fondamentaux de $a$.
\end{dfn}
\begin{exa}
	Soit $u\in\R^\N$ et $l\in\R$. On a $u_n\xrightarrow[n\to+\infty]{}l$ si, et seulement si, $\forall V\in\mathcal{V}(l),\ u_n\in V$ APCR.
\end{exa}
\begin{exa}
	$X$ est dense dans $\R$ si, et seulement si, $\forall a\in\R,\ \forall V\in\mathcal{V}(a),\ X\cap V\ne\emptyset$.
\end{exa}
\begin{dfn}[Adhérence dans $\R$]
	Soit $a\in\R$. $a$ est dit \textbf{adhérent à $X$} lorsque : $$\forall V\in\mathcal{V}(a),\ X\cap V\ne\emptyset$$ L'\textbf{adhérence de $X$ dans $\R$} est l'ensemble $\left\{a\in\R\ | \ a\text{ est adhérent à }X\right\}$.
\end{dfn}
\begin{exa}
	Tout point de $X$ est adhérent à $X$.
\end{exa}
\begin{exa}
	Personne n'est adhérent à l'ensemble vide.
\end{exa}
\begin{exa}
	$X$ est dense dans $\R$ si, et seulement si, l'adhérence de $X$ dans $\R$ vaut $\R$.
\end{exa}
\begin{thm}[Caractérisation séquentielle de l'adhérence dans $\R$]
	Soit $a\in\R$. $a$ est adhérent à $X$ si, et seulement si, $$\exists(x_n)_{n\in\N}\in X^\N,\ x_n\xrightarrow[n\to+\infty]{}a$$
\end{thm}
Désormais, on étend nos définitions à $\overline{\R}$.
\begin{dfn}[Voisinages fondamentaux dans $\overline{\R}$]
	Soit $a\in\overline{\R}$.\begin{itemize}
		\item Si $a\in\R$, la définition d'un voisinage fondamental reste la même.
		\item Si $a=+\infty$, on appelle voisinage fondamental de $a$ tout intervalle de la forme $[M,+\infty]$ avec $M\in\R$.
		\item Si $a=-\infty$, on appelle voisinage fondamental de $a$ tout intervalle de la forme $[-\infty,m]$ avec $m\in\R$.
	\end{itemize}
	Dans tous les cas, les voisinages fondamentaux sont des intervalles de $\overline{\R}$ qui contiennent $a$ (mais ce n'est pas suffisant, $\left\{a\right\}$ n'est pas un voisinage fondamental de $a$ mais contient $a$). On notera toujours $\mathcal{V}(a)$ l'ensemble des voisinages fondamentaux de $a$.
\end{dfn}
\begin{exa}
	Soit $u\in\R^\N$ et $l\in\overline{\R}$. On a $u_n\xrightarrow[n\to+\infty]{}l$ si, et seulement si, $\forall V\in\mathcal{V}(l),\ u_n\in V$ APCR.
\end{exa}
\begin{exa}
	Pour tout $a\in\overline{\R}$, $\mathcal{V}(a)$ est totalement ordonné pour $\subset$.
\end{exa}
\begin{prop}
	Pour tout $a\in\overline{\R}$, $\mathcal{V}(a)$ est stable par intersection finie (conséquence de la relation d'ordre totale).
\end{prop}
\begin{dfn}[Adhérence dans $\overline{\R}$]
	Soit $a\in\overline{\R}$. $a$ est dit \textbf{adhérent à $X$} lorsque : $$\forall V\in\mathcal{V}(a),\ X\cap V\ne\emptyset$$ L'\textbf{adhérence de $X$ dans $\overline{\R}$} est l'ensemble : $$adh_{\overline{\R}}(X):=\left\{a\in\overline{\R}\ | \ a\text{ est adhérent à }X\right\}$$ Jusqu'à la fin de l'année, implicitement, l'adhérence de $X$ sera toujours l'adhérence de $X$ dans $\overline{\R}$. On la note plus souvent $\overline{X}$.
\end{dfn}
\begin{thm}[Caractérisation séquentielle de l'adhérence dans $\overline{\R}$]
	Soit $a\in\overline{\R}$. $a$ est adhérent à $X$ si, et seulement si, $$\exists(x_n)_{n\in\N}\in X^\N,\ x_n\xrightarrow[n\to+\infty]{}a$$
\end{thm}
\begin{exa}
	Soit $u,v\in\overline{\R}$ tels que $u<v$. Alors $\overline{]u,v[}=[u,v]$
\end{exa}
\begin{dfn}[Propriété au voisinage de]
	Soit $f:X\rightarrow\R$ et $a\in\overline{X}$. On dit que \textbf{$f$ possède la propriété $P$ au voisinage de $a$} lorsqu'il existe un voisinage fondamental $V$ de $a$ tel que $f_{|X\cap V}$ possède la propriété $P$.
\end{dfn}
\begin{exa}
	$\cos$ est à valeurs strictement positives au voisinage de $0$. Il suffit de prendre $V=\displaystyle\left[-\frac{\pi}{4},\frac{\pi}{4}\right]$.
\end{exa}
\begin{exa}
	$\sin$ n'est pas à valeurs strictement positives au voisinage de $0$.
\end{exa}
\subsection{Premières définitions}
\begin{dfn}[Structure de $\R$-algèbre de $\R^X$, rappel]
	On rappelle que si $f:X\rightarrow\R$, $g:X\rightarrow\R$ et $\lambda\in\R$, on définit : $$\begin{array}{ccccc}f+g&:&X&\rightarrow&\R\\&&x&\mapsto&f(x)+g(x)\end{array}$$
	$$\begin{array}{ccccc}f\times g&:&X&\rightarrow&\R\\&&x&\mapsto&f(x)\times g(x)\end{array}$$ $$\begin{array}{ccccc}\lambda\cdot f&:&X&\rightarrow&\R\\&&x&\mapsto&\lambda\cdot f(x)\end{array}$$ Ainsi, $(\R^X,+,\times)$ est un anneau et $(\R^X,+,\times,\cdot)$ est une $\R$-algèbre.
\end{dfn}
\begin{dfn}[Relation d'ordre partielle sur $\R^X$, rappel]
	On définit une relation binaire sur $\R^X$ par : $$\forall(f,g)\in(\R^X)^2,\ f\leqslant g\iff \forall x\in X,\ f(x)\leqslant g(x)$$ On vérifie que $\leqslant$ est une relation d'ordre (partielle en général) sur $\R^X$.
\end{dfn}
\begin{dfn}[Fonction majorée, minorée, bornée, rappel]
	Soit $f\in\R^X$. On rappelle que :
	\begin{itemize}
		\item $f$ est dite \textbf{majorée} lorsque : $\exists M\in\R,\ \forall x\in X,\ f(x)\leqslant M$
		\item $f$ est dite \textbf{minorée} lorsque : $\exists m\in\R,\ \forall x\in X,\ f(x)\geqslant m$
		\item $f$ est dite \textbf{bornée} lorsqu'elle est à la fois majorée et minorée
	\end{itemize}
\end{dfn}
\begin{exa}
	L'ensemble des fonctions bornées de $X$ dans $\R$ forme une sous-algèbre de $\R^X$.
\end{exa}
\begin{prop}
	$f$ est bornée si, et seulement si, $\lvert f\rvert$ est majorée.
\end{prop}
\begin{dfn}[Maximum, minimum, maximum local, minimum local en un point]
	Soit $f:X\rightarrow\R$ et $a\in X$. 
	\begin{itemize}
		\item On dit que $f$ admet un \textbf{maximum en $a$} lorsque $\forall x\in X,\ f(x)\leqslant f(a)$.
		\item On dit que $f$ admet un \textbf{minimum en $a$} lorsque $\forall x\in X,\ f(x)\geqslant f(a)$.
		\item On dit que $f$ admet un \textbf{maximum local en $a$} lorsque $\exists V\in\mathcal{V}(a),\ \forall x\in X\cap V,\ f(x)\leqslant f(a)$.
		\item On dit que $f$ admet un \textbf{minimum local en $a$} lorsque $\exists V\in\mathcal{V}(a),\ \forall x\in X\cap V,\ f(x)\geqslant f(a)$.
	\end{itemize}
\end{dfn}
\begin{dfn}[Maximum et minimum d'une fonction]
	Soit $f:X\rightarrow\R$.
	\begin{itemize}
		\item S'il existe, le maximum de $f$ peut être noté $\underset{x\in X}{\max}(f(x))$.
		\item S'il existe, le minimum de $f$ peut être noté $\underset{x\in X}{\min}(f(x))$.
	\end{itemize}
\end{dfn}
\begin{dfn}[Sup et inf d'une fonction]
	Plus généralement, supposons que $X\ne\emptyset$ et soit $f:X\rightarrow\R$. Alors, $f(X)\ne\emptyset$ donc on peut définir :
	\begin{itemize}
		\item $\underset{x\in X}{\sup}(f(x))\in\R\cup\left\{+\infty\right\}$
		\item $\underset{x\in X}{\inf}(f(x))\in\R\cup\left\{-\infty\right\}$
	\end{itemize}
\end{dfn}
\begin{dfn}[Monotonies]
	Soit $f:X\rightarrow\R$.
	\begin{itemize}
		\item $f$ est dite \textbf{croissante} lorsque $\forall(x,y)\in X^2,\ x\leqslant y\implies f(x)\leqslant f(y)$
		\item $f$ est dite \textbf{décroissante} lorsque $\forall(x,y)\in X^2,\ x\leqslant y\implies f(x)\geqslant f(y)$
		\item $f$ est dite \textbf{strictement croissante} lorsque $\forall(x,y)\in X^2,\ x< y\implies f(x)< f(y)$
		\item $f$ est dite \textbf{strictement décroissante} lorsque $\forall(x,y)\in X^2,\ x< y\implies f(x)> f(y)$
		\item $f$ est dite \textbf{monotone} lorsqu'elle est croissante ou décroissante
		\item $f$ est dite \textbf{strictement monotone} lorsqu'elle est strictement croissante ou strictement décroissante
	\end{itemize}
\end{dfn}
\begin{exa}
	Si $X=\emptyset$, $f$ est la fois croissante, décroissante, strictement croissante et strictement décroissante par proposition vide. Si $X$ est un singleton, $f$ est croissante et décroissante car n'importe quoi implique le vrai, et $f$ est strictement croissante et strictement décroissante car le faux implique n'importe quoi.
\end{exa}
\begin{dfn}[Propriété sur]
	Soit $f:X\rightarrow\R$ et $Y\subset X$. On dit que $f$ \textbf{possède la propriété $P$ sur $Y$} lorsque $f_{|Y}$ possède la propriété $P$.
\end{dfn}
\begin{dfn}[Fonctions paires et impaires]
	Supposons que $X$ soit centré en $0$ ($ie$ que $\forall x\in X,\ -x\in X$). Soit $f:X\rightarrow\R$. On dit que $f$ est \textbf{paire} lorsque $$\forall x\in X,\ f(-x)=f(x)$$ On dit que $f$ est \textbf{impaire} lorsque $$\forall x\in X,\ f(-x)=-f(x)$$ L'ensemble des fonctions paires et l'ensemble des fonctions impaires de $X$ dans $\R$ forment des sev supplémentaires de $\R^X$.
\end{dfn}
\begin{rmq}
	Soit $a\in\R$ et supposons que $\forall x\in X,\ 2a-x\in X$. On peut alors dire que $f$ est paire par rapport à $a$ lorsque $\forall x\in X,\ f(2a-x)=f(x)$ et que $f$ est impaire par rapport à $a$ lorsque $\forall x\in X,\ f(2a-x)=-f(x)$. Là encore, ces ensembles forment des sev supplémentaires de $\R^X$.
\end{rmq}
\begin{dfn}[Périodicité]
	Soit $T>0$ et supposons que $\forall x\in X,\ x+t\in X$. Soit $f:X\rightarrow\R$. $f$ est dite \textbf{$T$-périodique} lorsque $$\forall x\in X,\ f(x+T)=f(x)$$
\end{dfn}
\begin{exa}
	Les fonctions $T$-périodiques forment une sous-algèbre de $\R^X$.
\end{exa}
\begin{dfn}[Fonctions lipschitziennes]
	Soit $\lambda\geqslant0$ et $f:X\rightarrow\R$. $f$ est dite \textbf{$\lambda$-lipschitzienne} lorsque $$\forall(x,y)\in X^2,\ \lvert f(x)-f(y)\rvert\leqslant\lambda\lvert x-y\rvert$$ Plus généralement, $f$ est dite \textbf{lipschitzienne} lorsqu'il existe $\lambda\geqslant0$ tel que $f$ soit $\lambda$-lipschitzienne. 
\end{dfn}
\begin{exa}
	Une fonction $0$-lipschitzienne est constante.
\end{exa}
\begin{rmq}
	La lipschitziannité est une condition très forte.
\end{rmq}
\begin{prop}[Interprétation géométrique des fonctions lipschitziennes]
	Supposons que $f$ soit $\lambda$-lipschitzienne. Imaginons un double cône de pentes $\lambda$ et $-\lambda$ qui coulisse sur la courbe. Alors, la courbe de $f$ se trouve toujours dans les parties Ouest et Est du double cône, peu importe la position de celui-ci.
\end{prop}
\begin{exa}
	La fonction racine carrée n'est pas lipschitzienne.
\end{exa}
\section{Étude locale d'une fonction}
Dans tout le paragraphe, on fixe $X\subset\R$.
\subsection{Limites}
\begin{dfn}[Limite d'une fonction]
	Soit $f:X\rightarrow\R$, $a\in\overline{X}$ et $b\in\overline{\R}$. On dit que \textbf{$f(x)$ tend vers $b$ quand $x$ tend vers $a$}, et on note $f(x)\xrightarrow[x\to a]{}b$, lorsque : $$\forall V_b\in\mathcal{V}(b),\ \exists V_a\in\mathcal{V}(a),\ \forall x\in X,\ x\in V_a\implies f(x)\in V_b$$
\end{dfn}
\begin{exa}[Fonctions constantes]
	Une fonction constante égale à $C$ tend vers $C$ en tout point de son adhérence.
\end{exa}
\begin{prop}
	Supposons que $(a,b)\in\R^2$. Les $4$ assertions suivantes sont équivalentes :
	\begin{enumerate}
		\item $f(x)\xrightarrow[x\to a]{}b$
		\item $f(x)-b\xrightarrow[x\to a]{}0$
		\item $\lvert f(x)-b\rvert\xrightarrow[x\to a]{}0$
		\item $f(a+h)\xrightarrow[h\to 0]{}b$
	\end{enumerate}
\end{prop}
\begin{dfn}
	La notation $f(x)\xrightarrow[x\to a]{}0^+$ signifie que $f(x)\xrightarrow[x\to a]{}0$ et $\exists V\in\mathcal{V}(a),\ \forall x\in X\cap V,\ f(x)>0$. La notation $f(x)\xrightarrow[x\to a]{}0^-$ signifie que $f(x)\xrightarrow[x\to a]{}0$ et $\exists V\in\mathcal{V}(a),\ \forall x\in X\cap V,\ f(x)<0$.
\end{dfn}
\begin{lem}[$\overline{\R}$ est séparé]
	On a : $$\forall(b,b')\in\overline{\R}^2,\ b\ne b'\implies\left(\exists(V,V')\in\mathcal{V}(b)\times\mathcal{V}(b'),\ V\cap V'=\emptyset\right)$$
\end{lem}
\begin{thm}[Unicité de la limité]
	On a : $$\forall a\in\overline{X},\ \forall(b,b')\in\overline{\R}^2,\ (f(x)\xrightarrow[x\to a]{}b\land f(x)\xrightarrow[x\to a]{}b')\implies b=b'$$ La limite d'une fonction en un point, si elle existe, est donc unique.
\end{thm}
\begin{dfn}[Limite d'une restriction]
	Soit $f:X\rightarrow\R$, $Y\subset\R$ et $a\in\overline{X\cap Y}$. Posons $X\st:=X\cap Y$  et $f\st:=f_{|X\st}$. Soit $b\in\overline{\R}$. Si $f\st(x)\xrightarrow[x\to a]{}b$, alors on note symboliquement $$f(x)\xrightarrow[\substack{x\to a\\ x\in Y}]{}b$$ On dit que $b$ est la \textbf{limite de $f$ en $a$ relativement à $Y$}. Voici trois cas particuliers : 
	\begin{itemize}
		\item $Y=]a,+\infty[$ On parle de \textbf{limite à droite} et on note $f(x)\xrightarrow[\substack{x\to a\\ x>a}]{}b$ ou $f(x)\xrightarrow[x\to a^+]{}b$, voire $f(a^+)=b$.
		\item $Y=]-\infty,a[$ On parle de \textbf{limite à gauche} et on note $f(x)\xrightarrow[\substack{x\to a\\ x<a}]{}b$ ou $f(x)\xrightarrow[x\to a^-]{}b$, voire $f(a^-)=b$.
		\item $Y=\R\setminus\left\{a\right\}$ On parle de \textbf{limite épointée} et on note $f(x)\xrightarrow[\substack{x\to a\\ x\ne a}]{}b$.
	\end{itemize}
\end{dfn}
\begin{rmq}
	Attention : dans le cas des limites à droite, à gauche et épointée, le point $a$ n'appartient pas à $X\st$, donc $f(a)$ peut exister et être différente de $b$. On pourra considérer la limite épointée en $0$ de l'indicatrice de $\left\{0\right\}$.
\end{rmq}
\begin{prop}[Limite induite]
	Si $f(x)\xrightarrow[x\to a]{}b$, alors $f(x)\xrightarrow[\substack{x\to a\\ x\in Y}]{}b$.
\end{prop}
\begin{rmq}
	Attention, la réciproque est fausse en générale : par exemple, si on a des limites à gauche et à droite égales, on ne peut en déduire que la limite épointée.
\end{rmq}
\begin{prop}[Réciproque partielle]
	Si $Y\in\mathcal{V}(a)$ et $f(x)\xrightarrow[\substack{x\to a\\ x\in Y}]{}b$, alors $f(x)\xrightarrow[x\to a]{}b$.
\end{prop}
\begin{exa}
	Si $f(x)\xrightarrow[\substack{x\to a\\ x\in Y_1}]{}b$, ... et $f(x)\xrightarrow[\substack{x\to a\\ x\in Y_n}]{}b$, alors $f(x)\xrightarrow[\substack{x\to a\\ x\in Y_1\cup\ldots\cup Y_n}]{}b$
\end{exa}
\begin{met}[Montrer que $f(x)$ ne tend pas vers $b$]
	Pour montrer que $f(x)$ ne tend pas vers $b$ lorsque $x$ tend vers $a$, il suffit d'exhiber $Y\subset\R$ tel que $a\in\overline{X\cap Y}$, ainsi que $b'\ne b$ tels que $f(x)\xrightarrow[\substack{x\to a\\ x\in Y}]{}b'$.
\end{met}
\begin{met}[Montrer que $f$ n'admet pas de limite en $a$]
	Pour montrer que $f$ n'admet pas de limite en $a$, il suffit d'exhiber $Y_1,Y_2\subset\R$ tels que $a\in\overline{X\cap Y_1}$ et $a\in\overline{X\cap Y_2}$ et $b_1\ne b_2$ tels que $f(x)\xrightarrow[\substack{x\to a\\ x\in Y_1}]{}b_1$ et $f(x)\xrightarrow[\substack{x\to a\\ x\in Y_2}]{}b$.
\end{met}
\begin{thm}[Caractère local de la limite]
	Supposons que $Y\in\mathcal{V}(a)$. Alors : $$f(x)\xrightarrow[x\to a]{}b \iff f(x)\xrightarrow[\substack{x\to a\\ x\in Y}]{}b$$
\end{thm}
\begin{exa}
	Si $f$ et $g$ coïncident au voisinage de $a$, alors elles ont le même comportement asymptotique en $a$.
\end{exa}
\begin{thm}[Caractérisation séquentielle de la limite]
	Soit $f:X\rightarrow\R$, $a\in\overline{X}$ et $b\in\overline{\R}$. Alors $$f(x)\xrightarrow[x\to a]{}b\iff \forall(x_n)\in X^\N,\ x_n\xrightarrow[n\to+\infty]{}a\implies f(x_n)\xrightarrow[n\to +\infty]{}b$$
\end{thm}
\begin{met}[Variante : montrer que $f(x)$ ne tend pas vers $b$]
	Pour montrer que $f(x)$ ne tend pas vers $b$ lorsque $x$ tend vers $a$, il suffit d'exhiber $(x_n)\in X^\N$ telle que $x_n\xrightarrow[n\to+\infty]{}a$ et $b'\ne b$ tels que $f(x_n)\xrightarrow[n\to+\infty]{}b'$.
\end{met}
\begin{met}[Variante : montrer que $f$ n'admet pas de limite en $a$]
	Pour montrer que $f$ n'admet pas de limite en $a$, il suffit d'exhiber $(x_n)\in X^\N$ et $(y_n)\in X^\N$ telles que $x_n\xrightarrow[n\to+\infty]{}a$ et $y_n\xrightarrow[n\to+\infty]{}a$ ainsi que $b\ne b'$ tels que $f(x_n)\xrightarrow[n\to+\infty]{}b$ et $f(y_n)\xrightarrow[n\to+\infty]{}b'$.
\end{met}
\begin{exa}
	$\cos$ n'admet pas de limite en $+\infty$. En effet, il suffit de considérer les suites $(2\pi n)_{n\in\N}$ et $(\pi+2\pi n)_{n\in\N}$
\end{exa}
\begin{prop}[Combinaison linéaire de limites]
	Soit $f,g:X\rightarrow\R$, $a\in\overline{X}$ et $b_1,b_2\in\overline{\R}$ tels que $f(x)\xrightarrow[x\to a]{}b_1$ et $g(x)\xrightarrow[x\to a]{}b_2$. Soit $\lambda,\mu\in\R$. Pourvu qu'on n'ait pas de forme indéterminée, on a : $$\lambda f(x)+\mu g(x)\xrightarrow[x\to a]{}\lambda b_1+\mu b_2$$
\end{prop}
\begin{prop}[Produit de limites]
	Soit $f,g:X\rightarrow\R$, $a\in\overline{X}$ et $b_1,b_2\in\overline{\R}$ tels que $f(x)\xrightarrow[x\to a]{}b_1$ et $g(x)\xrightarrow[x\to a]{}b_2$. Pourvu qu'on n'ait pas de forme indéterminée, on a : $$f(x)g(x)\xrightarrow[x\to a]{}b_1b_2$$
\end{prop}
\begin{prop}[Quotient de limites]
	Soit $f,g:X\rightarrow\R$, $a\in\overline{X}$ et $b_1,b_2\in\overline{\R}$ tels que $f(x)\xrightarrow[x\to a]{}b_1$ et $g(x)\xrightarrow[x\to a]{}b_2$. Supposons que $\forall x\in X,\ g(x)\ne0$ et $b_2\ne0$. Pourvu qu'on n'ait pas de forme indéterminée, on a : $$\frac{f(x)}{g(x)}\xrightarrow[x\to a]{}\frac{b_1}{b_2}$$ On a les raffinements habituels si $g$ ne s'annule pas au voisinage de $a$, selon le signe de $b_1$ et si $g(x)\xrightarrow[x\to a]{}0^+$ ou $g(x)\xrightarrow[x\to a]{}0^-$.
\end{prop}
\begin{prop}[Valeur absolue de limites]
	Soit $f:X\rightarrow\R$, $a\in\overline{X}$ et $b\in\overline{\R}$ tels que $f(x)\xrightarrow[x\to a]{}b$. Alors on a : $$\lvert f(x)\rvert \xrightarrow[x\to a]{}\lvert b\rvert$$
\end{prop}
\begin{prop}[Composition de limites]
	Soit $f:X\rightarrow\R$, $a\in\overline{X}$, $Y\subset\R$ tel que $f(X)\subset Y$, $g:Y\rightarrow\R$, $b\in \overline{Y}$ et $c\in\overline{\R}$. Si $f(x)\xrightarrow[x\to a]{}b$ et $g(y)\xrightarrow[y\to b]{}c$, alors : $$\left(g\circ f\right)(x)\xrightarrow[x\to a]{}c$$
\end{prop}
\begin{exa}
	Par caractérisation séquentielle, on obtient la limite de $x$ quand $x$ tend vers $+\infty$ (ou $-\infty$). On en déduit par ls propositions qui précèdent les limites des fonctions polynomiales et des fonctions rationnelles.
\end{exa}
\begin{thm}[Théorème d'encadrement/sandwich/des gendarmes]
	Soit $f,g,h:X\rightarrow\R$ telles que $f\leqslant g\leqslant h$. Soit $a\in\overline{X}$ et $b\in\overline{\R}$. Si $f(x)\xrightarrow[x\to a]{}b$ et $h(x)\xrightarrow[x\to a]{}b$, alors $g(x)\xrightarrow[x\to a]{}b$.
\end{thm}
\begin{cor}
	Soit $f,g:X\rightarrow\R$ telles que $\forall x\in X,\ \lvert f(x)\rvert\leqslant g(x)$ et $a\in\overline{X}$. Si $g(x)\xrightarrow[x\to a]{}0$, alors $f(x)\xrightarrow[x\to a]{}0$.
\end{cor}
\begin{cor}
	Soit $f,g:X\rightarrow\R$ avec $g$ bornée. Si $f(x)\xrightarrow[x\to a]{}0$, alors $f(x)g(x)\xrightarrow[x\to a]{}0$.
\end{cor}
\begin{thm}[Théorème de comparaison]
	Soit $f,g:X\rightarrow\R$ telles que $f\leqslant g$ et $a\in\overline{X}$.
	\begin{itemize}
		\item Principe de minoration : si $f(x)\xrightarrow[x\to a]{}+\infty$, alors $g(x)\xrightarrow[x\to a]{}+\infty$.
		\item Principe de majoration : si $g(x)\xrightarrow[x\to a]{}-\infty$, alors $f(x)\xrightarrow[x\to a]{}-\infty$.
	\end{itemize}
\end{thm}
\begin{thm}[Passage à la limite dans une inégalité large / PLIL]
	Soit $f,g:X\rightarrow\R$, $a\in\overline{X}$ et $b,b'\in\overline{\R}$ tels que $f(x)\xrightarrow[x\to a]{}b$ et $g(x)\xrightarrow[x\to a]{}b'$. Si $\forall x\in X,\ f(x)\leqslant g(x)$, alors $b\leqslant b'$.
\end{thm}
\begin{exa}[Cas particuliers fréquents]
	Si $\forall x\in X,\ f(x)\leqslant M$ avec $M$ indépendant de $x$, alors $b\leqslant M$. Si $\forall x\in X,\ m\leqslant g(x)$ avec $m$ indépendant de $x$, alors $m\leqslant b'$.
\end{exa}
\begin{rmq}
	Grâce au caractère local de la limite, tous les résultats précédents qui font appel à des inégalités sont encore vrais si les inégalités n'ont lieu qu'au voisinage de $a$.
\end{rmq}
\begin{thm}[Réciproque partielle du PLIL]
	Soit $f:X\rightarrow\R$, $a\in\overline{X}$ et $b\in\overline{\R}$ tels que $f(x)\xrightarrow[x\to a]{}b$. Soit $(\alpha,\beta)\in\R^2$.
	\begin{itemize}
		\item Si $b\in\R$, alors $f$ est bornée au voisinage de $a$.
		\item Si $b>\alpha$, alors $f(x)>\alpha$ au voisinage de $a$.
		\item Si $b<\beta$, alors $f(x)<\beta$ au voisinage de $a$.
	\end{itemize}
\end{thm}
\begin{thm}[Théorème de la limite monotone / TLM]
	Soit $(u,b)\in\overline{\R}^2$ tel que $u<v$ et $f:]u,v[\rightarrow\R$. Si $f$ est monotone, alors $f$ admet des limites en $u$ et en $v$. De manière plus précise : 
	\begin{itemize}
		\item Supposons que $f$ est croissante. Alors 
		\begin{itemize}
			\item[--] $f$ admet une limite $b\in\R\cup\left\{+\infty\right\}$ (dans $\R$ si $f$ est majorée, et valant $+\infty$ si $f$ n'est pas majorée) en $v$
			\item[--] $f$ admet une limite $b'\in\R\cup\left\{-\infty\right\}$ (dans $\R$ si $f$ est minorée, et valant $-\infty$ si $f$ n'est pas minorée) en $u$
		\end{itemize}
		\item Supposons que $f$ est décroissante. Alors 
		\begin{itemize}
			\item[--] $f$ admet une limite $b\in\R\cup\left\{-\infty\right\}$ (dans $\R$ si $f$ est minorée, et valant $-\infty$ si $f$ n'est pas minorée) en $v$
			\item[--] $f$ admet une limite $b'\in\R\cup\left\{+\infty\right\}$ (dans $\R$ si $f$ est majorée, et valant $+\infty$ si $f$ n'est pas majorée) en $u$
		\end{itemize}
	\end{itemize}
\end{thm}
\begin{exa}
	Soit $f:]0,1[\rightarrow\R$ croissante et majorée par un certain $M\in\R$. Alors elle admet une limite $l$ en $1$ d'après le TLM. Un PLIL montre que $l\leqslant M$. On montre que $\forall x\in]0,1[,\ f(x)\leqslant l$ en fixant $x$, en utilisant la croissance et en faisant un PLIL sur $y\geqslant x$. Si on suppose que $f$ est strictement croissante, on montre que $\forall x\in]0,1[,\ f(x)<l$ en fixant $y\in]x,1[$ et en utilisant le résultat qui précède sur $f(y)$ ainsi que la stricte croissance.
\end{exa}
\begin{cor}
	Soit $(u,b)\in\overline{\R}^2$ tel que $u<v$ et $f:]u,v[\rightarrow\R$. Si $f$ est monotone, alors en tout point $a$ de $]u,v[$, $f$ admet des limites à gauche et à droite qui encadrent $f(a)$. De manière précise, soit $a\in]u,v[$.
	\begin{itemize}
		\item Supposons que $f$ est croissante. Alors $f(a^-)$ et $f(a^+)$ existent et sont finies, et on a $f(a^-)\leqslant f(a)\leqslant f(a^+)$.
		\item Supposons que $f$ est décroissante. Alors $f(a^-)$ et $f(a^+)$ existent et sont finies, et on a $f(a^-)\geqslant f(a)\geqslant f(a^+)$.
	\end{itemize}
\end{cor}
\subsection{Continuité locale}
\begin{dfn}[Continuité en un point]
	Soit $f:X\rightarrow\R$ et $a\in X$. On dit que $f$ est \textbf{continue en $a$} lorsque $f(x)\xrightarrow[x\to a]{}f(a)$, autrement dit lorsque : $$\forall \varepsilon>0,\ \exists \delta>0,\ \forall x\in X,\ \lvert x-a\rvert\leqslant\delta\implies\lvert f(x)-f(a)\rvert\leqslant\varepsilon$$
\end{dfn}
\begin{exa}
	Une fonction constante est continue en chacun de ses points.
\end{exa}
\begin{exa}
	La fonction $\id_\R$ est continue en tout point de $\R$.
\end{exa}
\begin{prop}
	Soit $f:X\rightarrow\R$ et $a\in X$. Soit $b\in\overline{\R}$. Si $f(x)\xrightarrow[x\to a]{}b$, alors $b=f(a)$ si bien que $f$ est continue en $a$.
\end{prop}
\begin{rmq}
	Cela provient du fait que désormais, $a$ est dans $X$ donc ne peut plus être en bord d'intervalle, là d'où proviennent la majorité de nos contre-exemples.
\end{rmq}
\begin{dfn}[Continuité à gauche et à droite]
	Soit $f:X\rightarrow\R$ et $a\in X$. On dit que $f$ est \textbf{continue à gauche en $a$} lorsque $f(a^-)$ existe et vaut $f(a)$. On dit que $f$ est \textbf{continue à droite en $a$} lorsque $f(a^+)$ existe et vaut $f(a)$.
\end{dfn}
\begin{prop}
	Si $f$ est continue à gauche et à droite en $a$, alors $f$ est continue en $a$.
\end{prop}
\begin{exa}
	Soit $f:]u,v[\rightarrow\R$ croissante et $a\in]u,v[$. On rappelle que $f$ admet des limites finies à gauche et à droite en $a$ qui vérifient $f(a^-)\leqslant f(a)\leqslant f(a^+)$. Pour montrer que $f$ est continue en $a$, il suffit alors de montrer que $f(a^-)\geqslant f(a)$ et $f(a)\geqslant f(a^+)$.
\end{exa}
\begin{rmq}
	On peut adapter la proposition en bord de domaine, si une seule des deux limites existe. Elle reste alors vraie.
\end{rmq}
\begin{thm}[Caractérisation séquentielle de la continuité]
	Soit $f:X\rightarrow\R$ et $a\in X$. $f$ est continue en $a$ si, et seulement si, $$\forall (x_n)_{n\in\N}\in X^\N,\ x_n\xrightarrow[n\to+\infty]{}a\implies f(x_n)\xrightarrow[n\to+\infty]{}f(a)$$
\end{thm}
\begin{exa}
	Les suites récurrentes de la forme $u_{n+1}=f(u_n)$ avec $f$ continue. Voir le dernier paragraphe du chapitre "Réels et suites".
\end{exa}
\begin{dfn}[Prolongement par continuité]
	Soit $f:X\rightarrow\R$, $a\in \R$ et $b\in\R$ tels que $a\in\overline{X}\setminus X$ et $f(x)\xrightarrow[x\to a]{}b$. Alors $f$ admet un unique prolongement à $X\cup\left\{a\right\}$ qui soit continu en $a$. C'est la fonction : $$\begin{array}{ccccc}f\st&:&X\cup\left\{a\right\}&\rightarrow&\R\\ &&x&\mapsto&\begin{cases}
			f(x)\text{ si }x\in X \\ b\text{ si }x=a
	\end{cases}\end{array}$$ On l'appelle \textbf{le prolongement par continuité de $f$ en $a$}.
\end{dfn}
\begin{prop}[Combinaison linéaire]
	Soit $f,g:X\rightarrow\R$ et $a\in X$ tels que $f$ et $g$ soient continues en $a$. Soit $\lambda,\mu\in\R$. Alors $\lambda f+\mu g$ est continue en $a$.
\end{prop}
\begin{prop}[Produit]
	Soit $f,g:X\rightarrow\R$ et $a\in X$ tels que $f$ et $g$ soient continues en $a$. Alors $fg$ est continue en $a$.
\end{prop}
\begin{prop}[Quotient]
	Soit $f,g:X\rightarrow\R$ et $a\in X$ tels que $f$ et $g$ soient continues en $a$. Si $\forall x\in X,\ g(x)\ne0$, alors $\displaystyle\frac{f}{g}$ est continue en $a$.
\end{prop}
\begin{prop}[Valeur absolue]
	Soit $f:X\rightarrow\R$ et $a\in X$ tels que $f$ soient continue en $a$. Alors $\lvert f\rvert$ est continue en $a$.
\end{prop}
\begin{prop}[Minimum et maximum]
	Soit $f,g:X\rightarrow\R$ et $a\in X$ tels que $f$ et $g$ soient continues en $a$. Alors $x\mapsto\min(f(x),g(x))$ et $x\mapsto\max(f(x),g(x))$ sont continues en $a$.
\end{prop}
\begin{thm}[Composée]
	Soit $f:X\rightarrow\R$, $a\in X$, $Y\subset\R$ tel que $f(X)\subset Y$ et $g:Y\rightarrow\R$. Si $f$ est continue en $a$ et $g$ est continue en $f(a)$, alors $g\circ f$ est continue en $a$.
\end{thm}
\section{Extension globale}
Dans tout le paragraphe, on fixe $X\subset\R$.
\subsection{Continuité}
\begin{dfn}[Continuité globale]
	Soit $f:X\rightarrow\R$. $f$ est dite \textbf{continue} lorsque $f$ est continue en tout point de $X$.
\end{dfn}
\begin{dfn}[Continuité sur]
	Soit $f:X\rightarrow\R$ et $A\subset X$. $f$ est dite \textbf{continue sur $A$} lorsque $f_{|A}$ est continue.
\end{dfn}
\begin{prop}[Continuité induite]
	Soit $f:X\rightarrow\R$ et $A\subset X$. Si $f$ est continue, alors $f$ est continue sur $A$.
\end{prop}
\begin{rmq}
	Attention : la réciproque est fausse ! Par exemple, la partie entière est continue sur $[0,1[$ mais pas en tout point de $[0,1[$. Elle est discontinue en $0$. Toutefois, il suffit $A$ soit un intervalle ouvert pour que la réciproque soit vraie.
\end{rmq}
\begin{exa}
	Si $f$ est continue et admet un prolongement par continuité, alors ce prolongement est continu.
\end{exa}
\begin{thm}
	La continuité globale passe aux opérations algébriques suivantes : 
	\begin{itemize}
		\item Combinaison linéaire
		\item Produit
		\item Quotient dont le dénominateur ne s'annule pas
		\item Valeur absolue
		\item Minimum et maximum
		\item Composition
	\end{itemize}
\end{thm}
\begin{exa}
	$\mathcal{C}^0(X,\R)$ (l'ensemble des fonctions de $X$ dans $\R$ continues) est une sous-algèbre de $\R^X$.
\end{exa}
\begin{exa}
	Toute fonction polynomiale est continue. Toute fonction rationnelle est continue.
\end{exa}
\begin{dfn}[Uniforme continuité]
	Soit $f:X\rightarrow\R$. $f$ est dite \textbf{uniformément continue} lorsqu'on "peut choisir $\delta$ indépendamment de $a$" : $$\forall\varepsilon>0,\ \exists\delta>0,\ \forall(x,y)\i X^2,\ \lvert x-y\rvert\leqslant\delta\implies\lvert f(x)-f(y)\leqslant\varepsilon$$
\end{dfn}
\begin{thm}[Lipschitzienne $\implies$ Uniformément continue $\implies$ Continue]
	Si $f$ est lipschitzienne, alors est elle uniformément continue. Si $f$ est uniformément continue, alors elle est continue.
\end{thm}
\begin{exa}
	Toute fonction constante est $0$-lipschitzienne, donc uniformément continue?
\end{exa}
\begin{exa}
	D'après le seconde inégalité triangulaire, la valeur absolue est $1$-lipschitzienne donc elle est uniformément continue puis continue.
\end{exa}
\begin{rmq}
	Attention, les réciproques sont fausses. Par exemple, la racine carrée restreinte à $[0,1]$ est uniformément continue par le théorème de Heine (cf. la suite du chapitre), mais n'est pas lipschitzienne. Ensuite, la fonction carré est continue mais n'est pas uniformément continue.
\end{rmq}
\begin{thm}[Théorème de Heine]
	"Toute fonction continue sur un segment est uniformément continue." Soit $(a,b)\in\R^2$ tel que $a\leqslant b$ et $f:[a,b]\rightarrow\R$. Si $f$ est continue, alors $f$ est uniformément continue.
\end{thm}
\subsection{Théorèmes généraux}
\begin{thm}[Théorème des valeurs intermédiaires / TVI, version segment]
	Soit $(a,b)\in\R^2$ tel que $a\leqslant b$. Soit $f:[a,b]\rightarrow\R$ continue. Alors, $f$ "prend toutes les valeurs entre $f(a)$ et $f(b)$". Précisément : 
	\begin{itemize}
		\item Supposons que $f(a)\leqslant f(b)$. Alors $$\forall y\in\left[f(a),f(b)\right],\ \exists x\in[a,b],\ f(x)=y$$
		\item Supposons que $f(a)\geqslant f(b)$. Alors $$\forall y\in\left[f(b),f(a)\right],\ \exists x\in[a,b],\ f(x)=y$$
	\end{itemize}
\end{thm}
\begin{rmq}
	Bien que non constructif, ce théorème permet de prouver une surjectivité.
\end{rmq}
\begin{cor}[Variantes du TVI]
	On a les variantes suivantes du TVI : \begin{itemize}
		\item Soit $(a,b)\in\overline{\R}^2$ tel que $-\infty<a<b\leqslant +\infty$ et $f:[a,b[\rightarrow\R$ continue. Soit $l_b\in\overline{\R}$ tel que $f(x)\xrightarrow[x\to b]{}l_b$.
		\begin{itemize}
			\item[--] Supposons que $f(a)<l_b$. Alors $$\forall y\in[f(a),l_b[,\ \exists x\in[a,b[,\ y=f(x)$$
			\item[--] Supposons que $f(a)>l_b$. Alors $$\forall y\in]l_b,f(a)],\ \exists x\in[a,b[,\ y=f(x)$$
		\end{itemize}
		\item Soit $(a,b)\in\overline{\R}^2$ tel que $-\infty\leqslant a<b<+\infty$ et $f:]a,b]\rightarrow\R$ continue. Soit $l_a\in\overline{\R}$ tel que $f(x)\xrightarrow[x\to a]{}l_a$.
		\begin{itemize}
			\item[--] Supposons que $l_a<f(b)$. Alors $$\forall y\in]l_a,f(b)],\ \exists x\in]a,b],\ y=f(x)$$
			\item[--] Supposons que $l_a>f(b)$. Alors $$\forall y\in[f(b),l_a[,\ \exists x\in]a,b],\ y=f(x)$$
		\end{itemize}
		\item Soit $(a,b)\in\overline{\R}^2$ tel que $-\infty\leqslant a<b\leqslant +\infty$ et $f:]a,b[\rightarrow\R$ continue. Soit $l_b\in\overline{\R}$ tel que $f(x)\xrightarrow[x\to b]{}l_b$ et $l_a\in\overline{\R}$ tel que $f(x)\xrightarrow[x\to a]{}l_a$.
		\begin{itemize}
			\item[--] Supposons que $l_a<l_b$. Alors $$\forall y\in]l_a,l_b[,\ \exists x\in]a,b[,\ y=f(x)$$
			\item[--] Supposons que $l_a>l_b$. Alors $$\forall y\in]l_b,l_a[,\ \exists x\in]a,b[,\ y=f(x)$$
		\end{itemize}
	\end{itemize}
\end{cor}
\begin{rmq}
	Le TVI version segment et l'ensemble des ses variantes peuvent être appelés TVI.
\end{rmq}
\begin{cor}[Existence d'une racine réel pour les polynômes de degré impair]
	La troisième variante montre que toute fonction polynomiale de degré impair admet au moins une racine réelle.
\end{cor}
\begin{cor}[L'image d'un intervalle par une fonction continue est un intervalle]
	Soit $f:X\rightarrow\R$ continue et $I\subset X$ un intervalle. Alors $f(I)$ est un intervalle.
\end{cor}
\begin{rmq}
	Dans le cas où $f$ est continue et monotone, elle envoie les bornes de $I$ sur les bornes de $f(I)$, et $I$ et $f(I)$ ont alors même nature (ouvert, fermé, semi-ouvert). En pratique, un tableau de variations suffira pour justifier ce fait, mais on peut le montrer en faisant une double inclusion et en appliquant le TVI.
\end{rmq}
\begin{thm}[Théorème des bornes atteintes]
	Soit $(a,b)\in\R^2$ tel que $a\leqslant b$ et $f:[a,b]\rightarrow\R$. Si $f$ est continue, alors elle admet un minimum et un maximum. Autrement dit, toute fonction continue sur un segment est bornée et atteint ses bornes.
\end{thm}
\begin{rmq}
	On utilise plus souvent le fait que $f$ soit bornée que le fait que ses bornes soient atteintes, mais ce dernier fait peut toujours servir.
\end{rmq}
\begin{exa}
	En notant $m$ le minimum de $f$ et $M$ son maximum, on a par le TVI  $\im(f)=[m,M]$.
\end{exa}
\begin{thm}
	Soit $I$ un \textbf{intervalle}, $f:I\rightarrow\R$ une fonction \textbf{continue} et \textbf{injective}. Alors $f$ est strictement monotone.
\end{thm}
\begin{rmq}
	N'oublier aucune des trois hypothèses : un intervalle, continue et injective ! Dès que l'on enlève une des trois hypothèses, le théorème tombe en défaut.
\end{rmq}
\begin{thm}[Théorème de la bijection réciproque / TBR, cas continu]
	Soit $I$ et $J$ des intervalles de $\R$ non triviaux et $f:I\rightarrow J$ une bijection. Si $f$ est strictement monotone, alors $f\inv$ est continue.
\end{thm}
\begin{exa}
	La fonction racine carrée est donc continue en appliquant le TBR à la fonction carré.
\end{exa}
\begin{exa}
	En appliquant le théorème à $f\inv$ qui a les mêmes variations que $f$, si la bijection $f$ est strictement monotone, alors elle est continue.
\end{exa}
\section{Fonctions à valeurs complexes}
On garde $X\subset\R$.
\begin{dfn}[Fonctions de base]
	Soit $f:X\rightarrow\C$. On peut considérer : $$\begin{array}{ccccc}\text{Re}(f)&:&X&\rightarrow&\R\\&&x&\mapsto&\text{Re}(f(x))\end{array}$$ $$\begin{array}{ccccc}\text{Im}(f)&:&X&\rightarrow&\R\\&&x&\mapsto&\text{Im}(f(x))\end{array}$$
	$$\begin{array}{ccccc}\lvert f\rvert&:&X&\rightarrow&\R\\&&x&\mapsto&\lvert f(x)\rvert\end{array}$$
	$$\begin{array}{ccccc}\overline{f}&:&X&\rightarrow&\R\\&&x&\mapsto&\overline{f(x)}\end{array}$$
\end{dfn}
\begin{dfn}[Fonctions à valeurs complexes bornée]
	$f:X\rightarrow\C$ est dite \textbf{bornée} lorsque $\lvert f\rvert$ est majorée (au sens de $\R$).
\end{dfn}
\begin{rmq}
	Lorsque $f$ est à valeurs réelles, cela coïncide bien puisqu'alors $f$ est bornée si, et seulement si $\lvert f\rvert$ est majorée, et car le module prolonge la valeur absolue.
\end{rmq}
\begin{dfn}[Limite d'une fonction à valeurs complexes]
	Soit $a\in\overline{X}$, $b\in\C$ et $f:X\rightarrow\C$. On dit que $f(x)$ tend vers $b$ quand $x$ tend vers $a$, et on note $f(x)\xrightarrow[x\to a]{}b$ lorsque $$\forall\varepsilon>0,\ \exists V_a\in\mathcal{V}(a),\ \forall x\in X,\ x\in V_a\implies \lvert f(x)-b\rvert\leqslant\varepsilon$$
\end{dfn}
\begin{rmq}
	Attention : pour les fonctions à valeurs complexes, pas de limite infinie, $+\infty$ n'a pas de sens dans $\C$ (sauf si la fonction est à valeurs réelles bien sûr).
\end{rmq}
\begin{prop}
	Toute fonction à valeurs complexes qui tend vers $b\in\C$ lorsque $x$ tend vers $a\in\overline{X}$ est bornée au voisinage de $a$.
\end{prop}
\begin{thm}[Théorème passerelle]
	Soit $a\in\overline{X}$, $(u,v)\in\R^2$ et $f:C\rightarrow\C$. On a : $$f(x)\xrightarrow[x\to a]{}u+iv\iff\begin{cases}
		\text{Re}(f)(x)\xrightarrow[x\to a]{}u \\
		\text{Im}(f)(x)\xrightarrow[x\to a]{}v
	\end{cases}$$
\end{thm}
\begin{rmq}
	On déduit de ce théorème : \begin{itemize}
		\item \textbf{L'unicité de la limite}, quand elle existe : il suffit de passer par le cas réel et l'unicité de la forme algébrique.
		\item La \textbf{caractérisation séquentielle complexe de la limite} : il suffit de passer par la caractérisation séquentielle réelle et le théorème passerelle pour les suites complexes.
	\end{itemize}
\end{rmq}
\begin{dfn}[Continuité locale dans $\C$]
	Soit $a\in X$ et $f:X\rightarrow\C$. $f$ est dite \textbf{continue en $a$} lorsque $f(x)\xrightarrow[x\to a]{}f(a)$.
\end{dfn}
\begin{rmq}
	Cette définition prolonge celle donnée dans $\R$.
\end{rmq}
\begin{rmq}[Utile]
	S'il existe $b\in\C$ tel que $f(x)\xrightarrow[x\to a]{}b$, alors on a nécessairement $b=f(a)$, si bien que $f$ est continue en $a$. Il suffit d'appliquer la proposition analogue du cas réel et le théorème passerelle.
\end{rmq}
\begin{rmq}
	On prouve aisément les deux résultats suivants : 
	\begin{itemize}
		\item \textbf{Caractérisation séquentielle complexe de la continuité}
		\item \textbf{Théorème passerelle pour la continuité locale} : $f$ est continue en $a$ si, et seulement si, $\text{Re}(f)$ et $\text{Im}(f)$ sont continues en $a$.
	\end{itemize}
\end{rmq}
\begin{dfn}[Continuité globale dans $\C$]
	$f:X\rightarrow\C$ est dite \textbf{continue} lorsqu'elle est continue en tout point de $X$.
\end{dfn}
\begin{rmq}
	Cette définition prolonge celle donnée dans $\R$.
\end{rmq}
\begin{rmq}
	On prouve aisément le \textbf{théorème passerelle pour la continuité globale} : $f$ est continue si, et seulement si, $\text{Re}(f)$ et $\text{Im}(f)$ sont continues.
\end{rmq}
\begin{thm}[Opérations algébriques sur les fonctions à valeurs complexes]
	Les opérations algébriques suivantes sont valables pour les limites, la continuité locale et la continuité globale : 
	\begin{itemize}
		\item Combinaison linéaire
		\item Produit
		\item Quotient dont le dénominateur ne s'annule pas
		\item Module
		\item Conjugué
		\item Composition
	\end{itemize}
\end{thm}
\begin{rmq}
	Attention : pour la composition, seule la "dernière" fonction par laquelle on compose par la gauche peut être à valeurs complexes !
\end{rmq}

\newpage
\chapter{Dérivation}
$I$ est toujours un intervalle non trivial de $\R$ et sauf précision du contraire, $f$ est une fonction de $I$ dans $\R$
\section{Dérivation locale}
\begin{dfn}[Taux d'accroissement en un point]
	Soit $a\in I$. Le taux d'accroissement de $f$ en $a$ est la fonction
	$$
	\begin{array}[t]{lrcl}
		\tau_{a,f}: & I\setminus \left\{a\right\} & \longrightarrow & \R \\
		& x & \longmapsto & \displaystyle \frac{f(x)-f(a)}{x-a}
	\end{array}
	$$
\end{dfn}
\begin{dfn}[Dérivabilité en un point et nombre dérivé]
	Lorsque $\tau_{a,f}$ admet une limite finie $l$ en $a$, $f$ est dite dérivable en $a$. $l$ est appelé le nombre dérivé de $f$ en $a$ et on le note $f'(a)$. On dit que $f$ est dérivable à droite (resp. à gauche) en $a$ lorsque c$\tau_{a,f}$ admet une limite finie à droite (resp. à gauche). Cette limite est appelée nombre dérivé de $f$ en $a$ à droite (resp. à gauche) et est notée $f_d'(a)$ (resp. $f_g'(a)$)
\end{dfn}
\begin{prop}[Caractère local de la dérivabilité]
	Soit $f:I\rightarrow\R$ et $g:J\rightarrow\R$ et $a\in I\cap J$. Supposons que $f$ et $g$ coïncident au voisinage de $a$. Alors $f$ est dérivable en $a$ si, et seulement si, $g$ est dérivable en $a$. Dans ce cas, $g'(a)=f'(a)$.
\end{prop}
\begin{dfn}[Équation de la tangente]
	$$y-f(a)=f'(a)(x-a)$$
\end{dfn}
\begin{prop}[Dérivable $\Rightarrow$ continue]
	Si $f$ est dérivable en $a$, alors $f$ est continue en $a$.
\end{prop}
\begin{rmq} La réciproque est fausse. Considérer la valeur absolue en 0. \end{rmq}
\begin{thm}[Condition nécessaire d'extremum local]
	Soit $a\in\mathring{I}$ et $f$ dérivable en $a$. \textbf{Si} $f$ admet un extremum local en $a$, \textbf{alors} $f'(a)=0$. 
\end{thm}
\begin{rmq} La réciproque est fausse. Considérer la fonction cube en 0. \end{rmq}
\begin{thm}[Opérations algébriques sur la dérivabilité en un point]
	La dérivabilité en un point passe aux combinaisons linéaires, au produit, à l'inverse et au quotient si le dénominateur ne s'annule pas, et à la composition.
\end{thm}
\begin{thm}[Théorème de la bijection réciproque, cas dérivable en un point]
Soit $I$ et $J$ des intervalles non triviaux de $\R$ et $f:I\rightarrow J$ une bijection strictement monotone. Soit $a$ dans $I$. Supposons que $f$ est dérivable en $a$. Alors $f\inv$ est dérivable en $f(a)$ si, et seulement si, $f'(a)\ne 0$. Dans ce cas, on a $\displaystyle (f\inv)'(f(a))=\frac{1}{f'(a)}$
\end{thm}
\section{Dérivation globale}
\begin{dfn}[Dérivabilité globale]
	$f:I\rightarrow\R$ est dite dérivable lorsqu'elle est dérivable en tout point de $I$. On définit alors $f'$ la \textbf{fonction dérivée de $f$} qui à tout $x$ de $I$ associe $f'(x)$. On note parfois cette fonction $\displaystyle \frac{df}{dx}$, voire $Df$.
\end{dfn}
\begin{rmq} Plus généralement, on s'autorise à considérer des fonctions définies sur des unions d'INT. Les définitions restent les mêmes. Exemple : la fonction tangente. \end{rmq}
\begin{cor}[Dérivable $\Rightarrow$ continue]
	Toute fonction dérivable est continue
\end{cor}
\begin{dfn}[Dérivabilité sur]
	Soit $f:I\rightarrow\R$ et $J\subset I$ un autre INT. $f$ est dite dérivable sur $J$ lorsque la restriction de $f$ à $J$ est dérivable.
\end{dfn}
\begin{thm}[Opérations algébriques sur la dérivabilité globale]
	La dérivabilité globale passe aux opérations algébriques suivantes (sous réserve que les fonctions considérées soient dérivables et compatibles):
	\begin{enumerate}
		\item Combinaison linéaire : $(\lambda f +\mu g)'=\lambda f'+\mu g'$
		\item Produit : $(fg)'=f'g+fg'$
		\item Inverse (si $g$ ne s'annule pas) : $\displaystyle \left(\frac1f\right)'=-\frac{f'}{f^2}$
		\item Quotient (si $g$ ne s'annule pas) : $\displaystyle \left(\frac{f}{g}\right)'=\frac{f'g-fg'}{g^2}$
		\item Composée (si les domaines sont compatibles) : $(g\circ f)'=(g'\circ f)\times f'$
	\end{enumerate}
\end{thm}
\begin{cor}
	Soit $f:I\rightarrow\R$ dérivable.
	\begin{enumerate}
		\item Soit $n\in\N\st$. Alors $f^n$ est dérivable et $(f^n)'=nf^{n-1}f'$
		\item Soit $n\in\Z$. Supposons que $f$ est à valeurs dans $\R\st$. Alors $f^n$ est dérivable et $(f^n)'=nf^{n-1}f'$
		\item Soit $\alpha\in\R$. Supposons que $f$ est à valeurs dans $\R_+\st$. Alors $f^\alpha$ est dérivable et $(f^n)'=\alpha f^{\alpha-1}f'$
	\end{enumerate}
\end{cor}
\begin{thm}[Théorème de la bijection réciproque, cas dérivable]
	Soit $I$ et $J$ des intervalles non triviaux de $\R$ et $f:I\rightarrow J$ une bijection strictement monotone. Supposons que $f$ est dérivable et que $f'$ ne s'annule pas. Alors $f\inv$ est dérivable et $$(f\inv)'=\frac{1}{f'\circ f\inv}$$
\end{thm}
\section{Théorèmes généraux}
\begin{thm}[Théorème de Rolle]
	Soit $(a,b)\in\R^2$ tel que $a<b$ et $f:[a,b]\rightarrow\R$ telle que 
	\begin{enumerate}
		\item $f$ est continue en tout point de $[a,b]$
		\item $f$ est dérivable en tout point de $]a,b[$
		\item $f(a)=f(b)$
	\end{enumerate}
	Alors, il existe $c\in]a,b[$ tel que $f'(c)=0$.
\end{thm}
\begin{rmq} Ce théorème est donc $a$ $fortiori$ vrai lorsque la fonction est dérivable. Ces conditions sont typiques d'une \textbf{fonction non injective}. \end{rmq}
\begin{thm}[Théorème des accroissements finis / Égalité des accroissements finis]
	Soit $(a,b)\in\R^2$ tel que $a<b$ et $f:[a,b]\rightarrow\R$ telle que 
	\begin{enumerate}
		\item $f$ est continue en tout point de $[a,b]$
		\item $f$ est dérivable en tout point de $]a,b[$
	\end{enumerate}
	Alors, il existe $c\in]a,b[$ tel que $\displaystyle f'(c)=\frac{f(b)-f(a)}{b-a}$.
\end{thm}
\begin{rmq} Là encore, ce résultat est $a$ $fortiori$ vrai si la fonction est dérivable en tout point de l'intervalle. \end{rmq}
\begin{cor}[Inégalité des accroissements finis]
	Soit $I$ un INT de $\R$, $\lambda \geqslant 0$ et $f:I\rightarrow\R$ telle que 
	\begin{enumerate}
		\item $f$ est continue en tout point de $I$
		\item $f$ est dérivable en tout point de $\mathring{I}$
		\item $\forall x\in\mathring{I}, \ |f'(x)|\leqslant \lambda$
	\end{enumerate}
	Alors, $f$ est $\lambda$-lipschitzienne.
\end{cor}
\begin{exa}
	Permet de montrer que $\forall x \in\R, \ |\sin(x)|\leqslant|x|$. Utilisé dans la méthode de Héron sur l'approximation des racines carrées avec $u_{n+1}=\frac12(u_n+\frac{k}{u_n})$
\end{exa}
\begin{thm}[Équivalences entre signe de la dérivée et variations]
	Soit $I$ un INT de $\R$ et $f:I\rightarrow\R$. Supposons que $f$ est continue en tout point de $I$ et dérivable en tout point de $\mathring{I}$. Alors :
	\begin{enumerate}
		\item $f$ est croissante si, et seulement si, $\forall x \in\mathring{I}, \ f'(x)\geqslant 0$
		\item $f$ est décroissante si, et seulement si, $\forall x \in\mathring{I}, \ f'(x)\leqslant 0$
		\item $f$ est constante si, et seulement si, $\forall x \in\mathring{I}, \ f'(x)= 0$
	\end{enumerate}
\end{thm}
\begin{cor}[Condition suffisante de stricte montonie]
	Soit $I$ un INT de $\R$ et $f:I\rightarrow\R$. Supposons que $f$ est continue en tout point de $I$ et dérivable en tout point de $\mathring{I}$. Alors :
	\begin{enumerate}
		\item \textbf{Si} $\forall x \in\mathring{I}, \ f'(x)\geqslant0$, \textbf{alors} $f$ est strictement croissante.
		\item \textbf{Si} $\forall x \in\mathring{I}, \ f'(x)\leqslant0$, \textbf{alors} $f$ est strictement décroissante.
	\end{enumerate}
\end{cor}
\begin{cor}[Raffinement de la CS]
	Soit $I$ un INT de $\R$ et $f:I\rightarrow\R$. Supposons que $f$ est continue en tout point de $I$ et dérivable en tout point de $\mathring{I}$. Alors :
	\begin{enumerate}
		\item\textbf{Si} $\forall x \in\mathring{I}, \ f'(x)\geqslant 0$ et le nombre de points d'annulation de $f'$ est \textbf{fini}, \textbf{alors} $f$ est strictement croissante.
		\item\textbf{Si} $\forall x \in\mathring{I}, \ f'(x)\leqslant 0$ et le nombre de points d'annulation de $f'$ est \textbf{fini}, \textbf{alors} $f$ est strictement décroissante.
	\end{enumerate}
\end{cor}
\begin{thm}[Théorème de la limite de la dérivée]
	Soit $I$ un INT de $\R$, $f:I\rightarrow\R$ et $a\in I$. Supposons que $f$ est continue en tout point de $I$ et que $f$ est dérivable en tout point de $I\setminus \left\{a\right\}$. \textbf{Si} il existe $\lambda\in\R$ tel que $f'(x)\xrightarrow[x\ne a]{} \lambda$, \textbf{alors} :
	\begin{enumerate}
		\item $f$ est dérivable en $a$
		\item $f'(a)=\lambda$
	\end{enumerate}
\end{thm}
\begin{rmq} Réciproquement, si $f$ est dérivable en $a$ de dérivée $\lambda$, \textbf{on ne peut pas} en déduire que $f'(x)\xrightarrow[x\ne a]{} \lambda$. Comme contre-exemple, on pourra considérer la fonction qui à $x$ associe $x^2\sin(\frac1x)$ prolongée par 0 en 0. \end{rmq}
\section{Fonctions de classe $\mathcal{D}^n$ et de classe $\mathcal{C}^n$}
\begin{dfn}[Classe $\mathcal{D}^n$]
	Soit $f:I\rightarrow\R$ et $n\in\N$. $f$ est dite de classe $\mathcal{D}^n$ lorsqu'il existe une suite finie de fonctions de $I$ dans $\R$ $(\varphi_i)_{0 \leqslant i \leqslant n}$ telle que $\varphi_0=f$, chacune des $\varphi_k$ soit dérivable pour $0\leqslant k \leqslant n-1$ et que pour tout $0 \leqslant k \leqslant n-1$, on ait $\varphi_k'=\varphi_{k+1}$. On note $\mathcal{D}^n(I,\R)$ l'ensemble des fonctions de $I$ dans $\R$ de classe $\mathcal{D}^n$.
\end{dfn}
\begin{dfn}[Classe $\mathcal{D}^\infty$]
	Soit $f:I\rightarrow\R$ et $n\in\N$. $f$ est dite de classe $\mathcal{D}^\infty$ lorsque $f$ est de classe $\mathcal{D}^n$ quel que soit $n \in \N$. On note $\mathcal{D}^\infty(I,\R)$ l'ensemble des fonctions de $I$ dans $\R$ de classe $\mathcal{D}^\infty$. On a alors : $$\mathcal{D}^\infty=\bigcap_{n\in\N}\mathcal{D}^n(I,\R)$$
\end{dfn}
\begin{rmq} On a la chaîne infinie d'inclusion suivante : $$\mathcal{D}^\infty(I, \R) \subset ... \subset \mathcal{D}^1(I,\R) \subset \mathcal{D}^0(I,\R)$$ \end{rmq}
\begin{dfn}[Dérivée n-ème]
	Soit $f\in\mathcal{D}^n(I,\R)$. On définit la dérivée n-ème de $f$, notée $f^{(n)}$, ou bien encore $\displaystyle \frac{d^nf}{dx^n}$, voire plus rarement $D^nf$ en fixant une chaîne admissible quelconque (en fait, cette chaîne est unique) $(\varphi_0, \ ..., \ \varphi_n)$ puis en posant $f^{(n)}=\varphi_n$
\end{dfn}
\begin{prop}[Relation de Chasles sur la classe $\mathcal{D}^n$]
	Soit $(p,q)\in\N^2$ et $f:I\rightarrow\R$. Alors, $f$ est de classe $\mathcal{D}^{p+q}$ si, et seulement si, $f$ est de classe $\mathcal{D}^p$ et $f^{(p)}$ est de classe $\mathcal{D}^q$. Dans ce cas, on a $f^{(p+q)}=(f^{(p)})^{(q)}$
\end{prop}
\begin{cor}
	Soit $f:I\rightarrow\R$. $f$ est de classe $\mathcal{D}^\infty$ si, et seulement si, elle admet une chaîne admissible infinie.
\end{cor}
\begin{prop}
	Soit $f:I\rightarrow\R$. Si $f$ est de classe $\mathcal{D}^1$ et si $f'$ est de classe $\mathcal{D}^\infty$, alors $f$ est de classe $\mathcal{D}^\infty$.
\end{prop}
\begin{dfn}[Classe $\mathcal{D}^n$ sur]
	Soit $f:I\rightarrow\R$ et $J\subset I$ un autre intervalle non trivial. $f$ est dite de classe $\mathcal{D}^n$ sur $J$ lorsque la restriction de $f$ à $J$ est de classe $\mathcal{D}^n$.
\end{dfn}
\begin{prop}[Classe $\mathcal{D}^n$ induite]
	Soit $f\in\mathcal{D}^n(I,\R)$ et $J\subset I$ un autre INT. Alors $f$ est de classe $\mathcal{D}^n$ sur $J$ et toutes les dérivées k-èmes pour $0 \leqslant k \leqslant n$ de la restriction de $f$ à $J$ coïncident avec celles de $f$. 
\end{prop}
\begin{prop}[Caractère local de la classe $\mathcal{D}^n$]
	Soit $(I_\omega)_{\omega\in\Omega}$ une famille d'intervalles ouverts de $\R$ et $\displaystyle f:\bigcup_{\omega\in\Omega} \rightarrow \R$. Si pour tout $\omega\in\Omega$, f est de classe $\mathcal{D}^n$, alors $f$ est de classe $\mathcal{D}^n$.
\end{prop}
\begin{rmq} Les deux résultats qui précèdent sont aussi vrai pour la classe $\mathcal{D}^\infty$, cela s'en déduit immédiatement. \end{rmq}
\begin{dfn}[Classe $\mathcal{C}^n$]
	Soit $f:I\rightarrow\R$ et $n\in\N$. $f$ est dite de classe $\mathcal{C}^n$ lorsque $f$ est de classe $\mathcal{C}^n$ et $f^{(n)}$ est continue. On note $\mathcal{C}^n(I,\R)$ l'ensemble des fonctions de $I$ dans $\R$ de classe $\mathcal{C}^n$.
\end{dfn}
\begin{dfn}[Classe $\mathcal{C}^\infty$]
	On pose $$\mathcal{C}^\infty=\bigcap_{n\in\N}\mathcal{C}^n(I,\R)$$Les fonctions qui appartiennent à cet ensemble sont dites de classe $\mathcal{C}^\infty$.
\end{dfn}
\begin{rmq} On a la chaîne infinie d'inclusion suivante : $$\mathcal{C}^\infty(I, \R) \subset ... \subset \mathcal{C}^1(I,\R) \subset \mathcal{C}^0(I,\R)$$ \end{rmq}
\begin{rmq} On a $\mathcal{C}^\infty(I, \R) =\mathcal{D}^\infty(I, \R)$ \end{rmq}
\begin{prop}[Relation de Chasles sur la classe $\mathcal{C}^n$]
	Soit $(p,q)\in\N^2$ et $f:I\rightarrow\R$. Alors, $f$ est de classe $\mathcal{C}^{p+q}$ si, et seulement si, $f$ est de classe $\mathcal{D}^p$ et $f^{(p)}$ est de classe $\mathcal{C}^q$. Dans ce cas, on a $f^{(p+q)}=(f^{(p)})^{(q)}$
\end{prop}
\begin{thm}[Combinaison linéaire]
	Soit $f$ et $g$ des fonctions de $\mathcal{D}^n(I,\R)$ (resp. $\mathcal{C}^n(I,\R)$). Soit $\lambda$ et $\mu$ des réels. Alors 
	\begin{enumerate}
		\item $\lambda f+\mu g \in\mathcal{D}^n(I,\R)$ (resp. $\mathcal{C}^n(I,\R)$)
		\item $(\lambda f +\mu g)^{(n)}=\lambda f^{(n)}+\mu g^{(n)}$
	\end{enumerate}
	Le résultat est donc aussi valable pour la classe $\mathcal{C}^\infty$.
\end{thm}
\begin{thm}[Produit]
	Soit $f$ et $g$ des fonctions de $\mathcal{D}^n(I,\R)$ (resp. $\mathcal{C}^n(I,\R)$). Alors 
	\begin{enumerate}
		\item $fg \in\mathcal{D}^n(I,\R)$ (resp. $\mathcal{C}^n(I,\R)$)
		\item $\displaystyle (fg)^{(n)}=\sum_{k=0}^{n}\binom{n}{k}f^{(k)}g^{(n-k)}$ C'est ce qu'on appelle la \textbf{formule de Leibniz}.
	\end{enumerate}
	Le résultat est donc aussi valable pour la classe $\mathcal{C}^\infty$.
\end{thm}
\begin{thm}[Quotient]
	Soit $f$ et $g$ des fonctions de $\mathcal{D}^n(I,\R)$ (resp. $\mathcal{C}^n(I,\R)$) telles que $\forall x \in I, \ g(x)\ne 0$. Alors $\displaystyle \frac{f}{g} \in\mathcal{D}^n(I,\R)$ (resp. $\mathcal{C}^n(I,\R)$).
	Le résultat est donc aussi valable pour la classe $\mathcal{C}^\infty$.
\end{thm}
\begin{thm}[Composition]
	Soit $f \in \mathcal{D}^n(I,\R)$ (resp. $\mathcal{C}^n(I,\R)$) et $g \in \mathcal{D}^n(J,\R)$ (resp. $\mathcal{C}^n(I,\R)$) telles que $f(I)\subset J $. Alors $g \circ f \in\mathcal{D}^n(I,\R)$ (resp. $\mathcal{C}^n(I,\R)$).
	Le résultat est donc aussi valable pour la classe $\mathcal{C}^\infty$.
\end{thm}
\begin{thm}[Théorème de la bijection réciproque pour les fonctions de classe $\mathcal{C}^n$]
	Soit $I$ et $J$ des intervalles non triviaux de $\R$ et $f:I\rightarrow J$ une bijection strictement monotone. Soit $n\in\N\st$. Si $f$ est de classe $\mathcal{C}^n$ et $f'$ ne s'annule pas, alors $f\inv$ est de classe $\mathcal{C}^n$. Le résultat est donc aussi valable pour la classe  $\mathcal{C}^\infty$.
\end{thm}
\section{Extension aux fonctions à valeurs complexes}
\begin{dfn}[Dérivabilité des fonctions à valeurs complexes]
	Soit $f:I\rightarrow\C$. On définit de même la dérivabilité en un point, la dérivabilité à gauche et à droite en un point, et la dérivabilité globale. On définit de même les classes $\mathcal{C}^n$ et $\mathcal{D}^n$.
\end{dfn}
\begin{prop}[Dérivable $\Rightarrow$ continue]
	Soit $f:I\rightarrow\C$ $a\in I$. Si $f$ est dérivable en $a$, alors $f$ est continue en $a$.
\end{prop}
\begin{thm}[Théorème passerelle]
	On a les équivalences suivantes :
	\begin{enumerate}
		\item $f$ est dérivable en $a$ si, et seulement si, $\text{Re}(f)$ et $\text{Im}(f)$ sont dérivables en $a$. Dans ce cas, $f'(a)=[\text{Re}(f)]'(a)+i[\text{Im}(f)]'(a)$
		\item $f$ est dérivable si, et seulement si, $\text{Re}(f)$ et $\text{Im}(f)$ sont dérivables. Dans ce cas, $f'=[\text{Re}(f)]'+i[\text{Im}(f)]'$
		\item Soit $n\in\N$. $f$ est de classe $\mathcal{D}^n$ si, et seulement si, $\text{Re}(f)$ et $\text{Im}(f)$ sont de classe $\mathcal{D}^n$. Dans ce cas, $f^{(n)}=[\text{Re}(f)]^{(n)}+i[\text{Im}(f)]^{(n)}$
		\item Soit $n\in\N$. $f$ est de classe $\mathcal{C}^n$ si, et seulement si, $\text{Re}(f)$ et $\text{Im}(f)$ sont de classe $\mathcal{C}^n$. Dans ce cas, $f^{(n)}=[\text{Re}(f)]^{(n)}+i[\text{Im}(f)]^{(n)}$
		\item $f$ est de classe $\mathcal{C}^\infty$ si, et seulement si, $\text{Re}(f)$ et $\text{Im}(f)$ sont de classe $\mathcal{C}^\infty$.
	\end{enumerate}
\end{thm}
\begin{thm}[Opérations algébriques sur la dérivabilité des fonctions à valeurs complexes]
	Les opérations algébriques suivantes sur la dérivabilité des fonctions à valeurs complexes sont valables : combinaison linéaire, produit, inverse et quotient lorsque le dénominateur ne s'annule pas, passage au conjugué (via théorème passerelle) et composée $g\circ f$ lorsque $f$ est \textbf{à valeurs réelles}. Il en va de même pour les classes $\mathcal{D}^n$ et $\mathcal{C}^n$. La formule de Leibniz reste valable. \textbf{Attention}, le passage au module ne fonctionne pas, tout comme la valeur absolue sur $\R$.
\end{thm}
\begin{thm}[Dérivabilité de l'exponentielle complexe]
	Soit $f:I\rightarrow\C$ dérivable. Alors, $\exp\circ f:I\rightarrow\C$ est dérivable et on a $$(\exp\circ f)'=(\exp\circ f)\times f'$$
\end{thm}
\begin{cor}
	Soit $f\in\mathcal{D}^1(I,\C)$.
	\begin{enumerate}
		\item Soit $n\in\N\st$. Alors $f^n$ est dérivable et $(f^n)'=nf^{n-1}f'$
		\item Soit $n\in\Z$. Supposons que $f$ est à valeurs dans $\C\st$. Alors $f^n$ est dérivable et $(f^n)'=nf^{n-1}f'$
		\item Soit $\alpha\in\C$. Supposons que $f$ est à valeurs dans $\R_+\st$. Alors $f^\alpha$ est dérivable et $(f^n)'=\alpha f^{\alpha-1}f'$
	\end{enumerate}
\end{cor}
\begin{thm}[Inégalité des accroissements finis, cas complexe]
	Soit $f:I\rightarrow\C$ continue et dérivable en tout point intérieur de $I$. Supposons qu'il existe $\lambda\geqslant 0$ tel que pour tout $x$ de l'intérieur de $I$, $|f'(x)|\leqslant \lambda$. Alors, $f$ est $\lambda$-lipschitzienne.
\end{thm}
\begin{thm}
	Soit $f:I\rightarrow\C$ continue et dérivable en tout point intérieur de $I$. Alors $f$ est constante si, et seulement si, pour tout $x$ de l'intérieur de $I$, $f'(x)=0$.
\end{thm}
\section{Fonctions convexes}
\begin{dfn}[Fonction convexe]
	$f:I\rightarrow\R$ est dite convexe lorsque $$\forall(a,b)\in I^2, \ \forall t\in\left[0,1\right], \ f((1-t)a+tb)\leqslant(1-t)f(a)+tf(b)$$
\end{dfn}
\begin{thm}[Inégalité de Jensen]
	Soit $f:I\rightarrow\R$ convexe, $n\in\N\st$ et $\lambda_1, \ ..., \ \lambda_n$ des réels \textbf{positifs} tels que $\displaystyle \sum_{k=1}^{n}\lambda_k=1$. Alors $$\forall(a_1, \ ..., \ a_n)\in I^n, \ f\left(\sum_{k=1}^{n}\lambda_ka_k\right)\leqslant \sum_{k=1}^{n}\lambda_kf(a_k)$$
\end{thm}
\begin{lem}[Inégalité des pentes]
	Soit $f:I\rightarrow\R$ convexe et $(a,b,c)\in I^3$ tel que $a<b<c$. Alors $$\frac{f(b)-f(a)}{b-a}\leqslant\frac{f(c)-f(a)}{c-a}\leqslant\frac{f(c)-f(b)}{c-b}$$
\end{lem}
\begin{cor}[Caractérisation de la convexité par la croissance des taux d'accroissements]
	Soit $f:I\rightarrow\R$. $f$ est convexe si, et seulement, pour tout $a\in I$, la fonction $\tau_{a,f}$ (taux de variation de $f$ en $a$) est croissante. 
\end{cor}
\begin{cor}[Position relative des cordes pour les fonctions convexes]
	Une fonction convexe est située sous ses cordes entre les points d'intersection de la courbe avec les cordes et au-dessus en dehors.
\end{cor}
\begin{thm}[Dérivabilité à gauche et à droite des fonctions convexes]
	Soit $f:I\rightarrow\R$ convexe. Alors $f$ est dérivable à gauche et à droite en tout point de l'intérieur de $I$. En conséquence, $f$ est nécessairement continue en tout point de l'intérieur de $I$. De plus, l'ensemble des points de non dérivabilité est au plus dénombrable (en utilisant le lemme de l'inégalité des pentes et la densité de $\Q$ dans $\R$).
\end{thm}
\begin{rmq} On ne peut rien dire en bord d'intervalle ! Par exemple, la fonction indicatrice du singleton nul restreinte à $\R_+$ est convexe, mais n'est pas continue en 0. \end{rmq}
\begin{thm}[Caractérisation de la convexité pour les fonctions dérivables]
	Soit $f\in\mathcal{D}^1(I,\R)$. $f$ est convexe si, et seulement si, $f'$ est croissante. Dans ce cas, la courbe de $f$ est située au-dessus de toutes ses tangentes.
\end{thm}
\begin{cor}[Caractérisation de la convexité pour les fonctions deux fois dérivables]
	Soit $f\in\mathcal{D}^2(I,\R)$. $f$ est convexe si, et seulement si, $f''\geqslant 0$.
\end{cor}

\newpage
\chapter{Intégration}
\section{Fonctions en escalier et continues par morceaux}
Dans tout le paragraphe, on fixe $a$ et $b$ des réels tels que $a<b$.
\begin{dfn}[Subdivision]
	Une \textbf{subdivision} de $[a,b]$ est une famille $\sigma=(c_i)_{0\leqslant i\leqslant n}$ avec $n\in\N\st$ telle que $a=c_0<...<c_n=b$. Le \textbf{pas de la subdivision} est défini par $\underset{1\leqslant i\leqslant n}{\max}(c_i-c_{i-1})$.
\end{dfn}
\begin{dfn}[Subdivision régulière]
	Pour $n\in\N\st$, la subdivision $\displaystyle \left(a+i\frac{b-a}{n}\right)_{0\leqslant i\leqslant n}$ est appelée \textbf{subdivision régulière}.
\end{dfn}
\begin{dfn}[Subdivision plus fine]
	Soit $\sigma$ et $\sigma'$ des subdivisions de $[a,b]$. $\sigma$ est dite \textbf{plus fine} que $\sigma'$ lorsque tous les points de $\sigma'$ sont des points de $\sigma$.
\end{dfn}
\begin{prop}[Réunion de subdivisions]
	Soit $\sigma_1$ et $\sigma_2$ des subdivisions de $[a,b]$. En réunissant les points de $\sigma_1$ et $\sigma_2$, on obtient une subdivision de $[a,b]$ qu'on peut noter de façon NS $\sigma_1 \vee \sigma_2$. Cette subdivision est à la fois plus fine que $\sigma_1$ et que $\sigma_2$.
\end{prop}
\begin{dfn}[Fonction en escalier]
	Une fonction $\varphi:[a,b]\rightarrow\R$ est dite \textbf{en escalier} lorsqu'il existe une subdivision $\sigma=(c_i)_{0\leqslant i\leqslant n}$ telle que pour tout $i\in\llbracket1,n\rrbracket$, la restriction de $\varphi$ à $]c_{i-1},c_i[$ soit constante. On dit alors que $\sigma$ est une subdivision \textbf{adaptée} à $\varphi$. On notera $\mathcal{E}([a,b],\R)$ l'ensemble des fonctions de $[a,b]$ dans $\R$ en escalier.
\end{dfn}
\begin{rmq} Si on besoin d'expliciter la valeur constante de $\varphi$ sur $]c_{i-1},c_i[$, on peut choisir $\displaystyle \varphi \left(\frac{c_{i-1}+c_i}{2}\right)$ \end{rmq}
\begin{rmq} Toute subdivision plus fine d'une subdivision adaptée à une fonction en escalier reste adaptée à cette fonction en escalier. \end{rmq}
\begin{prop}
	$\mathcal{E}([a,b],\R)$ est une sous-algèbre de $\R^{[a,b]}$
\end{prop}
\begin{dfn}[Fonction continue par morceaux]
	Une fonction $f:[a,b]\rightarrow\R$ est dite \textbf{continue par morceaux} lorsqu'il existe une subdivision $\sigma=(c_i)_{0\leqslant i\leqslant n}$ telle que pour tout $i\in\llbracket1,n\rrbracket$, la restriction de $f$ à $]c_{i-1},c_i[$ soit continue et admette des limites finies à droite en $c_{i-1}$ et à gauche en $c_i$. On notera $CPM([a,b],\R)$ l'ensemble des fonctions de $[a,b]$ dans $\R$ continues par morceaux.
\end{dfn}
\begin{prop}(En escalier $\Rightarrow$ CPM)
	Toute fonction en escalier est continue par morceaux. Formellement, $\mathcal{E}([a,b],\R) \subset CPM([a,b],\R)$ 
\end{prop}
\begin{prop}[CPM $\Rightarrow$ bornée]
	Toute fonction continue par morceaux est bornée.
\end{prop}
\begin{dfn}[Fonctions continues par morceaux sur un INT quelconque]
	Soit $I$ un INT de $\R$. On dit qu'une fonction $f:I\rightarrow\R$ est continue par morceaux lorsque sa restriction à tout segment non trivial de $I$ est continue par morceaux.
\end{dfn}
\begin{thm}[Théorème d'approximation]
	Soit $f\in CPM([a,b],\R)$ et $\varepsilon>0$. Alors, il existe une fonction $\varphi\in\mathcal{E}([a,b],\R)$ telle que $$\forall x \in [a,b], \ |f(x)-\varphi(x)|\leqslant \varepsilon$$
\end{thm}
\begin{rmq} Utile pour démontrer le lemme de Riemann-Lebesgue dans le cas des fonctions CPM \end{rmq}
\begin{dfn}[Norme infinie]
	Soit $f:[a,b]\rightarrow\R$ bornée. Sa \textbf{norme infinie} est le réel défini par $$||f||_\infty= \underset{x\in[a,b]}{|f(x)|}$$
\end{dfn}
\begin{dfn}[Convergence uniforme]
	Soit $f:[a,b]\rightarrow\R$ bornée et $(f_n)_{n\in\N}$ une suite de fonctions de $[a,b]$ dans $\R$ bornées. Les deux assertions suivantes sont équivalentes :
	\begin{enumerate}
		\item $||f-f_n||_\infty \xrightarrow[n\to+\infty]{}0$
		\item $\forall \varepsilon>0, \ \exists N\in\N, \ \forall n\geqslant N, \ \forall x\in[a,b], \ |f(x)-f_n(x)|\leqslant \varepsilon$
	\end{enumerate}
	Lorsqu'une de ces assertion est vérifiée, on dit que la suite de fonctions $(f_n)$ \textbf{converge uniformément} vers $f$. On le note $$f_n \xrightarrow[n\to+\infty]{CVU} f$$
\end{dfn}
\begin{thm}[Approximation des fonctions continues par morceaux par des fonctions en escalier]
	Soit $f\in CPM([a,b],\R)$. Alors, il existe une suite de fonctions en escalier $(\varphi_n)_{n\in\N}$ telle que $$\varphi_n \xrightarrow[n\to+\infty]{CVU} f$$
\end{thm}

\section{Définition de l'intégrale}
\begin{dfn}[Intégrale d'une fonction en escalier]
	Soit $\varphi\in\mathcal{E}([a,b],\R)$. Fixons $\sigma=(c_i)_{0\leqslant i\leqslant n}$ une subdivision adaptée à $\varphi$. Pour tout $i\in\llbracket1,n\rrbracket$, on note $\lambda_i$ la valeur constante de $\varphi$ sur $]c_{i_1},c_i[$. On définit alors l'intégrale de $\varphi$ par $$I(\varphi,\sigma)=\sum_{i=1}^{n}\lambda_i(c_i-c_{i_1})$$ qui est indépendant du choix de la subdivision. On note plus souvent ce  scalaire $$\int_{[a,b]}\varphi$$
\end{dfn}
\begin{prop}[Linéarité de l'intégrale pour les fonctions en escalier]
	Soit $\lambda$ et $\mu$ des réels ainsi que $\varphi$ et $\psi$ des fonctions en escalier. Alors $$\int_{[a,b]}(\lambda\varphi+\mu\psi)=\lambda\int_{[a,b]}\varphi + \mu\int_{[a,b]}\psi$$
\end{prop}
\begin{dfn}[Intégrale "sur" pour les fonctions en escalier]
	Soit $\varphi\in\mathcal{E}([a,b],\R)$ et $(c,d)\in[a,b]^2$ tel que $c<d$. On définit l'intégrale de $\varphi$ sur $[c,d]$ par $$\int_{[c,d]}\varphi=\int_{[c,d]}\varphi_{|[c,d]}$$
\end{dfn}
\begin{prop}[Relation de Chasles pour les fonctions en escalier]
	Soit $\varphi\in\mathcal{E}([a,b],\R)$ et $c\in]a,b[$. On a $$\int_{[a,b]}\varphi=\int_{[a,c]}\varphi+\int_{[c,b]}\varphi$$
\end{prop}
\begin{prop}[Inégalité triangulaire pour les fonctions en escalier]
	Soit $\varphi\in\mathcal{E}([a,b],\R)$. On a
	\begin{enumerate}
		\item $|\varphi|\in\mathcal{E}([a,b],\R)$
		\item $$\left\lvert\int_{[a,b]}\varphi\right\rvert \leqslant \int_{[a,b]}\lvert\varphi\rvert$$
	\end{enumerate}
\end{prop}
\begin{prop}[Positivité et croissance de l'intégrale pour les fonctions en escalier]
	On a les deux propriétés suivantes 
	\begin{enumerate}
		\item Positivité : $$\forall\varphi\in\mathcal{E}([a,b],\R),\ \varphi\geqslant 0 \Rightarrow \int_{[a,b]}\varphi \geqslant 0$$
		\item Croissance :$$\forall(\varphi,\psi)\in(\mathcal{E}([a,b],\R))^2,\ \varphi\geqslant \psi \Rightarrow \int_{[a,b]}\varphi \geqslant \int_{[a,b]}\psi$$
	\end{enumerate}
\end{prop}
\begin{prop}[Insensibilité de l'intégrale pour les fonctions en escalier]
	Soit $(\varphi,\psi)\in(\mathcal{E}([a,b],\R))^2$. Si $\varphi$ et $\psi$ ne diffèrent qu'en un nombre fini de points, alors $$\int_{[a,b]}\varphi = \int_{[a,b]}\psi$$
\end{prop}
\begin{dfn}[Intégrale d'une fonction continue par morceaux]
	Soit $f\in CPM([a,b],\R)$. On fixe une suite $(\varphi_n)_{n\in\N}$ de fonctions en escalier telle que $\varphi_n \xrightarrow[n\to+\infty]{CVU} f$ puis on définit l'intégrale de $f$ par $$\int_{[a,b]}f=\underset{n\to+\infty}{\lim}\int_{[a,b]}\varphi_n$$
\end{dfn}
\begin{prop}[Linéarité de l'intégrale pour les fonctions CPM]
	Soit $\lambda$ et $\mu$ des réels ainsi que $f$ et $g$ des fonctions continues par morceaux. Alors $$\int_{[a,b]}(\lambda f+\mu g)=\lambda\int_{[a,b]}f + \mu\int_{[a,b]}g$$
\end{prop}
\begin{dfn}[Intégrale "sur" pour les fonctions CPM]
	Soit $f\in CPM([a,b],\R)$ et $(c,d)\in[a,b]^2$ tel que $c<d$. On définit l'intégrale de $f$ sur $[c,d]$ par $$\int_{[c,d]}f=\int_{[c,d]}f_{|[c,d]}$$
\end{dfn}
\begin{prop}[Relation de Chasles pour les fonctions CPM]
	Soit $f\in CPM([a,b],\R)$ et $c\in]a,b[$. On a $$\int_{[a,b]}f=\int_{[a,c]}f+\int_{[c,b]}f$$
\end{prop}
\begin{prop}[Inégalité triangulaire pour les fonctions CPM]
	Soit $f\in CPM([a,b],\R)$. On a
	\begin{enumerate}
		\item $|f|\in CPM([a,b],\R)$
		\item $$\left\lvert\int_{[a,b]}f\right\rvert \leqslant \int_{[a,b]}\lvert f\rvert$$
	\end{enumerate}
\end{prop}
\begin{prop}[Positivité et croissance de l'intégrale pour les fonctions CPM]
	On a les deux propriétés suivantes 
	\begin{enumerate}
		\item Positivité : $$\forall f\in CPM([a,b],\R),\ f\geqslant 0 \Rightarrow \int_{[a,b]}f \geqslant 0$$
		\item Croissance :$$\forall(f,g)\in(CPM([a,b],\R))^2,\ f\geqslant g \Rightarrow \int_{[a,b]}f \geqslant \int_{[a,b]}g$$
	\end{enumerate}
\end{prop}
\begin{prop}[Insensibilité de l'intégrale pour les fonctions CPM]
	Soit $(f,g)\in(CPM([a,b],\R))^2$. Si $\varphi$ et $\psi$ ne diffèrent qu'en un nombre fini de points, alors $$\int_{[a,b]}f = \int_{[a,b]}g$$
\end{prop}
\begin{prop}[Stricte positivité et stricte croissance de l'intégrale pour des fonctions continues.]
	On a les deux résultats suivant :
	\begin{enumerate}
		\item Stricte positivité : Soit $f\in\mathcal{C}^0([a,b],\R)$. Si $f>0$, alors $$\int_{[a,b]}f >0$$
		\item Stricte croissance : Soit $(f,g)\in(\mathcal{C}^0([a,b],\R))^2$. Si $f>g$, alors $$\int_{[a,b]}f >\int_{[a,b]}g$$
	\end{enumerate}
\end{prop}
\begin{prop}
	Soit $f\in\mathcal{C}^0([a,b],\R)$ de signe constant. Si $\displaystyle\int_{[a,b]}f=0$, alors $f=0$
\end{prop}
\begin{thm}[Sommes de Riemann]
	Soit $f\in CPM([a,b],\R)$. Posons $$\forall n\in\N\st, \ S_n=\frac{b-a}{n}\sum_{k=0}^{n-1}f\left(a+k\frac{b-a}{n}\right)$$ et $$\forall n\in\N\st, \ S_n'=\frac{b-a}{n}\sum_{k=1}^{n}f\left(a+k\frac{b-a}{n}\right)$$ Alors $$S_n \xrightarrow[n\to+\infty]{} \int_{[a,b]}f$$ et $$S_n' \xrightarrow[n\to+\infty]{} \int_{[a,b]}f$$
\end{thm}
\begin{dfn}[Valeur moyenne]
	Soit $f\in CPM([a,b],\R)$. On appelle \textbf{valeur moyenne} de $f$ le réel $$\mu(f)=\frac{1}{b-a} \int_{[a,b]}f$$
\end{dfn}
\begin{dfn}[Intégrale orientée]
	Soit $f\in CPM(I,\R)$ avec $I$ un INT de $\R$ et $(a,b)\in I^2$. On définit
	\begin{enumerate}
		\item Si $a<b$, $$\int_{a}^{b}f(t)dt=\int_{[a,b]}f$$
		\item Si $a=b$, $$\int_{a}^{b}f(t)dt=0$$
		\item Si $a>b$, $$\int_{a}^{b}f(t)dt=-\int_{[a,b]}f$$
	\end{enumerate}
\end{dfn}
\begin{prop}[Validité des propriétés de l'intégrale pour l'intégrale orientée]
	La linéarité, la relation de Chasles étendue à $(a,b,c)\in I^3$ et le théorème sur les sommes de Riemann restent valables pour l'intégrale orientée. En revanche, la positivité, la croissance, l'inégalité triangulaire ne sont valables que si $a\leqslant b$. La stricte positivité et la stricte croissance ne restent vraies que si $a<b$. Ces hypothèses devront \textbf{toujours} être mentionnées avant l'utilisation de ces théorèmes.
\end{prop}
\begin{prop}[Extension aux fonctions à valeurs complexes]
	On étend de manière naturelle les définitions aux fonctions à valeurs complexes. Tous les théorèmes restent vrais, sauf les inégalités, la seule inégalité qui reste vraie est l'inégalité triangulaire (car absence de relation d'ordre prolongeant celle de $\R$ sur $\C$)
\end{prop}
\begin{thm}[Théorème passerelle]
	Soit $f\in CPM([a,b],\C)$. Alors $\text{Re}(f)$ et $\text{Im}(f)$ sont continues par morceaux et à valeurs dans $\R$, et on a $$\int_{[a,b]}f=\int_{[a,b]}\text{Re}(f)+\text{i}\int_{[a,b]}\text{Im}(f)$$ Autrement dit, on a $\displaystyle \text{Re}\left(\int_{[a,b]}f\right)=\int_{[a,b]}\text{Re}(f)$ et $\displaystyle \text{Im}\left(\int_{[a,b]}f\right)=\int_{[a,b]}\text{Im}(f)$
\end{thm}

\section{Intégration et dérivation}
Dans tout le paragraphe, on fixe $I$ un intervalle non trivial de $\R$ et $f:I\rightarrow\K$ avec $\K=\R$ ou $\C$.

\begin{dfn}[Primitive]
	Une primitive de $f$ est une fonction $F\in\mathcal{D}^1(I,\K)$ telle que $F'=f$.
\end{dfn}
\begin{prop}[Unicité à constante près des primitives]
	Si elles existent, les primitives de $f$ diffèrent toutes d'une constante additive.
\end{prop}
\begin{thm}[Théorème fondamental de l'analyse]
	Supposons que $f$ est continue. Soit $a\in I$ quelconque. Alors la fonction
	$$
	\begin{array}[t]{lrcl}
		Fv: & I & \longrightarrow & \K \\
		& x & \longmapsto & \displaystyle \int_{a}^{x}f(t)dt
	\end{array}
	$$ est l'unique primitive de $f$ qui s'annule en $a$.
\end{thm}
\begin{met}[Calcul d'une primitive] Pour calculer une primitive d'une fonction, il suffit de s'assurer que celle-ci est continue, puis de choisir le $a$ le plus simple possible, de faire monter $x$ dans la borne puis de calculer l'intégrale grâce aux méthodes présentées plus loin. \end{met}
\begin{rmq} Attention, toute primitive ne s'obtient pas par le TFA. Par exemple, en considérant la fonction nulle, changer les bornes ne fera qu'ajouter 0, et ne nous donnera jamais les fonctions constantes. \end{rmq}
\begin{rmq} On peut trouver dommage de demander la forte hypothèse de continuité, mais on ne peut pas vraiment faire plus simple. Il existe même des fonctions discontinues qui admettent des primitives (considérer la dérivée de $x\mapsto x^2\sin(\frac1x))$ prolongée par 0 en 0, puisque la fonction précédente est de classe $\mathcal{D}^1$ mais pas de classe $\mathcal{C}^1$. \end{rmq}
\begin{rmq}[HP] Si $f$ est CPM et admet une primitive, alors elle est continue. Pour montrer cela, il faut appliquer le TLD aux bords de la subdivision. \end{rmq}

\begin{thm}[Calcul direct à l'aide d'une primitive]
	Soit $(a,b)\in I^2$. Supposons que $f$ soit continue et fixons $F$ une primitive de $f$. On a $$\int_{a}^{b}f(t)dt=F(b)-F(a)$$ On note, pour alléger l'écriture $[F(t)]_a^b$ : on parle de \textbf{terme tout intégré}.
\end{thm}
\begin{cor}
	Supposons que $f$ est de classe $\mathcal{C}^1$. Alors $$\forall (a,b)\in I^2, \ \int_{a}^{b}f'(t)dt=f(b)-f(a)$$
\end{cor}
\begin{rmq} Pour montrer l'invariance par $T$-translation d'une fonction continue $T$-périodique, il faut montrer que la fonction $x\mapsto\int_{x}^{x+T}f(t)dt$ est constante. Pour cela, l'exprimer à l'aide d'une primitive puis dériver pour obtenir la constance et évaluer en 0 pour conclure. Dans la cas CPM, il faut fixer une subdivision de l'intervalle de longueur $T$ puis coincer un multiple de $T$ dans cette subdivision. Ensuite, on travaille sur chacun des segments où la restriction est continue puis on recolle par la relation de Chasles. \end{rmq}
\begin{thm}[Intégration par parties]
	Soit $g$ ayant les mêmes hypothèses que $f$. Supposons que $f$ et $g$ soient de classe $\mathcal{C}^1$. Alors 
	$$\forall (a,b)\in I^2, \ \int_{a}^{b}f'(t)g(t)dt=[f(t)g(t)]_a^b-\int_{a}^{b}f(t)g'(t)dt$$
\end{thm}
\begin{thm}[Changement de variable]
	Soit $I$,$J$ des INT de $\R$, $f\in\mathcal{C}^0(I,\R)$ et $\varphi\in\mathcal{C}^1(J,I)$. Alors 
	$$\int_{\varphi(a)}^{\varphi(b)}f(t)dt=\int_{a}^{b}f(\varphi(t))\varphi'(t)dt$$
\end{thm}
\begin{rmq} On a les changements de variables usuels dans les cas suivants :
	\begin{enumerate}
		\item $\sqrt{1-t^2}$ : on utilise $t=\cos(u)$ ou $t=\sin(u)$
		\item $\sqrt{t^2-1}$ : on utilise $t=\cosh(u)$
		\item $\sqrt{t^2+1}$ : on utilise $t=\sinh(u)$
\end{enumerate} \end{rmq}
\begin{prop}[Règles de Bioche, HP]
	On souhaite calculer $\int F(\sin(\theta),\cos(\theta),\tan(theta))d\theta$ où $F$ est une fraction rationnelle. On pose $\omega(\theta)=F(\theta)d\theta$. Les règles de Bioche énoncent des changements de variable à effectuer pour accélérer le calcul :
	\begin{enumerate}
		\item Supposons que $\omega(-\theta)=\omega(\theta)$ : on pose alors $u=\cos(\theta)$
		\item Supposons que $\omega(\pi-\theta)=\omega(\theta)$ : on pose alors $u=\sin(\theta)$
		\item Supposons que $\omega(\pi+\theta)=\omega(\theta)$ : on pose alors $u=\tan(\theta)$
		\item Supposons que deux des propriétés précédentes soient vérifiées. Alors, la troisième l'est aussi et on pose $u=\cos(2\theta)$
		\item Si aucune des trois propriétés n'est vérifiée, on pose $\displaystyle u=\tan\left(\frac{\theta}{2}\right)$ et on utilise les formules $\displaystyle \cos(\theta)=\frac{1-u^2}{1+u^2}$, $\displaystyle \sin(\theta)=\frac{2u}{1+u^2}$ et $\displaystyle \tan(\theta)=\frac{2u}{1-u^2}$
	\end{enumerate}
	Dans tous les cas, on se ramène à un calcul d'intégrale fraction rationnelle qu'on sait de toute manière toujours réaliser. Les 4 premiers changements de variables sont en général beaucoup plus rapide que le dernier.
\end{prop}
\begin{rmq} Bien que HP, on peut toujours utiliser ces règles sans justification car ce n'est pas réellement un théorème mais une astuce. Pour retenir les trois premiers changements de variable, remarquer qu'il correspondent aux mêmes invariances trigonométriques que celles supposées sur la fonction $\omega$. Pour le dernier, il faut absolument connaître ses formules de trigonométrie. Le 4ème est en revanche beaucoup plus rare. \end{rmq}
\begin{cor}[Parité et imparité, HP]
	Supposons que $I$ soit centré en 0 et que $f$ soit continue.
	\begin{enumerate}
		\item Supposons que $f$ soit paire. Alors $$\forall a \in I, \ \int_{-a}^{a}f(t)dt=2\int_{0}^{a}f(t)dt$$
		\item Supposons que $f$ soit impaire. Alors $$\forall a \in I, \ \int_{-a}^{a}f(t)dt=0$$
	\end{enumerate}
\end{cor}
\begin{proof}
	Couper en deux par la relation de Chasles et utiliser un changement de variable $u=-t$.
\end{proof}
\begin{rmq} On a des résultats similaires pour les fonctions continues qui présentent des symétries en un certain réel $a_0$, qui se démontrent de la même manière. \end{rmq}
\begin{thm}[Formule de Taylor avec reste intégral]
	Soit $a\in I$. Soit $n\in\N$ et supposons que $f\in\mathcal{C}^{n+1}(I,\K)$. Alors $$\forall x \in I, \ f(x)=\sum_{k=0}^{n}\frac{f^{(k)}(a)}{k!}(x-a)^k + \int_{a}^{x}\frac{(x-t)^n}{n!}f^{(n+1)}(t)dt$$
\end{thm}
\begin{thm}[Inégalité de Taylor-Lagrange]
	Soit $a\in I$. Soit $n\in\N$ et supposons que $f\in\mathcal{C}^{n+1}(I,\K)$. Soit $x \in I$ et $M_{n+1}$ un majorant de $|f^{(n+1)}|$ sur $[min(a,x),max(a,x)]$. Alors 
	$$\left\lvert f(x)-\sum_{k=0}^{n}\frac{f^{(k)}(a)}{k!}(x-a)^k\right\rvert \leqslant M_{n+1}\frac{|x-a|^{n+1}}{(n+1)!}$$
\end{thm}
\begin{rmq} Lorsque $M_{n+1}$ ne dépend pas de $n$, on peut obtenir l'expression de $f(x)$ comme somme d'une série par croissances comparées. C'est par exemple le cas avec la fonction $\exp$ ou les fonctions $\sin$ et $\cos$. \end{rmq}
\section{Compléments, HP}
\subsection{Retour sur les sommes de Riemann}

Ici, nous nous proposons de démontrer le théorème sur les sommes de Riemann dans le cas le plus général, c'est-à-dire avec une fonction continue par morceaux et des subdivisions potentiellement irrégulières.

\begin{thm}
	Soit $f$ continue par morceaux sur $[a, b]$, et $\sigma = (a_k)_{0 \leqslant k \leqslant n}$ une subdivision de $[a, b]$. Si le pas de $\sigma$ tend vers 0 alors $$\sum_{k=0}^{n-1}(a_{k+1}-a_k)f(a_k) \rightarrow \int_{[a,b]}f$$
\end{thm}
\begin{proof}
	Dans un premier temps, prouvons le théorème pour une fonction en escalier. On se donne $(c_i)_{0 \leqslant i \leqslant N}$ une subdivision adaptée à $\varphi$. De la même manière que dans le cours, on obtient $$\left\lvert\int_{[a,b]}\varphi - \sum_{k=0}^{n-1}(a_{k+1}-a_k)\varphi(a_k)\right\rvert \leqslant \sum_{k=0}^{n-1}\int_{[a_k,a_{k+1}]}\left\lvert\varphi - \varphi(a_k)\right\rvert$$
	Pour $i \in \llbracket0,N\rrbracket$ fixé, si on considère $K(i) := \{k \in \llbracket0,n-1\rrbracket\ | \ a_k \leqslant c_i < a_{k+1}\}$, on a $|K(i)| \leqslant 1$ car les $[a_k,a_{k+1}[$ sont deux à deux disjoints. Donc si on pose : $$K := \bigcup_{0 \leqslant i \leqslant N}K(i)$$ on a $|K| \leqslant N+1.$ Maintenant, soit $k \in \llbracket0,n-1\rrbracket \backslash K$ fixé, et posons $$i_0 = \max\{i \in \llbracket0,N\rrbracket\ |\ c_i \leqslant a_k\}$$ Comme $[a_k,a_{k+1}[$ ne contient aucun $c_i,$ on a $c_{i_0}< a_k < a_{k+1} \leqslant c_{i_{0}+1}$. En particulier, $[a_k,a_{k+1}[\subset]C_{i_{0}},c_{i_{0}+1}[$ donc $\varphi$ est constante sur $[a_k,a_{k+1}]$ égale à $\varphi(a_k)$. Puisque $(a_k,a_{k+1})$ est adaptée à $|\varphi - \varphi(a_k)|$ on en déduit $$\int_{[a_k,a_{k+1}]}|\varphi - \varphi(a_k)| = 0$$
	De plus, $\varphi$ ne prend qu'un nombre fini de valeurs, donc $|\varphi|$ admet un majorant $M$. Finalement, si on note $p(\sigma)$ le pas de $\sigma$, alors $$\left\lvert\int_{[a,b]}\varphi - \sum_{k=0}^{n-1}(a_{k+1}-a_k)\varphi(a_k)\right\rvert \leqslant \sum_{k \in K}\int_{[a_k,a_{k+1}]}2M \leqslant 2M(N+1)p(\sigma) \rightarrow 0$$
	Dans un second temps, passons à $f$ continue par morceaux. Soit $\varepsilon > 0,$ et donnons-nous $\varphi$ en escalier telle que $\displaystyle|f-\varphi| \leqslant \frac{\varepsilon}{3(b-a)}$. On a 
	\begin{align*}
		\left\lvert \int_{[a,b]}f - \sum_{k=0}^{n-1}(a_{k+1}-a_k)f(a_k)\right\rvert &\leqslant \int_{[a,b]}|f-\varphi| + \left\lvert\left(\int_{[a,b]}\varphi\right) - \sum_{k=0}^{n-1}(a_{k+1}-a_k)\varphi(a_k)\right\rvert\\
		&+ \sum_{k=0}^{n-1}(a_{k+1}-a_k)|\varphi(a_k)-f(a_k)| \\
		&\leqslant \frac{2}{3}\varepsilon + \left\lvert\left(\int_{[a,b]}\varphi\right) - \sum_{k=0}^{n-1}(a_{k+1}-a_k)\varphi(a_k)\right\rvert
	\end{align*}
	Enfin, pour $p(\sigma)$ assez petit, la dernière valeur absolue est $\displaystyle\leqslant \frac{\varepsilon}{3}$.
\end{proof}
\subsection{Retour sur la construction de l'intégrale de Riemann}
La construction de l'intégrale de Riemann s'inscrit dans un cadre beaucoup plus vaste (et très fécond !) que nous proposons maintenant. Il s'agit d'une technique de construction maintes fois utilisée par les mathématiciens.\\ \par
Avant tout nous avons besoin de quelques définitions. Dans toute cette partie, $E$ est un espace vectoriel réel ou complexe.
\begin{dfn}[Norme]
	Une \textbf{norme} sur $E$ est une application $||\cdot|| : E \rightarrow \mathbb{R}$, qui vérifie les axiomes suivants :
\begin{itemize}
	\item Séparation : $\forall x \in E, ||x|| = 0 \Rightarrow x = 0$.
	\item Homogénéité : $\forall (\lambda, x) \in \mathbb{K} \times E, ||\lambda x|| = |\lambda| \cdot ||x||$.
	\item Inégalité triangulaire : $\forall (x, y) \in E^2, ||x+y|| \leqslant ||x|| + ||y||$.
\end{itemize}
\end{dfn}
\begin{exa}
	La valeur absolue et le module sont respectivement des normes sur $E = \mathbb{R}$ et sur $E = \mathbb{C}$.
\end{exa}
Dans le cadre d'un espace vectoriel normé (evn en abrégé), on peut ensuite définir les notions de suite convergente et de suite de Cauchy au sens usuel. Par exemple, on a le très important
\begin{dfn}
	Soit $[a, b]$ un segment de $\mathbb{R}$ et posons $\mathbb{K} = \mathbb{R} ou \mathbb{C}$. Alors l'ensemble $\mathcal{B}([a, b], \mathbb{K})$ des fonctions bornées de $[a, b]$ dans $\mathbb{K}$ est un $\mathbb{K}$-ev et peut le munir de la norme uniforme, définie de la manière suivante : $$||f||_{\infty} = \sup_{x \in [a, b]} |f(x)|$$
\end{dfn}
\begin{proof}
	Déjà le nombre $||f||_{\infty}$ est bien défini, car $[a, b]$ est non vide et $|f|$ est majorée. Ensuite on a bien $||f||_{\infty} \geqslant 0,$ avec égalité ssi $f$ est la fonction nulle. La semi-linéarité est facile à prouver. On commence par écrire $$\left(\forall x \in [a, b],\ |\lambda f(x)| \leqslant |\lambda| \cdot ||f||_{\infty}\right) \text{ donc } ||\lambda f||_{\infty} \leqslant |\lambda| \cdot ||f||_{\infty}$$
	Puis si $\lambda = 0$ on obtient $||0 \cdot f||_{\infty} = 0,$ et sinon on applique le résultat à $\frac{1}{\lambda}$ et $\lambda f$ pour obtenir l'inégalité inverse. Enfin pour l'inégalité triangulaire on écrit
	$$\left(\forall x \in [a, b],\ |f(x) + g(x)| \leqslant ||f||_{\infty} + ||g||_{\infty}\right) \text{ donc } ||f+g||_{\infty} \leqslant ||f||_{\infty} + ||g||_{\infty}$$
	Par ailleurs, on sait que $\mathcal{B}([a, b], \mathbb{K})$ est un sev de $\mathbb{K}^{[a, b]}$.
\end{proof}
\begin{dfn}[Adhérence d'une partie]
	Conformément au cas réel, nous dirons que l'adhérence $\overline{A}$ d'une partie $A \subset E$ est l'ensemble des limites de points de $A$.
\end{dfn}
\begin{rmq}
	On pourrait aussi donner une définition en termes de voisinages et montrer qu'elle est équivalente, mais ce n'est pas notre objet ici. Tout cela sera traité en détail en seconde année.
\end{rmq}
On arrive alors au fondamental
\begin{thm}
	Soit $E$ et $F$ des $\mathbb{K}$-ev, avec $F$ complet. Soit $A \subset E.$ et $f: A \rightarrow F$ une application uniformément continue. Alors il existe un unique prolongement continu de $f$ à $\overline{A}$.
\end{thm}
\begin{proof}
	On travaille en deux temps.
	\begin{itemize}
		\item Existence \\
		Soit $x \in \overline{A}$, alors on peut écrire $x = \lim x_n$ avec $x_n \in A$. Comme la suite $(x_n)$ converge elle est de Cauchy, puis comme $f$ est uniformément continue sur $A$, la suite $(f(x_n))$ est de Cauchy dans $F$ (considérer $\varepsilon > 0$ et $\delta > 0$ associé dans l'uniforme continuité de $f$). Puisque $F$ est complet, $(f(x_n))$ admet une limite qu'on note $\overline{f}(x)$. \\
		Montrons que cette définition de $\overline{f}$ est consistante. Si $(x_n')$ est une autre suite qui converge vers $x$, alors $\|x_n - x_n'\| \leqslant \|x_n - x\| + \|x - x_n'\| \rightarrow 0$. Par uniforme continuité de $f$ on a $\|f(x_n) - f(x_n')\| \rightarrow 0$. Notant $\overline{f}_2(x)$ la limite de $f(x_n')$, on a
		$$|\tilde{f}(x) - \tilde{f}_2(x)| \leqslant \|\tilde{f}(x) - f(x_n)\| + \|f(x_n) - f(x_n')\| + \|f(x_n') - \tilde{f}_2(x)\| \rightarrow 0$$
		donc $\overline{f}(x) = \overline{f}_2(x)$.
		Par ailleurs, $\overline{f}$ prolonge bien $f$, puisque si on se donne $x \in A$, il suffit de considérer la suite constante égale à $x$, la limite de la suite image étant égale à $f(x)$. \\
		Montrons enfin que $\overline{f}$ ainsi définie est continue. En fait, nous allons même montrer qu'elle est uniformément continue. Soit $\varepsilon > 0$ et $\delta > 0$ associé dans l'uniforme continuité de $f$. Soit ensuite $x, y \in \overline{A}$ tels que $\|x - y\| \leqslant \delta/2$. On choisit $x_n \rightarrow x, y_n \rightarrow y$, puis on écrit $$\|x_n - y_n\| \leqslant \|x_n - x\| + \|x - y\| + \|y - y_n\| \rightarrow \|x - y\|$$ donc à partir d'un certain rang on a $\|x_n - y_n\| \leqslant \delta$ puis $\|f(x_n) - f(y_n)\| \leqslant \varepsilon$. Or de même on a $$\|\tilde{f}(x) - \tilde{f}(y)\| \leqslant \|\tilde{f}(x) - f(x_n)\| + \varepsilon + \|f(y_n) - \tilde{f}(y)\|$$ puis par passage à la limite on obtient $\|\overline{f}(x) - \overline{f}(y)\| \leqslant \varepsilon$.
		\item Unicité. \\
		Supposons que $\overline{f}$ et $\overline{g}$ conviennent, et soit $x \in \overline{A}$. Écrivons $x = \lim x_n$ après quoi on a
		\begin{align*}
			|\tilde{f}(x) - \tilde{g}(x)| &\leqslant \|\tilde{f}(x) - f(x_n)\| + \|f(x_n) - g(x_n)\| + \|g(x_n) - \tilde{g}(x)\| \\ 
			&= \|\overline{f}(x) - \overline{f}(x_n)\| + \|\tilde{g}(x_n) - \tilde{g}(x)\|\\
			&\rightarrow 0
		\end{align*}
		par continuité de $\overline{f}$ et $\overline{g}$. Donc $|\tilde{f}(x) - \tilde{g}(x)| = 0$, puis comme cette égalité est valable pour tout $x$, on a $\tilde{f} = \tilde{g}$.
	\end{itemize}
Ainsi la preuve est achevée.
\end{proof}
Voici maintenant une application spectaculaire de ce procédé. On commence par considérer l'ensemble $A \subset \mathcal{B}([a,b], \mathbb{K})$ des fonctions en escalier, et dessus on construit une application $$I: \varphi \mapsto \int_{[a,b]} \varphi$$ dont on montre qu'elle est linéaire. L'inégalité triangulaire, quant à elle, montre que $I$ est $(b-a)$-lipschitzienne, et donc uniformément continue. Or, $\K=\R$ ou $\C$ est complet. Ainsi, il existe un unique prolongement continu de $I$ à l'ensemble $\overline{A}$ des fonctions réglées (qui contient en particulier les fonctions continues par morceaux), on l'appelle \textbf{intégrale de Riemann}. \par La linéarité de $I$ sur $A$ passe naturellement à la limite sur $\overline{A}$ (c'est un simple argument de continuité de $I$), et on obtient de la même manière l'inégalité triangulaire.


\newpage
\chapter{Analyse asymptotique}
\section{Relations de comparaison}
\subsection{Suites réelles}
Dans tout le paragraphe, on fixe $u$ et $v$ deux suites réelles indexées sur $\N$.
\begin{dfn}[Négligeabilité]
	$u$ est dite \textbf{négligeable devant $v$}, et on note $u_n\underset{n\to+\infty}{=}o(v_n)$ (voire $u_n=o(v_n)$) lorsque : $$\exists N\in\N,\ \exists(w_n)_{n\geqslant N}\in\R^{\llbracket N,+\infty\llbracket},\ \left(\forall n\geqslant N,\ u_n=v_nw_n\right) \ \ \text{et} \ \ w_n\xrightarrow[n\to+\infty]{} 0$$
\end{dfn}
\begin{dfn}[Domination]
	$u$ est dite \textbf{dominée par $v$}, et on note $u_n\underset{n\to+\infty}{=}O(v_n)$ (voire $u_n=O(v_n)$) lorsque : $$\exists N\in\N,\ \exists(w_n)_{n\geqslant N}\in\R^{\llbracket N,+\infty\llbracket},\ \left(\forall n\geqslant N,\ u_n=v_nw_n\right) \ \ \text{et} \ \ (w_n)_{n\geqslant N} \ \text{est bornée}$$
\end{dfn}
\begin{dfn}[Équivalence]
	$u$ est dite \textbf{équivalente à $v$}, et on note $u_n\underset{n\to+\infty}{\sim}v_n$ (voire $u_n\sim v_n$) lorsque : $$\exists N\in\N,\ \exists(w_n)_{n\geqslant N}\in\R^{\llbracket N,+\infty\llbracket},\ \left(\forall n\geqslant N,\ u_n=v_nw_n\right) \ \ \text{et} \ \ w_n\xrightarrow[n\to+\infty]{} 1$$
\end{dfn}
\begin{rmq}
	Attention, les relations $\cdot=o(\cdot)$ et $\cdot=O(\cdot)$ ne sont pas des relations d'équivalence. Elles n'ont aucune raison d'être symétriques notamment. Et ce n'est pas parce que deux suites sont un petit o d'une même suite qu'elles sont égales !
\end{rmq}
\begin{exa}[$u_n\sim 0$ ssi $(u_n)$ nulle APCR]
	On a toujours $0=o(v_n)$ mais la réciproque est beaucoup plus rare. En effet, on peut montrer que $u_n=o(0)$ si, et seulement si, $(u_n)$ est nulle APCR. Ce dernier fait est aussi vrai avec les grands O et l'équivalence.
\end{exa}
\begin{thm}
	On a : $$u_n\sim v_n \Longleftrightarrow u_n=v_n+o(v_n)$$
\end{thm}
\begin{thm}[Caractérisation par le quotient]
	Supposons que $v$ est non nulle APCR $N\i\N$. Alors :
	\begin{enumerate}
		\item $u_n=o(v_n)\Longleftrightarrow \displaystyle\frac{u_n}{v_n}\xrightarrow[n\to+\infty]{}0$
		\item $u_n=O(v_n)\Longleftrightarrow \displaystyle\left(\frac{u_n}{v_n}\right)_{n\geqslant N} \ \text{est bornée}$
		\item $u_n\sim v_n\Longleftrightarrow \displaystyle\frac{u_n}{v_n}\xrightarrow[n\to+\infty]{}1$
	\end{enumerate}
\end{thm}
\begin{exa}[$u_n=o(1)$]
	On a $u_n=o(1)\Longleftrightarrow u_n\xrightarrow[n\to+\infty]{}0$. Plus généralement, on a $$\forall \lambda\in\R\st, u_n=o(\lambda)\Longleftrightarrow u_n\xrightarrow[n\to+\infty]{}0$$
\end{exa}
\begin{exa}[$u_n=O(1)$]
	On a $u_n=O(1)\Longleftrightarrow (u_n) \ \text{est bornée}$. Plus généralement, on a $$\forall \lambda\in\R\st, u_n=O(\lambda)\Longleftrightarrow (u_n) \ \text{est bornée}$$
\end{exa}
\begin{exa}
	Soit $\alpha>\beta\geqslant0$ des réels. On a $n^\beta=o(n^\alpha)$ et $\displaystyle\frac{1}{n^\alpha}=o\left(\frac{1}{n^\beta}\right)$
\end{exa}
\begin{thm}[Obtention d'un équivalent par encadrement]
	Soit $u,v$ et $w$ des suites réelles et $\alpha$ une suite à valeurs dans $\R_+\st$ telles que $u_n\sim\alpha_n$, $w_n\sim\alpha_n$ et $u\leqslant v\leqslant w$ APCR. Alors, on a : $$v_n\sim\alpha_n$$
\end{thm}
\begin{rmq}
	On peut écrire les choses de façon plus condensée. Par exemple, si on veut écrire "Si $u_n=o(w_n)$ et $v_n=o(w_n)$, alors $u_n+v_n=o(w_n)$, on pourra écrire $o(w_n)+o(w_n)=o(w_n)$. Attention, cette notation est \textbf{symbolique} et \textbf{doit être lue de gauche à droite}. En particulier, le symbole égal n'est alors plus symétrique. Cependant, cela permet de condenser les écritures.
\end{rmq}
Dans ce contexte, on donne le théorème qui suit :
\begin{thm}[Composition des symboles]
	On a les relations \textbf{symboliques} (mais vraies si correctement reformulées) :
	\begin{enumerate}
		\item $o(u_n)=O(u_n)$
		\item $u_n\sim v_n \implies u_n=O(v_n)$
		\item $\alpha_n o(v_n)=o(\alpha_n v_n)$
		\item $o(\alpha_n v_n)=\alpha_n o(v_n)$
		\item $\alpha_n O(v_n)=O(\alpha_n v_n)$
		\item $O(\alpha_n v_n)=\alpha_n O(v_n)$
		\item $o(v_n)+o(v_n)=o(v_n)$
		\item $(v_n)+O(v_n)=O(v_n)$
		\item $O(v_n)+O(v_n)=O(v_n)$
		\item $o(o(v_n))=o(v_n)$
		\item $o(O(v_n))=o(v_n)$
		\item $O(o(v_n))=o(v_n)$
		\item $O(O(v_n))=O(v_n)$
		\item $o(u_n)o(v_n)=o(u_n v_n)$
		\item $O(u_n)O(v_n)=O(u_n v_n)$
	\end{enumerate}
\end{thm}
\begin{rmq}
	Attention ! $o(u_n)+o(v_n)\ne o(u_n+v_n)$ en général ! De même, $O(u_n)+O(v_n)\ne O(u_n+v_n)$ en général !
\end{rmq}
\begin{thm}
	$\sim$ est une relation d'équivalence sur $\R^\N$.
\end{thm}
\begin{rmq}
	\textbf{Il ne sert à rien d'écrire un équivalent à plus d'un terme !} Par exemple, écrire $u_n\sim 1+\displaystyle\frac1n$ revient à dire que $u_n\sim 1$ voire $u_n\sim 1-\displaystyle\frac1n$ par transitivité de $\sim$.
\end{rmq}
\begin{rmq}
	$\forall \lambda\in\R\st,\ u_n\sim\lambda \Longleftrightarrow u_n\xrightarrow[n\to+\infty]{}\lambda$
\end{rmq}
\begin{thm}[Équivalence et héritage]
	Supposons que $u_n\sim v_n$. Alors la limite s'hérite par équivalence : $$\forall l\in\overline{\R},\ u_n\xrightarrow[n\to+\infty]{}l \Longleftrightarrow v_n\xrightarrow[n\to+\infty]{}l$$ Aussi, le signe s'hérite par équivalence : \begin{enumerate}
		\item $u_n\ne0 \ \text{APCR} \Longleftrightarrow v_n\ne0 \ \text{APCR}$
		\item $u_n\geqslant0 \ \text{APCR} \Longleftrightarrow v_n\geqslant0 \ \text{APCR}$
		\item $u_n\leqslant0 \ \text{APCR} \Longleftrightarrow v_n\leqslant0 \ \text{APCR}$
		\item $u_n>0 \ \text{APCR} \Longleftrightarrow v_n>0 \ \text{APCR}$
		\item $u_n<0 \ \text{APCR} \Longleftrightarrow v_n<0 \ \text{APCR}$
	\end{enumerate}
	Enfin, soit $w\in\R^\N$. Alors, les relations de comparaison s'héritent :
	\begin{enumerate}
		\item $w_n=o(v_n)\Longleftrightarrow w_n=o(v_n)$
		\item $w_n=O(v_n)\Longleftrightarrow w_n=O(v_n)$
		\item $u_n=o(w_n)\Longleftrightarrow v_n=o(w_n)$
		\item $u_n=O(w_n)\Longleftrightarrow v_n=O(w_n)$
	\end{enumerate}
\end{thm}
\begin{thm}[Produit, quotient et puissances d'équivalents]
	Soit $a,b,c,d\in\R^\N$. Soit $\alpha\in\R$. Alors :
	\begin{enumerate}
		\item Si $a_n\sim b_n$ et $c_n\sim d_n$, alors $a_n c_n\sim b_n d_n$.
		\item Supposons que $c$ et $d$ soient non nulles APCR. Si $a_n\sim b_n$ et $c_n\sim d_n$, alors $\displaystyle\frac{a_n}{c_n}\sim \frac{b_n}{d_n}$.
		\item Supposons que $a$ et $b$ sont à valeurs dans $\R_+\st$ APCR. Si $a_n\sim b_n$, alors $a_n^\alpha\sim b_n^\alpha$.
	\end{enumerate}
\end{thm}
\begin{rmq}[Billet gratuit pour la 5/2]
	\textbf{Attention ! On ne somme pas les équivalents !!!}
\end{rmq}
\begin{thm}[HP]
	Si $u_n\sim v_n$ et $u_n\xrightarrow[n\to+\infty]{}+\infty$ (et donc de même pour $v_n$), alors : $$\ln(u_n)\sim\ln(v_n)$$
\end{thm}
\begin{proof}
	A refaire à chaque fois. On a $u_n=v_n+o(v_n)$. Donc en factorisant dans le $\ln$ et en utilisant la propriété fonctionnelle du $\ln$, on obtient : $$\ln(u_n)=\ln(v_n)+\ln(1+o(1))$$ Or, $1+o(1)$ tend vers $1$, donc par continuité du $\ln$ en $1$, $\ln(1+o(1))=o(1)$. Puisque $v_n$ tend vers $+\infty$, $\ln(v_n)$ aussi. A fortiori, on a donc $\ln(1+o(1))=o(\ln(v_n))$ (par une caractérisation par le quotient par exemple). D'où : $$\ln(u_n)=\ln(v_n)+o(\ln(v_n))$$ Autrement dit, $$\ln(u_n)\sim\ln(v_n)$$
\end{proof}
\begin{thm}[Formule de Stirling]
	On a :$$n!\sim \left(\frac{n}{e}\right)^n\sqrt{2\pi n}$$
\end{thm}
\begin{exa}
	Avec le théorème HP qui précède, on obtient facilement que $\ln(n!)\sim n\ln(n)$.
\end{exa}
\subsection{Brève extension aux suites complexes}
\begin{dfn}[Extension des définitions à $\C$]
	Soit $u,v\in\C^\N$. On définit de même $u_n=o(v_n)$, $u_n=O(v_n)$ et $u_n\sim v_n$ à ceci près que la suite $(w_n)$ est désormais à valeurs dans $\C$.
\end{dfn}
\begin{thm}[Extension des résultats à $\C$]
	Dans $\C$, les résultats suivants restent vrais :
	\begin{enumerate}
		\item $u_n\sim v_n \Longleftrightarrow u_n=v_n+o(v_n)$
		\item $\sim$ reste une relation d'équivalence
		\item Caractérisation par le quotient
		\item Composition des symboles
		\item Produit et quotient d'équivalents (les puissances restent vraies mais nous ramènent par hypothèse dans $\R_+\st$)
		\item Les héritages, exceptés ceux qui font intervenir des symboles d'inégalités et perdent alors leur sens.
	\end{enumerate}
\end{thm}
\subsection{Généralisation aux fonctions}
Dans cette section, on fixe $X\subset\R$ tel que $X\ne\emptyset$ ainsi que $a\in\overline{X}$. Les fonctions seront considérées de $X$ dans $\K=\R$ ou $\K=\C$.
\begin{dfn}
	Soit $f,g:X\rightarrow\K$. 
	\begin{enumerate}
		\item On écrira $f(x)\underset{x\to a}{=}o\left(g(x)\right)$ lorsque $$\exists V\in\mathcal{V}(a),\ \exists h:X\cap V\rightarrow\K,\ \left(\forall x\in X\cap V,\ f(x)=g(x)h(x)\right) \ \ \text{et} \ \ h(x)\xrightarrow[x\to a]{} 0$$
		\item On écrira $f(x)\underset{x\to a}{=}O\left(g(x)\right)$ lorsque $$\exists V\in\mathcal{V}(a),\ \exists h:X\cap V\rightarrow\K,\ \left(\forall x\in X\cap V,\ f(x)=g(x)h(x)\right) \ \ \text{et} \ \ h \ \text{est bornée}$$
		\item On écrira $f(x)\underset{x\to a}{\sim}o\left(g(x)\right)$ lorsque $$\exists V\in\mathcal{V}(a),\ \exists h:X\cap V\rightarrow\K,\ \left(\forall x\in X\cap V,\ f(x)=g(x)h(x)\right) \ \ \text{et} \ \ h(x)\xrightarrow[x\to a]{} 1$$
	\end{enumerate}
\end{dfn}
\begin{prop}[Caractère local des relations de comparaison]
	On garde $f$ et $g$ comme précédemment. Soit $\tilde{V}\in\mathcal{V}(a)$. Posons $\displaystyle \tilde{f}=f_{|X\cap\tilde{V}}$ et $\displaystyle \tilde{g}=g_{|X\cap\tilde{V}}$. Alors :
	\begin{enumerate}
		\item $f(x)\underset{x\to a}{=}o\left(g(x)\right) \Longleftrightarrow \tilde{f}(x)\underset{x\to a}{=}o\left(\tilde{g}(x)\right)$
		\item $f(x)\underset{x\to a}{=}O\left(g(x)\right) \Longleftrightarrow \tilde{f}(x)\underset{x\to a}{=}O\left(\tilde{g}(x)\right)$
		\item $f(x)\underset{x\to a}{\sim}g(x) \Longleftrightarrow \tilde{f}(x)\underset{x\to a}{\sim}\tilde{g}(x)$
	\end{enumerate}
\end{prop}
\begin{rmq}
	$f(x)\underset{x\to a}{\sim} 0$ si, et seulement si, $f$ est nulle au voisinage de $a$
\end{rmq}
\begin{rmq}
	Soit $\lambda\in\K\st$. On a $f(x)\underset{x\to a}{\sim}\lambda$ si, et seulement si, $f(x)\xrightarrow[x\to a]{}\lambda$
\end{rmq}
\begin{prop}
	On a $$f(x)\underset{x\to a}{\sim}g(x) \Longleftrightarrow f(x)\underset{x\to a}{=}g(x)+o\left(g(x)\right)$$
\end{prop}
\begin{rmq}
	$\sim$ est une relation d'équivalence.
\end{rmq}
\begin{rmq}[Caractérisation par le quotient]
	La caractérisation par le quotient reste valable tant que $g$ ne s'annule pas au voisinage de $a$.
\end{rmq}
\begin{rmq}
	Soit $\alpha\geqslant\beta>0$. On a $x^\beta \underset{x\to+\infty}{=}o(x^\alpha)$
\end{rmq}
\begin{exa}
	Soit $f:x\rightarrow a_nx^n+\ldots+a_0$ une fonction polynomiale (avec $n\in\N$ et $a_n\ne 0$). On a $$f(x)\underset{x\to+\infty}{\sim}a_nx^n$$ et $$f(x)\underset{x\to-\infty}{\sim}a_nx^n$$
\end{exa}
\begin{thm}
	Les résultats généraux suivants vus sur les suites restent valables :
	\begin{enumerate}
		\item Composition des symboles
		\item Équivalents qui passent au produit, au quotient et à une puissance \textbf{indépendante de $x$}
		\item Héritages par équivalence
		\item Obtention d'un équivalent par encadrement
	\end{enumerate}
\end{thm}
\begin{exa}
	Soit $(p,q)\in\N^2$ tel que $p<q$. En revenant à la définition, on montre que $$x^q\underset{x\to 0}{=}o(x^p)$$
\end{exa}
\begin{exa}
	Considérons la fonction rationnelle $\displaystyle f:x\rightarrow\frac{a_nx^n+\ldots+a_0}{b_px^p+\ldots+b_0}$ (avec $n,p\in\N$ et $a_n\ne0$ et $b_p\ne0$). On a $$f(x)\underset{x\to+\infty}{\sim}\frac{a_n}{b_p}x^{n-p}$$ et $$f(x)\underset{x\to-\infty}{\sim}\frac{a_n}{b_p}x^{n-p}$$
\end{exa}
\subsection{Petit retour sur les croissances comparées}
\begin{thm}[Croissances comparées pour les fonctions]
	Soit $\alpha,\beta,\gamma>0$ Au voisinage de $+\infty$, on a :
	\begin{enumerate}
		\item $\ln^\alpha(x)=o(x^\beta)$
		\item $x^\beta=o(\exp(\gamma x))$
	\end{enumerate}
	Et de plus, ces trois expressions tendent vers $+\infty$ quand $x$ tend vers $+\infty$.
\end{thm}
\begin{cor}
	On en déduit : $$\lvert\ln^\alpha(x)\rvert x^\beta \xrightarrow[x\to 0^+]{}0$$ et $$\lvert x^\beta\rvert \exp(\gamma x)\xrightarrow[x\to-\infty]{}0$$
\end{cor}
\begin{thm}[Croissances comparées pour les suites]
	Soit $\alpha,\beta,\gamma>0$ Au voisinage de $+\infty$, on a :
	\begin{enumerate}
		\item $\ln^\alpha(n)=o(n^\beta)$
		\item $n^\beta=o(\exp(\gamma n))$
		\item $\exp(\gamma n)=o(n!)$
		\item $n!=o(n^n)$
	\end{enumerate}
	Et de plus, ces cinq expressions tendent vers $+\infty$ quand $n$ tend vers $+\infty$.
\end{thm}
\section{Développements limités}
Dans tout le paragraphe, on fixe : $I$ un intervalle non trivial de $\R$, $a\in I$, $f:I\rightarrow\K$ avec $\K=\R$ ou $\K=\C$ ainsi que $n\in\N$.
\subsection{Premières définitions}
\begin{dfn}[Développement limité à l'ordre $n$ en $a$]
	On dit que $f$ admet un \textbf{développement limité à l'ordre $n$ en $a$/au point $a$} (noté $DL_n(a)$) lorsque : $$\exists(c_0,\ldots,c_n)\in\K^{n+1}, \ f(x)\underset{x\to a}{=}\sum_{k=0}^{n}c_k(x-a)^k+o\left((x-a)^n\right)$$
\end{dfn}
\begin{rmq}
	Sans perte de généralité, on pourra toujours supposer que la fonction $\varepsilon$ que contient le $o$ de la définition est définie sur $I$ tout entier, et non sur un voisinage de $a$, puisque de toute façon la définition du $o$ est \textbf{locale}.
\end{rmq}
\begin{thm}[Cas $n=0$, équivalence avec la continuité]
	$f$ admet un $DL_0(a)$ si, et seulement si, $f$ est continue en $a$. Dans ce cas, ce DL est nécessairement donnée par : $$f(x)\underset{x\to a}{=}f(a)+o(1)$$ ie on a $c_0=f(a)$ dans la définition.
\end{thm}
\begin{thm}[Cas $n=1$, équivalence avec la dérivabilité]
	$f$ admet un $DL_1(a)$ si, et seulement si, $f$ est dérivable en $a$. Dans ce cas, ce DL est nécessairement donnée par : $$f(x)\underset{x\to a}{=}f(a)+f'(a)(x-a)+o(x-a)$$ ie on a $c_0=f(a)$ et $c_1=f'(a)$ dans la définition.
\end{thm}
\begin{rmq}
	\textbf{Attention !} Pour $n\geqslant2$, on ne peut plus rien dire ! Par exemple, la fonction de $\R$ dans $\R$ qui à $x$ associe $\displaystyle x^3\sin\left(\frac{1}{x}\right)$ prolongée par $0$ en $0$ admet un $DL_2(0)$ (poser $c_0=c_1=c_2=0$) mais n'est pourtant pas deux fois dérivable en $0$.
\end{rmq}
\begin{thm}[Lemme de troncature]
	Supposons que $f$ admette un $DL_n(a)$ qu'on écrit sous la forme : $$f(x)\underset{x\to a}{=}\sum_{k=0}^{n}c_k(x-a)^k+o\left((x-a)^n\right)$$ Alors, pour tout $m\in\llbracket0,n\rrbracket$, $f$ admet un $DL_m(a)$ donné par : $$f(x)\underset{x\to a}{=}\sum_{k=0}^{m}c_k(x-a)^k+o\left((x-a)^m\right)$$
\end{thm}
\begin{thm}[Unicité]
	Supposons que $f$ admette un $DL_n(a)$ qu'on écrit sous la forme : $$f(x)\underset{x\to a}{=}\sum_{k=0}^{n}c_k(x-a)^k+o\left((x-a)^n\right)$$ Alors, les scalaires $c_0,\ldots,c_n$ sont uniques.
\end{thm}
\begin{dfn}[Partie régulière]
	La fonction polynomiale $$
	\begin{array}[t]{lrcl}
		I & \longrightarrow & \K \\
		x & \longmapsto & \displaystyle \sum_{k=0}^{n}c_k(x-a)^k
	\end{array}$$ est donc définie de façon unique. On l'appelle \textbf{la partie régulière du $DL_n(a)$}.
\end{dfn}
\begin{cor}[Fonctions paires]
	Supposons que $f$ est paire et que $f$ admet un $DL_n(0)$ qu'on écrit sous la forme : $$f(x)\underset{x\to 0}{=}\sum_{k=0}^{n}c_kx^k+o\left(x^n\right)$$ Alors, tous les $c_k$ pour $k$ \textbf{impair} sont nuls.
\end{cor}
\begin{cor}[Fonctions impaires]
	Supposons que $f$ est impaire et que $f$ admet un $DL_n(0)$ qu'on écrit sous la forme : $$f(x)\underset{x\to 0}{=}\sum_{k=0}^{n}c_kx^k+o\left(x^n\right)$$ Alors, tous les $c_k$ pour $k$ \textbf{pair} sont nuls.
\end{cor}
\subsection{Calculs pratiques}
\begin{thm}[Combinaison linéaire]
	Soit $f$ et $g$ des fonctions de $I$ dans $\K$ admettant des $DL_n(a)$ qu'on écrit sous la forme : $$f(x)\underset{x\to a}{=}\sum_{k=0}^{n}c_k(x-a)^k+o\left((x-a)^n\right)$$ et $$g(x)\underset{x\to a}{=}\sum_{k=0}^{n}d_k(x-a)^k+o\left((x-a)^n\right)$$ Soit $\lambda,\mu\in\K$. Alors, $\lambda f+\mu g$ admet un $DL_n(a)$ donné par : $$(\lambda f+\mu g)(x)\underset{x\to a}{=}\sum_{k=0}^{n}(\lambda c_k+\mu d_k)(x-a)^k+o\left((x-a)^n\right)$$
\end{thm}
\begin{thm}[Produit]
	Soit $f$ et $g$ des fonctions de $I$ dans $\K$ admettant des $DL_n(a)$ qu'on écrit sous la forme : $$f(x)\underset{x\to a}{=}\sum_{k=0}^{n}c_k(x-a)^k+o\left((x-a)^n\right)$$ et $$g(x)\underset{x\to a}{=}\sum_{k=0}^{n}d_k(x-a)^k+o\left((x-a)^n\right)$$ Alors $fg$ admet un $DL_n(a)$ donné par : Soit $f$ et $g$ des fonctions de $I$ dans $\K$ admettant des $DL_n(a)$ qu'on écrit sous la forme : $$(fg)(x)\underset{x\to a}{=}\sum_{k=0}^{n}\left(\sum_{i+j=k}c_i d_j\right)(x-a)^k+o\left((x-a)^n\right)$$
\end{thm}
Pour la composition, aucun résultat théorique n'est exigible, on se contente de faire de la substitution.
\begin{prop}[Substitution dans un DL]
	Supposons que $f$ possède un $DL(0)$ qu'on écrit sous la forme : $$f(u)\underset{u\to 0}{=}\sum_{k=0}^{n}c_ku^k+o\left(u^n\right)$$ et supposons qu'on possède une expression $u(x)$ telle que $u(x)\xrightarrow[x\to a]{} 0$. Alors on a par \textbf{substitution} : $$f(u(x))\underset{x\to a}{=}\sum_{k=0}^{n}c_k(u(x))^k+o\left((u(x))^n\right)$$ Sur une copie ou en khôlle, les mentions "$u(x)\xrightarrow[x\to a]{} 0$" et "substitution" devront \textbf{toujours} apparaître.
\end{prop}
\begin{prop}[Substitution dans les relations de comparaison]
	Supposons que $f(u)\underset{u\to 0}{\sim}g(u)$ et supposons qu'on possède une expression $u(x)$ telle que $u(x)\xrightarrow[x\to a]{} 0$. Alors, par substitution, on a : $$f(u(x))\underset{x\to a}{\sim}g(u(x))$$ On a les résultats symétriques avec les petits $o$ et les grands $O$.
\end{prop}
\begin{prop}[Substitution de suites dans des DL ou relations de comparaison]
	Avec les mêmes hypothèses, mais cette fois-ci des suites qui tendent vers $0$, on a les mêmes résultats de substitution.
\end{prop}
\begin{prop}[Quotient]
	Soit $u:I\rightarrow\K$ telle que $u(x)\xrightarrow[x\to a]{}0$ et $u$ admette un $DL_n(a)$. Alors, $\displaystyle \frac{1}{1-u}$ admet un $DL_n(a)$. En fait, on a : $$\frac{1}{1-u(x)}\underset{x\to a}{=}1+\ldots+(u(x))^n+o\left((x-a)^n\right)$$ et chacun des $u^k$ admet un $DL_n(a)$ par produit. On conclut par combinaison linéaire.
\end{prop}
\begin{rmq}
	\textbf{Attention !} Le résultat précédent n'est d'aucune utilité pratique, car si on anticipe mal les ordres, on se retrouve avec des ordres en trop. Moralité : il faut \textbf{anticiper les ordres} pour se faciliter le travail.
\end{rmq}
\begin{rmq}[Recentrage en $0$]
	Soit $(c_0,\ldots,c_n)\in\K^{n+1}$. L'assertion : $$f(x)\underset{x\to a}{=}\sum_{k=0}^{n}c_k(x-a)^k+o\left((x-a)^n\right)$$ est équivalente à l'assertion : $$f(a+h)\underset{h\to 0}{=}\sum_{k=0}^{n}c_kh^k+o\left(h^n\right)$$ On peut donc toujours se recentrer en $0$ pour se ramener à des DL usuels.
\end{rmq}
\begin{thm}[Forme normalisée, "règle du $c_mh^m$ (NS pour cette dernière)]
	On suppose que $$f(a+h)\underset{h\to 0}{=}\sum_{k=0}^{n}c_kh^k+o\left(h^n\right)$$ Supposons que $\exists k\in\llbracket0,n\rrbracket,\ c_k\ne0$. Considérons alors $m$ le plus petit entier tel que $c_m$ soit non nul. Alors on a : $$f(a+h)\underset{h\to 0}{=}h^m\left(c_m+c_{m+1}h+\ldots+c_nh^{n-m}+o\left(h^{n-m}\right)\right)$$ C'est ce qu'on appelle la \textbf{forme normalisée du DL}. De plus, par troncature, on a : $$f(a+h)\underset{h\to 0}{=}c_mh^m+o(h^m)$$ donc on en déduit que $$f(a+h)\underset{h\to 0}{\sim}c_mh^m$$
\end{thm}
\begin{exa}
	Soit $x$ au voisinage de $0$. On a : $$\forall \alpha\in\R,\ (1+x)^\alpha-1\sim\alpha x$$ $$\exp(x)-1\sim x$$ $$\ln(1+x)\sim x$$ $$\cos(x)-1\sim\frac{x^2}{2}$$ $$\sin(x)\sim x$$ $$\tan(x)\sim x$$ $$\arctan(x)\sim x$$
\end{exa}
\subsection{Résultats théoriques}
\begin{thm}
	Supposons que $f$ admette un $DL_n(a)$ qu'on écrit sous la forme : $$f(x)\underset{x\to a}{=}\sum_{k=0}^{n}c_k(x-a)^k+o\left((x-a)^n\right)$$ Supposons que $f$ admette une primitive $F$. Alors, $F$ admet un $DL_{n+1}(a)$ donné par : $$F(x)\underset{x\to a}{=}F(a)+\sum_{k=0}^{n}\frac{c_k}{k+1}(x-a)^{k+1}+o\left((x-a)^{n+1}\right)$$
\end{thm}
\begin{cor}[Formule de Taylor-Young]
	Soit $f\in\mathcal{C}^n(I,\K)$. Alors, quel que soit $a\in I$, $f$ admet un $DL_n(a)$ donné par : $$f(x)\underset{x\to a}{=}\sum_{k=0}^{n}\frac{f^{(k)}(a)}{k!}(x-a)^k+o\left((x-a)^n\right)$$ soit encore : $$f(a+h)\underset{h\to 0}{=}\sum_{k=0}^{n}\frac{f^{(k)}(a)}{k!}h^k+o\left(h^n\right)$$ La partie régulière de la première forme de ce DL s'appelle le \textbf{polynôme de Taylor de $f$ à l'ordre $n$ en $a$}.
\end{cor}
\begin{rmq}
	On rappelle qu'admettre un $DL_n$ ne rend pas $\mathcal{C}^n$ !!!
\end{rmq}
\begin{rmq}
	Une fonction de classe $\mathcal{C}^\infty$ admet donc des DL à tous les ordres.
\end{rmq}
\begin{prop}[Dérivée d'un $DL_n$ pour une fonction $\mathcal{C}^n$ avec $n\in\N\st$]
	En général, on ne dérive pas un DL. Cependant, supposons que $n\in\N\st$ et $f$ soit de classe $\mathcal{C}^n$. Alors, grâce à la formule de Taylor-Young, si on part d'un DL de $f$ sous la forme $$f(x)\underset{x\to a}{=}\sum_{k=0}^{n}c_k(x-a)^k+o\left((x-a)^n\right)$$ avec $$\forall k\in\llbracket0,n\rrbracket,\ c_k=\frac{f^{(k)}(a)}{k!}$$, on obtient alors $$f'(x)\underset{x\to a}{=}\sum_{k=1}^{n}kc_k(x-a)^{k-1}+o\left((x-a)^{n-1}\right)$$ sachant que pour mémoire, on avait $$f(x)\underset{x\to a}{=}f(a)+\sum_{k=1}^{n}c_k(x-a)^{k}+o\left((x-a)^{n}\right)$$ Virtuellement, tout se passe comme si on avait dérivé le $DL_n(a)$ de $f$ pour obtenir de le $DL_{n-1}(a)$ de $f'$. \textbf{Sur une copie, on pourra utiliser ce résultat à la condition d'écrire : $n\in\N\st$, $f$ est de classe $\mathcal{C}^n$ et "formule de Taylor-Young".}
\end{prop}
\begin{thm}[DL usuels au voisinage de $0$, à savoir par cœur]
	On a les DL usuels suivants au voisinage de $0$, \textbf{à connaître par cœur} : $$\frac{1}{1-x} \underset{x\to 0}{=} 1+x+...+x^n+o(x^n)$$ $$\frac{1}{1+x} \underset{x\to 0}{=} 1-x+...+(-1)^n x^n+o(x^n)$$ $$(1+x)^{\alpha} \underset{x\to 0}{=} 1 + \alpha x + \frac{\alpha(\alpha -1)}{2}x^2 + ... + \frac{\alpha(\alpha -1)...(\alpha -n+1)}{n!}x^n +o(x^n)$$ $$e^x \underset{x\to 0}{=} 1 + x + \frac{x^2}{2!}+...+\frac{x^n}{n!} + o(x^n)$$ $$\ln(1+x) \underset{x\to 0}{=} x - \frac{x^2}{2} + \frac{x^3}{3}+...+ (-1)^{n-1}\frac{x^n}{n} + o(x^n)$$ $$\cos(x)\underset{x\to 0}{=}1-\frac{x^2}{2!}+\frac{x^4}{4!}+\ldots+(-1)^n\frac{x^{2n}}{(2n)!}+o(x^{2n+1})$$ $$\sin(x)\underset{x\to 0}{=}x-\frac{x^3}{3!}+\frac{x^5}{5!}+\ldots+(-1)^n\frac{x^{2n+1}}{(2n+1)!}+o(x^{2n+2})$$ $$\cosh(x)\underset{x\to 0}{=}1+\frac{x^2}{2!}+\frac{x^4}{4!}+\ldots+\frac{x^{2n}}{(2n)!}+o(x^{2n+1})$$ $$\sinh(x)\underset{x\to 0}{=}x+\frac{x^3}{3!}+\frac{x^5}{5!}+\ldots+(-1)^n\frac{x^{2n+1}}{(2n+1)!}+o(x^{2n+2})$$ $$\arctan(x)\underset{x\to 0}{=}x-\frac{x^3}{3}+\frac{x^5}{5}+\ldots+(-1)^n\frac{x^{2n+1}}{2n+1}+o(x^{2n+2})$$ $$\tan(x)=x+\frac{1}{3}x^3+\frac{2}{15}x^5+\frac{17}{315}x^7+o(x^8)$$ $$\tanh(x)=x-\frac{1}{3}x^3+\frac{2}{15}x^5-\frac{17}{315}x^7+o(x^8)$$
\end{thm}
\begin{dfn}[DL au sens fort]
	Soit $n\in\N\st$ et $x$ au voisinage de $0$. On a : $$\ln(1+x) \underset{x\to 0}{=} x - \frac{x^2}{2} + \frac{x^3}{3}+...+ (-1)^{n-1}\frac{x^n}{n} + (-1)^{n}\frac{x^{n+1}}{n+1}+o(x^{n+1})$$ ce qu'on peut aussi réécrire : $$\ln(1+x) \underset{x\to 0}{=} x - \frac{x^2}{2} + \frac{x^3}{3}+...+ (-1)^{n-1}\frac{x^n}{n} + O(x^{n+1})$$ C'est ce qu'on appelle un \textbf{DL au sens fort}. En effet, $O(x^{n+1})\Longrightarrow o(x^n)$ mais pas l'inverse.
\end{dfn}
\begin{dfn}[Développement asymptotique]
	Il s'agit de la même idée que le DL, mais en bord de domaine avec $a\in\overline{I}\setminus I$. On n'a plus nécessairement des puissances de $x$, mais cela peut être des fonctions de forces croissantes avec par exemple les différents éléments des croissances comparées.
\end{dfn}
\begin{met}[DA de suites implicites]
	Voici la méthode à appliquer pour étudier le DA d'une suite définie de façon implicite (en vogue aux oraux des Mines) :
	\begin{enumerate}
		\item Prouver la bonne définition de la suite. Traditionnellement, un tableau de variation suffit pour prouver l'existence et l'unicité.
		\item On détermine la limite de la suite $(u_n)$ définie de façon implicite On peut par exemple l'interpréter comme la bijection réciproque de la fonction qui permet de la définir de façon implicite. Souvent, cette limite vaut $+\infty$.
		\item On cherche un moteur, c'est-à-dire une équation qui mette en jeu $u_n$ et un $o(u_n)$. Cette équation nous permettra de gagner peu à peu des ordres dans le DA. Souvent, ce moteur s'obtient soit avec la définition implicite de la suite, soit en passant cette définition au $\ln$ ou à d'autres fonctions du même genre.
		\item Il faut initialiser le moteur, donc chercher un premier terme de Dl ou de DA. Souvent, un équivalent obtenu directement ou indirectement à partir de la limite fait l'affaire.
	\end{enumerate}
\end{met}
\begin{exa}
	Pour mettre en œuvre cette méthode, se référer aux exercices 22, 23 et 24 du TD16.
\end{exa}
\subsection{Application aux études de courbes}
Soit $f:I\rightarrow\R$ et $a\in I$. Supposons que $f$ admette un $DL_1(a)$ : $f(x)=\underset{x\to a}{=}c_0+c_1(x-a)+o(x-a)$. Alors $f$ est dérivable en $a$ et on a $f(a)=c_0$ et $f'(a)=c_1$. On rappelle que l'équation de la tangente à la courbe représentative de $f$ en $a$ est alors donnée par $y-f(a)=f'(a)(x-a)$ soit encore $$y=c_0+c_1(x-a)$$ Ainsi, on "lit" l'équation de la tangente à la courbe de $f$ en $a$ dans le $DL_1(a)$. Pousser le DL plus loin nous donne un équivalent de $f-y_{tangente}$, ce qui nous donne localement les positions relatives de la courbe et de la tangente.
\begin{thm}[CN et CS sur la nature des extrema locaux]
	Supposons que $f$ soit $\mathcal{C}^2$ au voisinage de $a$ et que $f'(a)=0$. On a une condition nécessaire sur la nature de l'extremum local en $a$ :
	\begin{enumerate}
		\item Si $f$ admet un minimum local en $a$, alors $f''(a)\geqslant0$.
		\item Si $f$ admet un maximum local en $a$, alors $f''(a)\leqslant0$.
	\end{enumerate}
	On a ensuite une condition suffisante sur la nature d'un potentiel extremum local en $a$ :
	\begin{enumerate}
		\item Si $f''(a)>0$, alors $f$ admet un minimum local en $a$.
		\item Si $f''(a)<0$, alors $f$ admet un maximum local en $a$.
	\end{enumerate}
\end{thm}
\begin{rmq}
	Si $f''(a)=0$, on ne peut rien dire. En guise de contre-exemple, considérer $x\mapsto x^3$, $x\mapsto x^4$ et $x\mapsto -x^4$ en $0$.
\end{rmq}
\begin{dfn}[Asymptote au voisinage de $+\infty$ et $-\infty$]
	Supposons que $f$ admette un développement asymptotique au voisinage de $+\infty$ de la forme : $$f(x)\underset{x\to+\infty}{=}ax+b+o(1)$$ Alors on dit que la droite $y=ax+b$ est \textbf{asymptote à la courbe au voisinage de $+\infty$}. On définit de même les asymptotes au voisinage de $-\infty$.
\end{dfn}
\newpage
\chapter{Introduction aux équations différentielles}
Dans tout le chapitre, $I$ désigne un intervalle non trivial de $\R$ et toutes les fonctions vont de $I$ dans $\K=\R$ ou $\C$.

\section{Généralités}
\begin{dfn}[Équation différentielle linéaire d'ordre $n$]
	Une \textbf{équation différentielle linéaire d'ordre $n\in\N\st$} est une équation dont l'inconnue est une fonction $y\in\mathcal{D}^n(I,\K)$. Elle s'écrit sous la forme $$\sum_{k=a}^{n}a_k(t)y^{(k)}=b(t)$$ où $a_1, ...,a_n,b$ sont des fonctions continues de $I$ dans $\K$ et où $a_n$ n'est pas la fonction nulle. Plus précisément, cela signifie que $y$ vérifie $$\forall t \in I, \ \sum_{k=a}^{n}a_k(t)y^{(k)}(t)=b(t)$$
\end{dfn}
\begin{dfn}[EDL à coefficients constants]
	L'EDL est dite \textbf{"à coefficients constants"} lorsque les $a_k$ sont des fonctions constantes.
\end{dfn}
\begin{dfn}[Forme résolue]
	On dit que l'équation $$\sum_{k=a}^{n}a_k(t)y^{(k)}=b(t)$$ est sous \textbf{forme résolue} lorsque $a_n=1$. En pratique, on peut s'y ramener dès lors que $a_n$ ne s'annule pas : il suffit alors de diviser par $a_n$.
\end{dfn}
\begin{dfn}[Équation homogène]
	L'équation \textbf{homogène} associée à l'équation différentielle linéaire d'ordre $n$ $$\sum_{k=a}^{n}a_k(t)y^{(k)}=b(t)$$ est l'équation $$\sum_{k=a}^{n}a_k(t)y^{(k)}=0$$
\end{dfn}
\begin{prop}[Structure d'espace affin des solutions, si elles existent]
	La solution générale d'une équation différentielle linéaire, si elle existe, est la somme d'une solution particulière de l'équation avec second membre et de la solution générale de l'équation homogène associée.
\end{prop}
\begin{prop}
	On cherche une solution particulière de $$\sum_{k=a}^{n}a_k(t)y^{(k)}=\lambda b_1(t)+\mu b_2(t)$$ il suffit alors de trouver une solution particulière $y_1$ de $\displaystyle \sum_{k=a}^{n}a_k(t)y^{(k)}=b_1(t)$ et une solution particulière $y_2$ de $\displaystyle \sum_{k=a}^{n}a_k(t)y^{(k)}=b_2(t)$. Alors, puisque l'équation est linéaire, la fonction $y_p=\lambda y_1+\mu y_2$ convient pour le problème initial.
\end{prop}
\begin{rmq} Pour le $\cos$ et le $\sin$, il peut être plus rapide de les interpréter comme des parties réelle et imaginaire d'une exponentielle complexe. On résout alors l'équation de $\C$ puis on passe à la partie réelle ou à la partie imaginaire. \end{rmq}
\section{Équations différentielles linéaires d'ordre 1}
\begin{thm}[Résolution de $y'+a(t)y=0$]
	Soit $A$ une primitive quelconque de $a$. Les solutions sont données par $$\forall t \in I, \ y(t)=\lambda \exp(-A(t))$$ avec $\lambda \in\K$ une constante. Les solutions forment un $\K$-espace vectoriel de dimension 1.
\end{thm}
\begin{thm}[Résolution de $y'+a(t)y=b(t)$]
	Déjà, cette équation admet au moins une solution.
	On commence par résoudre l'équation homogène associée : $y'+a(t)y=0$. Ensuite, on fait varier la constante en écrivant $y$ sous la forme $y(t)=\lambda(t)\exp(-A(t))$. On détermine alors $\lambda(t)$ en réinjectant dans l'équation générale car on obtient $\lambda'(t)\exp(-A(t))=b(t)$ et il suffit alors d'isoler $\lambda'(t)$ et de primitiver. On remonte alors à $y$ et on obtient la solution générale.
\end{thm}
\begin{cor}[Existence et unicité d'une solution à un problème de Cauchy]
	Soit $t_0\in I$ et $C_0\in\K$. Alors il existe une unique fonction de $I$ dan $\K$ vérifiant $y'+a(t)y=b(t)$ \textbf{et} la condition de Cauchy $y(t_0)=C_0$.
\end{cor}
\begin{dfn}[Cas général $a(t)y'+b(t)y=c(t)$, recollement]
	Dans le cas général, $a$ peut s'annuler. Toutefois, par continuité, on peut trouver des INT maximaux sur lesquels $a$ ne s'annule pas. Sur ces intervalles, on divise par $a$ pour se ramener sous forme résolue et appliquer notre méthode sur ces sous-intervalles. Ensuite, lorsqu'on a résolu sur tous les sous-intervalles, on essaye de recoller nos solutions aux points d'annulation de $a$ en utilisant les hypothèses de continuité ou de dérivabilité de $y$ en ces points-là : c'est ce qu'on appelle le \textbf{recollement}.
\end{dfn}
\section{Équations différentielles linéaires d'ordre 2 à coefficients constants}
Le seul cas au programme pour les EDL d'ordre 2 est celui avec les coefficients constants. En particulier, le premier coefficient est donc non nul, et on peut se ramener à la forme générale :$$y''+ay'+by=f(t)$$ avec $f$ continue. On remarquera que toute solution est nécessairement de classe $\mathcal{C}^2$.
\begin{thm}[Équation homogène, cas complexe]
	Soit $(E_0)$ : $y''+ay'+b=0$ avec $(a,b)\in\C^2$. Voici comment trouver $S_0$. On considère le polynôme caractéristique $P=X^2+aX+b$ puis on distingue les cas suivants :
	\begin{enumerate}
		\item Supposons que $P$ possède deux racines complexes \textbf{distinctes} $r_1$ et $r_2$. Alors $$S_0=\left\{t\mapsto\lambda\exp(r_1t)+\mu\exp(r_2t) \ | \ (\lambda,\mu)\in\C^2\right\}$$
		\item Supposons que $P$ possède une racine \textbf{double} $r$. Alors $$S_0=\left\{t\mapsto(\lambda+\mu t)\exp(rt) \ | \ (\lambda,\mu)\in\C^2\right\}$$
	\end{enumerate}
\end{thm}
\begin{thm}[Équation homogène, cas réel]
	Soit $(E_0)$ : $y''+ay'+b=0$ avec $(a,b)\in\R^2$. Voici comment trouver $S_0$. On considère le polynôme caractéristique $P=X^2+aX+b$ puis on distingue les cas suivants :
	\begin{enumerate}
		\item Supposons que $P$ possède deux racines réelles \textbf{distinctes} $r_1$ et $r_2$. Alors $$S_0=\left\{t\mapsto\lambda\exp(r_1t)+\mu\exp(r_2t) \ | \ (\lambda,\mu)\in\R^2\right\}$$
		\item Supposons que $P$ possède une racine réelle \textbf{double} $r$. Alors $$S_0=\left\{t\mapsto(\lambda+\mu t)\exp(rt) \ | \ (\lambda,\mu)\in\R^2\right\}$$
		\item Supposons que $P$ possède 2 racines complexes non réelles conjuguées $s\pm i\omega$. Alors $$S_0=\left\{t\mapsto \exp(st)(\lambda\cos(\omega t)+\mu\sin(\omega t) \ | \ (\lambda,\mu)\in\R^2)\right\}$$
	\end{enumerate}
\end{thm}
\begin{cor}
	Que ce soit dans $\R$ ou dans $\C$, $S_0$ est un $\K$-espace vectoriel de dimension 2.
\end{cor}
\begin{thm}[Existence et unicité d'une solution à un problème de Cauchy]
	Soit $(C_0,C_1)\in\K^2$ et $t_0\in I$. Alors il existe une unique solution de $y''+ay'+by=f(t)$ qui vérifie $y(t_0)=C_0$ et $y'(t_0)=C_1$
\end{thm}
Concrètement, les seuls cas au programme pour $f$ sont les fonctions polynomiales et les exponentielles, ainsi que tout ce qui s'en déduit par superposition ou par passage aux parties réelle et imaginaire.
\begin{thm}[Cas polynomial]
	On garde l'équation $y''+ay'+by=f(t)$ et on suppose que $f(t)$ est de la forme $Q(t)$ où $Q\in\K[X]$ est de degré $n_in\N$. On cherche une solution polynomiale. On distingue alors différent cas :
	\begin{enumerate}
		\item Supposons que $b\ne 0$. Alors on cherche la solution sous la forme d'un polynôme de $\K_n[X]$ et on réinjecte dans l'équation pour obtenir un système linéaire.
		\item Supposons que $b=0$ et $a\ne 0$ On cherche la solution sous la forme d'un polynôme de $\K_{n+1}[X]$ et on réinjecte dans l'équation pour obtenir un système linéaire.
		\item Supposons que $b=0$ et $a=0$ On cherche la solution sous la forme d'un polynôme de $\K_{n+2}[X]$ et on réinjecte dans l'équation pour obtenir un système linéaire.
	\end{enumerate}
\end{thm}
\begin{proof}
	Tout cela est lié aux fonctions 
	$$
	\begin{array}[t]{lrcl}
		\varphi: & \K_p[X] & \longrightarrow & \K_n[X] \\
		& R & \longmapsto & R''+aR'+bR
	\end{array}
	$$ pour $p\in\left\{n,n+1,n+2\right\}$ dans nos différents cas (dans le même ordre). A chaque fois, on détermine le noyau par des considérations de degré. Puis par le théorème du rang ces application sont surjectives. On a alors l'existence d'une solution, puis on la cherche en réinjectant. Cela fonctionne de même avec des applications légèrement plus complexes pour des seconds membres de la forme $Q(t)\exp(rt)$.
\end{proof}
\begin{thm}
	On garde l'équation $y''+ay'+by=f(t)$ et on suppose que $f(t)$ est de la forme $\exp(rt)$ pour $r\in\K$. On note toujours $P=X^2+aX+b$ le polynôme caractéristique. Voici comment trouver une solution particulière. On distingue différents cas selon la multiplicité de $r$ en tant que racine de $P$ :
	\begin{enumerate}
		\item Supposons que $r$ n'est \textbf{pas racine} de $P$. Alors $$t\mapsto\frac{1}{P(r)}\exp(rt)$$ est une solution particulière.
		\item Supposons que $r$ est racine \textbf{simple} de $P$. Alors $$t\mapsto\frac{t}{P'(r)}\exp(rt)$$ est une solution particulière.
		\item Supposons que $r$ est racine \textbf{double} de $P$. Alors $$t\mapsto\frac{t^2}{P''(r)}\exp(rt)$$ est une solution particulière.
	\end{enumerate}
	Pour le retenir, on a de manière générique, en posant $m=m_P(r)$ la multiplicité de $r$ en tant que racine de $P$, la solution particulière $$t\mapsto \frac{t^m}{P^{(m)}(r)}\exp(rt)$$
\end{thm}
\section{Compléments, HP}
\subsection{Lien avec la science physique}
Il est très important pour l'étudiant de faire le lien entre ce chapitre et certains chapitres du programme de physique. En mécanique par exemple, l'équation du mouvement d'un mobile soumis à des forces aboutit à la résolution d'une équation différentielle. Le principe fondamental de la dynamique donne l'accélération, ce qui aboutit à une équation différentielle d'ordre 2 (avec pourquoi pas un terme à l'ordre 1 s'il y a des forces de frottement proportionnelles à la vitesse). \\ \par
Une autre application, et des plus notables, se rencontre en électricité dans l'étude des circuits RLC, RL ou RC. Prenons le circuit RLC par exemple, qui consiste à placer en série
\begin{itemize}
	\item un conducteur ohmique de résistance R ;
	\item une bobine d'inductance L; 
	\item un condensateur de capacité C.
\end{itemize}
Comme inconnue, nous choisissons $u$, la tension aux bornes du condensateur. La tension aux bornes du conducteur ohmique vaut alors $Ri = RCu'$, et celle aux bornes de la bobine vaut $Li' = LCu''$. On a ensuite $$u + RCu' + LCu'' = 0 \iff u'' + \frac{R}{L}u' + \frac{1}{LC}u = 0$$ Le polynôme caractéristique est $\displaystyle X^2 + \frac{R}{L}X + \frac{1}{LC}$, de discriminant $\displaystyle\frac{R^2}{L^2} - \frac{4}{LC}$
\begin{itemize}
	\item Un cas intéressant est celui où $R = 0$. On a alors deux racines complexes conjuguées $\displaystyle\pm \frac{i}{\sqrt{LC}}$. On posera $\displaystyle\omega = \frac{1}{\sqrt{LC}}$ si bien que les solutions de l'équation sont de la forme $$y(t) = \lambda \cos(\omega t) + \mu \sin(\omega t)$$
	On a ce qu'on appelle un \textbf{oscillateur harmonique} à la \textbf{pulsation} $\omega$.
	\item Plus généralement, si $\displaystyle R < 2\sqrt{\frac{L}{C}}$, alors en posant $$\omega := \frac{\sqrt{4L/C-R^2}}{2L}$$ on aura toujours un oscillateur harmonique à la pulsation $\omega$, mais il sera \textbf{amorti} au cours du temps : $$y(t) = (\lambda \cos(\omega t) + \mu \sin(\omega t))e^{-\frac{R}{2L}t}$$
	\item Si $\displaystyle R > 2\sqrt{\frac{L}{C}}$, notons $r_1$ et $r_2$ les deux racines du polynôme caractéristique ; les solutions sont de la forme $$y(t) = \lambda e^{r_1 t} + \mu e^{r_2 t}$$ Or, comme le produit des racines vaut $\displaystyle \frac{1}{LC} > 0$ et leur somme vaut $\displaystyle -\frac{R}{L} < 0$, on en déduit qu'elles sont toutes les deux strictement négatives. Donc $\lambda$ et $\mu$ ne peuvent être nuls sans risque d'obtenir des tensions infinies quand $t$ grandit. Quoi qu'il arrive, il n'y a plus d'oscillations du tout, seulement un amortissement qui tend vers une tension nulle.
	\item Enfin, si $R=\displaystyle2\sqrt{\frac{L}{C}}$, on est en "\textbf{régime critique}", à la limite des oscillations.
\end{itemize}
Si maintenant un générateur dans le circuit à la pulsation $\omega_2$, la solution générale de l'équation différentielle est la somme de la solution précédente et d'une solution particulière à la pulsation $\omega_2$. En régime permanent avec $R>0$ lorsque la composante homogène a été amortie, il ne reste plus que la partie à la pulsation $\omega_2$. On est sorti du "\textbf{régime transitoire}" pour entrer dans le "\textbf{régime forcé}". En passant en notations complexes (voir cours de physique), on démontre que l'amplitude du signal est maximale lorsque $\omega_2=\omega$ ; le circuit entre alors en \textbf{résonance}.

\subsection{Équations d'ordre 2}
On se propose ici de prouver un théorème admis dont la démonstration est instructive.
\begin{thm}
	Soit l'équation $y''+ay'+by=f(t)$ avec $a$ et $b$ constants et $f$ continue sur $I$.
Si on se donne des conditions $y(t_0)=C_0$ et $y'(t_0)=C_1$, alors l'équation admet une unique solution.
\end{thm}
\begin{proof}
On va adapter la technique de variation de la constante à l'ordre 2. Il y aura maintenant deux "constantes" à faire varier. La démonstration est assez longue, elle fonctionne en plusieurs points.
\begin{itemize}
	\item On rappelle que les solutions de l'équation sans second membre forment un $\K$-ev de dimension 2. Pour en trouver une base, on considère les racines du polynôme caractéristique $X^2+aX+b$.
	\begin{itemize}
		\item[--] Deux racines distinctes: une base est donnée par le couple $(t\mapsto e^{r_1t},t\mapsto e^{r_2t})$ ou par le couple $(t\mapsto \cos(\omega t)e^{st},t\mapsto \sin(\omega t)e^{st})$, selon qu'on est dans le cas complexe ou réel.
		\item[--] Une racine double: une base est donnée par $\langle t\mapsto e^{rt},t\mapsto te^{rt}\rangle$.
	\end{itemize}
	Dans tous les cas, nous noterons $(y_1,y_2)$ notre base. La solution générale de l'équation homogène s'écrit donc $\lambda_1y_1+\lambda_2y_2$.
	\item Ensuite, l'idée est d'associer à toute fonction $y$ deux fois dérivable des fonctions $\lambda_1$ et $\lambda_2$ telles que $$\forall t\in I,\ \begin{cases}y(t)=\lambda_1(t)y_1(t)+\lambda_2(t)y_2(t)\\ y'(t)=\lambda_1(t)y_1'(t)+\lambda_2(t)y_2'(t)\end{cases}$$
	Mais en avons-nous le droit? Pour bien mettre en évidence les inconnues, écrivons le
	système sous la forme $$\begin{cases}y_1\lambda_1+y_2\lambda_2=y\\ y_1'\lambda_1+y_2'\lambda_2=y'\end{cases}$$
	soit encore $\displaystyle A(t)\begin{pmatrix}\lambda_1\\ \lambda_2\end{pmatrix}=\begin{pmatrix}y\\ y'\end{pmatrix}$ avec $\displaystyle A(t):=\begin{pmatrix}y_1(t)&y_2(t)\\ y_1'(t)&y_2'(t)\end{pmatrix}$. \\
	Nous allons commencer par introduire la fonction $W:=y_1y_2'-y_1'y_2$ (c'est-à-dire le déterminant de notre système). La fonction $W$ est dérivable, et on a :
	\begin{align*}
		W'+aW&=y_1'y_2'+y_1y_2''-y_1''y_2-y_1'y_2'+ay_1y_2'-ay_1'y_2\\
		&=y_1(y_2''+ay_2')-y_2(y_1''+ay_1')\\
		&=y_1(-by_2)-y_2(-by_1)\\
		&=0
	\end{align*}
	La fonction $W$ joue un rôle très important dans l'étude des équations différentielles, on l'appelle le \textbf{wronskien}. Ici on a $\forall t \in I,\ W(t) = \mu e^{-at}$ avec $\mu \in \mathbb{K}$ fixé. Montrons à
	présent que $\mu = W(0) \neq 0$.
	\begin{itemize}
		\item[--]Si $y_1(t) = e^{r_1 t}$ et $y_2(t) = e^{r_2 t}$ on a $W(0) = r_2 - r_1 \neq 0$.
		\item[--]Si $y_1(t) = e^{rt}$ et $y_2(t) = te^{rt}$ on a $W(0) = 1 \neq 0$.
		\item[--]Si $y_1(t) = \cos(\omega t)e^{st}$ et $y_2(t) = \sin(\omega t)e^{st}$ on a $W(0) = \omega \neq 0$ (car les racines du polynôme caractéristique sont non réelles).
	\end{itemize}
	On en déduit que $$\forall t \in I, W(t) \neq 0$$ On vérifie ensuite que $$\forall t \in I, A(t) \begin{pmatrix} y_1 & -y_1' \\ -y_2 & y_2' \end{pmatrix} = W(t)I_2,$$
	Donc pour tout $t \in I,$ la matrice $A(t)$ est inversible et $$A(t)^{-1} = \frac{1}{W(t)} \begin{pmatrix} y_1(t) & -y_1'(t) \\ -y_2(t) & y_2'(t) \end{pmatrix}$$
	
	Si on note $\alpha(t), \beta(t), \gamma(t), \delta(t)$ les coefficients de $A(t)^{-1},$ on voit que $\alpha, \beta, \gamma, \delta$ sont des
	fonctions dérivables (comme quotients dont le dénominateur $W(t)$ ne s'annule pas). On
	peut alors poser, pour toute fonction $y$ deux fois dérivable,
	$$\begin{pmatrix} \lambda_1 \\ \lambda_2 \end{pmatrix} = A^{-1} \begin{pmatrix} y \\ y' \end{pmatrix}$$
	\item 
	$\lambda_1$ et $\lambda_2$ s'expriment comme sommes et produits à partir de $\alpha, \beta, \gamma, \delta, y, y'$, donc elles sont dérivables. Puis on a $$y' = (y_1 \lambda_1 + y_2 \lambda_2)' = y' + y_1 \lambda_1' + y_2 \lambda_2'$$ d'où $y_1 \lambda_1' + y_2 \lambda_2' = 0$ (notons que cette égalité est toujours satisfaite, que $y$ soit solution ou non). On a aussi :
	\begin{align*}
		y'' + ay' + by &= (y_1'' \lambda_1 + y_1' \lambda_1' + y_2'' \lambda_2 + y_2' \lambda_2') + a(y_1' \lambda_1 + y_2' \lambda_2) + b(y_1 \lambda_1 + y_2 \lambda_2)\\
		&= y_1' \lambda_1' + y_2' \lambda_2'
	\end{align*}
	Finalement, $y$ est solution ssi $\displaystyle A \begin{pmatrix} \lambda_1' \\ \lambda_2' \end{pmatrix} = \begin{pmatrix} 0 \\ f \end{pmatrix}$ ssi $\displaystyle\begin{pmatrix} \lambda_1' \\ \lambda_2' \end{pmatrix} = A^{-1} \binom{0}{f}$
	\item Or $\displaystyle A(t)^{-1} \begin{pmatrix} 0 \\ f(t) \end{pmatrix}$ est un vecteur-colonne dont les coordonnées s'expriment comme sommes et produits à partir de $\alpha, \beta, \gamma, \delta$ et $f$. On peut donc écrire $$\forall t\in I,\ A(t)^{-1} \begin{pmatrix} 0 \\ f(t) \end{pmatrix} = \begin{pmatrix} \mu_1(t) \\ \mu_2(t) \end{pmatrix}$$ où $\mu_1$ et $\mu_2$ sont continues. Finalement, $y$ est solution ssi $$\forall t \in I,\ \begin{cases} \lambda_1'(t) = \mu_1(t) \\ \lambda_2'(t) = \mu_2(t) \end{cases}$$
	Puisque toute fonction continue admet une primitive, on en déduit que l'équation possède des solutions.
	\item Si de plus on fixe une condition de Cauchy, on a $$\begin{cases}y(t_{0})=C_{0} \\ y'(t_{0})=C_{1} \end{cases} \iff A(t_{0}) \begin{pmatrix} \lambda_{1}(t_{0}) \\ \lambda_{2}(t_{0}) \end{pmatrix} = \begin{pmatrix} C_{0} \\ C_{1} \end{pmatrix}\iff  \begin{pmatrix} \lambda_{1}(t_{0}) \\ \lambda_{2}(t_{0}) \end{pmatrix} = A(t_{0})^{-1} \begin{pmatrix} C_{0} \\ C_{1} \end{pmatrix}$$
	donc $\lambda_{1}$ et $\lambda_{2}$ sont fixées de manière unique.
\end{itemize}
Ainsi la preuve est achevée.
\end{proof}

\newpage
\chapter{Fonctions de deux variables}
Dans tout le chapitre, on munit $\R^2$ de sa norme euclidienne canonique, notée $||\cdot||$.
\section{Fonctions continues sur un ouvert de $\R^2$}
\begin{dfn}[Disque ouvert]
	Soit $p\in\R^2$ et $r>0$ On appelle \textbf{disque ouvert} de centre $p$ et de rayon $r$ l'ensemble $$D(p,r)=\left\{z\in\R^2 \tq ||z-p||<r\right\}$$
\end{dfn}
\begin{dfn}[Ouvert de $\R^2$]
	Soit $U$ une partie de $\R^2$. On dit que $U$ est un \textbf{ouvert} de $\R^2$ lorsque $$\forall p \in U, \ \exists r>0, \ D(p,r) \subset U$$
\end{dfn}
Dans la suite, $U$ désigne toujours un ouvert non vide de $\R^2$.
\begin{dfn}[Limite, continuité locale et continuité globale]
	Soit $f:U\rightarrow\R$ et $p\in U$. On donne les définitions suivantes : \begin{enumerate}
		\item Soit $l\in\R$. On dit que $f$ \textbf{tend vers $l$ en $p$}, et on écrit $f(z) \xrightarrow[z\to p]{}l$, lorsque $$\forall \varepsilon>0, \ \exists\delta>0, \ \forall z \in U, \ ||z-p||\leqslant \delta \Longrightarrow|f(z)-l|\leqslant\varepsilon$$
		\item On dit que la fonction $f$ \textbf{est continue en $p$} lorsque $f(z) \xrightarrow[z\to p]{}f(p)$, c'est-à-dire lorsque $$\forall \varepsilon>0, \ \exists\delta>0, \ \forall z \in U, \ ||z-p||\leqslant \delta \Longrightarrow|f(z)-f(p)|\leqslant\varepsilon$$
		\item On dit que $f$ est \textbf{continue} lorsque $f$ est continue en tout point de $U$.
	\end{enumerate}
\end{dfn}
\begin{prop}[Caractérisation séquentielle de la limite, HP]
	On a $$f(x,y)\xrightarrow[(x,y)\to p=(a,b)]{}l$$ si, et seulement si, $$\forall(x_n,y_n)_{n\in\N}\in U^\N, \ \left[x_n\xrightarrow[n\to+\infty]{}a \ \text{et} \ y_n\xrightarrow[n\to+\infty]{}b \right] \Longrightarrow f(x_n,y_n)\xrightarrow[n\to+\infty]{}l$$
\end{prop}
\begin{cor}[Caractérisation séquentielle de la continuité, HP]
	$f$ est continue en $p=(a,b)$ si, et seulement si, $$\forall(x_n,y_n)_{n\in\N}\in U^\N, \ \left[x_n\xrightarrow[n\to+\infty]{}a \ \text{et} \ y_n\xrightarrow[n\to+\infty]{}b \right] \Longrightarrow f(x_n,y_n)\xrightarrow[n\to+\infty]{}f(a,b)$$
\end{cor}
\begin{cor}[Opérations algébriques sur la continuité des fonctions de deux variables]
	La continuité des fonctions de deux variables passe aux combinaisons linéaires, au produit et au quotient si le dénominateur ne s'annule pas.
\end{cor}
\begin{rmq} Soit $\theta\in\mathcal{C}^0(\R,\R)$. Les fonctions de $\R^2$ dans $\R$ $(x,y)\mapsto\theta(x)$ et $(x,y)\mapsto\theta(y)$ sont continues. En particulier, cela est vrai avec l'identité, donc les projections $(x,y)\mapsto x$ et $(x,y)\mapsto y$ sont continues, ce qui permet de montrer la continuité de nombreuses fonctions par opérations algébriques. \end{rmq}
\begin{prop}
	Soit $f:U\rightarrow\R$, $I$ un intervalle contenant $f(U)$ et $\theta:I\rightarrow\R$. Si $f$ est continue en $p\in U$, et si $\theta$ est continue en $f(p)$, alors la composée $\theta\circ f$ est continue en $p$.
\end{prop}
\begin{prop}
	Soit $f:U\rightarrow\R$, $I$ un intervalle ainsi que $\gamma_1$ et $\gamma_2$ deux fonctions de $I$ dans $\R$ telles que $$\forall t \in I, \ (\gamma_1(t),\gamma_2(t))\in U$$ Si $\gamma_1$ et $\gamma_2$ sont continues en $a\in I$, et si $f$ est continue en $(\gamma_1(a),\gamma_2(a))$, alors la composée $t\mapsto f(\gamma_1(t),\gamma_2(t))$ est continue en $a$.
\end{prop}
\begin{prop}
	Soit $f:U\rightarrow\R$, $V$ un ouvert de $\R^2$ et $\varphi,\psi : V\rightarrow \R$ deux fonctions telles que $$\forall z \in V, \ (\varphi(z),\psi(z))\in U$$ Si $\varphi$ et $\psi$ sont continues en $p\in V$, et si $f$ est continue en $(\varphi(p),\psi(p))$, alors la composée $z\mapsto f(\varphi(z),\psi(z))$ est continue en $p$.
\end{prop}

\section{Dérivation}
\begin{dfn}[Applications partielles]
	Étant donné un point $p=(a,b)\in U$, on considère les ensembles $$D_1=\left\{x\in\R\tq (x,b)\in U\right\} \ \text{et} \ D_2=\left\{y\in\R\tq(a,y)\in U\right\}$$ ainsi que les applications partielles $$
	\begin{array}[t]{lrcl}
		\varphi_1: & D_1 & \longrightarrow & \R \\
		& x & \longmapsto & \displaystyle f(x,b)
	\end{array}
	$$ et 
	$$
	\begin{array}[t]{lrcl}
		\varphi_2: & D_2 & \longrightarrow & \R \\
		& y & \longmapsto & \displaystyle f(a,y)
	\end{array}
	$$
\end{dfn}
\begin{dfn}[Dérivées partielles]
	Soit $p=(a,b)\in U$. \begin{enumerate}
		\item Si l'application partielle $\varphi_1$ est dérivable en $a$, on dit que $f$ admet une \textbf{première dérivée partielle} en $(a,b)$ et l'on pose $$\partial_1f(a,b)=\varphi_1'(a)$$
		\item Si l'application partielle $\varphi_2$ est dérivable en $a$, on dit que $f$ admet une \textbf{deuxième dérivée partielle} en $(a,b)$ et l'on pose $$\partial_2f(a,b)=\varphi_2'(b)$$
	\end{enumerate}
\end{dfn}
\begin{rmq} En pratique, pour une fonction $f:(x,y)\mapsto f(x,y)$, on utilise les notations plus parlantes $$\frac{\partial f}{\partial x}(a,b)=\partial_1f(a,b)$$ et $$\frac{\partial f}{\partial y}(a,b)=\partial_2f(a,b)$$ en s'adaptant au nom des variables utilisées dans la définition de $f$. \end{rmq}
\begin{prop}[Combinaison linéaire et produit]
	Soit $f$ et $g$ deux fonctions de $U$ dans $\R$, $\lambda$ et $\mu$ des réels et $p\in U$. Supposons que $f$ et $g$ admettent des dérivées partielles en $p$.
	\begin{enumerate}
		\item Alors $\lambda f+\mu g$ également, avec $$\frac{\partial(\lambda f+\mu g)}{\partial x}(p)=\lambda\frac{\partial f}{\partial x}(p)+\mu\frac{\partial g}{\partial x}(p)$$ et $$\frac{\partial(\lambda f+\mu g)}{\partial y}(p)=\lambda\frac{\partial f}{\partial y}(p)+\mu\frac{\partial g}{\partial y}(p)$$
		\item Et $fg$ également, avec $$\frac{\partial(fg)}{\partial x}(p)=\frac{\partial f}{\partial x}(p)g(p)+f(p)\frac{\partial g}{\partial x}(p)$$ et $$\frac{\partial(fg)}{\partial y}(p)=\frac{\partial f}{\partial y}(p)g(p)+f(p)\frac{\partial g}{\partial y}(p)$$
	\end{enumerate}
\end{prop}
\begin{rmq} \textbf{Attention}, l'existence de dérivées partielles, même en tout point de $U$, n'entraîne pas la continuité de $f$. Comme contre-exemple, on peut citer $\displaystyle (x,y)\mapsto\frac{xy}{x^2+y^2}$ prolongée par 0 en $(0,0)$ qui admet des dérivées partielles en tout point mais qui n'est pourtant pas continue en $(0,0)$. \end{rmq}
\begin{dfn}[Fonctions de classe $\mathcal{C}^1$]
	La fonction $f$ est dite \textbf{de classe $\mathcal{C}^1$} lorsqu'elle admet admet des dérivées partielles en tout point de $U$ et que les fonctions $\displaystyle \frac{\partial f}{\partial x}$ et $\displaystyle \frac{\partial f}{\partial y}$ définies alors sur $U$ sont continues. \\ On note $\mathcal{C}^1(U,\R)$ l'ensemble des fonctions de classe $\mathcal{C}^1$ définies sur $U$.
\end{dfn}
\begin{thm}[Développement limité à l'ordre 1]
	Soit $f\in\mathcal{C}^1(U,\R)$ et $p=(a,b)\in U$. Il existe une fonction $\varepsilon:U\rightarrow\R$ telle que $\varepsilon(z)\xrightarrow[z\to p]{}0$ et $$\forall(x,y)\in U, \ f(x,y)=f(a,b)+(x-a)\frac{\partial f}{\partial x}(a,b)+(y-b)\frac{\partial f}{\partial y}(a,b)+\varepsilon(x,y)||(x,y)-(a,b)||$$
\end{thm}
\begin{dfn}[Plan tangent]
	Le théorème précédent affirme que la fonction affine $$(x,y)\mapsto f(a,b)+(x-a)\frac{\partial f}{\partial x}(a,b)+(y-b)\frac{\partial f}{\partial y}(a,b)$$ approche au premier ordre $f$ au voisinage de $(a,b)$. On appelle alors \textbf{plan tangent du graphe de la fonction $f$ en $(a,b)$} le plan d'équation $$z=f(a,b)+(x-a)\frac{\partial f}{\partial x}(a,b)+(y-b)\frac{\partial f}{\partial y}(a,b)$$
\end{dfn}
\begin{cor}[$\mathcal{C}^1 \Longrightarrow \mathcal{C}^0$]
	Si $f\in\mathcal{C}^1(U,\R)$, alors $f$ est continue.
\end{cor}
\begin{cor}
	Soit $(f,g)\in(\mathcal{C}^1(U,\R))^2$ et $\lambda,\mu\in\R$. Alors $\lambda f +\mu g$ et $fg$ sont des fonctions de classe $\mathcal{C}^1$.
\end{cor}
\begin{dfn}[Gradient, champ de vecteurs]
	Soit $f\in\mathcal{C}^1(U,\R)$. Le \textbf{gradient} de $f$ est l'application (lire $\nabla$ "nabla")
	$$
	\begin{array}[t]{lrcl}
		\nabla f: & U & \longrightarrow & \R^2 \\
		& (a,b) & \longmapsto & \displaystyle \nabla f(a,b)=\begin{pmatrix}
			\dfrac{\partial f}{\partial x}(a,b) \\[6mm]
			\dfrac{\partial f}{\partial y}(a,b)
		\end{pmatrix}
	\end{array}
	$$ Le gradient est un \textbf{champ de vecteurs}, c'est-à-dire une application dont les valeurs sont des vecteurs.
\end{dfn}
\begin{rmq} On peut réécrire le DL à l'ordre 1 à l'aide du gradient et du produit scalaire usuel de $\R^2$ (avec les mêmes hypothèses que dans le théorème) $$\forall z \in U, \ f(z)=f(p)+\left\langle \nabla f(p),z-p\right\rangle + \varepsilon(z)||z-p||$$ \end{rmq}

\section{Dérivation des fonctions composées}
\begin{prop}
	Soit $f\in\mathcal{C}^1(U,\R)$, $I$ un intervalle non trivial contenant $f(U)$ et $\theta\in\mathcal{C}^1(I,\R)$. Alors $\theta\circ f$ est une fonction de classe $\mathcal{C}^1$ et on a $$\forall p \in U, \ \frac{\partial (\theta\circ f)}{\partial x}(p)=\theta'(f(p))\frac{\partial f}{\partial x}(p)$$ et$$\forall p \in U,\ \frac{\partial (\theta\circ f)}{\partial y}(p)=\theta'(f(p))\frac{\partial f}{\partial y}(p)$$ 
\end{prop}
\begin{thm}[Première règle de la chaîne]
	Soit $f\in\mathcal{C}^1(U,\R)$ et $(\gamma_1,\gamma_2)\in(\mathcal{C}^1(I,\R))^2$. On suppose que $$\forall t \in I, \ (\gamma_1(t),\gamma_2(t))\in U$$ Alors $$g:t\mapsto f(\gamma_1(t),\gamma_2(t))$$ est de classe $\mathcal{C}^1$ et on a $$\forall a \in I, \ g'(a)=\frac{\partial f}{\partial x}(\gamma_1(a),\gamma_2(a))\gamma_1'(a)+\frac{\partial f}{\partial y}(\gamma_1(a),\gamma_2(a))\gamma_2'(a)$$
\end{thm}
\begin{dfn}[Ligne de niveau, HP]
	On appelle \textbf{ligne de niveau} d'une fonction $f:\R^2\mapsto\R$ une partie de $\R^2$ d'équation $f(x,y)=c$ pour un certain $c\in\R$. \\ Si la fonction $\gamma$ paramètre une ligne de niveau, c'est-à-dire si la fonction $f\circ\gamma$ est constante, on obtient l'égalité $$\left\langle \nabla f(\gamma(a)),\gamma'(a)\right\rangle=0$$ c'est-à-dire que le gradient $\nabla f(\gamma(a))$ est orthogonal au vecteur dérivé $\gamma'(a)$ qui dirige la tangente au point $\gamma(a)$ de la ligne de niveau. On dit que le gradient $\nabla f$ est orthogonal aux lignes de niveau. 
\end{dfn}
\begin{dfn}[Dérivée selon un vecteur]
	Soit $f\in\mathcal{C}^1(U,\R)$, $p\in U$ et $v=\begin{pmatrix}
		h \\ k
	\end{pmatrix} \in \R^2$. L'application $$\varphi_v:t\mapsto f(p+tv)$$ est définie sur une union d'intervalle ouverts dont au moins un contient 0. De plus, elle est dérivable en $0$, de dérivée : $$\varphi_v'(0)=h\frac{\partial f}{\partial x}(p)+k\frac{\partial f}{\partial y}(p)=\left\langle \nabla f(p), v \right\rangle$$ On note $D_vf(p)=\varphi_v'(0)$ cette dérivée, et on l'appelle \textbf{dérivée de $f$ en $p$ selon le vecteur $v$}.
\end{dfn}
\begin{rmq} La notion de dérivée selon un vecteur généralise celle des dérivées partielles en ne faisant plus jouer un rôle particulier aux droites parallèles aux axes de coordonnées. \end{rmq}
\begin{rmq} Le gradient de $f$ en $p$ "pointe" dans la direction dans laquelle $f$ croît le plus vite. Cela correspond à la direction de plus grande pente sur le graphe $\Gamma_f$. \end{rmq}

\begin{thm}[Deuxième règle de la chaîne]
	Soit $f\in\mathcal{C}^1(U,\R)$. Soit $V$ un ouvert non vide de $\R^2$ et $(\varphi,\psi)\in(\mathcal{C}^1(V,\R))^2$ un couple de fonctions tel que $$\forall z \in V, \ (\varphi(z),\psi(z))\in U$$ Alors l'application composée 
	$$
	\begin{array}[t]{lrcl}
		g: & V & \longrightarrow & \R \\
		& z & \longmapsto & f(\varphi(z),\psi(z))
	\end{array}
	$$ est de classe $\mathcal{C}^1$ et vérifie $$\forall p \in V, \ \partial_1g(p)=\partial_1f(\varphi(p),\psi(p))\partial_1\varphi(p)+\partial_2f(\varphi(p),\psi(p))\partial_1\psi(p)$$ et $$\forall p\in V, \ \partial_2g(p)=\partial_1f(\varphi(p),\psi(p))\partial_2\varphi(p)+\partial_2f(\varphi(p),\psi(p))\partial_2\psi(p)$$
\end{thm}
\section{Extrema}
\begin{dfn}[Extrema]
	Soit $X$ une partie non vide de $\R^2$, $f:X\rightarrow \R$ et $p\in X$.
	\begin{enumerate}
		\item On dit que $f$ admet un \textbf{maximum} en $p$ lorsque $\forall z \in X, \ f(z)\leqslant f(p)$
		\item On dit que $f$ admet un \textbf{maximum local} lorsque $$\exists \delta>0, \ \forall z\in X\cap D(p,\delta), \ f(z)\leqslant f(p)$$
		\item On dit que $f$ admet un \textbf{minimum} en $p$ lorsque $\forall z \in X, \ f(z)\geqslant f(p)$
		\item On dit que $f$ admet un \textbf{minimum local} lorsque $$\exists \delta>0, \ \forall z\in X\cap D(p,\delta), \ f(z)\geqslant f(p)$$
		\item On dit que $f$ un \textbf{extremum} en $p$ en si $f$ admet un maximum ou un minimum en $p$.
		\item On dit que $f$ un \textbf{extremum local} en $p$ en si $f$ admet un maximum local ou un minimum local en $p$.
	\end{enumerate}
\end{dfn}
\begin{rmq} Pour insister sur la distinction avec les extrema locaux, les extrema sont parfois appelés \textbf{extrema globaux}. \end{rmq}
\begin{dfn}[Point critique]
	Soit $f\in\mathcal{C}^1(U,\R)$ et $p\in U$. On dit que $p$ est un \textbf{point critique} pour $f$ lorsque $$\frac{\partial f}{\partial x}(p)=\frac{\partial f}{\partial y}(p)=0$$ c'est-à-dire lorsque $\displaystyle \nabla f(p)=\begin{pmatrix} 0 \\ 0	\end{pmatrix}$.
\end{dfn}
\begin{thm}[Condition nécessaire d'extremum local]
	Soit $f\in\mathcal{C}^1(U,\R)$ et $p\in U$. Si $p$ admet un extremum local en $p$, alors $p$ est un point critique de $f$.
\end{thm}

\newpage
\chapter{Séries}
Dans tout le chapitre, les suites sont considérées à valeurs dans $\K=\R$ ou $\C$.
\section{Généralités}
\begin{dfn}[Série, somme partielle]
	Soit $(u_n)_{n\in\N}$ une suite à termes dans $\K$. La \textbf{série de terme général $u_n$} désigne en fait la suite $\displaystyle \left(\sum_{k=0}^{n}u_k\right)_{n\in\N}$. On la note $\displaystyle \sum u_n$. Le nombre $\displaystyle\sum_{k=0}^{n}u_k$ s'appelle la \textbf{somme partielle de rang $n$} de la série $\displaystyle \sum u_n$. Dans ce chapitre, nous utiliserons la notation très répandue $S_n$ pour la somme partielle de rang $n$.
\end{dfn}
\begin{dfn}[Série convergente, somme de la série]
	La série $\displaystyle \sum u_n$ est dite \textbf{convergente} lorsque la suite $(S_n)_{n\in\N}$, et la limite des sommes partielles s'appelle alors la \textbf{somme de la série}. On la note $$\sum_{n=0}^{+\infty}u_n$$
\end{dfn}
\begin{dfn}
	Plus généralement, si $n_0\in\N$ et si $(u_n)_{n\geqslant n_0}$ est une suite à termes réels ou complexes, la série $\displaystyle \sum u_n$ désigne la suite des sommes partielles $$\left(\sum_{k=0}^{n}u_k\right)_{n\geqslant n_0}$$ S'il y a ambiguïté, on pourra la noter $\displaystyle \sum_{n\geqslant n_0}u_n$. Si elle converge, sa somme est notée $$\sum_{n=n_0}^{+\infty}u_n$$
\end{dfn}
Les résultats qui suivent sont énoncés dans le cas où $n_0=0$, mais leurs énoncés symétriques dans le cas général sont également vrais.
\begin{prop}[Troncature]
	Soit $\displaystyle \sum u_n$ une série et $p\in\N\st$. Alors la série $\displaystyle \sum u_n$ converge si, et seulement si, la série tronquée $\displaystyle \sum_{n\geqslant p} u_n$ converge, et dans ce cas, on a : $$\sum_{n=0}^{+\infty}u_n=\sum_{n=0}^{p-1}u_n+\sum_{n=p}^{+\infty}u_n$$
\end{prop}
\begin{prop}[Décalage d'indice]
	Soit $p\in\N\st$ et $\displaystyle \sum_{n\geqslant p}u_n$ une série. Alors $\displaystyle \sum_{n\geqslant p}u_n$ converge si, et seulement si, $\displaystyle \sum_{n\geqslant 0}u_{n+p}$ converge, et dans ce cas on a $$\sum_{n=0}^{+\infty}u_{n+p}=\sum_{n=p}^{+\infty}u_n$$
\end{prop}
\begin{dfn}[Reste d'ordre $n$ d'une série convergente]
	Soit $\displaystyle \sum u_n$ une série \textbf{convergente} et $n\in\N$. Le \textbf{reste d'ordre $n$} de la série $\displaystyle \sum u_n$ est défini par $$R_n=\sum_{k=n+1}^{+\infty}u_k=\sum_{k>n}^{+\infty}u_k$$ Il vérifie $$S_n+R_n=\sum_{k=0}^{+\infty}u_k$$
\end{dfn}
\begin{cor}[Condition nécessaire de convergence sur les restes]
	\textbf{Si} $\displaystyle \sum u_n$ converge, \textbf{alors} $\displaystyle R_n \xrightarrow[n\to+\infty]{}0$.
\end{cor}
\begin{cor}[Condition nécessaire de convergence sur le terme général]
	\textbf{Si} $\displaystyle \sum u_n$ converge, \textbf{alors} $\displaystyle u_n \xrightarrow[n\to+\infty]{}0$.
\end{cor}
\begin{rmq} Attention, la réciproque est fausse ! Comme contre exemple, on pourra considérer la série harmonique. \end{rmq}
\begin{dfn}[Divergence]
	Une série est dite \textbf{divergente} lorsque la suite de ses sommes partielles diverge.
\end{dfn}
\begin{dfn}[Divergence grossière]
	La série $\displaystyle \sum u_n$ est dite \textbf{grossièrement divergente} lorsque $u_n$ ne tend pas vers $0$ lorsque $n$ tend vers $+\infty$.
\end{dfn}
\begin{thm}[Séries géométriques]
	Soit $z\in\C$. La série $\displaystyle \sum z^n$ converge si, et seulement si, $\lvert z\rvert<1$, et dans ce cas se somme vaut $\displaystyle \frac{1}{1-z}$
\end{thm}
\begin{thm}[Linéarité]
	Soit $\displaystyle \sum u_n$ et $\displaystyle \sum v_n$ des séries convergentes et $\lambda$ et $\mu$ des scalaires. Alors la série $\displaystyle \sum (\lambda u_n+\mu v_n)$ est convergente et sa somme vaut $$\sum_{n=0}^{+\infty}(\lambda u_n +\mu v_n)=\lambda \sum_{n=0}^{+\infty}u_n+\mu\sum_{n=0}^{+\infty}v_n$$
\end{thm}
\begin{thm}[Théorème passerelle pour les séries complexes]
	On suppose que $\displaystyle \sum u_n$ est à termes complexes. Alors $\displaystyle \sum u_n$ converge si, et seulement si, les séries réelles $\displaystyle \sum \text{Re}(u_n)$ et $\displaystyle \sum \text{Im}(u_n)$ convergent, et dans ce cas on a : $$\sum_{n=0}^{+\infty}u_n= \sum_{n=0}^{+\infty}\text{Re}(u_n)+\text{i}\sum_{n=0}^{+\infty}\text{Im}(u_n)$$
\end{thm}
\begin{prop}[Lien suite-série]
	La suite $(u_n)_{n\in\N}$ converge si, et seulement si, la série $\displaystyle(u_{n+1}-u_n)$ converge.
\end{prop}
\begin{rmq} Cette proposition évidente peut paraître inutile, mais elle permet d'étudier la nature d'une suite grâce à des outils sur les séries et peut donc se révéler très précieuse dans certains cas. \end{rmq}

\section{Séries réelles positives}
\begin{prop}
	Soit $\displaystyle \sum u_n$ une série à termes réels positifs. Alors elle est convergente si, et seulement si, la suite $(S_n)_{n\in\N}$ est majorée. Dans ce cas, on a : $$\forall n\in\N,\ S_n\leqslant S$$ où $S$ désigne la somme de la série $\displaystyle \sum u_n$.
\end{prop}
\begin{thm}
	Soit $\displaystyle \sum u_n$ et $\displaystyle \sum v_n$ des séries à termes réels positifs. Si $u_n\leqslant v_n$ pour tout $n\in\N$ et si $\displaystyle \sum v_n$, alors $\displaystyle \sum u_n$ converge, et dans ce cas on a : $$\sum_{n=0}^{+\infty}u_n\leqslant\sum_{n=0}^{+\infty}v_n$$
\end{thm}
\begin{rmq} Si l'inégalité $u_n\leqslant v_n$ n'a lieu qu'à partir d'un certain rang, alors $\displaystyle \sum u_n$ continue à converger, mais l'inégalité sur les sommes n'a plus aucune raison d'être valable. \end{rmq}
\begin{thm}
	Soit $\displaystyle \sum u_n$ et $\displaystyle \sum v_n$ des séries à termes réels positifs. Si $\displaystyle u_n \underset{n\to+\infty}{\sim}v_n$, alors $\displaystyle \sum u_n$ et $\displaystyle \sum v_n$ ont même nature.
\end{thm}
\begin{thm}[Critère spécial des séries alternées]
	Soit $(u_n)_{n\in\N}$ une suite réelle positive décroissante qui tend vers $0$. Alors 
	\begin{enumerate}
		\item D'une part, la série $\displaystyle \sum(-1)^nu_n$ converge
		\item (Contrôle de la somme et des restes) D'autre part, pour tout $n\in\N$, $\displaystyle \sum_{k=n}^{+\infty}(-1)^ku_k$ est du signe de $(-1)^n$ au sens large, et on a $$\left\lvert \sum_{k=n}^{+\infty}(-1)^ku_k\right\rvert \leqslant u_n$$
	\end{enumerate}
\end{thm}
\begin{rmq} On a le résultat symétrique avec la série $\displaystyle \sum (-1)^{n-1}u_n$. \end{rmq}
\section{Comparaison série-intégrale}
Donnons-nous $f:\R_+\rightarrow\R$ continue par morceaux. Lorsque l'on considère la série $\displaystyle \sum_{n\geqslant 1}f(n)$, la somme partielle de rang $n$, $\displaystyle \sum_{k=1}^{n}f(n)$, peut se réécrire un peu artificiellement $\displaystyle \sum_{k=1}^{n}\int_{k}^{k+1}f(k)dt$. Ici, on s'intéresse au cas où $f$ est monotone. A chaque fois, il est recommandé de faire un dessin avec les rectangles. \\ \par Par exemple, supposons que $f$ soit décroissante. On peut justifier à l'oral avec un dessin bien fait, mais formellement il faudrait écrire $$\forall k\in\N\st,\ \forall t\in[k-1,k],\ f(t)\geqslant f(k) \ \ \text{donc} \ \ f(k)\leqslant\int_{k-1}^{k}f(t)dt$$ et $$\forall k\in\N\st,\ \forall t\in[k,k+1],\ f(k)\geqslant f(t) \ \ \text{donc} \ \ f(k)\geqslant\int_{k}^{k+1}f(t)dt$$ d'où, en sommant et en utilisant la relation de Chasles, on obtient l'encadrement $$\int_{1}^{n+1}f(t)dt\leqslant S_n\leqslant\int_{0}^{n}f(t)dt$$ On obtient évidemment un résultat analogue dans le cas où $f$ est croissante. On peut adapter la méthode selon l'intervalle de définition de $f$.

\begin{thm}[Séries de Riemann]
	Soit $\alpha\in\R$. La série $\displaystyle \sum \frac{1}{n^\alpha}$ converge si, et seulement si $\alpha>1$.
\end{thm}
\begin{exa}
	La série harmonique $\displaystyle \sum_{n\geqslant 1} \frac{1}{n}$ diverge car elle correspond au cas où $\alpha=1$.
\end{exa}

\section{Convergence absolue}
\begin{dfn}[Convergence absolue]
	Soit $\displaystyle \sum u_n$ une série à termes complexes. On dit que la série $\displaystyle \sum u_n$ est absolument convergente lorsque la série à termes positifs $\displaystyle \sum \lvert u_n\rvert$ converge.
\end{dfn}
\begin{thm}[Convergence absolue $\Longrightarrow$ convergence]
	Soit $\displaystyle \sum u_n$ une série à termes complexes. \textbf{Si} $\displaystyle \sum \lvert u_n\rvert$ converge (ie si $\displaystyle \sum u_n$ converge absolument), \textbf{alors} $\displaystyle \sum u_n$ converge.
\end{thm}
\begin{dfn}[Prolongement de l'exponentielle à $\C$]
	Pour tout $z\in\C$, la série $\displaystyle \sum \frac{z^n}{n!}$ est absolument convergente. On \textbf{définit} alors $$\exp(z)=\sum_{n=0}^{+\infty}\frac{z^n}{n!}$$
\end{dfn}
\begin{cor}[Règle du grand O]
	Soit $\displaystyle \sum u_n$ une série à termes complexes et $\displaystyle \sum v_n$ une série à termes réels positifs convergente. Si $\displaystyle u_n\underset{n\to+\infty}{=}O(v_n)$, alors la série $\displaystyle \sum u_n$ converge absolument, donc converge.
\end{cor}
\begin{thm}[Règle de D'Alembert, HP]
	Soit $\displaystyle \sum u_n$ une série à termes complexes non nuls. Supposons que $$\left\lvert\frac{u_{n+1}}{u_n}\right\rvert\xrightarrow[n\to+\infty]{}\lambda\in[0,+\infty]$$ Si $\lambda<1$, alors la série converge absolument. Si $\lambda>1$, alors elle diverge grossièrement. Si $\lambda=1$, on ne peut rien dire.
\end{thm}
\begin{proof}
	Dans le premier cas, on prend $\alpha\in]\lambda,1[$ et le module du quotient est plus petit que $\alpha$ APCR. Alors, par une récurrence immédiate ou un télescopage, on montre que $u_n=O(\alpha^n)$ et alors la série converge absolument. Dans l'autre cas, $\lvert u_n\rvert \geqslant\lvert u_{n_0}\rvert$ APCR, donc $u_n$ ne tend pas vers $0$ et la série diverge grossièrement. Pour le cas $\lambda=1$, tous les cas de figure sont possibles.
\end{proof}
\section{Représentation décimale des réels, HP}
\begin{dfn}[Développement décimal]
	Soit $x$ un réel positif ou nul. Un \textbf{développement décimal} de $x$ est une écriture de la forme $$x=\sum_{n=0}^{+\infty}\frac{d_n}{10^n}$$ où $d_0\in\N$ et $\forall n\in\N\st,\ d_n\in\llbracket0,9\rrbracket$. Les $d_n$ (pour $n\in\N\st$) sont appelés les \textbf{décimales} de $x$, et on écrit : $$x=d_0,d_1d_2\ldots$$
\end{dfn}
\begin{rmq} Par convention, si $x<0$, l'écriture $x=-d_0,d_1d_2\ldots$ signifiera que $x$ est l'opposé de $d_0,d_1d_1\ldots$. \end{rmq}
\begin{dfn}[Développement décimal propre, impropre]
	Si tous les $d_n$ valent $9$ à partir d'un certain rang, le développement décimal est dit \textbf{impropre}. Sinon, il est dit \textbf{propre}.
\end{dfn}
\begin{thm}
	Tout réel $x\geqslant 0$ admet un unique développement décimal propre.
\end{thm}
\begin{rmq} On a le même résultat dans une base $b$ entière supérieure ou égale à 2 quelconque. \end{rmq}
\begin{cor}[Indénombrabilité de $\R$]
	$\R$ est indénombrable.
\end{cor}
\begin{proof}
	On utilise l'argument diagonal de Cantor qui utilise l'existence-unicité du développement décimal propre des réels positifs. L'argument est développé dans la partie "Dénombrabilité" du chapitre "Dénombrements".
\end{proof}

\newpage
\chapter{Sommabilité}
Dans tout ce qui suit, $\mathcal{P}_f(I)$ désigne l'ensemble des parties finies de $I$.
\section{Familles réelles positives}
Dans tout le paragraphe, les familles sont à termes dans $[0,+\infty]$.
\begin{dfn}[Ajout de $+\infty$ à $\R_+$]
	Considérons l'ensemble $[0,+\infty]=\R_+\cup\left\{+\infty\right\}$. On prolonge $\leqslant$, $+$ et $\times$ de la manière suivante :
	\begin{enumerate}
		\item $\forall x\in[0,+\infty],\ x\leqslant+\infty$
		\item $\forall x\in[0,+\infty],\ x+(+\infty)=(+\infty)+x=+\infty$
		\item $\forall x\in]0,+\infty],\ x\times(+\times)=(+\infty)\times x=+\infty$
		\item $0\times (+\infty)=(+\infty)\times 0=0$
	\end{enumerate}
\end{dfn}
\begin{prop}
	On vérifie immédiatement les points suivants :
	\begin{enumerate}
		\item $\leqslant$ reste une relation d'ordre totale
		\item $+$ reste associative et commutative
		\item $\times$ reste associative, commutative et distributive sur $+$
		\item $\leqslant$ reste compatible avec $+$ et $\times$
	\end{enumerate}
\end{prop}
\begin{prop}
	Dans $[0,+\infty]$, les résultats usuels sur les sommes et les produits finis restent valables : relation de Chasles, associativité, distributivité, etc.
\end{prop}
\begin{thm}[Propriété de la borne supérieure prolongée]
	Toute partie non vide de $[0,+\infty]$ admet une  borne supérieure dans $[0,+\infty]$.
\end{thm}
\begin{dfn}[Somme d'une famille, sommabilité]
	Soit $(a_i)_{i\in I}$ une famille à termes dans $[0,+\infty]$. l'ensemble des sommes finies $$\left\{\sum_{i\in J}a_i \tq J\in\mathcal{P}_f(I)\right\}$$ est non vide car il contient la somme vide qui vaut $0$. La \textbf{somme de la famille} $(a_i)_{i\in I}$ est définie par $$\sum_{i\in I}a_i=\underset{J\in\mathcal{P}_f(I)}{\sup}\left(\sum_{i\in J}a_i\right)$$ La famille est dite \textbf{sommable} lorsque $$\sum_{i\in I}a_i<+\infty$$
\end{dfn}
\begin{rmq} Dès que l'un des $a_i$ vaut $+\infty$, la somme de la famille vaut $+\infty$. \end{rmq}
\begin{rmq} La définition coïncide avec celle donnée pour les sommes finies et les sommes presque nulles. \end{rmq}
\begin{prop}
	Soit $(a_i)_{i\in I}$ une famille de somme $S$. Alors 
	\begin{enumerate}
		\item Toute sous-famille est de somme $\leqslant S$
		\item Toute famille majorée par la famille $(a_i)_{i\in I}$ est de somme $\leqslant S$
	\end{enumerate}
\end{prop}
\begin{dfn}
	Soit $(a_n)_{n\in\N}$ une famille à termes dans $\R_+$. Puisque la suite des sommes partielles de la série $\displaystyle \sum a_n$ admet de toute façon une limite, on s'autorisera à écrire $\displaystyle \sum_{n=0}^{+\infty}a_n$ dans tous les cas. Ainsi, si la série diverge, on écrira symboliquement $$\sum_{n=0}^{+\infty}a_n=+\infty$$
\end{dfn}
\begin{prop}[Cas des séries, cas positif]
	Soit $(a_n)_{n\in\N}$ une famille à termes dans $\R_+$. Alors "la somme de la famille est égale à la somme de la série". Pour essayer d'éviter les confusions, on prendra garde à utiliser autant que possible deux notations différentes. Par exemple : $$\sum_{n\in\N}a_n=\sum_{n=0}^{+\infty}a_n$$
\end{prop}
\begin{cor}
	La famille $(a_n)_{n\in\N}$ une famille à termes dans $\R_+$ est sommable si, et seulement si, la série $\displaystyle \sum a_n$ converge, et dans ce cas la somme de la famille est égale à la somme de la série.
\end{cor}
\begin{rmq} On a le même résultat en considérant une famille $(a_n)_{n\geqslant n_0}$ une famille à termes dans $\R_+$ avec $n_0\in\Z$. \end{rmq}
\begin{thm}[Dénombrabilité du support d'une famille sommable, HP]
	Une famille sommable possède un support au plus dénombrable.
\end{thm}
\begin{proof}
	Il suffit de remarquer que pour tout $n\in\N\st$, la famille ne peut prendre ses valeurs qu'un nombre fini de fois dans $\displaystyle \left]\frac{1}{n+1},\frac1n\right]$ et de même avec $]1,+\infty[$, sans quoi la famille ne serait pas sommable. Or ces ensembles partitionnent $\R_+\st$ et leur ensemble est dénombrable. Le support de la famille est donc l'union au plus dénombrable (dénombrable en fait) d'ensemble au plus dénombrables (finis même) : c'est donc un ensemble au plus dénombrable.
\end{proof}
\begin{exa}[Séries de Riemann]
	Si $s>1$, la famille $\displaystyle \left(\frac{1}{n^s}\right)_{n\in\N\st}$  est sommable.
\end{exa}
\begin{thm}[Changement d'indice, cas positif]
	Soit $(a_i)_{i\in I}$ une famille à termes dans $[0,+\infty]$ et $\sigma:J\rightarrow I$ une bijection. Alors $$\sum_{i\in I}a_i=\sum_{j\in J}a_{\sigma(j)}$$
\end{thm}
\begin{rmq} En particulier, la sommabilité est une propriété invariante par permutation de termes. \end{rmq}
\begin{thm}[Linéarité, cas positif]
	Soit $(a_i)_{i\in I}$ et $(b_i)_{i\in I}$ des familles à termes dans $[0,+\infty]$. Soit $\lambda,\mu \in[0,+\infty]$. On a $$\sum_{i\in I}(\lambda a_i+\mu b_i)=\lambda\sum_{i\in I}a_i+\mu\sum_{i\in I}b_i$$
\end{thm}
\begin{thm}[Relation de Chasles, cas positif]
	Soit $(a_i)_{i\in I}$ une famille à termes dans $[0,+\infty]$. Soient $I_1$ et $I_2$ deux ensembles disjoints tels que $I=I_1\sqcup I_2$. On a $$\sum_{i\in I}a_i=\sum_{i\in I_1}a_i+\sum_{i\in I_2}a_i$$
\end{thm}
\begin{rmq} Par une récurrence immédiate, on généralise à un nombre fini d'ensemble deux à deux disjoints. \end{rmq}
\begin{thm}[Sommation par paquets, cas positif]
	Soit $(a_i)_{i\in I}$ une famille à termes dans $[0,+\infty]$. Soit $(I_t)_{t\in T}$ une famille de parties de $I$ deux à deux disjointes telle que $\displaystyle I=\bigsqcup_{t\in T}I_t$. On a $$\sum_{i\in I}a_i=\sum_{t\in T}\left(\sum_{i\in I_t}a_i\right)$$
\end{thm}
\begin{cor}[Théorème de Fubini, cas positif]
	Soit $(a_{i,j})_{(i,j)\in I\times J}$  une famille à termes dans $[0,+\infty]$. On a : $$\sum_{i\in I}\left(\sum_{j\in J}a_{i,j}\right)=\sum_{j\in J}\left(\sum_{i\in I}a_{i,j}\right)=\sum_{(i,j)\in I\times J}a_{i,j}$$
\end{cor}
\begin{rmq} Représenter l'ensemble de sommation, notamment pour les sommations triangulaires peut aider : on peut passer d'une sommation ligne par ligne à une sommation colonne par colonne. Il vaut mieux retenir cela qu'apprendre par cœur  des formules complexes. \end{rmq}
\section{Famille réelles quelconques}
On procèdera toujours en deux temps : d'abord, on vérifie l'hypothèse de sommabilité en effectuant tous les calculs que l'on souhaite après être passées à la valeur absolue grâce au paragraphe précédent ; ensuite on applique les théorèmes du paragraphe. Il existe une variante qui consiste à faire du calcul formel, c'est-à-dire à calculer comme si on avait l'hypothèse de sommabilité, puis à constater que cette hypothèse est vérifiée en expliquant pourquoi avec un passage aux valeurs absolues.
\begin{dfn}[Sommabilité d'une famille réelle quelconque]
	Une famille réelle $(a_i)_{i\in I}$ est dite \textbf{sommable} lorsque $$\sum_{i\in I}\lvert a_i\rvert <+\infty$$autrement dit lorsque la famille réelle positive $(\lvert a_i\rvert)_{i\in I}$ est sommable. Dans ce cas, les familles à termes positifs $(a_i^+)_{i\in I}$ et $(a_i^-)_{i\in I}$ sont sommables (par comparaison) et on \textbf{pose} : $$\sum_{i\in I}a_i=\sum_{i\in I}a_i^+ - \sum_{i\in I}a_i^-$$
\end{dfn}
\begin{rmq} Dans le cas où la famille est positive, on retrouve bien la définition initiale. Même remarque dans le cas des familles finies ou presque nulles. \end{rmq}
\begin{thm}[Inégalité triangulaire, cas réel]
	Si la famille réelle $(a_i)_{i\in I}$ est sommable, alors $$\left\lvert\sum_{i\in I}a_i\right\rvert \leqslant\sum_{i\in I}\lvert a_i\rvert$$
\end{thm}
On observe désormais si les résultats du paragraphe précédent restent vrais.
\begin{prop}
	Toute sous-famille d'une famille sommable reste sommable, mais on ne peut rien dire sur sa somme. Attention, on ne peut rien dire d'une famille majorée, à moins de majorer la famille $(\lvert a_i\rvert)_{i\in I}$.
\end{prop}
\begin{thm}[Cas des séries, cas réel]
	Soit $(a_i)_{i\in\N}$ une famille réelle quelconque. La famille $(a_i)_{i\in\N}$ est sommable si, et seulement si, la série $\displaystyle \sum a_n$ est absolument convergente, et dans ce cas on a : $$\sum_{i\in\N}a_i=\sum_{i=0}^{+\infty}a_i$$
\end{thm}
\begin{thm}[Changement d'indice, cas réel]
	Soit $(a_i)_{i\in I}$ une famille réelle quelconque et $\sigma:J\rightarrow I$ une bijection. Alors la famille $(a_i)_{i\in I}$ est sommable si, et seulement si, la famille $(a_{\sigma(j)})_{j\in J}$, et dans ce cas on a : $$\sum_{i\in I}a_i=\sum_{j\in J}a_{\sigma(j)}$$
\end{thm}
\begin{thm}[Linéarité, cas réel]
	Soit $\lambda,\mu\in\R$. Si $(a_i)_{i\in I}$ et $(b_i)_{i\in I}$ réelles sont sommables, alors $(\lambda a_i+\mu b_i)_{i\in I}$ est sommable et on a : $$\sum_{i\in I}(\lambda a_i+\mu b_i)=\lambda\sum_{i\in I}a_i+\mu\sum_{i\in I}b_i$$
\end{thm}
\begin{thm}[Relation de Chasles, cas réel]
	Si $I_1$ et $I_2$ sont deux ensembles disjoints tels que $I=I_1\sqcup I_2$, alors $(a_i)_{i\in I}$ est sommable si, et seulement si, $(a_i)_{i\in I_1}$ et $(a_i)_{i\in I_2}$ son sommables, et dans ce cas on a : $$\sum_{i\in I}a_i=\sum_{i\in I_1}a_i+\sum_{i\in I_2}a_i$$ Par une récurrence immédiate, on généralise à un nombre fini d'ensembles deux à deux disjoints.
\end{thm}
\begin{thm}[Sommation par paquets, cas réel]
	Soit $(a_i)_{i\in I}$ une famille réelle quelconque. Soit $(I_t)_{t\in T}$ une famille de parties de $I$ deux à deux disjointes telle que $\displaystyle I=\bigsqcup_{t\in T}I_t$. \textbf{Si} $(a_i)_{i\in I}$ est sommable, \textbf{alors} chaque $(a_i)_{i\in I_t}$ est sommable et la famille $\displaystyle \left(\sum_{i\in I_t}a_i\right)_{t\in T}$ est sommable et on a : $$\sum_{i\in I}a_i=\sum_{t\in T}\left(\sum_{i\in I_t}a_i\right)$$
\end{thm}
\begin{cor}[Théorème de Fubini, cas réel]
	Soit $(a_{i,j})_{(i,j)\in I\times J}$  une famille réelle sommable. On a : $$\sum_{i\in I}\left(\sum_{j\in J}a_{i,j}\right)=\sum_{j\in J}\left(\sum_{i\in I}a_{i,j}\right)=\sum_{(i,j)\in I\times J}a_{i,j}$$
\end{cor}
\section{Familles complexes}
\begin{dfn}[Sommabilité d'une famille complexe]
	Une famille complexe $(a_i)_{i\in I}=(x_i+\text{i}y_i)_{i\in I}$ est dite \textbf{sommable} lorsque $$\sum_{i\in I}\lvert a_i\rvert <+\infty$$ Dans ce cas, les familles réelles $(x_i)_{i\in I}$ et $(y_i)_{i\in I}$ sont sommables et on \textbf{pose} : $$\sum_{i\in I}a_i=\sum_{i\in I}x_i +\text{i}\sum_{i\in I}y_i$$ L'ensemble des familles sommables est noté $l^1(I)$.
\end{dfn}
\begin{rmq} Dans le cas où la famille est réelle, on retrouve bien la définition initiale. Même remarque dans le cas des familles finies ou presque nulles.
	On observe désormais si les résultats du paragraphe précédent restent vrais. \end{rmq}
\begin{prop}
	Toute sous-famille d'une famille sommable reste sommable, mais on ne peut rien dire sur sa somme. Attention, on ne peut rien dire d'une famille majorée, à moins de majorer la famille $(\lvert a_i\rvert)_{i\in I}$.
\end{prop}
\begin{thm}[Inégalité triangulaire, cas complexe]
	Si la famille complexe $(a_i)_{i\in I}$ est sommable, alors $$\left\lvert\sum_{i\in I}a_i\right\rvert \leqslant\sum_{i\in I}\lvert a_i\rvert$$
\end{thm}
\begin{thm}[Cas des séries, cas complexe]
	Soit $(a_i)_{i\in\N}$ une famille complexe. La famille $(a_i)_{i\in\N}$ est sommable si, et seulement si, la série $\displaystyle \sum a_n$ est absolument convergente, et dans ce cas on a : $$\sum_{i\in\N}a_i=\sum_{i=0}^{+\infty}a_i$$
\end{thm}
\begin{thm}[Changement d'indice, cas complexe]
	Soit $(a_i)_{i\in I}$ une famille complexe et $\sigma:J\rightarrow I$ une bijection. Alors la famille $(a_i)_{i\in I}$ est sommable si, et seulement si, la famille $(a_{\sigma(j)})_{j\in J}$, et dans ce cas on a : $$\sum_{i\in I}a_i=\sum_{j\in J}a_{\sigma(j)}$$
\end{thm}
\begin{thm}[Linéarité, cas complexe]
	Soit $\lambda,\mu\in\C$. Si $(a_i)_{i\in I}$ et $(b_i)_{i\in I}$ complexes sont sommables, alors $(\lambda a_i+\mu b_i)_{i\in I}$ est sommable et on a : $$\sum_{i\in I}(\lambda a_i+\mu b_i)=\lambda\sum_{i\in I}a_i+\mu\sum_{i\in I}b_i$$
\end{thm}
\begin{thm}[Relation de Chasles, cas complexe]
	Si $I_1$ et $I_2$ sont deux ensembles disjoints tels que $I=I_1\sqcup I_2$, alors $(a_i)_{i\in I}$ complexe est sommable si, et seulement si, $(a_i)_{i\in I_1}$ et $(a_i)_{i\in I_2}$ sont sommables, et dans ce cas on a : $$\sum_{i\in I}a_i=\sum_{i\in I_1}a_i+\sum_{i\in I_2}a_i$$ Par une récurrence immédiate, on généralise à un nombre fini d'ensembles deux à deux disjoints.
\end{thm}
\begin{thm}[Sommation par paquets, cas complexe]
	Soit $(a_i)_{i\in I}$ une famille complexe. Soit $(I_t)_{t\in T}$ une famille de parties de $I$ deux à deux disjointes telle que $\displaystyle I=\bigsqcup_{t\in T}I_t$. \textbf{Si} $(a_i)_{i\in I}$ est sommable, \textbf{alors} chaque $(a_i)_{i\in I_t}$ est sommable et la famille $\displaystyle \left(\sum_{i\in I_t}a_i\right)_{t\in T}$ est sommable et on a : $$\sum_{i\in I}a_i=\sum_{t\in T}\left(\sum_{i\in I_t}a_i\right)$$
\end{thm}
\begin{cor}[Théorème de Fubini, cas complexe]
	Soit $(a_{i,j})_{(i,j)\in I\times J}$  une famille complexe sommable. On a : $$\sum_{i\in I}\left(\sum_{j\in J}a_{i,j}\right)=\sum_{j\in J}\left(\sum_{i\in I}a_{i,j}\right)=\sum_{(i,j)\in I\times J}a_{i,j}$$
\end{cor}
\section{Produit de familles}
\begin{thm}[Produit de familles, cas positif]
	Soit $(a_i)_{i\in I}$ et $(b_j)_{j\in J}$ deux familles à termes dans $[0,+\infty]$. Alors $$\sum_{(i,j)\in I\times J}a_ib_j=\left(\sum_{i\in I}a_i\right)\left(\sum_{j\in J}b_j\right)$$
\end{thm}
\begin{thm}[Produit de familles, cas complexe]
	Soit $(a_i)_{i\in I}$ et $(b_j)_{j\in J}$ deux familles complexes sommables. Alors la famille $(a_ib_j)_{(i,j)\in I\times J}$ est sommable et on a :$$\sum_{(i,j)\in I\times J}a_ib_j=\left(\sum_{i\in I}a_i\right)\left(\sum_{j\in J}b_j\right)$$
\end{thm}
\begin{rmq} Par une récurrence immédiate, on généralise le résultat à un produit fini de sommes. \end{rmq}
\begin{cor}[Produit de Cauchy de deux séries absolument convergentes]
	Soit $\displaystyle \sum u_n$ et $\displaystyle \sum v_n$ deux séries à termes complexes absolument convergentes. On pose $$\forall n\in\N,\ w_n=\sum_{k=0}^{n}u_kv_{n-k}$$ Alors, la série $\displaystyle \sum w_n$ est absolument convergente et on a :$$\sum_{n=0}^{+\infty}w_n=\left(\sum_{n=0}^{+\infty}u_n\right)\left(\sum_{n=0}^{+\infty}v_n\right)$$
\end{cor}
\begin{cor}
	$$\forall(z,z')\in\C^2,\ \exp(z+z')=\exp(z)\exp(z')$$
\end{cor}
\begin{proof}
	Utiliser le produit de Cauchy et faire apparaître un binôme de Newton en multipliant les $w_n$ par $\displaystyle \frac{n!}{n!}$.
\end{proof}

\newpage
\chapter{Dénombrements}
\section{Généralités}
\begin{dfn}[Cardinal d'un ensemble fini]
	Soit $E$ un ensemble. S'il existe une bijection de $E$ sur un certain $\llbracket1,n\rrbracket$, alors ce $n$ est unique. On dit que alors que $E$ est \textbf{fini} et son \textbf{cardinal} est alors défini par :
	$$\card(E)=n$$ On rencontre aussi souvent la notation $\left\lvert E\right\rvert$, voire plus rarement la notation $\#E$.
\end{dfn}
\begin{exa}
	Pour tout ensemble fini $E$ et pour tout $\lambda\in\C$ (marche plus généralement dans une structure algébrique comportant un + commutatif), on a $$\sum_{x\in E}\lambda = \card(E)\times\lambda$$
\end{exa}
\begin{thm}[Principe des tiroirs, cas général]
	Soit $E$ et $F$ des ensembles finis tels que $\card(E)>\card(F)$. Alors il n'existe pas d'injection de $E$ dans $F$.
\end{thm}
\begin{prop}
	Si $E$ est fini et si $E$ et $F$ peuvent être mis en bijection, alors $F$ est fini de même cardinal que $E$. Réciproquement, si $E$ et $F$ sont de même cardinal, alors ils peuvent être mis en bijection.
\end{prop}
\begin{thm}
	Soit $E$ un ensemble fini. Toute partie $A$ de $E$ est finie, et son cardinal vérifie : $$\card(A)\leqslant\card(E)$$ De plus, on égalité si, et seulement si, $A=E$.
\end{thm}
\begin{cor}
	Soit $E$ et $F$ finis de même cardinal ainsi que $f:E\rightarrow F$. Alors $f$ est bijective si, et seulement si, $f$ est injective si, et seulement si, $f$ est surjective.
\end{cor}
\begin{prop}[Une formule utile]
	Soit $E$ un ensemble fini et $A\subset E$. Alors $$\card(A)=\sum_{x\in E} \mathds{1}_A(x)$$
\end{prop}
\begin{rmq} Avec le formalisme des sommes presque nulles, la formule reste vraie si $A$ est une partie fini d'un ensemble quelconque. \end{rmq}
\begin{lem}[Simple additivité du cardinal]
	Soit $A$ et $B$ deux parties finies disjointes d'un ensemble $E$. Alors $A\cup B$ est finie de cardinal $$\card(A\cup B)=\card(A)+\card(B)$$
\end{lem}
\begin{exa}
	Si $E$ est fini, $A\subset$ et $\overline{A}$ désigne le complémentaire de $A$ dans $E$, alors $A$ et $\overline{A}$ sont finis et on a $$\card(\overline{A})=\card(E)-\card(A)$$ En pratique, on utilisera souvent cette formule car il est parfois beaucoup plus simple de dénombrer le complémentaire d'un ensemble.
\end{exa}
\begin{cor}[Généralisation]
	Soit $\mathcal{A}\subset\mathcal{P}(E)$ un ensemble fini de parties finies de deux à deux disjointes. Alors $\displaystyle \bigsqcup_{A\in\mathcal{A}}A$ est fini et on a $$\card\left(\bigsqcup_{A\in\mathcal{A}}A\right)=\sum_{A\in\mathcal{A}}\card(A)$$
\end{cor}
\begin{cor}[Lemme des bergers]
	Soit $f:E\rightarrow F$ avec $F$ fini. Si chacun des $n$ éléments de $F$ admet exactement $k$ antécédents, alors $E$ est fini de cardinal $$\card(E)=n\times k$$
\end{cor}
\begin{thm}[Cardinal d'une union]
	Soit $A$ et $B$ deux parties finies d'un ensemble $E$. Alors $A\cup B$ est finie et on a $$\card(A\cup B)=\card(A)+\card(B)-\card(A\cap B)$$
\end{thm}
\begin{prop}
	L'image directe d'un ensemble fini est finie.
\end{prop}
\begin{prop}[Inégalité de Boole pour les cardinaux]
	Soit $\mathcal{A}\subset\mathcal{P}(E)$ un ensemble fini de parties finies. Alors $\displaystyle \bigcup_{A\in\mathcal{A}}A$ est finie et on a : $$\card\left(\bigcup_{A\in\mathcal{A}}A\right)\leqslant\sum_{A\in\mathcal{A}}\card(A)$$
\end{prop}
\begin{thm}[Formule du crible de Poincaré / principe d'inclusion-exclusion, HP]
	Soit $A_1,\ldots,A_n$ des parties finies d'un ensemble $E$. Alors $\displaystyle \bigcup_{i=1}^{n}A_i$ est finie et on a la formule du \textbf{crible de Poincaré} : $$\card\left(\bigcup_{i=1}^{n}A_i\right)=\sum_{k=1}^{n}(-1)^{k-1}\sum_{1\leqslant i_1<\ldots<i_k\leqslant n}\card\left(\bigcap_{j=1}^{k}A_{i_j}\right)$$
\end{thm}
\begin{proof}
	Dénombrer le complémentaire en utilisant la distributivité généralisée sur les fonctions indicatrices des $A_i$
\end{proof}
\begin{thm}[Cardinal d'un produit cartésien]
	Si $E$ et $F$ sont des ensembles finis, alors l'ensemble $E\times F$ est fini de cardinal $$\card(E\times F)=\card(E)\times\card(F)$$
\end{thm}
\begin{cor}
	Si $E_1,\ldots,E_n$ sont des ensembles finis, alors $E_1\times\ldots\times E_n$ est fini de cardinal $$\card(E_1\times\ldots\times E_n)=\card(E_1)\times\ldots\times\card(E_n)$$
\end{cor}
\begin{thm}[Cardinal de l'ensemble des applications d'un ensemble fini vers un autre]
	Si $E$ et $F$ sont finis, alors $F^E$ est fini de cardinal $$\card\left(F^E\right)=\card(F)^{\card(E)}$$
\end{thm}
\begin{rmq} C'est ce résultat qui justifie a posteriori la notation $F^E$. \end{rmq}
\begin{thm}
	Si $E$ est un ensemble fini, alors $\mathcal{P}(E)$ est fini de cardinal $$\card\left(\mathcal{P}(E)\right)=2^{\card(E)}$$
\end{thm}
\section{Dénombrabilité, HP, au programme de spé}
\begin{dfn}[Dénombrabilité]
	Un ensemble $E$ est dit \textbf{dénombrable} s'il peut être mis en bijection avec $\N$. Il est dit \textbf{au plus dénombrable} s'il est fini ou dénombrable. il est dit \textbf{indénombrable} s'il n'est pas au plus dénombrable.
\end{dfn}
\begin{prop}
	$E$ est au plus dénombrable si, et seulement si, il existe une injection de $E$ dans $\N$.
\end{prop}
\begin{cor}
	Toute partie d'un ensemble au plus dénombrable est au plus dénombrable.
\end{cor}
\begin{lem}
	Soit $\varphi:E\rightarrow F$. Si $\varphi$ est injective et $F$ est au plus dénombrable, alors $E$ est au plus dénombrable. Si $\varphi$ est surjective et $E$ est au plus dénombrable, alors $F$ est au plus dénombrable.
\end{lem}
\begin{thm}
	Si $E_1,\ldots,E_n$ sont au plus dénombrables, alors $E_1\times\ldots\times E_n$ est au plus dénombrable.
\end{thm}
\begin{cor}
	$\Z$ et $\Q$ sont dénombrables.
\end{cor}
\begin{thm}
	Une union au plus dénombrable d'ensemble au plus dénombrables est au plus dénombrable.
\end{thm}
\begin{prop}
	$\mathcal{P}(\N)$ est indénombrable.
\end{prop}
\begin{proof}
	Il s'agit d'un cas particulier du théorème de Cantor qui stipule que pour tout ensemble $E$, il n'existe pas de surjection de $E$ dans $\mathcal{P}(E)$ (et donc $a$ $fortiori$ pas de bijection). Pour le démontrer, raisonner par l'absurde en considérant une éventuelle surjection puis en s'intéressant à l'ensemble des éléments de $E$ qui n'appartiennent pas à leur image par $f$. Par hypothèse, cette partie de $E$ admet un antécédent qui ne peut ni appartenir à son image, ni ne pas lui appartenir, ce qui est absurde.
\end{proof}
\begin{prop}
	$\R$ est indénombrable.
\end{prop}
\begin{proof}
	C'est l'argument diagonal de Cantor. Supposer que l'on puisse énumérer tous les réels, puis conclure à une absurdité en prenant le réel formé par des chiffres différents de 9 et différents des chiffres qui apparaissent sur la diagonale. Ce réel n'apparaît alors pas dans l'énumération, ce qui est absurde par existence et unicité du développement propre.
\end{proof}

\section{Choisir $p$ objets parmi $n$}
\subsection{Avec ordre et avec répétition}
On dénombre des $p$-uplets.
\begin{thm}
	Soit $F$ de cardinal $n$. Le nombre de $p$-uplets à valeurs dans $F$ est $n^p$.
\end{thm}
\subsection{Avec ordre et sans répétition}
On dénombre des $p$-uplets injectifs.
\begin{thm}[Nombre d'arrangements]
	Soit $F$ de cardinal $n\geqslant p$. Alors le nombre de $p$-uplets à valeurs distinctes dans $F$ vaut $$A_n^p=n(n-1)\ldots(n-p+1)=\frac{n!}{(n-p)!}$$ $A_n^p$ s'appelle un nombre d'\textbf{arrangements}.
\end{thm}
\begin{cor}[Nombre d'injections]
	Soit $E$ de cardinal $p$ et $F$ de cardinal $n\geqslant p$. Alors le nombre d'injections de $E$ dans $F$ vaut $A_n^p$.
\end{cor}
\begin{cor}[Nombre de permutations]
	Si $E$ est fini de cardinal $n$ alors il existe $n!$ permutations de $E$.
\end{cor}
\subsection{Sans ordre et sans répétition}
On dénombre des parties à $p$ éléments.
\begin{thm}[Nombre de combinaisons]
	Soit $F$ de cardinal $n$. Le nombre de parties de $F$ de cardinal $p$ vaut $$C_n^p=\binom{n}{p}=\frac{n(n-1)\ldots(n-p+1)}{p!}=\frac{n!}{p!(n-p)!}=\frac{A_n^p}{p!}$$ C'est pour cela que les coefficients binomiaux s'appellent aussi des nombres de \textbf{combinaisons}.
\end{thm}
\begin{cor}[Nombre d'applications strictement croissantes]
	Soit $E$ de cardinal $p$ et $F$ de cardinal $n\geqslant p$. Alors le nombre d'applications strictement croissantes de $E$ dans $F$ vaut $\binom{n}{k}$.
\end{cor}
\subsection{Sans ordre et avec répétition, HP, exo classique}
\begin{thm}
	Le nombre de façons de choisir $p$ objets parmi $n$ sans ordre et avec répétition est égal à $\displaystyle \binom{n+p-1}{p}$
\end{thm}
\begin{proof}
	 On se ramène à $F=\llbracket1,n\rrbracket$ pour plus de clarté. Si on choisit les éléments deux à deux distincts, on se ramène au cas précédent. Si les répétitions sont permises, tout revient à choisir un $p$-uplet croissant au sens large. On prend alors un $p$-uplet croissant $(x_1\leqslant\ldots,x_p)$ et on l'envoie sur le $p$-uplet $(x_1,x_2+1,\ldots,x_p+(p-1))$ qui est alors strictement croissant. On vérifie que cet "envoi" est bijectif vers les $p$-uplets à valeurs dans $\llbracket1,n+p-1\rrbracket$ puis on conclut alors par le premier sous-cas.
\end{proof}
\begin{cor}[Nombre d'applications croissantes]
	Le nombre d'applications croissantes d'un ensemble à $p$ éléments dans un ensemble à $n\leqslant p$ éléments vaut $\displaystyle \binom{n+p-1}{p}$.
\end{cor}
\begin{rmq} On peut retenir la méthode du "stars and bars" : lorsqu'on choisit $p$ objets parmi des objets de $n$ types sans ordre avec répétition, cela revient à choisir les $n-1$ cloisons entre les multiples choix des différents éléments, donc on a bien $\displaystyle \binom{n+p-1}{p}$ choix possibles. \end{rmq}

\newpage
\chapter{Probabilités}
\section{Probabilités sur univers fini}
\begin{dfn}[Univers, issue, événement, événement élémentaire]
	Soit $\Omega$ un ensemble non vide, appelé \textbf{univers}. Chaque élément $\omega\in\Omega$ correspond à une \textbf{issue} possible de l'expérience. Formellement, un \textbf{événement} est une partie $A\subset \Omega$. C'est donc un ensemble d'issues possibles. On dit que $A$ est \textbf{réalisé} lorsque l'issue de l'expérience est un $\omega\in A$. Dans toute la suite, nous nous restreindrons au cas où $\Omega$ est \textbf{fini} (et toujours non vide). Si $A=\left\{\omega\right\}$ est un singleton, on parle d'\textbf{événement élémentaire}.
\end{dfn}
\begin{rmq} Même si formellement un événement $A$ est défini comme une partie de $\Omega$, la plupart du temps on décrira $A$ par une proposition logique $p(\omega)$ qui dépend de $\omega\in\Omega$. Dans ce cas, $A$ est à comprendre comme l'ensemble $\left\{\omega\in\Omega\tq p(\omega)\right\}$ \end{rmq}
\begin{dfn}[Événements "ou", "et" et contraire]
	Soit $A$ et $B$ deux événements, décrits respectivement par des propositions logiques $p$ et $q$. L'événement $A$ \textbf{ou} $B$ désigne l'événement $A\cup B$. Il est décrit par $p\lor q$. L'événement $A$ \textbf{et} $B$ désigne l'événement $A\cap B$. Il est décrit par $p\land q$. L'événement \textbf{contraire} de $A$ est l'événement $\overline{A}=\Omega\setminus A$. Il est décrit par $\lnot p$.
\end{dfn}
\begin{dfn}[Événement impossible, événements incompatibles, système complet d'événements]
	L'événement \textbf{impossible} est l'événement $\emptyset$. Il doit son nom au fait qu'il ne peut jamais être réalisé. Si $A\cap B=\emptyset$, on dit que $A$ et $B$ sont \textbf{incompatibles}. Enfin, si $\displaystyle\Omega=\bigsqcup_{1\leqslant i\leqslant n} A_i$, on dit que les $A_i$ forment un \textbf{système complet d'événements}.
\end{dfn}
\begin{exa}
	Les événements élémentaires forment un système complet d'événements.
\end{exa}
\begin{exa}
	Si $A$ est un événement, alors $(A,\overline{A})$ est un système complet d'événements.
\end{exa}
\begin{dfn}[Probabilité, espace probabilisé]
	Une \textbf{probabilité} sur $\Omega$ est une application 
	$$
	\begin{array}[t]{lrcl}
		\P : & \mathcal{P}(\Omega) & \longrightarrow & [0,1] \\
		& A & \longmapsto & \P(A)
	\end{array}
	$$ telle que $\P(\Omega)=1$ et que pour toutes parties disjointes $A$ et $B$ de $\Omega$, on ait $$\P(A\cup B)=\P(A)+\P(B)$$ On dit que $(\Omega,\P)$ est un \textbf{espace probabilisé}.
\end{dfn}
\begin{prop}
	Puisque $\Omega$ et $\emptyset$ sont disjoints, on a immédiatement $\P(\emptyset)=0$.
\end{prop}
\begin{prop}
	Si $A_1,\ldots,A_n$ sont deux à deux disjoints, alors on a $$\P\left(\bigsqcup_{1\leqslant i\leqslant n} A_i\right)=\sum_{1\leqslant i\leqslant n}\P(A_i)$$
\end{prop}
\begin{prop}[Formule des probabilités totales, V1]
	Si les $A_i$ forment un système complet d'événements et si $B\subset \Omega$, on a $$\P(B)=\sum_{i=1}^{n}\P(B\cap A_i)$$
\end{prop}
\begin{thm}[Caractérisation par les probabilités élémentaires]
	Soit un univers fini $\Omega=\left\{\omega_1,\ldots,\omega_n\right\}$ et $p_1,\ldots,p_n$ des réels positifs dont la somme vaut $1$. Alors il existe une unique probabilité sur $\Omega$ telle que $$\forall i \in\llbracket1,n\rrbracket, \ \P(\left\{\omega_i\right\})=p_i$$
\end{thm}
\begin{cor}
	Si deux probabilités coïncident sur les événements élémentaires, alors coïncident partout.
\end{cor}
\begin{dfn}[Probabilité uniforme]
	La \textbf{probabilité uniforme} sur $\Omega$ est l'unique probabilité qui vaut $\displaystyle\frac{1}{\card(\Omega)}$ sur chaque événement élémentaire. Elle vérifie donc $$\forall A\subset\Omega, \ \P(A)=\frac{\card(A)}{\card(\Omega)}$$
\end{dfn}
\begin{thm}
	Soit $(\omega,\P)$ un espace probabilisé. On a les propriétés suivantes :
	\begin{enumerate}
		\item $\forall(A,B)\in\mathcal{P}(\Omega)^2, \ A\subset B \Longrightarrow \P(A)\leqslant \P(B)$ (croissance)
		\item $\forall A \in \mathcal{P}(\Omega), \ \P(\overline{A})=1-\P(A)$
		\item $\forall(A,B)\in\mathcal{P}(\Omega)^2, \ \P(A\cup B)=\P(A)+\P(B)-\P(A\cap B)$
	\end{enumerate}
\end{thm}
\begin{prop}[Inégalité de Boole]
	Soit $A_1,\ldots,A_n$ des événements. On a $$\P\left(\bigcup_{1\leqslant i\leqslant n}A_i\right)\leqslant\sum_{1\leqslant i\leqslant n}\P(A_i)$$
\end{prop}
\begin{dfn}[Événement presque sûr, HP]
	Un événement est dit \textbf{presque sûr} lorsqu'il est de probabilité 1. L'abréviation \textbf{p.s.} est standard.
\end{dfn}
\begin{prop}[HP]
	Une intersection finie d'événements presque sûrs est presque sûre.
\end{prop}
\begin{proof}
	Procéder par récurrence. On a simplement besoin du fait que si $A$ et $B$ sont presque sûrs, alors $A\cap B$ est presque sûr. Pour montrer cela, il suffit d'utiliser le fait que $\P(A\cap B)=\P(A)+\P(B)-\P(A\cup B)$. Or, par croissance de la probabilité, on a nécessairement $\P(A\cup B)=1$, ce qui achève la preuve.
\end{proof}
\begin{dfn}[Probabilité conditionnelle]
	Soit $B$ un événement tel que $\P(B)>0$. pour tout événement $A$, la \textbf{probabilité conditionnelle de $A$ sachant $B$} est définie par $$\P(A|B)=\P_B(A)=\frac{\P(A\cap B)}{\P(B)}$$
\end{dfn}
\begin{prop}
	$\P(\cdot| B)=\P_B$ est une probabilité sur $\Omega$.
\end{prop}
\begin{rmq} L'événement $B$ (et tout sur-événement de $B$) est $\P_B$-presque sûr. \end{rmq}
\begin{thm}[Formule des probabilités composées]
	Soit $A_1,\ldots,A_n$ des événements tels que $\P(A_1\cap\ldots\cap A_{n-1})>0$. Alors $$\P(A_1\cap\ldots\cap A_n)=\P(A_1)\times \P(A_2|A_1)\times\ldots\times \P(A_n|A_1\cap\ldots\cap A_{n-1})$$
\end{thm}
\begin{thm}[Formule des probabilités totales, V2]
	Soit $(A_i)_{1\leqslant i\leqslant n}$ un système complet d'événements tel que $\forall i \in\llbracket1,n\rrbracket, \ \P(A_i)>0$. Soit $B$ un événement. On a $$\P(B)=\sum_{i=1}^{n}\P_{A_i}(B)\P(A_i)$$
\end{thm}
\begin{rmq} Il existe une convention qui stipule que $\P_A(B)\P(A)=0$ lorsque $\P(A)=0$. Dans ce cas, la formule des probabilités totales V2 reste vraie en toute généralité. \end{rmq}
\begin{thm}[Première formule de Bayes]
	Soit $A$ et $B$ deux événements de probabilités non nulles. Alors $$\P_B(A)=\frac{\P_A(B)\P(A)}{\P(B)}$$
\end{thm} 
\begin{thm}[Seconde formule de Bayes]
	Soit $(A_i)_{1\leqslant i\leqslant n}$ un système complet d'événements tel que $\forall i \in\llbracket1,n\rrbracket, \ \P(A_i)>0$. Soit $B$ un événement de probabilité non nulle. On a $$\forall i\in\llbracket1,n\rrbracket, \ \P_B(A_i)=\frac{\P_{A_i}(B)\P(A_i)}{\displaystyle\sum_{j=1}^{n}\P_{A_j}(B)\P(A_j)}$$
\end{thm}
\begin{dfn}[Événements indépendants]
	On dit que les événements $A$ et $B$ sont \textbf{indépendants} lorsque $$\P(A\cap B)=\P(A)\P(B)$$
\end{dfn}
\begin{dfn}[Événements mutuellement indépendants]
	Soit $(A_i)_{1\leqslant i\leqslant n}$ une famille finie d'événements. Les $A_i$ sont dits \textbf{mutuellement indépendants} (ou plus simplement, \textbf{indépendants}) lorsque $$\forall I\subset\llbracket1,n\rrbracket, \ \P\left(\bigcap_{i\in I}A_i\right)=\prod_{i\in I}\P(A_i)$$ On rappelle qu'on a par convention $$\bigcap_{i\in\emptyset}A_i=\Omega$$
\end{dfn}
\begin{rmq} Si $I=\emptyset$ ou $I$ est un singleton, il n'y a rien à vérifier. En pratique, il y a donc $2^n-n-1$ vérifications à effectuer. \end{rmq}

\begin{prop}
	Soit $A_1,\ldots,A_n$ des événements mutuellement indépendants. Alors pour tout $\sigma\in\mathcal{S}_n$, les événements $A_{\sigma(i)}$ sont mutuellement indépendants. De plus, pour toute sous-famille $(A_{i_1},\ldots,A_{i_r})$, les événements restent indépendants.
\end{prop}
\begin{thm}
	Soit $A_1,\ldots,A_n$ des événements mutuellement indépendants. Soit $B_1,\ldots,B_n$ des événements que $$\forall i\in\llbracket1,n\rrbracket, \ B_i\in\left\{A_i,\overline{A_i}\right\}$$ Alors, les événements $B_1,\ldots,B_n$ sont mutuellement indépendants.
\end{thm}
\begin{thm}[Théorème de regroupement]
	Soit $A_1,\ldots,A_n$ des événements mutuellement indépendants. et soit $p\in\llbracket1,n-1\rrbracket$. Alors
	\begin{enumerate}
		\item $\displaystyle \bigcap_{1\leqslant i\leqslant p}A_i$ et $\displaystyle \bigcap_{p<i\leqslant n}A_i$ sont indépendants.
		\item $\displaystyle \bigcap_{1\leqslant i\leqslant p}A_i$ et $\displaystyle \bigcup_{p<i\leqslant n}A_i$ sont indépendants.
		\item $\displaystyle \bigcup_{1\leqslant i\leqslant p}A_i$ et $\displaystyle \bigcap_{p<i\leqslant n}A_i$ sont indépendants.
		\item $\displaystyle \bigcup_{1\leqslant i\leqslant p}A_i$ et $\displaystyle \bigcup_{p<i\leqslant n}A_i$ sont indépendants.
	\end{enumerate}
\end{thm}

\section{Variables aléatoires}
\begin{dfn}[Variable aléatoire]
	Soit $(\Omega,\P)$ un espace probabilisé fini. Une \textbf{variable aléatoire} définie sur $\Omega$ est une application
	$$
	\begin{array}[t]{lrcl}
		X : & \Omega & \longrightarrow & E \\
		& \omega & \longmapsto & X(\omega)
	\end{array}
	$$ En particulier, si $E\subset\R$, la variable aléatoire est dite \textbf{réelle}.
\end{dfn}
\begin{dfn}
	Soit $A$ un événement.  On définit : $$\left\{X\in A\right\}=\left\{\omega\in\Omega \ \tq \ X(\omega)\in A\right\}=X\inv(A)$$ On pourra par exemple considérer des événements du type $\left\{X=x\right\}$ où $x\in E$ est fixé. Si $X$ est réelle, on pourra considérer des événements du type $\left\{X\leqslant x\right\}$. La notation $\P(\left\{X\in A\right\})$ étant lourde, on lui préfèrera la notation suivante, légèrement abusive : $$\P(X\in A)$$
\end{dfn}
\begin{exa}
	Lorsque $x$ parcourt $X(\Omega)$, les événements $\left\{X=x\right\}$ forment un SCE.
\end{exa}
\begin{dfn}[Loi d'une variable aléatoire]
	La \textbf{loi de $X$} est définie sur $\mathcal{P}(E)$ par $$\forall A\subset E, \ \P_X(A)=\P(X\in A)=\P(X\inv(A))$$ Elle induit une loi de probabilité sur $X(\Omega)$ (définie sur $\mathcal{P}(X(\Omega))$ donc).
\end{dfn}
\begin{dfn}[Même loi]
	Soit $(\Omega,\P')$ un espace probabilisé et $Y:\Omega'\rightarrow E$ une variable aléatoire. On dit que $X$ et $Y$ \textbf{ont même loi} lorsque $$\forall A\subset E, \ \P(X\in A)=\P'(Y\in A)$$ On note alors $$X\sim Y$$Notons que le seul pré-requis est que $X$ et $Y$ pointent vers un même ensemble $E$, les espaces probabilisés peuvent être les mêmes comme différents.
\end{dfn}
\begin{prop}
	Si $A\subset E$, on a $$\P(X\in A)=\sum_{x\in A}\P(X=x)$$
\end{prop}
\begin{prop}[La loi de $X$ est caractérisée par les probabilités élémentaires $\P(X=x)$]
	$X$ et $Y$ ont même loi si, et seulement si, elles ont les mêmes probabilités élémentaires sur $E$. On dit que "la loi de $X$ est caractérisée par les probabilités élémentaires $\P(X=x)$".
\end{prop}
\begin{rmq} \textbf{Attention}, ce n'est pas parce que $X$ et $Y$ ont même loi qu'elles sont égales, il suffit de considérer les variables aléatoires qui correspondent au résultat pile et au résultat face d'un pièce équilibrée. \end{rmq}
\begin{dfn}[Variable aléatoire image]
	Soit $X$ à valeurs dans $E$ et $f$ une application de $E$ dans $F$. Alors $f\circ X$ est une variable aléatoire à valeurs dans $F$, appelée \textbf{image} de $X$ par $f$. On la note plus simplement $f(X)$ et elle vérifie $$\forall B\subset F, \ \P(f(X)\in B)=\P(X\in f\inv(B))$$ Autrement dit, on a $$\P_{f(X)}=\P_X\circ f\inv$$
\end{dfn}
\begin{prop}[Transfert de loi]
	Si $X\sim Y$, alors $f(X)\sim f(Y)$.
\end{prop}
\begin{dfn}[Couple de variables aléatoires]
	Soit $(\omega,\P$ un espace probabilisé et $X$ et $Y$ des variables aléatoires définies sur $\Omega$ à valeurs respectivement dans $E$ et $F$. le \textbf{couple de variables aléatoires} $(X,Y)$ est la variable aléatoire à valeurs dans $E\times F$ définie par 
	$$
	\begin{array}[t]{lrcl}
		(X,Y) : & \Omega & \longrightarrow & E\times F \\
		& \omega & \longmapsto & (X(\omega),Y(\omega))
	\end{array}
	$$ La notation $\left\{X=x,Y=y\right\}$ désignera en fait l'événement $\left\{X=x\right\}\cap\left\{Y=y\right\}$. De même, on peut imaginer des événements du type $\left\{X\leqslant x, Y\geqslant y\right\}$, etc.
\end{dfn}
\begin{dfn}[Loi conjointe, lois marginales]
	La \textbf{loi conjointe} de $X$ et $Y$ est la loi de $(X,Y)$. Réciproquement, les \textbf{lois marginales} de $(X,Y)$ sont les lois de $X$ et $Y$.
\end{dfn}
\begin{exa}
	Lorsque $x$ parcourt $X(\Omega)$ et $y$ parcourt $Y(\Omega)$, les événements $\left\{X=x,Y=y\right\}$ forment un SCE.
\end{exa}
\begin{rmq} En général, certains des événements $\left\{X=x,Y=y\right\}$ sont de probabilité nulle. En effet, si l'on note $Z=(X,Y)$, on a $Z(\Omega)\subset X(\Omega)\times Y(\Omega)$, mais l'inclusion réciproque n'a aucune raison d'être vraie. \end{rmq}
\begin{prop}
	On a $$\forall x \in E, \ \P(X=x)=\sum_{y\in F}\P(X=x,Y=y)$$ et $$\forall y \in F, \ \P(Y=y)=\sum_{x\in E}\P(X=x,Y=y)$$
\end{prop}
\begin{cor}[La loi conjointe détermine les lois marginales]
	Si $(X,Y)$ et $(X',Y')$ ont les mêmes probabilités élémentaires, alors $X\sim X'$ et $Y\sim Y'$.
\end{cor}
\begin{rmq} La réciproque est fausse. Il faut considérer un lancer de deux pièces équilibrées dont le résultat est recensé par $(X,Y)$, et d'un autre côté un lancer truqué où la deuxième pièce suit toujours le résultat de la première (équilibrée), dont le résultat est recensé par $(X',Y')$. Les lois marginales sont les mêmes, et pourtant les lois conjointes ne sont pas les mêmes. \end{rmq}
\begin{dfn}[Variables aléatoires indépendantes]
	Deux variables aléatoires $X$ et $Y$ définies sur le même espace probabilisé sont dites \textbf{indépendantes} lorsque $$\forall (x,y)\in X(\Omega)\times Y(\Omega), \ \P(X=x,Y=y)=\P(X=x)\P(Y=y)$$ Cela est équivalent au fait de vérifier $$\forall (x,y)\in E\times F, \ \P(X=x,Y=y)=\P(X=x)\P(Y=y)$$ Dans ce cas, on note alors $X\ind Y$
\end{dfn}
\begin{prop}
	Soit $X$ et $Y$ deux variables aléatoires indépendantes. Alors pour tout $A\subset E$ et tout $B\subset F$, les événements $\left\{X\in A\right\}$ et $\left\{Y\in B\right\}$ sont indépendants.
\end{prop}
\begin{dfn}[Loi conditionnelle]
	Pour tout $x\in E$ tel que $\P(X=x)>0$, on définit la \textbf{loi conditionnelle de $Y$ sachant $\left\{X=x\right\}$} par $$\forall B\subset F, \ \P_{\left\{X=x\right\}}(Y\in B)=\P(Y\in B \ | \ X=x)$$ On définirait de même la loi de $X$ sachant $\left\{Y=y\right\}$.
\end{dfn}
\begin{prop}[HP]
	Les trois propositions suivantes sont équivalentes :
	\begin{enumerate}
		\item $X$ et $Y$ sont indépendantes
		\item Pour tout $x\in E$ tel que $\P(X=x)>0$, la loi de $Y$ sachant $\left\{X=x\right\}$ vaut la loi de $Y$
		\item Pour tout $y\in F$ tel que $\P(Y=y)>0$, la loi de $X$ sachant $\left\{Y=y\right\}$ vaut la loi de $X$
	\end{enumerate}
\end{prop}
\begin{proof}
	Il suffit de montrer l'équivalence entre $1.$ et $2.$ car celle entre $1.$ et $3.$ est symétrique. Pour le sens direct, utiliser la définition de le fait que les événements $\left\{X\in A\right\}$ et $\left\{Y\in B\right\}$ sont indépendants puis revenir à la définition de la loi conditionnelle. Pour le sens réciproque, utiliser la définition de la loi conditionnelle. Pour les événements de probabilité nulle, la probabilité d'un événement "et" comportant l'un de ces événements est aussi nulle par croissance de la probabilité.
\end{proof}
\begin{prop}[Transfert d'indépendance]
	Soit $X$ et $Y$ deux variables aléatoires indépendantes à valeurs respectivement dans $E$ et $F$. Alors pour toutes fonctions $f:E\rightarrow E'$ et $g:F\rightarrow F'$, les variables aléatoires $f(X)$ et $g(Y)$ sont indépendantes.
\end{prop}
\begin{dfn}[Généralisation : $n$-uplet de variables aléatoires]
	On suppose que $n\geqslant2$. Si $X_1,\ldots,X_n$ sont des variables aléatoires définies sur un même espace probabilisé à valeurs respectivement dans $E_,\ldots,E_n$, alors le $n$-uplet de variables aléatoires $(X_1,\ldots,X_n)$ est une variable aléatoire à valeurs dans $E_1\times\ldots\times E_n$. \\ On définit de même la loi conjointe des $X_i$ et les lois marginales de $(X_1,\ldots,X_n)$. Les lois marginales continuent à être déterminées par la la loi conjointe. On définit aussi de même la loi de $X_i$ sachant $\left\{X_j=x_j\right\}$.
\end{dfn}
\begin{dfn}[Variables aléatoires mutuellement indépendantes]
	Les variables aléatoires $X_1,\ldots,X_n$ sont dites \textbf{mutuellement indépendantes} (ou plus simplement \textbf{indépendantes}) lorsque $$\forall(x_1,\ldots,x_n)\in X_1(\Omega)\times\ldots\times X_n(\Omega), \ \P(X_1=x_1,\ldots,X_n=x_n)=\P(X_1=x_1)\times\ldots\times\P(X_n=x_n)$$ Cela est équivalent au fait de vérifier$$\forall(x_1,\ldots,x_n)\in E_1\times\ldots\times E_n, \ \P(X_1=x_1,\ldots,X_n=x_n)=\P(X_1=x_1)\times\ldots\times\P(X_n=x_n)$$
\end{dfn}
\begin{prop}
	Soit $X_1,\ldots,X_n$ des variables aléatoires mutuellement indépendantes. Alors pour tous $A_1\subset E_1,\ldots,A_n\subset E_n$, les événements $\left\{X_1\in A_1\right\},\ldots,\left\{X_n\in A_n\right\}$ sont mutuellement indépendants.
\end{prop}
\begin{prop}
	Soit $X_1,\ldots,X_n$ des variables aléatoires mutuellement indépendantes. Alors pour toute permutation $\sigma\in\mathcal{S}_n$, les variables aléatoires $X_{\sigma(1)},\ldots,X_{\sigma(n)}$ sont mutuellement indépendantes. Pour toute sous-famille $(X_{i_1},\ldots,X_{i_r})$, les variables aléatoires restent mutuellement indépendantes. Les variables aléatoires sont deux à deux indépendantes, mais la réciproque est fausse dès que $n\geqslant3$.
\end{prop}
\begin{prop}
	Si $X_1,\ldots,X_n$ sont mutuellement indépendantes, alors $f_1(X_1),\ldots,f_n(X_n)$ sont mutuellement indépendantes.
\end{prop}
\begin{thm}[Lemme des coalitions]
	Soit $X_1,\ldots,X_n$ des variables aléatoires mutuellement indépendantes et $p\in\llbracket1,n-1\rrbracket$ Soit $f:E_1\times\ldots\times E_p \rightarrow F$ et $g:E_{P+1}\times\ldots\times E_n \rightarrow G$. Alors les variables aléatoires $f(X_1,\ldots,X_p)$ et $g(X_{P+1},\ldots,X_n)$ sont indépendantes. Le résultat s'étend de manière naturelle à plus de deux coalitions.
\end{thm}
\subsection{Lois usuelles}
\begin{dfn}
	Soit $(\Omega,\P)$ et $(E,\Q)$ des espaces probabilisés et $X:\Omega\rightarrow E$ une variable aléatoire. De manière logique, on dit que $X$ \textbf{suit} la loi $\Q$ lorsque $\P_X=\Q$ : $$\forall A\in E, \ \P_X(A)=\Q(A)$$ Dans ce cas, on note $X\sim \Q$.
\end{dfn}
\begin{rmq} Soit $X,Y:\Omega\rightarrow E$. Supposons que $X\sim \Q$. Alors $Y\sim \Q$ si, et seulement si, $X\sim Y$. Supposons que $X\sim Y$. Alors $X\sim \Q$ si, et seulement si, $Y\sim \Q$. \end{rmq}
\begin{prop}
	$X\sim \Q$ si, et seulement si, $\forall x\in E, \ \P(X=x)=\Q(\left\{x\right\})$
\end{prop}
\begin{dfn}[Loi uniforme]
	La \textbf{loi uniforme} est l'unique probabilité définie sur $E=\llbracket1,n\rrbracket$ caractérisée par $$\forall i\in\llbracket1,n\rrbracket, \ \Q(\left\{i\right\})=\frac1n$$ Ainsi, on aura $X\sim\mathcal{U}(\llbracket1,n\rrbracket)$ lorsque $X:\Omega\rightarrow\llbracket1,n\rrbracket$ et $$\forall i\in\llbracket1,n\rrbracket, \ \P(X=i)=\frac1n$$
\end{dfn}
\begin{dfn}[Loi de Bernoulli de paramètre $p$]
	Soit $p\in[0,1]$. La \textbf{loi de Bernoulli de paramètre $p$} est l'unique probabilité définie sur $E=\left\{0,1\right\}$ caractérisée par $$\Q(\left\{1\right\})=p \ \ \text{et} \ \ \Q(\left\{0\right\})=1-p$$ Ainsi, on aura $X\sim\mathcal{B}(p)$ lorsque $X:\Omega\rightarrow\left\{0,1\right\}$, $\P(X=1)=p$ et $\P(X=0)=1-p$.
\end{dfn}
\begin{exa}
	Soit $(\Omega,\P)$ un espace probabilisé et $A$ un événement. Alors, l'indicatrice $\mathds{1}_A$ suit la loi $\mathcal{B}(\P(A))$.
\end{exa}
\begin{dfn}[Loi binomiale de paramètres $n$ et $p$]
	Soit $n\in\N\st$ et $p\in[0,1]$. La \textbf{loi binomiale de paramètres $n$ et $p$} est l'unique probabilité sur $E=\llbracket0,n\rrbracket$ caractérisée par $$\forall k \in\llbracket0,n\rrbracket, \ \Q(\left\{k\right\})=\binom{n}{k}p^k(1-p)^{n-k}$$ Ainsi, on aura $X\sim \mathcal{B}(n,p)$ lorsque $X:\Omega\rightarrow\llbracket0,n\rrbracket$ et $$\forall k \in\llbracket0,n\rrbracket,\ \P(X=k)=\binom{n}{k}p^k(1-p)^{n-k}$$
\end{dfn}
\begin{thm}[Somme de $n$ Bernoulli indépendantes de même paramètre]
	Soit $X_1,\ldots,X_n$ des variables aléatoires indépendantes de loi $\mathcal{B}(p)$. Alors $$X_1+\ldots+X_n \sim \mathcal{B}(n,p)$$
\end{thm}
\subsection{Simulation de variables aléatoires}
\begin{thm}[Simulation de variables aléatoires]
	Soit $(E_i,\Q_i)_{1\leqslant i\leqslant n}$ des espaces probabilisés. Alors il existe un espace probabilisé $(\Omega,\P)$ et des variables aléatoires $X_i:\Omega\rightarrow E_i$ tels que 
	\begin{enumerate}
		\item les $X_i$ soient mutuellement indépendantes
		\item $\forall i\in\llbracket1,n\rrbracket,\ X_i\sim\Q_i$
	\end{enumerate}
\end{thm}
\begin{dfn}[Variables aléatoires indépendantes et identiquement distribuées]
	Lorsqu'on a $n$ copies identiques d'une même variable aléatoire qui sont mutuellement indépendantes, on dit que ces variables aléatoires sont \textbf{indépendantes et identiquement distribuées}. L'abréviation \textbf{iid} est standard. Le théorème précédent nous permet d'avoir autant de variables iid que l'on souhaite qui suivent les lois de notre choix.
\end{dfn}
\section{Espérance et variance}
Dans cette partie, sauf mention expresse du contraire, on considère des variables aléatoires réelles.
\subsection{Espérance}
\begin{dfn}[Espérance]
	Soit $X$ une variable aléatoire réelle. L'\textbf{espérance de $X$} est définie par $$\E[X]=\sum_{x\in X(\Omega)}\P(X=x)\times x$$ 
\end{dfn}
\begin{rmq} En vertu du formalisme des sommes presque nulles, pour tout $B$ tel que $X(\Omega)\subset B\subset\R$, on a $$\E[X]=\sum_{x\in B}\P(X=x)\times x$$ \end{rmq}
\begin{dfn}[Variable aléatoire centrée]
	On dit que la variable aléatoire $X$ est \textbf{centrée} lorsque $\E[X]=0$.
\end{dfn}
\begin{prop}[Espérance d'une indicatrice]
	Pour tout événement $A$, on a $$\E[\mathds{1}_A]=\P(A)$$ Cette formule peut s'avérer très utile lorsqu'on on interprète une variable aléatoire comme une fonction d'indicatrices, car l'espérance est linéaire (et multiplicative pour les variables aléatoires indépendantes).
\end{prop}
\begin{prop}
	On a $$\E[X]=\sum_{\omega\in\Omega}\P(\left\{\omega\right\})\times X(\Omega)$$
\end{prop}
\begin{prop}[Formule de transfert]
	Soit $X:\Omega\rightarrow E$ une variable aléatoire quelconque et $f:E\rightarrow\R$. On a $$\E[f(X)]=\sum_{x\in X(\Omega)}\P(X=x)f(x)$$
\end{prop}
\begin{rmq} De même, si $X(\Omega)\subset B\subset\R$, on a $$\E[f(X)]=\sum_{x\in B}\P(X=x)f(x)$$ \end{rmq}
\begin{cor}
	Soit $X:\Omega\rightarrow E$ et $Y:\Omega\rightarrow E$ deux variables aléatoires quelconques. Si $X\sim Y$, alors pour toute fonction $f:E\rightarrow\R$, on a $$\E[f(X)]=\E[f(Y)]$$
\end{cor}
\begin{rmq} Si $Y$ est définie sur un autre espace probabilisé $(\Omega',\P')$, le résultat reste vrai en écrivant $\E[f(X)]=\E'[f(Y)]$ \end{rmq}
\begin{cor}
	Deux variables aléatoires réelles de même loi ont même espérance.
\end{cor}
\begin{thm}[Propriétés intégrales de l'espérance]
	L'espérance vérifie les propriétés suivantes
	\begin{enumerate}
		\item Linéarité : $\E\left[\lambda X+\mu Y\right]=\lambda\E[X]+\mu\E[Y]$
		\item Inégalité triangulaire $\left\lvert\E[X]\right\rvert\leqslant\E\left[\lvert X\rvert\right]$
		\item Positivité : si $X\geqslant 0$, alors $\E[X]\geqslant 0$
		\item Croissance : si $X\leqslant Y$, alors $\E[X]\leqslant\E[Y]$
	\end{enumerate}
\end{thm}
\begin{rmq} Les deux derniers points restent vrais si l'inégalité n'a lieu que presque sûrement. \end{rmq}
\begin{prop}[Quelques espérances à connaître]
	On a les résultat suivants :
	\begin{enumerate}
		\item Si $X=x$, alors $\E[X]=x$
		\item Si $X\sim \mathcal{U}(\llbracket1,n\rrbracket)$, alors $\E[X]=\displaystyle\frac{n+1}{2}$
		\item Si $X\sim\mathcal{B}(p)$, alors $\E[X]=p$
		\item Si $X\sim\mathcal{B}(n,p)$ alors $\E[X]=np$
	\end{enumerate}
\end{prop}
\begin{rmq} Par conséquent, la variable aléatoire $X-\E[X]$ est toujours centrée. \end{rmq}
\begin{prop}[Formule du crible de Poincaré pour les probabilités, HP]
	Soit $(\Omega,\P)$ et $A_1,\ldots,A_n$ des événements. On a $$\P\left(\bigcup_{1\leqslant i\leqslant n}\right)=\sum_{1\leqslant k \leqslant n}(-1)^{k-1}\sum_{1\leqslant i_1<\ldots<i_k\leqslant n}\P\left(\bigcap_{1\leqslant j\leqslant k}A_{i_j}\right)$$
\end{prop}
\begin{proof}
	On passe par le complémentaire. Développer $(1-\mathds{1}_{A_1})\ldots(1-\mathds{1}_{A_n})$ par distributivité généralisée puis passer aux espérances.
\end{proof}
\begin{prop}[Formules sur des égalités ayant lieu presque sûrement, HP]
	Si $X=x$ p.s., alors $\E[X]=x$. Si $X=Y$ P.S., alors $\E[X]=\E[Y]$.
\end{prop}
\begin{proof}
	On calcule $$\E[X]=\P(X=x)\times x + \sum_{t\in X(\Omega)\setminus\left\{x\right\}}\P(X=t)\times t$$ Or, pour tout $t\in X(\Omega)\setminus\left\{x\right\}$, $\left\{X=t\right\}\subset \left\{X\ne x\right\}$ donc de probabilité nulle par croissance de la probabilité. D'où on tire $\E[X]=x$. La deuxième s'en déduit : $X-Y=0$ p.s., donc on en déduit le résultat par linéarité et par le premier cas.
\end{proof}
\begin{prop}[HP]
	Supposons que $X\geqslant 0$. On a $\E[X]=0$ si, et seulement si, $X=0$ p.s.
\end{prop}
\begin{proof}
	Dans le sens direct, on a pour tout $x\in X(\Omega)$, $\P(X=x)x$. Donc, pour tout $x\in\R_+\st$, $\P(X=x)=0$. Or $\Omega=X\inv(X(\Omega))=\displaystyle \bigsqcup_{x \in X(\Omega)}\left\{X=x\right\}$. Nécessairement, $0\in X(\Omega)$. Puis, par union disjointe, $\P(\Omega)=\P(X=0)=1$. Donc $X=0$ p.s. Réciproquement, on utilise la proposition précédente.
\end{proof}
\begin{thm}(Inégalité de Cauchy-Schwarz pour l'espérance)
	Soit $X$ et $Y$ des variables aléatoires réelles. On a $$\left\lvert\E[XY]\right\rvert \leqslant \sqrt{\E[X^2]}\sqrt{\E[Y^2]}$$ car la fonction $(X,Y)\mapsto\E[XY]$ est bilinéaire symétrique et positive, et ce sont les seuls hypothèses nécessaires pour l'inégalité de Cauchy-Schwarz. En revanche, le cas d'égalité devient faux. En effet, si $\E[X^2]=0$, on a seulement $X=0$ p.s.
\end{thm}
\begin{rmq}[HP] La solution consiste à quotienter le $\R$-ev $\R^\Omega$ par la relation d'équivalence "$X=Y$ p.s.". On vérifie que l'espérance passe au quotient et fait de l'espace quotient un espace préhilbertien réel. \end{rmq}
\begin{thm}[Espérance du produit de variables aléatoires indépendantes]
	Si $X$ et $Y$ sont indépendantes, alors $$\E[XY]=\E[X]\E[Y]$$
\end{thm}
\begin{rmq} La réciproque est fausse en général ! \end{rmq}
\begin{rmq} Le théorème se généralise au cas où $n$ variables aléatoires sont mutuellement indépendantes grâce au lemme des coalitions. \end{rmq}
\begin{thm}[Inégalité de Markov]
	Soit $X$ à valeurs dans $\R_+$. Alors $$\forall x>0, \ \P(X\geqslant x)\leqslant\frac{\E[X]}{x}$$
\end{thm}
\subsection{Autres moments}
\begin{dfn}[Moment d'ordre $k$]
	Si $X$ est une variable aléatoire réelle, son \textbf{moment d'ordre $k$} vaut $\E[X^k]$ pour $k\in\N$.
\end{dfn}
\begin{dfn}
	La \textbf{variance} de $X$ est définie par $$\V[X]=\E\left((X-\E[X])^2\right)$$ C'est un nombre positif ou nul, qui indique la "dispersion" de $X$.
\end{dfn}
\begin{prop}[HP]
	Soit $X$ réelle d'espérance $\mu$. On a $\V[X]=0$ si, et seulement si, $X=\mu$ p.s.
\end{prop}
\begin{proof}
	$(X-\E[X])^2$ est positive. En vertu d'une propriété précédente, $\V[X]=0$ ssi $(X-\E[X])^2=0$ p.s. ssi $(X-\E[X])=0$ p.s. ssi $X=\mu$ p.s.
\end{proof}
\begin{prop}[Formule de König-Huygens]
	On a aussi $$\V[X]=\E[X^2]-\E[X]^2$$
\end{prop}
\begin{prop}
	Si deux variables aléatoires ont même loi, alors elles ont même moment et donc elles ont même variance.
\end{prop}
\begin{prop}
	La variance vérifie $$\forall (a,b)\in\R^2,\ \V[aX+b]=a\V[X]$$
\end{prop}
\begin{dfn}[Variable aléatoire réduite]
	Si $\V[X]=1$, la variable aléatoire $X$ est dite \textbf{réduite}.
\end{dfn}
\begin{dfn}[Écart-type]
	L'\textbf{écart-type} de $X$ est défini par $$\sigma(X)=\sqrt{\V[X]}$$ Il vérifie la relation $$\sigma(aX+b)=\lvert a\rvert\sigma(X)$$
\end{dfn}
\begin{prop}
	Si $X$ est une variable aléatoire d'écart-type non nul, alors la variable aléatoire $\displaystyle \frac{X-\E[X]}{\sigma(X)}$ est centrée réduite.
\end{prop}
\begin{thm}[Inégalité de Bienaymé-Techebychev]
	Soit $X$ une variable aléatoire réelle. Alors $$\forall a>0,\ \P(\lvert X-\E[X]\rvert\geqslant a)\leqslant\frac{\V[X]}{a^2}$$
\end{thm}
\begin{dfn}[Covariance]
	La \textbf{covariance} de deux variables aléatoires $X$ et $Y$ est définie par $$\cov(X,Y)=\E\left[(X-\E[X])(Y-\E[Y])\right]$$
\end{dfn}
\begin{prop}[Formule de König-Huygens pour la covariance]
	On a aussi $$\cov(X,Y)=\E[XY]-\E[X]\E[Y]$$
\end{prop}
\begin{prop}[Inégalité de Cauchy-Schwarz pour la covariance]
	La covariance est bilinéaire, symétrique et positive, donc l'inégalité de Cauchy-Schwarz s'applique (mais pas la cas d'égalité !) : $$\lvert\cov(X,Y)\rvert \leqslant\sqrt{\V[X]}\sqrt{\V[Y]}$$ 
\end{prop}
\begin{cor}
	Si $X$ et $Y$ sont indépendantes, alors leur covariance est nulle. La réciproque est fausse !
\end{cor}
\begin{dfn}[Variables aléatoires décorrélées]
	Si $\cov(X,Y)=0$, on dit que $X$ et $y$ sont décorrélées.
\end{dfn}
\begin{thm}[Variance d'une somme]
	Soit $X_1,\ldots,X_n$ sont des variables aléatoires. Alors $$\V[X_1+\ldots+X_n]=\V[X_1]+\ldots+\V[X_n]+\sum_{i\ne j} \cov(X_i,X_j)$$ Par symétrie de la covariance, cela se réécrit $$\V[X_1+\ldots+X_n]=\V[X_1]+\ldots+\V[X_n]+2\sum_{i<j} \cov(X_i,X_j)$$
\end{thm}
\begin{cor}
	Si $X_1,\ldots,X_n$ sont deux à deux indépendantes, on a $$\V[X_1+\ldots+X_n]=\V[X_1]+\ldots+\V[X_n]$$
\end{cor}
\begin{rmq} A fortiori, le résultat reste vrai lorsque les variables aléatoires sont mutuellement indépendantes. \end{rmq}
\begin{cor}[Quelques espérances à connaître]
	On a les résultats suivants :
	\begin{enumerate}
		\item Si $X\sim \mathcal{B}(p)$, alors $\V[X]=p(1-p)$
		\item Si $X\sim \mathcal{B}(n,p)$, alors $\V[X]=np(1-p)$
	\end{enumerate}
\end{cor}
\begin{thm}[Loi faible des grands nombres, HP]
	Soit $(X_n)_{n\in\N\st}$ une suite de variables aléatoires réelles toutes définies sur le même univers $\Omega$, indépendantes et identiquement distribuées. Notons $\mu$ leur espérance (commune à toutes les variables aléatoires donc). On pose $$\forall n\in\N\st,\ Y_n=\frac{X_1+\ldots+X_n}{n}$$ Alors, "$Y_n$ a fortement tendance à se concentrer autour de $\mu$". Plus précisément : $$\forall \varepsilon >0,\ \P(\lvert Y_n-\mu\rvert >\varepsilon)\xrightarrow[n\to+\infty]{}0$$ Ce théorème justifie l'approche fréquentiste des probabilités.
\end{thm}
\begin{proof}
	L'espérance de $Y_n$ vaut toujours $\mu$. Utiliser ensuite l'inégalité de Bienaymé-Tchebychev et les propriétés de la variance pour conclure.
\end{proof}

\begin{thebibliography}{}
	\bibitem{FM} F. Morlot, Lycée Privé Sainte-Geneviève, Cours dispensés en MPSI1 durant l'année 2023-2024
	\bibitem{CL} C. Lafitte, Lycée Privé Sainte-Geneviève, \textit{Compléments sur les anneaux commutatifs}
\end{thebibliography}
\end{document}